\subsection{End-to-end results and their interpretation}
\label{sec:benchmark-results}

The following tables present the approved results available while the benchmark
campaign is still running. They compare the BioseqDB mapping pipeline with
minimap2 on one local Apple Silicon machine and four AWS machines. Each timing
covers a batch of 1,024 queries. A missing measurement remains visibly missing,
so the tables describe completed observations rather than an estimated full
matrix. The machine inventory precedes the result tables.

\begin{itemize}
  \item \textbf{Index.} A persisted index is constructed before measurement.
  Each repeat starts a fresh process and measures reading the index and mapping
  the queries. Construction of the persisted index is excluded from the current
  mapping interval. The index is available on disk, but is not retained in the
  previous process's memory. Operating-system page caches are not purged, so
  this is process-cold execution rather than a guaranteed cold disk read.
  \item \textbf{No index.} No pre-existing user index is available. The current
  measured interval includes transient target-index construction and mapping.
  The transient index can be reused within the batch, including by the remaining
  1,023 queries after its initial construction. Fixture generation and process
  restart setup lie outside the current measured interval.
  \item \textbf{Timing scopes.} Historical outer-E2E timings are treated as
  mismeasurements and excluded from the tables, coverage counts and comparisons.
  Standard observations use the current framework mapping-iteration interval.
  FAST fallbacks retain their distinct single-sample timing scope and carry
  an \textsuperscript{s} marker. They are never pooled with standard repeats.
  The final external-aligner table identifies its separate timing boundary explicitly.
  In particular, ``mapping interval'' still includes transient construction in
  no-index mode.
  \item \textbf{Two measurements per cell.} The first duration, A, is the
  nearest-rank P50 of the complete measured interval. The second, B, is the
  nearest-rank P50 after subtracting the physical union of instrumented CIGAR
  and MAPQ intervals separately in each repeat. With five repeats the P50 is the
  third sorted observation. For FAST entries, P50 is the single observed value.
  If $U_i$ is that interval union, the definitions are
  \[
    A_{50}=\operatorname{P50}(A_i),\qquad
    B_{50}=\operatorname{P50}(A_i-|U_i|).
  \]
  Overlapping intervals are subtracted only once. The published medians are
  used directly, without subtracting independently aggregated medians. On both
  BioseqDB KSW2 and Block Aligner paths, the CIGAR interval contains fused base
  alignment, trace storage, traceback and result construction. B therefore does
  not universally represent seeding and chaining alone. Likewise,
  $A_{50}-B_{50}$ is not a measured median CIGAR/MAPQ cost.
\end{itemize}

\paragraph{Input sequences and repetitions.}
The campaign configuration defines eight deterministic target fixtures, covering
both target sizes, both alphabets and both minimizer parameter pairs. Each
fixture contains one reference sequence and 1,024 derived queries. The same
fixture is reused across the relevant algorithm configurations and repeats.
The configured parameters are as follows.
\begin{itemize}
  \item \textbf{Target sizes and minimizers.} References contain exactly
  $2\times10^9$ or $4\times10^9$ base pairs, denoted 2 and 4 decimal Gbp.
  The parameter pairs are $(k,w)=(15,10)$ and $(20,15)$, where $k$ is the
  k-mer length and $w$ is the number of k-mers in a minimizer window.
  \item \textbf{Alphabets.} DNA4 uses equiprobable canonical bases. The DNA5
  generator assigns probability 1\% to N and 24.75\% to each canonical base.
  \item \textbf{Query lengths and strand.} The sampled pre-edit length has
  Gaussian mean 20,000 bp and standard deviation 500 bp, with rounding and
  clamping in the generator. Insertions and deletions can change the final
  length. A query is reverse complemented with probability 0.5.
  \item \textbf{Divergence parameters.} Substitution probability is 0.001.
  Insertion-event and deletion-event probabilities are each 0.0005. Their
  nominal sum is 0.002, or 0.2\%, rather than a measured edit-distance rate.
  Indel lengths have a configured mean parameter of 1 bp and the generator's
  default standard-deviation parameter of 1 bp, followed by rounding and
  clamping. They are not fixed single-base edits.
  \item \textbf{Occurrences.} Each query originates from one known sampled
  reference interval. The fixture does not specify an occurrence-count
  distribution or an occurrence standard deviation. This provenance does not
  imply that an aligner returns exactly one hit. The original five summaries certify
  exact baseline match-count equality but intentionally omit the actual counts.
  \item \textbf{Repetitions and variability.} Standard observations have five
  accepted repeats. FAST observations have one measured batch and cannot
  establish repeat-to-repeat variability. A repeat measures the complete query
  batch, not a single query. The original five summaries contain medians and repeat counts,
  but no raw timing samples, so timing standard deviations and confidence
  intervals cannot be recovered from these files.
  \item \textbf{FAST procedure.} Compatible persisted indexes are reused across
  index-mode configurations. PostgreSQL restarts before measurement, and the
  operating-system filesystem cache is not cleared. Index preparation and restart
  are outside the timer. No-index measurements still include transient index
  construction. FAST remote runs omit the additional plan and alternative-setting
  controls, while retaining measured-statement correctness and timing checks.
\end{itemize}

\paragraph{Execution rows.}
BioseqDB uses task-parallel execution with prefiltering disabled and both CIGAR
and MAPQ projection enabled. MPF and MPS denote minimizers per fragment and per
sequence. Each row identifies the requested minimizer, chaining algorithm and
base aligner in addition to this index representation. Scalar-DP and SIMD-DP
refer to the scalar and literal SIMD minimap2 chaining routes. RMQ denotes the
separate range-minimum-query algorithm. SIMD-ref is the reference minimizer,
Two-stack is the two-stack ntHash minimizer, and Bounded is bounded ntHash SIMD.
Block denotes Block Aligner. Vulkan rows can include CPU index handling,
backtracking and base alignment, and do not imply an entirely GPU pipeline.
If none of the Vulkan configurations has a measurement for the workloads and
machines in a table, those configurations share one ``Vulkan (all modes)'' row.
Recorded GPU unavailability is marked explicitly. A missing timing alone is
not taken as evidence of unavailable hardware.

\paragraph{Reading comparisons and highlights.}
Each cell places the A duration above its comparison, followed by the B duration
and its own comparison. Durations are rounded to tenths of a second for display.
All comparisons and minima use the original integer nanoseconds.
\begin{itemize}
  \item A BioseqDB value is compared with minimap2 for the same machine class,
  fixture and index mode. Marked FAST fallbacks may be compared with standard observations,
  which have a different repeat count and execution procedure.
  The minimap2 row compares with the fastest
  measured eligible BioseqDB effective path in that same column, independently
  for A and B. Thus its two reference algorithms can differ. Baselines are
  repeated on continuation pages solely to assist reading.
  \item A green upward triangle means less elapsed time. A red downward
  triangle means more elapsed time. A gray rightward triangle indicates a tie
  when the ratio of the larger to the smaller duration is at most 1.05.
  This is a display convention, not a statistical significance test.
  \item For current duration $t$ and reference $r$, the percentage is
  $100|t-r|/r$, describing relative elapsed-time reduction or increase.
  Differences above 30\% use ratio notation instead: the larger duration divided
  by the smaller, rounded to one decimal place. A green $2.5\times$ therefore
  means the reference takes 2.5 times as long, while a red $2.5\times$ means
  the displayed execution takes 2.5 times as long. Differences through 30\%
  retain percentage notation. Direction is always conveyed by the triangle.
  \item A blue background marks one smallest A and one smallest B down each
  machine/workload column, including minimap2 and all measured BioseqDB rows
  across continuation pages. Exact ties select the first displayed row.
  Repeated minimap2 baselines receive blue marks only on the first page, so
  each column has two winner marks across the complete table when measurements
  are available. Each AWS workload has four machine subcolumns separated by
  vertical rules.
  \item Light gray identifies every minimap2 baseline row. Pale yellow
  identifies the preferred BioseqDB configuration: MPF / Bounded / SIMD-DP /
  Block. This is the author's selected configuration for comparison with the
  external aligners. The yellow tint does not assert that it wins every
  workload. Blue timing highlights identify the measured column winners and
  remain visible within either row tint.
  \item A dot means there is no approved observation in the snapshot.
  A horizontal dash with \textsuperscript{g} identifies an explicitly recorded
  unavailable GPU configuration, which was not executed. Neither marker is a
  zero duration. The comparison reads n/a when no measured counterpart exists.
  Only execution rows with an approved observation or an explicit unavailable
  cell in a panel are printed. Entirely unmeasured row families remain omitted.
\end{itemize}

\paragraph{Limits of the comparison.}
These observations were accumulated across builds. Some retain an older sealed
minimap2 correctness control, which is not a new contemporaneous performance
run. A correctness-only index control is never substituted for a no-index
performance baseline. Comments in the source authorities also disclose missing process-priority
attestation, recovery after post-measurement wrapper failures and historical
callback-lifetime caveats. The \textsuperscript{p} marker identifies the local
observation retained by an explicit preservation decision without a run receipt.
Approval of its measurement does not certify the safety of its historical build.
No uniform elevated-priority or noise-free execution claim is made.

The remote C7G observations include severe paging, and other observations also
report resource pressure. Cross-machine differences therefore combine hardware,
resource availability and execution conditions. They are not an isolated test
of instruction-set performance. The table records physical RAM and provisioned
storage capacity. Selected row comments contain sampled available memory, but
there is no uniform per-repeat series of free-memory and free-disk measurements.
Missing C7G BioseqDB timings cannot establish either a win or a loss.

\paragraph{Reproducible table generation.}
The standalone Python generator reads only the normative
\texttt{approved\_results} rows and the remote
\texttt{hardware\_unavailable\_cells}, with their embedded machine metadata.
It validates identities and timing scopes, rejects duplicates within a source, and
records the SHA-256 of each input in the generated LaTeX source. Coverage totals
are recomputed from approved rows rather than copied from historical progress
summaries. Regeneration is manual and is not part of CMake. The fixture description
was checked against the campaign JSON and its sequence generator, together with
DOC0432 and DOC0433. The full row-specific qualifications remain in the
source YAML authorities. For each remote cell, the generator selects the first
available approved observation from the default summary, then TEST3, TEST4,
TEST7, TEST2 and finally FAST for the same machine class and workload. Local cells use
the default local summary followed by FAST LOCAL. Both A and B come from the
same selected observation. The generator checks that every unique eligible
tuple in these sources is rendered once logically, excluding the historical
mismeasurements and accounting for repeated baseline rows on continuation pages.
Missing fallback files are skipped. This selection also supplies the minimap2
and BioseqDB rows in the final comparison table.
The final table additionally reads the four external-aligner
authorities and explains their separate timing and validation rules.

% Generated by main/papers/scripts/generate_benchmark_results.py. Do not edit.
% Source byte hashes (regenerate manually after approved results change).
% ALL_RESULTS_TUPLES.yaml sha256=2df7cf1d9a066beb1de6cd2d061114031a8c2e7ef20b2e9044a424ca2642ec78
% ALL_RESULTS_TUPLES_FAST_LOCAL.yaml sha256=d42f70ff0c1970bd7e0fde70e345162e910a8ff9727a4599d4914a0c10258ca1
% ALL_RESULTS_TUPLES_REMOTE_M7A.yaml sha256=2a3bce5db04071b3b9e71227a0c69e7516d14d47c98fab6a0fcb7ca17f5eb8f7
% ALL_RESULTS_TUPLES_REMOTE_TEST3_M7A.yaml sha256=932bb904eebf35ce676fbb2f3015e49691a1b77c9ca0ceedb221f637bf184e75
% ALL_RESULTS_TUPLES_REMOTE_TEST4_M7A.yaml sha256=119bcf2c1e990513fd173b96df6c356a2992369e556fa5b2d64dc58f61d43a70
% ALL_RESULTS_TUPLES_REMOTE_TEST2_M7A.yaml sha256=c42a0b316b7660223e698c5b8a17d4b3c46dc22464b7c507d5629bdc3af2c366
% ALL_RESULTS_TUPLES_FAST_REMOTE_M7A.yaml sha256=9c0534090bbf0b315802fae26a009ddc7374a388e7df368cf5ef4fcdfec3152e
% ALL_RESULTS_TUPLES_REMOTE_M7I.yaml sha256=85f242b097fb4578bc9c236329c02014209c3f430507c49127b211a0170253c5
% ALL_RESULTS_TUPLES_REMOTE_TEST3_M7I.yaml sha256=6d43b87a7725e4da43774a66bb1b8771c8b993f1fc31d5b9473e42a59d2ab7d6
% ALL_RESULTS_TUPLES_REMOTE_TEST4_M7I.yaml sha256=77f31b33d8edbbf6d753be537668180b078dfa96b362d9f69673bc7cb6088d68
% ALL_RESULTS_TUPLES_REMOTE_TEST2_M7I.yaml sha256=c3da77d1e5e66058f3e939e6c741e8816b4f88b05b076eced6b5caf6a6c27978
% ALL_RESULTS_TUPLES_FAST_REMOTE_M7I.yaml sha256=9f18ad70caae5b4a13f6ff220b906815c8cb09309d34a43bac5752e80bd65a97
% ALL_RESULTS_TUPLES_REMOTE_M4.yaml sha256=0a0a5fb3c2d6e287ca50837f9eb6ac53f75ed4e97845fcab41f5fcbe1f262372
% ALL_RESULTS_TUPLES_REMOTE_TEST3_M4.yaml sha256=70f05642f6295b0e1f2631e1f5743a428169fd7cbf974e3885f79f0a9f6e023d
% ALL_RESULTS_TUPLES_REMOTE_TEST4_M4.yaml sha256=2e72745580632c795d027b25c9c8c550114c73faa8225fd5c45c7f435a57b71b
% ALL_RESULTS_TUPLES_REMOTE_TEST7_M4.yaml sha256=653edc0eeba2358189805850e4ab1a03248df4a1d3321dac2c5f063ab0f52c4b
% ALL_RESULTS_TUPLES_FAST_REMOTE_M4.yaml sha256=fc046174e1b1bc6472eba40c8ad76fcf9bb825c30659aa6fefa72b1f7d5f41db
% ALL_RESULTS_TUPLES_REMOTE_C7G.yaml sha256=0a86873dc0117c2024033478ad75eee15e50dfdfc2497965dc9f954421e4d8a9
% ALL_RESULTS_TUPLES_REMOTE_TEST3_C7G.yaml sha256=56867d988fd0ba74f76046af76971c599f84b4eb60a9b4024776c056c3166b76
% ALL_RESULTS_TUPLES_REMOTE_TEST4_C7G.yaml sha256=6b231dd87d6623a252214375ccfa4fc0e4f0d6881d23e986f62b4422b5e14da3
% ALL_RESULTS_TUPLES_REMOTE_TEST7_C7G.yaml sha256=d35ecafbe96dfa8fd3dc44bb63615cbc9db9667e8632dc2ab17ccf9e3607bc40
% ALL_RESULTS_TUPLES_FAST_REMOTE_C7G.yaml sha256=4754c86889581297b7239935d3f657fcb37fb4ea24acebfe377f731640735a54
% ALL_RESULTS_TUPLES.yaml: available=103, selected=103, shadowed=0
% ALL_RESULTS_TUPLES_FAST_LOCAL.yaml: available=8, selected=8, shadowed=0
% ALL_RESULTS_TUPLES_REMOTE_M7A.yaml: available=128, selected=128, shadowed=0
% ALL_RESULTS_TUPLES_REMOTE_TEST3_M7A.yaml: available=63, selected=59, shadowed=4
% ALL_RESULTS_TUPLES_REMOTE_TEST4_M7A.yaml: available=51, selected=45, shadowed=6
% ALL_RESULTS_TUPLES_REMOTE_TEST2_M7A.yaml: available=29, selected=29, shadowed=0
% ALL_RESULTS_TUPLES_FAST_REMOTE_M7A.yaml: available=0, selected=0, shadowed=0
% ALL_RESULTS_TUPLES_REMOTE_M7I.yaml: available=129, selected=129, shadowed=0
% ALL_RESULTS_TUPLES_REMOTE_TEST3_M7I.yaml: available=57, selected=52, shadowed=5
% ALL_RESULTS_TUPLES_REMOTE_TEST4_M7I.yaml: available=48, selected=42, shadowed=6
% ALL_RESULTS_TUPLES_REMOTE_TEST2_M7I.yaml: available=33, selected=33, shadowed=0
% ALL_RESULTS_TUPLES_FAST_REMOTE_M7I.yaml: available=0, selected=0, shadowed=0
% ALL_RESULTS_TUPLES_REMOTE_M4.yaml: available=75, selected=75, shadowed=0
% ALL_RESULTS_TUPLES_REMOTE_TEST3_M4.yaml: available=4, selected=0, shadowed=4
% ALL_RESULTS_TUPLES_REMOTE_TEST4_M4.yaml: available=5, selected=3, shadowed=2
% ALL_RESULTS_TUPLES_REMOTE_TEST7_M4.yaml: available=12, selected=12, shadowed=0
% ALL_RESULTS_TUPLES_FAST_REMOTE_M4.yaml: available=0, selected=0, shadowed=0
% ALL_RESULTS_TUPLES_REMOTE_C7G.yaml: available=48, selected=48, shadowed=0
% ALL_RESULTS_TUPLES_REMOTE_TEST3_C7G.yaml: available=23, selected=22, shadowed=1
% ALL_RESULTS_TUPLES_REMOTE_TEST4_C7G.yaml: available=10, selected=9, shadowed=1
% ALL_RESULTS_TUPLES_REMOTE_TEST7_C7G.yaml: available=5, selected=5, shadowed=0
% ALL_RESULTS_TUPLES_FAST_REMOTE_C7G.yaml: available=0, selected=0, shadowed=0
% ALL_RESULTS_TUPLES_E2E_EXTERNAL_REMOTE_M7A.yaml sha256=f62d914bb70f8a27814a896ed2ce3eb40d559229fae3c52d9bd2751b65716be3
% ALL_RESULTS_TUPLES_E2E_EXTERNAL_REMOTE_M7I.yaml sha256=3c221e48fbd3aab59e37ab2d394e8ce924689301ff60d93a1e6cb31d0eaa389d
% ALL_RESULTS_TUPLES_E2E_EXTERNAL_REMOTE_M4.yaml sha256=f32214788e4a798aa60b383c56ef5b2206d457ac6edcfea5b2e90d412da5fc2d
% ALL_RESULTS_TUPLES_E2E_EXTERNAL_REMOTE_C7G.yaml sha256=4ea69b323d95947e09273bf6a5c0cf733144342ab7e46ffa9170a9f28c54076c
\definecolor{BenchBest}{HTML}{DDECF8}
\definecolor{BenchBaseline}{HTML}{F4F4F4}
\definecolor{BenchPreferred}{HTML}{FFFBEA}
\newgeometry{left=10mm,right=10mm,top=12mm,bottom=12mm}
\begin{landscape}
\subsection{Recorded machines and completed coverage}
\begin{table}[H]\centering\small
\caption{Machine inventory captured with the benchmark summaries. Capacities are not free resources at execution time.}
\begin{tabularx}{\linewidth}{@{}p{24mm}*{5}{>{\raggedright\arraybackslash}X}@{}}\toprule
Property & Local & M7A & M7I & M4 & C7G\\\midrule
Machine & MacBook Pro & m7a.xlarge & m7i.xlarge & m4.xlarge & c7g.xlarge\\[3pt]
CPU & Apple M1 Pro & AMD EPYC 9R14 & Xeon Platinum 8488C & Xeon E5-2686 v4 @ 2.30GHz & Neoverse-V1\\[3pt]
Architecture & arm64 & x86\_64 & x86\_64 & x86\_64 & aarch64\\[3pt]
OS version / kernel & macOS 15.6 / Darwin 24.6.0 & Linux / 7.0.0-1012-aws & Linux / 7.0.0-1012-aws & Linux / 7.0.0-1012-aws & Linux / 7.0.0-1012-aws\\[3pt]
OS build / distribution & 24G5054d & Distribution not recorded & Distribution not recorded & Distribution not recorded & Distribution not recorded\\[3pt]
Logical / physical cores & 8 / 8 (6P + 2E) & 4 / 4 & 4 / 2 & 4 / 2 & 4 / 4\\[3pt]
Physical RAM & 16 GiB unified & 16 GiB & 16 GiB & 16 GiB & 8 GiB\\[3pt]
GPU & Apple M1 Pro, 14 cores & None available & None available & None available & None available\\[3pt]
Storage & APPLE SSD AP0512R & EBS gp3 & EBS gp3 & EBS gp3 & EBS gp3\\[3pt]
Disk capacity & 500.28 GB & 500 GiB & 500 GiB & 500 GiB & 500 GiB\\[3pt]
Provisioned disk I/O & Not recorded & 6,000 IOPS / 500 MiB/s & 6,000 IOPS / 500 MiB/s & 6,000 IOPS / 500 MiB/s & 6,000 IOPS / 500 MiB/s\\[3pt]
Free disk / free RAM & No uniform per-repeat record & No uniform per-repeat record & No uniform per-repeat record & No uniform per-repeat record & No uniform per-repeat record\\[3pt]
\bottomrule\end{tabularx}\end{table}
\noindent The AWS summaries record the Linux kernel but do not seal an OS distribution version. The local legacy capture does not seal the MoltenVK implementation identity. Unrecorded values are not inferred from the current machine.
\par Remote cells use the first available approved result in this order: the default summary, TEST3, TEST4, TEST7, TEST2, then FAST. Local cells use the default summary, then FAST LOCAL. A and B come from the same observation. The inventory describes the default hosts. Fallback measurements may come from other instances of the same machine class, and comparisons may combine campaigns.\par
\begin{table}[H]\centering\small
\caption{Published snapshot coverage after excluding mismeasured historical timings.}
\begin{tabular}{lrrrrr}\toprule Host & Logical tuples & Observations & Repeats & BioseqDB obs. & GPU unavailable\\\midrule
Local & 111 & 111 & 523 & 95 & 0\\
M7A & 261 & 261 & 1305 & 245 & 512\\
M7I & 256 & 256 & 1280 & 240 & 512\\
M4 & 90 & 90 & 450 & 74 & 512\\
C7G & 84 & 84 & 420 & 68 & 512\\
\bottomrule\end{tabular}\end{table}
The published BioseqDB campaign snapshot contains 802 timing observations and 3,978 accepted repeats. FAST entries retain their separate single-sample timing scope and are marked with a superscript s. The 66 historical outer-E2E observations are treated as mismeasurements and excluded from coverage and comparisons. Standard observations have five accepted repeats. FAST observations have one. The planned local matrix has 1,296 logical cells and 6,480 repeats. It remains incomplete. Unavailable-GPU counts refer only to the remote authorities' recorded eligible subsets.
\clearpage
\noindent{\large\bfseries Local machine: mapping interval\hfill 1/5}\par\smallskip
{\fontsize{7.1}{8.1}\selectfont\setlength{\tabcolsep}{1pt}\setlength{\fboxsep}{2pt}\renewcommand{\arraystretch}{1.12}
\noindent\begin{tabular}{@{}p{14mm}p{27mm}|p{13.55mm}|p{13.55mm}|p{13.55mm}|p{13.55mm}|p{13.55mm}|p{13.55mm}|p{13.55mm}|p{13.55mm}|p{13.55mm}|p{13.55mm}|p{13.55mm}|p{13.55mm}|p{13.55mm}|p{13.55mm}|p{13.55mm}|p{13.55mm}|@{}}\hline
Engine & Execution & \multicolumn{8}{c|}{\textbf{Index}} & \multicolumn{8}{c|}{\textbf{No index}}\\
 &  & \multicolumn{4}{c|}{\textbf{4 Gbp}} & \multicolumn{4}{c|}{\textbf{2 Gbp}} & \multicolumn{4}{c|}{\textbf{4 Gbp}} & \multicolumn{4}{c|}{\textbf{2 Gbp}}\\
 &  & \multicolumn{2}{c|}{\textbf{k=15, w=10}} & \multicolumn{2}{c|}{\textbf{k=20, w=15}} & \multicolumn{2}{c|}{\textbf{k=15, w=10}} & \multicolumn{2}{c|}{\textbf{k=20, w=15}} & \multicolumn{2}{c|}{\textbf{k=15, w=10}} & \multicolumn{2}{c|}{\textbf{k=20, w=15}} & \multicolumn{2}{c|}{\textbf{k=15, w=10}} & \multicolumn{2}{c|}{\textbf{k=20, w=15}}\\
 &  & \multicolumn{1}{c|}{\textbf{DNA4}} & \multicolumn{1}{c|}{\textbf{DNA5}} & \multicolumn{1}{c|}{\textbf{DNA4}} & \multicolumn{1}{c|}{\textbf{DNA5}} & \multicolumn{1}{c|}{\textbf{DNA4}} & \multicolumn{1}{c|}{\textbf{DNA5}} & \multicolumn{1}{c|}{\textbf{DNA4}} & \multicolumn{1}{c|}{\textbf{DNA5}} & \multicolumn{1}{c|}{\textbf{DNA4}} & \multicolumn{1}{c|}{\textbf{DNA5}} & \multicolumn{1}{c|}{\textbf{DNA4}} & \multicolumn{1}{c|}{\textbf{DNA5}} & \multicolumn{1}{c|}{\textbf{DNA4}} & \multicolumn{1}{c|}{\textbf{DNA5}} & \multicolumn{1}{c|}{\textbf{DNA4}} & \multicolumn{1}{c|}{\textbf{DNA5}}\\
\hline
\rowcolor{BenchBaseline}[0.5pt][0.5pt]\textbf{minimap2} & \textbf{Baseline} & \hspace*{1pt}\begin{tabular}[t]{@{}c@{}}\rule{0pt}{12pt}39.7s\\[-.2pt]{\fontsize{5.8}{6.4}\selectfont \textcolor{BrickRed}{$ \blacktriangledown\,2.2\times $}}\\[-.2pt]36.8s\\[-.2pt]{\fontsize{5.8}{6.4}\selectfont \textcolor{BrickRed}{$ \blacktriangledown\,2.2\times $}}\end{tabular} & \hspace*{1pt}\begin{tabular}[t]{@{}c@{}}\rule{0pt}{12pt}25.2s\\[-.2pt]{\fontsize{5.8}{6.4}\selectfont \textcolor{BrickRed}{$ \blacktriangledown\,17\% $}}\\[-.2pt]22.7s\\[-.2pt]{\fontsize{5.8}{6.4}\selectfont \textcolor{BrickRed}{$ \blacktriangledown\,12\% $}}\end{tabular} & \hspace*{1pt}\begin{tabular}[t]{@{}c@{}}\rule{0pt}{12pt}1m\,56.9s\\[-.2pt]{\fontsize{5.8}{6.4}\selectfont \textcolor{BrickRed}{$ \blacktriangledown\,8.0\times $}}\\[-.2pt]1m\,53.2s\\[-.2pt]{\fontsize{5.8}{6.4}\selectfont \textcolor{BrickRed}{$ \blacktriangledown\,8.2\times $}}\end{tabular} & \hspace*{1pt}\begin{tabular}[t]{@{}c@{}}\rule{0pt}{12pt}24.9s\\[-.2pt]{\fontsize{5.8}{6.4}\selectfont \textcolor{BrickRed}{$ \blacktriangledown\,19\% $}}\\[-.2pt]22.2s\\[-.2pt]{\fontsize{5.8}{6.4}\selectfont \textcolor{BrickRed}{$ \blacktriangledown\,12\% $}}\end{tabular} & \hspace*{1pt}\begin{tabular}[t]{@{}c@{}}\rule{0pt}{12pt}28.6s\\[-.2pt]{\fontsize{5.8}{6.4}\selectfont \textcolor{BrickRed}{$ \blacktriangledown\,11.5\times $}}\\[-.2pt]25.8s\\[-.2pt]{\fontsize{5.8}{6.4}\selectfont \textcolor{BrickRed}{$ \blacktriangledown\,15.6\times $}}\end{tabular} & \hspace*{1pt}\begin{tabular}[t]{@{}c@{}}\rule{0pt}{12pt}15.3s\\[-.2pt]{\fontsize{5.8}{6.4}\selectfont \textcolor{BrickRed}{$ \blacktriangledown\,2.8\times $}}\\[-.2pt]12.7s\\[-.2pt]{\fontsize{5.8}{6.4}\selectfont \textcolor{BrickRed}{$ \blacktriangledown\,3.3\times $}}\end{tabular} & \hspace*{1pt}\begin{tabular}[t]{@{}c@{}}\rule{0pt}{12pt}29.3s\\[-.2pt]{\fontsize{5.8}{6.4}\selectfont \textcolor{BrickRed}{$ \blacktriangledown\,15.0\times $}}\\[-.2pt]26.2s\\[-.2pt]{\fontsize{5.8}{6.4}\selectfont \textcolor{BrickRed}{$ \blacktriangledown\,22.9\times $}}\end{tabular} & \hspace*{1pt}\begin{tabular}[t]{@{}c@{}}\rule{0pt}{12pt}8.7s\\[-.2pt]{\fontsize{5.8}{6.4}\selectfont \textcolor{BrickRed}{$ \blacktriangledown\,1.7\times $}}\\[-.2pt]8.1s\\[-.2pt]{\fontsize{5.8}{6.4}\selectfont \textcolor{BrickRed}{$ \blacktriangledown\,2.3\times $}}\end{tabular} & \hspace*{1pt}\begin{tabular}[t]{@{}c@{}}\rule{0pt}{12pt}\colorbox{BenchBest}{\strut 2m\,52.4s}\\[-.2pt]{\fontsize{5.8}{6.4}\selectfont \textcolor{ForestGreen}{$ \blacktriangle\,19\% $}}\\[-.2pt]\colorbox{BenchBest}{\strut 2m\,48.5s}\\[-.2pt]{\fontsize{5.8}{6.4}\selectfont \textcolor{ForestGreen}{$ \blacktriangle\,15\% $}}\end{tabular} & \hspace*{1pt}\begin{tabular}[t]{@{}c@{}}\rule{0pt}{12pt}\colorbox{BenchBest}{\strut 2m\,17.1s}\\[-.2pt]{\fontsize{5.8}{6.4}\selectfont \textcolor{ForestGreen}{$ \blacktriangle\,3.9\times $}}\\[-.2pt]\colorbox{BenchBest}{\strut 2m\,15.9s}\\[-.2pt]{\fontsize{5.8}{6.4}\selectfont \textcolor{ForestGreen}{$ \blacktriangle\,1.7\times $}}\end{tabular} & \hspace*{1pt}\begin{tabular}[t]{@{}c@{}}\rule{0pt}{12pt}2m\,55s\\[-.2pt]{\fontsize{5.8}{6.4}\selectfont \textcolor{BrickRed}{$ \blacktriangledown\,6\% $}}\\[-.2pt]2m\,51s\\[-.2pt]{\fontsize{5.8}{6.4}\selectfont \textcolor{BrickRed}{$ \blacktriangledown\,12\% $}}\end{tabular} & \hspace*{1pt}\begin{tabular}[t]{@{}c@{}}\rule{0pt}{12pt}\colorbox{BenchBest}{\strut 55.9s}\\[-.2pt]{\fontsize{5.8}{6.4}\selectfont \textcolor{ForestGreen}{$ \blacktriangle\,8.6\times $}}\\[-.2pt]\colorbox{BenchBest}{\strut 55.1s}\\[-.2pt]{\fontsize{5.8}{6.4}\selectfont \textcolor{ForestGreen}{$ \blacktriangle\,3.1\times $}}\end{tabular} & \hspace*{1pt}\begin{tabular}[t]{@{}c@{}}\rule{0pt}{12pt}\colorbox{BenchBest}{\strut 43.3s}\\[-.2pt]{\fontsize{5.8}{6.4}\selectfont \textcolor{ForestGreen}{$ \blacktriangle\,2.6\times $}}\\[-.2pt]\colorbox{BenchBest}{\strut 42.5s}\\[-.2pt]{\fontsize{5.8}{6.4}\selectfont \textcolor{ForestGreen}{$ \blacktriangle\,2.4\times $}}\end{tabular} & \hspace*{1pt}\begin{tabular}[t]{@{}c@{}}\rule{0pt}{12pt}\colorbox{BenchBest}{\strut 38.9s}\\[-.2pt]{\fontsize{5.8}{6.4}\selectfont \textcolor{ForestGreen}{$ \blacktriangle\,8.0\times $}}\\[-.2pt]\colorbox{BenchBest}{\strut 38.2s}\\[-.2pt]{\fontsize{5.8}{6.4}\selectfont \textcolor{ForestGreen}{$ \blacktriangle\,2.8\times $}}\end{tabular} & \hspace*{1pt}\begin{tabular}[t]{@{}c@{}}\rule{0pt}{12pt}\colorbox{BenchBest}{\strut 53.1s}\\[-.2pt]{\fontsize{5.8}{6.4}\selectfont \textcolor{ForestGreen}{$ \blacktriangle\,1.6\times $}}\\[-.2pt]\colorbox{BenchBest}{\strut 48.2s}\\[-.2pt]{\fontsize{5.8}{6.4}\selectfont \textcolor{ForestGreen}{$ \blacktriangle\,1.5\times $}}\end{tabular} & \hspace*{1pt}\begin{tabular}[t]{@{}c@{}}\rule{0pt}{12pt}\colorbox{BenchBest}{\strut 28.7s}\\[-.2pt]{\fontsize{5.8}{6.4}\selectfont \textcolor{ForestGreen}{$ \blacktriangle\,9.8\times $}}\\[-.2pt]\colorbox{BenchBest}{\strut 28.2s}\\[-.2pt]{\fontsize{5.8}{6.4}\selectfont \textcolor{ForestGreen}{$ \blacktriangle\,2.9\times $}}\end{tabular}\\[3pt]\hline
\textbf{BioseqDB} & \begin{tabular}[t]{@{}l@{}}\rule{0pt}{12pt}MPF / Bounded\\RMQ / Block\end{tabular} & \textcolor{gray}{\ensuremath{\cdot}} & \textcolor{gray}{\ensuremath{\cdot}} & \textcolor{gray}{\ensuremath{\cdot}} & \textcolor{gray}{\ensuremath{\cdot}} & \textcolor{gray}{\ensuremath{\cdot}} & \hspace*{1pt}\begin{tabular}[t]{@{}c@{}}\rule{0pt}{12pt}5.9s\\[-.2pt]{\fontsize{5.8}{6.4}\selectfont \textcolor{ForestGreen}{$ \blacktriangle\,2.6\times $}\textsuperscript{s}}\\[-.2pt]4.6s\\[-.2pt]{\fontsize{5.8}{6.4}\selectfont \textcolor{ForestGreen}{$ \blacktriangle\,2.8\times $}\textsuperscript{s}}\end{tabular} & \hspace*{1pt}\begin{tabular}[t]{@{}c@{}}\rule{0pt}{12pt}9.5s\\[-.2pt]{\fontsize{5.8}{6.4}\selectfont \textcolor{ForestGreen}{$ \blacktriangle\,3.1\times $}}\\[-.2pt]9s\\[-.2pt]{\fontsize{5.8}{6.4}\selectfont \textcolor{ForestGreen}{$ \blacktriangle\,2.9\times $}}\end{tabular} & \textcolor{gray}{\ensuremath{\cdot}} & \textcolor{gray}{\ensuremath{\cdot}} & \textcolor{gray}{\ensuremath{\cdot}} & \textcolor{gray}{\ensuremath{\cdot}} & \textcolor{gray}{\ensuremath{\cdot}} & \textcolor{gray}{\ensuremath{\cdot}} & \textcolor{gray}{\ensuremath{\cdot}} & \textcolor{gray}{\ensuremath{\cdot}} & \textcolor{gray}{\ensuremath{\cdot}}\\[3pt]\hline
 & \begin{tabular}[t]{@{}l@{}}\rule{0pt}{12pt}MPF / Bounded\\RMQ / KSW2\end{tabular} & \textcolor{gray}{\ensuremath{\cdot}} & \textcolor{gray}{\ensuremath{\cdot}} & \textcolor{gray}{\ensuremath{\cdot}} & \textcolor{gray}{\ensuremath{\cdot}} & \textcolor{gray}{\ensuremath{\cdot}} & \hspace*{1pt}\begin{tabular}[t]{@{}c@{}}\rule{0pt}{12pt}5.8s\\[-.2pt]{\fontsize{5.8}{6.4}\selectfont \textcolor{ForestGreen}{$ \blacktriangle\,2.6\times $}\textsuperscript{s}}\\[-.2pt]\colorbox{BenchBest}{\strut 3.9s}\\[-.2pt]{\fontsize{5.8}{6.4}\selectfont \textcolor{ForestGreen}{$ \blacktriangle\,3.3\times $}\textsuperscript{s}}\end{tabular} & \textcolor{gray}{\ensuremath{\cdot}} & \textcolor{gray}{\ensuremath{\cdot}} & \textcolor{gray}{\ensuremath{\cdot}} & \textcolor{gray}{\ensuremath{\cdot}} & \textcolor{gray}{\ensuremath{\cdot}} & \textcolor{gray}{\ensuremath{\cdot}} & \textcolor{gray}{\ensuremath{\cdot}} & \textcolor{gray}{\ensuremath{\cdot}} & \textcolor{gray}{\ensuremath{\cdot}} & \textcolor{gray}{\ensuremath{\cdot}}\\[3pt]\hline
 & \begin{tabular}[t]{@{}l@{}}\rule{0pt}{12pt}MPF / Bounded\\Scalar-DP / Block\end{tabular} & \textcolor{gray}{\ensuremath{\cdot}} & \textcolor{gray}{\ensuremath{\cdot}} & \textcolor{gray}{\ensuremath{\cdot}} & \textcolor{gray}{\ensuremath{\cdot}} & \textcolor{gray}{\ensuremath{\cdot}} & \hspace*{1pt}\begin{tabular}[t]{@{}c@{}}\rule{0pt}{12pt}5.4s\\[-.2pt]{\fontsize{5.8}{6.4}\selectfont \textcolor{ForestGreen}{$ \blacktriangle\,2.8\times $}\textsuperscript{s}}\\[-.2pt]4.3s\\[-.2pt]{\fontsize{5.8}{6.4}\selectfont \textcolor{ForestGreen}{$ \blacktriangle\,3.0\times $}\textsuperscript{s}}\end{tabular} & \hspace*{1pt}\begin{tabular}[t]{@{}c@{}}\rule{0pt}{12pt}2.4s\\[-.2pt]{\fontsize{5.8}{6.4}\selectfont \textcolor{ForestGreen}{$ \blacktriangle\,12.0\times $}}\\[-.2pt]1.9s\\[-.2pt]{\fontsize{5.8}{6.4}\selectfont \textcolor{ForestGreen}{$ \blacktriangle\,13.8\times $}}\end{tabular} & \textcolor{gray}{\ensuremath{\cdot}} & \textcolor{gray}{\ensuremath{\cdot}} & \textcolor{gray}{\ensuremath{\cdot}} & \textcolor{gray}{\ensuremath{\cdot}} & \textcolor{gray}{\ensuremath{\cdot}} & \textcolor{gray}{\ensuremath{\cdot}} & \textcolor{gray}{\ensuremath{\cdot}} & \textcolor{gray}{\ensuremath{\cdot}} & \textcolor{gray}{\ensuremath{\cdot}}\\[3pt]\hline
 & \begin{tabular}[t]{@{}l@{}}\rule{0pt}{12pt}MPF / Bounded\\Scalar-DP / KSW2\end{tabular} & \textcolor{gray}{\ensuremath{\cdot}} & \textcolor{gray}{\ensuremath{\cdot}} & \textcolor{gray}{\ensuremath{\cdot}} & \textcolor{gray}{\ensuremath{\cdot}} & \textcolor{gray}{\ensuremath{\cdot}} & \hspace*{1pt}\begin{tabular}[t]{@{}c@{}}\rule{0pt}{12pt}5.8s\\[-.2pt]{\fontsize{5.8}{6.4}\selectfont \textcolor{ForestGreen}{$ \blacktriangle\,2.6\times $}\textsuperscript{s}}\\[-.2pt]4.2s\\[-.2pt]{\fontsize{5.8}{6.4}\selectfont \textcolor{ForestGreen}{$ \blacktriangle\,3.1\times $}\textsuperscript{s}}\end{tabular} & \hspace*{1pt}\begin{tabular}[t]{@{}c@{}}\rule{0pt}{12pt}2.4s\\[-.2pt]{\fontsize{5.8}{6.4}\selectfont \textcolor{ForestGreen}{$ \blacktriangle\,12.3\times $}}\\[-.2pt]\colorbox{BenchBest}{\strut 1.1s}\\[-.2pt]{\fontsize{5.8}{6.4}\selectfont \textcolor{ForestGreen}{$ \blacktriangle\,22.9\times $}}\end{tabular} & \textcolor{gray}{\ensuremath{\cdot}} & \textcolor{gray}{\ensuremath{\cdot}} & \textcolor{gray}{\ensuremath{\cdot}} & \textcolor{gray}{\ensuremath{\cdot}} & \textcolor{gray}{\ensuremath{\cdot}} & \textcolor{gray}{\ensuremath{\cdot}} & \textcolor{gray}{\ensuremath{\cdot}} & \textcolor{gray}{\ensuremath{\cdot}} & \textcolor{gray}{\ensuremath{\cdot}}\\[3pt]\hline
\rowcolor{BenchPreferred}[0.5pt][0.5pt] & \begin{tabular}[t]{@{}l@{}}\rule{0pt}{12pt}MPF / Bounded\\SIMD-DP / Block\end{tabular} & \hspace*{1pt}\begin{tabular}[t]{@{}c@{}}\rule{0pt}{12pt}\colorbox{BenchBest}{\strut 17.7s}\\[-.2pt]{\fontsize{5.8}{6.4}\selectfont \textcolor{ForestGreen}{$ \blacktriangle\,2.2\times $}}\\[-.2pt]\colorbox{BenchBest}{\strut 17s}\\[-.2pt]{\fontsize{5.8}{6.4}\selectfont \textcolor{ForestGreen}{$ \blacktriangle\,2.2\times $}}\end{tabular} & \hspace*{1pt}\begin{tabular}[t]{@{}c@{}}\rule{0pt}{12pt}\colorbox{BenchBest}{\strut 21.5s}\\[-.2pt]{\fontsize{5.8}{6.4}\selectfont \textcolor{ForestGreen}{$ \blacktriangle\,15\% $}}\\[-.2pt]\colorbox{BenchBest}{\strut 20.2s}\\[-.2pt]{\fontsize{5.8}{6.4}\selectfont \textcolor{ForestGreen}{$ \blacktriangle\,11\% $}}\end{tabular} & \hspace*{1pt}\begin{tabular}[t]{@{}c@{}}\rule{0pt}{12pt}\colorbox{BenchBest}{\strut 14.6s}\\[-.2pt]{\fontsize{5.8}{6.4}\selectfont \textcolor{ForestGreen}{$ \blacktriangle\,8.0\times $}}\\[-.2pt]\colorbox{BenchBest}{\strut 13.8s}\\[-.2pt]{\fontsize{5.8}{6.4}\selectfont \textcolor{ForestGreen}{$ \blacktriangle\,8.2\times $}}\end{tabular} & \hspace*{1pt}\begin{tabular}[t]{@{}c@{}}\rule{0pt}{12pt}\colorbox{BenchBest}{\strut 21s}\\[-.2pt]{\fontsize{5.8}{6.4}\selectfont \textcolor{ForestGreen}{$ \blacktriangle\,16\% $}}\\[-.2pt]\colorbox{BenchBest}{\strut 19.8s}\\[-.2pt]{\fontsize{5.8}{6.4}\selectfont \textcolor{ForestGreen}{$ \blacktriangle\,11\% $}}\end{tabular} & \hspace*{1pt}\begin{tabular}[t]{@{}c@{}}\rule{0pt}{12pt}\colorbox{BenchBest}{\strut 2.5s}\\[-.2pt]{\fontsize{5.8}{6.4}\selectfont \textcolor{ForestGreen}{$ \blacktriangle\,11.5\times $}}\\[-.2pt]\colorbox{BenchBest}{\strut 1.7s}\\[-.2pt]{\fontsize{5.8}{6.4}\selectfont \textcolor{ForestGreen}{$ \blacktriangle\,15.6\times $}}\end{tabular} & \hspace*{1pt}\begin{tabular}[t]{@{}c@{}}\rule{0pt}{12pt}\colorbox{BenchBest}{\strut 5.4s}\\[-.2pt]{\fontsize{5.8}{6.4}\selectfont \textcolor{ForestGreen}{$ \blacktriangle\,2.8\times $}}\\[-.2pt]4.2s\\[-.2pt]{\fontsize{5.8}{6.4}\selectfont \textcolor{ForestGreen}{$ \blacktriangle\,3.1\times $}}\end{tabular} & \hspace*{1pt}\begin{tabular}[t]{@{}c@{}}\rule{0pt}{12pt}\colorbox{BenchBest}{\strut 2s}\\[-.2pt]{\fontsize{5.8}{6.4}\selectfont \textcolor{ForestGreen}{$ \blacktriangle\,15.0\times $}}\\[-.2pt]1.3s\\[-.2pt]{\fontsize{5.8}{6.4}\selectfont \textcolor{ForestGreen}{$ \blacktriangle\,20.4\times $}}\end{tabular} & \hspace*{1pt}\begin{tabular}[t]{@{}c@{}}\rule{0pt}{12pt}5.3s\\[-.2pt]{\fontsize{5.8}{6.4}\selectfont \textcolor{ForestGreen}{$ \blacktriangle\,1.6\times $}}\\[-.2pt]3.9s\\[-.2pt]{\fontsize{5.8}{6.4}\selectfont \textcolor{ForestGreen}{$ \blacktriangle\,2.1\times $}}\end{tabular} & \hspace*{1pt}\begin{tabular}[t]{@{}c@{}}\rule{0pt}{12pt}3m\,48.1s\\[-.2pt]{\fontsize{5.8}{6.4}\selectfont \textcolor{BrickRed}{$ \blacktriangledown\,1.3\times $}}\\[-.2pt]3m\,35s\\[-.2pt]{\fontsize{5.8}{6.4}\selectfont \textcolor{BrickRed}{$ \blacktriangledown\,28\% $}}\end{tabular} & \hspace*{1pt}\begin{tabular}[t]{@{}c@{}}\rule{0pt}{12pt}8m\,57.6s\\[-.2pt]{\fontsize{5.8}{6.4}\selectfont \textcolor{BrickRed}{$ \blacktriangledown\,3.9\times $}}\\[-.2pt]3m\,49.2s\\[-.2pt]{\fontsize{5.8}{6.4}\selectfont \textcolor{BrickRed}{$ \blacktriangledown\,1.7\times $}}\end{tabular} & \hspace*{1pt}\begin{tabular}[t]{@{}c@{}}\rule{0pt}{12pt}2m\,52s\\[-.2pt]{\fontsize{5.8}{6.4}\selectfont \textcolor{gray}{$ \blacktriangleright\,2\% $}}\\[-.2pt]2m\,39.3s\\[-.2pt]{\fontsize{5.8}{6.4}\selectfont \textcolor{ForestGreen}{$ \blacktriangle\,7\% $}}\end{tabular} & \hspace*{1pt}\begin{tabular}[t]{@{}c@{}}\rule{0pt}{12pt}8m\,13.1s\\[-.2pt]{\fontsize{5.8}{6.4}\selectfont \textcolor{BrickRed}{$ \blacktriangledown\,8.8\times $}}\\[-.2pt]2m\,54.5s\\[-.2pt]{\fontsize{5.8}{6.4}\selectfont \textcolor{BrickRed}{$ \blacktriangledown\,3.2\times $}}\end{tabular} & \hspace*{1pt}\begin{tabular}[t]{@{}c@{}}\rule{0pt}{12pt}1m\,54.3s\\[-.2pt]{\fontsize{5.8}{6.4}\selectfont \textcolor{BrickRed}{$ \blacktriangledown\,2.6\times $}}\\[-.2pt]1m\,45.5s\\[-.2pt]{\fontsize{5.8}{6.4}\selectfont \textcolor{BrickRed}{$ \blacktriangledown\,2.5\times $}}\end{tabular} & \hspace*{1pt}\begin{tabular}[t]{@{}c@{}}\rule{0pt}{12pt}5m\,17.3s\\[-.2pt]{\fontsize{5.8}{6.4}\selectfont \textcolor{BrickRed}{$ \blacktriangledown\,8.2\times $}}\\[-.2pt]1m\,54.9s\\[-.2pt]{\fontsize{5.8}{6.4}\selectfont \textcolor{BrickRed}{$ \blacktriangledown\,3.0\times $}}\end{tabular} & \hspace*{1pt}\begin{tabular}[t]{@{}c@{}}\rule{0pt}{12pt}1m\,28.3s\\[-.2pt]{\fontsize{5.8}{6.4}\selectfont \textcolor{BrickRed}{$ \blacktriangledown\,1.7\times $}}\\[-.2pt]1m\,16.8s\\[-.2pt]{\fontsize{5.8}{6.4}\selectfont \textcolor{BrickRed}{$ \blacktriangledown\,1.6\times $}}\end{tabular} & \hspace*{1pt}\begin{tabular}[t]{@{}c@{}}\rule{0pt}{12pt}4m\,45.7s\\[-.2pt]{\fontsize{5.8}{6.4}\selectfont \textcolor{BrickRed}{$ \blacktriangledown\,9.9\times $}}\\[-.2pt]1m\,26.6s\\[-.2pt]{\fontsize{5.8}{6.4}\selectfont \textcolor{BrickRed}{$ \blacktriangledown\,3.1\times $}}\end{tabular}\\[3pt]\hline
 & \begin{tabular}[t]{@{}l@{}}\rule{0pt}{12pt}MPF / Bounded\\SIMD-DP / KSW2\end{tabular} & \hspace*{1pt}\begin{tabular}[t]{@{}c@{}}\rule{0pt}{12pt}20.8s\\[-.2pt]{\fontsize{5.8}{6.4}\selectfont \textcolor{ForestGreen}{$ \blacktriangle\,1.9\times $}}\\[-.2pt]19.1s\\[-.2pt]{\fontsize{5.8}{6.4}\selectfont \textcolor{ForestGreen}{$ \blacktriangle\,1.9\times $}}\end{tabular} & \textcolor{gray}{\ensuremath{\cdot}} & \hspace*{1pt}\begin{tabular}[t]{@{}c@{}}\rule{0pt}{12pt}19.8s\\[-.2pt]{\fontsize{5.8}{6.4}\selectfont \textcolor{ForestGreen}{$ \blacktriangle\,5.9\times $}}\\[-.2pt]18.2s\\[-.2pt]{\fontsize{5.8}{6.4}\selectfont \textcolor{ForestGreen}{$ \blacktriangle\,6.2\times $}}\end{tabular} & \hspace*{1pt}\begin{tabular}[t]{@{}c@{}}\rule{0pt}{12pt}30.1s\\[-.2pt]{\fontsize{5.8}{6.4}\selectfont \textcolor{BrickRed}{$ \blacktriangledown\,21\% $}}\\[-.2pt]28.5s\\[-.2pt]{\fontsize{5.8}{6.4}\selectfont \textcolor{BrickRed}{$ \blacktriangledown\,28\% $}}\end{tabular} & \hspace*{1pt}\begin{tabular}[t]{@{}c@{}}\rule{0pt}{12pt}3.7s\\[-.2pt]{\fontsize{5.8}{6.4}\selectfont \textcolor{ForestGreen}{$ \blacktriangle\,7.8\times $}}\\[-.2pt]2.4s\\[-.2pt]{\fontsize{5.8}{6.4}\selectfont \textcolor{ForestGreen}{$ \blacktriangle\,10.7\times $}}\end{tabular} & \hspace*{1pt}\begin{tabular}[t]{@{}c@{}}\rule{0pt}{12pt}6s\\[-.2pt]{\fontsize{5.8}{6.4}\selectfont \textcolor{ForestGreen}{$ \blacktriangle\,2.5\times $}}\\[-.2pt]4.5s\\[-.2pt]{\fontsize{5.8}{6.4}\selectfont \textcolor{ForestGreen}{$ \blacktriangle\,2.9\times $}}\end{tabular} & \hspace*{1pt}\begin{tabular}[t]{@{}c@{}}\rule{0pt}{12pt}2.5s\\[-.2pt]{\fontsize{5.8}{6.4}\selectfont \textcolor{ForestGreen}{$ \blacktriangle\,11.7\times $}}\\[-.2pt]1.2s\\[-.2pt]{\fontsize{5.8}{6.4}\selectfont \textcolor{ForestGreen}{$ \blacktriangle\,22.6\times $}}\end{tabular} & \hspace*{1pt}\begin{tabular}[t]{@{}c@{}}\rule{0pt}{12pt}\colorbox{BenchBest}{\strut 5s}\\[-.2pt]{\fontsize{5.8}{6.4}\selectfont \textcolor{ForestGreen}{$ \blacktriangle\,1.7\times $}}\\[-.2pt]\colorbox{BenchBest}{\strut 3.6s}\\[-.2pt]{\fontsize{5.8}{6.4}\selectfont \textcolor{ForestGreen}{$ \blacktriangle\,2.3\times $}}\end{tabular} & \hspace*{1pt}\begin{tabular}[t]{@{}c@{}}\rule{0pt}{12pt}3m\,32.7s\\[-.2pt]{\fontsize{5.8}{6.4}\selectfont \textcolor{BrickRed}{$ \blacktriangledown\,23\% $}}\\[-.2pt]3m\,17.7s\\[-.2pt]{\fontsize{5.8}{6.4}\selectfont \textcolor{BrickRed}{$ \blacktriangledown\,17\% $}}\end{tabular} & \textcolor{gray}{\ensuremath{\cdot}} & \hspace*{1pt}\begin{tabular}[t]{@{}c@{}}\rule{0pt}{12pt}\colorbox{BenchBest}{\strut 2m\,44.9s}\\[-.2pt]{\fontsize{5.8}{6.4}\selectfont \textcolor{ForestGreen}{$ \blacktriangle\,6\% $}}\\[-.2pt]\colorbox{BenchBest}{\strut 2m\,32.2s}\\[-.2pt]{\fontsize{5.8}{6.4}\selectfont \textcolor{ForestGreen}{$ \blacktriangle\,11\% $}}\end{tabular} & \textcolor{gray}{\ensuremath{\cdot}} & \hspace*{1pt}\begin{tabular}[t]{@{}c@{}}\rule{0pt}{12pt}1m\,52.7s\\[-.2pt]{\fontsize{5.8}{6.4}\selectfont \textcolor{BrickRed}{$ \blacktriangledown\,2.6\times $}}\\[-.2pt]1m\,40s\\[-.2pt]{\fontsize{5.8}{6.4}\selectfont \textcolor{BrickRed}{$ \blacktriangledown\,2.4\times $}}\end{tabular} & \hspace*{1pt}\begin{tabular}[t]{@{}c@{}}\rule{0pt}{12pt}5m\,12.1s\\[-.2pt]{\fontsize{5.8}{6.4}\selectfont \textcolor{BrickRed}{$ \blacktriangledown\,8.0\times $}}\\[-.2pt]1m\,46.7s\\[-.2pt]{\fontsize{5.8}{6.4}\selectfont \textcolor{BrickRed}{$ \blacktriangledown\,2.8\times $}}\end{tabular} & \hspace*{1pt}\begin{tabular}[t]{@{}c@{}}\rule{0pt}{12pt}1m\,26.1s\\[-.2pt]{\fontsize{5.8}{6.4}\selectfont \textcolor{BrickRed}{$ \blacktriangledown\,1.6\times $}}\\[-.2pt]1m\,13.9s\\[-.2pt]{\fontsize{5.8}{6.4}\selectfont \textcolor{BrickRed}{$ \blacktriangledown\,1.5\times $}}\end{tabular} & \hspace*{1pt}\begin{tabular}[t]{@{}c@{}}\rule{0pt}{12pt}4m\,41.4s\\[-.2pt]{\fontsize{5.8}{6.4}\selectfont \textcolor{BrickRed}{$ \blacktriangledown\,9.8\times $}}\\[-.2pt]1m\,21.9s\\[-.2pt]{\fontsize{5.8}{6.4}\selectfont \textcolor{BrickRed}{$ \blacktriangledown\,2.9\times $}}\end{tabular}\\[3pt]\hline
\end{tabular}\par}
\smallskip{\footnotesize\noindent Each cell: A, comparison, B, comparison. Times cover 1,024 queries. $\cdot$ = no approved observation. ---\textsuperscript{g} = GPU unavailable, not executed. n/a = no measured comparison.\par
\noindent \textsuperscript{f}SIMD-ref request observed as scalar, compared only as that effective path. \textsuperscript{p}Accepted by explicit preservation after wrapper postflight failure, no run receipt claimed.\par
\noindent MPF/MPS = minimizers per fragment/sequence. Row order: index representation / minimizer, then chaining / base alignment. SIMD-ref is the reference minimizer, Bounded is bounded ntHash.\par
\noindent\textsuperscript{s}FAST: one measured batch of 1,024 queries. The reported P50 is that single datapoint, not a five-repeat estimate. Samples are not pooled with standard runs.\par
\noindent Blue = one lowest A and one lowest B down each machine/workload column, across continuation pages. Exact ties select the first row. Repeated baselines are marked only on page 1.\par
\noindent Light gray = minimap2 baseline. Pale yellow = preferred MPF / Bounded / SIMD-DP / Block configuration, not a claim of winning every column.\par
\noindent When all Vulkan modes are unmeasured, they share one row. Baseline comparisons use the fastest measured BioseqDB effective path separately for A and B.\par}
\clearpage
\noindent{\large\bfseries Local machine: mapping interval\hfill 2/5}\par\smallskip
{\fontsize{7.1}{8.1}\selectfont\setlength{\tabcolsep}{1pt}\setlength{\fboxsep}{2pt}\renewcommand{\arraystretch}{1.12}
\noindent\begin{tabular}{@{}p{14mm}p{27mm}|p{13.55mm}|p{13.55mm}|p{13.55mm}|p{13.55mm}|p{13.55mm}|p{13.55mm}|p{13.55mm}|p{13.55mm}|p{13.55mm}|p{13.55mm}|p{13.55mm}|p{13.55mm}|p{13.55mm}|p{13.55mm}|p{13.55mm}|p{13.55mm}|@{}}\hline
Engine & Execution & \multicolumn{8}{c|}{\textbf{Index}} & \multicolumn{8}{c|}{\textbf{No index}}\\
 &  & \multicolumn{4}{c|}{\textbf{4 Gbp}} & \multicolumn{4}{c|}{\textbf{2 Gbp}} & \multicolumn{4}{c|}{\textbf{4 Gbp}} & \multicolumn{4}{c|}{\textbf{2 Gbp}}\\
 &  & \multicolumn{2}{c|}{\textbf{k=15, w=10}} & \multicolumn{2}{c|}{\textbf{k=20, w=15}} & \multicolumn{2}{c|}{\textbf{k=15, w=10}} & \multicolumn{2}{c|}{\textbf{k=20, w=15}} & \multicolumn{2}{c|}{\textbf{k=15, w=10}} & \multicolumn{2}{c|}{\textbf{k=20, w=15}} & \multicolumn{2}{c|}{\textbf{k=15, w=10}} & \multicolumn{2}{c|}{\textbf{k=20, w=15}}\\
 &  & \multicolumn{1}{c|}{\textbf{DNA4}} & \multicolumn{1}{c|}{\textbf{DNA5}} & \multicolumn{1}{c|}{\textbf{DNA4}} & \multicolumn{1}{c|}{\textbf{DNA5}} & \multicolumn{1}{c|}{\textbf{DNA4}} & \multicolumn{1}{c|}{\textbf{DNA5}} & \multicolumn{1}{c|}{\textbf{DNA4}} & \multicolumn{1}{c|}{\textbf{DNA5}} & \multicolumn{1}{c|}{\textbf{DNA4}} & \multicolumn{1}{c|}{\textbf{DNA5}} & \multicolumn{1}{c|}{\textbf{DNA4}} & \multicolumn{1}{c|}{\textbf{DNA5}} & \multicolumn{1}{c|}{\textbf{DNA4}} & \multicolumn{1}{c|}{\textbf{DNA5}} & \multicolumn{1}{c|}{\textbf{DNA4}} & \multicolumn{1}{c|}{\textbf{DNA5}}\\
\hline
\rowcolor{BenchBaseline}[0.5pt][0.5pt]\textbf{minimap2} & \textbf{Baseline} & \hspace*{1pt}\begin{tabular}[t]{@{}c@{}}\rule{0pt}{12pt}39.7s\\[-.2pt]{\fontsize{5.8}{6.4}\selectfont \textcolor{BrickRed}{$ \blacktriangledown\,2.2\times $}}\\[-.2pt]36.8s\\[-.2pt]{\fontsize{5.8}{6.4}\selectfont \textcolor{BrickRed}{$ \blacktriangledown\,2.2\times $}}\end{tabular} & \hspace*{1pt}\begin{tabular}[t]{@{}c@{}}\rule{0pt}{12pt}25.2s\\[-.2pt]{\fontsize{5.8}{6.4}\selectfont \textcolor{BrickRed}{$ \blacktriangledown\,17\% $}}\\[-.2pt]22.7s\\[-.2pt]{\fontsize{5.8}{6.4}\selectfont \textcolor{BrickRed}{$ \blacktriangledown\,12\% $}}\end{tabular} & \hspace*{1pt}\begin{tabular}[t]{@{}c@{}}\rule{0pt}{12pt}1m\,56.9s\\[-.2pt]{\fontsize{5.8}{6.4}\selectfont \textcolor{BrickRed}{$ \blacktriangledown\,8.0\times $}}\\[-.2pt]1m\,53.2s\\[-.2pt]{\fontsize{5.8}{6.4}\selectfont \textcolor{BrickRed}{$ \blacktriangledown\,8.2\times $}}\end{tabular} & \hspace*{1pt}\begin{tabular}[t]{@{}c@{}}\rule{0pt}{12pt}24.9s\\[-.2pt]{\fontsize{5.8}{6.4}\selectfont \textcolor{BrickRed}{$ \blacktriangledown\,19\% $}}\\[-.2pt]22.2s\\[-.2pt]{\fontsize{5.8}{6.4}\selectfont \textcolor{BrickRed}{$ \blacktriangledown\,12\% $}}\end{tabular} & \hspace*{1pt}\begin{tabular}[t]{@{}c@{}}\rule{0pt}{12pt}28.6s\\[-.2pt]{\fontsize{5.8}{6.4}\selectfont \textcolor{BrickRed}{$ \blacktriangledown\,11.5\times $}}\\[-.2pt]25.8s\\[-.2pt]{\fontsize{5.8}{6.4}\selectfont \textcolor{BrickRed}{$ \blacktriangledown\,15.6\times $}}\end{tabular} & \hspace*{1pt}\begin{tabular}[t]{@{}c@{}}\rule{0pt}{12pt}15.3s\\[-.2pt]{\fontsize{5.8}{6.4}\selectfont \textcolor{BrickRed}{$ \blacktriangledown\,2.8\times $}}\\[-.2pt]12.7s\\[-.2pt]{\fontsize{5.8}{6.4}\selectfont \textcolor{BrickRed}{$ \blacktriangledown\,3.3\times $}}\end{tabular} & \hspace*{1pt}\begin{tabular}[t]{@{}c@{}}\rule{0pt}{12pt}29.3s\\[-.2pt]{\fontsize{5.8}{6.4}\selectfont \textcolor{BrickRed}{$ \blacktriangledown\,15.0\times $}}\\[-.2pt]26.2s\\[-.2pt]{\fontsize{5.8}{6.4}\selectfont \textcolor{BrickRed}{$ \blacktriangledown\,22.9\times $}}\end{tabular} & \hspace*{1pt}\begin{tabular}[t]{@{}c@{}}\rule{0pt}{12pt}8.7s\\[-.2pt]{\fontsize{5.8}{6.4}\selectfont \textcolor{BrickRed}{$ \blacktriangledown\,1.7\times $}}\\[-.2pt]8.1s\\[-.2pt]{\fontsize{5.8}{6.4}\selectfont \textcolor{BrickRed}{$ \blacktriangledown\,2.3\times $}}\end{tabular} & \hspace*{1pt}\begin{tabular}[t]{@{}c@{}}\rule{0pt}{12pt}2m\,52.4s\\[-.2pt]{\fontsize{5.8}{6.4}\selectfont \textcolor{ForestGreen}{$ \blacktriangle\,19\% $}}\\[-.2pt]2m\,48.5s\\[-.2pt]{\fontsize{5.8}{6.4}\selectfont \textcolor{ForestGreen}{$ \blacktriangle\,15\% $}}\end{tabular} & \hspace*{1pt}\begin{tabular}[t]{@{}c@{}}\rule{0pt}{12pt}2m\,17.1s\\[-.2pt]{\fontsize{5.8}{6.4}\selectfont \textcolor{ForestGreen}{$ \blacktriangle\,3.9\times $}}\\[-.2pt]2m\,15.9s\\[-.2pt]{\fontsize{5.8}{6.4}\selectfont \textcolor{ForestGreen}{$ \blacktriangle\,1.7\times $}}\end{tabular} & \hspace*{1pt}\begin{tabular}[t]{@{}c@{}}\rule{0pt}{12pt}2m\,55s\\[-.2pt]{\fontsize{5.8}{6.4}\selectfont \textcolor{BrickRed}{$ \blacktriangledown\,6\% $}}\\[-.2pt]2m\,51s\\[-.2pt]{\fontsize{5.8}{6.4}\selectfont \textcolor{BrickRed}{$ \blacktriangledown\,12\% $}}\end{tabular} & \hspace*{1pt}\begin{tabular}[t]{@{}c@{}}\rule{0pt}{12pt}55.9s\\[-.2pt]{\fontsize{5.8}{6.4}\selectfont \textcolor{ForestGreen}{$ \blacktriangle\,8.6\times $}}\\[-.2pt]55.1s\\[-.2pt]{\fontsize{5.8}{6.4}\selectfont \textcolor{ForestGreen}{$ \blacktriangle\,3.1\times $}}\end{tabular} & \hspace*{1pt}\begin{tabular}[t]{@{}c@{}}\rule{0pt}{12pt}43.3s\\[-.2pt]{\fontsize{5.8}{6.4}\selectfont \textcolor{ForestGreen}{$ \blacktriangle\,2.6\times $}}\\[-.2pt]42.5s\\[-.2pt]{\fontsize{5.8}{6.4}\selectfont \textcolor{ForestGreen}{$ \blacktriangle\,2.4\times $}}\end{tabular} & \hspace*{1pt}\begin{tabular}[t]{@{}c@{}}\rule{0pt}{12pt}38.9s\\[-.2pt]{\fontsize{5.8}{6.4}\selectfont \textcolor{ForestGreen}{$ \blacktriangle\,8.0\times $}}\\[-.2pt]38.2s\\[-.2pt]{\fontsize{5.8}{6.4}\selectfont \textcolor{ForestGreen}{$ \blacktriangle\,2.8\times $}}\end{tabular} & \hspace*{1pt}\begin{tabular}[t]{@{}c@{}}\rule{0pt}{12pt}53.1s\\[-.2pt]{\fontsize{5.8}{6.4}\selectfont \textcolor{ForestGreen}{$ \blacktriangle\,1.6\times $}}\\[-.2pt]48.2s\\[-.2pt]{\fontsize{5.8}{6.4}\selectfont \textcolor{ForestGreen}{$ \blacktriangle\,1.5\times $}}\end{tabular} & \hspace*{1pt}\begin{tabular}[t]{@{}c@{}}\rule{0pt}{12pt}28.7s\\[-.2pt]{\fontsize{5.8}{6.4}\selectfont \textcolor{ForestGreen}{$ \blacktriangle\,9.8\times $}}\\[-.2pt]28.2s\\[-.2pt]{\fontsize{5.8}{6.4}\selectfont \textcolor{ForestGreen}{$ \blacktriangle\,2.9\times $}}\end{tabular}\\[3pt]\hline
\textbf{BioseqDB} & \begin{tabular}[t]{@{}l@{}}\rule{0pt}{12pt}MPF / Bounded\\Vulkan-DP / Block\end{tabular} & \textcolor{gray}{\ensuremath{\cdot}} & \textcolor{gray}{\ensuremath{\cdot}} & \textcolor{gray}{\ensuremath{\cdot}} & \textcolor{gray}{\ensuremath{\cdot}} & \hspace*{1pt}\begin{tabular}[t]{@{}c@{}}\rule{0pt}{12pt}37.9s\\[-.2pt]{\fontsize{5.8}{6.4}\selectfont \textcolor{BrickRed}{$ \blacktriangledown\,1.3\times $}\textsuperscript{s}}\\[-.2pt]25.3s\\[-.2pt]{\fontsize{5.8}{6.4}\selectfont \textcolor{gray}{$ \blacktriangleright\,2\% $}\textsuperscript{s}}\end{tabular} & \hspace*{1pt}\begin{tabular}[t]{@{}c@{}}\rule{0pt}{12pt}4m\,4.2s\\[-.2pt]{\fontsize{5.8}{6.4}\selectfont \textcolor{BrickRed}{$ \blacktriangledown\,16.0\times $}\textsuperscript{s}}\\[-.2pt]38.6s\\[-.2pt]{\fontsize{5.8}{6.4}\selectfont \textcolor{BrickRed}{$ \blacktriangledown\,3.0\times $}\textsuperscript{s}}\end{tabular} & \textcolor{gray}{\ensuremath{\cdot}} & \textcolor{gray}{\ensuremath{\cdot}} & \textcolor{gray}{\ensuremath{\cdot}} & \textcolor{gray}{\ensuremath{\cdot}} & \textcolor{gray}{\ensuremath{\cdot}} & \textcolor{gray}{\ensuremath{\cdot}} & \textcolor{gray}{\ensuremath{\cdot}} & \textcolor{gray}{\ensuremath{\cdot}} & \textcolor{gray}{\ensuremath{\cdot}} & \textcolor{gray}{\ensuremath{\cdot}}\\[3pt]\hline
 & \begin{tabular}[t]{@{}l@{}}\rule{0pt}{12pt}MPF / Bounded\\Vulkan-DP / KSW2\end{tabular} & \textcolor{gray}{\ensuremath{\cdot}} & \textcolor{gray}{\ensuremath{\cdot}} & \textcolor{gray}{\ensuremath{\cdot}} & \textcolor{gray}{\ensuremath{\cdot}} & \hspace*{1pt}\begin{tabular}[t]{@{}c@{}}\rule{0pt}{12pt}42s\\[-.2pt]{\fontsize{5.8}{6.4}\selectfont \textcolor{BrickRed}{$ \blacktriangledown\,1.5\times $}\textsuperscript{s}}\\[-.2pt]27.9s\\[-.2pt]{\fontsize{5.8}{6.4}\selectfont \textcolor{BrickRed}{$ \blacktriangledown\,8\% $}\textsuperscript{s}}\end{tabular} & \hspace*{1pt}\begin{tabular}[t]{@{}c@{}}\rule{0pt}{12pt}3m\,59.2s\\[-.2pt]{\fontsize{5.8}{6.4}\selectfont \textcolor{BrickRed}{$ \blacktriangledown\,15.7\times $}\textsuperscript{s}}\\[-.2pt]30.8s\\[-.2pt]{\fontsize{5.8}{6.4}\selectfont \textcolor{BrickRed}{$ \blacktriangledown\,2.4\times $}\textsuperscript{s}}\end{tabular} & \textcolor{gray}{\ensuremath{\cdot}} & \textcolor{gray}{\ensuremath{\cdot}} & \textcolor{gray}{\ensuremath{\cdot}} & \textcolor{gray}{\ensuremath{\cdot}} & \textcolor{gray}{\ensuremath{\cdot}} & \textcolor{gray}{\ensuremath{\cdot}} & \textcolor{gray}{\ensuremath{\cdot}} & \textcolor{gray}{\ensuremath{\cdot}} & \textcolor{gray}{\ensuremath{\cdot}} & \textcolor{gray}{\ensuremath{\cdot}}\\[3pt]\hline
 & \begin{tabular}[t]{@{}l@{}}\rule{0pt}{12pt}MPF / Scalar\\RMQ / Block\end{tabular} & \textcolor{gray}{\ensuremath{\cdot}} & \textcolor{gray}{\ensuremath{\cdot}} & \textcolor{gray}{\ensuremath{\cdot}} & \textcolor{gray}{\ensuremath{\cdot}} & \textcolor{gray}{\ensuremath{\cdot}} & \textcolor{gray}{\ensuremath{\cdot}} & \hspace*{1pt}\begin{tabular}[t]{@{}c@{}}\rule{0pt}{12pt}1m\,50s\\[-.2pt]{\fontsize{5.8}{6.4}\selectfont \textcolor{BrickRed}{$ \blacktriangledown\,3.7\times $}}\\[-.2pt]1m\,49.6s\\[-.2pt]{\fontsize{5.8}{6.4}\selectfont \textcolor{BrickRed}{$ \blacktriangledown\,4.2\times $}}\end{tabular} & \textcolor{gray}{\ensuremath{\cdot}} & \textcolor{gray}{\ensuremath{\cdot}} & \textcolor{gray}{\ensuremath{\cdot}} & \textcolor{gray}{\ensuremath{\cdot}} & \textcolor{gray}{\ensuremath{\cdot}} & \textcolor{gray}{\ensuremath{\cdot}} & \textcolor{gray}{\ensuremath{\cdot}} & \textcolor{gray}{\ensuremath{\cdot}} & \textcolor{gray}{\ensuremath{\cdot}}\\[3pt]\hline
 & \begin{tabular}[t]{@{}l@{}}\rule{0pt}{12pt}MPF / Scalar\\Scalar-DP / Block\end{tabular} & \textcolor{gray}{\ensuremath{\cdot}} & \textcolor{gray}{\ensuremath{\cdot}} & \textcolor{gray}{\ensuremath{\cdot}} & \textcolor{gray}{\ensuremath{\cdot}} & \textcolor{gray}{\ensuremath{\cdot}} & \textcolor{gray}{\ensuremath{\cdot}} & \hspace*{1pt}\begin{tabular}[t]{@{}c@{}}\rule{0pt}{12pt}1m\,48.8s\\[-.2pt]{\fontsize{5.8}{6.4}\selectfont \textcolor{BrickRed}{$ \blacktriangledown\,3.7\times $}}\\[-.2pt]1m\,48.4s\\[-.2pt]{\fontsize{5.8}{6.4}\selectfont \textcolor{BrickRed}{$ \blacktriangledown\,4.1\times $}}\end{tabular} & \textcolor{gray}{\ensuremath{\cdot}} & \textcolor{gray}{\ensuremath{\cdot}} & \textcolor{gray}{\ensuremath{\cdot}} & \textcolor{gray}{\ensuremath{\cdot}} & \textcolor{gray}{\ensuremath{\cdot}} & \textcolor{gray}{\ensuremath{\cdot}} & \textcolor{gray}{\ensuremath{\cdot}} & \textcolor{gray}{\ensuremath{\cdot}} & \textcolor{gray}{\ensuremath{\cdot}}\\[3pt]\hline
 & \begin{tabular}[t]{@{}l@{}}\rule{0pt}{12pt}MPF / Scalar\\SIMD-DP / Block\end{tabular} & \textcolor{gray}{\ensuremath{\cdot}} & \textcolor{gray}{\ensuremath{\cdot}} & \textcolor{gray}{\ensuremath{\cdot}} & \textcolor{gray}{\ensuremath{\cdot}} & \textcolor{gray}{\ensuremath{\cdot}} & \textcolor{gray}{\ensuremath{\cdot}} & \hspace*{1pt}\begin{tabular}[t]{@{}c@{}}\rule{0pt}{12pt}1m\,52.3s\\[-.2pt]{\fontsize{5.8}{6.4}\selectfont \textcolor{BrickRed}{$ \blacktriangledown\,3.8\times $}}\\[-.2pt]1m\,51.9s\\[-.2pt]{\fontsize{5.8}{6.4}\selectfont \textcolor{BrickRed}{$ \blacktriangledown\,4.3\times $}}\end{tabular} & \textcolor{gray}{\ensuremath{\cdot}} & \textcolor{gray}{\ensuremath{\cdot}} & \textcolor{gray}{\ensuremath{\cdot}} & \textcolor{gray}{\ensuremath{\cdot}} & \textcolor{gray}{\ensuremath{\cdot}} & \textcolor{gray}{\ensuremath{\cdot}} & \textcolor{gray}{\ensuremath{\cdot}} & \textcolor{gray}{\ensuremath{\cdot}} & \textcolor{gray}{\ensuremath{\cdot}}\\[3pt]\hline
 & \begin{tabular}[t]{@{}l@{}}\rule{0pt}{12pt}MPF / Scalar\\Vulkan-DP / Block\end{tabular} & \textcolor{gray}{\ensuremath{\cdot}} & \textcolor{gray}{\ensuremath{\cdot}} & \textcolor{gray}{\ensuremath{\cdot}} & \textcolor{gray}{\ensuremath{\cdot}} & \textcolor{gray}{\ensuremath{\cdot}} & \textcolor{gray}{\ensuremath{\cdot}} & \hspace*{1pt}\begin{tabular}[t]{@{}c@{}}\rule{0pt}{12pt}12m\,37.6s\\[-.2pt]{\fontsize{5.8}{6.4}\selectfont \textcolor{BrickRed}{$ \blacktriangledown\,25.8\times $}}\\[-.2pt]12m\,26.5s\\[-.2pt]{\fontsize{5.8}{6.4}\selectfont \textcolor{BrickRed}{$ \blacktriangledown\,28.5\times $}}\end{tabular} & \textcolor{gray}{\ensuremath{\cdot}} & \textcolor{gray}{\ensuremath{\cdot}} & \textcolor{gray}{\ensuremath{\cdot}} & \textcolor{gray}{\ensuremath{\cdot}} & \textcolor{gray}{\ensuremath{\cdot}} & \textcolor{gray}{\ensuremath{\cdot}} & \textcolor{gray}{\ensuremath{\cdot}} & \textcolor{gray}{\ensuremath{\cdot}} & \textcolor{gray}{\ensuremath{\cdot}}\\[3pt]\hline
 & \begin{tabular}[t]{@{}l@{}}\rule{0pt}{12pt}MPF / SIMD-ref\\RMQ / Block\end{tabular} & \textcolor{gray}{\ensuremath{\cdot}} & \textcolor{gray}{\ensuremath{\cdot}} & \textcolor{gray}{\ensuremath{\cdot}} & \textcolor{gray}{\ensuremath{\cdot}} & \textcolor{gray}{\ensuremath{\cdot}} & \textcolor{gray}{\ensuremath{\cdot}} & \hspace*{1pt}\begin{tabular}[t]{@{}c@{}}\rule{0pt}{12pt}1m\,40.4s\\[-.2pt]{\fontsize{5.8}{6.4}\selectfont \textcolor{BrickRed}{$ \blacktriangledown\,3.4\times $}}\\[-.2pt]1m\,40s\\[-.2pt]{\fontsize{5.8}{6.4}\selectfont \textcolor{BrickRed}{$ \blacktriangledown\,3.8\times $}}\end{tabular} & \textcolor{gray}{\ensuremath{\cdot}} & \textcolor{gray}{\ensuremath{\cdot}} & \textcolor{gray}{\ensuremath{\cdot}} & \textcolor{gray}{\ensuremath{\cdot}} & \textcolor{gray}{\ensuremath{\cdot}} & \textcolor{gray}{\ensuremath{\cdot}} & \textcolor{gray}{\ensuremath{\cdot}} & \textcolor{gray}{\ensuremath{\cdot}} & \textcolor{gray}{\ensuremath{\cdot}}\\[3pt]\hline
\end{tabular}\par}
\smallskip{\footnotesize\noindent Each cell: A, comparison, B, comparison. Times cover 1,024 queries. $\cdot$ = no approved observation. ---\textsuperscript{g} = GPU unavailable, not executed. n/a = no measured comparison.\par
\noindent \textsuperscript{f}SIMD-ref request observed as scalar, compared only as that effective path. \textsuperscript{p}Accepted by explicit preservation after wrapper postflight failure, no run receipt claimed.\par
\noindent MPF/MPS = minimizers per fragment/sequence. Row order: index representation / minimizer, then chaining / base alignment. SIMD-ref is the reference minimizer, Bounded is bounded ntHash.\par
\noindent\textsuperscript{s}FAST: one measured batch of 1,024 queries. The reported P50 is that single datapoint, not a five-repeat estimate. Samples are not pooled with standard runs.\par
\noindent Blue = one lowest A and one lowest B down each machine/workload column, across continuation pages. Exact ties select the first row. Repeated baselines are marked only on page 1.\par
\noindent Light gray = minimap2 baseline. Pale yellow = preferred MPF / Bounded / SIMD-DP / Block configuration, not a claim of winning every column.\par
\noindent When all Vulkan modes are unmeasured, they share one row. Baseline comparisons use the fastest measured BioseqDB effective path separately for A and B.\par}
\clearpage
\noindent{\large\bfseries Local machine: mapping interval\hfill 3/5}\par\smallskip
{\fontsize{7.1}{8.1}\selectfont\setlength{\tabcolsep}{1pt}\setlength{\fboxsep}{2pt}\renewcommand{\arraystretch}{1.12}
\noindent\begin{tabular}{@{}p{14mm}p{27mm}|p{13.55mm}|p{13.55mm}|p{13.55mm}|p{13.55mm}|p{13.55mm}|p{13.55mm}|p{13.55mm}|p{13.55mm}|p{13.55mm}|p{13.55mm}|p{13.55mm}|p{13.55mm}|p{13.55mm}|p{13.55mm}|p{13.55mm}|p{13.55mm}|@{}}\hline
Engine & Execution & \multicolumn{8}{c|}{\textbf{Index}} & \multicolumn{8}{c|}{\textbf{No index}}\\
 &  & \multicolumn{4}{c|}{\textbf{4 Gbp}} & \multicolumn{4}{c|}{\textbf{2 Gbp}} & \multicolumn{4}{c|}{\textbf{4 Gbp}} & \multicolumn{4}{c|}{\textbf{2 Gbp}}\\
 &  & \multicolumn{2}{c|}{\textbf{k=15, w=10}} & \multicolumn{2}{c|}{\textbf{k=20, w=15}} & \multicolumn{2}{c|}{\textbf{k=15, w=10}} & \multicolumn{2}{c|}{\textbf{k=20, w=15}} & \multicolumn{2}{c|}{\textbf{k=15, w=10}} & \multicolumn{2}{c|}{\textbf{k=20, w=15}} & \multicolumn{2}{c|}{\textbf{k=15, w=10}} & \multicolumn{2}{c|}{\textbf{k=20, w=15}}\\
 &  & \multicolumn{1}{c|}{\textbf{DNA4}} & \multicolumn{1}{c|}{\textbf{DNA5}} & \multicolumn{1}{c|}{\textbf{DNA4}} & \multicolumn{1}{c|}{\textbf{DNA5}} & \multicolumn{1}{c|}{\textbf{DNA4}} & \multicolumn{1}{c|}{\textbf{DNA5}} & \multicolumn{1}{c|}{\textbf{DNA4}} & \multicolumn{1}{c|}{\textbf{DNA5}} & \multicolumn{1}{c|}{\textbf{DNA4}} & \multicolumn{1}{c|}{\textbf{DNA5}} & \multicolumn{1}{c|}{\textbf{DNA4}} & \multicolumn{1}{c|}{\textbf{DNA5}} & \multicolumn{1}{c|}{\textbf{DNA4}} & \multicolumn{1}{c|}{\textbf{DNA5}} & \multicolumn{1}{c|}{\textbf{DNA4}} & \multicolumn{1}{c|}{\textbf{DNA5}}\\
\hline
\rowcolor{BenchBaseline}[0.5pt][0.5pt]\textbf{minimap2} & \textbf{Baseline} & \hspace*{1pt}\begin{tabular}[t]{@{}c@{}}\rule{0pt}{12pt}39.7s\\[-.2pt]{\fontsize{5.8}{6.4}\selectfont \textcolor{BrickRed}{$ \blacktriangledown\,2.2\times $}}\\[-.2pt]36.8s\\[-.2pt]{\fontsize{5.8}{6.4}\selectfont \textcolor{BrickRed}{$ \blacktriangledown\,2.2\times $}}\end{tabular} & \hspace*{1pt}\begin{tabular}[t]{@{}c@{}}\rule{0pt}{12pt}25.2s\\[-.2pt]{\fontsize{5.8}{6.4}\selectfont \textcolor{BrickRed}{$ \blacktriangledown\,17\% $}}\\[-.2pt]22.7s\\[-.2pt]{\fontsize{5.8}{6.4}\selectfont \textcolor{BrickRed}{$ \blacktriangledown\,12\% $}}\end{tabular} & \hspace*{1pt}\begin{tabular}[t]{@{}c@{}}\rule{0pt}{12pt}1m\,56.9s\\[-.2pt]{\fontsize{5.8}{6.4}\selectfont \textcolor{BrickRed}{$ \blacktriangledown\,8.0\times $}}\\[-.2pt]1m\,53.2s\\[-.2pt]{\fontsize{5.8}{6.4}\selectfont \textcolor{BrickRed}{$ \blacktriangledown\,8.2\times $}}\end{tabular} & \hspace*{1pt}\begin{tabular}[t]{@{}c@{}}\rule{0pt}{12pt}24.9s\\[-.2pt]{\fontsize{5.8}{6.4}\selectfont \textcolor{BrickRed}{$ \blacktriangledown\,19\% $}}\\[-.2pt]22.2s\\[-.2pt]{\fontsize{5.8}{6.4}\selectfont \textcolor{BrickRed}{$ \blacktriangledown\,12\% $}}\end{tabular} & \hspace*{1pt}\begin{tabular}[t]{@{}c@{}}\rule{0pt}{12pt}28.6s\\[-.2pt]{\fontsize{5.8}{6.4}\selectfont \textcolor{BrickRed}{$ \blacktriangledown\,11.5\times $}}\\[-.2pt]25.8s\\[-.2pt]{\fontsize{5.8}{6.4}\selectfont \textcolor{BrickRed}{$ \blacktriangledown\,15.6\times $}}\end{tabular} & \hspace*{1pt}\begin{tabular}[t]{@{}c@{}}\rule{0pt}{12pt}15.3s\\[-.2pt]{\fontsize{5.8}{6.4}\selectfont \textcolor{BrickRed}{$ \blacktriangledown\,2.8\times $}}\\[-.2pt]12.7s\\[-.2pt]{\fontsize{5.8}{6.4}\selectfont \textcolor{BrickRed}{$ \blacktriangledown\,3.3\times $}}\end{tabular} & \hspace*{1pt}\begin{tabular}[t]{@{}c@{}}\rule{0pt}{12pt}29.3s\\[-.2pt]{\fontsize{5.8}{6.4}\selectfont \textcolor{BrickRed}{$ \blacktriangledown\,15.0\times $}}\\[-.2pt]26.2s\\[-.2pt]{\fontsize{5.8}{6.4}\selectfont \textcolor{BrickRed}{$ \blacktriangledown\,22.9\times $}}\end{tabular} & \hspace*{1pt}\begin{tabular}[t]{@{}c@{}}\rule{0pt}{12pt}8.7s\\[-.2pt]{\fontsize{5.8}{6.4}\selectfont \textcolor{BrickRed}{$ \blacktriangledown\,1.7\times $}}\\[-.2pt]8.1s\\[-.2pt]{\fontsize{5.8}{6.4}\selectfont \textcolor{BrickRed}{$ \blacktriangledown\,2.3\times $}}\end{tabular} & \hspace*{1pt}\begin{tabular}[t]{@{}c@{}}\rule{0pt}{12pt}2m\,52.4s\\[-.2pt]{\fontsize{5.8}{6.4}\selectfont \textcolor{ForestGreen}{$ \blacktriangle\,19\% $}}\\[-.2pt]2m\,48.5s\\[-.2pt]{\fontsize{5.8}{6.4}\selectfont \textcolor{ForestGreen}{$ \blacktriangle\,15\% $}}\end{tabular} & \hspace*{1pt}\begin{tabular}[t]{@{}c@{}}\rule{0pt}{12pt}2m\,17.1s\\[-.2pt]{\fontsize{5.8}{6.4}\selectfont \textcolor{ForestGreen}{$ \blacktriangle\,3.9\times $}}\\[-.2pt]2m\,15.9s\\[-.2pt]{\fontsize{5.8}{6.4}\selectfont \textcolor{ForestGreen}{$ \blacktriangle\,1.7\times $}}\end{tabular} & \hspace*{1pt}\begin{tabular}[t]{@{}c@{}}\rule{0pt}{12pt}2m\,55s\\[-.2pt]{\fontsize{5.8}{6.4}\selectfont \textcolor{BrickRed}{$ \blacktriangledown\,6\% $}}\\[-.2pt]2m\,51s\\[-.2pt]{\fontsize{5.8}{6.4}\selectfont \textcolor{BrickRed}{$ \blacktriangledown\,12\% $}}\end{tabular} & \hspace*{1pt}\begin{tabular}[t]{@{}c@{}}\rule{0pt}{12pt}55.9s\\[-.2pt]{\fontsize{5.8}{6.4}\selectfont \textcolor{ForestGreen}{$ \blacktriangle\,8.6\times $}}\\[-.2pt]55.1s\\[-.2pt]{\fontsize{5.8}{6.4}\selectfont \textcolor{ForestGreen}{$ \blacktriangle\,3.1\times $}}\end{tabular} & \hspace*{1pt}\begin{tabular}[t]{@{}c@{}}\rule{0pt}{12pt}43.3s\\[-.2pt]{\fontsize{5.8}{6.4}\selectfont \textcolor{ForestGreen}{$ \blacktriangle\,2.6\times $}}\\[-.2pt]42.5s\\[-.2pt]{\fontsize{5.8}{6.4}\selectfont \textcolor{ForestGreen}{$ \blacktriangle\,2.4\times $}}\end{tabular} & \hspace*{1pt}\begin{tabular}[t]{@{}c@{}}\rule{0pt}{12pt}38.9s\\[-.2pt]{\fontsize{5.8}{6.4}\selectfont \textcolor{ForestGreen}{$ \blacktriangle\,8.0\times $}}\\[-.2pt]38.2s\\[-.2pt]{\fontsize{5.8}{6.4}\selectfont \textcolor{ForestGreen}{$ \blacktriangle\,2.8\times $}}\end{tabular} & \hspace*{1pt}\begin{tabular}[t]{@{}c@{}}\rule{0pt}{12pt}53.1s\\[-.2pt]{\fontsize{5.8}{6.4}\selectfont \textcolor{ForestGreen}{$ \blacktriangle\,1.6\times $}}\\[-.2pt]48.2s\\[-.2pt]{\fontsize{5.8}{6.4}\selectfont \textcolor{ForestGreen}{$ \blacktriangle\,1.5\times $}}\end{tabular} & \hspace*{1pt}\begin{tabular}[t]{@{}c@{}}\rule{0pt}{12pt}28.7s\\[-.2pt]{\fontsize{5.8}{6.4}\selectfont \textcolor{ForestGreen}{$ \blacktriangle\,9.8\times $}}\\[-.2pt]28.2s\\[-.2pt]{\fontsize{5.8}{6.4}\selectfont \textcolor{ForestGreen}{$ \blacktriangle\,2.9\times $}}\end{tabular}\\[3pt]\hline
\textbf{BioseqDB} & \begin{tabular}[t]{@{}l@{}}\rule{0pt}{12pt}MPF / SIMD-ref\\Scalar-DP / Block\end{tabular} & \textcolor{gray}{\ensuremath{\cdot}} & \textcolor{gray}{\ensuremath{\cdot}} & \textcolor{gray}{\ensuremath{\cdot}} & \textcolor{gray}{\ensuremath{\cdot}} & \textcolor{gray}{\ensuremath{\cdot}} & \textcolor{gray}{\ensuremath{\cdot}} & \hspace*{1pt}\begin{tabular}[t]{@{}c@{}}\rule{0pt}{12pt}1m\,34.2s\\[-.2pt]{\fontsize{5.8}{6.4}\selectfont \textcolor{BrickRed}{$ \blacktriangledown\,3.2\times $}}\\[-.2pt]1m\,34s\\[-.2pt]{\fontsize{5.8}{6.4}\selectfont \textcolor{BrickRed}{$ \blacktriangledown\,3.6\times $}}\end{tabular} & \textcolor{gray}{\ensuremath{\cdot}} & \textcolor{gray}{\ensuremath{\cdot}} & \textcolor{gray}{\ensuremath{\cdot}} & \textcolor{gray}{\ensuremath{\cdot}} & \textcolor{gray}{\ensuremath{\cdot}} & \textcolor{gray}{\ensuremath{\cdot}} & \textcolor{gray}{\ensuremath{\cdot}} & \textcolor{gray}{\ensuremath{\cdot}} & \textcolor{gray}{\ensuremath{\cdot}}\\[3pt]\hline
 & \begin{tabular}[t]{@{}l@{}}\rule{0pt}{12pt}MPF / SIMD-ref\\SIMD-DP / Block\end{tabular} & \textcolor{gray}{\ensuremath{\cdot}} & \textcolor{gray}{\ensuremath{\cdot}} & \textcolor{gray}{\ensuremath{\cdot}} & \textcolor{gray}{\ensuremath{\cdot}} & \textcolor{gray}{\ensuremath{\cdot}} & \textcolor{gray}{\ensuremath{\cdot}} & \hspace*{1pt}\begin{tabular}[t]{@{}c@{}}\rule{0pt}{12pt}1m\,42.9s\\[-.2pt]{\fontsize{5.8}{6.4}\selectfont \textcolor{BrickRed}{$ \blacktriangledown\,3.5\times $}}\\[-.2pt]1m\,42.6s\\[-.2pt]{\fontsize{5.8}{6.4}\selectfont \textcolor{BrickRed}{$ \blacktriangledown\,3.9\times $}}\end{tabular} & \textcolor{gray}{\ensuremath{\cdot}} & \textcolor{gray}{\ensuremath{\cdot}} & \textcolor{gray}{\ensuremath{\cdot}} & \textcolor{gray}{\ensuremath{\cdot}} & \textcolor{gray}{\ensuremath{\cdot}} & \textcolor{gray}{\ensuremath{\cdot}} & \textcolor{gray}{\ensuremath{\cdot}} & \textcolor{gray}{\ensuremath{\cdot}} & \textcolor{gray}{\ensuremath{\cdot}}\\[3pt]\hline
 & \begin{tabular}[t]{@{}l@{}}\rule{0pt}{12pt}MPF / Two-stack\\RMQ / Block\end{tabular} & \textcolor{gray}{\ensuremath{\cdot}} & \textcolor{gray}{\ensuremath{\cdot}} & \textcolor{gray}{\ensuremath{\cdot}} & \textcolor{gray}{\ensuremath{\cdot}} & \textcolor{gray}{\ensuremath{\cdot}} & \textcolor{gray}{\ensuremath{\cdot}} & \hspace*{1pt}\begin{tabular}[t]{@{}c@{}}\rule{0pt}{12pt}1m\,35.2s\\[-.2pt]{\fontsize{5.8}{6.4}\selectfont \textcolor{BrickRed}{$ \blacktriangledown\,3.2\times $}}\\[-.2pt]1m\,34.5s\\[-.2pt]{\fontsize{5.8}{6.4}\selectfont \textcolor{BrickRed}{$ \blacktriangledown\,3.6\times $}}\end{tabular} & \textcolor{gray}{\ensuremath{\cdot}} & \textcolor{gray}{\ensuremath{\cdot}} & \textcolor{gray}{\ensuremath{\cdot}} & \textcolor{gray}{\ensuremath{\cdot}} & \textcolor{gray}{\ensuremath{\cdot}} & \textcolor{gray}{\ensuremath{\cdot}} & \textcolor{gray}{\ensuremath{\cdot}} & \textcolor{gray}{\ensuremath{\cdot}} & \textcolor{gray}{\ensuremath{\cdot}}\\[3pt]\hline
 & \begin{tabular}[t]{@{}l@{}}\rule{0pt}{12pt}MPF / Two-stack\\RMQ / KSW2\end{tabular} & \textcolor{gray}{\ensuremath{\cdot}} & \textcolor{gray}{\ensuremath{\cdot}} & \textcolor{gray}{\ensuremath{\cdot}} & \textcolor{gray}{\ensuremath{\cdot}} & \textcolor{gray}{\ensuremath{\cdot}} & \textcolor{gray}{\ensuremath{\cdot}} & \hspace*{1pt}\begin{tabular}[t]{@{}c@{}}\rule{0pt}{12pt}1m\,12.2s\\[-.2pt]{\fontsize{5.8}{6.4}\selectfont \textcolor{BrickRed}{$ \blacktriangledown\,2.5\times $}}\\[-.2pt]1m\,10.5s\\[-.2pt]{\fontsize{5.8}{6.4}\selectfont \textcolor{BrickRed}{$ \blacktriangledown\,2.7\times $}}\end{tabular} & \textcolor{gray}{\ensuremath{\cdot}} & \textcolor{gray}{\ensuremath{\cdot}} & \textcolor{gray}{\ensuremath{\cdot}} & \textcolor{gray}{\ensuremath{\cdot}} & \textcolor{gray}{\ensuremath{\cdot}} & \textcolor{gray}{\ensuremath{\cdot}} & \textcolor{gray}{\ensuremath{\cdot}} & \textcolor{gray}{\ensuremath{\cdot}} & \textcolor{gray}{\ensuremath{\cdot}}\\[3pt]\hline
 & \begin{tabular}[t]{@{}l@{}}\rule{0pt}{12pt}MPF / Two-stack\\Scalar-DP / Block\end{tabular} & \textcolor{gray}{\ensuremath{\cdot}} & \textcolor{gray}{\ensuremath{\cdot}} & \textcolor{gray}{\ensuremath{\cdot}} & \textcolor{gray}{\ensuremath{\cdot}} & \textcolor{gray}{\ensuremath{\cdot}} & \textcolor{gray}{\ensuremath{\cdot}} & \hspace*{1pt}\begin{tabular}[t]{@{}c@{}}\rule{0pt}{12pt}1m\,21s\\[-.2pt]{\fontsize{5.8}{6.4}\selectfont \textcolor{BrickRed}{$ \blacktriangledown\,2.8\times $}}\\[-.2pt]1m\,20.6s\\[-.2pt]{\fontsize{5.8}{6.4}\selectfont \textcolor{BrickRed}{$ \blacktriangledown\,3.1\times $}}\end{tabular} & \textcolor{gray}{\ensuremath{\cdot}} & \textcolor{gray}{\ensuremath{\cdot}} & \textcolor{gray}{\ensuremath{\cdot}} & \textcolor{gray}{\ensuremath{\cdot}} & \textcolor{gray}{\ensuremath{\cdot}} & \textcolor{gray}{\ensuremath{\cdot}} & \textcolor{gray}{\ensuremath{\cdot}} & \textcolor{gray}{\ensuremath{\cdot}} & \textcolor{gray}{\ensuremath{\cdot}}\\[3pt]\hline
 & \begin{tabular}[t]{@{}l@{}}\rule{0pt}{12pt}MPF / Two-stack\\Scalar-DP / KSW2\end{tabular} & \textcolor{gray}{\ensuremath{\cdot}} & \textcolor{gray}{\ensuremath{\cdot}} & \textcolor{gray}{\ensuremath{\cdot}} & \textcolor{gray}{\ensuremath{\cdot}} & \textcolor{gray}{\ensuremath{\cdot}} & \textcolor{gray}{\ensuremath{\cdot}} & \hspace*{1pt}\begin{tabular}[t]{@{}c@{}}\rule{0pt}{12pt}1m\,21.1s\\[-.2pt]{\fontsize{5.8}{6.4}\selectfont \textcolor{BrickRed}{$ \blacktriangledown\,2.8\times $}}\\[-.2pt]1m\,19.3s\\[-.2pt]{\fontsize{5.8}{6.4}\selectfont \textcolor{BrickRed}{$ \blacktriangledown\,3.0\times $}}\end{tabular} & \textcolor{gray}{\ensuremath{\cdot}} & \textcolor{gray}{\ensuremath{\cdot}} & \textcolor{gray}{\ensuremath{\cdot}} & \textcolor{gray}{\ensuremath{\cdot}} & \textcolor{gray}{\ensuremath{\cdot}} & \textcolor{gray}{\ensuremath{\cdot}} & \textcolor{gray}{\ensuremath{\cdot}} & \textcolor{gray}{\ensuremath{\cdot}} & \textcolor{gray}{\ensuremath{\cdot}}\\[3pt]\hline
\end{tabular}\par}
\smallskip{\footnotesize\noindent Each cell: A, comparison, B, comparison. Times cover 1,024 queries. $\cdot$ = no approved observation. ---\textsuperscript{g} = GPU unavailable, not executed. n/a = no measured comparison.\par
\noindent \textsuperscript{f}SIMD-ref request observed as scalar, compared only as that effective path. \textsuperscript{p}Accepted by explicit preservation after wrapper postflight failure, no run receipt claimed.\par
\noindent MPF/MPS = minimizers per fragment/sequence. Row order: index representation / minimizer, then chaining / base alignment. SIMD-ref is the reference minimizer, Bounded is bounded ntHash.\par
\noindent Blue = one lowest A and one lowest B down each machine/workload column, across continuation pages. Exact ties select the first row. Repeated baselines are marked only on page 1.\par
\noindent Light gray = minimap2 baseline. Pale yellow = preferred MPF / Bounded / SIMD-DP / Block configuration, not a claim of winning every column.\par
\noindent When all Vulkan modes are unmeasured, they share one row. Baseline comparisons use the fastest measured BioseqDB effective path separately for A and B.\par}
\clearpage
\noindent{\large\bfseries Local machine: mapping interval\hfill 4/5}\par\smallskip
{\fontsize{7.1}{8.1}\selectfont\setlength{\tabcolsep}{1pt}\setlength{\fboxsep}{2pt}\renewcommand{\arraystretch}{1.12}
\noindent\begin{tabular}{@{}p{14mm}p{27mm}|p{13.55mm}|p{13.55mm}|p{13.55mm}|p{13.55mm}|p{13.55mm}|p{13.55mm}|p{13.55mm}|p{13.55mm}|p{13.55mm}|p{13.55mm}|p{13.55mm}|p{13.55mm}|p{13.55mm}|p{13.55mm}|p{13.55mm}|p{13.55mm}|@{}}\hline
Engine & Execution & \multicolumn{8}{c|}{\textbf{Index}} & \multicolumn{8}{c|}{\textbf{No index}}\\
 &  & \multicolumn{4}{c|}{\textbf{4 Gbp}} & \multicolumn{4}{c|}{\textbf{2 Gbp}} & \multicolumn{4}{c|}{\textbf{4 Gbp}} & \multicolumn{4}{c|}{\textbf{2 Gbp}}\\
 &  & \multicolumn{2}{c|}{\textbf{k=15, w=10}} & \multicolumn{2}{c|}{\textbf{k=20, w=15}} & \multicolumn{2}{c|}{\textbf{k=15, w=10}} & \multicolumn{2}{c|}{\textbf{k=20, w=15}} & \multicolumn{2}{c|}{\textbf{k=15, w=10}} & \multicolumn{2}{c|}{\textbf{k=20, w=15}} & \multicolumn{2}{c|}{\textbf{k=15, w=10}} & \multicolumn{2}{c|}{\textbf{k=20, w=15}}\\
 &  & \multicolumn{1}{c|}{\textbf{DNA4}} & \multicolumn{1}{c|}{\textbf{DNA5}} & \multicolumn{1}{c|}{\textbf{DNA4}} & \multicolumn{1}{c|}{\textbf{DNA5}} & \multicolumn{1}{c|}{\textbf{DNA4}} & \multicolumn{1}{c|}{\textbf{DNA5}} & \multicolumn{1}{c|}{\textbf{DNA4}} & \multicolumn{1}{c|}{\textbf{DNA5}} & \multicolumn{1}{c|}{\textbf{DNA4}} & \multicolumn{1}{c|}{\textbf{DNA5}} & \multicolumn{1}{c|}{\textbf{DNA4}} & \multicolumn{1}{c|}{\textbf{DNA5}} & \multicolumn{1}{c|}{\textbf{DNA4}} & \multicolumn{1}{c|}{\textbf{DNA5}} & \multicolumn{1}{c|}{\textbf{DNA4}} & \multicolumn{1}{c|}{\textbf{DNA5}}\\
\hline
\rowcolor{BenchBaseline}[0.5pt][0.5pt]\textbf{minimap2} & \textbf{Baseline} & \hspace*{1pt}\begin{tabular}[t]{@{}c@{}}\rule{0pt}{12pt}39.7s\\[-.2pt]{\fontsize{5.8}{6.4}\selectfont \textcolor{BrickRed}{$ \blacktriangledown\,2.2\times $}}\\[-.2pt]36.8s\\[-.2pt]{\fontsize{5.8}{6.4}\selectfont \textcolor{BrickRed}{$ \blacktriangledown\,2.2\times $}}\end{tabular} & \hspace*{1pt}\begin{tabular}[t]{@{}c@{}}\rule{0pt}{12pt}25.2s\\[-.2pt]{\fontsize{5.8}{6.4}\selectfont \textcolor{BrickRed}{$ \blacktriangledown\,17\% $}}\\[-.2pt]22.7s\\[-.2pt]{\fontsize{5.8}{6.4}\selectfont \textcolor{BrickRed}{$ \blacktriangledown\,12\% $}}\end{tabular} & \hspace*{1pt}\begin{tabular}[t]{@{}c@{}}\rule{0pt}{12pt}1m\,56.9s\\[-.2pt]{\fontsize{5.8}{6.4}\selectfont \textcolor{BrickRed}{$ \blacktriangledown\,8.0\times $}}\\[-.2pt]1m\,53.2s\\[-.2pt]{\fontsize{5.8}{6.4}\selectfont \textcolor{BrickRed}{$ \blacktriangledown\,8.2\times $}}\end{tabular} & \hspace*{1pt}\begin{tabular}[t]{@{}c@{}}\rule{0pt}{12pt}24.9s\\[-.2pt]{\fontsize{5.8}{6.4}\selectfont \textcolor{BrickRed}{$ \blacktriangledown\,19\% $}}\\[-.2pt]22.2s\\[-.2pt]{\fontsize{5.8}{6.4}\selectfont \textcolor{BrickRed}{$ \blacktriangledown\,12\% $}}\end{tabular} & \hspace*{1pt}\begin{tabular}[t]{@{}c@{}}\rule{0pt}{12pt}28.6s\\[-.2pt]{\fontsize{5.8}{6.4}\selectfont \textcolor{BrickRed}{$ \blacktriangledown\,11.5\times $}}\\[-.2pt]25.8s\\[-.2pt]{\fontsize{5.8}{6.4}\selectfont \textcolor{BrickRed}{$ \blacktriangledown\,15.6\times $}}\end{tabular} & \hspace*{1pt}\begin{tabular}[t]{@{}c@{}}\rule{0pt}{12pt}15.3s\\[-.2pt]{\fontsize{5.8}{6.4}\selectfont \textcolor{BrickRed}{$ \blacktriangledown\,2.8\times $}}\\[-.2pt]12.7s\\[-.2pt]{\fontsize{5.8}{6.4}\selectfont \textcolor{BrickRed}{$ \blacktriangledown\,3.3\times $}}\end{tabular} & \hspace*{1pt}\begin{tabular}[t]{@{}c@{}}\rule{0pt}{12pt}29.3s\\[-.2pt]{\fontsize{5.8}{6.4}\selectfont \textcolor{BrickRed}{$ \blacktriangledown\,15.0\times $}}\\[-.2pt]26.2s\\[-.2pt]{\fontsize{5.8}{6.4}\selectfont \textcolor{BrickRed}{$ \blacktriangledown\,22.9\times $}}\end{tabular} & \hspace*{1pt}\begin{tabular}[t]{@{}c@{}}\rule{0pt}{12pt}8.7s\\[-.2pt]{\fontsize{5.8}{6.4}\selectfont \textcolor{BrickRed}{$ \blacktriangledown\,1.7\times $}}\\[-.2pt]8.1s\\[-.2pt]{\fontsize{5.8}{6.4}\selectfont \textcolor{BrickRed}{$ \blacktriangledown\,2.3\times $}}\end{tabular} & \hspace*{1pt}\begin{tabular}[t]{@{}c@{}}\rule{0pt}{12pt}2m\,52.4s\\[-.2pt]{\fontsize{5.8}{6.4}\selectfont \textcolor{ForestGreen}{$ \blacktriangle\,19\% $}}\\[-.2pt]2m\,48.5s\\[-.2pt]{\fontsize{5.8}{6.4}\selectfont \textcolor{ForestGreen}{$ \blacktriangle\,15\% $}}\end{tabular} & \hspace*{1pt}\begin{tabular}[t]{@{}c@{}}\rule{0pt}{12pt}2m\,17.1s\\[-.2pt]{\fontsize{5.8}{6.4}\selectfont \textcolor{ForestGreen}{$ \blacktriangle\,3.9\times $}}\\[-.2pt]2m\,15.9s\\[-.2pt]{\fontsize{5.8}{6.4}\selectfont \textcolor{ForestGreen}{$ \blacktriangle\,1.7\times $}}\end{tabular} & \hspace*{1pt}\begin{tabular}[t]{@{}c@{}}\rule{0pt}{12pt}2m\,55s\\[-.2pt]{\fontsize{5.8}{6.4}\selectfont \textcolor{BrickRed}{$ \blacktriangledown\,6\% $}}\\[-.2pt]2m\,51s\\[-.2pt]{\fontsize{5.8}{6.4}\selectfont \textcolor{BrickRed}{$ \blacktriangledown\,12\% $}}\end{tabular} & \hspace*{1pt}\begin{tabular}[t]{@{}c@{}}\rule{0pt}{12pt}55.9s\\[-.2pt]{\fontsize{5.8}{6.4}\selectfont \textcolor{ForestGreen}{$ \blacktriangle\,8.6\times $}}\\[-.2pt]55.1s\\[-.2pt]{\fontsize{5.8}{6.4}\selectfont \textcolor{ForestGreen}{$ \blacktriangle\,3.1\times $}}\end{tabular} & \hspace*{1pt}\begin{tabular}[t]{@{}c@{}}\rule{0pt}{12pt}43.3s\\[-.2pt]{\fontsize{5.8}{6.4}\selectfont \textcolor{ForestGreen}{$ \blacktriangle\,2.6\times $}}\\[-.2pt]42.5s\\[-.2pt]{\fontsize{5.8}{6.4}\selectfont \textcolor{ForestGreen}{$ \blacktriangle\,2.4\times $}}\end{tabular} & \hspace*{1pt}\begin{tabular}[t]{@{}c@{}}\rule{0pt}{12pt}38.9s\\[-.2pt]{\fontsize{5.8}{6.4}\selectfont \textcolor{ForestGreen}{$ \blacktriangle\,8.0\times $}}\\[-.2pt]38.2s\\[-.2pt]{\fontsize{5.8}{6.4}\selectfont \textcolor{ForestGreen}{$ \blacktriangle\,2.8\times $}}\end{tabular} & \hspace*{1pt}\begin{tabular}[t]{@{}c@{}}\rule{0pt}{12pt}53.1s\\[-.2pt]{\fontsize{5.8}{6.4}\selectfont \textcolor{ForestGreen}{$ \blacktriangle\,1.6\times $}}\\[-.2pt]48.2s\\[-.2pt]{\fontsize{5.8}{6.4}\selectfont \textcolor{ForestGreen}{$ \blacktriangle\,1.5\times $}}\end{tabular} & \hspace*{1pt}\begin{tabular}[t]{@{}c@{}}\rule{0pt}{12pt}28.7s\\[-.2pt]{\fontsize{5.8}{6.4}\selectfont \textcolor{ForestGreen}{$ \blacktriangle\,9.8\times $}}\\[-.2pt]28.2s\\[-.2pt]{\fontsize{5.8}{6.4}\selectfont \textcolor{ForestGreen}{$ \blacktriangle\,2.9\times $}}\end{tabular}\\[3pt]\hline
\textbf{BioseqDB} & \begin{tabular}[t]{@{}l@{}}\rule{0pt}{12pt}MPF / Two-stack\\SIMD-DP / Block\end{tabular} & \textcolor{gray}{\ensuremath{\cdot}} & \textcolor{gray}{\ensuremath{\cdot}} & \textcolor{gray}{\ensuremath{\cdot}} & \textcolor{gray}{\ensuremath{\cdot}} & \textcolor{gray}{\ensuremath{\cdot}} & \textcolor{gray}{\ensuremath{\cdot}} & \hspace*{1pt}\begin{tabular}[t]{@{}c@{}}\rule{0pt}{12pt}1m\,45s\\[-.2pt]{\fontsize{5.8}{6.4}\selectfont \textcolor{BrickRed}{$ \blacktriangledown\,3.6\times $}}\\[-.2pt]1m\,44.4s\\[-.2pt]{\fontsize{5.8}{6.4}\selectfont \textcolor{BrickRed}{$ \blacktriangledown\,4.0\times $}}\end{tabular} & \textcolor{gray}{\ensuremath{\cdot}} & \textcolor{gray}{\ensuremath{\cdot}} & \textcolor{gray}{\ensuremath{\cdot}} & \textcolor{gray}{\ensuremath{\cdot}} & \textcolor{gray}{\ensuremath{\cdot}} & \textcolor{gray}{\ensuremath{\cdot}} & \textcolor{gray}{\ensuremath{\cdot}} & \textcolor{gray}{\ensuremath{\cdot}} & \textcolor{gray}{\ensuremath{\cdot}}\\[3pt]\hline
 & \begin{tabular}[t]{@{}l@{}}\rule{0pt}{12pt}MPF / Two-stack\\SIMD-DP / KSW2\end{tabular} & \textcolor{gray}{\ensuremath{\cdot}} & \textcolor{gray}{\ensuremath{\cdot}} & \textcolor{gray}{\ensuremath{\cdot}} & \textcolor{gray}{\ensuremath{\cdot}} & \textcolor{gray}{\ensuremath{\cdot}} & \textcolor{gray}{\ensuremath{\cdot}} & \hspace*{1pt}\begin{tabular}[t]{@{}c@{}}\rule{0pt}{12pt}1m\,21.3s\\[-.2pt]{\fontsize{5.8}{6.4}\selectfont \textcolor{BrickRed}{$ \blacktriangledown\,2.8\times $}}\\[-.2pt]1m\,19.5s\\[-.2pt]{\fontsize{5.8}{6.4}\selectfont \textcolor{BrickRed}{$ \blacktriangledown\,3.0\times $}}\end{tabular} & \textcolor{gray}{\ensuremath{\cdot}} & \textcolor{gray}{\ensuremath{\cdot}} & \textcolor{gray}{\ensuremath{\cdot}} & \textcolor{gray}{\ensuremath{\cdot}} & \textcolor{gray}{\ensuremath{\cdot}} & \textcolor{gray}{\ensuremath{\cdot}} & \textcolor{gray}{\ensuremath{\cdot}} & \textcolor{gray}{\ensuremath{\cdot}} & \textcolor{gray}{\ensuremath{\cdot}}\\[3pt]\hline
 & \begin{tabular}[t]{@{}l@{}}\rule{0pt}{12pt}MPF / Two-stack\\Vulkan-DP / Block\end{tabular} & \textcolor{gray}{\ensuremath{\cdot}} & \textcolor{gray}{\ensuremath{\cdot}} & \textcolor{gray}{\ensuremath{\cdot}} & \textcolor{gray}{\ensuremath{\cdot}} & \textcolor{gray}{\ensuremath{\cdot}} & \textcolor{gray}{\ensuremath{\cdot}} & \hspace*{1pt}\begin{tabular}[t]{@{}c@{}}\rule{0pt}{12pt}12m\,53.5s\\[-.2pt]{\fontsize{5.8}{6.4}\selectfont \textcolor{BrickRed}{$ \blacktriangledown\,26.4\times $}}\\[-.2pt]12m\,25.3s\\[-.2pt]{\fontsize{5.8}{6.4}\selectfont \textcolor{BrickRed}{$ \blacktriangledown\,28.4\times $}}\end{tabular} & \textcolor{gray}{\ensuremath{\cdot}} & \textcolor{gray}{\ensuremath{\cdot}} & \textcolor{gray}{\ensuremath{\cdot}} & \textcolor{gray}{\ensuremath{\cdot}} & \textcolor{gray}{\ensuremath{\cdot}} & \textcolor{gray}{\ensuremath{\cdot}} & \textcolor{gray}{\ensuremath{\cdot}} & \textcolor{gray}{\ensuremath{\cdot}} & \textcolor{gray}{\ensuremath{\cdot}}\\[3pt]\hline
 & \begin{tabular}[t]{@{}l@{}}\rule{0pt}{12pt}MPF / Two-stack\\Vulkan-DP / KSW2\end{tabular} & \textcolor{gray}{\ensuremath{\cdot}} & \textcolor{gray}{\ensuremath{\cdot}} & \textcolor{gray}{\ensuremath{\cdot}} & \textcolor{gray}{\ensuremath{\cdot}} & \textcolor{gray}{\ensuremath{\cdot}} & \textcolor{gray}{\ensuremath{\cdot}} & \hspace*{1pt}\begin{tabular}[t]{@{}c@{}}\rule{0pt}{12pt}13m\,56.9s\\[-.2pt]{\fontsize{5.8}{6.4}\selectfont \textcolor{BrickRed}{$ \blacktriangledown\,28.5\times $}}\\[-.2pt]13m\,25.1s\\[-.2pt]{\fontsize{5.8}{6.4}\selectfont \textcolor{BrickRed}{$ \blacktriangledown\,30.7\times $}}\end{tabular} & \textcolor{gray}{\ensuremath{\cdot}} & \textcolor{gray}{\ensuremath{\cdot}} & \textcolor{gray}{\ensuremath{\cdot}} & \textcolor{gray}{\ensuremath{\cdot}} & \textcolor{gray}{\ensuremath{\cdot}} & \textcolor{gray}{\ensuremath{\cdot}} & \textcolor{gray}{\ensuremath{\cdot}} & \textcolor{gray}{\ensuremath{\cdot}} & \textcolor{gray}{\ensuremath{\cdot}}\\[3pt]\hline
 & \begin{tabular}[t]{@{}l@{}}\rule{0pt}{12pt}MPF / Vulkan\\Scalar-DP / Block\end{tabular} & \textcolor{gray}{\ensuremath{\cdot}} & \textcolor{gray}{\ensuremath{\cdot}} & \textcolor{gray}{\ensuremath{\cdot}} & \textcolor{gray}{\ensuremath{\cdot}} & \textcolor{gray}{\ensuremath{\cdot}} & \textcolor{gray}{\ensuremath{\cdot}} & \hspace*{1pt}\begin{tabular}[t]{@{}c@{}}\rule{0pt}{12pt}3m\,41.9s\\[-.2pt]{\fontsize{5.8}{6.4}\selectfont \textcolor{BrickRed}{$ \blacktriangledown\,7.6\times $}}\\[-.2pt]3m\,31.8s\\[-.2pt]{\fontsize{5.8}{6.4}\selectfont \textcolor{BrickRed}{$ \blacktriangledown\,8.1\times $}}\end{tabular} & \textcolor{gray}{\ensuremath{\cdot}} & \textcolor{gray}{\ensuremath{\cdot}} & \textcolor{gray}{\ensuremath{\cdot}} & \textcolor{gray}{\ensuremath{\cdot}} & \textcolor{gray}{\ensuremath{\cdot}} & \textcolor{gray}{\ensuremath{\cdot}} & \textcolor{gray}{\ensuremath{\cdot}} & \textcolor{gray}{\ensuremath{\cdot}} & \textcolor{gray}{\ensuremath{\cdot}}\\[3pt]\hline
 & \begin{tabular}[t]{@{}l@{}}\rule{0pt}{12pt}MPF / Vulkan\\Scalar-DP / KSW2\end{tabular} & \textcolor{gray}{\ensuremath{\cdot}} & \textcolor{gray}{\ensuremath{\cdot}} & \textcolor{gray}{\ensuremath{\cdot}} & \textcolor{gray}{\ensuremath{\cdot}} & \textcolor{gray}{\ensuremath{\cdot}} & \textcolor{gray}{\ensuremath{\cdot}} & \hspace*{1pt}\begin{tabular}[t]{@{}c@{}}\rule{0pt}{12pt}3m\,25.8s\\[-.2pt]{\fontsize{5.8}{6.4}\selectfont \textcolor{BrickRed}{$ \blacktriangledown\,7.0\times $}}\\[-.2pt]3m\,14.4s\\[-.2pt]{\fontsize{5.8}{6.4}\selectfont \textcolor{BrickRed}{$ \blacktriangledown\,7.4\times $}}\end{tabular} & \textcolor{gray}{\ensuremath{\cdot}} & \textcolor{gray}{\ensuremath{\cdot}} & \textcolor{gray}{\ensuremath{\cdot}} & \textcolor{gray}{\ensuremath{\cdot}} & \textcolor{gray}{\ensuremath{\cdot}} & \textcolor{gray}{\ensuremath{\cdot}} & \textcolor{gray}{\ensuremath{\cdot}} & \textcolor{gray}{\ensuremath{\cdot}} & \textcolor{gray}{\ensuremath{\cdot}}\\[3pt]\hline
 & \begin{tabular}[t]{@{}l@{}}\rule{0pt}{12pt}MPF / Vulkan\\SIMD-DP / Block\end{tabular} & \textcolor{gray}{\ensuremath{\cdot}} & \textcolor{gray}{\ensuremath{\cdot}} & \textcolor{gray}{\ensuremath{\cdot}} & \textcolor{gray}{\ensuremath{\cdot}} & \textcolor{gray}{\ensuremath{\cdot}} & \textcolor{gray}{\ensuremath{\cdot}} & \hspace*{1pt}\begin{tabular}[t]{@{}c@{}}\rule{0pt}{12pt}4m\,42.2s\\[-.2pt]{\fontsize{5.8}{6.4}\selectfont \textcolor{BrickRed}{$ \blacktriangledown\,9.6\times $}}\\[-.2pt]4m\,30.5s\\[-.2pt]{\fontsize{5.8}{6.4}\selectfont \textcolor{BrickRed}{$ \blacktriangledown\,10.3\times $}}\end{tabular} & \textcolor{gray}{\ensuremath{\cdot}} & \textcolor{gray}{\ensuremath{\cdot}} & \textcolor{gray}{\ensuremath{\cdot}} & \textcolor{gray}{\ensuremath{\cdot}} & \textcolor{gray}{\ensuremath{\cdot}} & \textcolor{gray}{\ensuremath{\cdot}} & \textcolor{gray}{\ensuremath{\cdot}} & \textcolor{gray}{\ensuremath{\cdot}} & \textcolor{gray}{\ensuremath{\cdot}}\\[3pt]\hline
\end{tabular}\par}
\smallskip{\footnotesize\noindent Each cell: A, comparison, B, comparison. Times cover 1,024 queries. $\cdot$ = no approved observation. ---\textsuperscript{g} = GPU unavailable, not executed. n/a = no measured comparison.\par
\noindent \textsuperscript{f}SIMD-ref request observed as scalar, compared only as that effective path. \textsuperscript{p}Accepted by explicit preservation after wrapper postflight failure, no run receipt claimed.\par
\noindent MPF/MPS = minimizers per fragment/sequence. Row order: index representation / minimizer, then chaining / base alignment. SIMD-ref is the reference minimizer, Bounded is bounded ntHash.\par
\noindent Blue = one lowest A and one lowest B down each machine/workload column, across continuation pages. Exact ties select the first row. Repeated baselines are marked only on page 1.\par
\noindent Light gray = minimap2 baseline. Pale yellow = preferred MPF / Bounded / SIMD-DP / Block configuration, not a claim of winning every column.\par
\noindent When all Vulkan modes are unmeasured, they share one row. Baseline comparisons use the fastest measured BioseqDB effective path separately for A and B.\par}
\clearpage
\noindent{\large\bfseries Local machine: mapping interval\hfill 5/5}\par\smallskip
{\fontsize{7.1}{8.1}\selectfont\setlength{\tabcolsep}{1pt}\setlength{\fboxsep}{2pt}\renewcommand{\arraystretch}{1.12}
\noindent\begin{tabular}{@{}p{14mm}p{27mm}|p{13.55mm}|p{13.55mm}|p{13.55mm}|p{13.55mm}|p{13.55mm}|p{13.55mm}|p{13.55mm}|p{13.55mm}|p{13.55mm}|p{13.55mm}|p{13.55mm}|p{13.55mm}|p{13.55mm}|p{13.55mm}|p{13.55mm}|p{13.55mm}|@{}}\hline
Engine & Execution & \multicolumn{8}{c|}{\textbf{Index}} & \multicolumn{8}{c|}{\textbf{No index}}\\
 &  & \multicolumn{4}{c|}{\textbf{4 Gbp}} & \multicolumn{4}{c|}{\textbf{2 Gbp}} & \multicolumn{4}{c|}{\textbf{4 Gbp}} & \multicolumn{4}{c|}{\textbf{2 Gbp}}\\
 &  & \multicolumn{2}{c|}{\textbf{k=15, w=10}} & \multicolumn{2}{c|}{\textbf{k=20, w=15}} & \multicolumn{2}{c|}{\textbf{k=15, w=10}} & \multicolumn{2}{c|}{\textbf{k=20, w=15}} & \multicolumn{2}{c|}{\textbf{k=15, w=10}} & \multicolumn{2}{c|}{\textbf{k=20, w=15}} & \multicolumn{2}{c|}{\textbf{k=15, w=10}} & \multicolumn{2}{c|}{\textbf{k=20, w=15}}\\
 &  & \multicolumn{1}{c|}{\textbf{DNA4}} & \multicolumn{1}{c|}{\textbf{DNA5}} & \multicolumn{1}{c|}{\textbf{DNA4}} & \multicolumn{1}{c|}{\textbf{DNA5}} & \multicolumn{1}{c|}{\textbf{DNA4}} & \multicolumn{1}{c|}{\textbf{DNA5}} & \multicolumn{1}{c|}{\textbf{DNA4}} & \multicolumn{1}{c|}{\textbf{DNA5}} & \multicolumn{1}{c|}{\textbf{DNA4}} & \multicolumn{1}{c|}{\textbf{DNA5}} & \multicolumn{1}{c|}{\textbf{DNA4}} & \multicolumn{1}{c|}{\textbf{DNA5}} & \multicolumn{1}{c|}{\textbf{DNA4}} & \multicolumn{1}{c|}{\textbf{DNA5}} & \multicolumn{1}{c|}{\textbf{DNA4}} & \multicolumn{1}{c|}{\textbf{DNA5}}\\
\hline
\rowcolor{BenchBaseline}[0.5pt][0.5pt]\textbf{minimap2} & \textbf{Baseline} & \hspace*{1pt}\begin{tabular}[t]{@{}c@{}}\rule{0pt}{12pt}39.7s\\[-.2pt]{\fontsize{5.8}{6.4}\selectfont \textcolor{BrickRed}{$ \blacktriangledown\,2.2\times $}}\\[-.2pt]36.8s\\[-.2pt]{\fontsize{5.8}{6.4}\selectfont \textcolor{BrickRed}{$ \blacktriangledown\,2.2\times $}}\end{tabular} & \hspace*{1pt}\begin{tabular}[t]{@{}c@{}}\rule{0pt}{12pt}25.2s\\[-.2pt]{\fontsize{5.8}{6.4}\selectfont \textcolor{BrickRed}{$ \blacktriangledown\,17\% $}}\\[-.2pt]22.7s\\[-.2pt]{\fontsize{5.8}{6.4}\selectfont \textcolor{BrickRed}{$ \blacktriangledown\,12\% $}}\end{tabular} & \hspace*{1pt}\begin{tabular}[t]{@{}c@{}}\rule{0pt}{12pt}1m\,56.9s\\[-.2pt]{\fontsize{5.8}{6.4}\selectfont \textcolor{BrickRed}{$ \blacktriangledown\,8.0\times $}}\\[-.2pt]1m\,53.2s\\[-.2pt]{\fontsize{5.8}{6.4}\selectfont \textcolor{BrickRed}{$ \blacktriangledown\,8.2\times $}}\end{tabular} & \hspace*{1pt}\begin{tabular}[t]{@{}c@{}}\rule{0pt}{12pt}24.9s\\[-.2pt]{\fontsize{5.8}{6.4}\selectfont \textcolor{BrickRed}{$ \blacktriangledown\,19\% $}}\\[-.2pt]22.2s\\[-.2pt]{\fontsize{5.8}{6.4}\selectfont \textcolor{BrickRed}{$ \blacktriangledown\,12\% $}}\end{tabular} & \hspace*{1pt}\begin{tabular}[t]{@{}c@{}}\rule{0pt}{12pt}28.6s\\[-.2pt]{\fontsize{5.8}{6.4}\selectfont \textcolor{BrickRed}{$ \blacktriangledown\,11.5\times $}}\\[-.2pt]25.8s\\[-.2pt]{\fontsize{5.8}{6.4}\selectfont \textcolor{BrickRed}{$ \blacktriangledown\,15.6\times $}}\end{tabular} & \hspace*{1pt}\begin{tabular}[t]{@{}c@{}}\rule{0pt}{12pt}15.3s\\[-.2pt]{\fontsize{5.8}{6.4}\selectfont \textcolor{BrickRed}{$ \blacktriangledown\,2.8\times $}}\\[-.2pt]12.7s\\[-.2pt]{\fontsize{5.8}{6.4}\selectfont \textcolor{BrickRed}{$ \blacktriangledown\,3.3\times $}}\end{tabular} & \hspace*{1pt}\begin{tabular}[t]{@{}c@{}}\rule{0pt}{12pt}29.3s\\[-.2pt]{\fontsize{5.8}{6.4}\selectfont \textcolor{BrickRed}{$ \blacktriangledown\,15.0\times $}}\\[-.2pt]26.2s\\[-.2pt]{\fontsize{5.8}{6.4}\selectfont \textcolor{BrickRed}{$ \blacktriangledown\,22.9\times $}}\end{tabular} & \hspace*{1pt}\begin{tabular}[t]{@{}c@{}}\rule{0pt}{12pt}8.7s\\[-.2pt]{\fontsize{5.8}{6.4}\selectfont \textcolor{BrickRed}{$ \blacktriangledown\,1.7\times $}}\\[-.2pt]8.1s\\[-.2pt]{\fontsize{5.8}{6.4}\selectfont \textcolor{BrickRed}{$ \blacktriangledown\,2.3\times $}}\end{tabular} & \hspace*{1pt}\begin{tabular}[t]{@{}c@{}}\rule{0pt}{12pt}2m\,52.4s\\[-.2pt]{\fontsize{5.8}{6.4}\selectfont \textcolor{ForestGreen}{$ \blacktriangle\,19\% $}}\\[-.2pt]2m\,48.5s\\[-.2pt]{\fontsize{5.8}{6.4}\selectfont \textcolor{ForestGreen}{$ \blacktriangle\,15\% $}}\end{tabular} & \hspace*{1pt}\begin{tabular}[t]{@{}c@{}}\rule{0pt}{12pt}2m\,17.1s\\[-.2pt]{\fontsize{5.8}{6.4}\selectfont \textcolor{ForestGreen}{$ \blacktriangle\,3.9\times $}}\\[-.2pt]2m\,15.9s\\[-.2pt]{\fontsize{5.8}{6.4}\selectfont \textcolor{ForestGreen}{$ \blacktriangle\,1.7\times $}}\end{tabular} & \hspace*{1pt}\begin{tabular}[t]{@{}c@{}}\rule{0pt}{12pt}2m\,55s\\[-.2pt]{\fontsize{5.8}{6.4}\selectfont \textcolor{BrickRed}{$ \blacktriangledown\,6\% $}}\\[-.2pt]2m\,51s\\[-.2pt]{\fontsize{5.8}{6.4}\selectfont \textcolor{BrickRed}{$ \blacktriangledown\,12\% $}}\end{tabular} & \hspace*{1pt}\begin{tabular}[t]{@{}c@{}}\rule{0pt}{12pt}55.9s\\[-.2pt]{\fontsize{5.8}{6.4}\selectfont \textcolor{ForestGreen}{$ \blacktriangle\,8.6\times $}}\\[-.2pt]55.1s\\[-.2pt]{\fontsize{5.8}{6.4}\selectfont \textcolor{ForestGreen}{$ \blacktriangle\,3.1\times $}}\end{tabular} & \hspace*{1pt}\begin{tabular}[t]{@{}c@{}}\rule{0pt}{12pt}43.3s\\[-.2pt]{\fontsize{5.8}{6.4}\selectfont \textcolor{ForestGreen}{$ \blacktriangle\,2.6\times $}}\\[-.2pt]42.5s\\[-.2pt]{\fontsize{5.8}{6.4}\selectfont \textcolor{ForestGreen}{$ \blacktriangle\,2.4\times $}}\end{tabular} & \hspace*{1pt}\begin{tabular}[t]{@{}c@{}}\rule{0pt}{12pt}38.9s\\[-.2pt]{\fontsize{5.8}{6.4}\selectfont \textcolor{ForestGreen}{$ \blacktriangle\,8.0\times $}}\\[-.2pt]38.2s\\[-.2pt]{\fontsize{5.8}{6.4}\selectfont \textcolor{ForestGreen}{$ \blacktriangle\,2.8\times $}}\end{tabular} & \hspace*{1pt}\begin{tabular}[t]{@{}c@{}}\rule{0pt}{12pt}53.1s\\[-.2pt]{\fontsize{5.8}{6.4}\selectfont \textcolor{ForestGreen}{$ \blacktriangle\,1.6\times $}}\\[-.2pt]48.2s\\[-.2pt]{\fontsize{5.8}{6.4}\selectfont \textcolor{ForestGreen}{$ \blacktriangle\,1.5\times $}}\end{tabular} & \hspace*{1pt}\begin{tabular}[t]{@{}c@{}}\rule{0pt}{12pt}28.7s\\[-.2pt]{\fontsize{5.8}{6.4}\selectfont \textcolor{ForestGreen}{$ \blacktriangle\,9.8\times $}}\\[-.2pt]28.2s\\[-.2pt]{\fontsize{5.8}{6.4}\selectfont \textcolor{ForestGreen}{$ \blacktriangle\,2.9\times $}}\end{tabular}\\[3pt]\hline
\textbf{BioseqDB} & \begin{tabular}[t]{@{}l@{}}\rule{0pt}{12pt}MPF / Vulkan\\SIMD-DP / KSW2\end{tabular} & \textcolor{gray}{\ensuremath{\cdot}} & \textcolor{gray}{\ensuremath{\cdot}} & \textcolor{gray}{\ensuremath{\cdot}} & \textcolor{gray}{\ensuremath{\cdot}} & \textcolor{gray}{\ensuremath{\cdot}} & \textcolor{gray}{\ensuremath{\cdot}} & \hspace*{1pt}\begin{tabular}[t]{@{}c@{}}\rule{0pt}{12pt}3m\,33.4s\\[-.2pt]{\fontsize{5.8}{6.4}\selectfont \textcolor{BrickRed}{$ \blacktriangledown\,7.3\times $}}\\[-.2pt]3m\,19s\\[-.2pt]{\fontsize{5.8}{6.4}\selectfont \textcolor{BrickRed}{$ \blacktriangledown\,7.6\times $}}\end{tabular} & \textcolor{gray}{\ensuremath{\cdot}} & \textcolor{gray}{\ensuremath{\cdot}} & \textcolor{gray}{\ensuremath{\cdot}} & \textcolor{gray}{\ensuremath{\cdot}} & \textcolor{gray}{\ensuremath{\cdot}} & \textcolor{gray}{\ensuremath{\cdot}} & \textcolor{gray}{\ensuremath{\cdot}} & \textcolor{gray}{\ensuremath{\cdot}} & \textcolor{gray}{\ensuremath{\cdot}}\\[3pt]\hline
 & \begin{tabular}[t]{@{}l@{}}\rule{0pt}{12pt}MPF / Vulkan\\Vulkan-DP / Block\end{tabular} & \hspace*{1pt}\begin{tabular}[t]{@{}c@{}}\rule{0pt}{12pt}9m\,36.5s\\[-.2pt]{\fontsize{5.8}{6.4}\selectfont \textcolor{BrickRed}{$ \blacktriangledown\,14.5\times $}}\\[-.2pt]9m\,22.3s\\[-.2pt]{\fontsize{5.8}{6.4}\selectfont \textcolor{BrickRed}{$ \blacktriangledown\,15.3\times $}}\end{tabular} & \hspace*{1pt}\begin{tabular}[t]{@{}c@{}}\rule{0pt}{12pt}13m\,49.8s\\[-.2pt]{\fontsize{5.8}{6.4}\selectfont \textcolor{BrickRed}{$ \blacktriangledown\,32.9\times $}}\\[-.2pt]8m\,45.2s\\[-.2pt]{\fontsize{5.8}{6.4}\selectfont \textcolor{BrickRed}{$ \blacktriangledown\,23.2\times $}}\end{tabular} & \hspace*{1pt}\begin{tabular}[t]{@{}c@{}}\rule{0pt}{12pt}4m\,57.7s\\[-.2pt]{\fontsize{5.8}{6.4}\selectfont \textcolor{BrickRed}{$ \blacktriangledown\,2.5\times $}}\\[-.2pt]4m\,42.7s\\[-.2pt]{\fontsize{5.8}{6.4}\selectfont \textcolor{BrickRed}{$ \blacktriangledown\,2.5\times $}}\end{tabular} & \hspace*{1pt}\begin{tabular}[t]{@{}c@{}}\rule{0pt}{12pt}10m\,50.3s\\[-.2pt]{\fontsize{5.8}{6.4}\selectfont \textcolor{BrickRed}{$ \blacktriangledown\,26.1\times $}}\\[-.2pt]5m\,42.9s\\[-.2pt]{\fontsize{5.8}{6.4}\selectfont \textcolor{BrickRed}{$ \blacktriangledown\,15.4\times $}}\end{tabular} & \hspace*{1pt}\begin{tabular}[t]{@{}c@{}}\rule{0pt}{12pt}6m\,46.4s\\[-.2pt]{\fontsize{5.8}{6.4}\selectfont \textcolor{BrickRed}{$ \blacktriangledown\,14.2\times $}}\\[-.2pt]6m\,34.1s\\[-.2pt]{\fontsize{5.8}{6.4}\selectfont \textcolor{BrickRed}{$ \blacktriangledown\,15.3\times $}}\end{tabular} & \hspace*{1pt}\begin{tabular}[t]{@{}c@{}}\rule{0pt}{12pt}8m\,46.4s\\[-.2pt]{\fontsize{5.8}{6.4}\selectfont \textcolor{BrickRed}{$ \blacktriangledown\,34.5\times $}}\\[-.2pt]5m\,22.9s\\[-.2pt]{\fontsize{5.8}{6.4}\selectfont \textcolor{BrickRed}{$ \blacktriangledown\,25.4\times $}}\end{tabular} & \hspace*{1pt}\begin{tabular}[t]{@{}c@{}}\rule{0pt}{12pt}3m\,11.5s\\[-.2pt]{\fontsize{5.8}{6.4}\selectfont \textcolor{BrickRed}{$ \blacktriangledown\,6.5\times $}}\\[-.2pt]3m\,2.2s\\[-.2pt]{\fontsize{5.8}{6.4}\selectfont \textcolor{BrickRed}{$ \blacktriangledown\,6.9\times $}}\end{tabular} & \hspace*{1pt}\begin{tabular}[t]{@{}c@{}}\rule{0pt}{12pt}6m\,30.8s\\[-.2pt]{\fontsize{5.8}{6.4}\selectfont \textcolor{BrickRed}{$ \blacktriangledown\,45.0\times $}}\\[-.2pt]3m\,13.4s\\[-.2pt]{\fontsize{5.8}{6.4}\selectfont \textcolor{BrickRed}{$ \blacktriangledown\,23.8\times $}}\end{tabular} & \hspace*{1pt}\begin{tabular}[t]{@{}c@{}}\rule{0pt}{12pt}9m\,50.5s\\[-.2pt]{\fontsize{5.8}{6.4}\selectfont \textcolor{BrickRed}{$ \blacktriangledown\,3.4\times $}}\\[-.2pt]9m\,37.9s\\[-.2pt]{\fontsize{5.8}{6.4}\selectfont \textcolor{BrickRed}{$ \blacktriangledown\,3.4\times $}}\end{tabular} & \hspace*{1pt}\begin{tabular}[t]{@{}c@{}}\rule{0pt}{12pt}14m\,21.8s\\[-.2pt]{\fontsize{5.8}{6.4}\selectfont \textcolor{BrickRed}{$ \blacktriangledown\,6.3\times $}}\\[-.2pt]9m\,10.7s\\[-.2pt]{\fontsize{5.8}{6.4}\selectfont \textcolor{BrickRed}{$ \blacktriangledown\,4.1\times $}}\end{tabular} & \hspace*{1pt}\begin{tabular}[t]{@{}c@{}}\rule{0pt}{12pt}4m\,10.1s\\[-.2pt]{\fontsize{5.8}{6.4}\selectfont \textcolor{BrickRed}{$ \blacktriangledown\,1.4\times $}}\\[-.2pt]4m\,1s\\[-.2pt]{\fontsize{5.8}{6.4}\selectfont \textcolor{BrickRed}{$ \blacktriangledown\,1.4\times $}}\end{tabular} & \hspace*{1pt}\begin{tabular}[t]{@{}c@{}}\rule{0pt}{12pt}9m\,51.7s\\[-.2pt]{\fontsize{5.8}{6.4}\selectfont \textcolor{BrickRed}{$ \blacktriangledown\,10.6\times $}}\\[-.2pt]4m\,35.5s\\[-.2pt]{\fontsize{5.8}{6.4}\selectfont \textcolor{BrickRed}{$ \blacktriangledown\,5.0\times $}}\end{tabular} & \hspace*{1pt}\begin{tabular}[t]{@{}c@{}}\rule{0pt}{12pt}6m\,16.9s\\[-.2pt]{\fontsize{5.8}{6.4}\selectfont \textcolor{BrickRed}{$ \blacktriangledown\,8.7\times $}}\\[-.2pt]6m\,6.9s\\[-.2pt]{\fontsize{5.8}{6.4}\selectfont \textcolor{BrickRed}{$ \blacktriangledown\,8.6\times $}}\end{tabular} & \hspace*{1pt}\begin{tabular}[t]{@{}c@{}}\rule{0pt}{12pt}9m\,20.4s\\[-.2pt]{\fontsize{5.8}{6.4}\selectfont \textcolor{BrickRed}{$ \blacktriangledown\,14.4\times $}}\\[-.2pt]5m\,57.2s\\[-.2pt]{\fontsize{5.8}{6.4}\selectfont \textcolor{BrickRed}{$ \blacktriangledown\,9.4\times $}}\end{tabular} & \hspace*{1pt}\begin{tabular}[t]{@{}c@{}}\rule{0pt}{12pt}3m\,10.7s\\[-.2pt]{\fontsize{5.8}{6.4}\selectfont \textcolor{BrickRed}{$ \blacktriangledown\,3.6\times $}}\\[-.2pt]3m\,0.2s\\[-.2pt]{\fontsize{5.8}{6.4}\selectfont \textcolor{BrickRed}{$ \blacktriangledown\,3.7\times $}}\end{tabular} & \hspace*{1pt}\begin{tabular}[t]{@{}c@{}}\rule{0pt}{12pt}6m\,41.2s\\[-.2pt]{\fontsize{5.8}{6.4}\selectfont \textcolor{BrickRed}{$ \blacktriangledown\,14.0\times $}}\\[-.2pt]3m\,21.5s\\[-.2pt]{\fontsize{5.8}{6.4}\selectfont \textcolor{BrickRed}{$ \blacktriangledown\,7.1\times $}}\end{tabular}\\[3pt]\hline
 & \begin{tabular}[t]{@{}l@{}}\rule{0pt}{12pt}MPF / Vulkan\\Vulkan-DP / KSW2\end{tabular} & \textcolor{gray}{\ensuremath{\cdot}} & \textcolor{gray}{\ensuremath{\cdot}} & \textcolor{gray}{\ensuremath{\cdot}} & \textcolor{gray}{\ensuremath{\cdot}} & \textcolor{gray}{\ensuremath{\cdot}} & \textcolor{gray}{\ensuremath{\cdot}} & \hspace*{1pt}\begin{tabular}[t]{@{}c@{}}\rule{0pt}{12pt}3m\,46.3s\\[-.2pt]{\fontsize{5.8}{6.4}\selectfont \textcolor{BrickRed}{$ \blacktriangledown\,7.7\times $}}\\[-.2pt]3m\,33.1s\\[-.2pt]{\fontsize{5.8}{6.4}\selectfont \textcolor{BrickRed}{$ \blacktriangledown\,8.1\times $}}\end{tabular} & \textcolor{gray}{\ensuremath{\cdot}} & \textcolor{gray}{\ensuremath{\cdot}} & \textcolor{gray}{\ensuremath{\cdot}} & \textcolor{gray}{\ensuremath{\cdot}} & \textcolor{gray}{\ensuremath{\cdot}} & \textcolor{gray}{\ensuremath{\cdot}} & \textcolor{gray}{\ensuremath{\cdot}} & \textcolor{gray}{\ensuremath{\cdot}} & \textcolor{gray}{\ensuremath{\cdot}}\\[3pt]\hline
 & \begin{tabular}[t]{@{}l@{}}\rule{0pt}{12pt}MPS / Bounded\\Scalar-DP / KSW2\end{tabular} & \textcolor{gray}{\ensuremath{\cdot}} & \textcolor{gray}{\ensuremath{\cdot}} & \textcolor{gray}{\ensuremath{\cdot}} & \textcolor{gray}{\ensuremath{\cdot}} & \textcolor{gray}{\ensuremath{\cdot}} & \textcolor{gray}{\ensuremath{\cdot}} & \hspace*{1pt}\begin{tabular}[t]{@{}c@{}}\rule{0pt}{12pt}5.5s\\[-.2pt]{\fontsize{5.8}{6.4}\selectfont \textcolor{ForestGreen}{$ \blacktriangle\,5.3\times $}}\\[-.2pt]3.9s\\[-.2pt]{\fontsize{5.8}{6.4}\selectfont \textcolor{ForestGreen}{$ \blacktriangle\,6.7\times $}}\end{tabular} & \textcolor{gray}{\ensuremath{\cdot}} & \textcolor{gray}{\ensuremath{\cdot}} & \textcolor{gray}{\ensuremath{\cdot}} & \textcolor{gray}{\ensuremath{\cdot}} & \textcolor{gray}{\ensuremath{\cdot}} & \textcolor{gray}{\ensuremath{\cdot}} & \textcolor{gray}{\ensuremath{\cdot}} & \textcolor{gray}{\ensuremath{\cdot}} & \textcolor{gray}{\ensuremath{\cdot}}\\[3pt]\hline
 & \begin{tabular}[t]{@{}l@{}}\rule{0pt}{12pt}MPS / Bounded\\SIMD-DP / Block\end{tabular} & \textcolor{gray}{\ensuremath{\cdot}} & \textcolor{gray}{\ensuremath{\cdot}} & \textcolor{gray}{\ensuremath{\cdot}} & \textcolor{gray}{\ensuremath{\cdot}} & \textcolor{gray}{\ensuremath{\cdot}} & \textcolor{gray}{\ensuremath{\cdot}} & \hspace*{1pt}\begin{tabular}[t]{@{}c@{}}\rule{0pt}{12pt}4.5s\\[-.2pt]{\fontsize{5.8}{6.4}\selectfont \textcolor{ForestGreen}{$ \blacktriangle\,6.6\times $}}\\[-.2pt]3.8s\\[-.2pt]{\fontsize{5.8}{6.4}\selectfont \textcolor{ForestGreen}{$ \blacktriangle\,6.9\times $}}\end{tabular} & \textcolor{gray}{\ensuremath{\cdot}} & \hspace*{1pt}\begin{tabular}[t]{@{}c@{}}\rule{0pt}{12pt}3m\,53.6s\\[-.2pt]{\fontsize{5.8}{6.4}\selectfont \textcolor{BrickRed}{$ \blacktriangledown\,1.4\times $}}\\[-.2pt]3m\,40.3s\\[-.2pt]{\fontsize{5.8}{6.4}\selectfont \textcolor{BrickRed}{$ \blacktriangledown\,1.3\times $}}\end{tabular} & \hspace*{1pt}\begin{tabular}[t]{@{}c@{}}\rule{0pt}{12pt}8m\,50.7s\\[-.2pt]{\fontsize{5.8}{6.4}\selectfont \textcolor{BrickRed}{$ \blacktriangledown\,3.9\times $}}\\[-.2pt]3m\,46s\\[-.2pt]{\fontsize{5.8}{6.4}\selectfont \textcolor{BrickRed}{$ \blacktriangledown\,1.7\times $}}\end{tabular} & \hspace*{1pt}\begin{tabular}[t]{@{}c@{}}\rule{0pt}{12pt}2m\,51.3s\\[-.2pt]{\fontsize{5.8}{6.4}\selectfont \textcolor{gray}{$ \blacktriangleright\,2\% $}}\\[-.2pt]2m\,38.1s\\[-.2pt]{\fontsize{5.8}{6.4}\selectfont \textcolor{ForestGreen}{$ \blacktriangle\,8\% $}}\end{tabular} & \hspace*{1pt}\begin{tabular}[t]{@{}c@{}}\rule{0pt}{12pt}7m\,58.5s\\[-.2pt]{\fontsize{5.8}{6.4}\selectfont \textcolor{BrickRed}{$ \blacktriangledown\,8.6\times $}}\\[-.2pt]2m\,48.9s\\[-.2pt]{\fontsize{5.8}{6.4}\selectfont \textcolor{BrickRed}{$ \blacktriangledown\,3.1\times $}}\end{tabular} & \hspace*{1pt}\begin{tabular}[t]{@{}c@{}}\rule{0pt}{12pt}1m\,59.6s\\[-.2pt]{\fontsize{5.8}{6.4}\selectfont \textcolor{BrickRed}{$ \blacktriangledown\,2.8\times $}}\\[-.2pt]1m\,49s\\[-.2pt]{\fontsize{5.8}{6.4}\selectfont \textcolor{BrickRed}{$ \blacktriangledown\,2.6\times $}}\end{tabular} & \hspace*{1pt}\begin{tabular}[t]{@{}c@{}}\rule{0pt}{12pt}5m\,15.7s\\[-.2pt]{\fontsize{5.8}{6.4}\selectfont \textcolor{BrickRed}{$ \blacktriangledown\,8.1\times $}}\\[-.2pt]1m\,54.7s\\[-.2pt]{\fontsize{5.8}{6.4}\selectfont \textcolor{BrickRed}{$ \blacktriangledown\,3.0\times $}}\end{tabular} & \hspace*{1pt}\begin{tabular}[t]{@{}c@{}}\rule{0pt}{12pt}1m\,28s\\[-.2pt]{\fontsize{5.8}{6.4}\selectfont \textcolor{BrickRed}{$ \blacktriangledown\,1.7\times $}}\\[-.2pt]1m\,18.3s\\[-.2pt]{\fontsize{5.8}{6.4}\selectfont \textcolor{BrickRed}{$ \blacktriangledown\,1.6\times $}}\end{tabular} & \hspace*{1pt}\begin{tabular}[t]{@{}c@{}}\rule{0pt}{12pt}4m\,46.9s\\[-.2pt]{\fontsize{5.8}{6.4}\selectfont \textcolor{BrickRed}{$ \blacktriangledown\,10.0\times $}}\\[-.2pt]1m\,25.3s\\[-.2pt]{\fontsize{5.8}{6.4}\selectfont \textcolor{BrickRed}{$ \blacktriangledown\,3.0\times $}}\end{tabular}\\[3pt]\hline
 & \begin{tabular}[t]{@{}l@{}}\rule{0pt}{12pt}MPS / Bounded\\SIMD-DP / KSW2\end{tabular} & \textcolor{gray}{\ensuremath{\cdot}} & \textcolor{gray}{\ensuremath{\cdot}} & \textcolor{gray}{\ensuremath{\cdot}} & \textcolor{gray}{\ensuremath{\cdot}} & \hspace*{1pt}\begin{tabular}[t]{@{}c@{}}\rule{0pt}{12pt}5.2s\\[-.2pt]{\fontsize{5.8}{6.4}\selectfont \textcolor{ForestGreen}{$ \blacktriangle\,5.5\times $}}\\[-.2pt]3.2s\\[-.2pt]{\fontsize{5.8}{6.4}\selectfont \textcolor{ForestGreen}{$ \blacktriangle\,8.1\times $}}\end{tabular} & \textcolor{gray}{\ensuremath{\cdot}} & \hspace*{1pt}\begin{tabular}[t]{@{}c@{}}\rule{0pt}{12pt}5.3s\\[-.2pt]{\fontsize{5.8}{6.4}\selectfont \textcolor{ForestGreen}{$ \blacktriangle\,5.5\times $}}\\[-.2pt]3.8s\\[-.2pt]{\fontsize{5.8}{6.4}\selectfont \textcolor{ForestGreen}{$ \blacktriangle\,7.0\times $}}\end{tabular} & \textcolor{gray}{\ensuremath{\cdot}} & \hspace*{1pt}\begin{tabular}[t]{@{}c@{}}\rule{0pt}{12pt}3m\,54s\\[-.2pt]{\fontsize{5.8}{6.4}\selectfont \textcolor{BrickRed}{$ \blacktriangledown\,1.4\times $}}\\[-.2pt]3m\,39s\\[-.2pt]{\fontsize{5.8}{6.4}\selectfont \textcolor{BrickRed}{$ \blacktriangledown\,30\% $}}\end{tabular} & \textcolor{gray}{\ensuremath{\cdot}} & \hspace*{1pt}\begin{tabular}[t]{@{}c@{}}\rule{0pt}{12pt}2m\,52.6s\\[-.2pt]{\fontsize{5.8}{6.4}\selectfont \textcolor{gray}{$ \blacktriangleright\,1\% $}}\\[-.2pt]2m\,36.8s\\[-.2pt]{\fontsize{5.8}{6.4}\selectfont \textcolor{ForestGreen}{$ \blacktriangle\,8\% $}}\end{tabular} & \textcolor{gray}{\ensuremath{\cdot}} & \hspace*{1pt}\begin{tabular}[t]{@{}c@{}}\rule{0pt}{12pt}2m\,3.1s\\[-.2pt]{\fontsize{5.8}{6.4}\selectfont \textcolor{BrickRed}{$ \blacktriangledown\,2.8\times $}}\\[-.2pt]1m\,50.4s\\[-.2pt]{\fontsize{5.8}{6.4}\selectfont \textcolor{BrickRed}{$ \blacktriangledown\,2.6\times $}}\end{tabular} & \hspace*{1pt}\begin{tabular}[t]{@{}c@{}}\rule{0pt}{12pt}5m\,13.3s\\[-.2pt]{\fontsize{5.8}{6.4}\selectfont \textcolor{BrickRed}{$ \blacktriangledown\,8.1\times $}}\\[-.2pt]1m\,50.9s\\[-.2pt]{\fontsize{5.8}{6.4}\selectfont \textcolor{BrickRed}{$ \blacktriangledown\,2.9\times $}}\end{tabular} & \hspace*{1pt}\begin{tabular}[t]{@{}c@{}}\rule{0pt}{12pt}1m\,30.6s\\[-.2pt]{\fontsize{5.8}{6.4}\selectfont \textcolor{BrickRed}{$ \blacktriangledown\,1.7\times $}}\\[-.2pt]1m\,17.7s\\[-.2pt]{\fontsize{5.8}{6.4}\selectfont \textcolor{BrickRed}{$ \blacktriangledown\,1.6\times $}}\end{tabular} & \hspace*{1pt}\begin{tabular}[t]{@{}c@{}}\rule{0pt}{12pt}4m\,44.7s\\[-.2pt]{\fontsize{5.8}{6.4}\selectfont \textcolor{BrickRed}{$ \blacktriangledown\,9.9\times $}}\\[-.2pt]1m\,25.3s\\[-.2pt]{\fontsize{5.8}{6.4}\selectfont \textcolor{BrickRed}{$ \blacktriangledown\,3.0\times $}}\end{tabular}\\[3pt]\hline
 & \begin{tabular}[t]{@{}l@{}}\rule{0pt}{12pt}MPS / Two-stack\\Scalar-DP / KSW2\end{tabular} & \textcolor{gray}{\ensuremath{\cdot}} & \textcolor{gray}{\ensuremath{\cdot}} & \textcolor{gray}{\ensuremath{\cdot}} & \textcolor{gray}{\ensuremath{\cdot}} & \textcolor{gray}{\ensuremath{\cdot}} & \textcolor{gray}{\ensuremath{\cdot}} & \hspace*{1pt}\begin{tabular}[t]{@{}c@{}}\rule{0pt}{12pt}21.8s\\[-.2pt]{\fontsize{5.8}{6.4}\selectfont \textcolor{ForestGreen}{$ \blacktriangle\,26\% $}}\\[-.2pt]4.8s\\[-.2pt]{\fontsize{5.8}{6.4}\selectfont \textcolor{ForestGreen}{$ \blacktriangle\,5.5\times $}}\end{tabular} & \textcolor{gray}{\ensuremath{\cdot}} & \textcolor{gray}{\ensuremath{\cdot}} & \textcolor{gray}{\ensuremath{\cdot}} & \textcolor{gray}{\ensuremath{\cdot}} & \textcolor{gray}{\ensuremath{\cdot}} & \textcolor{gray}{\ensuremath{\cdot}} & \textcolor{gray}{\ensuremath{\cdot}} & \textcolor{gray}{\ensuremath{\cdot}} & \textcolor{gray}{\ensuremath{\cdot}}\\[3pt]\hline
\end{tabular}\par}
\smallskip{\footnotesize\noindent Each cell: A, comparison, B, comparison. Times cover 1,024 queries. $\cdot$ = no approved observation. ---\textsuperscript{g} = GPU unavailable, not executed. n/a = no measured comparison.\par
\noindent \textsuperscript{f}SIMD-ref request observed as scalar, compared only as that effective path. \textsuperscript{p}Accepted by explicit preservation after wrapper postflight failure, no run receipt claimed.\par
\noindent MPF/MPS = minimizers per fragment/sequence. Row order: index representation / minimizer, then chaining / base alignment. SIMD-ref is the reference minimizer, Bounded is bounded ntHash.\par
\noindent Blue = one lowest A and one lowest B down each machine/workload column, across continuation pages. Exact ties select the first row. Repeated baselines are marked only on page 1.\par
\noindent Light gray = minimap2 baseline. Pale yellow = preferred MPF / Bounded / SIMD-DP / Block configuration, not a claim of winning every column.\par
\noindent When all Vulkan modes are unmeasured, they share one row. Baseline comparisons use the fastest measured BioseqDB effective path separately for A and B.\par}
\clearpage
\noindent{\large\bfseries AWS: 4 Gbp, k=15, w=10 (mapping interval)\hfill 1/6}\par\smallskip
{\fontsize{7.1}{8.1}\selectfont\setlength{\tabcolsep}{1pt}\setlength{\fboxsep}{2pt}\renewcommand{\arraystretch}{1.12}
\noindent\begin{tabular}{@{}p{14mm}p{27mm}|p{13.55mm}|p{13.55mm}|p{13.55mm}|p{13.55mm}|p{13.55mm}|p{13.55mm}|p{13.55mm}|p{13.55mm}|p{13.55mm}|p{13.55mm}|p{13.55mm}|p{13.55mm}|p{13.55mm}|p{13.55mm}|p{13.55mm}|p{13.55mm}|@{}}\hline
Engine & Execution & \multicolumn{8}{c|}{\textbf{Index}} & \multicolumn{8}{c|}{\textbf{No index}}\\
 &  & \multicolumn{8}{c|}{\textbf{4 Gbp}} & \multicolumn{8}{c|}{\textbf{4 Gbp}}\\
 &  & \multicolumn{8}{c|}{\textbf{k=15, w=10}} & \multicolumn{8}{c|}{\textbf{k=15, w=10}}\\
 &  & \multicolumn{4}{c|}{\textbf{DNA4}} & \multicolumn{4}{c|}{\textbf{DNA5}} & \multicolumn{4}{c|}{\textbf{DNA4}} & \multicolumn{4}{c|}{\textbf{DNA5}}\\
 &  & \multicolumn{1}{c|}{\textbf{M7A}} & \multicolumn{1}{c|}{\textbf{M7I}} & \multicolumn{1}{c|}{\textbf{M4}} & \multicolumn{1}{c|}{\textbf{C7G}} & \multicolumn{1}{c|}{\textbf{M7A}} & \multicolumn{1}{c|}{\textbf{M7I}} & \multicolumn{1}{c|}{\textbf{M4}} & \multicolumn{1}{c|}{\textbf{C7G}} & \multicolumn{1}{c|}{\textbf{M7A}} & \multicolumn{1}{c|}{\textbf{M7I}} & \multicolumn{1}{c|}{\textbf{M4}} & \multicolumn{1}{c|}{\textbf{C7G}} & \multicolumn{1}{c|}{\textbf{M7A}} & \multicolumn{1}{c|}{\textbf{M7I}} & \multicolumn{1}{c|}{\textbf{M4}} & \multicolumn{1}{c|}{\textbf{C7G}}\\
\hline
\rowcolor{BenchBaseline}[0.5pt][0.5pt]\textbf{minimap2} & \textbf{Baseline} & \hspace*{1pt}\begin{tabular}[t]{@{}c@{}}\rule{0pt}{12pt}36.2s\\[-.2pt]{\fontsize{5.8}{6.4}\selectfont \textcolor{BrickRed}{$ \blacktriangledown\,6.9\times $}}\\[-.2pt]35s\\[-.2pt]{\fontsize{5.8}{6.4}\selectfont \textcolor{BrickRed}{$ \blacktriangledown\,8.0\times $}}\end{tabular} & \hspace*{1pt}\begin{tabular}[t]{@{}c@{}}\rule{0pt}{12pt}38.1s\\[-.2pt]{\fontsize{5.8}{6.4}\selectfont \textcolor{BrickRed}{$ \blacktriangledown\,7.8\times $}}\\[-.2pt]36s\\[-.2pt]{\fontsize{5.8}{6.4}\selectfont \textcolor{BrickRed}{$ \blacktriangledown\,9.3\times $}}\end{tabular} & \hspace*{1pt}\begin{tabular}[t]{@{}c@{}}\rule{0pt}{12pt}1m\,33.5s\\[-.2pt]{\fontsize{5.8}{6.4}\selectfont \textcolor{BrickRed}{$ \blacktriangledown\,10.8\times $}}\\[-.2pt]1m\,29.7s\\[-.2pt]{\fontsize{5.8}{6.4}\selectfont \textcolor{BrickRed}{$ \blacktriangledown\,13.7\times $}}\end{tabular} & \hspace*{1pt}\begin{tabular}[t]{@{}c@{}}\rule{0pt}{12pt}\colorbox{BenchBest}{\strut 6m\,31.5s}\\[-.2pt]{\fontsize{5.8}{6.4}\selectfont \textcolor{gray}{n/a}}\\[-.2pt]\colorbox{BenchBest}{\strut 6m\,26.3s}\\[-.2pt]{\fontsize{5.8}{6.4}\selectfont \textcolor{gray}{n/a}}\end{tabular} & \hspace*{1pt}\begin{tabular}[t]{@{}c@{}}\rule{0pt}{12pt}\colorbox{BenchBest}{\strut 31.6s}\\[-.2pt]{\fontsize{5.8}{6.4}\selectfont \textcolor{gray}{n/a}}\\[-.2pt]\colorbox{BenchBest}{\strut 30.3s}\\[-.2pt]{\fontsize{5.8}{6.4}\selectfont \textcolor{gray}{n/a}}\end{tabular} & \hspace*{1pt}\begin{tabular}[t]{@{}c@{}}\rule{0pt}{12pt}\colorbox{BenchBest}{\strut 34.1s}\\[-.2pt]{\fontsize{5.8}{6.4}\selectfont \textcolor{gray}{n/a}}\\[-.2pt]\colorbox{BenchBest}{\strut 32s}\\[-.2pt]{\fontsize{5.8}{6.4}\selectfont \textcolor{gray}{n/a}}\end{tabular} & \hspace*{1pt}\begin{tabular}[t]{@{}c@{}}\rule{0pt}{12pt}\colorbox{BenchBest}{\strut 1m\,9.5s}\\[-.2pt]{\fontsize{5.8}{6.4}\selectfont \textcolor{gray}{n/a}}\\[-.2pt]\colorbox{BenchBest}{\strut 1m\,6.1s}\\[-.2pt]{\fontsize{5.8}{6.4}\selectfont \textcolor{gray}{n/a}}\end{tabular} & \hspace*{1pt}\begin{tabular}[t]{@{}c@{}}\rule{0pt}{12pt}\colorbox{BenchBest}{\strut 3m\,1s}\\[-.2pt]{\fontsize{5.8}{6.4}\selectfont \textcolor{gray}{n/a}}\\[-.2pt]\colorbox{BenchBest}{\strut 2m\,51.5s}\\[-.2pt]{\fontsize{5.8}{6.4}\selectfont \textcolor{gray}{n/a}}\end{tabular} & \hspace*{1pt}\begin{tabular}[t]{@{}c@{}}\rule{0pt}{12pt}\colorbox{BenchBest}{\strut 2m\,55.5s}\\[-.2pt]{\fontsize{5.8}{6.4}\selectfont \textcolor{ForestGreen}{$ \blacktriangle\,20\% $}}\\[-.2pt]\colorbox{BenchBest}{\strut 2m\,54s}\\[-.2pt]{\fontsize{5.8}{6.4}\selectfont \textcolor{ForestGreen}{$ \blacktriangle\,17\% $}}\end{tabular} & \hspace*{1pt}\begin{tabular}[t]{@{}c@{}}\rule{0pt}{12pt}\colorbox{BenchBest}{\strut 2m\,50.4s}\\[-.2pt]{\fontsize{5.8}{6.4}\selectfont \textcolor{ForestGreen}{$ \blacktriangle\,23\% $}}\\[-.2pt]\colorbox{BenchBest}{\strut 2m\,48.2s}\\[-.2pt]{\fontsize{5.8}{6.4}\selectfont \textcolor{ForestGreen}{$ \blacktriangle\,20\% $}}\end{tabular} & \hspace*{1pt}\begin{tabular}[t]{@{}c@{}}\rule{0pt}{12pt}\colorbox{BenchBest}{\strut 6m\,12.2s}\\[-.2pt]{\fontsize{5.8}{6.4}\selectfont \textcolor{ForestGreen}{$ \blacktriangle\,9\% $}}\\[-.2pt]\colorbox{BenchBest}{\strut 6m\,8.2s}\\[-.2pt]{\fontsize{5.8}{6.4}\selectfont \textcolor{gray}{$ \blacktriangleright\,5\% $}}\end{tabular} & \hspace*{1pt}\begin{tabular}[t]{@{}c@{}}\rule{0pt}{12pt}\colorbox{BenchBest}{\strut 44m\,20.5s}\\[-.2pt]{\fontsize{5.8}{6.4}\selectfont \textcolor{gray}{n/a}}\\[-.2pt]\colorbox{BenchBest}{\strut 44m\,15.6s}\\[-.2pt]{\fontsize{5.8}{6.4}\selectfont \textcolor{gray}{n/a}}\end{tabular} & \hspace*{1pt}\begin{tabular}[t]{@{}c@{}}\rule{0pt}{12pt}\colorbox{BenchBest}{\strut 1m\,47.1s}\\[-.2pt]{\fontsize{5.8}{6.4}\selectfont \textcolor{ForestGreen}{$ \blacktriangle\,5.5\times $}}\\[-.2pt]\colorbox{BenchBest}{\strut 1m\,45.8s}\\[-.2pt]{\fontsize{5.8}{6.4}\selectfont \textcolor{ForestGreen}{$ \blacktriangle\,2.1\times $}}\end{tabular} & \hspace*{1pt}\begin{tabular}[t]{@{}c@{}}\rule{0pt}{12pt}\colorbox{BenchBest}{\strut 1m\,47.2s}\\[-.2pt]{\fontsize{5.8}{6.4}\selectfont \textcolor{ForestGreen}{$ \blacktriangle\,4.6\times $}}\\[-.2pt]\colorbox{BenchBest}{\strut 1m\,45.1s}\\[-.2pt]{\fontsize{5.8}{6.4}\selectfont \textcolor{ForestGreen}{$ \blacktriangle\,2.2\times $}}\end{tabular} & \hspace*{1pt}\begin{tabular}[t]{@{}c@{}}\rule{0pt}{12pt}\colorbox{BenchBest}{\strut 3m\,40.3s}\\[-.2pt]{\fontsize{5.8}{6.4}\selectfont \textcolor{ForestGreen}{$ \blacktriangle\,4.2\times $}}\\[-.2pt]\colorbox{BenchBest}{\strut 3m\,36.6s}\\[-.2pt]{\fontsize{5.8}{6.4}\selectfont \textcolor{ForestGreen}{$ \blacktriangle\,1.7\times $}}\end{tabular} & \hspace*{1pt}\begin{tabular}[t]{@{}c@{}}\rule{0pt}{12pt}\colorbox{BenchBest}{\strut 29m\,43.4s}\\[-.2pt]{\fontsize{5.8}{6.4}\selectfont \textcolor{gray}{n/a}}\\[-.2pt]\colorbox{BenchBest}{\strut 29m\,38.5s}\\[-.2pt]{\fontsize{5.8}{6.4}\selectfont \textcolor{gray}{n/a}}\end{tabular}\\[3pt]\hline
\textbf{BioseqDB} & \begin{tabular}[t]{@{}l@{}}\rule{0pt}{12pt}MPF / Bounded\\RMQ / Block\end{tabular} & \hspace*{1pt}\begin{tabular}[t]{@{}c@{}}\rule{0pt}{12pt}5.4s\\[-.2pt]{\fontsize{5.8}{6.4}\selectfont \textcolor{ForestGreen}{$ \blacktriangle\,6.8\times $}}\\[-.2pt]4.6s\\[-.2pt]{\fontsize{5.8}{6.4}\selectfont \textcolor{ForestGreen}{$ \blacktriangle\,7.6\times $}}\end{tabular} & \hspace*{1pt}\begin{tabular}[t]{@{}c@{}}\rule{0pt}{12pt}9.9s\\[-.2pt]{\fontsize{5.8}{6.4}\selectfont \textcolor{ForestGreen}{$ \blacktriangle\,3.9\times $}}\\[-.2pt]9.1s\\[-.2pt]{\fontsize{5.8}{6.4}\selectfont \textcolor{ForestGreen}{$ \blacktriangle\,4.0\times $}}\end{tabular} & \textcolor{gray}{\ensuremath{\cdot}} & \textcolor{gray}{\ensuremath{\cdot}} & \textcolor{gray}{\ensuremath{\cdot}} & \textcolor{gray}{\ensuremath{\cdot}} & \textcolor{gray}{\ensuremath{\cdot}} & \textcolor{gray}{\ensuremath{\cdot}} & \hspace*{1pt}\begin{tabular}[t]{@{}c@{}}\rule{0pt}{12pt}4m\,2.1s\\[-.2pt]{\fontsize{5.8}{6.4}\selectfont \textcolor{BrickRed}{$ \blacktriangledown\,1.4\times $}}\\[-.2pt]3m\,53.3s\\[-.2pt]{\fontsize{5.8}{6.4}\selectfont \textcolor{BrickRed}{$ \blacktriangledown\,1.3\times $}}\end{tabular} & \hspace*{1pt}\begin{tabular}[t]{@{}c@{}}\rule{0pt}{12pt}3m\,42.6s\\[-.2pt]{\fontsize{5.8}{6.4}\selectfont \textcolor{BrickRed}{$ \blacktriangledown\,1.3\times $}}\\[-.2pt]3m\,31.1s\\[-.2pt]{\fontsize{5.8}{6.4}\selectfont \textcolor{BrickRed}{$ \blacktriangledown\,25\% $}}\end{tabular} & \textcolor{gray}{\ensuremath{\cdot}} & \textcolor{gray}{\ensuremath{\cdot}} & \textcolor{gray}{\ensuremath{\cdot}} & \textcolor{gray}{\ensuremath{\cdot}} & \textcolor{gray}{\ensuremath{\cdot}} & \textcolor{gray}{\ensuremath{\cdot}}\\[3pt]\hline
 & \begin{tabular}[t]{@{}l@{}}\rule{0pt}{12pt}MPF / Bounded\\RMQ / KSW2\end{tabular} & \hspace*{1pt}\begin{tabular}[t]{@{}c@{}}\rule{0pt}{12pt}6.9s\\[-.2pt]{\fontsize{5.8}{6.4}\selectfont \textcolor{ForestGreen}{$ \blacktriangle\,5.2\times $}}\\[-.2pt]4.5s\\[-.2pt]{\fontsize{5.8}{6.4}\selectfont \textcolor{ForestGreen}{$ \blacktriangle\,7.9\times $}}\end{tabular} & \hspace*{1pt}\begin{tabular}[t]{@{}c@{}}\rule{0pt}{12pt}9.6s\\[-.2pt]{\fontsize{5.8}{6.4}\selectfont \textcolor{ForestGreen}{$ \blacktriangle\,3.9\times $}}\\[-.2pt]7.2s\\[-.2pt]{\fontsize{5.8}{6.4}\selectfont \textcolor{ForestGreen}{$ \blacktriangle\,5.0\times $}}\end{tabular} & \textcolor{gray}{\ensuremath{\cdot}} & \textcolor{gray}{\ensuremath{\cdot}} & \textcolor{gray}{\ensuremath{\cdot}} & \textcolor{gray}{\ensuremath{\cdot}} & \textcolor{gray}{\ensuremath{\cdot}} & \textcolor{gray}{\ensuremath{\cdot}} & \hspace*{1pt}\begin{tabular}[t]{@{}c@{}}\rule{0pt}{12pt}3m\,41.3s\\[-.2pt]{\fontsize{5.8}{6.4}\selectfont \textcolor{BrickRed}{$ \blacktriangledown\,26\% $}}\\[-.2pt]3m\,30.7s\\[-.2pt]{\fontsize{5.8}{6.4}\selectfont \textcolor{BrickRed}{$ \blacktriangledown\,21\% $}}\end{tabular} & \hspace*{1pt}\begin{tabular}[t]{@{}c@{}}\rule{0pt}{12pt}3m\,59.5s\\[-.2pt]{\fontsize{5.8}{6.4}\selectfont \textcolor{BrickRed}{$ \blacktriangledown\,1.4\times $}}\\[-.2pt]3m\,45.5s\\[-.2pt]{\fontsize{5.8}{6.4}\selectfont \textcolor{BrickRed}{$ \blacktriangledown\,1.3\times $}}\end{tabular} & \textcolor{gray}{\ensuremath{\cdot}} & \textcolor{gray}{\ensuremath{\cdot}} & \textcolor{gray}{\ensuremath{\cdot}} & \textcolor{gray}{\ensuremath{\cdot}} & \textcolor{gray}{\ensuremath{\cdot}} & \textcolor{gray}{\ensuremath{\cdot}}\\[3pt]\hline
 & \begin{tabular}[t]{@{}l@{}}\rule{0pt}{12pt}MPF / Bounded\\Scalar-DP / Block\end{tabular} & \hspace*{1pt}\begin{tabular}[t]{@{}c@{}}\rule{0pt}{12pt}5.3s\\[-.2pt]{\fontsize{5.8}{6.4}\selectfont \textcolor{ForestGreen}{$ \blacktriangle\,6.8\times $}}\\[-.2pt]4.6s\\[-.2pt]{\fontsize{5.8}{6.4}\selectfont \textcolor{ForestGreen}{$ \blacktriangle\,7.7\times $}}\end{tabular} & \hspace*{1pt}\begin{tabular}[t]{@{}c@{}}\rule{0pt}{12pt}5.5s\\[-.2pt]{\fontsize{5.8}{6.4}\selectfont \textcolor{ForestGreen}{$ \blacktriangle\,6.9\times $}}\\[-.2pt]4.6s\\[-.2pt]{\fontsize{5.8}{6.4}\selectfont \textcolor{ForestGreen}{$ \blacktriangle\,7.8\times $}}\end{tabular} & \textcolor{gray}{\ensuremath{\cdot}} & \textcolor{gray}{\ensuremath{\cdot}} & \textcolor{gray}{\ensuremath{\cdot}} & \textcolor{gray}{\ensuremath{\cdot}} & \textcolor{gray}{\ensuremath{\cdot}} & \textcolor{gray}{\ensuremath{\cdot}} & \hspace*{1pt}\begin{tabular}[t]{@{}c@{}}\rule{0pt}{12pt}3m\,40.6s\\[-.2pt]{\fontsize{5.8}{6.4}\selectfont \textcolor{BrickRed}{$ \blacktriangledown\,26\% $}}\\[-.2pt]3m\,31.8s\\[-.2pt]{\fontsize{5.8}{6.4}\selectfont \textcolor{BrickRed}{$ \blacktriangledown\,22\% $}}\end{tabular} & \hspace*{1pt}\begin{tabular}[t]{@{}c@{}}\rule{0pt}{12pt}4m\,25.9s\\[-.2pt]{\fontsize{5.8}{6.4}\selectfont \textcolor{BrickRed}{$ \blacktriangledown\,1.6\times $}}\\[-.2pt]4m\,13.9s\\[-.2pt]{\fontsize{5.8}{6.4}\selectfont \textcolor{BrickRed}{$ \blacktriangledown\,1.5\times $}}\end{tabular} & \textcolor{gray}{\ensuremath{\cdot}} & \textcolor{gray}{\ensuremath{\cdot}} & \textcolor{gray}{\ensuremath{\cdot}} & \textcolor{gray}{\ensuremath{\cdot}} & \hspace*{1pt}\begin{tabular}[t]{@{}c@{}}\rule{0pt}{12pt}17m\,59.1s\\[-.2pt]{\fontsize{5.8}{6.4}\selectfont \textcolor{BrickRed}{$ \blacktriangledown\,4.9\times $}}\\[-.2pt]9m\,9.4s\\[-.2pt]{\fontsize{5.8}{6.4}\selectfont \textcolor{BrickRed}{$ \blacktriangledown\,2.5\times $}}\end{tabular} & \textcolor{gray}{\ensuremath{\cdot}}\\[3pt]\hline
 & \begin{tabular}[t]{@{}l@{}}\rule{0pt}{12pt}MPF / Bounded\\Scalar-DP / KSW2\end{tabular} & \hspace*{1pt}\begin{tabular}[t]{@{}c@{}}\rule{0pt}{12pt}6.9s\\[-.2pt]{\fontsize{5.8}{6.4}\selectfont \textcolor{ForestGreen}{$ \blacktriangle\,5.3\times $}}\\[-.2pt]4.4s\\[-.2pt]{\fontsize{5.8}{6.4}\selectfont \textcolor{ForestGreen}{$ \blacktriangle\,7.9\times $}}\end{tabular} & \hspace*{1pt}\begin{tabular}[t]{@{}c@{}}\rule{0pt}{12pt}7.6s\\[-.2pt]{\fontsize{5.8}{6.4}\selectfont \textcolor{ForestGreen}{$ \blacktriangle\,5.0\times $}}\\[-.2pt]5.2s\\[-.2pt]{\fontsize{5.8}{6.4}\selectfont \textcolor{ForestGreen}{$ \blacktriangle\,6.9\times $}}\end{tabular} & \textcolor{gray}{\ensuremath{\cdot}} & \textcolor{gray}{\ensuremath{\cdot}} & \textcolor{gray}{\ensuremath{\cdot}} & \textcolor{gray}{\ensuremath{\cdot}} & \textcolor{gray}{\ensuremath{\cdot}} & \textcolor{gray}{\ensuremath{\cdot}} & \hspace*{1pt}\begin{tabular}[t]{@{}c@{}}\rule{0pt}{12pt}3m\,56.3s\\[-.2pt]{\fontsize{5.8}{6.4}\selectfont \textcolor{BrickRed}{$ \blacktriangledown\,1.3\times $}}\\[-.2pt]3m\,44.9s\\[-.2pt]{\fontsize{5.8}{6.4}\selectfont \textcolor{BrickRed}{$ \blacktriangledown\,29\% $}}\end{tabular} & \hspace*{1pt}\begin{tabular}[t]{@{}c@{}}\rule{0pt}{12pt}4m\,6.8s\\[-.2pt]{\fontsize{5.8}{6.4}\selectfont \textcolor{BrickRed}{$ \blacktriangledown\,1.4\times $}}\\[-.2pt]3m\,52.3s\\[-.2pt]{\fontsize{5.8}{6.4}\selectfont \textcolor{BrickRed}{$ \blacktriangledown\,1.4\times $}}\end{tabular} & \textcolor{gray}{\ensuremath{\cdot}} & \textcolor{gray}{\ensuremath{\cdot}} & \textcolor{gray}{\ensuremath{\cdot}} & \textcolor{gray}{\ensuremath{\cdot}} & \textcolor{gray}{\ensuremath{\cdot}} & \textcolor{gray}{\ensuremath{\cdot}}\\[3pt]\hline
\rowcolor{BenchPreferred}[0.5pt][0.5pt] & \begin{tabular}[t]{@{}l@{}}\rule{0pt}{12pt}MPF / Bounded\\SIMD-DP / Block\end{tabular} & \hspace*{1pt}\begin{tabular}[t]{@{}c@{}}\rule{0pt}{12pt}9.9s\\[-.2pt]{\fontsize{5.8}{6.4}\selectfont \textcolor{ForestGreen}{$ \blacktriangle\,3.6\times $}}\\[-.2pt]9s\\[-.2pt]{\fontsize{5.8}{6.4}\selectfont \textcolor{ForestGreen}{$ \blacktriangle\,3.9\times $}}\end{tabular} & \hspace*{1pt}\begin{tabular}[t]{@{}c@{}}\rule{0pt}{12pt}13s\\[-.2pt]{\fontsize{5.8}{6.4}\selectfont \textcolor{ForestGreen}{$ \blacktriangle\,2.9\times $}}\\[-.2pt]11.9s\\[-.2pt]{\fontsize{5.8}{6.4}\selectfont \textcolor{ForestGreen}{$ \blacktriangle\,3.0\times $}}\end{tabular} & \hspace*{1pt}\begin{tabular}[t]{@{}c@{}}\rule{0pt}{12pt}9.4s\\[-.2pt]{\fontsize{5.8}{6.4}\selectfont \textcolor{ForestGreen}{$ \blacktriangle\,9.9\times $}}\\[-.2pt]7s\\[-.2pt]{\fontsize{5.8}{6.4}\selectfont \textcolor{ForestGreen}{$ \blacktriangle\,12.8\times $}}\end{tabular} & \textcolor{gray}{\ensuremath{\cdot}} & \textcolor{gray}{\ensuremath{\cdot}} & \textcolor{gray}{\ensuremath{\cdot}} & \textcolor{gray}{\ensuremath{\cdot}} & \textcolor{gray}{\ensuremath{\cdot}} & \hspace*{1pt}\begin{tabular}[t]{@{}c@{}}\rule{0pt}{12pt}3m\,44.5s\\[-.2pt]{\fontsize{5.8}{6.4}\selectfont \textcolor{BrickRed}{$ \blacktriangledown\,28\% $}}\\[-.2pt]3m\,35.9s\\[-.2pt]{\fontsize{5.8}{6.4}\selectfont \textcolor{BrickRed}{$ \blacktriangledown\,24\% $}}\end{tabular} & \hspace*{1pt}\begin{tabular}[t]{@{}c@{}}\rule{0pt}{12pt}3m\,57.4s\\[-.2pt]{\fontsize{5.8}{6.4}\selectfont \textcolor{BrickRed}{$ \blacktriangledown\,1.4\times $}}\\[-.2pt]3m\,45.7s\\[-.2pt]{\fontsize{5.8}{6.4}\selectfont \textcolor{BrickRed}{$ \blacktriangledown\,1.3\times $}}\end{tabular} & \hspace*{1pt}\begin{tabular}[t]{@{}c@{}}\rule{0pt}{12pt}6m\,50.6s\\[-.2pt]{\fontsize{5.8}{6.4}\selectfont \textcolor{BrickRed}{$ \blacktriangledown\,10\% $}}\\[-.2pt]6m\,30.7s\\[-.2pt]{\fontsize{5.8}{6.4}\selectfont \textcolor{BrickRed}{$ \blacktriangledown\,6\% $}}\end{tabular} & \textcolor{gray}{\ensuremath{\cdot}} & \hspace*{1pt}\begin{tabular}[t]{@{}c@{}}\rule{0pt}{12pt}9m\,45.5s\\[-.2pt]{\fontsize{5.8}{6.4}\selectfont \textcolor{BrickRed}{$ \blacktriangledown\,5.5\times $}}\\[-.2pt]3m\,41.5s\\[-.2pt]{\fontsize{5.8}{6.4}\selectfont \textcolor{BrickRed}{$ \blacktriangledown\,2.1\times $}}\end{tabular} & \hspace*{1pt}\begin{tabular}[t]{@{}c@{}}\rule{0pt}{12pt}8m\,25.6s\\[-.2pt]{\fontsize{5.8}{6.4}\selectfont \textcolor{BrickRed}{$ \blacktriangledown\,4.7\times $}}\\[-.2pt]3m\,56.7s\\[-.2pt]{\fontsize{5.8}{6.4}\selectfont \textcolor{BrickRed}{$ \blacktriangledown\,2.3\times $}}\end{tabular} & \hspace*{1pt}\begin{tabular}[t]{@{}c@{}}\rule{0pt}{12pt}15m\,24.9s\\[-.2pt]{\fontsize{5.8}{6.4}\selectfont \textcolor{BrickRed}{$ \blacktriangledown\,4.2\times $}}\\[-.2pt]6m\,0.2s\\[-.2pt]{\fontsize{5.8}{6.4}\selectfont \textcolor{BrickRed}{$ \blacktriangledown\,1.7\times $}}\end{tabular} & \textcolor{gray}{\ensuremath{\cdot}}\\[3pt]\hline
 & \begin{tabular}[t]{@{}l@{}}\rule{0pt}{12pt}MPF / Bounded\\SIMD-DP / KSW2\end{tabular} & \hspace*{1pt}\begin{tabular}[t]{@{}c@{}}\rule{0pt}{12pt}14.7s\\[-.2pt]{\fontsize{5.8}{6.4}\selectfont \textcolor{ForestGreen}{$ \blacktriangle\,2.5\times $}}\\[-.2pt]12.1s\\[-.2pt]{\fontsize{5.8}{6.4}\selectfont \textcolor{ForestGreen}{$ \blacktriangle\,2.9\times $}}\end{tabular} & \hspace*{1pt}\begin{tabular}[t]{@{}c@{}}\rule{0pt}{12pt}8.1s\\[-.2pt]{\fontsize{5.8}{6.4}\selectfont \textcolor{ForestGreen}{$ \blacktriangle\,4.7\times $}}\\[-.2pt]5.5s\\[-.2pt]{\fontsize{5.8}{6.4}\selectfont \textcolor{ForestGreen}{$ \blacktriangle\,6.5\times $}}\end{tabular} & \textcolor{gray}{\ensuremath{\cdot}} & \textcolor{gray}{\ensuremath{\cdot}} & \textcolor{gray}{\ensuremath{\cdot}} & \textcolor{gray}{\ensuremath{\cdot}} & \textcolor{gray}{\ensuremath{\cdot}} & \textcolor{gray}{\ensuremath{\cdot}} & \hspace*{1pt}\begin{tabular}[t]{@{}c@{}}\rule{0pt}{12pt}3m\,57.7s\\[-.2pt]{\fontsize{5.8}{6.4}\selectfont \textcolor{BrickRed}{$ \blacktriangledown\,1.4\times $}}\\[-.2pt]3m\,46.9s\\[-.2pt]{\fontsize{5.8}{6.4}\selectfont \textcolor{BrickRed}{$ \blacktriangledown\,1.3\times $}}\end{tabular} & \hspace*{1pt}\begin{tabular}[t]{@{}c@{}}\rule{0pt}{12pt}4m\,0.1s\\[-.2pt]{\fontsize{5.8}{6.4}\selectfont \textcolor{BrickRed}{$ \blacktriangledown\,1.4\times $}}\\[-.2pt]3m\,46s\\[-.2pt]{\fontsize{5.8}{6.4}\selectfont \textcolor{BrickRed}{$ \blacktriangledown\,1.3\times $}}\end{tabular} & \textcolor{gray}{\ensuremath{\cdot}} & \textcolor{gray}{\ensuremath{\cdot}} & \textcolor{gray}{\ensuremath{\cdot}} & \textcolor{gray}{\ensuremath{\cdot}} & \textcolor{gray}{\ensuremath{\cdot}} & \textcolor{gray}{\ensuremath{\cdot}}\\[3pt]\hline
 & \begin{tabular}[t]{@{}l@{}}\rule{0pt}{12pt}MPF / Scalar\\RMQ / Block\end{tabular} & \hspace*{1pt}\begin{tabular}[t]{@{}c@{}}\rule{0pt}{12pt}5.5s\\[-.2pt]{\fontsize{5.8}{6.4}\selectfont \textcolor{ForestGreen}{$ \blacktriangle\,6.6\times $}}\\[-.2pt]4.6s\\[-.2pt]{\fontsize{5.8}{6.4}\selectfont \textcolor{ForestGreen}{$ \blacktriangle\,7.6\times $}}\end{tabular} & \hspace*{1pt}\begin{tabular}[t]{@{}c@{}}\rule{0pt}{12pt}5s\\[-.2pt]{\fontsize{5.8}{6.4}\selectfont \textcolor{ForestGreen}{$ \blacktriangle\,7.6\times $}}\\[-.2pt]4.1s\\[-.2pt]{\fontsize{5.8}{6.4}\selectfont \textcolor{ForestGreen}{$ \blacktriangle\,8.8\times $}}\end{tabular} & \textcolor{gray}{\ensuremath{\cdot}} & \textcolor{gray}{\ensuremath{\cdot}} & \textcolor{gray}{\ensuremath{\cdot}} & \textcolor{gray}{\ensuremath{\cdot}} & \textcolor{gray}{\ensuremath{\cdot}} & \textcolor{gray}{\ensuremath{\cdot}} & \hspace*{1pt}\begin{tabular}[t]{@{}c@{}}\rule{0pt}{12pt}5m\,26.3s\\[-.2pt]{\fontsize{5.8}{6.4}\selectfont \textcolor{BrickRed}{$ \blacktriangledown\,1.9\times $}}\\[-.2pt]5m\,17.1s\\[-.2pt]{\fontsize{5.8}{6.4}\selectfont \textcolor{BrickRed}{$ \blacktriangledown\,1.8\times $}}\end{tabular} & \hspace*{1pt}\begin{tabular}[t]{@{}c@{}}\rule{0pt}{12pt}5m\,17.7s\\[-.2pt]{\fontsize{5.8}{6.4}\selectfont \textcolor{BrickRed}{$ \blacktriangledown\,1.9\times $}}\\[-.2pt]5m\,5.9s\\[-.2pt]{\fontsize{5.8}{6.4}\selectfont \textcolor{BrickRed}{$ \blacktriangledown\,1.8\times $}}\end{tabular} & \textcolor{gray}{\ensuremath{\cdot}} & \textcolor{gray}{\ensuremath{\cdot}} & \textcolor{gray}{\ensuremath{\cdot}} & \textcolor{gray}{\ensuremath{\cdot}} & \textcolor{gray}{\ensuremath{\cdot}} & \textcolor{gray}{\ensuremath{\cdot}}\\[3pt]\hline
\end{tabular}\par}
\smallskip{\footnotesize\noindent Each cell: A, comparison, B, comparison. Times cover 1,024 queries. $\cdot$ = no approved observation. ---\textsuperscript{g} = GPU unavailable, not executed. n/a = no measured comparison.\par
\noindent \textsuperscript{f}SIMD-ref request observed as scalar, compared only as that effective path. \textsuperscript{p}Accepted by explicit preservation after wrapper postflight failure, no run receipt claimed.\par
\noindent MPF/MPS = minimizers per fragment/sequence. Row order: index representation / minimizer, then chaining / base alignment. SIMD-ref is the reference minimizer, Bounded is bounded ntHash.\par
\noindent Blue = one lowest A and one lowest B down each machine/workload column, across continuation pages. Exact ties select the first row. Repeated baselines are marked only on page 1.\par
\noindent Light gray = minimap2 baseline. Pale yellow = preferred MPF / Bounded / SIMD-DP / Block configuration, not a claim of winning every column.\par
\noindent When all Vulkan modes are unmeasured, they share one row. Baseline comparisons use the fastest measured BioseqDB effective path separately for A and B.\par}
\clearpage
\noindent{\large\bfseries AWS: 4 Gbp, k=15, w=10 (mapping interval)\hfill 2/6}\par\smallskip
{\fontsize{7.1}{8.1}\selectfont\setlength{\tabcolsep}{1pt}\setlength{\fboxsep}{2pt}\renewcommand{\arraystretch}{1.12}
\noindent\begin{tabular}{@{}p{14mm}p{27mm}|p{13.55mm}|p{13.55mm}|p{13.55mm}|p{13.55mm}|p{13.55mm}|p{13.55mm}|p{13.55mm}|p{13.55mm}|p{13.55mm}|p{13.55mm}|p{13.55mm}|p{13.55mm}|p{13.55mm}|p{13.55mm}|p{13.55mm}|p{13.55mm}|@{}}\hline
Engine & Execution & \multicolumn{8}{c|}{\textbf{Index}} & \multicolumn{8}{c|}{\textbf{No index}}\\
 &  & \multicolumn{8}{c|}{\textbf{4 Gbp}} & \multicolumn{8}{c|}{\textbf{4 Gbp}}\\
 &  & \multicolumn{8}{c|}{\textbf{k=15, w=10}} & \multicolumn{8}{c|}{\textbf{k=15, w=10}}\\
 &  & \multicolumn{4}{c|}{\textbf{DNA4}} & \multicolumn{4}{c|}{\textbf{DNA5}} & \multicolumn{4}{c|}{\textbf{DNA4}} & \multicolumn{4}{c|}{\textbf{DNA5}}\\
 &  & \multicolumn{1}{c|}{\textbf{M7A}} & \multicolumn{1}{c|}{\textbf{M7I}} & \multicolumn{1}{c|}{\textbf{M4}} & \multicolumn{1}{c|}{\textbf{C7G}} & \multicolumn{1}{c|}{\textbf{M7A}} & \multicolumn{1}{c|}{\textbf{M7I}} & \multicolumn{1}{c|}{\textbf{M4}} & \multicolumn{1}{c|}{\textbf{C7G}} & \multicolumn{1}{c|}{\textbf{M7A}} & \multicolumn{1}{c|}{\textbf{M7I}} & \multicolumn{1}{c|}{\textbf{M4}} & \multicolumn{1}{c|}{\textbf{C7G}} & \multicolumn{1}{c|}{\textbf{M7A}} & \multicolumn{1}{c|}{\textbf{M7I}} & \multicolumn{1}{c|}{\textbf{M4}} & \multicolumn{1}{c|}{\textbf{C7G}}\\
\hline
\rowcolor{BenchBaseline}[0.5pt][0.5pt]\textbf{minimap2} & \textbf{Baseline} & \hspace*{1pt}\begin{tabular}[t]{@{}c@{}}\rule{0pt}{12pt}36.2s\\[-.2pt]{\fontsize{5.8}{6.4}\selectfont \textcolor{BrickRed}{$ \blacktriangledown\,6.9\times $}}\\[-.2pt]35s\\[-.2pt]{\fontsize{5.8}{6.4}\selectfont \textcolor{BrickRed}{$ \blacktriangledown\,8.0\times $}}\end{tabular} & \hspace*{1pt}\begin{tabular}[t]{@{}c@{}}\rule{0pt}{12pt}38.1s\\[-.2pt]{\fontsize{5.8}{6.4}\selectfont \textcolor{BrickRed}{$ \blacktriangledown\,7.8\times $}}\\[-.2pt]36s\\[-.2pt]{\fontsize{5.8}{6.4}\selectfont \textcolor{BrickRed}{$ \blacktriangledown\,9.3\times $}}\end{tabular} & \hspace*{1pt}\begin{tabular}[t]{@{}c@{}}\rule{0pt}{12pt}1m\,33.5s\\[-.2pt]{\fontsize{5.8}{6.4}\selectfont \textcolor{BrickRed}{$ \blacktriangledown\,10.8\times $}}\\[-.2pt]1m\,29.7s\\[-.2pt]{\fontsize{5.8}{6.4}\selectfont \textcolor{BrickRed}{$ \blacktriangledown\,13.7\times $}}\end{tabular} & \hspace*{1pt}\begin{tabular}[t]{@{}c@{}}\rule{0pt}{12pt}6m\,31.5s\\[-.2pt]{\fontsize{5.8}{6.4}\selectfont \textcolor{gray}{n/a}}\\[-.2pt]6m\,26.3s\\[-.2pt]{\fontsize{5.8}{6.4}\selectfont \textcolor{gray}{n/a}}\end{tabular} & \hspace*{1pt}\begin{tabular}[t]{@{}c@{}}\rule{0pt}{12pt}31.6s\\[-.2pt]{\fontsize{5.8}{6.4}\selectfont \textcolor{gray}{n/a}}\\[-.2pt]30.3s\\[-.2pt]{\fontsize{5.8}{6.4}\selectfont \textcolor{gray}{n/a}}\end{tabular} & \hspace*{1pt}\begin{tabular}[t]{@{}c@{}}\rule{0pt}{12pt}34.1s\\[-.2pt]{\fontsize{5.8}{6.4}\selectfont \textcolor{gray}{n/a}}\\[-.2pt]32s\\[-.2pt]{\fontsize{5.8}{6.4}\selectfont \textcolor{gray}{n/a}}\end{tabular} & \hspace*{1pt}\begin{tabular}[t]{@{}c@{}}\rule{0pt}{12pt}1m\,9.5s\\[-.2pt]{\fontsize{5.8}{6.4}\selectfont \textcolor{gray}{n/a}}\\[-.2pt]1m\,6.1s\\[-.2pt]{\fontsize{5.8}{6.4}\selectfont \textcolor{gray}{n/a}}\end{tabular} & \hspace*{1pt}\begin{tabular}[t]{@{}c@{}}\rule{0pt}{12pt}3m\,1s\\[-.2pt]{\fontsize{5.8}{6.4}\selectfont \textcolor{gray}{n/a}}\\[-.2pt]2m\,51.5s\\[-.2pt]{\fontsize{5.8}{6.4}\selectfont \textcolor{gray}{n/a}}\end{tabular} & \hspace*{1pt}\begin{tabular}[t]{@{}c@{}}\rule{0pt}{12pt}2m\,55.5s\\[-.2pt]{\fontsize{5.8}{6.4}\selectfont \textcolor{ForestGreen}{$ \blacktriangle\,20\% $}}\\[-.2pt]2m\,54s\\[-.2pt]{\fontsize{5.8}{6.4}\selectfont \textcolor{ForestGreen}{$ \blacktriangle\,17\% $}}\end{tabular} & \hspace*{1pt}\begin{tabular}[t]{@{}c@{}}\rule{0pt}{12pt}2m\,50.4s\\[-.2pt]{\fontsize{5.8}{6.4}\selectfont \textcolor{ForestGreen}{$ \blacktriangle\,23\% $}}\\[-.2pt]2m\,48.2s\\[-.2pt]{\fontsize{5.8}{6.4}\selectfont \textcolor{ForestGreen}{$ \blacktriangle\,20\% $}}\end{tabular} & \hspace*{1pt}\begin{tabular}[t]{@{}c@{}}\rule{0pt}{12pt}6m\,12.2s\\[-.2pt]{\fontsize{5.8}{6.4}\selectfont \textcolor{ForestGreen}{$ \blacktriangle\,9\% $}}\\[-.2pt]6m\,8.2s\\[-.2pt]{\fontsize{5.8}{6.4}\selectfont \textcolor{gray}{$ \blacktriangleright\,5\% $}}\end{tabular} & \hspace*{1pt}\begin{tabular}[t]{@{}c@{}}\rule{0pt}{12pt}44m\,20.5s\\[-.2pt]{\fontsize{5.8}{6.4}\selectfont \textcolor{gray}{n/a}}\\[-.2pt]44m\,15.6s\\[-.2pt]{\fontsize{5.8}{6.4}\selectfont \textcolor{gray}{n/a}}\end{tabular} & \hspace*{1pt}\begin{tabular}[t]{@{}c@{}}\rule{0pt}{12pt}1m\,47.1s\\[-.2pt]{\fontsize{5.8}{6.4}\selectfont \textcolor{ForestGreen}{$ \blacktriangle\,5.5\times $}}\\[-.2pt]1m\,45.8s\\[-.2pt]{\fontsize{5.8}{6.4}\selectfont \textcolor{ForestGreen}{$ \blacktriangle\,2.1\times $}}\end{tabular} & \hspace*{1pt}\begin{tabular}[t]{@{}c@{}}\rule{0pt}{12pt}1m\,47.2s\\[-.2pt]{\fontsize{5.8}{6.4}\selectfont \textcolor{ForestGreen}{$ \blacktriangle\,4.6\times $}}\\[-.2pt]1m\,45.1s\\[-.2pt]{\fontsize{5.8}{6.4}\selectfont \textcolor{ForestGreen}{$ \blacktriangle\,2.2\times $}}\end{tabular} & \hspace*{1pt}\begin{tabular}[t]{@{}c@{}}\rule{0pt}{12pt}3m\,40.3s\\[-.2pt]{\fontsize{5.8}{6.4}\selectfont \textcolor{ForestGreen}{$ \blacktriangle\,4.2\times $}}\\[-.2pt]3m\,36.6s\\[-.2pt]{\fontsize{5.8}{6.4}\selectfont \textcolor{ForestGreen}{$ \blacktriangle\,1.7\times $}}\end{tabular} & \hspace*{1pt}\begin{tabular}[t]{@{}c@{}}\rule{0pt}{12pt}29m\,43.4s\\[-.2pt]{\fontsize{5.8}{6.4}\selectfont \textcolor{gray}{n/a}}\\[-.2pt]29m\,38.5s\\[-.2pt]{\fontsize{5.8}{6.4}\selectfont \textcolor{gray}{n/a}}\end{tabular}\\[3pt]\hline
\textbf{BioseqDB} & \begin{tabular}[t]{@{}l@{}}\rule{0pt}{12pt}MPF / Scalar\\RMQ / KSW2\end{tabular} & \hspace*{1pt}\begin{tabular}[t]{@{}c@{}}\rule{0pt}{12pt}7s\\[-.2pt]{\fontsize{5.8}{6.4}\selectfont \textcolor{ForestGreen}{$ \blacktriangle\,5.1\times $}}\\[-.2pt]4.5s\\[-.2pt]{\fontsize{5.8}{6.4}\selectfont \textcolor{ForestGreen}{$ \blacktriangle\,7.7\times $}}\end{tabular} & \hspace*{1pt}\begin{tabular}[t]{@{}c@{}}\rule{0pt}{12pt}6.5s\\[-.2pt]{\fontsize{5.8}{6.4}\selectfont \textcolor{ForestGreen}{$ \blacktriangle\,5.9\times $}}\\[-.2pt]3.9s\\[-.2pt]{\fontsize{5.8}{6.4}\selectfont \textcolor{ForestGreen}{$ \blacktriangle\,9.1\times $}}\end{tabular} & \textcolor{gray}{\ensuremath{\cdot}} & \textcolor{gray}{\ensuremath{\cdot}} & \textcolor{gray}{\ensuremath{\cdot}} & \textcolor{gray}{\ensuremath{\cdot}} & \textcolor{gray}{\ensuremath{\cdot}} & \textcolor{gray}{\ensuremath{\cdot}} & \textcolor{gray}{\ensuremath{\cdot}} & \textcolor{gray}{\ensuremath{\cdot}} & \textcolor{gray}{\ensuremath{\cdot}} & \textcolor{gray}{\ensuremath{\cdot}} & \hspace*{1pt}\begin{tabular}[t]{@{}c@{}}\rule{0pt}{12pt}11m\,18.9s\\[-.2pt]{\fontsize{5.8}{6.4}\selectfont \textcolor{BrickRed}{$ \blacktriangledown\,6.3\times $}}\\[-.2pt]5m\,12.1s\\[-.2pt]{\fontsize{5.8}{6.4}\selectfont \textcolor{BrickRed}{$ \blacktriangledown\,2.9\times $}}\end{tabular} & \hspace*{1pt}\begin{tabular}[t]{@{}c@{}}\rule{0pt}{12pt}9m\,52.6s\\[-.2pt]{\fontsize{5.8}{6.4}\selectfont \textcolor{BrickRed}{$ \blacktriangledown\,5.5\times $}}\\[-.2pt]5m\,20s\\[-.2pt]{\fontsize{5.8}{6.4}\selectfont \textcolor{BrickRed}{$ \blacktriangledown\,3.0\times $}}\end{tabular} & \textcolor{gray}{\ensuremath{\cdot}} & \textcolor{gray}{\ensuremath{\cdot}}\\[3pt]\hline
 & \begin{tabular}[t]{@{}l@{}}\rule{0pt}{12pt}MPF / Scalar\\Scalar-DP / Block\end{tabular} & \hspace*{1pt}\begin{tabular}[t]{@{}c@{}}\rule{0pt}{12pt}5.4s\\[-.2pt]{\fontsize{5.8}{6.4}\selectfont \textcolor{ForestGreen}{$ \blacktriangle\,6.7\times $}}\\[-.2pt]4.6s\\[-.2pt]{\fontsize{5.8}{6.4}\selectfont \textcolor{ForestGreen}{$ \blacktriangle\,7.6\times $}}\end{tabular} & \hspace*{1pt}\begin{tabular}[t]{@{}c@{}}\rule{0pt}{12pt}4.9s\\[-.2pt]{\fontsize{5.8}{6.4}\selectfont \textcolor{ForestGreen}{$ \blacktriangle\,7.7\times $}}\\[-.2pt]4s\\[-.2pt]{\fontsize{5.8}{6.4}\selectfont \textcolor{ForestGreen}{$ \blacktriangle\,8.9\times $}}\end{tabular} & \textcolor{gray}{\ensuremath{\cdot}} & \textcolor{gray}{\ensuremath{\cdot}} & \textcolor{gray}{\ensuremath{\cdot}} & \textcolor{gray}{\ensuremath{\cdot}} & \textcolor{gray}{\ensuremath{\cdot}} & \textcolor{gray}{\ensuremath{\cdot}} & \textcolor{gray}{\ensuremath{\cdot}} & \textcolor{gray}{\ensuremath{\cdot}} & \textcolor{gray}{\ensuremath{\cdot}} & \textcolor{gray}{\ensuremath{\cdot}} & \hspace*{1pt}\begin{tabular}[t]{@{}c@{}}\rule{0pt}{12pt}11m\,9.4s\\[-.2pt]{\fontsize{5.8}{6.4}\selectfont \textcolor{BrickRed}{$ \blacktriangledown\,6.3\times $}}\\[-.2pt]5m\,3.9s\\[-.2pt]{\fontsize{5.8}{6.4}\selectfont \textcolor{BrickRed}{$ \blacktriangledown\,2.9\times $}}\end{tabular} & \hspace*{1pt}\begin{tabular}[t]{@{}c@{}}\rule{0pt}{12pt}9m\,52.5s\\[-.2pt]{\fontsize{5.8}{6.4}\selectfont \textcolor{BrickRed}{$ \blacktriangledown\,5.5\times $}}\\[-.2pt]5m\,21.2s\\[-.2pt]{\fontsize{5.8}{6.4}\selectfont \textcolor{BrickRed}{$ \blacktriangledown\,3.1\times $}}\end{tabular} & \textcolor{gray}{\ensuremath{\cdot}} & \textcolor{gray}{\ensuremath{\cdot}}\\[3pt]\hline
 & \begin{tabular}[t]{@{}l@{}}\rule{0pt}{12pt}MPF / Scalar\\Scalar-DP / KSW2\end{tabular} & \hspace*{1pt}\begin{tabular}[t]{@{}c@{}}\rule{0pt}{12pt}15s\\[-.2pt]{\fontsize{5.8}{6.4}\selectfont \textcolor{ForestGreen}{$ \blacktriangle\,2.4\times $}}\\[-.2pt]12.5s\\[-.2pt]{\fontsize{5.8}{6.4}\selectfont \textcolor{ForestGreen}{$ \blacktriangle\,2.8\times $}}\end{tabular} & \hspace*{1pt}\begin{tabular}[t]{@{}c@{}}\rule{0pt}{12pt}10.6s\\[-.2pt]{\fontsize{5.8}{6.4}\selectfont \textcolor{ForestGreen}{$ \blacktriangle\,3.6\times $}}\\[-.2pt]8.2s\\[-.2pt]{\fontsize{5.8}{6.4}\selectfont \textcolor{ForestGreen}{$ \blacktriangle\,4.4\times $}}\end{tabular} & \textcolor{gray}{\ensuremath{\cdot}} & \textcolor{gray}{\ensuremath{\cdot}} & \textcolor{gray}{\ensuremath{\cdot}} & \textcolor{gray}{\ensuremath{\cdot}} & \textcolor{gray}{\ensuremath{\cdot}} & \textcolor{gray}{\ensuremath{\cdot}} & \textcolor{gray}{\ensuremath{\cdot}} & \textcolor{gray}{\ensuremath{\cdot}} & \textcolor{gray}{\ensuremath{\cdot}} & \textcolor{gray}{\ensuremath{\cdot}} & \hspace*{1pt}\begin{tabular}[t]{@{}c@{}}\rule{0pt}{12pt}11m\,14.1s\\[-.2pt]{\fontsize{5.8}{6.4}\selectfont \textcolor{BrickRed}{$ \blacktriangledown\,6.3\times $}}\\[-.2pt]5m\,5.9s\\[-.2pt]{\fontsize{5.8}{6.4}\selectfont \textcolor{BrickRed}{$ \blacktriangledown\,2.9\times $}}\end{tabular} & \hspace*{1pt}\begin{tabular}[t]{@{}c@{}}\rule{0pt}{12pt}9m\,54.4s\\[-.2pt]{\fontsize{5.8}{6.4}\selectfont \textcolor{BrickRed}{$ \blacktriangledown\,5.5\times $}}\\[-.2pt]5m\,20.4s\\[-.2pt]{\fontsize{5.8}{6.4}\selectfont \textcolor{BrickRed}{$ \blacktriangledown\,3.0\times $}}\end{tabular} & \textcolor{gray}{\ensuremath{\cdot}} & \textcolor{gray}{\ensuremath{\cdot}}\\[3pt]\hline
 & \begin{tabular}[t]{@{}l@{}}\rule{0pt}{12pt}MPF / Scalar\\SIMD-DP / Block\end{tabular} & \hspace*{1pt}\begin{tabular}[t]{@{}c@{}}\rule{0pt}{12pt}5.5s\\[-.2pt]{\fontsize{5.8}{6.4}\selectfont \textcolor{ForestGreen}{$ \blacktriangle\,6.6\times $}}\\[-.2pt]4.6s\\[-.2pt]{\fontsize{5.8}{6.4}\selectfont \textcolor{ForestGreen}{$ \blacktriangle\,7.6\times $}}\end{tabular} & \hspace*{1pt}\begin{tabular}[t]{@{}c@{}}\rule{0pt}{12pt}5.1s\\[-.2pt]{\fontsize{5.8}{6.4}\selectfont \textcolor{ForestGreen}{$ \blacktriangle\,7.5\times $}}\\[-.2pt]4.1s\\[-.2pt]{\fontsize{5.8}{6.4}\selectfont \textcolor{ForestGreen}{$ \blacktriangle\,8.8\times $}}\end{tabular} & \textcolor{gray}{\ensuremath{\cdot}} & \textcolor{gray}{\ensuremath{\cdot}} & \textcolor{gray}{\ensuremath{\cdot}} & \textcolor{gray}{\ensuremath{\cdot}} & \textcolor{gray}{\ensuremath{\cdot}} & \textcolor{gray}{\ensuremath{\cdot}} & \hspace*{1pt}\begin{tabular}[t]{@{}c@{}}\rule{0pt}{12pt}5m\,24.9s\\[-.2pt]{\fontsize{5.8}{6.4}\selectfont \textcolor{BrickRed}{$ \blacktriangledown\,1.9\times $}}\\[-.2pt]5m\,16s\\[-.2pt]{\fontsize{5.8}{6.4}\selectfont \textcolor{BrickRed}{$ \blacktriangledown\,1.8\times $}}\end{tabular} & \textcolor{gray}{\ensuremath{\cdot}} & \textcolor{gray}{\ensuremath{\cdot}} & \textcolor{gray}{\ensuremath{\cdot}} & \textcolor{gray}{\ensuremath{\cdot}} & \hspace*{1pt}\begin{tabular}[t]{@{}c@{}}\rule{0pt}{12pt}9m\,48.5s\\[-.2pt]{\fontsize{5.8}{6.4}\selectfont \textcolor{BrickRed}{$ \blacktriangledown\,5.5\times $}}\\[-.2pt]5m\,17.7s\\[-.2pt]{\fontsize{5.8}{6.4}\selectfont \textcolor{BrickRed}{$ \blacktriangledown\,3.0\times $}}\end{tabular} & \textcolor{gray}{\ensuremath{\cdot}} & \textcolor{gray}{\ensuremath{\cdot}}\\[3pt]\hline
 & \begin{tabular}[t]{@{}l@{}}\rule{0pt}{12pt}MPF / Scalar\\SIMD-DP / KSW2\end{tabular} & \hspace*{1pt}\begin{tabular}[t]{@{}c@{}}\rule{0pt}{12pt}8.2s\\[-.2pt]{\fontsize{5.8}{6.4}\selectfont \textcolor{ForestGreen}{$ \blacktriangle\,4.4\times $}}\\[-.2pt]5.6s\\[-.2pt]{\fontsize{5.8}{6.4}\selectfont \textcolor{ForestGreen}{$ \blacktriangle\,6.2\times $}}\end{tabular} & \hspace*{1pt}\begin{tabular}[t]{@{}c@{}}\rule{0pt}{12pt}6.5s\\[-.2pt]{\fontsize{5.8}{6.4}\selectfont \textcolor{ForestGreen}{$ \blacktriangle\,5.9\times $}}\\[-.2pt]3.9s\\[-.2pt]{\fontsize{5.8}{6.4}\selectfont \textcolor{ForestGreen}{$ \blacktriangle\,9.1\times $}}\end{tabular} & \textcolor{gray}{\ensuremath{\cdot}} & \textcolor{gray}{\ensuremath{\cdot}} & \textcolor{gray}{\ensuremath{\cdot}} & \textcolor{gray}{\ensuremath{\cdot}} & \textcolor{gray}{\ensuremath{\cdot}} & \textcolor{gray}{\ensuremath{\cdot}} & \hspace*{1pt}\begin{tabular}[t]{@{}c@{}}\rule{0pt}{12pt}5m\,27.2s\\[-.2pt]{\fontsize{5.8}{6.4}\selectfont \textcolor{BrickRed}{$ \blacktriangledown\,1.9\times $}}\\[-.2pt]5m\,16.6s\\[-.2pt]{\fontsize{5.8}{6.4}\selectfont \textcolor{BrickRed}{$ \blacktriangledown\,1.8\times $}}\end{tabular} & \hspace*{1pt}\begin{tabular}[t]{@{}c@{}}\rule{0pt}{12pt}5m\,50s\\[-.2pt]{\fontsize{5.8}{6.4}\selectfont \textcolor{BrickRed}{$ \blacktriangledown\,2.1\times $}}\\[-.2pt]5m\,36.1s\\[-.2pt]{\fontsize{5.8}{6.4}\selectfont \textcolor{BrickRed}{$ \blacktriangledown\,2.0\times $}}\end{tabular} & \textcolor{gray}{\ensuremath{\cdot}} & \textcolor{gray}{\ensuremath{\cdot}} & \hspace*{1pt}\begin{tabular}[t]{@{}c@{}}\rule{0pt}{12pt}11m\,14.3s\\[-.2pt]{\fontsize{5.8}{6.4}\selectfont \textcolor{BrickRed}{$ \blacktriangledown\,6.3\times $}}\\[-.2pt]5m\,8.1s\\[-.2pt]{\fontsize{5.8}{6.4}\selectfont \textcolor{BrickRed}{$ \blacktriangledown\,2.9\times $}}\end{tabular} & \hspace*{1pt}\begin{tabular}[t]{@{}c@{}}\rule{0pt}{12pt}9m\,48.4s\\[-.2pt]{\fontsize{5.8}{6.4}\selectfont \textcolor{BrickRed}{$ \blacktriangledown\,5.5\times $}}\\[-.2pt]5m\,16.4s\\[-.2pt]{\fontsize{5.8}{6.4}\selectfont \textcolor{BrickRed}{$ \blacktriangledown\,3.0\times $}}\end{tabular} & \textcolor{gray}{\ensuremath{\cdot}} & \textcolor{gray}{\ensuremath{\cdot}}\\[3pt]\hline
 & \begin{tabular}[t]{@{}l@{}}\rule{0pt}{12pt}MPF / SIMD-ref\\RMQ / Block\end{tabular} & \hspace*{1pt}\begin{tabular}[t]{@{}c@{}}\rule{0pt}{12pt}5.4s\\[-.2pt]{\fontsize{5.8}{6.4}\selectfont \textcolor{ForestGreen}{$ \blacktriangle\,6.7\times $}}\\[-.2pt]4.5s\\[-.2pt]{\fontsize{5.8}{6.4}\selectfont \textcolor{ForestGreen}{$ \blacktriangle\,7.8\times $}}\end{tabular} & \hspace*{1pt}\begin{tabular}[t]{@{}c@{}}\rule{0pt}{12pt}5s\\[-.2pt]{\fontsize{5.8}{6.4}\selectfont \textcolor{ForestGreen}{$ \blacktriangle\,7.7\times $}}\\[-.2pt]4.1s\\[-.2pt]{\fontsize{5.8}{6.4}\selectfont \textcolor{ForestGreen}{$ \blacktriangle\,8.8\times $}}\end{tabular} & \textcolor{gray}{\ensuremath{\cdot}} & \textcolor{gray}{\ensuremath{\cdot}} & \textcolor{gray}{\ensuremath{\cdot}} & \textcolor{gray}{\ensuremath{\cdot}} & \textcolor{gray}{\ensuremath{\cdot}} & \textcolor{gray}{\ensuremath{\cdot}} & \textcolor{gray}{\ensuremath{\cdot}} & \textcolor{gray}{\ensuremath{\cdot}} & \textcolor{gray}{\ensuremath{\cdot}} & \textcolor{gray}{\ensuremath{\cdot}} & \textcolor{gray}{\ensuremath{\cdot}} & \textcolor{gray}{\ensuremath{\cdot}} & \textcolor{gray}{\ensuremath{\cdot}} & \textcolor{gray}{\ensuremath{\cdot}}\\[3pt]\hline
 & \begin{tabular}[t]{@{}l@{}}\rule{0pt}{12pt}MPF / SIMD-ref\\RMQ / KSW2\end{tabular} & \hspace*{1pt}\begin{tabular}[t]{@{}c@{}}\rule{0pt}{12pt}6.9s\\[-.2pt]{\fontsize{5.8}{6.4}\selectfont \textcolor{ForestGreen}{$ \blacktriangle\,5.3\times $}}\\[-.2pt]4.4s\\[-.2pt]{\fontsize{5.8}{6.4}\selectfont \textcolor{ForestGreen}{$ \blacktriangle\,8.0\times $}}\end{tabular} & \hspace*{1pt}\begin{tabular}[t]{@{}c@{}}\rule{0pt}{12pt}6.5s\\[-.2pt]{\fontsize{5.8}{6.4}\selectfont \textcolor{ForestGreen}{$ \blacktriangle\,5.9\times $}}\\[-.2pt]3.9s\\[-.2pt]{\fontsize{5.8}{6.4}\selectfont \textcolor{ForestGreen}{$ \blacktriangle\,9.3\times $}}\end{tabular} & \textcolor{gray}{\ensuremath{\cdot}} & \textcolor{gray}{\ensuremath{\cdot}} & \textcolor{gray}{\ensuremath{\cdot}} & \textcolor{gray}{\ensuremath{\cdot}} & \textcolor{gray}{\ensuremath{\cdot}} & \textcolor{gray}{\ensuremath{\cdot}} & \textcolor{gray}{\ensuremath{\cdot}} & \textcolor{gray}{\ensuremath{\cdot}} & \textcolor{gray}{\ensuremath{\cdot}} & \textcolor{gray}{\ensuremath{\cdot}} & \textcolor{gray}{\ensuremath{\cdot}} & \textcolor{gray}{\ensuremath{\cdot}} & \textcolor{gray}{\ensuremath{\cdot}} & \textcolor{gray}{\ensuremath{\cdot}}\\[3pt]\hline
\end{tabular}\par}
\smallskip{\footnotesize\noindent Each cell: A, comparison, B, comparison. Times cover 1,024 queries. $\cdot$ = no approved observation. ---\textsuperscript{g} = GPU unavailable, not executed. n/a = no measured comparison.\par
\noindent \textsuperscript{f}SIMD-ref request observed as scalar, compared only as that effective path. \textsuperscript{p}Accepted by explicit preservation after wrapper postflight failure, no run receipt claimed.\par
\noindent MPF/MPS = minimizers per fragment/sequence. Row order: index representation / minimizer, then chaining / base alignment. SIMD-ref is the reference minimizer, Bounded is bounded ntHash.\par
\noindent Blue = one lowest A and one lowest B down each machine/workload column, across continuation pages. Exact ties select the first row. Repeated baselines are marked only on page 1.\par
\noindent Light gray = minimap2 baseline. Pale yellow = preferred MPF / Bounded / SIMD-DP / Block configuration, not a claim of winning every column.\par
\noindent When all Vulkan modes are unmeasured, they share one row. Baseline comparisons use the fastest measured BioseqDB effective path separately for A and B.\par}
\clearpage
\noindent{\large\bfseries AWS: 4 Gbp, k=15, w=10 (mapping interval)\hfill 3/6}\par\smallskip
{\fontsize{7.1}{8.1}\selectfont\setlength{\tabcolsep}{1pt}\setlength{\fboxsep}{2pt}\renewcommand{\arraystretch}{1.12}
\noindent\begin{tabular}{@{}p{14mm}p{27mm}|p{13.55mm}|p{13.55mm}|p{13.55mm}|p{13.55mm}|p{13.55mm}|p{13.55mm}|p{13.55mm}|p{13.55mm}|p{13.55mm}|p{13.55mm}|p{13.55mm}|p{13.55mm}|p{13.55mm}|p{13.55mm}|p{13.55mm}|p{13.55mm}|@{}}\hline
Engine & Execution & \multicolumn{8}{c|}{\textbf{Index}} & \multicolumn{8}{c|}{\textbf{No index}}\\
 &  & \multicolumn{8}{c|}{\textbf{4 Gbp}} & \multicolumn{8}{c|}{\textbf{4 Gbp}}\\
 &  & \multicolumn{8}{c|}{\textbf{k=15, w=10}} & \multicolumn{8}{c|}{\textbf{k=15, w=10}}\\
 &  & \multicolumn{4}{c|}{\textbf{DNA4}} & \multicolumn{4}{c|}{\textbf{DNA5}} & \multicolumn{4}{c|}{\textbf{DNA4}} & \multicolumn{4}{c|}{\textbf{DNA5}}\\
 &  & \multicolumn{1}{c|}{\textbf{M7A}} & \multicolumn{1}{c|}{\textbf{M7I}} & \multicolumn{1}{c|}{\textbf{M4}} & \multicolumn{1}{c|}{\textbf{C7G}} & \multicolumn{1}{c|}{\textbf{M7A}} & \multicolumn{1}{c|}{\textbf{M7I}} & \multicolumn{1}{c|}{\textbf{M4}} & \multicolumn{1}{c|}{\textbf{C7G}} & \multicolumn{1}{c|}{\textbf{M7A}} & \multicolumn{1}{c|}{\textbf{M7I}} & \multicolumn{1}{c|}{\textbf{M4}} & \multicolumn{1}{c|}{\textbf{C7G}} & \multicolumn{1}{c|}{\textbf{M7A}} & \multicolumn{1}{c|}{\textbf{M7I}} & \multicolumn{1}{c|}{\textbf{M4}} & \multicolumn{1}{c|}{\textbf{C7G}}\\
\hline
\rowcolor{BenchBaseline}[0.5pt][0.5pt]\textbf{minimap2} & \textbf{Baseline} & \hspace*{1pt}\begin{tabular}[t]{@{}c@{}}\rule{0pt}{12pt}36.2s\\[-.2pt]{\fontsize{5.8}{6.4}\selectfont \textcolor{BrickRed}{$ \blacktriangledown\,6.9\times $}}\\[-.2pt]35s\\[-.2pt]{\fontsize{5.8}{6.4}\selectfont \textcolor{BrickRed}{$ \blacktriangledown\,8.0\times $}}\end{tabular} & \hspace*{1pt}\begin{tabular}[t]{@{}c@{}}\rule{0pt}{12pt}38.1s\\[-.2pt]{\fontsize{5.8}{6.4}\selectfont \textcolor{BrickRed}{$ \blacktriangledown\,7.8\times $}}\\[-.2pt]36s\\[-.2pt]{\fontsize{5.8}{6.4}\selectfont \textcolor{BrickRed}{$ \blacktriangledown\,9.3\times $}}\end{tabular} & \hspace*{1pt}\begin{tabular}[t]{@{}c@{}}\rule{0pt}{12pt}1m\,33.5s\\[-.2pt]{\fontsize{5.8}{6.4}\selectfont \textcolor{BrickRed}{$ \blacktriangledown\,10.8\times $}}\\[-.2pt]1m\,29.7s\\[-.2pt]{\fontsize{5.8}{6.4}\selectfont \textcolor{BrickRed}{$ \blacktriangledown\,13.7\times $}}\end{tabular} & \hspace*{1pt}\begin{tabular}[t]{@{}c@{}}\rule{0pt}{12pt}6m\,31.5s\\[-.2pt]{\fontsize{5.8}{6.4}\selectfont \textcolor{gray}{n/a}}\\[-.2pt]6m\,26.3s\\[-.2pt]{\fontsize{5.8}{6.4}\selectfont \textcolor{gray}{n/a}}\end{tabular} & \hspace*{1pt}\begin{tabular}[t]{@{}c@{}}\rule{0pt}{12pt}31.6s\\[-.2pt]{\fontsize{5.8}{6.4}\selectfont \textcolor{gray}{n/a}}\\[-.2pt]30.3s\\[-.2pt]{\fontsize{5.8}{6.4}\selectfont \textcolor{gray}{n/a}}\end{tabular} & \hspace*{1pt}\begin{tabular}[t]{@{}c@{}}\rule{0pt}{12pt}34.1s\\[-.2pt]{\fontsize{5.8}{6.4}\selectfont \textcolor{gray}{n/a}}\\[-.2pt]32s\\[-.2pt]{\fontsize{5.8}{6.4}\selectfont \textcolor{gray}{n/a}}\end{tabular} & \hspace*{1pt}\begin{tabular}[t]{@{}c@{}}\rule{0pt}{12pt}1m\,9.5s\\[-.2pt]{\fontsize{5.8}{6.4}\selectfont \textcolor{gray}{n/a}}\\[-.2pt]1m\,6.1s\\[-.2pt]{\fontsize{5.8}{6.4}\selectfont \textcolor{gray}{n/a}}\end{tabular} & \hspace*{1pt}\begin{tabular}[t]{@{}c@{}}\rule{0pt}{12pt}3m\,1s\\[-.2pt]{\fontsize{5.8}{6.4}\selectfont \textcolor{gray}{n/a}}\\[-.2pt]2m\,51.5s\\[-.2pt]{\fontsize{5.8}{6.4}\selectfont \textcolor{gray}{n/a}}\end{tabular} & \hspace*{1pt}\begin{tabular}[t]{@{}c@{}}\rule{0pt}{12pt}2m\,55.5s\\[-.2pt]{\fontsize{5.8}{6.4}\selectfont \textcolor{ForestGreen}{$ \blacktriangle\,20\% $}}\\[-.2pt]2m\,54s\\[-.2pt]{\fontsize{5.8}{6.4}\selectfont \textcolor{ForestGreen}{$ \blacktriangle\,17\% $}}\end{tabular} & \hspace*{1pt}\begin{tabular}[t]{@{}c@{}}\rule{0pt}{12pt}2m\,50.4s\\[-.2pt]{\fontsize{5.8}{6.4}\selectfont \textcolor{ForestGreen}{$ \blacktriangle\,23\% $}}\\[-.2pt]2m\,48.2s\\[-.2pt]{\fontsize{5.8}{6.4}\selectfont \textcolor{ForestGreen}{$ \blacktriangle\,20\% $}}\end{tabular} & \hspace*{1pt}\begin{tabular}[t]{@{}c@{}}\rule{0pt}{12pt}6m\,12.2s\\[-.2pt]{\fontsize{5.8}{6.4}\selectfont \textcolor{ForestGreen}{$ \blacktriangle\,9\% $}}\\[-.2pt]6m\,8.2s\\[-.2pt]{\fontsize{5.8}{6.4}\selectfont \textcolor{gray}{$ \blacktriangleright\,5\% $}}\end{tabular} & \hspace*{1pt}\begin{tabular}[t]{@{}c@{}}\rule{0pt}{12pt}44m\,20.5s\\[-.2pt]{\fontsize{5.8}{6.4}\selectfont \textcolor{gray}{n/a}}\\[-.2pt]44m\,15.6s\\[-.2pt]{\fontsize{5.8}{6.4}\selectfont \textcolor{gray}{n/a}}\end{tabular} & \hspace*{1pt}\begin{tabular}[t]{@{}c@{}}\rule{0pt}{12pt}1m\,47.1s\\[-.2pt]{\fontsize{5.8}{6.4}\selectfont \textcolor{ForestGreen}{$ \blacktriangle\,5.5\times $}}\\[-.2pt]1m\,45.8s\\[-.2pt]{\fontsize{5.8}{6.4}\selectfont \textcolor{ForestGreen}{$ \blacktriangle\,2.1\times $}}\end{tabular} & \hspace*{1pt}\begin{tabular}[t]{@{}c@{}}\rule{0pt}{12pt}1m\,47.2s\\[-.2pt]{\fontsize{5.8}{6.4}\selectfont \textcolor{ForestGreen}{$ \blacktriangle\,4.6\times $}}\\[-.2pt]1m\,45.1s\\[-.2pt]{\fontsize{5.8}{6.4}\selectfont \textcolor{ForestGreen}{$ \blacktriangle\,2.2\times $}}\end{tabular} & \hspace*{1pt}\begin{tabular}[t]{@{}c@{}}\rule{0pt}{12pt}3m\,40.3s\\[-.2pt]{\fontsize{5.8}{6.4}\selectfont \textcolor{ForestGreen}{$ \blacktriangle\,4.2\times $}}\\[-.2pt]3m\,36.6s\\[-.2pt]{\fontsize{5.8}{6.4}\selectfont \textcolor{ForestGreen}{$ \blacktriangle\,1.7\times $}}\end{tabular} & \hspace*{1pt}\begin{tabular}[t]{@{}c@{}}\rule{0pt}{12pt}29m\,43.4s\\[-.2pt]{\fontsize{5.8}{6.4}\selectfont \textcolor{gray}{n/a}}\\[-.2pt]29m\,38.5s\\[-.2pt]{\fontsize{5.8}{6.4}\selectfont \textcolor{gray}{n/a}}\end{tabular}\\[3pt]\hline
\textbf{BioseqDB} & \begin{tabular}[t]{@{}l@{}}\rule{0pt}{12pt}MPF / SIMD-ref\\Scalar-DP / Block\end{tabular} & \hspace*{1pt}\begin{tabular}[t]{@{}c@{}}\rule{0pt}{12pt}\colorbox{BenchBest}{\strut 5.3s}\\[-.2pt]{\fontsize{5.8}{6.4}\selectfont \textcolor{ForestGreen}{$ \blacktriangle\,6.9\times $}}\\[-.2pt]4.5s\\[-.2pt]{\fontsize{5.8}{6.4}\selectfont \textcolor{ForestGreen}{$ \blacktriangle\,7.8\times $}}\end{tabular} & \hspace*{1pt}\begin{tabular}[t]{@{}c@{}}\rule{0pt}{12pt}\colorbox{BenchBest}{\strut 4.9s}\\[-.2pt]{\fontsize{5.8}{6.4}\selectfont \textcolor{ForestGreen}{$ \blacktriangle\,7.8\times $}}\\[-.2pt]4s\\[-.2pt]{\fontsize{5.8}{6.4}\selectfont \textcolor{ForestGreen}{$ \blacktriangle\,9.0\times $}}\end{tabular} & \textcolor{gray}{\ensuremath{\cdot}} & \textcolor{gray}{\ensuremath{\cdot}} & \textcolor{gray}{\ensuremath{\cdot}} & \textcolor{gray}{\ensuremath{\cdot}} & \textcolor{gray}{\ensuremath{\cdot}} & \textcolor{gray}{\ensuremath{\cdot}} & \textcolor{gray}{\ensuremath{\cdot}} & \textcolor{gray}{\ensuremath{\cdot}} & \textcolor{gray}{\ensuremath{\cdot}} & \textcolor{gray}{\ensuremath{\cdot}} & \textcolor{gray}{\ensuremath{\cdot}} & \textcolor{gray}{\ensuremath{\cdot}} & \textcolor{gray}{\ensuremath{\cdot}} & \textcolor{gray}{\ensuremath{\cdot}}\\[3pt]\hline
 & \begin{tabular}[t]{@{}l@{}}\rule{0pt}{12pt}MPF / SIMD-ref\\Scalar-DP / KSW2\end{tabular} & \hspace*{1pt}\begin{tabular}[t]{@{}c@{}}\rule{0pt}{12pt}13.8s\\[-.2pt]{\fontsize{5.8}{6.4}\selectfont \textcolor{ForestGreen}{$ \blacktriangle\,2.6\times $}}\\[-.2pt]11.3s\\[-.2pt]{\fontsize{5.8}{6.4}\selectfont \textcolor{ForestGreen}{$ \blacktriangle\,3.1\times $}}\end{tabular} & \hspace*{1pt}\begin{tabular}[t]{@{}c@{}}\rule{0pt}{12pt}6.3s\\[-.2pt]{\fontsize{5.8}{6.4}\selectfont \textcolor{ForestGreen}{$ \blacktriangle\,6.0\times $}}\\[-.2pt]3.9s\\[-.2pt]{\fontsize{5.8}{6.4}\selectfont \textcolor{ForestGreen}{$ \blacktriangle\,9.3\times $}}\end{tabular} & \textcolor{gray}{\ensuremath{\cdot}} & \textcolor{gray}{\ensuremath{\cdot}} & \textcolor{gray}{\ensuremath{\cdot}} & \textcolor{gray}{\ensuremath{\cdot}} & \textcolor{gray}{\ensuremath{\cdot}} & \textcolor{gray}{\ensuremath{\cdot}} & \textcolor{gray}{\ensuremath{\cdot}} & \textcolor{gray}{\ensuremath{\cdot}} & \textcolor{gray}{\ensuremath{\cdot}} & \textcolor{gray}{\ensuremath{\cdot}} & \textcolor{gray}{\ensuremath{\cdot}} & \textcolor{gray}{\ensuremath{\cdot}} & \textcolor{gray}{\ensuremath{\cdot}} & \textcolor{gray}{\ensuremath{\cdot}}\\[3pt]\hline
 & \begin{tabular}[t]{@{}l@{}}\rule{0pt}{12pt}MPF / SIMD-ref\\SIMD-DP / Block\end{tabular} & \hspace*{1pt}\begin{tabular}[t]{@{}c@{}}\rule{0pt}{12pt}5.4s\\[-.2pt]{\fontsize{5.8}{6.4}\selectfont \textcolor{ForestGreen}{$ \blacktriangle\,6.7\times $}}\\[-.2pt]4.5s\\[-.2pt]{\fontsize{5.8}{6.4}\selectfont \textcolor{ForestGreen}{$ \blacktriangle\,7.8\times $}}\end{tabular} & \hspace*{1pt}\begin{tabular}[t]{@{}c@{}}\rule{0pt}{12pt}17.6s\\[-.2pt]{\fontsize{5.8}{6.4}\selectfont \textcolor{ForestGreen}{$ \blacktriangle\,2.2\times $}}\\[-.2pt]16.5s\\[-.2pt]{\fontsize{5.8}{6.4}\selectfont \textcolor{ForestGreen}{$ \blacktriangle\,2.2\times $}}\end{tabular} & \hspace*{1pt}\begin{tabular}[t]{@{}c@{}}\rule{0pt}{12pt}\colorbox{BenchBest}{\strut 8.7s}\\[-.2pt]{\fontsize{5.8}{6.4}\selectfont \textcolor{ForestGreen}{$ \blacktriangle\,10.8\times $}}\\[-.2pt]\colorbox{BenchBest}{\strut 6.5s}\\[-.2pt]{\fontsize{5.8}{6.4}\selectfont \textcolor{ForestGreen}{$ \blacktriangle\,13.7\times $}}\end{tabular} & \textcolor{gray}{\ensuremath{\cdot}} & \textcolor{gray}{\ensuremath{\cdot}} & \textcolor{gray}{\ensuremath{\cdot}} & \textcolor{gray}{\ensuremath{\cdot}} & \textcolor{gray}{\ensuremath{\cdot}} & \textcolor{gray}{\ensuremath{\cdot}} & \textcolor{gray}{\ensuremath{\cdot}} & \textcolor{gray}{\ensuremath{\cdot}} & \textcolor{gray}{\ensuremath{\cdot}} & \textcolor{gray}{\ensuremath{\cdot}} & \textcolor{gray}{\ensuremath{\cdot}} & \textcolor{gray}{\ensuremath{\cdot}} & \textcolor{gray}{\ensuremath{\cdot}}\\[3pt]\hline
 & \begin{tabular}[t]{@{}l@{}}\rule{0pt}{12pt}MPF / SIMD-ref\\SIMD-DP / KSW2\end{tabular} & \hspace*{1pt}\begin{tabular}[t]{@{}c@{}}\rule{0pt}{12pt}6.9s\\[-.2pt]{\fontsize{5.8}{6.4}\selectfont \textcolor{ForestGreen}{$ \blacktriangle\,5.2\times $}}\\[-.2pt]4.4s\\[-.2pt]{\fontsize{5.8}{6.4}\selectfont \textcolor{ForestGreen}{$ \blacktriangle\,8.0\times $}}\end{tabular} & \hspace*{1pt}\begin{tabular}[t]{@{}c@{}}\rule{0pt}{12pt}6.5s\\[-.2pt]{\fontsize{5.8}{6.4}\selectfont \textcolor{ForestGreen}{$ \blacktriangle\,5.9\times $}}\\[-.2pt]3.9s\\[-.2pt]{\fontsize{5.8}{6.4}\selectfont \textcolor{ForestGreen}{$ \blacktriangle\,9.2\times $}}\end{tabular} & \textcolor{gray}{\ensuremath{\cdot}} & \textcolor{gray}{\ensuremath{\cdot}} & \textcolor{gray}{\ensuremath{\cdot}} & \textcolor{gray}{\ensuremath{\cdot}} & \textcolor{gray}{\ensuremath{\cdot}} & \textcolor{gray}{\ensuremath{\cdot}} & \textcolor{gray}{\ensuremath{\cdot}} & \textcolor{gray}{\ensuremath{\cdot}} & \textcolor{gray}{\ensuremath{\cdot}} & \textcolor{gray}{\ensuremath{\cdot}} & \textcolor{gray}{\ensuremath{\cdot}} & \textcolor{gray}{\ensuremath{\cdot}} & \textcolor{gray}{\ensuremath{\cdot}} & \textcolor{gray}{\ensuremath{\cdot}}\\[3pt]\hline
 & \begin{tabular}[t]{@{}l@{}}\rule{0pt}{12pt}MPF / Two-stack\\RMQ / Block\end{tabular} & \hspace*{1pt}\begin{tabular}[t]{@{}c@{}}\rule{0pt}{12pt}5.3s\\[-.2pt]{\fontsize{5.8}{6.4}\selectfont \textcolor{ForestGreen}{$ \blacktriangle\,6.8\times $}}\\[-.2pt]4.4s\\[-.2pt]{\fontsize{5.8}{6.4}\selectfont \textcolor{ForestGreen}{$ \blacktriangle\,7.9\times $}}\end{tabular} & \hspace*{1pt}\begin{tabular}[t]{@{}c@{}}\rule{0pt}{12pt}5s\\[-.2pt]{\fontsize{5.8}{6.4}\selectfont \textcolor{ForestGreen}{$ \blacktriangle\,7.7\times $}}\\[-.2pt]4.1s\\[-.2pt]{\fontsize{5.8}{6.4}\selectfont \textcolor{ForestGreen}{$ \blacktriangle\,8.9\times $}}\end{tabular} & \textcolor{gray}{\ensuremath{\cdot}} & \textcolor{gray}{\ensuremath{\cdot}} & \textcolor{gray}{\ensuremath{\cdot}} & \textcolor{gray}{\ensuremath{\cdot}} & \textcolor{gray}{\ensuremath{\cdot}} & \textcolor{gray}{\ensuremath{\cdot}} & \hspace*{1pt}\begin{tabular}[t]{@{}c@{}}\rule{0pt}{12pt}4m\,31.5s\\[-.2pt]{\fontsize{5.8}{6.4}\selectfont \textcolor{BrickRed}{$ \blacktriangledown\,1.5\times $}}\\[-.2pt]4m\,22.4s\\[-.2pt]{\fontsize{5.8}{6.4}\selectfont \textcolor{BrickRed}{$ \blacktriangledown\,1.5\times $}}\end{tabular} & \hspace*{1pt}\begin{tabular}[t]{@{}c@{}}\rule{0pt}{12pt}4m\,45s\\[-.2pt]{\fontsize{5.8}{6.4}\selectfont \textcolor{BrickRed}{$ \blacktriangledown\,1.7\times $}}\\[-.2pt]4m\,33.2s\\[-.2pt]{\fontsize{5.8}{6.4}\selectfont \textcolor{BrickRed}{$ \blacktriangledown\,1.6\times $}}\end{tabular} & \textcolor{gray}{\ensuremath{\cdot}} & \textcolor{gray}{\ensuremath{\cdot}} & \textcolor{gray}{\ensuremath{\cdot}} & \textcolor{gray}{\ensuremath{\cdot}} & \textcolor{gray}{\ensuremath{\cdot}} & \textcolor{gray}{\ensuremath{\cdot}}\\[3pt]\hline
 & \begin{tabular}[t]{@{}l@{}}\rule{0pt}{12pt}MPF / Two-stack\\RMQ / KSW2\end{tabular} & \hspace*{1pt}\begin{tabular}[t]{@{}c@{}}\rule{0pt}{12pt}6.9s\\[-.2pt]{\fontsize{5.8}{6.4}\selectfont \textcolor{ForestGreen}{$ \blacktriangle\,5.2\times $}}\\[-.2pt]4.4s\\[-.2pt]{\fontsize{5.8}{6.4}\selectfont \textcolor{ForestGreen}{$ \blacktriangle\,8.0\times $}}\end{tabular} & \hspace*{1pt}\begin{tabular}[t]{@{}c@{}}\rule{0pt}{12pt}6.4s\\[-.2pt]{\fontsize{5.8}{6.4}\selectfont \textcolor{ForestGreen}{$ \blacktriangle\,5.9\times $}}\\[-.2pt]\colorbox{BenchBest}{\strut 3.9s}\\[-.2pt]{\fontsize{5.8}{6.4}\selectfont \textcolor{ForestGreen}{$ \blacktriangle\,9.3\times $}}\end{tabular} & \textcolor{gray}{\ensuremath{\cdot}} & \textcolor{gray}{\ensuremath{\cdot}} & \textcolor{gray}{\ensuremath{\cdot}} & \textcolor{gray}{\ensuremath{\cdot}} & \textcolor{gray}{\ensuremath{\cdot}} & \textcolor{gray}{\ensuremath{\cdot}} & \textcolor{gray}{\ensuremath{\cdot}} & \textcolor{gray}{\ensuremath{\cdot}} & \textcolor{gray}{\ensuremath{\cdot}} & \textcolor{gray}{\ensuremath{\cdot}} & \textcolor{gray}{\ensuremath{\cdot}} & \textcolor{gray}{\ensuremath{\cdot}} & \textcolor{gray}{\ensuremath{\cdot}} & \textcolor{gray}{\ensuremath{\cdot}}\\[3pt]\hline
 & \begin{tabular}[t]{@{}l@{}}\rule{0pt}{12pt}MPF / Two-stack\\Scalar-DP / Block\end{tabular} & \hspace*{1pt}\begin{tabular}[t]{@{}c@{}}\rule{0pt}{12pt}5.3s\\[-.2pt]{\fontsize{5.8}{6.4}\selectfont \textcolor{ForestGreen}{$ \blacktriangle\,6.8\times $}}\\[-.2pt]4.5s\\[-.2pt]{\fontsize{5.8}{6.4}\selectfont \textcolor{ForestGreen}{$ \blacktriangle\,7.7\times $}}\end{tabular} & \hspace*{1pt}\begin{tabular}[t]{@{}c@{}}\rule{0pt}{12pt}4.9s\\[-.2pt]{\fontsize{5.8}{6.4}\selectfont \textcolor{ForestGreen}{$ \blacktriangle\,7.8\times $}}\\[-.2pt]4s\\[-.2pt]{\fontsize{5.8}{6.4}\selectfont \textcolor{ForestGreen}{$ \blacktriangle\,8.9\times $}}\end{tabular} & \textcolor{gray}{\ensuremath{\cdot}} & \textcolor{gray}{\ensuremath{\cdot}} & \textcolor{gray}{\ensuremath{\cdot}} & \textcolor{gray}{\ensuremath{\cdot}} & \textcolor{gray}{\ensuremath{\cdot}} & \textcolor{gray}{\ensuremath{\cdot}} & \textcolor{gray}{\ensuremath{\cdot}} & \textcolor{gray}{\ensuremath{\cdot}} & \textcolor{gray}{\ensuremath{\cdot}} & \textcolor{gray}{\ensuremath{\cdot}} & \textcolor{gray}{\ensuremath{\cdot}} & \textcolor{gray}{\ensuremath{\cdot}} & \textcolor{gray}{\ensuremath{\cdot}} & \textcolor{gray}{\ensuremath{\cdot}}\\[3pt]\hline
\end{tabular}\par}
\smallskip{\footnotesize\noindent Each cell: A, comparison, B, comparison. Times cover 1,024 queries. $\cdot$ = no approved observation. ---\textsuperscript{g} = GPU unavailable, not executed. n/a = no measured comparison.\par
\noindent \textsuperscript{f}SIMD-ref request observed as scalar, compared only as that effective path. \textsuperscript{p}Accepted by explicit preservation after wrapper postflight failure, no run receipt claimed.\par
\noindent MPF/MPS = minimizers per fragment/sequence. Row order: index representation / minimizer, then chaining / base alignment. SIMD-ref is the reference minimizer, Bounded is bounded ntHash.\par
\noindent Blue = one lowest A and one lowest B down each machine/workload column, across continuation pages. Exact ties select the first row. Repeated baselines are marked only on page 1.\par
\noindent Light gray = minimap2 baseline. Pale yellow = preferred MPF / Bounded / SIMD-DP / Block configuration, not a claim of winning every column.\par
\noindent When all Vulkan modes are unmeasured, they share one row. Baseline comparisons use the fastest measured BioseqDB effective path separately for A and B.\par}
\clearpage
\noindent{\large\bfseries AWS: 4 Gbp, k=15, w=10 (mapping interval)\hfill 4/6}\par\smallskip
{\fontsize{7.1}{8.1}\selectfont\setlength{\tabcolsep}{1pt}\setlength{\fboxsep}{2pt}\renewcommand{\arraystretch}{1.12}
\noindent\begin{tabular}{@{}p{14mm}p{27mm}|p{13.55mm}|p{13.55mm}|p{13.55mm}|p{13.55mm}|p{13.55mm}|p{13.55mm}|p{13.55mm}|p{13.55mm}|p{13.55mm}|p{13.55mm}|p{13.55mm}|p{13.55mm}|p{13.55mm}|p{13.55mm}|p{13.55mm}|p{13.55mm}|@{}}\hline
Engine & Execution & \multicolumn{8}{c|}{\textbf{Index}} & \multicolumn{8}{c|}{\textbf{No index}}\\
 &  & \multicolumn{8}{c|}{\textbf{4 Gbp}} & \multicolumn{8}{c|}{\textbf{4 Gbp}}\\
 &  & \multicolumn{8}{c|}{\textbf{k=15, w=10}} & \multicolumn{8}{c|}{\textbf{k=15, w=10}}\\
 &  & \multicolumn{4}{c|}{\textbf{DNA4}} & \multicolumn{4}{c|}{\textbf{DNA5}} & \multicolumn{4}{c|}{\textbf{DNA4}} & \multicolumn{4}{c|}{\textbf{DNA5}}\\
 &  & \multicolumn{1}{c|}{\textbf{M7A}} & \multicolumn{1}{c|}{\textbf{M7I}} & \multicolumn{1}{c|}{\textbf{M4}} & \multicolumn{1}{c|}{\textbf{C7G}} & \multicolumn{1}{c|}{\textbf{M7A}} & \multicolumn{1}{c|}{\textbf{M7I}} & \multicolumn{1}{c|}{\textbf{M4}} & \multicolumn{1}{c|}{\textbf{C7G}} & \multicolumn{1}{c|}{\textbf{M7A}} & \multicolumn{1}{c|}{\textbf{M7I}} & \multicolumn{1}{c|}{\textbf{M4}} & \multicolumn{1}{c|}{\textbf{C7G}} & \multicolumn{1}{c|}{\textbf{M7A}} & \multicolumn{1}{c|}{\textbf{M7I}} & \multicolumn{1}{c|}{\textbf{M4}} & \multicolumn{1}{c|}{\textbf{C7G}}\\
\hline
\rowcolor{BenchBaseline}[0.5pt][0.5pt]\textbf{minimap2} & \textbf{Baseline} & \hspace*{1pt}\begin{tabular}[t]{@{}c@{}}\rule{0pt}{12pt}36.2s\\[-.2pt]{\fontsize{5.8}{6.4}\selectfont \textcolor{BrickRed}{$ \blacktriangledown\,6.9\times $}}\\[-.2pt]35s\\[-.2pt]{\fontsize{5.8}{6.4}\selectfont \textcolor{BrickRed}{$ \blacktriangledown\,8.0\times $}}\end{tabular} & \hspace*{1pt}\begin{tabular}[t]{@{}c@{}}\rule{0pt}{12pt}38.1s\\[-.2pt]{\fontsize{5.8}{6.4}\selectfont \textcolor{BrickRed}{$ \blacktriangledown\,7.8\times $}}\\[-.2pt]36s\\[-.2pt]{\fontsize{5.8}{6.4}\selectfont \textcolor{BrickRed}{$ \blacktriangledown\,9.3\times $}}\end{tabular} & \hspace*{1pt}\begin{tabular}[t]{@{}c@{}}\rule{0pt}{12pt}1m\,33.5s\\[-.2pt]{\fontsize{5.8}{6.4}\selectfont \textcolor{BrickRed}{$ \blacktriangledown\,10.8\times $}}\\[-.2pt]1m\,29.7s\\[-.2pt]{\fontsize{5.8}{6.4}\selectfont \textcolor{BrickRed}{$ \blacktriangledown\,13.7\times $}}\end{tabular} & \hspace*{1pt}\begin{tabular}[t]{@{}c@{}}\rule{0pt}{12pt}6m\,31.5s\\[-.2pt]{\fontsize{5.8}{6.4}\selectfont \textcolor{gray}{n/a}}\\[-.2pt]6m\,26.3s\\[-.2pt]{\fontsize{5.8}{6.4}\selectfont \textcolor{gray}{n/a}}\end{tabular} & \hspace*{1pt}\begin{tabular}[t]{@{}c@{}}\rule{0pt}{12pt}31.6s\\[-.2pt]{\fontsize{5.8}{6.4}\selectfont \textcolor{gray}{n/a}}\\[-.2pt]30.3s\\[-.2pt]{\fontsize{5.8}{6.4}\selectfont \textcolor{gray}{n/a}}\end{tabular} & \hspace*{1pt}\begin{tabular}[t]{@{}c@{}}\rule{0pt}{12pt}34.1s\\[-.2pt]{\fontsize{5.8}{6.4}\selectfont \textcolor{gray}{n/a}}\\[-.2pt]32s\\[-.2pt]{\fontsize{5.8}{6.4}\selectfont \textcolor{gray}{n/a}}\end{tabular} & \hspace*{1pt}\begin{tabular}[t]{@{}c@{}}\rule{0pt}{12pt}1m\,9.5s\\[-.2pt]{\fontsize{5.8}{6.4}\selectfont \textcolor{gray}{n/a}}\\[-.2pt]1m\,6.1s\\[-.2pt]{\fontsize{5.8}{6.4}\selectfont \textcolor{gray}{n/a}}\end{tabular} & \hspace*{1pt}\begin{tabular}[t]{@{}c@{}}\rule{0pt}{12pt}3m\,1s\\[-.2pt]{\fontsize{5.8}{6.4}\selectfont \textcolor{gray}{n/a}}\\[-.2pt]2m\,51.5s\\[-.2pt]{\fontsize{5.8}{6.4}\selectfont \textcolor{gray}{n/a}}\end{tabular} & \hspace*{1pt}\begin{tabular}[t]{@{}c@{}}\rule{0pt}{12pt}2m\,55.5s\\[-.2pt]{\fontsize{5.8}{6.4}\selectfont \textcolor{ForestGreen}{$ \blacktriangle\,20\% $}}\\[-.2pt]2m\,54s\\[-.2pt]{\fontsize{5.8}{6.4}\selectfont \textcolor{ForestGreen}{$ \blacktriangle\,17\% $}}\end{tabular} & \hspace*{1pt}\begin{tabular}[t]{@{}c@{}}\rule{0pt}{12pt}2m\,50.4s\\[-.2pt]{\fontsize{5.8}{6.4}\selectfont \textcolor{ForestGreen}{$ \blacktriangle\,23\% $}}\\[-.2pt]2m\,48.2s\\[-.2pt]{\fontsize{5.8}{6.4}\selectfont \textcolor{ForestGreen}{$ \blacktriangle\,20\% $}}\end{tabular} & \hspace*{1pt}\begin{tabular}[t]{@{}c@{}}\rule{0pt}{12pt}6m\,12.2s\\[-.2pt]{\fontsize{5.8}{6.4}\selectfont \textcolor{ForestGreen}{$ \blacktriangle\,9\% $}}\\[-.2pt]6m\,8.2s\\[-.2pt]{\fontsize{5.8}{6.4}\selectfont \textcolor{gray}{$ \blacktriangleright\,5\% $}}\end{tabular} & \hspace*{1pt}\begin{tabular}[t]{@{}c@{}}\rule{0pt}{12pt}44m\,20.5s\\[-.2pt]{\fontsize{5.8}{6.4}\selectfont \textcolor{gray}{n/a}}\\[-.2pt]44m\,15.6s\\[-.2pt]{\fontsize{5.8}{6.4}\selectfont \textcolor{gray}{n/a}}\end{tabular} & \hspace*{1pt}\begin{tabular}[t]{@{}c@{}}\rule{0pt}{12pt}1m\,47.1s\\[-.2pt]{\fontsize{5.8}{6.4}\selectfont \textcolor{ForestGreen}{$ \blacktriangle\,5.5\times $}}\\[-.2pt]1m\,45.8s\\[-.2pt]{\fontsize{5.8}{6.4}\selectfont \textcolor{ForestGreen}{$ \blacktriangle\,2.1\times $}}\end{tabular} & \hspace*{1pt}\begin{tabular}[t]{@{}c@{}}\rule{0pt}{12pt}1m\,47.2s\\[-.2pt]{\fontsize{5.8}{6.4}\selectfont \textcolor{ForestGreen}{$ \blacktriangle\,4.6\times $}}\\[-.2pt]1m\,45.1s\\[-.2pt]{\fontsize{5.8}{6.4}\selectfont \textcolor{ForestGreen}{$ \blacktriangle\,2.2\times $}}\end{tabular} & \hspace*{1pt}\begin{tabular}[t]{@{}c@{}}\rule{0pt}{12pt}3m\,40.3s\\[-.2pt]{\fontsize{5.8}{6.4}\selectfont \textcolor{ForestGreen}{$ \blacktriangle\,4.2\times $}}\\[-.2pt]3m\,36.6s\\[-.2pt]{\fontsize{5.8}{6.4}\selectfont \textcolor{ForestGreen}{$ \blacktriangle\,1.7\times $}}\end{tabular} & \hspace*{1pt}\begin{tabular}[t]{@{}c@{}}\rule{0pt}{12pt}29m\,43.4s\\[-.2pt]{\fontsize{5.8}{6.4}\selectfont \textcolor{gray}{n/a}}\\[-.2pt]29m\,38.5s\\[-.2pt]{\fontsize{5.8}{6.4}\selectfont \textcolor{gray}{n/a}}\end{tabular}\\[3pt]\hline
\textbf{BioseqDB} & \begin{tabular}[t]{@{}l@{}}\rule{0pt}{12pt}MPF / Two-stack\\Scalar-DP / KSW2\end{tabular} & \hspace*{1pt}\begin{tabular}[t]{@{}c@{}}\rule{0pt}{12pt}20.1s\\[-.2pt]{\fontsize{5.8}{6.4}\selectfont \textcolor{ForestGreen}{$ \blacktriangle\,1.8\times $}}\\[-.2pt]17.6s\\[-.2pt]{\fontsize{5.8}{6.4}\selectfont \textcolor{ForestGreen}{$ \blacktriangle\,2.0\times $}}\end{tabular} & \hspace*{1pt}\begin{tabular}[t]{@{}c@{}}\rule{0pt}{12pt}16.8s\\[-.2pt]{\fontsize{5.8}{6.4}\selectfont \textcolor{ForestGreen}{$ \blacktriangle\,2.3\times $}}\\[-.2pt]14.4s\\[-.2pt]{\fontsize{5.8}{6.4}\selectfont \textcolor{ForestGreen}{$ \blacktriangle\,2.5\times $}}\end{tabular} & \textcolor{gray}{\ensuremath{\cdot}} & \textcolor{gray}{\ensuremath{\cdot}} & \textcolor{gray}{\ensuremath{\cdot}} & \textcolor{gray}{\ensuremath{\cdot}} & \textcolor{gray}{\ensuremath{\cdot}} & \textcolor{gray}{\ensuremath{\cdot}} & \textcolor{gray}{\ensuremath{\cdot}} & \textcolor{gray}{\ensuremath{\cdot}} & \textcolor{gray}{\ensuremath{\cdot}} & \textcolor{gray}{\ensuremath{\cdot}} & \hspace*{1pt}\begin{tabular}[t]{@{}c@{}}\rule{0pt}{12pt}10m\,1.7s\\[-.2pt]{\fontsize{5.8}{6.4}\selectfont \textcolor{BrickRed}{$ \blacktriangledown\,5.6\times $}}\\[-.2pt]3m\,53.1s\\[-.2pt]{\fontsize{5.8}{6.4}\selectfont \textcolor{BrickRed}{$ \blacktriangledown\,2.2\times $}}\end{tabular} & \hspace*{1pt}\begin{tabular}[t]{@{}c@{}}\rule{0pt}{12pt}8m\,16.7s\\[-.2pt]{\fontsize{5.8}{6.4}\selectfont \textcolor{BrickRed}{$ \blacktriangledown\,4.6\times $}}\\[-.2pt]3m\,46.2s\\[-.2pt]{\fontsize{5.8}{6.4}\selectfont \textcolor{BrickRed}{$ \blacktriangledown\,2.2\times $}}\end{tabular} & \textcolor{gray}{\ensuremath{\cdot}} & \textcolor{gray}{\ensuremath{\cdot}}\\[3pt]\hline
 & \begin{tabular}[t]{@{}l@{}}\rule{0pt}{12pt}MPF / Two-stack\\SIMD-DP / Block\end{tabular} & \hspace*{1pt}\begin{tabular}[t]{@{}c@{}}\rule{0pt}{12pt}5.4s\\[-.2pt]{\fontsize{5.8}{6.4}\selectfont \textcolor{ForestGreen}{$ \blacktriangle\,6.7\times $}}\\[-.2pt]4.5s\\[-.2pt]{\fontsize{5.8}{6.4}\selectfont \textcolor{ForestGreen}{$ \blacktriangle\,7.8\times $}}\end{tabular} & \hspace*{1pt}\begin{tabular}[t]{@{}c@{}}\rule{0pt}{12pt}5s\\[-.2pt]{\fontsize{5.8}{6.4}\selectfont \textcolor{ForestGreen}{$ \blacktriangle\,7.6\times $}}\\[-.2pt]4s\\[-.2pt]{\fontsize{5.8}{6.4}\selectfont \textcolor{ForestGreen}{$ \blacktriangle\,8.9\times $}}\end{tabular} & \textcolor{gray}{\ensuremath{\cdot}} & \textcolor{gray}{\ensuremath{\cdot}} & \textcolor{gray}{\ensuremath{\cdot}} & \textcolor{gray}{\ensuremath{\cdot}} & \textcolor{gray}{\ensuremath{\cdot}} & \textcolor{gray}{\ensuremath{\cdot}} & \textcolor{gray}{\ensuremath{\cdot}} & \textcolor{gray}{\ensuremath{\cdot}} & \textcolor{gray}{\ensuremath{\cdot}} & \textcolor{gray}{\ensuremath{\cdot}} & \textcolor{gray}{\ensuremath{\cdot}} & \textcolor{gray}{\ensuremath{\cdot}} & \textcolor{gray}{\ensuremath{\cdot}} & \textcolor{gray}{\ensuremath{\cdot}}\\[3pt]\hline
 & \begin{tabular}[t]{@{}l@{}}\rule{0pt}{12pt}MPF / Two-stack\\SIMD-DP / KSW2\end{tabular} & \hspace*{1pt}\begin{tabular}[t]{@{}c@{}}\rule{0pt}{12pt}6.9s\\[-.2pt]{\fontsize{5.8}{6.4}\selectfont \textcolor{ForestGreen}{$ \blacktriangle\,5.3\times $}}\\[-.2pt]\colorbox{BenchBest}{\strut 4.4s}\\[-.2pt]{\fontsize{5.8}{6.4}\selectfont \textcolor{ForestGreen}{$ \blacktriangle\,8.0\times $}}\end{tabular} & \hspace*{1pt}\begin{tabular}[t]{@{}c@{}}\rule{0pt}{12pt}7.3s\\[-.2pt]{\fontsize{5.8}{6.4}\selectfont \textcolor{ForestGreen}{$ \blacktriangle\,5.2\times $}}\\[-.2pt]4.8s\\[-.2pt]{\fontsize{5.8}{6.4}\selectfont \textcolor{ForestGreen}{$ \blacktriangle\,7.5\times $}}\end{tabular} & \hspace*{1pt}\begin{tabular}[t]{@{}c@{}}\rule{0pt}{12pt}12.5s\\[-.2pt]{\fontsize{5.8}{6.4}\selectfont \textcolor{ForestGreen}{$ \blacktriangle\,7.5\times $}}\\[-.2pt]6.8s\\[-.2pt]{\fontsize{5.8}{6.4}\selectfont \textcolor{ForestGreen}{$ \blacktriangle\,13.3\times $}}\end{tabular} & \textcolor{gray}{\ensuremath{\cdot}} & \textcolor{gray}{\ensuremath{\cdot}} & \textcolor{gray}{\ensuremath{\cdot}} & \textcolor{gray}{\ensuremath{\cdot}} & \textcolor{gray}{\ensuremath{\cdot}} & \hspace*{1pt}\begin{tabular}[t]{@{}c@{}}\rule{0pt}{12pt}4m\,37.5s\\[-.2pt]{\fontsize{5.8}{6.4}\selectfont \textcolor{BrickRed}{$ \blacktriangledown\,1.6\times $}}\\[-.2pt]4m\,25.9s\\[-.2pt]{\fontsize{5.8}{6.4}\selectfont \textcolor{BrickRed}{$ \blacktriangledown\,1.5\times $}}\end{tabular} & \hspace*{1pt}\begin{tabular}[t]{@{}c@{}}\rule{0pt}{12pt}4m\,48.3s\\[-.2pt]{\fontsize{5.8}{6.4}\selectfont \textcolor{BrickRed}{$ \blacktriangledown\,1.7\times $}}\\[-.2pt]4m\,34.5s\\[-.2pt]{\fontsize{5.8}{6.4}\selectfont \textcolor{BrickRed}{$ \blacktriangledown\,1.6\times $}}\end{tabular} & \textcolor{gray}{\ensuremath{\cdot}} & \textcolor{gray}{\ensuremath{\cdot}} & \textcolor{gray}{\ensuremath{\cdot}} & \textcolor{gray}{\ensuremath{\cdot}} & \textcolor{gray}{\ensuremath{\cdot}} & \textcolor{gray}{\ensuremath{\cdot}}\\[3pt]\hline
 & \begin{tabular}[t]{@{}l@{}}\rule{0pt}{12pt}MPS / Bounded\\RMQ / Block\end{tabular} & \hspace*{1pt}\begin{tabular}[t]{@{}c@{}}\rule{0pt}{12pt}11s\\[-.2pt]{\fontsize{5.8}{6.4}\selectfont \textcolor{ForestGreen}{$ \blacktriangle\,3.3\times $}}\\[-.2pt]10.3s\\[-.2pt]{\fontsize{5.8}{6.4}\selectfont \textcolor{ForestGreen}{$ \blacktriangle\,3.4\times $}}\end{tabular} & \hspace*{1pt}\begin{tabular}[t]{@{}c@{}}\rule{0pt}{12pt}11.5s\\[-.2pt]{\fontsize{5.8}{6.4}\selectfont \textcolor{ForestGreen}{$ \blacktriangle\,3.3\times $}}\\[-.2pt]10.7s\\[-.2pt]{\fontsize{5.8}{6.4}\selectfont \textcolor{ForestGreen}{$ \blacktriangle\,3.4\times $}}\end{tabular} & \textcolor{gray}{\ensuremath{\cdot}} & \textcolor{gray}{\ensuremath{\cdot}} & \textcolor{gray}{\ensuremath{\cdot}} & \textcolor{gray}{\ensuremath{\cdot}} & \textcolor{gray}{\ensuremath{\cdot}} & \textcolor{gray}{\ensuremath{\cdot}} & \hspace*{1pt}\begin{tabular}[t]{@{}c@{}}\rule{0pt}{12pt}4m\,10.2s\\[-.2pt]{\fontsize{5.8}{6.4}\selectfont \textcolor{BrickRed}{$ \blacktriangledown\,1.4\times $}}\\[-.2pt]4m\,1s\\[-.2pt]{\fontsize{5.8}{6.4}\selectfont \textcolor{BrickRed}{$ \blacktriangledown\,1.4\times $}}\end{tabular} & \hspace*{1pt}\begin{tabular}[t]{@{}c@{}}\rule{0pt}{12pt}4m\,11.6s\\[-.2pt]{\fontsize{5.8}{6.4}\selectfont \textcolor{BrickRed}{$ \blacktriangledown\,1.5\times $}}\\[-.2pt]3m\,59.9s\\[-.2pt]{\fontsize{5.8}{6.4}\selectfont \textcolor{BrickRed}{$ \blacktriangledown\,1.4\times $}}\end{tabular} & \hspace*{1pt}\begin{tabular}[t]{@{}c@{}}\rule{0pt}{12pt}6m\,51.1s\\[-.2pt]{\fontsize{5.8}{6.4}\selectfont \textcolor{BrickRed}{$ \blacktriangledown\,10\% $}}\\[-.2pt]6m\,31.1s\\[-.2pt]{\fontsize{5.8}{6.4}\selectfont \textcolor{BrickRed}{$ \blacktriangledown\,6\% $}}\end{tabular} & \textcolor{gray}{\ensuremath{\cdot}} & \textcolor{gray}{\ensuremath{\cdot}} & \textcolor{gray}{\ensuremath{\cdot}} & \textcolor{gray}{\ensuremath{\cdot}} & \textcolor{gray}{\ensuremath{\cdot}}\\[3pt]\hline
 & \begin{tabular}[t]{@{}l@{}}\rule{0pt}{12pt}MPS / Bounded\\RMQ / KSW2\end{tabular} & \hspace*{1pt}\begin{tabular}[t]{@{}c@{}}\rule{0pt}{12pt}12.6s\\[-.2pt]{\fontsize{5.8}{6.4}\selectfont \textcolor{ForestGreen}{$ \blacktriangle\,2.9\times $}}\\[-.2pt]9.7s\\[-.2pt]{\fontsize{5.8}{6.4}\selectfont \textcolor{ForestGreen}{$ \blacktriangle\,3.6\times $}}\end{tabular} & \hspace*{1pt}\begin{tabular}[t]{@{}c@{}}\rule{0pt}{12pt}13.7s\\[-.2pt]{\fontsize{5.8}{6.4}\selectfont \textcolor{ForestGreen}{$ \blacktriangle\,2.8\times $}}\\[-.2pt]9.5s\\[-.2pt]{\fontsize{5.8}{6.4}\selectfont \textcolor{ForestGreen}{$ \blacktriangle\,3.8\times $}}\end{tabular} & \textcolor{gray}{\ensuremath{\cdot}} & \textcolor{gray}{\ensuremath{\cdot}} & \textcolor{gray}{\ensuremath{\cdot}} & \textcolor{gray}{\ensuremath{\cdot}} & \textcolor{gray}{\ensuremath{\cdot}} & \textcolor{gray}{\ensuremath{\cdot}} & \hspace*{1pt}\begin{tabular}[t]{@{}c@{}}\rule{0pt}{12pt}4m\,20.8s\\[-.2pt]{\fontsize{5.8}{6.4}\selectfont \textcolor{BrickRed}{$ \blacktriangledown\,1.5\times $}}\\[-.2pt]4m\,9.9s\\[-.2pt]{\fontsize{5.8}{6.4}\selectfont \textcolor{BrickRed}{$ \blacktriangledown\,1.4\times $}}\end{tabular} & \hspace*{1pt}\begin{tabular}[t]{@{}c@{}}\rule{0pt}{12pt}3m\,58.2s\\[-.2pt]{\fontsize{5.8}{6.4}\selectfont \textcolor{BrickRed}{$ \blacktriangledown\,1.4\times $}}\\[-.2pt]3m\,44.9s\\[-.2pt]{\fontsize{5.8}{6.4}\selectfont \textcolor{BrickRed}{$ \blacktriangledown\,1.3\times $}}\end{tabular} & \hspace*{1pt}\begin{tabular}[t]{@{}c@{}}\rule{0pt}{12pt}6m\,49.2s\\[-.2pt]{\fontsize{5.8}{6.4}\selectfont \textcolor{BrickRed}{$ \blacktriangledown\,10\% $}}\\[-.2pt]6m\,25.6s\\[-.2pt]{\fontsize{5.8}{6.4}\selectfont \textcolor{gray}{$ \blacktriangleright\,5\% $}}\end{tabular} & \textcolor{gray}{\ensuremath{\cdot}} & \textcolor{gray}{\ensuremath{\cdot}} & \textcolor{gray}{\ensuremath{\cdot}} & \textcolor{gray}{\ensuremath{\cdot}} & \textcolor{gray}{\ensuremath{\cdot}}\\[3pt]\hline
 & \begin{tabular}[t]{@{}l@{}}\rule{0pt}{12pt}MPS / Bounded\\Scalar-DP / Block\end{tabular} & \hspace*{1pt}\begin{tabular}[t]{@{}c@{}}\rule{0pt}{12pt}10.9s\\[-.2pt]{\fontsize{5.8}{6.4}\selectfont \textcolor{ForestGreen}{$ \blacktriangle\,3.3\times $}}\\[-.2pt]10.3s\\[-.2pt]{\fontsize{5.8}{6.4}\selectfont \textcolor{ForestGreen}{$ \blacktriangle\,3.4\times $}}\end{tabular} & \hspace*{1pt}\begin{tabular}[t]{@{}c@{}}\rule{0pt}{12pt}11.1s\\[-.2pt]{\fontsize{5.8}{6.4}\selectfont \textcolor{ForestGreen}{$ \blacktriangle\,3.4\times $}}\\[-.2pt]10.3s\\[-.2pt]{\fontsize{5.8}{6.4}\selectfont \textcolor{ForestGreen}{$ \blacktriangle\,3.5\times $}}\end{tabular} & \textcolor{gray}{\ensuremath{\cdot}} & \textcolor{gray}{\ensuremath{\cdot}} & \textcolor{gray}{\ensuremath{\cdot}} & \textcolor{gray}{\ensuremath{\cdot}} & \textcolor{gray}{\ensuremath{\cdot}} & \textcolor{gray}{\ensuremath{\cdot}} & \hspace*{1pt}\begin{tabular}[t]{@{}c@{}}\rule{0pt}{12pt}4m\,17.5s\\[-.2pt]{\fontsize{5.8}{6.4}\selectfont \textcolor{BrickRed}{$ \blacktriangledown\,1.5\times $}}\\[-.2pt]4m\,8.2s\\[-.2pt]{\fontsize{5.8}{6.4}\selectfont \textcolor{BrickRed}{$ \blacktriangledown\,1.4\times $}}\end{tabular} & \hspace*{1pt}\begin{tabular}[t]{@{}c@{}}\rule{0pt}{12pt}4m\,7.1s\\[-.2pt]{\fontsize{5.8}{6.4}\selectfont \textcolor{BrickRed}{$ \blacktriangledown\,1.4\times $}}\\[-.2pt]3m\,55.3s\\[-.2pt]{\fontsize{5.8}{6.4}\selectfont \textcolor{BrickRed}{$ \blacktriangledown\,1.4\times $}}\end{tabular} & \textcolor{gray}{\ensuremath{\cdot}} & \textcolor{gray}{\ensuremath{\cdot}} & \textcolor{gray}{\ensuremath{\cdot}} & \textcolor{gray}{\ensuremath{\cdot}} & \textcolor{gray}{\ensuremath{\cdot}} & \textcolor{gray}{\ensuremath{\cdot}}\\[3pt]\hline
 & \begin{tabular}[t]{@{}l@{}}\rule{0pt}{12pt}MPS / Bounded\\Scalar-DP / KSW2\end{tabular} & \hspace*{1pt}\begin{tabular}[t]{@{}c@{}}\rule{0pt}{12pt}12.3s\\[-.2pt]{\fontsize{5.8}{6.4}\selectfont \textcolor{ForestGreen}{$ \blacktriangle\,2.9\times $}}\\[-.2pt]9.7s\\[-.2pt]{\fontsize{5.8}{6.4}\selectfont \textcolor{ForestGreen}{$ \blacktriangle\,3.6\times $}}\end{tabular} & \hspace*{1pt}\begin{tabular}[t]{@{}c@{}}\rule{0pt}{12pt}13.5s\\[-.2pt]{\fontsize{5.8}{6.4}\selectfont \textcolor{ForestGreen}{$ \blacktriangle\,2.8\times $}}\\[-.2pt]9.6s\\[-.2pt]{\fontsize{5.8}{6.4}\selectfont \textcolor{ForestGreen}{$ \blacktriangle\,3.7\times $}}\end{tabular} & \textcolor{gray}{\ensuremath{\cdot}} & \textcolor{gray}{\ensuremath{\cdot}} & \textcolor{gray}{\ensuremath{\cdot}} & \textcolor{gray}{\ensuremath{\cdot}} & \textcolor{gray}{\ensuremath{\cdot}} & \textcolor{gray}{\ensuremath{\cdot}} & \hspace*{1pt}\begin{tabular}[t]{@{}c@{}}\rule{0pt}{12pt}4m\,19.1s\\[-.2pt]{\fontsize{5.8}{6.4}\selectfont \textcolor{BrickRed}{$ \blacktriangledown\,1.5\times $}}\\[-.2pt]4m\,8.7s\\[-.2pt]{\fontsize{5.8}{6.4}\selectfont \textcolor{BrickRed}{$ \blacktriangledown\,1.4\times $}}\end{tabular} & \hspace*{1pt}\begin{tabular}[t]{@{}c@{}}\rule{0pt}{12pt}4m\,29.2s\\[-.2pt]{\fontsize{5.8}{6.4}\selectfont \textcolor{BrickRed}{$ \blacktriangledown\,1.6\times $}}\\[-.2pt]4m\,15s\\[-.2pt]{\fontsize{5.8}{6.4}\selectfont \textcolor{BrickRed}{$ \blacktriangledown\,1.5\times $}}\end{tabular} & \textcolor{gray}{\ensuremath{\cdot}} & \textcolor{gray}{\ensuremath{\cdot}} & \textcolor{gray}{\ensuremath{\cdot}} & \textcolor{gray}{\ensuremath{\cdot}} & \textcolor{gray}{\ensuremath{\cdot}} & \textcolor{gray}{\ensuremath{\cdot}}\\[3pt]\hline
\end{tabular}\par}
\smallskip{\footnotesize\noindent Each cell: A, comparison, B, comparison. Times cover 1,024 queries. $\cdot$ = no approved observation. ---\textsuperscript{g} = GPU unavailable, not executed. n/a = no measured comparison.\par
\noindent \textsuperscript{f}SIMD-ref request observed as scalar, compared only as that effective path. \textsuperscript{p}Accepted by explicit preservation after wrapper postflight failure, no run receipt claimed.\par
\noindent MPF/MPS = minimizers per fragment/sequence. Row order: index representation / minimizer, then chaining / base alignment. SIMD-ref is the reference minimizer, Bounded is bounded ntHash.\par
\noindent Blue = one lowest A and one lowest B down each machine/workload column, across continuation pages. Exact ties select the first row. Repeated baselines are marked only on page 1.\par
\noindent Light gray = minimap2 baseline. Pale yellow = preferred MPF / Bounded / SIMD-DP / Block configuration, not a claim of winning every column.\par
\noindent When all Vulkan modes are unmeasured, they share one row. Baseline comparisons use the fastest measured BioseqDB effective path separately for A and B.\par}
\clearpage
\noindent{\large\bfseries AWS: 4 Gbp, k=15, w=10 (mapping interval)\hfill 5/6}\par\smallskip
{\fontsize{7.1}{8.1}\selectfont\setlength{\tabcolsep}{1pt}\setlength{\fboxsep}{2pt}\renewcommand{\arraystretch}{1.12}
\noindent\begin{tabular}{@{}p{14mm}p{27mm}|p{13.55mm}|p{13.55mm}|p{13.55mm}|p{13.55mm}|p{13.55mm}|p{13.55mm}|p{13.55mm}|p{13.55mm}|p{13.55mm}|p{13.55mm}|p{13.55mm}|p{13.55mm}|p{13.55mm}|p{13.55mm}|p{13.55mm}|p{13.55mm}|@{}}\hline
Engine & Execution & \multicolumn{8}{c|}{\textbf{Index}} & \multicolumn{8}{c|}{\textbf{No index}}\\
 &  & \multicolumn{8}{c|}{\textbf{4 Gbp}} & \multicolumn{8}{c|}{\textbf{4 Gbp}}\\
 &  & \multicolumn{8}{c|}{\textbf{k=15, w=10}} & \multicolumn{8}{c|}{\textbf{k=15, w=10}}\\
 &  & \multicolumn{4}{c|}{\textbf{DNA4}} & \multicolumn{4}{c|}{\textbf{DNA5}} & \multicolumn{4}{c|}{\textbf{DNA4}} & \multicolumn{4}{c|}{\textbf{DNA5}}\\
 &  & \multicolumn{1}{c|}{\textbf{M7A}} & \multicolumn{1}{c|}{\textbf{M7I}} & \multicolumn{1}{c|}{\textbf{M4}} & \multicolumn{1}{c|}{\textbf{C7G}} & \multicolumn{1}{c|}{\textbf{M7A}} & \multicolumn{1}{c|}{\textbf{M7I}} & \multicolumn{1}{c|}{\textbf{M4}} & \multicolumn{1}{c|}{\textbf{C7G}} & \multicolumn{1}{c|}{\textbf{M7A}} & \multicolumn{1}{c|}{\textbf{M7I}} & \multicolumn{1}{c|}{\textbf{M4}} & \multicolumn{1}{c|}{\textbf{C7G}} & \multicolumn{1}{c|}{\textbf{M7A}} & \multicolumn{1}{c|}{\textbf{M7I}} & \multicolumn{1}{c|}{\textbf{M4}} & \multicolumn{1}{c|}{\textbf{C7G}}\\
\hline
\rowcolor{BenchBaseline}[0.5pt][0.5pt]\textbf{minimap2} & \textbf{Baseline} & \hspace*{1pt}\begin{tabular}[t]{@{}c@{}}\rule{0pt}{12pt}36.2s\\[-.2pt]{\fontsize{5.8}{6.4}\selectfont \textcolor{BrickRed}{$ \blacktriangledown\,6.9\times $}}\\[-.2pt]35s\\[-.2pt]{\fontsize{5.8}{6.4}\selectfont \textcolor{BrickRed}{$ \blacktriangledown\,8.0\times $}}\end{tabular} & \hspace*{1pt}\begin{tabular}[t]{@{}c@{}}\rule{0pt}{12pt}38.1s\\[-.2pt]{\fontsize{5.8}{6.4}\selectfont \textcolor{BrickRed}{$ \blacktriangledown\,7.8\times $}}\\[-.2pt]36s\\[-.2pt]{\fontsize{5.8}{6.4}\selectfont \textcolor{BrickRed}{$ \blacktriangledown\,9.3\times $}}\end{tabular} & \hspace*{1pt}\begin{tabular}[t]{@{}c@{}}\rule{0pt}{12pt}1m\,33.5s\\[-.2pt]{\fontsize{5.8}{6.4}\selectfont \textcolor{BrickRed}{$ \blacktriangledown\,10.8\times $}}\\[-.2pt]1m\,29.7s\\[-.2pt]{\fontsize{5.8}{6.4}\selectfont \textcolor{BrickRed}{$ \blacktriangledown\,13.7\times $}}\end{tabular} & \hspace*{1pt}\begin{tabular}[t]{@{}c@{}}\rule{0pt}{12pt}6m\,31.5s\\[-.2pt]{\fontsize{5.8}{6.4}\selectfont \textcolor{gray}{n/a}}\\[-.2pt]6m\,26.3s\\[-.2pt]{\fontsize{5.8}{6.4}\selectfont \textcolor{gray}{n/a}}\end{tabular} & \hspace*{1pt}\begin{tabular}[t]{@{}c@{}}\rule{0pt}{12pt}31.6s\\[-.2pt]{\fontsize{5.8}{6.4}\selectfont \textcolor{gray}{n/a}}\\[-.2pt]30.3s\\[-.2pt]{\fontsize{5.8}{6.4}\selectfont \textcolor{gray}{n/a}}\end{tabular} & \hspace*{1pt}\begin{tabular}[t]{@{}c@{}}\rule{0pt}{12pt}34.1s\\[-.2pt]{\fontsize{5.8}{6.4}\selectfont \textcolor{gray}{n/a}}\\[-.2pt]32s\\[-.2pt]{\fontsize{5.8}{6.4}\selectfont \textcolor{gray}{n/a}}\end{tabular} & \hspace*{1pt}\begin{tabular}[t]{@{}c@{}}\rule{0pt}{12pt}1m\,9.5s\\[-.2pt]{\fontsize{5.8}{6.4}\selectfont \textcolor{gray}{n/a}}\\[-.2pt]1m\,6.1s\\[-.2pt]{\fontsize{5.8}{6.4}\selectfont \textcolor{gray}{n/a}}\end{tabular} & \hspace*{1pt}\begin{tabular}[t]{@{}c@{}}\rule{0pt}{12pt}3m\,1s\\[-.2pt]{\fontsize{5.8}{6.4}\selectfont \textcolor{gray}{n/a}}\\[-.2pt]2m\,51.5s\\[-.2pt]{\fontsize{5.8}{6.4}\selectfont \textcolor{gray}{n/a}}\end{tabular} & \hspace*{1pt}\begin{tabular}[t]{@{}c@{}}\rule{0pt}{12pt}2m\,55.5s\\[-.2pt]{\fontsize{5.8}{6.4}\selectfont \textcolor{ForestGreen}{$ \blacktriangle\,20\% $}}\\[-.2pt]2m\,54s\\[-.2pt]{\fontsize{5.8}{6.4}\selectfont \textcolor{ForestGreen}{$ \blacktriangle\,17\% $}}\end{tabular} & \hspace*{1pt}\begin{tabular}[t]{@{}c@{}}\rule{0pt}{12pt}2m\,50.4s\\[-.2pt]{\fontsize{5.8}{6.4}\selectfont \textcolor{ForestGreen}{$ \blacktriangle\,23\% $}}\\[-.2pt]2m\,48.2s\\[-.2pt]{\fontsize{5.8}{6.4}\selectfont \textcolor{ForestGreen}{$ \blacktriangle\,20\% $}}\end{tabular} & \hspace*{1pt}\begin{tabular}[t]{@{}c@{}}\rule{0pt}{12pt}6m\,12.2s\\[-.2pt]{\fontsize{5.8}{6.4}\selectfont \textcolor{ForestGreen}{$ \blacktriangle\,9\% $}}\\[-.2pt]6m\,8.2s\\[-.2pt]{\fontsize{5.8}{6.4}\selectfont \textcolor{gray}{$ \blacktriangleright\,5\% $}}\end{tabular} & \hspace*{1pt}\begin{tabular}[t]{@{}c@{}}\rule{0pt}{12pt}44m\,20.5s\\[-.2pt]{\fontsize{5.8}{6.4}\selectfont \textcolor{gray}{n/a}}\\[-.2pt]44m\,15.6s\\[-.2pt]{\fontsize{5.8}{6.4}\selectfont \textcolor{gray}{n/a}}\end{tabular} & \hspace*{1pt}\begin{tabular}[t]{@{}c@{}}\rule{0pt}{12pt}1m\,47.1s\\[-.2pt]{\fontsize{5.8}{6.4}\selectfont \textcolor{ForestGreen}{$ \blacktriangle\,5.5\times $}}\\[-.2pt]1m\,45.8s\\[-.2pt]{\fontsize{5.8}{6.4}\selectfont \textcolor{ForestGreen}{$ \blacktriangle\,2.1\times $}}\end{tabular} & \hspace*{1pt}\begin{tabular}[t]{@{}c@{}}\rule{0pt}{12pt}1m\,47.2s\\[-.2pt]{\fontsize{5.8}{6.4}\selectfont \textcolor{ForestGreen}{$ \blacktriangle\,4.6\times $}}\\[-.2pt]1m\,45.1s\\[-.2pt]{\fontsize{5.8}{6.4}\selectfont \textcolor{ForestGreen}{$ \blacktriangle\,2.2\times $}}\end{tabular} & \hspace*{1pt}\begin{tabular}[t]{@{}c@{}}\rule{0pt}{12pt}3m\,40.3s\\[-.2pt]{\fontsize{5.8}{6.4}\selectfont \textcolor{ForestGreen}{$ \blacktriangle\,4.2\times $}}\\[-.2pt]3m\,36.6s\\[-.2pt]{\fontsize{5.8}{6.4}\selectfont \textcolor{ForestGreen}{$ \blacktriangle\,1.7\times $}}\end{tabular} & \hspace*{1pt}\begin{tabular}[t]{@{}c@{}}\rule{0pt}{12pt}29m\,43.4s\\[-.2pt]{\fontsize{5.8}{6.4}\selectfont \textcolor{gray}{n/a}}\\[-.2pt]29m\,38.5s\\[-.2pt]{\fontsize{5.8}{6.4}\selectfont \textcolor{gray}{n/a}}\end{tabular}\\[3pt]\hline
\textbf{BioseqDB} & \begin{tabular}[t]{@{}l@{}}\rule{0pt}{12pt}MPS / Bounded\\SIMD-DP / Block\end{tabular} & \hspace*{1pt}\begin{tabular}[t]{@{}c@{}}\rule{0pt}{12pt}11.1s\\[-.2pt]{\fontsize{5.8}{6.4}\selectfont \textcolor{ForestGreen}{$ \blacktriangle\,3.3\times $}}\\[-.2pt]10.5s\\[-.2pt]{\fontsize{5.8}{6.4}\selectfont \textcolor{ForestGreen}{$ \blacktriangle\,3.3\times $}}\end{tabular} & \hspace*{1pt}\begin{tabular}[t]{@{}c@{}}\rule{0pt}{12pt}11.4s\\[-.2pt]{\fontsize{5.8}{6.4}\selectfont \textcolor{ForestGreen}{$ \blacktriangle\,3.4\times $}}\\[-.2pt]10.5s\\[-.2pt]{\fontsize{5.8}{6.4}\selectfont \textcolor{ForestGreen}{$ \blacktriangle\,3.4\times $}}\end{tabular} & \textcolor{gray}{\ensuremath{\cdot}} & \textcolor{gray}{\ensuremath{\cdot}} & \textcolor{gray}{\ensuremath{\cdot}} & \textcolor{gray}{\ensuremath{\cdot}} & \textcolor{gray}{\ensuremath{\cdot}} & \textcolor{gray}{\ensuremath{\cdot}} & \hspace*{1pt}\begin{tabular}[t]{@{}c@{}}\rule{0pt}{12pt}4m\,19.3s\\[-.2pt]{\fontsize{5.8}{6.4}\selectfont \textcolor{BrickRed}{$ \blacktriangledown\,1.5\times $}}\\[-.2pt]4m\,10.3s\\[-.2pt]{\fontsize{5.8}{6.4}\selectfont \textcolor{BrickRed}{$ \blacktriangledown\,1.4\times $}}\end{tabular} & \hspace*{1pt}\begin{tabular}[t]{@{}c@{}}\rule{0pt}{12pt}4m\,29.2s\\[-.2pt]{\fontsize{5.8}{6.4}\selectfont \textcolor{BrickRed}{$ \blacktriangledown\,1.6\times $}}\\[-.2pt]4m\,17.1s\\[-.2pt]{\fontsize{5.8}{6.4}\selectfont \textcolor{BrickRed}{$ \blacktriangledown\,1.5\times $}}\end{tabular} & \textcolor{gray}{\ensuremath{\cdot}} & \textcolor{gray}{\ensuremath{\cdot}} & \textcolor{gray}{\ensuremath{\cdot}} & \textcolor{gray}{\ensuremath{\cdot}} & \textcolor{gray}{\ensuremath{\cdot}} & \textcolor{gray}{\ensuremath{\cdot}}\\[3pt]\hline
 & \begin{tabular}[t]{@{}l@{}}\rule{0pt}{12pt}MPS / Bounded\\SIMD-DP / KSW2\end{tabular} & \hspace*{1pt}\begin{tabular}[t]{@{}c@{}}\rule{0pt}{12pt}12.3s\\[-.2pt]{\fontsize{5.8}{6.4}\selectfont \textcolor{ForestGreen}{$ \blacktriangle\,3.0\times $}}\\[-.2pt]9.7s\\[-.2pt]{\fontsize{5.8}{6.4}\selectfont \textcolor{ForestGreen}{$ \blacktriangle\,3.6\times $}}\end{tabular} & \hspace*{1pt}\begin{tabular}[t]{@{}c@{}}\rule{0pt}{12pt}12.1s\\[-.2pt]{\fontsize{5.8}{6.4}\selectfont \textcolor{ForestGreen}{$ \blacktriangle\,3.1\times $}}\\[-.2pt]9.1s\\[-.2pt]{\fontsize{5.8}{6.4}\selectfont \textcolor{ForestGreen}{$ \blacktriangle\,4.0\times $}}\end{tabular} & \textcolor{gray}{\ensuremath{\cdot}} & \textcolor{gray}{\ensuremath{\cdot}} & \textcolor{gray}{\ensuremath{\cdot}} & \textcolor{gray}{\ensuremath{\cdot}} & \textcolor{gray}{\ensuremath{\cdot}} & \textcolor{gray}{\ensuremath{\cdot}} & \hspace*{1pt}\begin{tabular}[t]{@{}c@{}}\rule{0pt}{12pt}4m\,19.7s\\[-.2pt]{\fontsize{5.8}{6.4}\selectfont \textcolor{BrickRed}{$ \blacktriangledown\,1.5\times $}}\\[-.2pt]4m\,8.9s\\[-.2pt]{\fontsize{5.8}{6.4}\selectfont \textcolor{BrickRed}{$ \blacktriangledown\,1.4\times $}}\end{tabular} & \hspace*{1pt}\begin{tabular}[t]{@{}c@{}}\rule{0pt}{12pt}3m\,58.7s\\[-.2pt]{\fontsize{5.8}{6.4}\selectfont \textcolor{BrickRed}{$ \blacktriangledown\,1.4\times $}}\\[-.2pt]3m\,45.4s\\[-.2pt]{\fontsize{5.8}{6.4}\selectfont \textcolor{BrickRed}{$ \blacktriangledown\,1.3\times $}}\end{tabular} & \hspace*{1pt}\begin{tabular}[t]{@{}c@{}}\rule{0pt}{12pt}10m\,17.4s\\[-.2pt]{\fontsize{5.8}{6.4}\selectfont \textcolor{BrickRed}{$ \blacktriangledown\,1.7\times $}}\\[-.2pt]9m\,56s\\[-.2pt]{\fontsize{5.8}{6.4}\selectfont \textcolor{BrickRed}{$ \blacktriangledown\,1.6\times $}}\end{tabular} & \textcolor{gray}{\ensuremath{\cdot}} & \textcolor{gray}{\ensuremath{\cdot}} & \textcolor{gray}{\ensuremath{\cdot}} & \textcolor{gray}{\ensuremath{\cdot}} & \textcolor{gray}{\ensuremath{\cdot}}\\[3pt]\hline
 & \begin{tabular}[t]{@{}l@{}}\rule{0pt}{12pt}MPS / Scalar\\RMQ / Block\end{tabular} & \hspace*{1pt}\begin{tabular}[t]{@{}c@{}}\rule{0pt}{12pt}11.5s\\[-.2pt]{\fontsize{5.8}{6.4}\selectfont \textcolor{ForestGreen}{$ \blacktriangle\,3.2\times $}}\\[-.2pt]10.9s\\[-.2pt]{\fontsize{5.8}{6.4}\selectfont \textcolor{ForestGreen}{$ \blacktriangle\,3.2\times $}}\end{tabular} & \hspace*{1pt}\begin{tabular}[t]{@{}c@{}}\rule{0pt}{12pt}11.9s\\[-.2pt]{\fontsize{5.8}{6.4}\selectfont \textcolor{ForestGreen}{$ \blacktriangle\,3.2\times $}}\\[-.2pt]11s\\[-.2pt]{\fontsize{5.8}{6.4}\selectfont \textcolor{ForestGreen}{$ \blacktriangle\,3.3\times $}}\end{tabular} & \textcolor{gray}{\ensuremath{\cdot}} & \textcolor{gray}{\ensuremath{\cdot}} & \textcolor{gray}{\ensuremath{\cdot}} & \textcolor{gray}{\ensuremath{\cdot}} & \textcolor{gray}{\ensuremath{\cdot}} & \textcolor{gray}{\ensuremath{\cdot}} & \hspace*{1pt}\begin{tabular}[t]{@{}c@{}}\rule{0pt}{12pt}5m\,40.1s\\[-.2pt]{\fontsize{5.8}{6.4}\selectfont \textcolor{BrickRed}{$ \blacktriangledown\,1.9\times $}}\\[-.2pt]5m\,31s\\[-.2pt]{\fontsize{5.8}{6.4}\selectfont \textcolor{BrickRed}{$ \blacktriangledown\,1.9\times $}}\end{tabular} & \hspace*{1pt}\begin{tabular}[t]{@{}c@{}}\rule{0pt}{12pt}6m\,11.3s\\[-.2pt]{\fontsize{5.8}{6.4}\selectfont \textcolor{BrickRed}{$ \blacktriangledown\,2.2\times $}}\\[-.2pt]5m\,59.5s\\[-.2pt]{\fontsize{5.8}{6.4}\selectfont \textcolor{BrickRed}{$ \blacktriangledown\,2.1\times $}}\end{tabular} & \textcolor{gray}{\ensuremath{\cdot}} & \textcolor{gray}{\ensuremath{\cdot}} & \textcolor{gray}{\ensuremath{\cdot}} & \textcolor{gray}{\ensuremath{\cdot}} & \textcolor{gray}{\ensuremath{\cdot}} & \textcolor{gray}{\ensuremath{\cdot}}\\[3pt]\hline
 & \begin{tabular}[t]{@{}l@{}}\rule{0pt}{12pt}MPS / Scalar\\RMQ / KSW2\end{tabular} & \hspace*{1pt}\begin{tabular}[t]{@{}c@{}}\rule{0pt}{12pt}13.3s\\[-.2pt]{\fontsize{5.8}{6.4}\selectfont \textcolor{ForestGreen}{$ \blacktriangle\,2.7\times $}}\\[-.2pt]10.8s\\[-.2pt]{\fontsize{5.8}{6.4}\selectfont \textcolor{ForestGreen}{$ \blacktriangle\,3.3\times $}}\end{tabular} & \textcolor{gray}{\ensuremath{\cdot}} & \textcolor{gray}{\ensuremath{\cdot}} & \textcolor{gray}{\ensuremath{\cdot}} & \textcolor{gray}{\ensuremath{\cdot}} & \textcolor{gray}{\ensuremath{\cdot}} & \textcolor{gray}{\ensuremath{\cdot}} & \textcolor{gray}{\ensuremath{\cdot}} & \textcolor{gray}{\ensuremath{\cdot}} & \textcolor{gray}{\ensuremath{\cdot}} & \textcolor{gray}{\ensuremath{\cdot}} & \textcolor{gray}{\ensuremath{\cdot}} & \textcolor{gray}{\ensuremath{\cdot}} & \textcolor{gray}{\ensuremath{\cdot}} & \textcolor{gray}{\ensuremath{\cdot}} & \textcolor{gray}{\ensuremath{\cdot}}\\[3pt]\hline
 & \begin{tabular}[t]{@{}l@{}}\rule{0pt}{12pt}MPS / Scalar\\Scalar-DP / Block\end{tabular} & \textcolor{gray}{\ensuremath{\cdot}} & \textcolor{gray}{\ensuremath{\cdot}} & \textcolor{gray}{\ensuremath{\cdot}} & \textcolor{gray}{\ensuremath{\cdot}} & \textcolor{gray}{\ensuremath{\cdot}} & \textcolor{gray}{\ensuremath{\cdot}} & \textcolor{gray}{\ensuremath{\cdot}} & \textcolor{gray}{\ensuremath{\cdot}} & \textcolor{gray}{\ensuremath{\cdot}} & \textcolor{gray}{\ensuremath{\cdot}} & \textcolor{gray}{\ensuremath{\cdot}} & \textcolor{gray}{\ensuremath{\cdot}} & \hspace*{1pt}\begin{tabular}[t]{@{}c@{}}\rule{0pt}{12pt}11m\,44.1s\\[-.2pt]{\fontsize{5.8}{6.4}\selectfont \textcolor{BrickRed}{$ \blacktriangledown\,6.6\times $}}\\[-.2pt]5m\,20.6s\\[-.2pt]{\fontsize{5.8}{6.4}\selectfont \textcolor{BrickRed}{$ \blacktriangledown\,3.0\times $}}\end{tabular} & \hspace*{1pt}\begin{tabular}[t]{@{}c@{}}\rule{0pt}{12pt}10m\,13.1s\\[-.2pt]{\fontsize{5.8}{6.4}\selectfont \textcolor{BrickRed}{$ \blacktriangledown\,5.7\times $}}\\[-.2pt]5m\,32.9s\\[-.2pt]{\fontsize{5.8}{6.4}\selectfont \textcolor{BrickRed}{$ \blacktriangledown\,3.2\times $}}\end{tabular} & \textcolor{gray}{\ensuremath{\cdot}} & \textcolor{gray}{\ensuremath{\cdot}}\\[3pt]\hline
 & \begin{tabular}[t]{@{}l@{}}\rule{0pt}{12pt}MPS / Scalar\\SIMD-DP / Block\end{tabular} & \hspace*{1pt}\begin{tabular}[t]{@{}c@{}}\rule{0pt}{12pt}11.3s\\[-.2pt]{\fontsize{5.8}{6.4}\selectfont \textcolor{ForestGreen}{$ \blacktriangle\,3.2\times $}}\\[-.2pt]10.7s\\[-.2pt]{\fontsize{5.8}{6.4}\selectfont \textcolor{ForestGreen}{$ \blacktriangle\,3.3\times $}}\end{tabular} & \hspace*{1pt}\begin{tabular}[t]{@{}c@{}}\rule{0pt}{12pt}11.4s\\[-.2pt]{\fontsize{5.8}{6.4}\selectfont \textcolor{ForestGreen}{$ \blacktriangle\,3.3\times $}}\\[-.2pt]10.7s\\[-.2pt]{\fontsize{5.8}{6.4}\selectfont \textcolor{ForestGreen}{$ \blacktriangle\,3.4\times $}}\end{tabular} & \textcolor{gray}{\ensuremath{\cdot}} & \textcolor{gray}{\ensuremath{\cdot}} & \textcolor{gray}{\ensuremath{\cdot}} & \textcolor{gray}{\ensuremath{\cdot}} & \textcolor{gray}{\ensuremath{\cdot}} & \textcolor{gray}{\ensuremath{\cdot}} & \textcolor{gray}{\ensuremath{\cdot}} & \hspace*{1pt}\begin{tabular}[t]{@{}c@{}}\rule{0pt}{12pt}5m\,51.7s\\[-.2pt]{\fontsize{5.8}{6.4}\selectfont \textcolor{BrickRed}{$ \blacktriangledown\,2.1\times $}}\\[-.2pt]5m\,39.9s\\[-.2pt]{\fontsize{5.8}{6.4}\selectfont \textcolor{BrickRed}{$ \blacktriangledown\,2.0\times $}}\end{tabular} & \textcolor{gray}{\ensuremath{\cdot}} & \textcolor{gray}{\ensuremath{\cdot}} & \hspace*{1pt}\begin{tabular}[t]{@{}c@{}}\rule{0pt}{12pt}11m\,45.1s\\[-.2pt]{\fontsize{5.8}{6.4}\selectfont \textcolor{BrickRed}{$ \blacktriangledown\,6.6\times $}}\\[-.2pt]5m\,20.1s\\[-.2pt]{\fontsize{5.8}{6.4}\selectfont \textcolor{BrickRed}{$ \blacktriangledown\,3.0\times $}}\end{tabular} & \hspace*{1pt}\begin{tabular}[t]{@{}c@{}}\rule{0pt}{12pt}10m\,10.3s\\[-.2pt]{\fontsize{5.8}{6.4}\selectfont \textcolor{BrickRed}{$ \blacktriangledown\,5.7\times $}}\\[-.2pt]5m\,30.2s\\[-.2pt]{\fontsize{5.8}{6.4}\selectfont \textcolor{BrickRed}{$ \blacktriangledown\,3.1\times $}}\end{tabular} & \textcolor{gray}{\ensuremath{\cdot}} & \textcolor{gray}{\ensuremath{\cdot}}\\[3pt]\hline
 & \begin{tabular}[t]{@{}l@{}}\rule{0pt}{12pt}MPS / Scalar\\SIMD-DP / KSW2\end{tabular} & \hspace*{1pt}\begin{tabular}[t]{@{}c@{}}\rule{0pt}{12pt}12.7s\\[-.2pt]{\fontsize{5.8}{6.4}\selectfont \textcolor{ForestGreen}{$ \blacktriangle\,2.9\times $}}\\[-.2pt]9.9s\\[-.2pt]{\fontsize{5.8}{6.4}\selectfont \textcolor{ForestGreen}{$ \blacktriangle\,3.6\times $}}\end{tabular} & \hspace*{1pt}\begin{tabular}[t]{@{}c@{}}\rule{0pt}{12pt}13.2s\\[-.2pt]{\fontsize{5.8}{6.4}\selectfont \textcolor{ForestGreen}{$ \blacktriangle\,2.9\times $}}\\[-.2pt]9.7s\\[-.2pt]{\fontsize{5.8}{6.4}\selectfont \textcolor{ForestGreen}{$ \blacktriangle\,3.7\times $}}\end{tabular} & \textcolor{gray}{\ensuremath{\cdot}} & \textcolor{gray}{\ensuremath{\cdot}} & \textcolor{gray}{\ensuremath{\cdot}} & \textcolor{gray}{\ensuremath{\cdot}} & \textcolor{gray}{\ensuremath{\cdot}} & \textcolor{gray}{\ensuremath{\cdot}} & \hspace*{1pt}\begin{tabular}[t]{@{}c@{}}\rule{0pt}{12pt}5m\,39.7s\\[-.2pt]{\fontsize{5.8}{6.4}\selectfont \textcolor{BrickRed}{$ \blacktriangledown\,1.9\times $}}\\[-.2pt]5m\,28.2s\\[-.2pt]{\fontsize{5.8}{6.4}\selectfont \textcolor{BrickRed}{$ \blacktriangledown\,1.9\times $}}\end{tabular} & \hspace*{1pt}\begin{tabular}[t]{@{}c@{}}\rule{0pt}{12pt}5m\,42.5s\\[-.2pt]{\fontsize{5.8}{6.4}\selectfont \textcolor{BrickRed}{$ \blacktriangledown\,2.0\times $}}\\[-.2pt]5m\,29s\\[-.2pt]{\fontsize{5.8}{6.4}\selectfont \textcolor{BrickRed}{$ \blacktriangledown\,2.0\times $}}\end{tabular} & \textcolor{gray}{\ensuremath{\cdot}} & \textcolor{gray}{\ensuremath{\cdot}} & \hspace*{1pt}\begin{tabular}[t]{@{}c@{}}\rule{0pt}{12pt}11m\,45s\\[-.2pt]{\fontsize{5.8}{6.4}\selectfont \textcolor{BrickRed}{$ \blacktriangledown\,6.6\times $}}\\[-.2pt]5m\,19.1s\\[-.2pt]{\fontsize{5.8}{6.4}\selectfont \textcolor{BrickRed}{$ \blacktriangledown\,3.0\times $}}\end{tabular} & \hspace*{1pt}\begin{tabular}[t]{@{}c@{}}\rule{0pt}{12pt}10m\,13.5s\\[-.2pt]{\fontsize{5.8}{6.4}\selectfont \textcolor{BrickRed}{$ \blacktriangledown\,5.7\times $}}\\[-.2pt]5m\,31.5s\\[-.2pt]{\fontsize{5.8}{6.4}\selectfont \textcolor{BrickRed}{$ \blacktriangledown\,3.2\times $}}\end{tabular} & \textcolor{gray}{\ensuremath{\cdot}} & \textcolor{gray}{\ensuremath{\cdot}}\\[3pt]\hline
\end{tabular}\par}
\smallskip{\footnotesize\noindent Each cell: A, comparison, B, comparison. Times cover 1,024 queries. $\cdot$ = no approved observation. ---\textsuperscript{g} = GPU unavailable, not executed. n/a = no measured comparison.\par
\noindent \textsuperscript{f}SIMD-ref request observed as scalar, compared only as that effective path. \textsuperscript{p}Accepted by explicit preservation after wrapper postflight failure, no run receipt claimed.\par
\noindent MPF/MPS = minimizers per fragment/sequence. Row order: index representation / minimizer, then chaining / base alignment. SIMD-ref is the reference minimizer, Bounded is bounded ntHash.\par
\noindent Blue = one lowest A and one lowest B down each machine/workload column, across continuation pages. Exact ties select the first row. Repeated baselines are marked only on page 1.\par
\noindent Light gray = minimap2 baseline. Pale yellow = preferred MPF / Bounded / SIMD-DP / Block configuration, not a claim of winning every column.\par
\noindent When all Vulkan modes are unmeasured, they share one row. Baseline comparisons use the fastest measured BioseqDB effective path separately for A and B.\par}
\clearpage
\noindent{\large\bfseries AWS: 4 Gbp, k=15, w=10 (mapping interval)\hfill 6/6}\par\smallskip
{\fontsize{7.1}{8.1}\selectfont\setlength{\tabcolsep}{1pt}\setlength{\fboxsep}{2pt}\renewcommand{\arraystretch}{1.12}
\noindent\begin{tabular}{@{}p{14mm}p{27mm}|p{13.55mm}|p{13.55mm}|p{13.55mm}|p{13.55mm}|p{13.55mm}|p{13.55mm}|p{13.55mm}|p{13.55mm}|p{13.55mm}|p{13.55mm}|p{13.55mm}|p{13.55mm}|p{13.55mm}|p{13.55mm}|p{13.55mm}|p{13.55mm}|@{}}\hline
Engine & Execution & \multicolumn{8}{c|}{\textbf{Index}} & \multicolumn{8}{c|}{\textbf{No index}}\\
 &  & \multicolumn{8}{c|}{\textbf{4 Gbp}} & \multicolumn{8}{c|}{\textbf{4 Gbp}}\\
 &  & \multicolumn{8}{c|}{\textbf{k=15, w=10}} & \multicolumn{8}{c|}{\textbf{k=15, w=10}}\\
 &  & \multicolumn{4}{c|}{\textbf{DNA4}} & \multicolumn{4}{c|}{\textbf{DNA5}} & \multicolumn{4}{c|}{\textbf{DNA4}} & \multicolumn{4}{c|}{\textbf{DNA5}}\\
 &  & \multicolumn{1}{c|}{\textbf{M7A}} & \multicolumn{1}{c|}{\textbf{M7I}} & \multicolumn{1}{c|}{\textbf{M4}} & \multicolumn{1}{c|}{\textbf{C7G}} & \multicolumn{1}{c|}{\textbf{M7A}} & \multicolumn{1}{c|}{\textbf{M7I}} & \multicolumn{1}{c|}{\textbf{M4}} & \multicolumn{1}{c|}{\textbf{C7G}} & \multicolumn{1}{c|}{\textbf{M7A}} & \multicolumn{1}{c|}{\textbf{M7I}} & \multicolumn{1}{c|}{\textbf{M4}} & \multicolumn{1}{c|}{\textbf{C7G}} & \multicolumn{1}{c|}{\textbf{M7A}} & \multicolumn{1}{c|}{\textbf{M7I}} & \multicolumn{1}{c|}{\textbf{M4}} & \multicolumn{1}{c|}{\textbf{C7G}}\\
\hline
\rowcolor{BenchBaseline}[0.5pt][0.5pt]\textbf{minimap2} & \textbf{Baseline} & \hspace*{1pt}\begin{tabular}[t]{@{}c@{}}\rule{0pt}{12pt}36.2s\\[-.2pt]{\fontsize{5.8}{6.4}\selectfont \textcolor{BrickRed}{$ \blacktriangledown\,6.9\times $}}\\[-.2pt]35s\\[-.2pt]{\fontsize{5.8}{6.4}\selectfont \textcolor{BrickRed}{$ \blacktriangledown\,8.0\times $}}\end{tabular} & \hspace*{1pt}\begin{tabular}[t]{@{}c@{}}\rule{0pt}{12pt}38.1s\\[-.2pt]{\fontsize{5.8}{6.4}\selectfont \textcolor{BrickRed}{$ \blacktriangledown\,7.8\times $}}\\[-.2pt]36s\\[-.2pt]{\fontsize{5.8}{6.4}\selectfont \textcolor{BrickRed}{$ \blacktriangledown\,9.3\times $}}\end{tabular} & \hspace*{1pt}\begin{tabular}[t]{@{}c@{}}\rule{0pt}{12pt}1m\,33.5s\\[-.2pt]{\fontsize{5.8}{6.4}\selectfont \textcolor{BrickRed}{$ \blacktriangledown\,10.8\times $}}\\[-.2pt]1m\,29.7s\\[-.2pt]{\fontsize{5.8}{6.4}\selectfont \textcolor{BrickRed}{$ \blacktriangledown\,13.7\times $}}\end{tabular} & \hspace*{1pt}\begin{tabular}[t]{@{}c@{}}\rule{0pt}{12pt}6m\,31.5s\\[-.2pt]{\fontsize{5.8}{6.4}\selectfont \textcolor{gray}{n/a}}\\[-.2pt]6m\,26.3s\\[-.2pt]{\fontsize{5.8}{6.4}\selectfont \textcolor{gray}{n/a}}\end{tabular} & \hspace*{1pt}\begin{tabular}[t]{@{}c@{}}\rule{0pt}{12pt}31.6s\\[-.2pt]{\fontsize{5.8}{6.4}\selectfont \textcolor{gray}{n/a}}\\[-.2pt]30.3s\\[-.2pt]{\fontsize{5.8}{6.4}\selectfont \textcolor{gray}{n/a}}\end{tabular} & \hspace*{1pt}\begin{tabular}[t]{@{}c@{}}\rule{0pt}{12pt}34.1s\\[-.2pt]{\fontsize{5.8}{6.4}\selectfont \textcolor{gray}{n/a}}\\[-.2pt]32s\\[-.2pt]{\fontsize{5.8}{6.4}\selectfont \textcolor{gray}{n/a}}\end{tabular} & \hspace*{1pt}\begin{tabular}[t]{@{}c@{}}\rule{0pt}{12pt}1m\,9.5s\\[-.2pt]{\fontsize{5.8}{6.4}\selectfont \textcolor{gray}{n/a}}\\[-.2pt]1m\,6.1s\\[-.2pt]{\fontsize{5.8}{6.4}\selectfont \textcolor{gray}{n/a}}\end{tabular} & \hspace*{1pt}\begin{tabular}[t]{@{}c@{}}\rule{0pt}{12pt}3m\,1s\\[-.2pt]{\fontsize{5.8}{6.4}\selectfont \textcolor{gray}{n/a}}\\[-.2pt]2m\,51.5s\\[-.2pt]{\fontsize{5.8}{6.4}\selectfont \textcolor{gray}{n/a}}\end{tabular} & \hspace*{1pt}\begin{tabular}[t]{@{}c@{}}\rule{0pt}{12pt}2m\,55.5s\\[-.2pt]{\fontsize{5.8}{6.4}\selectfont \textcolor{ForestGreen}{$ \blacktriangle\,20\% $}}\\[-.2pt]2m\,54s\\[-.2pt]{\fontsize{5.8}{6.4}\selectfont \textcolor{ForestGreen}{$ \blacktriangle\,17\% $}}\end{tabular} & \hspace*{1pt}\begin{tabular}[t]{@{}c@{}}\rule{0pt}{12pt}2m\,50.4s\\[-.2pt]{\fontsize{5.8}{6.4}\selectfont \textcolor{ForestGreen}{$ \blacktriangle\,23\% $}}\\[-.2pt]2m\,48.2s\\[-.2pt]{\fontsize{5.8}{6.4}\selectfont \textcolor{ForestGreen}{$ \blacktriangle\,20\% $}}\end{tabular} & \hspace*{1pt}\begin{tabular}[t]{@{}c@{}}\rule{0pt}{12pt}6m\,12.2s\\[-.2pt]{\fontsize{5.8}{6.4}\selectfont \textcolor{ForestGreen}{$ \blacktriangle\,9\% $}}\\[-.2pt]6m\,8.2s\\[-.2pt]{\fontsize{5.8}{6.4}\selectfont \textcolor{gray}{$ \blacktriangleright\,5\% $}}\end{tabular} & \hspace*{1pt}\begin{tabular}[t]{@{}c@{}}\rule{0pt}{12pt}44m\,20.5s\\[-.2pt]{\fontsize{5.8}{6.4}\selectfont \textcolor{gray}{n/a}}\\[-.2pt]44m\,15.6s\\[-.2pt]{\fontsize{5.8}{6.4}\selectfont \textcolor{gray}{n/a}}\end{tabular} & \hspace*{1pt}\begin{tabular}[t]{@{}c@{}}\rule{0pt}{12pt}1m\,47.1s\\[-.2pt]{\fontsize{5.8}{6.4}\selectfont \textcolor{ForestGreen}{$ \blacktriangle\,5.5\times $}}\\[-.2pt]1m\,45.8s\\[-.2pt]{\fontsize{5.8}{6.4}\selectfont \textcolor{ForestGreen}{$ \blacktriangle\,2.1\times $}}\end{tabular} & \hspace*{1pt}\begin{tabular}[t]{@{}c@{}}\rule{0pt}{12pt}1m\,47.2s\\[-.2pt]{\fontsize{5.8}{6.4}\selectfont \textcolor{ForestGreen}{$ \blacktriangle\,4.6\times $}}\\[-.2pt]1m\,45.1s\\[-.2pt]{\fontsize{5.8}{6.4}\selectfont \textcolor{ForestGreen}{$ \blacktriangle\,2.2\times $}}\end{tabular} & \hspace*{1pt}\begin{tabular}[t]{@{}c@{}}\rule{0pt}{12pt}3m\,40.3s\\[-.2pt]{\fontsize{5.8}{6.4}\selectfont \textcolor{ForestGreen}{$ \blacktriangle\,4.2\times $}}\\[-.2pt]3m\,36.6s\\[-.2pt]{\fontsize{5.8}{6.4}\selectfont \textcolor{ForestGreen}{$ \blacktriangle\,1.7\times $}}\end{tabular} & \hspace*{1pt}\begin{tabular}[t]{@{}c@{}}\rule{0pt}{12pt}29m\,43.4s\\[-.2pt]{\fontsize{5.8}{6.4}\selectfont \textcolor{gray}{n/a}}\\[-.2pt]29m\,38.5s\\[-.2pt]{\fontsize{5.8}{6.4}\selectfont \textcolor{gray}{n/a}}\end{tabular}\\[3pt]\hline
\textbf{BioseqDB} & \begin{tabular}[t]{@{}l@{}}\rule{0pt}{12pt}MPS / SIMD-ref\\SIMD-DP / Block\end{tabular} & \textcolor{gray}{\ensuremath{\cdot}} & \textcolor{gray}{\ensuremath{\cdot}} & \textcolor{gray}{\ensuremath{\cdot}} & \textcolor{gray}{\ensuremath{\cdot}} & \textcolor{gray}{\ensuremath{\cdot}} & \textcolor{gray}{\ensuremath{\cdot}} & \textcolor{gray}{\ensuremath{\cdot}} & \textcolor{gray}{\ensuremath{\cdot}} & \textcolor{gray}{\ensuremath{\cdot}} & \textcolor{gray}{\ensuremath{\cdot}} & \hspace*{1pt}\begin{tabular}[t]{@{}c@{}}\rule{0pt}{12pt}8m\,12.5s\\[-.2pt]{\fontsize{5.8}{6.4}\selectfont \textcolor{BrickRed}{$ \blacktriangledown\,1.3\times $}}\\[-.2pt]7m\,55.2s\\[-.2pt]{\fontsize{5.8}{6.4}\selectfont \textcolor{BrickRed}{$ \blacktriangledown\,29\% $}}\end{tabular} & \textcolor{gray}{\ensuremath{\cdot}} & \textcolor{gray}{\ensuremath{\cdot}} & \textcolor{gray}{\ensuremath{\cdot}} & \textcolor{gray}{\ensuremath{\cdot}} & \textcolor{gray}{\ensuremath{\cdot}}\\[3pt]\hline
 & \begin{tabular}[t]{@{}l@{}}\rule{0pt}{12pt}MPS / Two-stack\\RMQ / Block\end{tabular} & \hspace*{1pt}\begin{tabular}[t]{@{}c@{}}\rule{0pt}{12pt}11.4s\\[-.2pt]{\fontsize{5.8}{6.4}\selectfont \textcolor{ForestGreen}{$ \blacktriangle\,3.2\times $}}\\[-.2pt]10.9s\\[-.2pt]{\fontsize{5.8}{6.4}\selectfont \textcolor{ForestGreen}{$ \blacktriangle\,3.2\times $}}\end{tabular} & \hspace*{1pt}\begin{tabular}[t]{@{}c@{}}\rule{0pt}{12pt}11.8s\\[-.2pt]{\fontsize{5.8}{6.4}\selectfont \textcolor{ForestGreen}{$ \blacktriangle\,3.2\times $}}\\[-.2pt]10.9s\\[-.2pt]{\fontsize{5.8}{6.4}\selectfont \textcolor{ForestGreen}{$ \blacktriangle\,3.3\times $}}\end{tabular} & \textcolor{gray}{\ensuremath{\cdot}} & \textcolor{gray}{\ensuremath{\cdot}} & \textcolor{gray}{\ensuremath{\cdot}} & \textcolor{gray}{\ensuremath{\cdot}} & \textcolor{gray}{\ensuremath{\cdot}} & \textcolor{gray}{\ensuremath{\cdot}} & \hspace*{1pt}\begin{tabular}[t]{@{}c@{}}\rule{0pt}{12pt}4m\,44.1s\\[-.2pt]{\fontsize{5.8}{6.4}\selectfont \textcolor{BrickRed}{$ \blacktriangledown\,1.6\times $}}\\[-.2pt]4m\,35.1s\\[-.2pt]{\fontsize{5.8}{6.4}\selectfont \textcolor{BrickRed}{$ \blacktriangledown\,1.6\times $}}\end{tabular} & \hspace*{1pt}\begin{tabular}[t]{@{}c@{}}\rule{0pt}{12pt}4m\,42.8s\\[-.2pt]{\fontsize{5.8}{6.4}\selectfont \textcolor{BrickRed}{$ \blacktriangledown\,1.7\times $}}\\[-.2pt]4m\,30.9s\\[-.2pt]{\fontsize{5.8}{6.4}\selectfont \textcolor{BrickRed}{$ \blacktriangledown\,1.6\times $}}\end{tabular} & \textcolor{gray}{\ensuremath{\cdot}} & \textcolor{gray}{\ensuremath{\cdot}} & \textcolor{gray}{\ensuremath{\cdot}} & \textcolor{gray}{\ensuremath{\cdot}} & \textcolor{gray}{\ensuremath{\cdot}} & \textcolor{gray}{\ensuremath{\cdot}}\\[3pt]\hline
 & \begin{tabular}[t]{@{}l@{}}\rule{0pt}{12pt}MPS / Two-stack\\RMQ / KSW2\end{tabular} & \hspace*{1pt}\begin{tabular}[t]{@{}c@{}}\rule{0pt}{12pt}13.3s\\[-.2pt]{\fontsize{5.8}{6.4}\selectfont \textcolor{ForestGreen}{$ \blacktriangle\,2.7\times $}}\\[-.2pt]10.8s\\[-.2pt]{\fontsize{5.8}{6.4}\selectfont \textcolor{ForestGreen}{$ \blacktriangle\,3.2\times $}}\end{tabular} & \textcolor{gray}{\ensuremath{\cdot}} & \textcolor{gray}{\ensuremath{\cdot}} & \textcolor{gray}{\ensuremath{\cdot}} & \textcolor{gray}{\ensuremath{\cdot}} & \textcolor{gray}{\ensuremath{\cdot}} & \textcolor{gray}{\ensuremath{\cdot}} & \textcolor{gray}{\ensuremath{\cdot}} & \textcolor{gray}{\ensuremath{\cdot}} & \textcolor{gray}{\ensuremath{\cdot}} & \textcolor{gray}{\ensuremath{\cdot}} & \textcolor{gray}{\ensuremath{\cdot}} & \textcolor{gray}{\ensuremath{\cdot}} & \textcolor{gray}{\ensuremath{\cdot}} & \textcolor{gray}{\ensuremath{\cdot}} & \textcolor{gray}{\ensuremath{\cdot}}\\[3pt]\hline
 & \begin{tabular}[t]{@{}l@{}}\rule{0pt}{12pt}MPS / Two-stack\\Scalar-DP / KSW2\end{tabular} & \textcolor{gray}{\ensuremath{\cdot}} & \textcolor{gray}{\ensuremath{\cdot}} & \textcolor{gray}{\ensuremath{\cdot}} & \textcolor{gray}{\ensuremath{\cdot}} & \textcolor{gray}{\ensuremath{\cdot}} & \textcolor{gray}{\ensuremath{\cdot}} & \textcolor{gray}{\ensuremath{\cdot}} & \textcolor{gray}{\ensuremath{\cdot}} & \textcolor{gray}{\ensuremath{\cdot}} & \textcolor{gray}{\ensuremath{\cdot}} & \textcolor{gray}{\ensuremath{\cdot}} & \textcolor{gray}{\ensuremath{\cdot}} & \hspace*{1pt}\begin{tabular}[t]{@{}c@{}}\rule{0pt}{12pt}10m\,24.5s\\[-.2pt]{\fontsize{5.8}{6.4}\selectfont \textcolor{BrickRed}{$ \blacktriangledown\,5.8\times $}}\\[-.2pt]4m\,4.1s\\[-.2pt]{\fontsize{5.8}{6.4}\selectfont \textcolor{BrickRed}{$ \blacktriangledown\,2.3\times $}}\end{tabular} & \hspace*{1pt}\begin{tabular}[t]{@{}c@{}}\rule{0pt}{12pt}8m\,30.6s\\[-.2pt]{\fontsize{5.8}{6.4}\selectfont \textcolor{BrickRed}{$ \blacktriangledown\,4.8\times $}}\\[-.2pt]3m\,55.3s\\[-.2pt]{\fontsize{5.8}{6.4}\selectfont \textcolor{BrickRed}{$ \blacktriangledown\,2.2\times $}}\end{tabular} & \textcolor{gray}{\ensuremath{\cdot}} & \textcolor{gray}{\ensuremath{\cdot}}\\[3pt]\hline
 & \begin{tabular}[t]{@{}l@{}}\rule{0pt}{12pt}MPS / Two-stack\\SIMD-DP / Block\end{tabular} & \hspace*{1pt}\begin{tabular}[t]{@{}c@{}}\rule{0pt}{12pt}11.5s\\[-.2pt]{\fontsize{5.8}{6.4}\selectfont \textcolor{ForestGreen}{$ \blacktriangle\,3.1\times $}}\\[-.2pt]10.9s\\[-.2pt]{\fontsize{5.8}{6.4}\selectfont \textcolor{ForestGreen}{$ \blacktriangle\,3.2\times $}}\end{tabular} & \hspace*{1pt}\begin{tabular}[t]{@{}c@{}}\rule{0pt}{12pt}11.2s\\[-.2pt]{\fontsize{5.8}{6.4}\selectfont \textcolor{ForestGreen}{$ \blacktriangle\,3.4\times $}}\\[-.2pt]10.4s\\[-.2pt]{\fontsize{5.8}{6.4}\selectfont \textcolor{ForestGreen}{$ \blacktriangle\,3.5\times $}}\end{tabular} & \textcolor{gray}{\ensuremath{\cdot}} & \textcolor{gray}{\ensuremath{\cdot}} & \textcolor{gray}{\ensuremath{\cdot}} & \textcolor{gray}{\ensuremath{\cdot}} & \textcolor{gray}{\ensuremath{\cdot}} & \textcolor{gray}{\ensuremath{\cdot}} & \hspace*{1pt}\begin{tabular}[t]{@{}c@{}}\rule{0pt}{12pt}4m\,39.5s\\[-.2pt]{\fontsize{5.8}{6.4}\selectfont \textcolor{BrickRed}{$ \blacktriangledown\,1.6\times $}}\\[-.2pt]4m\,30.1s\\[-.2pt]{\fontsize{5.8}{6.4}\selectfont \textcolor{BrickRed}{$ \blacktriangledown\,1.6\times $}}\end{tabular} & \hspace*{1pt}\begin{tabular}[t]{@{}c@{}}\rule{0pt}{12pt}4m\,24.7s\\[-.2pt]{\fontsize{5.8}{6.4}\selectfont \textcolor{BrickRed}{$ \blacktriangledown\,1.6\times $}}\\[-.2pt]4m\,13s\\[-.2pt]{\fontsize{5.8}{6.4}\selectfont \textcolor{BrickRed}{$ \blacktriangledown\,1.5\times $}}\end{tabular} & \textcolor{gray}{\ensuremath{\cdot}} & \textcolor{gray}{\ensuremath{\cdot}} & \textcolor{gray}{\ensuremath{\cdot}} & \textcolor{gray}{\ensuremath{\cdot}} & \textcolor{gray}{\ensuremath{\cdot}} & \textcolor{gray}{\ensuremath{\cdot}}\\[3pt]\hline
 & \begin{tabular}[t]{@{}l@{}}\rule{0pt}{12pt}MPS / Two-stack\\SIMD-DP / KSW2\end{tabular} & \hspace*{1pt}\begin{tabular}[t]{@{}c@{}}\rule{0pt}{12pt}13.1s\\[-.2pt]{\fontsize{5.8}{6.4}\selectfont \textcolor{ForestGreen}{$ \blacktriangle\,2.8\times $}}\\[-.2pt]10.5s\\[-.2pt]{\fontsize{5.8}{6.4}\selectfont \textcolor{ForestGreen}{$ \blacktriangle\,3.3\times $}}\end{tabular} & \hspace*{1pt}\begin{tabular}[t]{@{}c@{}}\rule{0pt}{12pt}13.4s\\[-.2pt]{\fontsize{5.8}{6.4}\selectfont \textcolor{ForestGreen}{$ \blacktriangle\,2.8\times $}}\\[-.2pt]9.7s\\[-.2pt]{\fontsize{5.8}{6.4}\selectfont \textcolor{ForestGreen}{$ \blacktriangle\,3.7\times $}}\end{tabular} & \textcolor{gray}{\ensuremath{\cdot}} & \textcolor{gray}{\ensuremath{\cdot}} & \textcolor{gray}{\ensuremath{\cdot}} & \textcolor{gray}{\ensuremath{\cdot}} & \textcolor{gray}{\ensuremath{\cdot}} & \textcolor{gray}{\ensuremath{\cdot}} & \hspace*{1pt}\begin{tabular}[t]{@{}c@{}}\rule{0pt}{12pt}4m\,54.4s\\[-.2pt]{\fontsize{5.8}{6.4}\selectfont \textcolor{BrickRed}{$ \blacktriangledown\,1.7\times $}}\\[-.2pt]4m\,44.1s\\[-.2pt]{\fontsize{5.8}{6.4}\selectfont \textcolor{BrickRed}{$ \blacktriangledown\,1.6\times $}}\end{tabular} & \hspace*{1pt}\begin{tabular}[t]{@{}c@{}}\rule{0pt}{12pt}4m\,39.9s\\[-.2pt]{\fontsize{5.8}{6.4}\selectfont \textcolor{BrickRed}{$ \blacktriangledown\,1.6\times $}}\\[-.2pt]4m\,26.4s\\[-.2pt]{\fontsize{5.8}{6.4}\selectfont \textcolor{BrickRed}{$ \blacktriangledown\,1.6\times $}}\end{tabular} & \textcolor{gray}{\ensuremath{\cdot}} & \textcolor{gray}{\ensuremath{\cdot}} & \textcolor{gray}{\ensuremath{\cdot}} & \textcolor{gray}{\ensuremath{\cdot}} & \textcolor{gray}{\ensuremath{\cdot}} & \textcolor{gray}{\ensuremath{\cdot}}\\[3pt]\hline
\rule{0pt}{12pt}\textbf{BioseqDB} & Vulkan (all modes) & \textcolor{gray}{---\textsuperscript{g}} & \textcolor{gray}{---\textsuperscript{g}} & \textcolor{gray}{---\textsuperscript{g}} & \textcolor{gray}{---\textsuperscript{g}} & \textcolor{gray}{---\textsuperscript{g}} & \textcolor{gray}{---\textsuperscript{g}} & \textcolor{gray}{---\textsuperscript{g}} & \textcolor{gray}{---\textsuperscript{g}} & \textcolor{gray}{---\textsuperscript{g}} & \textcolor{gray}{---\textsuperscript{g}} & \textcolor{gray}{---\textsuperscript{g}} & \textcolor{gray}{---\textsuperscript{g}} & \textcolor{gray}{---\textsuperscript{g}} & \textcolor{gray}{---\textsuperscript{g}} & \textcolor{gray}{---\textsuperscript{g}} & \textcolor{gray}{---\textsuperscript{g}}\\[3pt]\hline
\end{tabular}\par}
\smallskip{\footnotesize\noindent Each cell: A, comparison, B, comparison. Times cover 1,024 queries. $\cdot$ = no approved observation. ---\textsuperscript{g} = GPU unavailable, not executed. n/a = no measured comparison.\par
\noindent \textsuperscript{f}SIMD-ref request observed as scalar, compared only as that effective path. \textsuperscript{p}Accepted by explicit preservation after wrapper postflight failure, no run receipt claimed.\par
\noindent MPF/MPS = minimizers per fragment/sequence. Row order: index representation / minimizer, then chaining / base alignment. SIMD-ref is the reference minimizer, Bounded is bounded ntHash.\par
\noindent Blue = one lowest A and one lowest B down each machine/workload column, across continuation pages. Exact ties select the first row. Repeated baselines are marked only on page 1.\par
\noindent Light gray = minimap2 baseline. Pale yellow = preferred MPF / Bounded / SIMD-DP / Block configuration, not a claim of winning every column.\par
\noindent When all Vulkan modes are unmeasured, they share one row. Baseline comparisons use the fastest measured BioseqDB effective path separately for A and B.\par}
\clearpage
\noindent{\large\bfseries AWS: 4 Gbp, k=20, w=15 (mapping interval)\hfill 1/4}\par\smallskip
{\fontsize{7.1}{8.1}\selectfont\setlength{\tabcolsep}{1pt}\setlength{\fboxsep}{2pt}\renewcommand{\arraystretch}{1.12}
\noindent\begin{tabular}{@{}p{14mm}p{27mm}|p{13.55mm}|p{13.55mm}|p{13.55mm}|p{13.55mm}|p{13.55mm}|p{13.55mm}|p{13.55mm}|p{13.55mm}|p{13.55mm}|p{13.55mm}|p{13.55mm}|p{13.55mm}|p{13.55mm}|p{13.55mm}|p{13.55mm}|p{13.55mm}|@{}}\hline
Engine & Execution & \multicolumn{8}{c|}{\textbf{Index}} & \multicolumn{8}{c|}{\textbf{No index}}\\
 &  & \multicolumn{8}{c|}{\textbf{4 Gbp}} & \multicolumn{8}{c|}{\textbf{4 Gbp}}\\
 &  & \multicolumn{8}{c|}{\textbf{k=20, w=15}} & \multicolumn{8}{c|}{\textbf{k=20, w=15}}\\
 &  & \multicolumn{4}{c|}{\textbf{DNA4}} & \multicolumn{4}{c|}{\textbf{DNA5}} & \multicolumn{4}{c|}{\textbf{DNA4}} & \multicolumn{4}{c|}{\textbf{DNA5}}\\
 &  & \multicolumn{1}{c|}{\textbf{M7A}} & \multicolumn{1}{c|}{\textbf{M7I}} & \multicolumn{1}{c|}{\textbf{M4}} & \multicolumn{1}{c|}{\textbf{C7G}} & \multicolumn{1}{c|}{\textbf{M7A}} & \multicolumn{1}{c|}{\textbf{M7I}} & \multicolumn{1}{c|}{\textbf{M4}} & \multicolumn{1}{c|}{\textbf{C7G}} & \multicolumn{1}{c|}{\textbf{M7A}} & \multicolumn{1}{c|}{\textbf{M7I}} & \multicolumn{1}{c|}{\textbf{M4}} & \multicolumn{1}{c|}{\textbf{C7G}} & \multicolumn{1}{c|}{\textbf{M7A}} & \multicolumn{1}{c|}{\textbf{M7I}} & \multicolumn{1}{c|}{\textbf{M4}} & \multicolumn{1}{c|}{\textbf{C7G}}\\
\hline
\rowcolor{BenchBaseline}[0.5pt][0.5pt]\textbf{minimap2} & \textbf{Baseline} & \hspace*{1pt}\begin{tabular}[t]{@{}c@{}}\rule{0pt}{12pt}2m\,48.1s\\[-.2pt]{\fontsize{5.8}{6.4}\selectfont \textcolor{BrickRed}{$ \blacktriangledown\,20.4\times $}}\\[-.2pt]2m\,41.7s\\[-.2pt]{\fontsize{5.8}{6.4}\selectfont \textcolor{BrickRed}{$ \blacktriangledown\,22.6\times $}}\end{tabular} & \hspace*{1pt}\begin{tabular}[t]{@{}c@{}}\rule{0pt}{12pt}2m\,50.3s\\[-.2pt]{\fontsize{5.8}{6.4}\selectfont \textcolor{BrickRed}{$ \blacktriangledown\,19.8\times $}}\\[-.2pt]2m\,43s\\[-.2pt]{\fontsize{5.8}{6.4}\selectfont \textcolor{BrickRed}{$ \blacktriangledown\,22.2\times $}}\end{tabular} & \hspace*{1pt}\begin{tabular}[t]{@{}c@{}}\rule{0pt}{12pt}4m\,54.3s\\[-.2pt]{\fontsize{5.8}{6.4}\selectfont \textcolor{BrickRed}{$ \blacktriangledown\,43.2\times $}}\\[-.2pt]4m\,39s\\[-.2pt]{\fontsize{5.8}{6.4}\selectfont \textcolor{BrickRed}{$ \blacktriangledown\,54.6\times $}}\end{tabular} & \hspace*{1pt}\begin{tabular}[t]{@{}c@{}}\rule{0pt}{12pt}\colorbox{BenchBest}{\strut 25m\,5.5s}\\[-.2pt]{\fontsize{5.8}{6.4}\selectfont \textcolor{gray}{n/a}}\\[-.2pt]\colorbox{BenchBest}{\strut 25m\,0.7s}\\[-.2pt]{\fontsize{5.8}{6.4}\selectfont \textcolor{gray}{n/a}}\end{tabular} & \hspace*{1pt}\begin{tabular}[t]{@{}c@{}}\rule{0pt}{12pt}33.9s\\[-.2pt]{\fontsize{5.8}{6.4}\selectfont \textcolor{BrickRed}{$ \blacktriangledown\,1.9\times $}}\\[-.2pt]32.9s\\[-.2pt]{\fontsize{5.8}{6.4}\selectfont \textcolor{BrickRed}{$ \blacktriangledown\,2.1\times $}}\end{tabular} & \hspace*{1pt}\begin{tabular}[t]{@{}c@{}}\rule{0pt}{12pt}34.4s\\[-.2pt]{\fontsize{5.8}{6.4}\selectfont \textcolor{BrickRed}{$ \blacktriangledown\,2.6\times $}}\\[-.2pt]32.7s\\[-.2pt]{\fontsize{5.8}{6.4}\selectfont \textcolor{BrickRed}{$ \blacktriangledown\,2.9\times $}}\end{tabular} & \hspace*{1pt}\begin{tabular}[t]{@{}c@{}}\rule{0pt}{12pt}1m\,8.2s\\[-.2pt]{\fontsize{5.8}{6.4}\selectfont \textcolor{BrickRed}{$ \blacktriangledown\,3.7\times $}}\\[-.2pt]1m\,5.1s\\[-.2pt]{\fontsize{5.8}{6.4}\selectfont \textcolor{BrickRed}{$ \blacktriangledown\,4.6\times $}}\end{tabular} & \hspace*{1pt}\begin{tabular}[t]{@{}c@{}}\rule{0pt}{12pt}\colorbox{BenchBest}{\strut 3m\,9.8s}\\[-.2pt]{\fontsize{5.8}{6.4}\selectfont \textcolor{gray}{n/a}}\\[-.2pt]\colorbox{BenchBest}{\strut 3m\,3.7s}\\[-.2pt]{\fontsize{5.8}{6.4}\selectfont \textcolor{gray}{n/a}}\end{tabular} & \hspace*{1pt}\begin{tabular}[t]{@{}c@{}}\rule{0pt}{12pt}3m\,28.3s\\[-.2pt]{\fontsize{5.8}{6.4}\selectfont \textcolor{BrickRed}{$ \blacktriangledown\,1.3\times $}}\\[-.2pt]3m\,22.8s\\[-.2pt]{\fontsize{5.8}{6.4}\selectfont \textcolor{BrickRed}{$ \blacktriangledown\,1.4\times $}}\end{tabular} & \hspace*{1pt}\begin{tabular}[t]{@{}c@{}}\rule{0pt}{12pt}3m\,53.2s\\[-.2pt]{\fontsize{5.8}{6.4}\selectfont \textcolor{BrickRed}{$ \blacktriangledown\,1.4\times $}}\\[-.2pt]3m\,43.5s\\[-.2pt]{\fontsize{5.8}{6.4}\selectfont \textcolor{BrickRed}{$ \blacktriangledown\,1.4\times $}}\end{tabular} & \hspace*{1pt}\begin{tabular}[t]{@{}c@{}}\rule{0pt}{12pt}7m\,14.6s\\[-.2pt]{\fontsize{5.8}{6.4}\selectfont \textcolor{BrickRed}{$ \blacktriangledown\,1.5\times $}}\\[-.2pt]6m\,57.2s\\[-.2pt]{\fontsize{5.8}{6.4}\selectfont \textcolor{BrickRed}{$ \blacktriangledown\,1.6\times $}}\end{tabular} & \hspace*{1pt}\begin{tabular}[t]{@{}c@{}}\rule{0pt}{12pt}\colorbox{BenchBest}{\strut 51m\,22s}\\[-.2pt]{\fontsize{5.8}{6.4}\selectfont \textcolor{gray}{n/a}}\\[-.2pt]\colorbox{BenchBest}{\strut 51m\,17.2s}\\[-.2pt]{\fontsize{5.8}{6.4}\selectfont \textcolor{gray}{n/a}}\end{tabular} & \hspace*{1pt}\begin{tabular}[t]{@{}c@{}}\rule{0pt}{12pt}\colorbox{BenchBest}{\strut 1m\,10.3s}\\[-.2pt]{\fontsize{5.8}{6.4}\selectfont \textcolor{ForestGreen}{$ \blacktriangle\,7.6\times $}}\\[-.2pt]\colorbox{BenchBest}{\strut 1m\,9s}\\[-.2pt]{\fontsize{5.8}{6.4}\selectfont \textcolor{ForestGreen}{$ \blacktriangle\,2.4\times $}}\end{tabular} & \hspace*{1pt}\begin{tabular}[t]{@{}c@{}}\rule{0pt}{12pt}\colorbox{BenchBest}{\strut 1m\,8.6s}\\[-.2pt]{\fontsize{5.8}{6.4}\selectfont \textcolor{ForestGreen}{$ \blacktriangle\,6.4\times $}}\\[-.2pt]\colorbox{BenchBest}{\strut 1m\,6.9s}\\[-.2pt]{\fontsize{5.8}{6.4}\selectfont \textcolor{ForestGreen}{$ \blacktriangle\,2.5\times $}}\end{tabular} & \hspace*{1pt}\begin{tabular}[t]{@{}c@{}}\rule{0pt}{12pt}\colorbox{BenchBest}{\strut 2m\,21.5s}\\[-.2pt]{\fontsize{5.8}{6.4}\selectfont \textcolor{ForestGreen}{$ \blacktriangle\,6.0\times $}}\\[-.2pt]\colorbox{BenchBest}{\strut 2m\,18.4s}\\[-.2pt]{\fontsize{5.8}{6.4}\selectfont \textcolor{ForestGreen}{$ \blacktriangle\,2.2\times $}}\end{tabular} & \hspace*{1pt}\begin{tabular}[t]{@{}c@{}}\rule{0pt}{12pt}\colorbox{BenchBest}{\strut 15m\,56.8s}\\[-.2pt]{\fontsize{5.8}{6.4}\selectfont \textcolor{gray}{n/a}}\\[-.2pt]\colorbox{BenchBest}{\strut 15m\,51.9s}\\[-.2pt]{\fontsize{5.8}{6.4}\selectfont \textcolor{gray}{n/a}}\end{tabular}\\[3pt]\hline
\textbf{BioseqDB} & \begin{tabular}[t]{@{}l@{}}\rule{0pt}{12pt}MPF / Bounded\\RMQ / Block\end{tabular} & \hspace*{1pt}\begin{tabular}[t]{@{}c@{}}\rule{0pt}{12pt}\colorbox{BenchBest}{\strut 8.3s}\\[-.2pt]{\fontsize{5.8}{6.4}\selectfont \textcolor{ForestGreen}{$ \blacktriangle\,20.4\times $}}\\[-.2pt]7.6s\\[-.2pt]{\fontsize{5.8}{6.4}\selectfont \textcolor{ForestGreen}{$ \blacktriangle\,21.3\times $}}\end{tabular} & \hspace*{1pt}\begin{tabular}[t]{@{}c@{}}\rule{0pt}{12pt}\colorbox{BenchBest}{\strut 8.6s}\\[-.2pt]{\fontsize{5.8}{6.4}\selectfont \textcolor{ForestGreen}{$ \blacktriangle\,19.8\times $}}\\[-.2pt]7.8s\\[-.2pt]{\fontsize{5.8}{6.4}\selectfont \textcolor{ForestGreen}{$ \blacktriangle\,20.8\times $}}\end{tabular} & \textcolor{gray}{\ensuremath{\cdot}} & \textcolor{gray}{\ensuremath{\cdot}} & \textcolor{gray}{\ensuremath{\cdot}} & \textcolor{gray}{\ensuremath{\cdot}} & \textcolor{gray}{\ensuremath{\cdot}} & \textcolor{gray}{\ensuremath{\cdot}} & \hspace*{1pt}\begin{tabular}[t]{@{}c@{}}\rule{0pt}{12pt}2m\,53.8s\\[-.2pt]{\fontsize{5.8}{6.4}\selectfont \textcolor{ForestGreen}{$ \blacktriangle\,17\% $}}\\[-.2pt]2m\,43.7s\\[-.2pt]{\fontsize{5.8}{6.4}\selectfont \textcolor{ForestGreen}{$ \blacktriangle\,19\% $}}\end{tabular} & \hspace*{1pt}\begin{tabular}[t]{@{}c@{}}\rule{0pt}{12pt}3m\,4.5s\\[-.2pt]{\fontsize{5.8}{6.4}\selectfont \textcolor{ForestGreen}{$ \blacktriangle\,21\% $}}\\[-.2pt]2m\,52.1s\\[-.2pt]{\fontsize{5.8}{6.4}\selectfont \textcolor{ForestGreen}{$ \blacktriangle\,23\% $}}\end{tabular} & \textcolor{gray}{\ensuremath{\cdot}} & \textcolor{gray}{\ensuremath{\cdot}} & \textcolor{gray}{\ensuremath{\cdot}} & \textcolor{gray}{\ensuremath{\cdot}} & \textcolor{gray}{\ensuremath{\cdot}} & \textcolor{gray}{\ensuremath{\cdot}}\\[3pt]\hline
 & \begin{tabular}[t]{@{}l@{}}\rule{0pt}{12pt}MPF / Bounded\\RMQ / KSW2\end{tabular} & \hspace*{1pt}\begin{tabular}[t]{@{}c@{}}\rule{0pt}{12pt}10.6s\\[-.2pt]{\fontsize{5.8}{6.4}\selectfont \textcolor{ForestGreen}{$ \blacktriangle\,15.9\times $}}\\[-.2pt]8.3s\\[-.2pt]{\fontsize{5.8}{6.4}\selectfont \textcolor{ForestGreen}{$ \blacktriangle\,19.5\times $}}\end{tabular} & \hspace*{1pt}\begin{tabular}[t]{@{}c@{}}\rule{0pt}{12pt}14.6s\\[-.2pt]{\fontsize{5.8}{6.4}\selectfont \textcolor{ForestGreen}{$ \blacktriangle\,11.7\times $}}\\[-.2pt]12.5s\\[-.2pt]{\fontsize{5.8}{6.4}\selectfont \textcolor{ForestGreen}{$ \blacktriangle\,13.0\times $}}\end{tabular} & \textcolor{gray}{\ensuremath{\cdot}} & \textcolor{gray}{\ensuremath{\cdot}} & \textcolor{gray}{\ensuremath{\cdot}} & \textcolor{gray}{\ensuremath{\cdot}} & \textcolor{gray}{\ensuremath{\cdot}} & \textcolor{gray}{\ensuremath{\cdot}} & \hspace*{1pt}\begin{tabular}[t]{@{}c@{}}\rule{0pt}{12pt}2m\,51.8s\\[-.2pt]{\fontsize{5.8}{6.4}\selectfont \textcolor{ForestGreen}{$ \blacktriangle\,18\% $}}\\[-.2pt]2m\,40.6s\\[-.2pt]{\fontsize{5.8}{6.4}\selectfont \textcolor{ForestGreen}{$ \blacktriangle\,21\% $}}\end{tabular} & \hspace*{1pt}\begin{tabular}[t]{@{}c@{}}\rule{0pt}{12pt}2m\,50.2s\\[-.2pt]{\fontsize{5.8}{6.4}\selectfont \textcolor{ForestGreen}{$ \blacktriangle\,27\% $}}\\[-.2pt]2m\,36.4s\\[-.2pt]{\fontsize{5.8}{6.4}\selectfont \textcolor{ForestGreen}{$ \blacktriangle\,1.4\times $}}\end{tabular} & \textcolor{gray}{\ensuremath{\cdot}} & \textcolor{gray}{\ensuremath{\cdot}} & \textcolor{gray}{\ensuremath{\cdot}} & \textcolor{gray}{\ensuremath{\cdot}} & \textcolor{gray}{\ensuremath{\cdot}} & \textcolor{gray}{\ensuremath{\cdot}}\\[3pt]\hline
 & \begin{tabular}[t]{@{}l@{}}\rule{0pt}{12pt}MPF / Bounded\\Scalar-DP / Block\end{tabular} & \hspace*{1pt}\begin{tabular}[t]{@{}c@{}}\rule{0pt}{12pt}10.7s\\[-.2pt]{\fontsize{5.8}{6.4}\selectfont \textcolor{ForestGreen}{$ \blacktriangle\,15.7\times $}}\\[-.2pt]10.1s\\[-.2pt]{\fontsize{5.8}{6.4}\selectfont \textcolor{ForestGreen}{$ \blacktriangle\,16.1\times $}}\end{tabular} & \hspace*{1pt}\begin{tabular}[t]{@{}c@{}}\rule{0pt}{12pt}11.5s\\[-.2pt]{\fontsize{5.8}{6.4}\selectfont \textcolor{ForestGreen}{$ \blacktriangle\,14.8\times $}}\\[-.2pt]10.8s\\[-.2pt]{\fontsize{5.8}{6.4}\selectfont \textcolor{ForestGreen}{$ \blacktriangle\,15.2\times $}}\end{tabular} & \textcolor{gray}{\ensuremath{\cdot}} & \textcolor{gray}{\ensuremath{\cdot}} & \textcolor{gray}{\ensuremath{\cdot}} & \textcolor{gray}{\ensuremath{\cdot}} & \textcolor{gray}{\ensuremath{\cdot}} & \textcolor{gray}{\ensuremath{\cdot}} & \hspace*{1pt}\begin{tabular}[t]{@{}c@{}}\rule{0pt}{12pt}2m\,50s\\[-.2pt]{\fontsize{5.8}{6.4}\selectfont \textcolor{ForestGreen}{$ \blacktriangle\,18\% $}}\\[-.2pt]2m\,41s\\[-.2pt]{\fontsize{5.8}{6.4}\selectfont \textcolor{ForestGreen}{$ \blacktriangle\,21\% $}}\end{tabular} & \hspace*{1pt}\begin{tabular}[t]{@{}c@{}}\rule{0pt}{12pt}\colorbox{BenchBest}{\strut 2m\,48.3s}\\[-.2pt]{\fontsize{5.8}{6.4}\selectfont \textcolor{ForestGreen}{$ \blacktriangle\,28\% $}}\\[-.2pt]\colorbox{BenchBest}{\strut 2m\,35.9s}\\[-.2pt]{\fontsize{5.8}{6.4}\selectfont \textcolor{ForestGreen}{$ \blacktriangle\,1.4\times $}}\end{tabular} & \hspace*{1pt}\begin{tabular}[t]{@{}c@{}}\rule{0pt}{12pt}\colorbox{BenchBest}{\strut 4m\,47.5s}\\[-.2pt]{\fontsize{5.8}{6.4}\selectfont \textcolor{ForestGreen}{$ \blacktriangle\,1.5\times $}}\\[-.2pt]\colorbox{BenchBest}{\strut 4m\,26.9s}\\[-.2pt]{\fontsize{5.8}{6.4}\selectfont \textcolor{ForestGreen}{$ \blacktriangle\,1.6\times $}}\end{tabular} & \textcolor{gray}{\ensuremath{\cdot}} & \textcolor{gray}{\ensuremath{\cdot}} & \textcolor{gray}{\ensuremath{\cdot}} & \textcolor{gray}{\ensuremath{\cdot}} & \textcolor{gray}{\ensuremath{\cdot}}\\[3pt]\hline
 & \begin{tabular}[t]{@{}l@{}}\rule{0pt}{12pt}MPF / Bounded\\Scalar-DP / KSW2\end{tabular} & \hspace*{1pt}\begin{tabular}[t]{@{}c@{}}\rule{0pt}{12pt}14.5s\\[-.2pt]{\fontsize{5.8}{6.4}\selectfont \textcolor{ForestGreen}{$ \blacktriangle\,11.6\times $}}\\[-.2pt]12.2s\\[-.2pt]{\fontsize{5.8}{6.4}\selectfont \textcolor{ForestGreen}{$ \blacktriangle\,13.3\times $}}\end{tabular} & \hspace*{1pt}\begin{tabular}[t]{@{}c@{}}\rule{0pt}{12pt}14.2s\\[-.2pt]{\fontsize{5.8}{6.4}\selectfont \textcolor{ForestGreen}{$ \blacktriangle\,12.0\times $}}\\[-.2pt]11.9s\\[-.2pt]{\fontsize{5.8}{6.4}\selectfont \textcolor{ForestGreen}{$ \blacktriangle\,13.7\times $}}\end{tabular} & \textcolor{gray}{\ensuremath{\cdot}} & \textcolor{gray}{\ensuremath{\cdot}} & \textcolor{gray}{\ensuremath{\cdot}} & \textcolor{gray}{\ensuremath{\cdot}} & \textcolor{gray}{\ensuremath{\cdot}} & \textcolor{gray}{\ensuremath{\cdot}} & \hspace*{1pt}\begin{tabular}[t]{@{}c@{}}\rule{0pt}{12pt}2m\,40.6s\\[-.2pt]{\fontsize{5.8}{6.4}\selectfont \textcolor{ForestGreen}{$ \blacktriangle\,23\% $}}\\[-.2pt]2m\,29.9s\\[-.2pt]{\fontsize{5.8}{6.4}\selectfont \textcolor{ForestGreen}{$ \blacktriangle\,26\% $}}\end{tabular} & \hspace*{1pt}\begin{tabular}[t]{@{}c@{}}\rule{0pt}{12pt}3m\,2.9s\\[-.2pt]{\fontsize{5.8}{6.4}\selectfont \textcolor{ForestGreen}{$ \blacktriangle\,22\% $}}\\[-.2pt]2m\,48.3s\\[-.2pt]{\fontsize{5.8}{6.4}\selectfont \textcolor{ForestGreen}{$ \blacktriangle\,25\% $}}\end{tabular} & \textcolor{gray}{\ensuremath{\cdot}} & \textcolor{gray}{\ensuremath{\cdot}} & \textcolor{gray}{\ensuremath{\cdot}} & \textcolor{gray}{\ensuremath{\cdot}} & \textcolor{gray}{\ensuremath{\cdot}} & \textcolor{gray}{\ensuremath{\cdot}}\\[3pt]\hline
\rowcolor{BenchPreferred}[0.5pt][0.5pt] & \begin{tabular}[t]{@{}l@{}}\rule{0pt}{12pt}MPF / Bounded\\SIMD-DP / Block\end{tabular} & \hspace*{1pt}\begin{tabular}[t]{@{}c@{}}\rule{0pt}{12pt}9.2s\\[-.2pt]{\fontsize{5.8}{6.4}\selectfont \textcolor{ForestGreen}{$ \blacktriangle\,18.3\times $}}\\[-.2pt]8.5s\\[-.2pt]{\fontsize{5.8}{6.4}\selectfont \textcolor{ForestGreen}{$ \blacktriangle\,19.0\times $}}\end{tabular} & \hspace*{1pt}\begin{tabular}[t]{@{}c@{}}\rule{0pt}{12pt}11.4s\\[-.2pt]{\fontsize{5.8}{6.4}\selectfont \textcolor{ForestGreen}{$ \blacktriangle\,15.0\times $}}\\[-.2pt]10.7s\\[-.2pt]{\fontsize{5.8}{6.4}\selectfont \textcolor{ForestGreen}{$ \blacktriangle\,15.3\times $}}\end{tabular} & \hspace*{1pt}\begin{tabular}[t]{@{}c@{}}\rule{0pt}{12pt}\colorbox{BenchBest}{\strut 6.8s}\\[-.2pt]{\fontsize{5.8}{6.4}\selectfont \textcolor{ForestGreen}{$ \blacktriangle\,43.2\times $}}\\[-.2pt]\colorbox{BenchBest}{\strut 5.1s}\\[-.2pt]{\fontsize{5.8}{6.4}\selectfont \textcolor{ForestGreen}{$ \blacktriangle\,54.6\times $}}\end{tabular} & \textcolor{gray}{\ensuremath{\cdot}} & \hspace*{1pt}\begin{tabular}[t]{@{}c@{}}\rule{0pt}{12pt}\colorbox{BenchBest}{\strut 18.2s}\\[-.2pt]{\fontsize{5.8}{6.4}\selectfont \textcolor{ForestGreen}{$ \blacktriangle\,1.9\times $}}\\[-.2pt]\colorbox{BenchBest}{\strut 16s}\\[-.2pt]{\fontsize{5.8}{6.4}\selectfont \textcolor{ForestGreen}{$ \blacktriangle\,2.1\times $}}\end{tabular} & \hspace*{1pt}\begin{tabular}[t]{@{}c@{}}\rule{0pt}{12pt}\colorbox{BenchBest}{\strut 13.2s}\\[-.2pt]{\fontsize{5.8}{6.4}\selectfont \textcolor{ForestGreen}{$ \blacktriangle\,2.6\times $}}\\[-.2pt]\colorbox{BenchBest}{\strut 11.4s}\\[-.2pt]{\fontsize{5.8}{6.4}\selectfont \textcolor{ForestGreen}{$ \blacktriangle\,2.9\times $}}\end{tabular} & \hspace*{1pt}\begin{tabular}[t]{@{}c@{}}\rule{0pt}{12pt}\colorbox{BenchBest}{\strut 18.6s}\\[-.2pt]{\fontsize{5.8}{6.4}\selectfont \textcolor{ForestGreen}{$ \blacktriangle\,3.7\times $}}\\[-.2pt]\colorbox{BenchBest}{\strut 14.3s}\\[-.2pt]{\fontsize{5.8}{6.4}\selectfont \textcolor{ForestGreen}{$ \blacktriangle\,4.6\times $}}\end{tabular} & \textcolor{gray}{\ensuremath{\cdot}} & \hspace*{1pt}\begin{tabular}[t]{@{}c@{}}\rule{0pt}{12pt}\colorbox{BenchBest}{\strut 2m\,38.5s}\\[-.2pt]{\fontsize{5.8}{6.4}\selectfont \textcolor{ForestGreen}{$ \blacktriangle\,24\% $}}\\[-.2pt]\colorbox{BenchBest}{\strut 2m\,29.8s}\\[-.2pt]{\fontsize{5.8}{6.4}\selectfont \textcolor{ForestGreen}{$ \blacktriangle\,26\% $}}\end{tabular} & \hspace*{1pt}\begin{tabular}[t]{@{}c@{}}\rule{0pt}{12pt}2m\,49s\\[-.2pt]{\fontsize{5.8}{6.4}\selectfont \textcolor{ForestGreen}{$ \blacktriangle\,28\% $}}\\[-.2pt]2m\,37s\\[-.2pt]{\fontsize{5.8}{6.4}\selectfont \textcolor{ForestGreen}{$ \blacktriangle\,30\% $}}\end{tabular} & \hspace*{1pt}\begin{tabular}[t]{@{}c@{}}\rule{0pt}{12pt}5m\,3.3s\\[-.2pt]{\fontsize{5.8}{6.4}\selectfont \textcolor{ForestGreen}{$ \blacktriangle\,1.4\times $}}\\[-.2pt]4m\,42.7s\\[-.2pt]{\fontsize{5.8}{6.4}\selectfont \textcolor{ForestGreen}{$ \blacktriangle\,1.5\times $}}\end{tabular} & \textcolor{gray}{\ensuremath{\cdot}} & \hspace*{1pt}\begin{tabular}[t]{@{}c@{}}\rule{0pt}{12pt}8m\,52.3s\\[-.2pt]{\fontsize{5.8}{6.4}\selectfont \textcolor{BrickRed}{$ \blacktriangledown\,7.6\times $}}\\[-.2pt]2m\,45.5s\\[-.2pt]{\fontsize{5.8}{6.4}\selectfont \textcolor{BrickRed}{$ \blacktriangledown\,2.4\times $}}\end{tabular} & \hspace*{1pt}\begin{tabular}[t]{@{}c@{}}\rule{0pt}{12pt}7m\,22.2s\\[-.2pt]{\fontsize{5.8}{6.4}\selectfont \textcolor{BrickRed}{$ \blacktriangledown\,6.4\times $}}\\[-.2pt]2m\,48.2s\\[-.2pt]{\fontsize{5.8}{6.4}\selectfont \textcolor{BrickRed}{$ \blacktriangledown\,2.5\times $}}\end{tabular} & \hspace*{1pt}\begin{tabular}[t]{@{}c@{}}\rule{0pt}{12pt}14m\,43.3s\\[-.2pt]{\fontsize{5.8}{6.4}\selectfont \textcolor{BrickRed}{$ \blacktriangledown\,6.2\times $}}\\[-.2pt]5m\,9.5s\\[-.2pt]{\fontsize{5.8}{6.4}\selectfont \textcolor{BrickRed}{$ \blacktriangledown\,2.2\times $}}\end{tabular} & \textcolor{gray}{\ensuremath{\cdot}}\\[3pt]\hline
 & \begin{tabular}[t]{@{}l@{}}\rule{0pt}{12pt}MPF / Bounded\\SIMD-DP / KSW2\end{tabular} & \hspace*{1pt}\begin{tabular}[t]{@{}c@{}}\rule{0pt}{12pt}9.6s\\[-.2pt]{\fontsize{5.8}{6.4}\selectfont \textcolor{ForestGreen}{$ \blacktriangle\,17.6\times $}}\\[-.2pt]\colorbox{BenchBest}{\strut 7.1s}\\[-.2pt]{\fontsize{5.8}{6.4}\selectfont \textcolor{ForestGreen}{$ \blacktriangle\,22.6\times $}}\end{tabular} & \hspace*{1pt}\begin{tabular}[t]{@{}c@{}}\rule{0pt}{12pt}9.6s\\[-.2pt]{\fontsize{5.8}{6.4}\selectfont \textcolor{ForestGreen}{$ \blacktriangle\,17.7\times $}}\\[-.2pt]\colorbox{BenchBest}{\strut 7.3s}\\[-.2pt]{\fontsize{5.8}{6.4}\selectfont \textcolor{ForestGreen}{$ \blacktriangle\,22.2\times $}}\end{tabular} & \textcolor{gray}{\ensuremath{\cdot}} & \textcolor{gray}{\ensuremath{\cdot}} & \textcolor{gray}{\ensuremath{\cdot}} & \textcolor{gray}{\ensuremath{\cdot}} & \textcolor{gray}{\ensuremath{\cdot}} & \textcolor{gray}{\ensuremath{\cdot}} & \hspace*{1pt}\begin{tabular}[t]{@{}c@{}}\rule{0pt}{12pt}2m\,51.2s\\[-.2pt]{\fontsize{5.8}{6.4}\selectfont \textcolor{ForestGreen}{$ \blacktriangle\,18\% $}}\\[-.2pt]2m\,39.2s\\[-.2pt]{\fontsize{5.8}{6.4}\selectfont \textcolor{ForestGreen}{$ \blacktriangle\,21\% $}}\end{tabular} & \hspace*{1pt}\begin{tabular}[t]{@{}c@{}}\rule{0pt}{12pt}3m\,0s\\[-.2pt]{\fontsize{5.8}{6.4}\selectfont \textcolor{ForestGreen}{$ \blacktriangle\,23\% $}}\\[-.2pt]2m\,45.6s\\[-.2pt]{\fontsize{5.8}{6.4}\selectfont \textcolor{ForestGreen}{$ \blacktriangle\,26\% $}}\end{tabular} & \textcolor{gray}{\ensuremath{\cdot}} & \textcolor{gray}{\ensuremath{\cdot}} & \textcolor{gray}{\ensuremath{\cdot}} & \textcolor{gray}{\ensuremath{\cdot}} & \textcolor{gray}{\ensuremath{\cdot}} & \textcolor{gray}{\ensuremath{\cdot}}\\[3pt]\hline
\end{tabular}\par}
\smallskip{\footnotesize\noindent Each cell: A, comparison, B, comparison. Times cover 1,024 queries. $\cdot$ = no approved observation. ---\textsuperscript{g} = GPU unavailable, not executed. n/a = no measured comparison.\par
\noindent \textsuperscript{f}SIMD-ref request observed as scalar, compared only as that effective path. \textsuperscript{p}Accepted by explicit preservation after wrapper postflight failure, no run receipt claimed.\par
\noindent MPF/MPS = minimizers per fragment/sequence. Row order: index representation / minimizer, then chaining / base alignment. SIMD-ref is the reference minimizer, Bounded is bounded ntHash.\par
\noindent Blue = one lowest A and one lowest B down each machine/workload column, across continuation pages. Exact ties select the first row. Repeated baselines are marked only on page 1.\par
\noindent Light gray = minimap2 baseline. Pale yellow = preferred MPF / Bounded / SIMD-DP / Block configuration, not a claim of winning every column.\par
\noindent When all Vulkan modes are unmeasured, they share one row. Baseline comparisons use the fastest measured BioseqDB effective path separately for A and B.\par}
\clearpage
\noindent{\large\bfseries AWS: 4 Gbp, k=20, w=15 (mapping interval)\hfill 2/4}\par\smallskip
{\fontsize{7.1}{8.1}\selectfont\setlength{\tabcolsep}{1pt}\setlength{\fboxsep}{2pt}\renewcommand{\arraystretch}{1.12}
\noindent\begin{tabular}{@{}p{14mm}p{27mm}|p{13.55mm}|p{13.55mm}|p{13.55mm}|p{13.55mm}|p{13.55mm}|p{13.55mm}|p{13.55mm}|p{13.55mm}|p{13.55mm}|p{13.55mm}|p{13.55mm}|p{13.55mm}|p{13.55mm}|p{13.55mm}|p{13.55mm}|p{13.55mm}|@{}}\hline
Engine & Execution & \multicolumn{8}{c|}{\textbf{Index}} & \multicolumn{8}{c|}{\textbf{No index}}\\
 &  & \multicolumn{8}{c|}{\textbf{4 Gbp}} & \multicolumn{8}{c|}{\textbf{4 Gbp}}\\
 &  & \multicolumn{8}{c|}{\textbf{k=20, w=15}} & \multicolumn{8}{c|}{\textbf{k=20, w=15}}\\
 &  & \multicolumn{4}{c|}{\textbf{DNA4}} & \multicolumn{4}{c|}{\textbf{DNA5}} & \multicolumn{4}{c|}{\textbf{DNA4}} & \multicolumn{4}{c|}{\textbf{DNA5}}\\
 &  & \multicolumn{1}{c|}{\textbf{M7A}} & \multicolumn{1}{c|}{\textbf{M7I}} & \multicolumn{1}{c|}{\textbf{M4}} & \multicolumn{1}{c|}{\textbf{C7G}} & \multicolumn{1}{c|}{\textbf{M7A}} & \multicolumn{1}{c|}{\textbf{M7I}} & \multicolumn{1}{c|}{\textbf{M4}} & \multicolumn{1}{c|}{\textbf{C7G}} & \multicolumn{1}{c|}{\textbf{M7A}} & \multicolumn{1}{c|}{\textbf{M7I}} & \multicolumn{1}{c|}{\textbf{M4}} & \multicolumn{1}{c|}{\textbf{C7G}} & \multicolumn{1}{c|}{\textbf{M7A}} & \multicolumn{1}{c|}{\textbf{M7I}} & \multicolumn{1}{c|}{\textbf{M4}} & \multicolumn{1}{c|}{\textbf{C7G}}\\
\hline
\rowcolor{BenchBaseline}[0.5pt][0.5pt]\textbf{minimap2} & \textbf{Baseline} & \hspace*{1pt}\begin{tabular}[t]{@{}c@{}}\rule{0pt}{12pt}2m\,48.1s\\[-.2pt]{\fontsize{5.8}{6.4}\selectfont \textcolor{BrickRed}{$ \blacktriangledown\,20.4\times $}}\\[-.2pt]2m\,41.7s\\[-.2pt]{\fontsize{5.8}{6.4}\selectfont \textcolor{BrickRed}{$ \blacktriangledown\,22.6\times $}}\end{tabular} & \hspace*{1pt}\begin{tabular}[t]{@{}c@{}}\rule{0pt}{12pt}2m\,50.3s\\[-.2pt]{\fontsize{5.8}{6.4}\selectfont \textcolor{BrickRed}{$ \blacktriangledown\,19.8\times $}}\\[-.2pt]2m\,43s\\[-.2pt]{\fontsize{5.8}{6.4}\selectfont \textcolor{BrickRed}{$ \blacktriangledown\,22.2\times $}}\end{tabular} & \hspace*{1pt}\begin{tabular}[t]{@{}c@{}}\rule{0pt}{12pt}4m\,54.3s\\[-.2pt]{\fontsize{5.8}{6.4}\selectfont \textcolor{BrickRed}{$ \blacktriangledown\,43.2\times $}}\\[-.2pt]4m\,39s\\[-.2pt]{\fontsize{5.8}{6.4}\selectfont \textcolor{BrickRed}{$ \blacktriangledown\,54.6\times $}}\end{tabular} & \hspace*{1pt}\begin{tabular}[t]{@{}c@{}}\rule{0pt}{12pt}25m\,5.5s\\[-.2pt]{\fontsize{5.8}{6.4}\selectfont \textcolor{gray}{n/a}}\\[-.2pt]25m\,0.7s\\[-.2pt]{\fontsize{5.8}{6.4}\selectfont \textcolor{gray}{n/a}}\end{tabular} & \hspace*{1pt}\begin{tabular}[t]{@{}c@{}}\rule{0pt}{12pt}33.9s\\[-.2pt]{\fontsize{5.8}{6.4}\selectfont \textcolor{BrickRed}{$ \blacktriangledown\,1.9\times $}}\\[-.2pt]32.9s\\[-.2pt]{\fontsize{5.8}{6.4}\selectfont \textcolor{BrickRed}{$ \blacktriangledown\,2.1\times $}}\end{tabular} & \hspace*{1pt}\begin{tabular}[t]{@{}c@{}}\rule{0pt}{12pt}34.4s\\[-.2pt]{\fontsize{5.8}{6.4}\selectfont \textcolor{BrickRed}{$ \blacktriangledown\,2.6\times $}}\\[-.2pt]32.7s\\[-.2pt]{\fontsize{5.8}{6.4}\selectfont \textcolor{BrickRed}{$ \blacktriangledown\,2.9\times $}}\end{tabular} & \hspace*{1pt}\begin{tabular}[t]{@{}c@{}}\rule{0pt}{12pt}1m\,8.2s\\[-.2pt]{\fontsize{5.8}{6.4}\selectfont \textcolor{BrickRed}{$ \blacktriangledown\,3.7\times $}}\\[-.2pt]1m\,5.1s\\[-.2pt]{\fontsize{5.8}{6.4}\selectfont \textcolor{BrickRed}{$ \blacktriangledown\,4.6\times $}}\end{tabular} & \hspace*{1pt}\begin{tabular}[t]{@{}c@{}}\rule{0pt}{12pt}3m\,9.8s\\[-.2pt]{\fontsize{5.8}{6.4}\selectfont \textcolor{gray}{n/a}}\\[-.2pt]3m\,3.7s\\[-.2pt]{\fontsize{5.8}{6.4}\selectfont \textcolor{gray}{n/a}}\end{tabular} & \hspace*{1pt}\begin{tabular}[t]{@{}c@{}}\rule{0pt}{12pt}3m\,28.3s\\[-.2pt]{\fontsize{5.8}{6.4}\selectfont \textcolor{BrickRed}{$ \blacktriangledown\,1.3\times $}}\\[-.2pt]3m\,22.8s\\[-.2pt]{\fontsize{5.8}{6.4}\selectfont \textcolor{BrickRed}{$ \blacktriangledown\,1.4\times $}}\end{tabular} & \hspace*{1pt}\begin{tabular}[t]{@{}c@{}}\rule{0pt}{12pt}3m\,53.2s\\[-.2pt]{\fontsize{5.8}{6.4}\selectfont \textcolor{BrickRed}{$ \blacktriangledown\,1.4\times $}}\\[-.2pt]3m\,43.5s\\[-.2pt]{\fontsize{5.8}{6.4}\selectfont \textcolor{BrickRed}{$ \blacktriangledown\,1.4\times $}}\end{tabular} & \hspace*{1pt}\begin{tabular}[t]{@{}c@{}}\rule{0pt}{12pt}7m\,14.6s\\[-.2pt]{\fontsize{5.8}{6.4}\selectfont \textcolor{BrickRed}{$ \blacktriangledown\,1.5\times $}}\\[-.2pt]6m\,57.2s\\[-.2pt]{\fontsize{5.8}{6.4}\selectfont \textcolor{BrickRed}{$ \blacktriangledown\,1.6\times $}}\end{tabular} & \hspace*{1pt}\begin{tabular}[t]{@{}c@{}}\rule{0pt}{12pt}51m\,22s\\[-.2pt]{\fontsize{5.8}{6.4}\selectfont \textcolor{gray}{n/a}}\\[-.2pt]51m\,17.2s\\[-.2pt]{\fontsize{5.8}{6.4}\selectfont \textcolor{gray}{n/a}}\end{tabular} & \hspace*{1pt}\begin{tabular}[t]{@{}c@{}}\rule{0pt}{12pt}1m\,10.3s\\[-.2pt]{\fontsize{5.8}{6.4}\selectfont \textcolor{ForestGreen}{$ \blacktriangle\,7.6\times $}}\\[-.2pt]1m\,9s\\[-.2pt]{\fontsize{5.8}{6.4}\selectfont \textcolor{ForestGreen}{$ \blacktriangle\,2.4\times $}}\end{tabular} & \hspace*{1pt}\begin{tabular}[t]{@{}c@{}}\rule{0pt}{12pt}1m\,8.6s\\[-.2pt]{\fontsize{5.8}{6.4}\selectfont \textcolor{ForestGreen}{$ \blacktriangle\,6.4\times $}}\\[-.2pt]1m\,6.9s\\[-.2pt]{\fontsize{5.8}{6.4}\selectfont \textcolor{ForestGreen}{$ \blacktriangle\,2.5\times $}}\end{tabular} & \hspace*{1pt}\begin{tabular}[t]{@{}c@{}}\rule{0pt}{12pt}2m\,21.5s\\[-.2pt]{\fontsize{5.8}{6.4}\selectfont \textcolor{ForestGreen}{$ \blacktriangle\,6.0\times $}}\\[-.2pt]2m\,18.4s\\[-.2pt]{\fontsize{5.8}{6.4}\selectfont \textcolor{ForestGreen}{$ \blacktriangle\,2.2\times $}}\end{tabular} & \hspace*{1pt}\begin{tabular}[t]{@{}c@{}}\rule{0pt}{12pt}15m\,56.8s\\[-.2pt]{\fontsize{5.8}{6.4}\selectfont \textcolor{gray}{n/a}}\\[-.2pt]15m\,51.9s\\[-.2pt]{\fontsize{5.8}{6.4}\selectfont \textcolor{gray}{n/a}}\end{tabular}\\[3pt]\hline
\textbf{BioseqDB} & \begin{tabular}[t]{@{}l@{}}\rule{0pt}{12pt}MPF / Scalar\\RMQ / Block\end{tabular} & \hspace*{1pt}\begin{tabular}[t]{@{}c@{}}\rule{0pt}{12pt}11m\,22.5s\\[-.2pt]{\fontsize{5.8}{6.4}\selectfont \textcolor{BrickRed}{$ \blacktriangledown\,4.1\times $}}\\[-.2pt]11m\,22s\\[-.2pt]{\fontsize{5.8}{6.4}\selectfont \textcolor{BrickRed}{$ \blacktriangledown\,4.2\times $}}\end{tabular} & \hspace*{1pt}\begin{tabular}[t]{@{}c@{}}\rule{0pt}{12pt}24m\,52.9s\\[-.2pt]{\fontsize{5.8}{6.4}\selectfont \textcolor{BrickRed}{$ \blacktriangledown\,8.8\times $}}\\[-.2pt]24m\,52.4s\\[-.2pt]{\fontsize{5.8}{6.4}\selectfont \textcolor{BrickRed}{$ \blacktriangledown\,9.2\times $}}\end{tabular} & \textcolor{gray}{\ensuremath{\cdot}} & \textcolor{gray}{\ensuremath{\cdot}} & \textcolor{gray}{\ensuremath{\cdot}} & \textcolor{gray}{\ensuremath{\cdot}} & \textcolor{gray}{\ensuremath{\cdot}} & \textcolor{gray}{\ensuremath{\cdot}} & \hspace*{1pt}\begin{tabular}[t]{@{}c@{}}\rule{0pt}{12pt}5m\,16.6s\\[-.2pt]{\fontsize{5.8}{6.4}\selectfont \textcolor{BrickRed}{$ \blacktriangledown\,1.5\times $}}\\[-.2pt]4m\,33.9s\\[-.2pt]{\fontsize{5.8}{6.4}\selectfont \textcolor{BrickRed}{$ \blacktriangledown\,1.4\times $}}\end{tabular} & \hspace*{1pt}\begin{tabular}[t]{@{}c@{}}\rule{0pt}{12pt}4m\,53.5s\\[-.2pt]{\fontsize{5.8}{6.4}\selectfont \textcolor{BrickRed}{$ \blacktriangledown\,26\% $}}\\[-.2pt]4m\,17.9s\\[-.2pt]{\fontsize{5.8}{6.4}\selectfont \textcolor{BrickRed}{$ \blacktriangledown\,15\% $}}\end{tabular} & \textcolor{gray}{\ensuremath{\cdot}} & \textcolor{gray}{\ensuremath{\cdot}} & \textcolor{gray}{\ensuremath{\cdot}} & \textcolor{gray}{\ensuremath{\cdot}} & \textcolor{gray}{\ensuremath{\cdot}} & \textcolor{gray}{\ensuremath{\cdot}}\\[3pt]\hline
 & \begin{tabular}[t]{@{}l@{}}\rule{0pt}{12pt}MPF / Scalar\\SIMD-DP / Block\end{tabular} & \textcolor{gray}{\ensuremath{\cdot}} & \textcolor{gray}{\ensuremath{\cdot}} & \textcolor{gray}{\ensuremath{\cdot}} & \textcolor{gray}{\ensuremath{\cdot}} & \textcolor{gray}{\ensuremath{\cdot}} & \textcolor{gray}{\ensuremath{\cdot}} & \textcolor{gray}{\ensuremath{\cdot}} & \textcolor{gray}{\ensuremath{\cdot}} & \hspace*{1pt}\begin{tabular}[t]{@{}c@{}}\rule{0pt}{12pt}4m\,51s\\[-.2pt]{\fontsize{5.8}{6.4}\selectfont \textcolor{BrickRed}{$ \blacktriangledown\,1.4\times $}}\\[-.2pt]4m\,8.2s\\[-.2pt]{\fontsize{5.8}{6.4}\selectfont \textcolor{BrickRed}{$ \blacktriangledown\,22\% $}}\end{tabular} & \hspace*{1pt}\begin{tabular}[t]{@{}c@{}}\rule{0pt}{12pt}4m\,40s\\[-.2pt]{\fontsize{5.8}{6.4}\selectfont \textcolor{BrickRed}{$ \blacktriangledown\,20\% $}}\\[-.2pt]4m\,5.7s\\[-.2pt]{\fontsize{5.8}{6.4}\selectfont \textcolor{BrickRed}{$ \blacktriangledown\,10\% $}}\end{tabular} & \textcolor{gray}{\ensuremath{\cdot}} & \textcolor{gray}{\ensuremath{\cdot}} & \textcolor{gray}{\ensuremath{\cdot}} & \textcolor{gray}{\ensuremath{\cdot}} & \textcolor{gray}{\ensuremath{\cdot}} & \textcolor{gray}{\ensuremath{\cdot}}\\[3pt]\hline
 & \begin{tabular}[t]{@{}l@{}}\rule{0pt}{12pt}MPF / Scalar\\SIMD-DP / KSW2\end{tabular} & \hspace*{1pt}\begin{tabular}[t]{@{}c@{}}\rule{0pt}{12pt}11m\,14.3s\\[-.2pt]{\fontsize{5.8}{6.4}\selectfont \textcolor{BrickRed}{$ \blacktriangledown\,4.0\times $}}\\[-.2pt]11m\,12.2s\\[-.2pt]{\fontsize{5.8}{6.4}\selectfont \textcolor{BrickRed}{$ \blacktriangledown\,4.2\times $}}\end{tabular} & \hspace*{1pt}\begin{tabular}[t]{@{}c@{}}\rule{0pt}{12pt}24m\,37.5s\\[-.2pt]{\fontsize{5.8}{6.4}\selectfont \textcolor{BrickRed}{$ \blacktriangledown\,8.7\times $}}\\[-.2pt]24m\,35.6s\\[-.2pt]{\fontsize{5.8}{6.4}\selectfont \textcolor{BrickRed}{$ \blacktriangledown\,9.1\times $}}\end{tabular} & \textcolor{gray}{\ensuremath{\cdot}} & \textcolor{gray}{\ensuremath{\cdot}} & \textcolor{gray}{\ensuremath{\cdot}} & \textcolor{gray}{\ensuremath{\cdot}} & \textcolor{gray}{\ensuremath{\cdot}} & \textcolor{gray}{\ensuremath{\cdot}} & \hspace*{1pt}\begin{tabular}[t]{@{}c@{}}\rule{0pt}{12pt}4m\,53s\\[-.2pt]{\fontsize{5.8}{6.4}\selectfont \textcolor{BrickRed}{$ \blacktriangledown\,1.4\times $}}\\[-.2pt]4m\,8.2s\\[-.2pt]{\fontsize{5.8}{6.4}\selectfont \textcolor{BrickRed}{$ \blacktriangledown\,22\% $}}\end{tabular} & \hspace*{1pt}\begin{tabular}[t]{@{}c@{}}\rule{0pt}{12pt}5m\,9.3s\\[-.2pt]{\fontsize{5.8}{6.4}\selectfont \textcolor{BrickRed}{$ \blacktriangledown\,1.3\times $}}\\[-.2pt]4m\,31.7s\\[-.2pt]{\fontsize{5.8}{6.4}\selectfont \textcolor{BrickRed}{$ \blacktriangledown\,22\% $}}\end{tabular} & \textcolor{gray}{\ensuremath{\cdot}} & \textcolor{gray}{\ensuremath{\cdot}} & \textcolor{gray}{\ensuremath{\cdot}} & \textcolor{gray}{\ensuremath{\cdot}} & \textcolor{gray}{\ensuremath{\cdot}} & \textcolor{gray}{\ensuremath{\cdot}}\\[3pt]\hline
 & \begin{tabular}[t]{@{}l@{}}\rule{0pt}{12pt}MPF / Two-stack\\RMQ / Block\end{tabular} & \hspace*{1pt}\begin{tabular}[t]{@{}c@{}}\rule{0pt}{12pt}9m\,11s\\[-.2pt]{\fontsize{5.8}{6.4}\selectfont \textcolor{BrickRed}{$ \blacktriangledown\,3.3\times $}}\\[-.2pt]9m\,10.5s\\[-.2pt]{\fontsize{5.8}{6.4}\selectfont \textcolor{BrickRed}{$ \blacktriangledown\,3.4\times $}}\end{tabular} & \hspace*{1pt}\begin{tabular}[t]{@{}c@{}}\rule{0pt}{12pt}13m\,20.9s\\[-.2pt]{\fontsize{5.8}{6.4}\selectfont \textcolor{BrickRed}{$ \blacktriangledown\,4.7\times $}}\\[-.2pt]13m\,20.1s\\[-.2pt]{\fontsize{5.8}{6.4}\selectfont \textcolor{BrickRed}{$ \blacktriangledown\,4.9\times $}}\end{tabular} & \textcolor{gray}{\ensuremath{\cdot}} & \textcolor{gray}{\ensuremath{\cdot}} & \textcolor{gray}{\ensuremath{\cdot}} & \textcolor{gray}{\ensuremath{\cdot}} & \textcolor{gray}{\ensuremath{\cdot}} & \textcolor{gray}{\ensuremath{\cdot}} & \textcolor{gray}{\ensuremath{\cdot}} & \textcolor{gray}{\ensuremath{\cdot}} & \textcolor{gray}{\ensuremath{\cdot}} & \textcolor{gray}{\ensuremath{\cdot}} & \textcolor{gray}{\ensuremath{\cdot}} & \textcolor{gray}{\ensuremath{\cdot}} & \textcolor{gray}{\ensuremath{\cdot}} & \textcolor{gray}{\ensuremath{\cdot}}\\[3pt]\hline
 & \begin{tabular}[t]{@{}l@{}}\rule{0pt}{12pt}MPF / Two-stack\\RMQ / KSW2\end{tabular} & \hspace*{1pt}\begin{tabular}[t]{@{}c@{}}\rule{0pt}{12pt}9m\,9.8s\\[-.2pt]{\fontsize{5.8}{6.4}\selectfont \textcolor{BrickRed}{$ \blacktriangledown\,3.3\times $}}\\[-.2pt]9m\,7.7s\\[-.2pt]{\fontsize{5.8}{6.4}\selectfont \textcolor{BrickRed}{$ \blacktriangledown\,3.4\times $}}\end{tabular} & \textcolor{gray}{\ensuremath{\cdot}} & \textcolor{gray}{\ensuremath{\cdot}} & \textcolor{gray}{\ensuremath{\cdot}} & \textcolor{gray}{\ensuremath{\cdot}} & \textcolor{gray}{\ensuremath{\cdot}} & \textcolor{gray}{\ensuremath{\cdot}} & \textcolor{gray}{\ensuremath{\cdot}} & \textcolor{gray}{\ensuremath{\cdot}} & \textcolor{gray}{\ensuremath{\cdot}} & \textcolor{gray}{\ensuremath{\cdot}} & \textcolor{gray}{\ensuremath{\cdot}} & \textcolor{gray}{\ensuremath{\cdot}} & \textcolor{gray}{\ensuremath{\cdot}} & \textcolor{gray}{\ensuremath{\cdot}} & \textcolor{gray}{\ensuremath{\cdot}}\\[3pt]\hline
 & \begin{tabular}[t]{@{}l@{}}\rule{0pt}{12pt}MPF / Two-stack\\SIMD-DP / Block\end{tabular} & \hspace*{1pt}\begin{tabular}[t]{@{}c@{}}\rule{0pt}{12pt}9m\,11.2s\\[-.2pt]{\fontsize{5.8}{6.4}\selectfont \textcolor{BrickRed}{$ \blacktriangledown\,3.3\times $}}\\[-.2pt]9m\,10.7s\\[-.2pt]{\fontsize{5.8}{6.4}\selectfont \textcolor{BrickRed}{$ \blacktriangledown\,3.4\times $}}\end{tabular} & \textcolor{gray}{\ensuremath{\cdot}} & \textcolor{gray}{\ensuremath{\cdot}} & \textcolor{gray}{\ensuremath{\cdot}} & \textcolor{gray}{\ensuremath{\cdot}} & \textcolor{gray}{\ensuremath{\cdot}} & \textcolor{gray}{\ensuremath{\cdot}} & \textcolor{gray}{\ensuremath{\cdot}} & \textcolor{gray}{\ensuremath{\cdot}} & \textcolor{gray}{\ensuremath{\cdot}} & \textcolor{gray}{\ensuremath{\cdot}} & \textcolor{gray}{\ensuremath{\cdot}} & \textcolor{gray}{\ensuremath{\cdot}} & \textcolor{gray}{\ensuremath{\cdot}} & \textcolor{gray}{\ensuremath{\cdot}} & \textcolor{gray}{\ensuremath{\cdot}}\\[3pt]\hline
 & \begin{tabular}[t]{@{}l@{}}\rule{0pt}{12pt}MPF / Two-stack\\SIMD-DP / KSW2\end{tabular} & \hspace*{1pt}\begin{tabular}[t]{@{}c@{}}\rule{0pt}{12pt}9m\,11.4s\\[-.2pt]{\fontsize{5.8}{6.4}\selectfont \textcolor{BrickRed}{$ \blacktriangledown\,3.3\times $}}\\[-.2pt]9m\,9.5s\\[-.2pt]{\fontsize{5.8}{6.4}\selectfont \textcolor{BrickRed}{$ \blacktriangledown\,3.4\times $}}\end{tabular} & \hspace*{1pt}\begin{tabular}[t]{@{}c@{}}\rule{0pt}{12pt}13m\,22s\\[-.2pt]{\fontsize{5.8}{6.4}\selectfont \textcolor{BrickRed}{$ \blacktriangledown\,4.7\times $}}\\[-.2pt]13m\,19.3s\\[-.2pt]{\fontsize{5.8}{6.4}\selectfont \textcolor{BrickRed}{$ \blacktriangledown\,4.9\times $}}\end{tabular} & \textcolor{gray}{\ensuremath{\cdot}} & \textcolor{gray}{\ensuremath{\cdot}} & \textcolor{gray}{\ensuremath{\cdot}} & \textcolor{gray}{\ensuremath{\cdot}} & \textcolor{gray}{\ensuremath{\cdot}} & \textcolor{gray}{\ensuremath{\cdot}} & \textcolor{gray}{\ensuremath{\cdot}} & \textcolor{gray}{\ensuremath{\cdot}} & \textcolor{gray}{\ensuremath{\cdot}} & \textcolor{gray}{\ensuremath{\cdot}} & \textcolor{gray}{\ensuremath{\cdot}} & \textcolor{gray}{\ensuremath{\cdot}} & \textcolor{gray}{\ensuremath{\cdot}} & \textcolor{gray}{\ensuremath{\cdot}}\\[3pt]\hline
\end{tabular}\par}
\smallskip{\footnotesize\noindent Each cell: A, comparison, B, comparison. Times cover 1,024 queries. $\cdot$ = no approved observation. ---\textsuperscript{g} = GPU unavailable, not executed. n/a = no measured comparison.\par
\noindent \textsuperscript{f}SIMD-ref request observed as scalar, compared only as that effective path. \textsuperscript{p}Accepted by explicit preservation after wrapper postflight failure, no run receipt claimed.\par
\noindent MPF/MPS = minimizers per fragment/sequence. Row order: index representation / minimizer, then chaining / base alignment. SIMD-ref is the reference minimizer, Bounded is bounded ntHash.\par
\noindent Blue = one lowest A and one lowest B down each machine/workload column, across continuation pages. Exact ties select the first row. Repeated baselines are marked only on page 1.\par
\noindent Light gray = minimap2 baseline. Pale yellow = preferred MPF / Bounded / SIMD-DP / Block configuration, not a claim of winning every column.\par
\noindent When all Vulkan modes are unmeasured, they share one row. Baseline comparisons use the fastest measured BioseqDB effective path separately for A and B.\par}
\clearpage
\noindent{\large\bfseries AWS: 4 Gbp, k=20, w=15 (mapping interval)\hfill 3/4}\par\smallskip
{\fontsize{7.1}{8.1}\selectfont\setlength{\tabcolsep}{1pt}\setlength{\fboxsep}{2pt}\renewcommand{\arraystretch}{1.12}
\noindent\begin{tabular}{@{}p{14mm}p{27mm}|p{13.55mm}|p{13.55mm}|p{13.55mm}|p{13.55mm}|p{13.55mm}|p{13.55mm}|p{13.55mm}|p{13.55mm}|p{13.55mm}|p{13.55mm}|p{13.55mm}|p{13.55mm}|p{13.55mm}|p{13.55mm}|p{13.55mm}|p{13.55mm}|@{}}\hline
Engine & Execution & \multicolumn{8}{c|}{\textbf{Index}} & \multicolumn{8}{c|}{\textbf{No index}}\\
 &  & \multicolumn{8}{c|}{\textbf{4 Gbp}} & \multicolumn{8}{c|}{\textbf{4 Gbp}}\\
 &  & \multicolumn{8}{c|}{\textbf{k=20, w=15}} & \multicolumn{8}{c|}{\textbf{k=20, w=15}}\\
 &  & \multicolumn{4}{c|}{\textbf{DNA4}} & \multicolumn{4}{c|}{\textbf{DNA5}} & \multicolumn{4}{c|}{\textbf{DNA4}} & \multicolumn{4}{c|}{\textbf{DNA5}}\\
 &  & \multicolumn{1}{c|}{\textbf{M7A}} & \multicolumn{1}{c|}{\textbf{M7I}} & \multicolumn{1}{c|}{\textbf{M4}} & \multicolumn{1}{c|}{\textbf{C7G}} & \multicolumn{1}{c|}{\textbf{M7A}} & \multicolumn{1}{c|}{\textbf{M7I}} & \multicolumn{1}{c|}{\textbf{M4}} & \multicolumn{1}{c|}{\textbf{C7G}} & \multicolumn{1}{c|}{\textbf{M7A}} & \multicolumn{1}{c|}{\textbf{M7I}} & \multicolumn{1}{c|}{\textbf{M4}} & \multicolumn{1}{c|}{\textbf{C7G}} & \multicolumn{1}{c|}{\textbf{M7A}} & \multicolumn{1}{c|}{\textbf{M7I}} & \multicolumn{1}{c|}{\textbf{M4}} & \multicolumn{1}{c|}{\textbf{C7G}}\\
\hline
\rowcolor{BenchBaseline}[0.5pt][0.5pt]\textbf{minimap2} & \textbf{Baseline} & \hspace*{1pt}\begin{tabular}[t]{@{}c@{}}\rule{0pt}{12pt}2m\,48.1s\\[-.2pt]{\fontsize{5.8}{6.4}\selectfont \textcolor{BrickRed}{$ \blacktriangledown\,20.4\times $}}\\[-.2pt]2m\,41.7s\\[-.2pt]{\fontsize{5.8}{6.4}\selectfont \textcolor{BrickRed}{$ \blacktriangledown\,22.6\times $}}\end{tabular} & \hspace*{1pt}\begin{tabular}[t]{@{}c@{}}\rule{0pt}{12pt}2m\,50.3s\\[-.2pt]{\fontsize{5.8}{6.4}\selectfont \textcolor{BrickRed}{$ \blacktriangledown\,19.8\times $}}\\[-.2pt]2m\,43s\\[-.2pt]{\fontsize{5.8}{6.4}\selectfont \textcolor{BrickRed}{$ \blacktriangledown\,22.2\times $}}\end{tabular} & \hspace*{1pt}\begin{tabular}[t]{@{}c@{}}\rule{0pt}{12pt}4m\,54.3s\\[-.2pt]{\fontsize{5.8}{6.4}\selectfont \textcolor{BrickRed}{$ \blacktriangledown\,43.2\times $}}\\[-.2pt]4m\,39s\\[-.2pt]{\fontsize{5.8}{6.4}\selectfont \textcolor{BrickRed}{$ \blacktriangledown\,54.6\times $}}\end{tabular} & \hspace*{1pt}\begin{tabular}[t]{@{}c@{}}\rule{0pt}{12pt}25m\,5.5s\\[-.2pt]{\fontsize{5.8}{6.4}\selectfont \textcolor{gray}{n/a}}\\[-.2pt]25m\,0.7s\\[-.2pt]{\fontsize{5.8}{6.4}\selectfont \textcolor{gray}{n/a}}\end{tabular} & \hspace*{1pt}\begin{tabular}[t]{@{}c@{}}\rule{0pt}{12pt}33.9s\\[-.2pt]{\fontsize{5.8}{6.4}\selectfont \textcolor{BrickRed}{$ \blacktriangledown\,1.9\times $}}\\[-.2pt]32.9s\\[-.2pt]{\fontsize{5.8}{6.4}\selectfont \textcolor{BrickRed}{$ \blacktriangledown\,2.1\times $}}\end{tabular} & \hspace*{1pt}\begin{tabular}[t]{@{}c@{}}\rule{0pt}{12pt}34.4s\\[-.2pt]{\fontsize{5.8}{6.4}\selectfont \textcolor{BrickRed}{$ \blacktriangledown\,2.6\times $}}\\[-.2pt]32.7s\\[-.2pt]{\fontsize{5.8}{6.4}\selectfont \textcolor{BrickRed}{$ \blacktriangledown\,2.9\times $}}\end{tabular} & \hspace*{1pt}\begin{tabular}[t]{@{}c@{}}\rule{0pt}{12pt}1m\,8.2s\\[-.2pt]{\fontsize{5.8}{6.4}\selectfont \textcolor{BrickRed}{$ \blacktriangledown\,3.7\times $}}\\[-.2pt]1m\,5.1s\\[-.2pt]{\fontsize{5.8}{6.4}\selectfont \textcolor{BrickRed}{$ \blacktriangledown\,4.6\times $}}\end{tabular} & \hspace*{1pt}\begin{tabular}[t]{@{}c@{}}\rule{0pt}{12pt}3m\,9.8s\\[-.2pt]{\fontsize{5.8}{6.4}\selectfont \textcolor{gray}{n/a}}\\[-.2pt]3m\,3.7s\\[-.2pt]{\fontsize{5.8}{6.4}\selectfont \textcolor{gray}{n/a}}\end{tabular} & \hspace*{1pt}\begin{tabular}[t]{@{}c@{}}\rule{0pt}{12pt}3m\,28.3s\\[-.2pt]{\fontsize{5.8}{6.4}\selectfont \textcolor{BrickRed}{$ \blacktriangledown\,1.3\times $}}\\[-.2pt]3m\,22.8s\\[-.2pt]{\fontsize{5.8}{6.4}\selectfont \textcolor{BrickRed}{$ \blacktriangledown\,1.4\times $}}\end{tabular} & \hspace*{1pt}\begin{tabular}[t]{@{}c@{}}\rule{0pt}{12pt}3m\,53.2s\\[-.2pt]{\fontsize{5.8}{6.4}\selectfont \textcolor{BrickRed}{$ \blacktriangledown\,1.4\times $}}\\[-.2pt]3m\,43.5s\\[-.2pt]{\fontsize{5.8}{6.4}\selectfont \textcolor{BrickRed}{$ \blacktriangledown\,1.4\times $}}\end{tabular} & \hspace*{1pt}\begin{tabular}[t]{@{}c@{}}\rule{0pt}{12pt}7m\,14.6s\\[-.2pt]{\fontsize{5.8}{6.4}\selectfont \textcolor{BrickRed}{$ \blacktriangledown\,1.5\times $}}\\[-.2pt]6m\,57.2s\\[-.2pt]{\fontsize{5.8}{6.4}\selectfont \textcolor{BrickRed}{$ \blacktriangledown\,1.6\times $}}\end{tabular} & \hspace*{1pt}\begin{tabular}[t]{@{}c@{}}\rule{0pt}{12pt}51m\,22s\\[-.2pt]{\fontsize{5.8}{6.4}\selectfont \textcolor{gray}{n/a}}\\[-.2pt]51m\,17.2s\\[-.2pt]{\fontsize{5.8}{6.4}\selectfont \textcolor{gray}{n/a}}\end{tabular} & \hspace*{1pt}\begin{tabular}[t]{@{}c@{}}\rule{0pt}{12pt}1m\,10.3s\\[-.2pt]{\fontsize{5.8}{6.4}\selectfont \textcolor{ForestGreen}{$ \blacktriangle\,7.6\times $}}\\[-.2pt]1m\,9s\\[-.2pt]{\fontsize{5.8}{6.4}\selectfont \textcolor{ForestGreen}{$ \blacktriangle\,2.4\times $}}\end{tabular} & \hspace*{1pt}\begin{tabular}[t]{@{}c@{}}\rule{0pt}{12pt}1m\,8.6s\\[-.2pt]{\fontsize{5.8}{6.4}\selectfont \textcolor{ForestGreen}{$ \blacktriangle\,6.4\times $}}\\[-.2pt]1m\,6.9s\\[-.2pt]{\fontsize{5.8}{6.4}\selectfont \textcolor{ForestGreen}{$ \blacktriangle\,2.5\times $}}\end{tabular} & \hspace*{1pt}\begin{tabular}[t]{@{}c@{}}\rule{0pt}{12pt}2m\,21.5s\\[-.2pt]{\fontsize{5.8}{6.4}\selectfont \textcolor{ForestGreen}{$ \blacktriangle\,6.0\times $}}\\[-.2pt]2m\,18.4s\\[-.2pt]{\fontsize{5.8}{6.4}\selectfont \textcolor{ForestGreen}{$ \blacktriangle\,2.2\times $}}\end{tabular} & \hspace*{1pt}\begin{tabular}[t]{@{}c@{}}\rule{0pt}{12pt}15m\,56.8s\\[-.2pt]{\fontsize{5.8}{6.4}\selectfont \textcolor{gray}{n/a}}\\[-.2pt]15m\,51.9s\\[-.2pt]{\fontsize{5.8}{6.4}\selectfont \textcolor{gray}{n/a}}\end{tabular}\\[3pt]\hline
\textbf{BioseqDB} & \begin{tabular}[t]{@{}l@{}}\rule{0pt}{12pt}MPS / Bounded\\RMQ / Block\end{tabular} & \hspace*{1pt}\begin{tabular}[t]{@{}c@{}}\rule{0pt}{12pt}11.7s\\[-.2pt]{\fontsize{5.8}{6.4}\selectfont \textcolor{ForestGreen}{$ \blacktriangle\,14.3\times $}}\\[-.2pt]11s\\[-.2pt]{\fontsize{5.8}{6.4}\selectfont \textcolor{ForestGreen}{$ \blacktriangle\,14.6\times $}}\end{tabular} & \hspace*{1pt}\begin{tabular}[t]{@{}c@{}}\rule{0pt}{12pt}11.8s\\[-.2pt]{\fontsize{5.8}{6.4}\selectfont \textcolor{ForestGreen}{$ \blacktriangle\,14.5\times $}}\\[-.2pt]11.1s\\[-.2pt]{\fontsize{5.8}{6.4}\selectfont \textcolor{ForestGreen}{$ \blacktriangle\,14.7\times $}}\end{tabular} & \textcolor{gray}{\ensuremath{\cdot}} & \textcolor{gray}{\ensuremath{\cdot}} & \textcolor{gray}{\ensuremath{\cdot}} & \textcolor{gray}{\ensuremath{\cdot}} & \textcolor{gray}{\ensuremath{\cdot}} & \textcolor{gray}{\ensuremath{\cdot}} & \hspace*{1pt}\begin{tabular}[t]{@{}c@{}}\rule{0pt}{12pt}2m\,43.8s\\[-.2pt]{\fontsize{5.8}{6.4}\selectfont \textcolor{ForestGreen}{$ \blacktriangle\,21\% $}}\\[-.2pt]2m\,34.7s\\[-.2pt]{\fontsize{5.8}{6.4}\selectfont \textcolor{ForestGreen}{$ \blacktriangle\,24\% $}}\end{tabular} & \hspace*{1pt}\begin{tabular}[t]{@{}c@{}}\rule{0pt}{12pt}2m\,52.1s\\[-.2pt]{\fontsize{5.8}{6.4}\selectfont \textcolor{ForestGreen}{$ \blacktriangle\,26\% $}}\\[-.2pt]2m\,39.6s\\[-.2pt]{\fontsize{5.8}{6.4}\selectfont \textcolor{ForestGreen}{$ \blacktriangle\,29\% $}}\end{tabular} & \textcolor{gray}{\ensuremath{\cdot}} & \textcolor{gray}{\ensuremath{\cdot}} & \textcolor{gray}{\ensuremath{\cdot}} & \textcolor{gray}{\ensuremath{\cdot}} & \hspace*{1pt}\begin{tabular}[t]{@{}c@{}}\rule{0pt}{12pt}14m\,16s\\[-.2pt]{\fontsize{5.8}{6.4}\selectfont \textcolor{BrickRed}{$ \blacktriangledown\,6.0\times $}}\\[-.2pt]5m\,26.4s\\[-.2pt]{\fontsize{5.8}{6.4}\selectfont \textcolor{BrickRed}{$ \blacktriangledown\,2.4\times $}}\end{tabular} & \textcolor{gray}{\ensuremath{\cdot}}\\[3pt]\hline
 & \begin{tabular}[t]{@{}l@{}}\rule{0pt}{12pt}MPS / Bounded\\RMQ / KSW2\end{tabular} & \hspace*{1pt}\begin{tabular}[t]{@{}c@{}}\rule{0pt}{12pt}13.5s\\[-.2pt]{\fontsize{5.8}{6.4}\selectfont \textcolor{ForestGreen}{$ \blacktriangle\,12.4\times $}}\\[-.2pt]11s\\[-.2pt]{\fontsize{5.8}{6.4}\selectfont \textcolor{ForestGreen}{$ \blacktriangle\,14.7\times $}}\end{tabular} & \hspace*{1pt}\begin{tabular}[t]{@{}c@{}}\rule{0pt}{12pt}12.8s\\[-.2pt]{\fontsize{5.8}{6.4}\selectfont \textcolor{ForestGreen}{$ \blacktriangle\,13.3\times $}}\\[-.2pt]10.6s\\[-.2pt]{\fontsize{5.8}{6.4}\selectfont \textcolor{ForestGreen}{$ \blacktriangle\,15.4\times $}}\end{tabular} & \textcolor{gray}{\ensuremath{\cdot}} & \textcolor{gray}{\ensuremath{\cdot}} & \textcolor{gray}{\ensuremath{\cdot}} & \textcolor{gray}{\ensuremath{\cdot}} & \textcolor{gray}{\ensuremath{\cdot}} & \textcolor{gray}{\ensuremath{\cdot}} & \hspace*{1pt}\begin{tabular}[t]{@{}c@{}}\rule{0pt}{12pt}2m\,56.8s\\[-.2pt]{\fontsize{5.8}{6.4}\selectfont \textcolor{ForestGreen}{$ \blacktriangle\,15\% $}}\\[-.2pt]2m\,44.8s\\[-.2pt]{\fontsize{5.8}{6.4}\selectfont \textcolor{ForestGreen}{$ \blacktriangle\,19\% $}}\end{tabular} & \hspace*{1pt}\begin{tabular}[t]{@{}c@{}}\rule{0pt}{12pt}3m\,11.4s\\[-.2pt]{\fontsize{5.8}{6.4}\selectfont \textcolor{ForestGreen}{$ \blacktriangle\,18\% $}}\\[-.2pt]2m\,56.8s\\[-.2pt]{\fontsize{5.8}{6.4}\selectfont \textcolor{ForestGreen}{$ \blacktriangle\,21\% $}}\end{tabular} & \textcolor{gray}{\ensuremath{\cdot}} & \textcolor{gray}{\ensuremath{\cdot}} & \textcolor{gray}{\ensuremath{\cdot}} & \textcolor{gray}{\ensuremath{\cdot}} & \textcolor{gray}{\ensuremath{\cdot}} & \textcolor{gray}{\ensuremath{\cdot}}\\[3pt]\hline
 & \begin{tabular}[t]{@{}l@{}}\rule{0pt}{12pt}MPS / Bounded\\Scalar-DP / Block\end{tabular} & \hspace*{1pt}\begin{tabular}[t]{@{}c@{}}\rule{0pt}{12pt}11.8s\\[-.2pt]{\fontsize{5.8}{6.4}\selectfont \textcolor{ForestGreen}{$ \blacktriangle\,14.2\times $}}\\[-.2pt]11.1s\\[-.2pt]{\fontsize{5.8}{6.4}\selectfont \textcolor{ForestGreen}{$ \blacktriangle\,14.6\times $}}\end{tabular} & \hspace*{1pt}\begin{tabular}[t]{@{}c@{}}\rule{0pt}{12pt}12.7s\\[-.2pt]{\fontsize{5.8}{6.4}\selectfont \textcolor{ForestGreen}{$ \blacktriangle\,13.4\times $}}\\[-.2pt]11.9s\\[-.2pt]{\fontsize{5.8}{6.4}\selectfont \textcolor{ForestGreen}{$ \blacktriangle\,13.7\times $}}\end{tabular} & \textcolor{gray}{\ensuremath{\cdot}} & \textcolor{gray}{\ensuremath{\cdot}} & \textcolor{gray}{\ensuremath{\cdot}} & \textcolor{gray}{\ensuremath{\cdot}} & \textcolor{gray}{\ensuremath{\cdot}} & \textcolor{gray}{\ensuremath{\cdot}} & \hspace*{1pt}\begin{tabular}[t]{@{}c@{}}\rule{0pt}{12pt}2m\,52.6s\\[-.2pt]{\fontsize{5.8}{6.4}\selectfont \textcolor{ForestGreen}{$ \blacktriangle\,17\% $}}\\[-.2pt]2m\,42.8s\\[-.2pt]{\fontsize{5.8}{6.4}\selectfont \textcolor{ForestGreen}{$ \blacktriangle\,20\% $}}\end{tabular} & \hspace*{1pt}\begin{tabular}[t]{@{}c@{}}\rule{0pt}{12pt}2m\,58.9s\\[-.2pt]{\fontsize{5.8}{6.4}\selectfont \textcolor{ForestGreen}{$ \blacktriangle\,23\% $}}\\[-.2pt]2m\,46.9s\\[-.2pt]{\fontsize{5.8}{6.4}\selectfont \textcolor{ForestGreen}{$ \blacktriangle\,25\% $}}\end{tabular} & \textcolor{gray}{\ensuremath{\cdot}} & \textcolor{gray}{\ensuremath{\cdot}} & \textcolor{gray}{\ensuremath{\cdot}} & \textcolor{gray}{\ensuremath{\cdot}} & \textcolor{gray}{\ensuremath{\cdot}} & \textcolor{gray}{\ensuremath{\cdot}}\\[3pt]\hline
 & \begin{tabular}[t]{@{}l@{}}\rule{0pt}{12pt}MPS / Bounded\\Scalar-DP / KSW2\end{tabular} & \hspace*{1pt}\begin{tabular}[t]{@{}c@{}}\rule{0pt}{12pt}13.4s\\[-.2pt]{\fontsize{5.8}{6.4}\selectfont \textcolor{ForestGreen}{$ \blacktriangle\,12.6\times $}}\\[-.2pt]10.9s\\[-.2pt]{\fontsize{5.8}{6.4}\selectfont \textcolor{ForestGreen}{$ \blacktriangle\,14.9\times $}}\end{tabular} & \hspace*{1pt}\begin{tabular}[t]{@{}c@{}}\rule{0pt}{12pt}13s\\[-.2pt]{\fontsize{5.8}{6.4}\selectfont \textcolor{ForestGreen}{$ \blacktriangle\,13.1\times $}}\\[-.2pt]10.7s\\[-.2pt]{\fontsize{5.8}{6.4}\selectfont \textcolor{ForestGreen}{$ \blacktriangle\,15.2\times $}}\end{tabular} & \textcolor{gray}{\ensuremath{\cdot}} & \textcolor{gray}{\ensuremath{\cdot}} & \textcolor{gray}{\ensuremath{\cdot}} & \textcolor{gray}{\ensuremath{\cdot}} & \textcolor{gray}{\ensuremath{\cdot}} & \textcolor{gray}{\ensuremath{\cdot}} & \hspace*{1pt}\begin{tabular}[t]{@{}c@{}}\rule{0pt}{12pt}2m\,59.1s\\[-.2pt]{\fontsize{5.8}{6.4}\selectfont \textcolor{ForestGreen}{$ \blacktriangle\,14\% $}}\\[-.2pt]2m\,47.7s\\[-.2pt]{\fontsize{5.8}{6.4}\selectfont \textcolor{ForestGreen}{$ \blacktriangle\,17\% $}}\end{tabular} & \hspace*{1pt}\begin{tabular}[t]{@{}c@{}}\rule{0pt}{12pt}2m\,50.8s\\[-.2pt]{\fontsize{5.8}{6.4}\selectfont \textcolor{ForestGreen}{$ \blacktriangle\,27\% $}}\\[-.2pt]2m\,36.9s\\[-.2pt]{\fontsize{5.8}{6.4}\selectfont \textcolor{ForestGreen}{$ \blacktriangle\,30\% $}}\end{tabular} & \textcolor{gray}{\ensuremath{\cdot}} & \textcolor{gray}{\ensuremath{\cdot}} & \textcolor{gray}{\ensuremath{\cdot}} & \textcolor{gray}{\ensuremath{\cdot}} & \textcolor{gray}{\ensuremath{\cdot}} & \textcolor{gray}{\ensuremath{\cdot}}\\[3pt]\hline
 & \begin{tabular}[t]{@{}l@{}}\rule{0pt}{12pt}MPS / Bounded\\SIMD-DP / Block\end{tabular} & \hspace*{1pt}\begin{tabular}[t]{@{}c@{}}\rule{0pt}{12pt}11.7s\\[-.2pt]{\fontsize{5.8}{6.4}\selectfont \textcolor{ForestGreen}{$ \blacktriangle\,14.3\times $}}\\[-.2pt]11s\\[-.2pt]{\fontsize{5.8}{6.4}\selectfont \textcolor{ForestGreen}{$ \blacktriangle\,14.7\times $}}\end{tabular} & \hspace*{1pt}\begin{tabular}[t]{@{}c@{}}\rule{0pt}{12pt}11.9s\\[-.2pt]{\fontsize{5.8}{6.4}\selectfont \textcolor{ForestGreen}{$ \blacktriangle\,14.3\times $}}\\[-.2pt]11.2s\\[-.2pt]{\fontsize{5.8}{6.4}\selectfont \textcolor{ForestGreen}{$ \blacktriangle\,14.6\times $}}\end{tabular} & \textcolor{gray}{\ensuremath{\cdot}} & \textcolor{gray}{\ensuremath{\cdot}} & \textcolor{gray}{\ensuremath{\cdot}} & \textcolor{gray}{\ensuremath{\cdot}} & \textcolor{gray}{\ensuremath{\cdot}} & \textcolor{gray}{\ensuremath{\cdot}} & \hspace*{1pt}\begin{tabular}[t]{@{}c@{}}\rule{0pt}{12pt}2m\,43.8s\\[-.2pt]{\fontsize{5.8}{6.4}\selectfont \textcolor{ForestGreen}{$ \blacktriangle\,21\% $}}\\[-.2pt]2m\,34.7s\\[-.2pt]{\fontsize{5.8}{6.4}\selectfont \textcolor{ForestGreen}{$ \blacktriangle\,24\% $}}\end{tabular} & \hspace*{1pt}\begin{tabular}[t]{@{}c@{}}\rule{0pt}{12pt}2m\,58.9s\\[-.2pt]{\fontsize{5.8}{6.4}\selectfont \textcolor{ForestGreen}{$ \blacktriangle\,23\% $}}\\[-.2pt]2m\,46.3s\\[-.2pt]{\fontsize{5.8}{6.4}\selectfont \textcolor{ForestGreen}{$ \blacktriangle\,26\% $}}\end{tabular} & \textcolor{gray}{\ensuremath{\cdot}} & \textcolor{gray}{\ensuremath{\cdot}} & \textcolor{gray}{\ensuremath{\cdot}} & \textcolor{gray}{\ensuremath{\cdot}} & \textcolor{gray}{\ensuremath{\cdot}} & \textcolor{gray}{\ensuremath{\cdot}}\\[3pt]\hline
 & \begin{tabular}[t]{@{}l@{}}\rule{0pt}{12pt}MPS / Bounded\\SIMD-DP / KSW2\end{tabular} & \hspace*{1pt}\begin{tabular}[t]{@{}c@{}}\rule{0pt}{12pt}13.3s\\[-.2pt]{\fontsize{5.8}{6.4}\selectfont \textcolor{ForestGreen}{$ \blacktriangle\,12.7\times $}}\\[-.2pt]10.8s\\[-.2pt]{\fontsize{5.8}{6.4}\selectfont \textcolor{ForestGreen}{$ \blacktriangle\,15.0\times $}}\end{tabular} & \hspace*{1pt}\begin{tabular}[t]{@{}c@{}}\rule{0pt}{12pt}13.6s\\[-.2pt]{\fontsize{5.8}{6.4}\selectfont \textcolor{ForestGreen}{$ \blacktriangle\,12.5\times $}}\\[-.2pt]11.1s\\[-.2pt]{\fontsize{5.8}{6.4}\selectfont \textcolor{ForestGreen}{$ \blacktriangle\,14.7\times $}}\end{tabular} & \textcolor{gray}{\ensuremath{\cdot}} & \textcolor{gray}{\ensuremath{\cdot}} & \textcolor{gray}{\ensuremath{\cdot}} & \textcolor{gray}{\ensuremath{\cdot}} & \textcolor{gray}{\ensuremath{\cdot}} & \textcolor{gray}{\ensuremath{\cdot}} & \hspace*{1pt}\begin{tabular}[t]{@{}c@{}}\rule{0pt}{12pt}2m\,59s\\[-.2pt]{\fontsize{5.8}{6.4}\selectfont \textcolor{ForestGreen}{$ \blacktriangle\,14\% $}}\\[-.2pt]2m\,48.5s\\[-.2pt]{\fontsize{5.8}{6.4}\selectfont \textcolor{ForestGreen}{$ \blacktriangle\,17\% $}}\end{tabular} & \hspace*{1pt}\begin{tabular}[t]{@{}c@{}}\rule{0pt}{12pt}2m\,52.5s\\[-.2pt]{\fontsize{5.8}{6.4}\selectfont \textcolor{ForestGreen}{$ \blacktriangle\,26\% $}}\\[-.2pt]2m\,38.8s\\[-.2pt]{\fontsize{5.8}{6.4}\selectfont \textcolor{ForestGreen}{$ \blacktriangle\,29\% $}}\end{tabular} & \textcolor{gray}{\ensuremath{\cdot}} & \textcolor{gray}{\ensuremath{\cdot}} & \textcolor{gray}{\ensuremath{\cdot}} & \textcolor{gray}{\ensuremath{\cdot}} & \textcolor{gray}{\ensuremath{\cdot}} & \textcolor{gray}{\ensuremath{\cdot}}\\[3pt]\hline
 & \begin{tabular}[t]{@{}l@{}}\rule{0pt}{12pt}MPS / Scalar\\RMQ / Block\end{tabular} & \hspace*{1pt}\begin{tabular}[t]{@{}c@{}}\rule{0pt}{12pt}57.4s\\[-.2pt]{\fontsize{5.8}{6.4}\selectfont \textcolor{ForestGreen}{$ \blacktriangle\,2.9\times $}}\\[-.2pt]14.6s\\[-.2pt]{\fontsize{5.8}{6.4}\selectfont \textcolor{ForestGreen}{$ \blacktriangle\,11.1\times $}}\end{tabular} & \hspace*{1pt}\begin{tabular}[t]{@{}c@{}}\rule{0pt}{12pt}51.1s\\[-.2pt]{\fontsize{5.8}{6.4}\selectfont \textcolor{ForestGreen}{$ \blacktriangle\,3.3\times $}}\\[-.2pt]15.9s\\[-.2pt]{\fontsize{5.8}{6.4}\selectfont \textcolor{ForestGreen}{$ \blacktriangle\,10.3\times $}}\end{tabular} & \textcolor{gray}{\ensuremath{\cdot}} & \textcolor{gray}{\ensuremath{\cdot}} & \textcolor{gray}{\ensuremath{\cdot}} & \textcolor{gray}{\ensuremath{\cdot}} & \textcolor{gray}{\ensuremath{\cdot}} & \textcolor{gray}{\ensuremath{\cdot}} & \hspace*{1pt}\begin{tabular}[t]{@{}c@{}}\rule{0pt}{12pt}5m\,10.6s\\[-.2pt]{\fontsize{5.8}{6.4}\selectfont \textcolor{BrickRed}{$ \blacktriangledown\,1.5\times $}}\\[-.2pt]4m\,28.1s\\[-.2pt]{\fontsize{5.8}{6.4}\selectfont \textcolor{BrickRed}{$ \blacktriangledown\,1.3\times $}}\end{tabular} & \hspace*{1pt}\begin{tabular}[t]{@{}c@{}}\rule{0pt}{12pt}5m\,37s\\[-.2pt]{\fontsize{5.8}{6.4}\selectfont \textcolor{BrickRed}{$ \blacktriangledown\,1.4\times $}}\\[-.2pt]5m\,0.7s\\[-.2pt]{\fontsize{5.8}{6.4}\selectfont \textcolor{BrickRed}{$ \blacktriangledown\,1.3\times $}}\end{tabular} & \textcolor{gray}{\ensuremath{\cdot}} & \textcolor{gray}{\ensuremath{\cdot}} & \textcolor{gray}{\ensuremath{\cdot}} & \textcolor{gray}{\ensuremath{\cdot}} & \textcolor{gray}{\ensuremath{\cdot}} & \textcolor{gray}{\ensuremath{\cdot}}\\[3pt]\hline
\end{tabular}\par}
\smallskip{\footnotesize\noindent Each cell: A, comparison, B, comparison. Times cover 1,024 queries. $\cdot$ = no approved observation. ---\textsuperscript{g} = GPU unavailable, not executed. n/a = no measured comparison.\par
\noindent \textsuperscript{f}SIMD-ref request observed as scalar, compared only as that effective path. \textsuperscript{p}Accepted by explicit preservation after wrapper postflight failure, no run receipt claimed.\par
\noindent MPF/MPS = minimizers per fragment/sequence. Row order: index representation / minimizer, then chaining / base alignment. SIMD-ref is the reference minimizer, Bounded is bounded ntHash.\par
\noindent Blue = one lowest A and one lowest B down each machine/workload column, across continuation pages. Exact ties select the first row. Repeated baselines are marked only on page 1.\par
\noindent Light gray = minimap2 baseline. Pale yellow = preferred MPF / Bounded / SIMD-DP / Block configuration, not a claim of winning every column.\par
\noindent When all Vulkan modes are unmeasured, they share one row. Baseline comparisons use the fastest measured BioseqDB effective path separately for A and B.\par}
\clearpage
\noindent{\large\bfseries AWS: 4 Gbp, k=20, w=15 (mapping interval)\hfill 4/4}\par\smallskip
{\fontsize{7.1}{8.1}\selectfont\setlength{\tabcolsep}{1pt}\setlength{\fboxsep}{2pt}\renewcommand{\arraystretch}{1.12}
\noindent\begin{tabular}{@{}p{14mm}p{27mm}|p{13.55mm}|p{13.55mm}|p{13.55mm}|p{13.55mm}|p{13.55mm}|p{13.55mm}|p{13.55mm}|p{13.55mm}|p{13.55mm}|p{13.55mm}|p{13.55mm}|p{13.55mm}|p{13.55mm}|p{13.55mm}|p{13.55mm}|p{13.55mm}|@{}}\hline
Engine & Execution & \multicolumn{8}{c|}{\textbf{Index}} & \multicolumn{8}{c|}{\textbf{No index}}\\
 &  & \multicolumn{8}{c|}{\textbf{4 Gbp}} & \multicolumn{8}{c|}{\textbf{4 Gbp}}\\
 &  & \multicolumn{8}{c|}{\textbf{k=20, w=15}} & \multicolumn{8}{c|}{\textbf{k=20, w=15}}\\
 &  & \multicolumn{4}{c|}{\textbf{DNA4}} & \multicolumn{4}{c|}{\textbf{DNA5}} & \multicolumn{4}{c|}{\textbf{DNA4}} & \multicolumn{4}{c|}{\textbf{DNA5}}\\
 &  & \multicolumn{1}{c|}{\textbf{M7A}} & \multicolumn{1}{c|}{\textbf{M7I}} & \multicolumn{1}{c|}{\textbf{M4}} & \multicolumn{1}{c|}{\textbf{C7G}} & \multicolumn{1}{c|}{\textbf{M7A}} & \multicolumn{1}{c|}{\textbf{M7I}} & \multicolumn{1}{c|}{\textbf{M4}} & \multicolumn{1}{c|}{\textbf{C7G}} & \multicolumn{1}{c|}{\textbf{M7A}} & \multicolumn{1}{c|}{\textbf{M7I}} & \multicolumn{1}{c|}{\textbf{M4}} & \multicolumn{1}{c|}{\textbf{C7G}} & \multicolumn{1}{c|}{\textbf{M7A}} & \multicolumn{1}{c|}{\textbf{M7I}} & \multicolumn{1}{c|}{\textbf{M4}} & \multicolumn{1}{c|}{\textbf{C7G}}\\
\hline
\rowcolor{BenchBaseline}[0.5pt][0.5pt]\textbf{minimap2} & \textbf{Baseline} & \hspace*{1pt}\begin{tabular}[t]{@{}c@{}}\rule{0pt}{12pt}2m\,48.1s\\[-.2pt]{\fontsize{5.8}{6.4}\selectfont \textcolor{BrickRed}{$ \blacktriangledown\,20.4\times $}}\\[-.2pt]2m\,41.7s\\[-.2pt]{\fontsize{5.8}{6.4}\selectfont \textcolor{BrickRed}{$ \blacktriangledown\,22.6\times $}}\end{tabular} & \hspace*{1pt}\begin{tabular}[t]{@{}c@{}}\rule{0pt}{12pt}2m\,50.3s\\[-.2pt]{\fontsize{5.8}{6.4}\selectfont \textcolor{BrickRed}{$ \blacktriangledown\,19.8\times $}}\\[-.2pt]2m\,43s\\[-.2pt]{\fontsize{5.8}{6.4}\selectfont \textcolor{BrickRed}{$ \blacktriangledown\,22.2\times $}}\end{tabular} & \hspace*{1pt}\begin{tabular}[t]{@{}c@{}}\rule{0pt}{12pt}4m\,54.3s\\[-.2pt]{\fontsize{5.8}{6.4}\selectfont \textcolor{BrickRed}{$ \blacktriangledown\,43.2\times $}}\\[-.2pt]4m\,39s\\[-.2pt]{\fontsize{5.8}{6.4}\selectfont \textcolor{BrickRed}{$ \blacktriangledown\,54.6\times $}}\end{tabular} & \hspace*{1pt}\begin{tabular}[t]{@{}c@{}}\rule{0pt}{12pt}25m\,5.5s\\[-.2pt]{\fontsize{5.8}{6.4}\selectfont \textcolor{gray}{n/a}}\\[-.2pt]25m\,0.7s\\[-.2pt]{\fontsize{5.8}{6.4}\selectfont \textcolor{gray}{n/a}}\end{tabular} & \hspace*{1pt}\begin{tabular}[t]{@{}c@{}}\rule{0pt}{12pt}33.9s\\[-.2pt]{\fontsize{5.8}{6.4}\selectfont \textcolor{BrickRed}{$ \blacktriangledown\,1.9\times $}}\\[-.2pt]32.9s\\[-.2pt]{\fontsize{5.8}{6.4}\selectfont \textcolor{BrickRed}{$ \blacktriangledown\,2.1\times $}}\end{tabular} & \hspace*{1pt}\begin{tabular}[t]{@{}c@{}}\rule{0pt}{12pt}34.4s\\[-.2pt]{\fontsize{5.8}{6.4}\selectfont \textcolor{BrickRed}{$ \blacktriangledown\,2.6\times $}}\\[-.2pt]32.7s\\[-.2pt]{\fontsize{5.8}{6.4}\selectfont \textcolor{BrickRed}{$ \blacktriangledown\,2.9\times $}}\end{tabular} & \hspace*{1pt}\begin{tabular}[t]{@{}c@{}}\rule{0pt}{12pt}1m\,8.2s\\[-.2pt]{\fontsize{5.8}{6.4}\selectfont \textcolor{BrickRed}{$ \blacktriangledown\,3.7\times $}}\\[-.2pt]1m\,5.1s\\[-.2pt]{\fontsize{5.8}{6.4}\selectfont \textcolor{BrickRed}{$ \blacktriangledown\,4.6\times $}}\end{tabular} & \hspace*{1pt}\begin{tabular}[t]{@{}c@{}}\rule{0pt}{12pt}3m\,9.8s\\[-.2pt]{\fontsize{5.8}{6.4}\selectfont \textcolor{gray}{n/a}}\\[-.2pt]3m\,3.7s\\[-.2pt]{\fontsize{5.8}{6.4}\selectfont \textcolor{gray}{n/a}}\end{tabular} & \hspace*{1pt}\begin{tabular}[t]{@{}c@{}}\rule{0pt}{12pt}3m\,28.3s\\[-.2pt]{\fontsize{5.8}{6.4}\selectfont \textcolor{BrickRed}{$ \blacktriangledown\,1.3\times $}}\\[-.2pt]3m\,22.8s\\[-.2pt]{\fontsize{5.8}{6.4}\selectfont \textcolor{BrickRed}{$ \blacktriangledown\,1.4\times $}}\end{tabular} & \hspace*{1pt}\begin{tabular}[t]{@{}c@{}}\rule{0pt}{12pt}3m\,53.2s\\[-.2pt]{\fontsize{5.8}{6.4}\selectfont \textcolor{BrickRed}{$ \blacktriangledown\,1.4\times $}}\\[-.2pt]3m\,43.5s\\[-.2pt]{\fontsize{5.8}{6.4}\selectfont \textcolor{BrickRed}{$ \blacktriangledown\,1.4\times $}}\end{tabular} & \hspace*{1pt}\begin{tabular}[t]{@{}c@{}}\rule{0pt}{12pt}7m\,14.6s\\[-.2pt]{\fontsize{5.8}{6.4}\selectfont \textcolor{BrickRed}{$ \blacktriangledown\,1.5\times $}}\\[-.2pt]6m\,57.2s\\[-.2pt]{\fontsize{5.8}{6.4}\selectfont \textcolor{BrickRed}{$ \blacktriangledown\,1.6\times $}}\end{tabular} & \hspace*{1pt}\begin{tabular}[t]{@{}c@{}}\rule{0pt}{12pt}51m\,22s\\[-.2pt]{\fontsize{5.8}{6.4}\selectfont \textcolor{gray}{n/a}}\\[-.2pt]51m\,17.2s\\[-.2pt]{\fontsize{5.8}{6.4}\selectfont \textcolor{gray}{n/a}}\end{tabular} & \hspace*{1pt}\begin{tabular}[t]{@{}c@{}}\rule{0pt}{12pt}1m\,10.3s\\[-.2pt]{\fontsize{5.8}{6.4}\selectfont \textcolor{ForestGreen}{$ \blacktriangle\,7.6\times $}}\\[-.2pt]1m\,9s\\[-.2pt]{\fontsize{5.8}{6.4}\selectfont \textcolor{ForestGreen}{$ \blacktriangle\,2.4\times $}}\end{tabular} & \hspace*{1pt}\begin{tabular}[t]{@{}c@{}}\rule{0pt}{12pt}1m\,8.6s\\[-.2pt]{\fontsize{5.8}{6.4}\selectfont \textcolor{ForestGreen}{$ \blacktriangle\,6.4\times $}}\\[-.2pt]1m\,6.9s\\[-.2pt]{\fontsize{5.8}{6.4}\selectfont \textcolor{ForestGreen}{$ \blacktriangle\,2.5\times $}}\end{tabular} & \hspace*{1pt}\begin{tabular}[t]{@{}c@{}}\rule{0pt}{12pt}2m\,21.5s\\[-.2pt]{\fontsize{5.8}{6.4}\selectfont \textcolor{ForestGreen}{$ \blacktriangle\,6.0\times $}}\\[-.2pt]2m\,18.4s\\[-.2pt]{\fontsize{5.8}{6.4}\selectfont \textcolor{ForestGreen}{$ \blacktriangle\,2.2\times $}}\end{tabular} & \hspace*{1pt}\begin{tabular}[t]{@{}c@{}}\rule{0pt}{12pt}15m\,56.8s\\[-.2pt]{\fontsize{5.8}{6.4}\selectfont \textcolor{gray}{n/a}}\\[-.2pt]15m\,51.9s\\[-.2pt]{\fontsize{5.8}{6.4}\selectfont \textcolor{gray}{n/a}}\end{tabular}\\[3pt]\hline
\textbf{BioseqDB} & \begin{tabular}[t]{@{}l@{}}\rule{0pt}{12pt}MPS / Scalar\\SIMD-DP / Block\end{tabular} & \hspace*{1pt}\begin{tabular}[t]{@{}c@{}}\rule{0pt}{12pt}57s\\[-.2pt]{\fontsize{5.8}{6.4}\selectfont \textcolor{ForestGreen}{$ \blacktriangle\,3.0\times $}}\\[-.2pt]14.3s\\[-.2pt]{\fontsize{5.8}{6.4}\selectfont \textcolor{ForestGreen}{$ \blacktriangle\,11.3\times $}}\end{tabular} & \hspace*{1pt}\begin{tabular}[t]{@{}c@{}}\rule{0pt}{12pt}50.6s\\[-.2pt]{\fontsize{5.8}{6.4}\selectfont \textcolor{ForestGreen}{$ \blacktriangle\,3.4\times $}}\\[-.2pt]15.1s\\[-.2pt]{\fontsize{5.8}{6.4}\selectfont \textcolor{ForestGreen}{$ \blacktriangle\,10.8\times $}}\end{tabular} & \textcolor{gray}{\ensuremath{\cdot}} & \textcolor{gray}{\ensuremath{\cdot}} & \textcolor{gray}{\ensuremath{\cdot}} & \textcolor{gray}{\ensuremath{\cdot}} & \textcolor{gray}{\ensuremath{\cdot}} & \textcolor{gray}{\ensuremath{\cdot}} & \hspace*{1pt}\begin{tabular}[t]{@{}c@{}}\rule{0pt}{12pt}5m\,10.1s\\[-.2pt]{\fontsize{5.8}{6.4}\selectfont \textcolor{BrickRed}{$ \blacktriangledown\,1.5\times $}}\\[-.2pt]4m\,27.6s\\[-.2pt]{\fontsize{5.8}{6.4}\selectfont \textcolor{BrickRed}{$ \blacktriangledown\,1.3\times $}}\end{tabular} & \hspace*{1pt}\begin{tabular}[t]{@{}c@{}}\rule{0pt}{12pt}5m\,3.5s\\[-.2pt]{\fontsize{5.8}{6.4}\selectfont \textcolor{BrickRed}{$ \blacktriangledown\,1.3\times $}}\\[-.2pt]4m\,28.9s\\[-.2pt]{\fontsize{5.8}{6.4}\selectfont \textcolor{BrickRed}{$ \blacktriangledown\,20\% $}}\end{tabular} & \textcolor{gray}{\ensuremath{\cdot}} & \textcolor{gray}{\ensuremath{\cdot}} & \textcolor{gray}{\ensuremath{\cdot}} & \textcolor{gray}{\ensuremath{\cdot}} & \textcolor{gray}{\ensuremath{\cdot}} & \textcolor{gray}{\ensuremath{\cdot}}\\[3pt]\hline
 & \begin{tabular}[t]{@{}l@{}}\rule{0pt}{12pt}MPS / Scalar\\SIMD-DP / KSW2\end{tabular} & \hspace*{1pt}\begin{tabular}[t]{@{}c@{}}\rule{0pt}{12pt}58.6s\\[-.2pt]{\fontsize{5.8}{6.4}\selectfont \textcolor{ForestGreen}{$ \blacktriangle\,2.9\times $}}\\[-.2pt]14.4s\\[-.2pt]{\fontsize{5.8}{6.4}\selectfont \textcolor{ForestGreen}{$ \blacktriangle\,11.2\times $}}\end{tabular} & \hspace*{1pt}\begin{tabular}[t]{@{}c@{}}\rule{0pt}{12pt}52.3s\\[-.2pt]{\fontsize{5.8}{6.4}\selectfont \textcolor{ForestGreen}{$ \blacktriangle\,3.3\times $}}\\[-.2pt]15.4s\\[-.2pt]{\fontsize{5.8}{6.4}\selectfont \textcolor{ForestGreen}{$ \blacktriangle\,10.6\times $}}\end{tabular} & \textcolor{gray}{\ensuremath{\cdot}} & \textcolor{gray}{\ensuremath{\cdot}} & \textcolor{gray}{\ensuremath{\cdot}} & \textcolor{gray}{\ensuremath{\cdot}} & \textcolor{gray}{\ensuremath{\cdot}} & \textcolor{gray}{\ensuremath{\cdot}} & \hspace*{1pt}\begin{tabular}[t]{@{}c@{}}\rule{0pt}{12pt}5m\,10.5s\\[-.2pt]{\fontsize{5.8}{6.4}\selectfont \textcolor{BrickRed}{$ \blacktriangledown\,1.5\times $}}\\[-.2pt]4m\,25.5s\\[-.2pt]{\fontsize{5.8}{6.4}\selectfont \textcolor{BrickRed}{$ \blacktriangledown\,1.3\times $}}\end{tabular} & \hspace*{1pt}\begin{tabular}[t]{@{}c@{}}\rule{0pt}{12pt}4m\,45.4s\\[-.2pt]{\fontsize{5.8}{6.4}\selectfont \textcolor{BrickRed}{$ \blacktriangledown\,22\% $}}\\[-.2pt]4m\,9.3s\\[-.2pt]{\fontsize{5.8}{6.4}\selectfont \textcolor{BrickRed}{$ \blacktriangledown\,12\% $}}\end{tabular} & \textcolor{gray}{\ensuremath{\cdot}} & \textcolor{gray}{\ensuremath{\cdot}} & \textcolor{gray}{\ensuremath{\cdot}} & \textcolor{gray}{\ensuremath{\cdot}} & \textcolor{gray}{\ensuremath{\cdot}} & \textcolor{gray}{\ensuremath{\cdot}}\\[3pt]\hline
 & \begin{tabular}[t]{@{}l@{}}\rule{0pt}{12pt}MPS / SIMD-ref\\SIMD-DP / Block\end{tabular} & \textcolor{gray}{\ensuremath{\cdot}} & \textcolor{gray}{\ensuremath{\cdot}} & \textcolor{gray}{\ensuremath{\cdot}} & \textcolor{gray}{\ensuremath{\cdot}} & \textcolor{gray}{\ensuremath{\cdot}} & \textcolor{gray}{\ensuremath{\cdot}} & \textcolor{gray}{\ensuremath{\cdot}} & \textcolor{gray}{\ensuremath{\cdot}} & \textcolor{gray}{\ensuremath{\cdot}} & \textcolor{gray}{\ensuremath{\cdot}} & \hspace*{1pt}\begin{tabular}[t]{@{}c@{}}\rule{0pt}{12pt}6m\,45.8s\\[-.2pt]{\fontsize{5.8}{6.4}\selectfont \textcolor{ForestGreen}{$ \blacktriangle\,7\% $}}\\[-.2pt]5m\,46.6s\\[-.2pt]{\fontsize{5.8}{6.4}\selectfont \textcolor{ForestGreen}{$ \blacktriangle\,17\% $}}\end{tabular} & \textcolor{gray}{\ensuremath{\cdot}} & \textcolor{gray}{\ensuremath{\cdot}} & \textcolor{gray}{\ensuremath{\cdot}} & \textcolor{gray}{\ensuremath{\cdot}} & \textcolor{gray}{\ensuremath{\cdot}}\\[3pt]\hline
 & \begin{tabular}[t]{@{}l@{}}\rule{0pt}{12pt}MPS / Two-stack\\RMQ / Block\end{tabular} & \hspace*{1pt}\begin{tabular}[t]{@{}c@{}}\rule{0pt}{12pt}50.7s\\[-.2pt]{\fontsize{5.8}{6.4}\selectfont \textcolor{ForestGreen}{$ \blacktriangle\,3.3\times $}}\\[-.2pt]9.4s\\[-.2pt]{\fontsize{5.8}{6.4}\selectfont \textcolor{ForestGreen}{$ \blacktriangle\,17.2\times $}}\end{tabular} & \hspace*{1pt}\begin{tabular}[t]{@{}c@{}}\rule{0pt}{12pt}44.9s\\[-.2pt]{\fontsize{5.8}{6.4}\selectfont \textcolor{ForestGreen}{$ \blacktriangle\,3.8\times $}}\\[-.2pt]9.9s\\[-.2pt]{\fontsize{5.8}{6.4}\selectfont \textcolor{ForestGreen}{$ \blacktriangle\,16.5\times $}}\end{tabular} & \textcolor{gray}{\ensuremath{\cdot}} & \textcolor{gray}{\ensuremath{\cdot}} & \textcolor{gray}{\ensuremath{\cdot}} & \textcolor{gray}{\ensuremath{\cdot}} & \textcolor{gray}{\ensuremath{\cdot}} & \textcolor{gray}{\ensuremath{\cdot}} & \textcolor{gray}{\ensuremath{\cdot}} & \textcolor{gray}{\ensuremath{\cdot}} & \textcolor{gray}{\ensuremath{\cdot}} & \textcolor{gray}{\ensuremath{\cdot}} & \textcolor{gray}{\ensuremath{\cdot}} & \textcolor{gray}{\ensuremath{\cdot}} & \textcolor{gray}{\ensuremath{\cdot}} & \textcolor{gray}{\ensuremath{\cdot}}\\[3pt]\hline
 & \begin{tabular}[t]{@{}l@{}}\rule{0pt}{12pt}MPS / Two-stack\\SIMD-DP / Block\end{tabular} & \hspace*{1pt}\begin{tabular}[t]{@{}c@{}}\rule{0pt}{12pt}51.6s\\[-.2pt]{\fontsize{5.8}{6.4}\selectfont \textcolor{ForestGreen}{$ \blacktriangle\,3.3\times $}}\\[-.2pt]9.8s\\[-.2pt]{\fontsize{5.8}{6.4}\selectfont \textcolor{ForestGreen}{$ \blacktriangle\,16.6\times $}}\end{tabular} & \hspace*{1pt}\begin{tabular}[t]{@{}c@{}}\rule{0pt}{12pt}44.7s\\[-.2pt]{\fontsize{5.8}{6.4}\selectfont \textcolor{ForestGreen}{$ \blacktriangle\,3.8\times $}}\\[-.2pt]9.6s\\[-.2pt]{\fontsize{5.8}{6.4}\selectfont \textcolor{ForestGreen}{$ \blacktriangle\,17.0\times $}}\end{tabular} & \hspace*{1pt}\begin{tabular}[t]{@{}c@{}}\rule{0pt}{12pt}1m\,16.1s\\[-.2pt]{\fontsize{5.8}{6.4}\selectfont \textcolor{ForestGreen}{$ \blacktriangle\,3.9\times $}}\\[-.2pt]15.7s\\[-.2pt]{\fontsize{5.8}{6.4}\selectfont \textcolor{ForestGreen}{$ \blacktriangle\,17.7\times $}}\end{tabular} & \textcolor{gray}{\ensuremath{\cdot}} & \textcolor{gray}{\ensuremath{\cdot}} & \textcolor{gray}{\ensuremath{\cdot}} & \textcolor{gray}{\ensuremath{\cdot}} & \textcolor{gray}{\ensuremath{\cdot}} & \textcolor{gray}{\ensuremath{\cdot}} & \textcolor{gray}{\ensuremath{\cdot}} & \textcolor{gray}{\ensuremath{\cdot}} & \textcolor{gray}{\ensuremath{\cdot}} & \textcolor{gray}{\ensuremath{\cdot}} & \textcolor{gray}{\ensuremath{\cdot}} & \textcolor{gray}{\ensuremath{\cdot}} & \textcolor{gray}{\ensuremath{\cdot}}\\[3pt]\hline
 & \begin{tabular}[t]{@{}l@{}}\rule{0pt}{12pt}MPS / Two-stack\\SIMD-DP / KSW2\end{tabular} & \hspace*{1pt}\begin{tabular}[t]{@{}c@{}}\rule{0pt}{12pt}55.3s\\[-.2pt]{\fontsize{5.8}{6.4}\selectfont \textcolor{ForestGreen}{$ \blacktriangle\,3.0\times $}}\\[-.2pt]9.6s\\[-.2pt]{\fontsize{5.8}{6.4}\selectfont \textcolor{ForestGreen}{$ \blacktriangle\,16.8\times $}}\end{tabular} & \hspace*{1pt}\begin{tabular}[t]{@{}c@{}}\rule{0pt}{12pt}48.6s\\[-.2pt]{\fontsize{5.8}{6.4}\selectfont \textcolor{ForestGreen}{$ \blacktriangle\,3.5\times $}}\\[-.2pt]9.6s\\[-.2pt]{\fontsize{5.8}{6.4}\selectfont \textcolor{ForestGreen}{$ \blacktriangle\,17.0\times $}}\end{tabular} & \textcolor{gray}{\ensuremath{\cdot}} & \textcolor{gray}{\ensuremath{\cdot}} & \textcolor{gray}{\ensuremath{\cdot}} & \textcolor{gray}{\ensuremath{\cdot}} & \textcolor{gray}{\ensuremath{\cdot}} & \textcolor{gray}{\ensuremath{\cdot}} & \textcolor{gray}{\ensuremath{\cdot}} & \textcolor{gray}{\ensuremath{\cdot}} & \textcolor{gray}{\ensuremath{\cdot}} & \textcolor{gray}{\ensuremath{\cdot}} & \textcolor{gray}{\ensuremath{\cdot}} & \textcolor{gray}{\ensuremath{\cdot}} & \textcolor{gray}{\ensuremath{\cdot}} & \textcolor{gray}{\ensuremath{\cdot}}\\[3pt]\hline
\rule{0pt}{12pt}\textbf{BioseqDB} & Vulkan (all modes) & \textcolor{gray}{---\textsuperscript{g}} & \textcolor{gray}{---\textsuperscript{g}} & \textcolor{gray}{---\textsuperscript{g}} & \textcolor{gray}{---\textsuperscript{g}} & \textcolor{gray}{---\textsuperscript{g}} & \textcolor{gray}{---\textsuperscript{g}} & \textcolor{gray}{---\textsuperscript{g}} & \textcolor{gray}{---\textsuperscript{g}} & \textcolor{gray}{---\textsuperscript{g}} & \textcolor{gray}{---\textsuperscript{g}} & \textcolor{gray}{---\textsuperscript{g}} & \textcolor{gray}{---\textsuperscript{g}} & \textcolor{gray}{---\textsuperscript{g}} & \textcolor{gray}{---\textsuperscript{g}} & \textcolor{gray}{---\textsuperscript{g}} & \textcolor{gray}{---\textsuperscript{g}}\\[3pt]\hline
\end{tabular}\par}
\smallskip{\footnotesize\noindent Each cell: A, comparison, B, comparison. Times cover 1,024 queries. $\cdot$ = no approved observation. ---\textsuperscript{g} = GPU unavailable, not executed. n/a = no measured comparison.\par
\noindent \textsuperscript{f}SIMD-ref request observed as scalar, compared only as that effective path. \textsuperscript{p}Accepted by explicit preservation after wrapper postflight failure, no run receipt claimed.\par
\noindent MPF/MPS = minimizers per fragment/sequence. Row order: index representation / minimizer, then chaining / base alignment. SIMD-ref is the reference minimizer, Bounded is bounded ntHash.\par
\noindent Blue = one lowest A and one lowest B down each machine/workload column, across continuation pages. Exact ties select the first row. Repeated baselines are marked only on page 1.\par
\noindent Light gray = minimap2 baseline. Pale yellow = preferred MPF / Bounded / SIMD-DP / Block configuration, not a claim of winning every column.\par
\noindent When all Vulkan modes are unmeasured, they share one row. Baseline comparisons use the fastest measured BioseqDB effective path separately for A and B.\par}
\clearpage
\noindent{\large\bfseries AWS: 2 Gbp, k=15, w=10 (mapping interval)\hfill 1/7}\par\smallskip
{\fontsize{7.1}{8.1}\selectfont\setlength{\tabcolsep}{1pt}\setlength{\fboxsep}{2pt}\renewcommand{\arraystretch}{1.12}
\noindent\begin{tabular}{@{}p{14mm}p{27mm}|p{13.55mm}|p{13.55mm}|p{13.55mm}|p{13.55mm}|p{13.55mm}|p{13.55mm}|p{13.55mm}|p{13.55mm}|p{13.55mm}|p{13.55mm}|p{13.55mm}|p{13.55mm}|p{13.55mm}|p{13.55mm}|p{13.55mm}|p{13.55mm}|@{}}\hline
Engine & Execution & \multicolumn{8}{c|}{\textbf{Index}} & \multicolumn{8}{c|}{\textbf{No index}}\\
 &  & \multicolumn{8}{c|}{\textbf{2 Gbp}} & \multicolumn{8}{c|}{\textbf{2 Gbp}}\\
 &  & \multicolumn{8}{c|}{\textbf{k=15, w=10}} & \multicolumn{8}{c|}{\textbf{k=15, w=10}}\\
 &  & \multicolumn{4}{c|}{\textbf{DNA4}} & \multicolumn{4}{c|}{\textbf{DNA5}} & \multicolumn{4}{c|}{\textbf{DNA4}} & \multicolumn{4}{c|}{\textbf{DNA5}}\\
 &  & \multicolumn{1}{c|}{\textbf{M7A}} & \multicolumn{1}{c|}{\textbf{M7I}} & \multicolumn{1}{c|}{\textbf{M4}} & \multicolumn{1}{c|}{\textbf{C7G}} & \multicolumn{1}{c|}{\textbf{M7A}} & \multicolumn{1}{c|}{\textbf{M7I}} & \multicolumn{1}{c|}{\textbf{M4}} & \multicolumn{1}{c|}{\textbf{C7G}} & \multicolumn{1}{c|}{\textbf{M7A}} & \multicolumn{1}{c|}{\textbf{M7I}} & \multicolumn{1}{c|}{\textbf{M4}} & \multicolumn{1}{c|}{\textbf{C7G}} & \multicolumn{1}{c|}{\textbf{M7A}} & \multicolumn{1}{c|}{\textbf{M7I}} & \multicolumn{1}{c|}{\textbf{M4}} & \multicolumn{1}{c|}{\textbf{C7G}}\\
\hline
\rowcolor{BenchBaseline}[0.5pt][0.5pt]\textbf{minimap2} & \textbf{Baseline} & \hspace*{1pt}\begin{tabular}[t]{@{}c@{}}\rule{0pt}{12pt}11.2s\\[-.2pt]{\fontsize{5.8}{6.4}\selectfont \textcolor{BrickRed}{$ \blacktriangledown\,3.1\times $}}\\[-.2pt]10s\\[-.2pt]{\fontsize{5.8}{6.4}\selectfont \textcolor{BrickRed}{$ \blacktriangledown\,3.6\times $}}\end{tabular} & \hspace*{1pt}\begin{tabular}[t]{@{}c@{}}\rule{0pt}{12pt}10.7s\\[-.2pt]{\fontsize{5.8}{6.4}\selectfont \textcolor{BrickRed}{$ \blacktriangledown\,3.2\times $}}\\[-.2pt]8.6s\\[-.2pt]{\fontsize{5.8}{6.4}\selectfont \textcolor{BrickRed}{$ \blacktriangledown\,3.5\times $}}\end{tabular} & \hspace*{1pt}\begin{tabular}[t]{@{}c@{}}\rule{0pt}{12pt}22.5s\\[-.2pt]{\fontsize{5.8}{6.4}\selectfont \textcolor{BrickRed}{$ \blacktriangledown\,3.6\times $}}\\[-.2pt]18.8s\\[-.2pt]{\fontsize{5.8}{6.4}\selectfont \textcolor{BrickRed}{$ \blacktriangledown\,4.4\times $}}\end{tabular} & \hspace*{1pt}\begin{tabular}[t]{@{}c@{}}\rule{0pt}{12pt}1m\,13.1s\\[-.2pt]{\fontsize{5.8}{6.4}\selectfont \textcolor{BrickRed}{$ \blacktriangledown\,22.5\times $}}\\[-.2pt]1m\,5.7s\\[-.2pt]{\fontsize{5.8}{6.4}\selectfont \textcolor{BrickRed}{$ \blacktriangledown\,27.7\times $}}\end{tabular} & \hspace*{1pt}\begin{tabular}[t]{@{}c@{}}\rule{0pt}{12pt}9.9s\\[-.2pt]{\fontsize{5.8}{6.4}\selectfont \textcolor{BrickRed}{$ \blacktriangledown\,21\% $}}\\[-.2pt]8.8s\\[-.2pt]{\fontsize{5.8}{6.4}\selectfont \textcolor{BrickRed}{$ \blacktriangledown\,1.5\times $}}\end{tabular} & \hspace*{1pt}\begin{tabular}[t]{@{}c@{}}\rule{0pt}{12pt}9.5s\\[-.2pt]{\fontsize{5.8}{6.4}\selectfont \textcolor{BrickRed}{$ \blacktriangledown\,1.4\times $}}\\[-.2pt]7.6s\\[-.2pt]{\fontsize{5.8}{6.4}\selectfont \textcolor{BrickRed}{$ \blacktriangledown\,1.6\times $}}\end{tabular} & \hspace*{1pt}\begin{tabular}[t]{@{}c@{}}\rule{0pt}{12pt}16.5s\\[-.2pt]{\fontsize{5.8}{6.4}\selectfont \textcolor{BrickRed}{$ \blacktriangledown\,1.7\times $}}\\[-.2pt]12.9s\\[-.2pt]{\fontsize{5.8}{6.4}\selectfont \textcolor{BrickRed}{$ \blacktriangledown\,2.0\times $}}\end{tabular} & \hspace*{1pt}\begin{tabular}[t]{@{}c@{}}\rule{0pt}{12pt}28.2s\\[-.2pt]{\fontsize{5.8}{6.4}\selectfont \textcolor{BrickRed}{$ \blacktriangledown\,2.1\times $}}\\[-.2pt]26.2s\\[-.2pt]{\fontsize{5.8}{6.4}\selectfont \textcolor{BrickRed}{$ \blacktriangledown\,2.4\times $}}\end{tabular} & \hspace*{1pt}\begin{tabular}[t]{@{}c@{}}\rule{0pt}{12pt}\colorbox{BenchBest}{\strut 53s}\\[-.2pt]{\fontsize{5.8}{6.4}\selectfont \textcolor{ForestGreen}{$ \blacktriangle\,2.2\times $}}\\[-.2pt]\colorbox{BenchBest}{\strut 51.9s}\\[-.2pt]{\fontsize{5.8}{6.4}\selectfont \textcolor{ForestGreen}{$ \blacktriangle\,2.1\times $}}\end{tabular} & \hspace*{1pt}\begin{tabular}[t]{@{}c@{}}\rule{0pt}{12pt}\colorbox{BenchBest}{\strut 56.6s}\\[-.2pt]{\fontsize{5.8}{6.4}\selectfont \textcolor{ForestGreen}{$ \blacktriangle\,2.1\times $}}\\[-.2pt]\colorbox{BenchBest}{\strut 54.6s}\\[-.2pt]{\fontsize{5.8}{6.4}\selectfont \textcolor{ForestGreen}{$ \blacktriangle\,1.9\times $}}\end{tabular} & \hspace*{1pt}\begin{tabular}[t]{@{}c@{}}\rule{0pt}{12pt}\colorbox{BenchBest}{\strut 1m\,55.5s}\\[-.2pt]{\fontsize{5.8}{6.4}\selectfont \textcolor{ForestGreen}{$ \blacktriangle\,1.8\times $}}\\[-.2pt]\colorbox{BenchBest}{\strut 1m\,51.8s}\\[-.2pt]{\fontsize{5.8}{6.4}\selectfont \textcolor{ForestGreen}{$ \blacktriangle\,1.5\times $}}\end{tabular} & \hspace*{1pt}\begin{tabular}[t]{@{}c@{}}\rule{0pt}{12pt}\colorbox{BenchBest}{\strut 6m\,1.3s}\\[-.2pt]{\fontsize{5.8}{6.4}\selectfont \textcolor{gray}{n/a}}\\[-.2pt]\colorbox{BenchBest}{\strut 5m\,55.2s}\\[-.2pt]{\fontsize{5.8}{6.4}\selectfont \textcolor{gray}{n/a}}\end{tabular} & \hspace*{1pt}\begin{tabular}[t]{@{}c@{}}\rule{0pt}{12pt}\colorbox{BenchBest}{\strut 48.1s}\\[-.2pt]{\fontsize{5.8}{6.4}\selectfont \textcolor{ForestGreen}{$ \blacktriangle\,7.3\times $}}\\[-.2pt]\colorbox{BenchBest}{\strut 47s}\\[-.2pt]{\fontsize{5.8}{6.4}\selectfont \textcolor{ForestGreen}{$ \blacktriangle\,2.4\times $}}\end{tabular} & \hspace*{1pt}\begin{tabular}[t]{@{}c@{}}\rule{0pt}{12pt}\colorbox{BenchBest}{\strut 44.5s}\\[-.2pt]{\fontsize{5.8}{6.4}\selectfont \textcolor{ForestGreen}{$ \blacktriangle\,6.5\times $}}\\[-.2pt]\colorbox{BenchBest}{\strut 42.6s}\\[-.2pt]{\fontsize{5.8}{6.4}\selectfont \textcolor{ForestGreen}{$ \blacktriangle\,2.7\times $}}\end{tabular} & \hspace*{1pt}\begin{tabular}[t]{@{}c@{}}\rule{0pt}{12pt}\colorbox{BenchBest}{\strut 1m\,30.2s}\\[-.2pt]{\fontsize{5.8}{6.4}\selectfont \textcolor{ForestGreen}{$ \blacktriangle\,6.4\times $}}\\[-.2pt]\colorbox{BenchBest}{\strut 1m\,26.7s}\\[-.2pt]{\fontsize{5.8}{6.4}\selectfont \textcolor{ForestGreen}{$ \blacktriangle\,2.2\times $}}\end{tabular} & \hspace*{1pt}\begin{tabular}[t]{@{}c@{}}\rule{0pt}{12pt}\colorbox{BenchBest}{\strut 2m\,46.4s}\\[-.2pt]{\fontsize{5.8}{6.4}\selectfont \textcolor{gray}{n/a}}\\[-.2pt]\colorbox{BenchBest}{\strut 2m\,42.7s}\\[-.2pt]{\fontsize{5.8}{6.4}\selectfont \textcolor{gray}{n/a}}\end{tabular}\\[3pt]\hline
\textbf{BioseqDB} & \begin{tabular}[t]{@{}l@{}}\rule{0pt}{12pt}MPF / Bounded\\RMQ / Block\end{tabular} & \hspace*{1pt}\begin{tabular}[t]{@{}c@{}}\rule{0pt}{12pt}3.7s\\[-.2pt]{\fontsize{5.8}{6.4}\selectfont \textcolor{ForestGreen}{$ \blacktriangle\,3.0\times $}}\\[-.2pt]2.9s\\[-.2pt]{\fontsize{5.8}{6.4}\selectfont \textcolor{ForestGreen}{$ \blacktriangle\,3.4\times $}}\end{tabular} & \hspace*{1pt}\begin{tabular}[t]{@{}c@{}}\rule{0pt}{12pt}3.5s\\[-.2pt]{\fontsize{5.8}{6.4}\selectfont \textcolor{ForestGreen}{$ \blacktriangle\,3.0\times $}}\\[-.2pt]2.7s\\[-.2pt]{\fontsize{5.8}{6.4}\selectfont \textcolor{ForestGreen}{$ \blacktriangle\,3.2\times $}}\end{tabular} & \hspace*{1pt}\begin{tabular}[t]{@{}c@{}}\rule{0pt}{12pt}6.5s\\[-.2pt]{\fontsize{5.8}{6.4}\selectfont \textcolor{ForestGreen}{$ \blacktriangle\,3.5\times $}}\\[-.2pt]4.5s\\[-.2pt]{\fontsize{5.8}{6.4}\selectfont \textcolor{ForestGreen}{$ \blacktriangle\,4.2\times $}}\end{tabular} & \hspace*{1pt}\begin{tabular}[t]{@{}c@{}}\rule{0pt}{12pt}3.4s\\[-.2pt]{\fontsize{5.8}{6.4}\selectfont \textcolor{ForestGreen}{$ \blacktriangle\,21.2\times $}}\\[-.2pt]2.5s\\[-.2pt]{\fontsize{5.8}{6.4}\selectfont \textcolor{ForestGreen}{$ \blacktriangle\,26.1\times $}}\end{tabular} & \textcolor{gray}{\ensuremath{\cdot}} & \textcolor{gray}{\ensuremath{\cdot}} & \textcolor{gray}{\ensuremath{\cdot}} & \textcolor{gray}{\ensuremath{\cdot}} & \hspace*{1pt}\begin{tabular}[t]{@{}c@{}}\rule{0pt}{12pt}1m\,55.5s\\[-.2pt]{\fontsize{5.8}{6.4}\selectfont \textcolor{BrickRed}{$ \blacktriangledown\,2.2\times $}}\\[-.2pt]1m\,47.1s\\[-.2pt]{\fontsize{5.8}{6.4}\selectfont \textcolor{BrickRed}{$ \blacktriangledown\,2.1\times $}}\end{tabular} & \hspace*{1pt}\begin{tabular}[t]{@{}c@{}}\rule{0pt}{12pt}2m\,21s\\[-.2pt]{\fontsize{5.8}{6.4}\selectfont \textcolor{BrickRed}{$ \blacktriangledown\,2.5\times $}}\\[-.2pt]2m\,10.3s\\[-.2pt]{\fontsize{5.8}{6.4}\selectfont \textcolor{BrickRed}{$ \blacktriangledown\,2.4\times $}}\end{tabular} & \textcolor{gray}{\ensuremath{\cdot}} & \textcolor{gray}{\ensuremath{\cdot}} & \textcolor{gray}{\ensuremath{\cdot}} & \textcolor{gray}{\ensuremath{\cdot}} & \textcolor{gray}{\ensuremath{\cdot}} & \textcolor{gray}{\ensuremath{\cdot}}\\[3pt]\hline
 & \begin{tabular}[t]{@{}l@{}}\rule{0pt}{12pt}MPF / Bounded\\RMQ / KSW2\end{tabular} & \hspace*{1pt}\begin{tabular}[t]{@{}c@{}}\rule{0pt}{12pt}5.3s\\[-.2pt]{\fontsize{5.8}{6.4}\selectfont \textcolor{ForestGreen}{$ \blacktriangle\,2.1\times $}}\\[-.2pt]2.8s\\[-.2pt]{\fontsize{5.8}{6.4}\selectfont \textcolor{ForestGreen}{$ \blacktriangle\,3.6\times $}}\end{tabular} & \hspace*{1pt}\begin{tabular}[t]{@{}c@{}}\rule{0pt}{12pt}4.8s\\[-.2pt]{\fontsize{5.8}{6.4}\selectfont \textcolor{ForestGreen}{$ \blacktriangle\,2.2\times $}}\\[-.2pt]2.5s\\[-.2pt]{\fontsize{5.8}{6.4}\selectfont \textcolor{ForestGreen}{$ \blacktriangle\,3.5\times $}}\end{tabular} & \hspace*{1pt}\begin{tabular}[t]{@{}c@{}}\rule{0pt}{12pt}9.8s\\[-.2pt]{\fontsize{5.8}{6.4}\selectfont \textcolor{ForestGreen}{$ \blacktriangle\,2.3\times $}}\\[-.2pt]\colorbox{BenchBest}{\strut 4.2s}\\[-.2pt]{\fontsize{5.8}{6.4}\selectfont \textcolor{ForestGreen}{$ \blacktriangle\,4.4\times $}}\end{tabular} & \hspace*{1pt}\begin{tabular}[t]{@{}c@{}}\rule{0pt}{12pt}7.7s\\[-.2pt]{\fontsize{5.8}{6.4}\selectfont \textcolor{ForestGreen}{$ \blacktriangle\,9.5\times $}}\\[-.2pt]4.6s\\[-.2pt]{\fontsize{5.8}{6.4}\selectfont \textcolor{ForestGreen}{$ \blacktriangle\,14.2\times $}}\end{tabular} & \textcolor{gray}{\ensuremath{\cdot}} & \textcolor{gray}{\ensuremath{\cdot}} & \textcolor{gray}{\ensuremath{\cdot}} & \textcolor{gray}{\ensuremath{\cdot}} & \hspace*{1pt}\begin{tabular}[t]{@{}c@{}}\rule{0pt}{12pt}2m\,17.7s\\[-.2pt]{\fontsize{5.8}{6.4}\selectfont \textcolor{BrickRed}{$ \blacktriangledown\,2.6\times $}}\\[-.2pt]2m\,7s\\[-.2pt]{\fontsize{5.8}{6.4}\selectfont \textcolor{BrickRed}{$ \blacktriangledown\,2.4\times $}}\end{tabular} & \hspace*{1pt}\begin{tabular}[t]{@{}c@{}}\rule{0pt}{12pt}1m\,59.4s\\[-.2pt]{\fontsize{5.8}{6.4}\selectfont \textcolor{BrickRed}{$ \blacktriangledown\,2.1\times $}}\\[-.2pt]1m\,47.5s\\[-.2pt]{\fontsize{5.8}{6.4}\selectfont \textcolor{BrickRed}{$ \blacktriangledown\,2.0\times $}}\end{tabular} & \textcolor{gray}{\ensuremath{\cdot}} & \textcolor{gray}{\ensuremath{\cdot}} & \textcolor{gray}{\ensuremath{\cdot}} & \textcolor{gray}{\ensuremath{\cdot}} & \textcolor{gray}{\ensuremath{\cdot}} & \textcolor{gray}{\ensuremath{\cdot}}\\[3pt]\hline
 & \begin{tabular}[t]{@{}l@{}}\rule{0pt}{12pt}MPF / Bounded\\Scalar-DP / Block\end{tabular} & \hspace*{1pt}\begin{tabular}[t]{@{}c@{}}\rule{0pt}{12pt}3.6s\\[-.2pt]{\fontsize{5.8}{6.4}\selectfont \textcolor{ForestGreen}{$ \blacktriangle\,3.1\times $}}\\[-.2pt]2.8s\\[-.2pt]{\fontsize{5.8}{6.4}\selectfont \textcolor{ForestGreen}{$ \blacktriangle\,3.5\times $}}\end{tabular} & \hspace*{1pt}\begin{tabular}[t]{@{}c@{}}\rule{0pt}{12pt}3.4s\\[-.2pt]{\fontsize{5.8}{6.4}\selectfont \textcolor{ForestGreen}{$ \blacktriangle\,3.1\times $}}\\[-.2pt]2.6s\\[-.2pt]{\fontsize{5.8}{6.4}\selectfont \textcolor{ForestGreen}{$ \blacktriangle\,3.3\times $}}\end{tabular} & \hspace*{1pt}\begin{tabular}[t]{@{}c@{}}\rule{0pt}{12pt}6.4s\\[-.2pt]{\fontsize{5.8}{6.4}\selectfont \textcolor{ForestGreen}{$ \blacktriangle\,3.5\times $}}\\[-.2pt]4.4s\\[-.2pt]{\fontsize{5.8}{6.4}\selectfont \textcolor{ForestGreen}{$ \blacktriangle\,4.2\times $}}\end{tabular} & \hspace*{1pt}\begin{tabular}[t]{@{}c@{}}\rule{0pt}{12pt}\colorbox{BenchBest}{\strut 3.2s}\\[-.2pt]{\fontsize{5.8}{6.4}\selectfont \textcolor{ForestGreen}{$ \blacktriangle\,22.5\times $}}\\[-.2pt]\colorbox{BenchBest}{\strut 2.4s}\\[-.2pt]{\fontsize{5.8}{6.4}\selectfont \textcolor{ForestGreen}{$ \blacktriangle\,27.7\times $}}\end{tabular} & \textcolor{gray}{\ensuremath{\cdot}} & \textcolor{gray}{\ensuremath{\cdot}} & \textcolor{gray}{\ensuremath{\cdot}} & \textcolor{gray}{\ensuremath{\cdot}} & \hspace*{1pt}\begin{tabular}[t]{@{}c@{}}\rule{0pt}{12pt}1m\,56.2s\\[-.2pt]{\fontsize{5.8}{6.4}\selectfont \textcolor{BrickRed}{$ \blacktriangledown\,2.2\times $}}\\[-.2pt]1m\,47.8s\\[-.2pt]{\fontsize{5.8}{6.4}\selectfont \textcolor{BrickRed}{$ \blacktriangledown\,2.1\times $}}\end{tabular} & \hspace*{1pt}\begin{tabular}[t]{@{}c@{}}\rule{0pt}{12pt}2m\,7.4s\\[-.2pt]{\fontsize{5.8}{6.4}\selectfont \textcolor{BrickRed}{$ \blacktriangledown\,2.3\times $}}\\[-.2pt]1m\,56.8s\\[-.2pt]{\fontsize{5.8}{6.4}\selectfont \textcolor{BrickRed}{$ \blacktriangledown\,2.1\times $}}\end{tabular} & \textcolor{gray}{\ensuremath{\cdot}} & \textcolor{gray}{\ensuremath{\cdot}} & \textcolor{gray}{\ensuremath{\cdot}} & \textcolor{gray}{\ensuremath{\cdot}} & \textcolor{gray}{\ensuremath{\cdot}} & \textcolor{gray}{\ensuremath{\cdot}}\\[3pt]\hline
 & \begin{tabular}[t]{@{}l@{}}\rule{0pt}{12pt}MPF / Bounded\\Scalar-DP / KSW2\end{tabular} & \hspace*{1pt}\begin{tabular}[t]{@{}c@{}}\rule{0pt}{12pt}5.2s\\[-.2pt]{\fontsize{5.8}{6.4}\selectfont \textcolor{ForestGreen}{$ \blacktriangle\,2.1\times $}}\\[-.2pt]2.8s\\[-.2pt]{\fontsize{5.8}{6.4}\selectfont \textcolor{ForestGreen}{$ \blacktriangle\,3.6\times $}}\end{tabular} & \hspace*{1pt}\begin{tabular}[t]{@{}c@{}}\rule{0pt}{12pt}4.7s\\[-.2pt]{\fontsize{5.8}{6.4}\selectfont \textcolor{ForestGreen}{$ \blacktriangle\,2.3\times $}}\\[-.2pt]\colorbox{BenchBest}{\strut 2.5s}\\[-.2pt]{\fontsize{5.8}{6.4}\selectfont \textcolor{ForestGreen}{$ \blacktriangle\,3.5\times $}}\end{tabular} & \hspace*{1pt}\begin{tabular}[t]{@{}c@{}}\rule{0pt}{12pt}9.7s\\[-.2pt]{\fontsize{5.8}{6.4}\selectfont \textcolor{ForestGreen}{$ \blacktriangle\,2.3\times $}}\\[-.2pt]4.3s\\[-.2pt]{\fontsize{5.8}{6.4}\selectfont \textcolor{ForestGreen}{$ \blacktriangle\,4.4\times $}}\end{tabular} & \hspace*{1pt}\begin{tabular}[t]{@{}c@{}}\rule{0pt}{12pt}7.2s\\[-.2pt]{\fontsize{5.8}{6.4}\selectfont \textcolor{ForestGreen}{$ \blacktriangle\,10.1\times $}}\\[-.2pt]4.3s\\[-.2pt]{\fontsize{5.8}{6.4}\selectfont \textcolor{ForestGreen}{$ \blacktriangle\,15.4\times $}}\end{tabular} & \textcolor{gray}{\ensuremath{\cdot}} & \textcolor{gray}{\ensuremath{\cdot}} & \textcolor{gray}{\ensuremath{\cdot}} & \textcolor{gray}{\ensuremath{\cdot}} & \hspace*{1pt}\begin{tabular}[t]{@{}c@{}}\rule{0pt}{12pt}2m\,8.2s\\[-.2pt]{\fontsize{5.8}{6.4}\selectfont \textcolor{BrickRed}{$ \blacktriangledown\,2.4\times $}}\\[-.2pt]1m\,58.3s\\[-.2pt]{\fontsize{5.8}{6.4}\selectfont \textcolor{BrickRed}{$ \blacktriangledown\,2.3\times $}}\end{tabular} & \hspace*{1pt}\begin{tabular}[t]{@{}c@{}}\rule{0pt}{12pt}2m\,3.7s\\[-.2pt]{\fontsize{5.8}{6.4}\selectfont \textcolor{BrickRed}{$ \blacktriangledown\,2.2\times $}}\\[-.2pt]1m\,51.2s\\[-.2pt]{\fontsize{5.8}{6.4}\selectfont \textcolor{BrickRed}{$ \blacktriangledown\,2.0\times $}}\end{tabular} & \textcolor{gray}{\ensuremath{\cdot}} & \textcolor{gray}{\ensuremath{\cdot}} & \textcolor{gray}{\ensuremath{\cdot}} & \textcolor{gray}{\ensuremath{\cdot}} & \hspace*{1pt}\begin{tabular}[t]{@{}c@{}}\rule{0pt}{12pt}9m\,39.3s\\[-.2pt]{\fontsize{5.8}{6.4}\selectfont \textcolor{BrickRed}{$ \blacktriangledown\,6.4\times $}}\\[-.2pt]3m\,13.2s\\[-.2pt]{\fontsize{5.8}{6.4}\selectfont \textcolor{BrickRed}{$ \blacktriangledown\,2.2\times $}}\end{tabular} & \textcolor{gray}{\ensuremath{\cdot}}\\[3pt]\hline
\rowcolor{BenchPreferred}[0.5pt][0.5pt] & \begin{tabular}[t]{@{}l@{}}\rule{0pt}{12pt}MPF / Bounded\\SIMD-DP / Block\end{tabular} & \hspace*{1pt}\begin{tabular}[t]{@{}c@{}}\rule{0pt}{12pt}3.7s\\[-.2pt]{\fontsize{5.8}{6.4}\selectfont \textcolor{ForestGreen}{$ \blacktriangle\,3.0\times $}}\\[-.2pt]3s\\[-.2pt]{\fontsize{5.8}{6.4}\selectfont \textcolor{ForestGreen}{$ \blacktriangle\,3.4\times $}}\end{tabular} & \hspace*{1pt}\begin{tabular}[t]{@{}c@{}}\rule{0pt}{12pt}3.5s\\[-.2pt]{\fontsize{5.8}{6.4}\selectfont \textcolor{ForestGreen}{$ \blacktriangle\,3.0\times $}}\\[-.2pt]2.7s\\[-.2pt]{\fontsize{5.8}{6.4}\selectfont \textcolor{ForestGreen}{$ \blacktriangle\,3.2\times $}}\end{tabular} & \hspace*{1pt}\begin{tabular}[t]{@{}c@{}}\rule{0pt}{12pt}6.7s\\[-.2pt]{\fontsize{5.8}{6.4}\selectfont \textcolor{ForestGreen}{$ \blacktriangle\,3.3\times $}}\\[-.2pt]4.5s\\[-.2pt]{\fontsize{5.8}{6.4}\selectfont \textcolor{ForestGreen}{$ \blacktriangle\,4.2\times $}}\end{tabular} & \hspace*{1pt}\begin{tabular}[t]{@{}c@{}}\rule{0pt}{12pt}8.8s\\[-.2pt]{\fontsize{5.8}{6.4}\selectfont \textcolor{ForestGreen}{$ \blacktriangle\,8.3\times $}}\\[-.2pt]7.9s\\[-.2pt]{\fontsize{5.8}{6.4}\selectfont \textcolor{ForestGreen}{$ \blacktriangle\,8.3\times $}}\end{tabular} & \hspace*{1pt}\begin{tabular}[t]{@{}c@{}}\rule{0pt}{12pt}\colorbox{BenchBest}{\strut 8.2s}\\[-.2pt]{\fontsize{5.8}{6.4}\selectfont \textcolor{ForestGreen}{$ \blacktriangle\,17\% $}}\\[-.2pt]\colorbox{BenchBest}{\strut 5.8s}\\[-.2pt]{\fontsize{5.8}{6.4}\selectfont \textcolor{ForestGreen}{$ \blacktriangle\,1.5\times $}}\end{tabular} & \hspace*{1pt}\begin{tabular}[t]{@{}c@{}}\rule{0pt}{12pt}\colorbox{BenchBest}{\strut 6.6s}\\[-.2pt]{\fontsize{5.8}{6.4}\selectfont \textcolor{ForestGreen}{$ \blacktriangle\,1.4\times $}}\\[-.2pt]\colorbox{BenchBest}{\strut 4.6s}\\[-.2pt]{\fontsize{5.8}{6.4}\selectfont \textcolor{ForestGreen}{$ \blacktriangle\,1.6\times $}}\end{tabular} & \hspace*{1pt}\begin{tabular}[t]{@{}c@{}}\rule{0pt}{12pt}13.7s\\[-.2pt]{\fontsize{5.8}{6.4}\selectfont \textcolor{ForestGreen}{$ \blacktriangle\,17\% $}}\\[-.2pt]8.8s\\[-.2pt]{\fontsize{5.8}{6.4}\selectfont \textcolor{ForestGreen}{$ \blacktriangle\,1.5\times $}}\end{tabular} & \textcolor{gray}{\ensuremath{\cdot}} & \hspace*{1pt}\begin{tabular}[t]{@{}c@{}}\rule{0pt}{12pt}1m\,56.4s\\[-.2pt]{\fontsize{5.8}{6.4}\selectfont \textcolor{BrickRed}{$ \blacktriangledown\,2.2\times $}}\\[-.2pt]1m\,48.1s\\[-.2pt]{\fontsize{5.8}{6.4}\selectfont \textcolor{BrickRed}{$ \blacktriangledown\,2.1\times $}}\end{tabular} & \hspace*{1pt}\begin{tabular}[t]{@{}c@{}}\rule{0pt}{12pt}1m\,56.2s\\[-.2pt]{\fontsize{5.8}{6.4}\selectfont \textcolor{BrickRed}{$ \blacktriangledown\,2.1\times $}}\\[-.2pt]1m\,46.2s\\[-.2pt]{\fontsize{5.8}{6.4}\selectfont \textcolor{BrickRed}{$ \blacktriangledown\,1.9\times $}}\end{tabular} & \hspace*{1pt}\begin{tabular}[t]{@{}c@{}}\rule{0pt}{12pt}3m\,42.9s\\[-.2pt]{\fontsize{5.8}{6.4}\selectfont \textcolor{BrickRed}{$ \blacktriangledown\,1.9\times $}}\\[-.2pt]3m\,25s\\[-.2pt]{\fontsize{5.8}{6.4}\selectfont \textcolor{BrickRed}{$ \blacktriangledown\,1.8\times $}}\end{tabular} & \textcolor{gray}{\ensuremath{\cdot}} & \hspace*{1pt}\begin{tabular}[t]{@{}c@{}}\rule{0pt}{12pt}5m\,52.1s\\[-.2pt]{\fontsize{5.8}{6.4}\selectfont \textcolor{BrickRed}{$ \blacktriangledown\,7.3\times $}}\\[-.2pt]1m\,53.5s\\[-.2pt]{\fontsize{5.8}{6.4}\selectfont \textcolor{BrickRed}{$ \blacktriangledown\,2.4\times $}}\end{tabular} & \hspace*{1pt}\begin{tabular}[t]{@{}c@{}}\rule{0pt}{12pt}4m\,49.2s\\[-.2pt]{\fontsize{5.8}{6.4}\selectfont \textcolor{BrickRed}{$ \blacktriangledown\,6.5\times $}}\\[-.2pt]1m\,56.8s\\[-.2pt]{\fontsize{5.8}{6.4}\selectfont \textcolor{BrickRed}{$ \blacktriangledown\,2.7\times $}}\end{tabular} & \hspace*{1pt}\begin{tabular}[t]{@{}c@{}}\rule{0pt}{12pt}9m\,39.5s\\[-.2pt]{\fontsize{5.8}{6.4}\selectfont \textcolor{BrickRed}{$ \blacktriangledown\,6.4\times $}}\\[-.2pt]3m\,18.1s\\[-.2pt]{\fontsize{5.8}{6.4}\selectfont \textcolor{BrickRed}{$ \blacktriangledown\,2.3\times $}}\end{tabular} & \textcolor{gray}{\ensuremath{\cdot}}\\[3pt]\hline
 & \begin{tabular}[t]{@{}l@{}}\rule{0pt}{12pt}MPF / Bounded\\SIMD-DP / KSW2\end{tabular} & \hspace*{1pt}\begin{tabular}[t]{@{}c@{}}\rule{0pt}{12pt}5.2s\\[-.2pt]{\fontsize{5.8}{6.4}\selectfont \textcolor{ForestGreen}{$ \blacktriangle\,2.1\times $}}\\[-.2pt]\colorbox{BenchBest}{\strut 2.8s}\\[-.2pt]{\fontsize{5.8}{6.4}\selectfont \textcolor{ForestGreen}{$ \blacktriangle\,3.6\times $}}\end{tabular} & \hspace*{1pt}\begin{tabular}[t]{@{}c@{}}\rule{0pt}{12pt}4.7s\\[-.2pt]{\fontsize{5.8}{6.4}\selectfont \textcolor{ForestGreen}{$ \blacktriangle\,2.3\times $}}\\[-.2pt]2.5s\\[-.2pt]{\fontsize{5.8}{6.4}\selectfont \textcolor{ForestGreen}{$ \blacktriangle\,3.5\times $}}\end{tabular} & \hspace*{1pt}\begin{tabular}[t]{@{}c@{}}\rule{0pt}{12pt}9.9s\\[-.2pt]{\fontsize{5.8}{6.4}\selectfont \textcolor{ForestGreen}{$ \blacktriangle\,2.3\times $}}\\[-.2pt]4.3s\\[-.2pt]{\fontsize{5.8}{6.4}\selectfont \textcolor{ForestGreen}{$ \blacktriangle\,4.4\times $}}\end{tabular} & \hspace*{1pt}\begin{tabular}[t]{@{}c@{}}\rule{0pt}{12pt}7.4s\\[-.2pt]{\fontsize{5.8}{6.4}\selectfont \textcolor{ForestGreen}{$ \blacktriangle\,9.9\times $}}\\[-.2pt]4.4s\\[-.2pt]{\fontsize{5.8}{6.4}\selectfont \textcolor{ForestGreen}{$ \blacktriangle\,15.1\times $}}\end{tabular} & \textcolor{gray}{\ensuremath{\cdot}} & \textcolor{gray}{\ensuremath{\cdot}} & \hspace*{1pt}\begin{tabular}[t]{@{}c@{}}\rule{0pt}{12pt}14.1s\\[-.2pt]{\fontsize{5.8}{6.4}\selectfont \textcolor{ForestGreen}{$ \blacktriangle\,14\% $}}\\[-.2pt]8.5s\\[-.2pt]{\fontsize{5.8}{6.4}\selectfont \textcolor{ForestGreen}{$ \blacktriangle\,1.5\times $}}\end{tabular} & \textcolor{gray}{\ensuremath{\cdot}} & \hspace*{1pt}\begin{tabular}[t]{@{}c@{}}\rule{0pt}{12pt}2m\,11.6s\\[-.2pt]{\fontsize{5.8}{6.4}\selectfont \textcolor{BrickRed}{$ \blacktriangledown\,2.5\times $}}\\[-.2pt]2m\,1.1s\\[-.2pt]{\fontsize{5.8}{6.4}\selectfont \textcolor{BrickRed}{$ \blacktriangledown\,2.3\times $}}\end{tabular} & \hspace*{1pt}\begin{tabular}[t]{@{}c@{}}\rule{0pt}{12pt}2m\,7.5s\\[-.2pt]{\fontsize{5.8}{6.4}\selectfont \textcolor{BrickRed}{$ \blacktriangledown\,2.3\times $}}\\[-.2pt]1m\,54.8s\\[-.2pt]{\fontsize{5.8}{6.4}\selectfont \textcolor{BrickRed}{$ \blacktriangledown\,2.1\times $}}\end{tabular} & \textcolor{gray}{\ensuremath{\cdot}} & \textcolor{gray}{\ensuremath{\cdot}} & \textcolor{gray}{\ensuremath{\cdot}} & \textcolor{gray}{\ensuremath{\cdot}} & \textcolor{gray}{\ensuremath{\cdot}} & \textcolor{gray}{\ensuremath{\cdot}}\\[3pt]\hline
\end{tabular}\par}
\smallskip{\footnotesize\noindent Each cell: A, comparison, B, comparison. Times cover 1,024 queries. $\cdot$ = no approved observation. ---\textsuperscript{g} = GPU unavailable, not executed. n/a = no measured comparison.\par
\noindent \textsuperscript{f}SIMD-ref request observed as scalar, compared only as that effective path. \textsuperscript{p}Accepted by explicit preservation after wrapper postflight failure, no run receipt claimed.\par
\noindent MPF/MPS = minimizers per fragment/sequence. Row order: index representation / minimizer, then chaining / base alignment. SIMD-ref is the reference minimizer, Bounded is bounded ntHash.\par
\noindent Blue = one lowest A and one lowest B down each machine/workload column, across continuation pages. Exact ties select the first row. Repeated baselines are marked only on page 1.\par
\noindent Light gray = minimap2 baseline. Pale yellow = preferred MPF / Bounded / SIMD-DP / Block configuration, not a claim of winning every column.\par
\noindent When all Vulkan modes are unmeasured, they share one row. Baseline comparisons use the fastest measured BioseqDB effective path separately for A and B.\par}
\clearpage
\noindent{\large\bfseries AWS: 2 Gbp, k=15, w=10 (mapping interval)\hfill 2/7}\par\smallskip
{\fontsize{7.1}{8.1}\selectfont\setlength{\tabcolsep}{1pt}\setlength{\fboxsep}{2pt}\renewcommand{\arraystretch}{1.12}
\noindent\begin{tabular}{@{}p{14mm}p{27mm}|p{13.55mm}|p{13.55mm}|p{13.55mm}|p{13.55mm}|p{13.55mm}|p{13.55mm}|p{13.55mm}|p{13.55mm}|p{13.55mm}|p{13.55mm}|p{13.55mm}|p{13.55mm}|p{13.55mm}|p{13.55mm}|p{13.55mm}|p{13.55mm}|@{}}\hline
Engine & Execution & \multicolumn{8}{c|}{\textbf{Index}} & \multicolumn{8}{c|}{\textbf{No index}}\\
 &  & \multicolumn{8}{c|}{\textbf{2 Gbp}} & \multicolumn{8}{c|}{\textbf{2 Gbp}}\\
 &  & \multicolumn{8}{c|}{\textbf{k=15, w=10}} & \multicolumn{8}{c|}{\textbf{k=15, w=10}}\\
 &  & \multicolumn{4}{c|}{\textbf{DNA4}} & \multicolumn{4}{c|}{\textbf{DNA5}} & \multicolumn{4}{c|}{\textbf{DNA4}} & \multicolumn{4}{c|}{\textbf{DNA5}}\\
 &  & \multicolumn{1}{c|}{\textbf{M7A}} & \multicolumn{1}{c|}{\textbf{M7I}} & \multicolumn{1}{c|}{\textbf{M4}} & \multicolumn{1}{c|}{\textbf{C7G}} & \multicolumn{1}{c|}{\textbf{M7A}} & \multicolumn{1}{c|}{\textbf{M7I}} & \multicolumn{1}{c|}{\textbf{M4}} & \multicolumn{1}{c|}{\textbf{C7G}} & \multicolumn{1}{c|}{\textbf{M7A}} & \multicolumn{1}{c|}{\textbf{M7I}} & \multicolumn{1}{c|}{\textbf{M4}} & \multicolumn{1}{c|}{\textbf{C7G}} & \multicolumn{1}{c|}{\textbf{M7A}} & \multicolumn{1}{c|}{\textbf{M7I}} & \multicolumn{1}{c|}{\textbf{M4}} & \multicolumn{1}{c|}{\textbf{C7G}}\\
\hline
\rowcolor{BenchBaseline}[0.5pt][0.5pt]\textbf{minimap2} & \textbf{Baseline} & \hspace*{1pt}\begin{tabular}[t]{@{}c@{}}\rule{0pt}{12pt}11.2s\\[-.2pt]{\fontsize{5.8}{6.4}\selectfont \textcolor{BrickRed}{$ \blacktriangledown\,3.1\times $}}\\[-.2pt]10s\\[-.2pt]{\fontsize{5.8}{6.4}\selectfont \textcolor{BrickRed}{$ \blacktriangledown\,3.6\times $}}\end{tabular} & \hspace*{1pt}\begin{tabular}[t]{@{}c@{}}\rule{0pt}{12pt}10.7s\\[-.2pt]{\fontsize{5.8}{6.4}\selectfont \textcolor{BrickRed}{$ \blacktriangledown\,3.2\times $}}\\[-.2pt]8.6s\\[-.2pt]{\fontsize{5.8}{6.4}\selectfont \textcolor{BrickRed}{$ \blacktriangledown\,3.5\times $}}\end{tabular} & \hspace*{1pt}\begin{tabular}[t]{@{}c@{}}\rule{0pt}{12pt}22.5s\\[-.2pt]{\fontsize{5.8}{6.4}\selectfont \textcolor{BrickRed}{$ \blacktriangledown\,3.6\times $}}\\[-.2pt]18.8s\\[-.2pt]{\fontsize{5.8}{6.4}\selectfont \textcolor{BrickRed}{$ \blacktriangledown\,4.4\times $}}\end{tabular} & \hspace*{1pt}\begin{tabular}[t]{@{}c@{}}\rule{0pt}{12pt}1m\,13.1s\\[-.2pt]{\fontsize{5.8}{6.4}\selectfont \textcolor{BrickRed}{$ \blacktriangledown\,22.5\times $}}\\[-.2pt]1m\,5.7s\\[-.2pt]{\fontsize{5.8}{6.4}\selectfont \textcolor{BrickRed}{$ \blacktriangledown\,27.7\times $}}\end{tabular} & \hspace*{1pt}\begin{tabular}[t]{@{}c@{}}\rule{0pt}{12pt}9.9s\\[-.2pt]{\fontsize{5.8}{6.4}\selectfont \textcolor{BrickRed}{$ \blacktriangledown\,21\% $}}\\[-.2pt]8.8s\\[-.2pt]{\fontsize{5.8}{6.4}\selectfont \textcolor{BrickRed}{$ \blacktriangledown\,1.5\times $}}\end{tabular} & \hspace*{1pt}\begin{tabular}[t]{@{}c@{}}\rule{0pt}{12pt}9.5s\\[-.2pt]{\fontsize{5.8}{6.4}\selectfont \textcolor{BrickRed}{$ \blacktriangledown\,1.4\times $}}\\[-.2pt]7.6s\\[-.2pt]{\fontsize{5.8}{6.4}\selectfont \textcolor{BrickRed}{$ \blacktriangledown\,1.6\times $}}\end{tabular} & \hspace*{1pt}\begin{tabular}[t]{@{}c@{}}\rule{0pt}{12pt}16.5s\\[-.2pt]{\fontsize{5.8}{6.4}\selectfont \textcolor{BrickRed}{$ \blacktriangledown\,1.7\times $}}\\[-.2pt]12.9s\\[-.2pt]{\fontsize{5.8}{6.4}\selectfont \textcolor{BrickRed}{$ \blacktriangledown\,2.0\times $}}\end{tabular} & \hspace*{1pt}\begin{tabular}[t]{@{}c@{}}\rule{0pt}{12pt}28.2s\\[-.2pt]{\fontsize{5.8}{6.4}\selectfont \textcolor{BrickRed}{$ \blacktriangledown\,2.1\times $}}\\[-.2pt]26.2s\\[-.2pt]{\fontsize{5.8}{6.4}\selectfont \textcolor{BrickRed}{$ \blacktriangledown\,2.4\times $}}\end{tabular} & \hspace*{1pt}\begin{tabular}[t]{@{}c@{}}\rule{0pt}{12pt}53s\\[-.2pt]{\fontsize{5.8}{6.4}\selectfont \textcolor{ForestGreen}{$ \blacktriangle\,2.2\times $}}\\[-.2pt]51.9s\\[-.2pt]{\fontsize{5.8}{6.4}\selectfont \textcolor{ForestGreen}{$ \blacktriangle\,2.1\times $}}\end{tabular} & \hspace*{1pt}\begin{tabular}[t]{@{}c@{}}\rule{0pt}{12pt}56.6s\\[-.2pt]{\fontsize{5.8}{6.4}\selectfont \textcolor{ForestGreen}{$ \blacktriangle\,2.1\times $}}\\[-.2pt]54.6s\\[-.2pt]{\fontsize{5.8}{6.4}\selectfont \textcolor{ForestGreen}{$ \blacktriangle\,1.9\times $}}\end{tabular} & \hspace*{1pt}\begin{tabular}[t]{@{}c@{}}\rule{0pt}{12pt}1m\,55.5s\\[-.2pt]{\fontsize{5.8}{6.4}\selectfont \textcolor{ForestGreen}{$ \blacktriangle\,1.8\times $}}\\[-.2pt]1m\,51.8s\\[-.2pt]{\fontsize{5.8}{6.4}\selectfont \textcolor{ForestGreen}{$ \blacktriangle\,1.5\times $}}\end{tabular} & \hspace*{1pt}\begin{tabular}[t]{@{}c@{}}\rule{0pt}{12pt}6m\,1.3s\\[-.2pt]{\fontsize{5.8}{6.4}\selectfont \textcolor{gray}{n/a}}\\[-.2pt]5m\,55.2s\\[-.2pt]{\fontsize{5.8}{6.4}\selectfont \textcolor{gray}{n/a}}\end{tabular} & \hspace*{1pt}\begin{tabular}[t]{@{}c@{}}\rule{0pt}{12pt}48.1s\\[-.2pt]{\fontsize{5.8}{6.4}\selectfont \textcolor{ForestGreen}{$ \blacktriangle\,7.3\times $}}\\[-.2pt]47s\\[-.2pt]{\fontsize{5.8}{6.4}\selectfont \textcolor{ForestGreen}{$ \blacktriangle\,2.4\times $}}\end{tabular} & \hspace*{1pt}\begin{tabular}[t]{@{}c@{}}\rule{0pt}{12pt}44.5s\\[-.2pt]{\fontsize{5.8}{6.4}\selectfont \textcolor{ForestGreen}{$ \blacktriangle\,6.5\times $}}\\[-.2pt]42.6s\\[-.2pt]{\fontsize{5.8}{6.4}\selectfont \textcolor{ForestGreen}{$ \blacktriangle\,2.7\times $}}\end{tabular} & \hspace*{1pt}\begin{tabular}[t]{@{}c@{}}\rule{0pt}{12pt}1m\,30.2s\\[-.2pt]{\fontsize{5.8}{6.4}\selectfont \textcolor{ForestGreen}{$ \blacktriangle\,6.4\times $}}\\[-.2pt]1m\,26.7s\\[-.2pt]{\fontsize{5.8}{6.4}\selectfont \textcolor{ForestGreen}{$ \blacktriangle\,2.2\times $}}\end{tabular} & \hspace*{1pt}\begin{tabular}[t]{@{}c@{}}\rule{0pt}{12pt}2m\,46.4s\\[-.2pt]{\fontsize{5.8}{6.4}\selectfont \textcolor{gray}{n/a}}\\[-.2pt]2m\,42.7s\\[-.2pt]{\fontsize{5.8}{6.4}\selectfont \textcolor{gray}{n/a}}\end{tabular}\\[3pt]\hline
\textbf{BioseqDB} & \begin{tabular}[t]{@{}l@{}}\rule{0pt}{12pt}MPF / Scalar\\RMQ / Block\end{tabular} & \hspace*{1pt}\begin{tabular}[t]{@{}c@{}}\rule{0pt}{12pt}4s\\[-.2pt]{\fontsize{5.8}{6.4}\selectfont \textcolor{ForestGreen}{$ \blacktriangle\,2.8\times $}}\\[-.2pt]3.1s\\[-.2pt]{\fontsize{5.8}{6.4}\selectfont \textcolor{ForestGreen}{$ \blacktriangle\,3.2\times $}}\end{tabular} & \hspace*{1pt}\begin{tabular}[t]{@{}c@{}}\rule{0pt}{12pt}3.7s\\[-.2pt]{\fontsize{5.8}{6.4}\selectfont \textcolor{ForestGreen}{$ \blacktriangle\,2.9\times $}}\\[-.2pt]2.8s\\[-.2pt]{\fontsize{5.8}{6.4}\selectfont \textcolor{ForestGreen}{$ \blacktriangle\,3.1\times $}}\end{tabular} & \hspace*{1pt}\begin{tabular}[t]{@{}c@{}}\rule{0pt}{12pt}7.1s\\[-.2pt]{\fontsize{5.8}{6.4}\selectfont \textcolor{ForestGreen}{$ \blacktriangle\,3.2\times $}}\\[-.2pt]4.8s\\[-.2pt]{\fontsize{5.8}{6.4}\selectfont \textcolor{ForestGreen}{$ \blacktriangle\,3.9\times $}}\end{tabular} & \hspace*{1pt}\begin{tabular}[t]{@{}c@{}}\rule{0pt}{12pt}3.5s\\[-.2pt]{\fontsize{5.8}{6.4}\selectfont \textcolor{ForestGreen}{$ \blacktriangle\,20.6\times $}}\\[-.2pt]2.6s\\[-.2pt]{\fontsize{5.8}{6.4}\selectfont \textcolor{ForestGreen}{$ \blacktriangle\,25.3\times $}}\end{tabular} & \textcolor{gray}{\ensuremath{\cdot}} & \textcolor{gray}{\ensuremath{\cdot}} & \hspace*{1pt}\begin{tabular}[t]{@{}c@{}}\rule{0pt}{12pt}\colorbox{BenchBest}{\strut 9.5s}\\[-.2pt]{\fontsize{5.8}{6.4}\selectfont \textcolor{ForestGreen}{$ \blacktriangle\,1.7\times $}}\\[-.2pt]\colorbox{BenchBest}{\strut 6.5s}\\[-.2pt]{\fontsize{5.8}{6.4}\selectfont \textcolor{ForestGreen}{$ \blacktriangle\,2.0\times $}}\end{tabular} & \textcolor{gray}{\ensuremath{\cdot}} & \hspace*{1pt}\begin{tabular}[t]{@{}c@{}}\rule{0pt}{12pt}2m\,51.5s\\[-.2pt]{\fontsize{5.8}{6.4}\selectfont \textcolor{BrickRed}{$ \blacktriangledown\,3.2\times $}}\\[-.2pt]2m\,29.1s\\[-.2pt]{\fontsize{5.8}{6.4}\selectfont \textcolor{BrickRed}{$ \blacktriangledown\,2.9\times $}}\end{tabular} & \hspace*{1pt}\begin{tabular}[t]{@{}c@{}}\rule{0pt}{12pt}3m\,4.8s\\[-.2pt]{\fontsize{5.8}{6.4}\selectfont \textcolor{BrickRed}{$ \blacktriangledown\,3.3\times $}}\\[-.2pt]2m\,44.5s\\[-.2pt]{\fontsize{5.8}{6.4}\selectfont \textcolor{BrickRed}{$ \blacktriangledown\,3.0\times $}}\end{tabular} & \textcolor{gray}{\ensuremath{\cdot}} & \textcolor{gray}{\ensuremath{\cdot}} & \textcolor{gray}{\ensuremath{\cdot}} & \textcolor{gray}{\ensuremath{\cdot}} & \textcolor{gray}{\ensuremath{\cdot}} & \textcolor{gray}{\ensuremath{\cdot}}\\[3pt]\hline
 & \begin{tabular}[t]{@{}l@{}}\rule{0pt}{12pt}MPF / Scalar\\RMQ / KSW2\end{tabular} & \hspace*{1pt}\begin{tabular}[t]{@{}c@{}}\rule{0pt}{12pt}5.5s\\[-.2pt]{\fontsize{5.8}{6.4}\selectfont \textcolor{ForestGreen}{$ \blacktriangle\,2.0\times $}}\\[-.2pt]2.9s\\[-.2pt]{\fontsize{5.8}{6.4}\selectfont \textcolor{ForestGreen}{$ \blacktriangle\,3.4\times $}}\end{tabular} & \hspace*{1pt}\begin{tabular}[t]{@{}c@{}}\rule{0pt}{12pt}5.1s\\[-.2pt]{\fontsize{5.8}{6.4}\selectfont \textcolor{ForestGreen}{$ \blacktriangle\,2.1\times $}}\\[-.2pt]2.7s\\[-.2pt]{\fontsize{5.8}{6.4}\selectfont \textcolor{ForestGreen}{$ \blacktriangle\,3.2\times $}}\end{tabular} & \hspace*{1pt}\begin{tabular}[t]{@{}c@{}}\rule{0pt}{12pt}10.3s\\[-.2pt]{\fontsize{5.8}{6.4}\selectfont \textcolor{ForestGreen}{$ \blacktriangle\,2.2\times $}}\\[-.2pt]4.6s\\[-.2pt]{\fontsize{5.8}{6.4}\selectfont \textcolor{ForestGreen}{$ \blacktriangle\,4.1\times $}}\end{tabular} & \hspace*{1pt}\begin{tabular}[t]{@{}c@{}}\rule{0pt}{12pt}5.6s\\[-.2pt]{\fontsize{5.8}{6.4}\selectfont \textcolor{ForestGreen}{$ \blacktriangle\,13.1\times $}}\\[-.2pt]2.5s\\[-.2pt]{\fontsize{5.8}{6.4}\selectfont \textcolor{ForestGreen}{$ \blacktriangle\,26.3\times $}}\end{tabular} & \textcolor{gray}{\ensuremath{\cdot}} & \textcolor{gray}{\ensuremath{\cdot}} & \textcolor{gray}{\ensuremath{\cdot}} & \textcolor{gray}{\ensuremath{\cdot}} & \hspace*{1pt}\begin{tabular}[t]{@{}c@{}}\rule{0pt}{12pt}3m\,1.5s\\[-.2pt]{\fontsize{5.8}{6.4}\selectfont \textcolor{BrickRed}{$ \blacktriangledown\,3.4\times $}}\\[-.2pt]2m\,37.7s\\[-.2pt]{\fontsize{5.8}{6.4}\selectfont \textcolor{BrickRed}{$ \blacktriangledown\,3.0\times $}}\end{tabular} & \hspace*{1pt}\begin{tabular}[t]{@{}c@{}}\rule{0pt}{12pt}2m\,54.6s\\[-.2pt]{\fontsize{5.8}{6.4}\selectfont \textcolor{BrickRed}{$ \blacktriangledown\,3.1\times $}}\\[-.2pt]2m\,33.3s\\[-.2pt]{\fontsize{5.8}{6.4}\selectfont \textcolor{BrickRed}{$ \blacktriangledown\,2.8\times $}}\end{tabular} & \textcolor{gray}{\ensuremath{\cdot}} & \textcolor{gray}{\ensuremath{\cdot}} & \textcolor{gray}{\ensuremath{\cdot}} & \textcolor{gray}{\ensuremath{\cdot}} & \textcolor{gray}{\ensuremath{\cdot}} & \textcolor{gray}{\ensuremath{\cdot}}\\[3pt]\hline
 & \begin{tabular}[t]{@{}l@{}}\rule{0pt}{12pt}MPF / Scalar\\Scalar-DP / Block\end{tabular} & \hspace*{1pt}\begin{tabular}[t]{@{}c@{}}\rule{0pt}{12pt}3.8s\\[-.2pt]{\fontsize{5.8}{6.4}\selectfont \textcolor{ForestGreen}{$ \blacktriangle\,3.0\times $}}\\[-.2pt]3.1s\\[-.2pt]{\fontsize{5.8}{6.4}\selectfont \textcolor{ForestGreen}{$ \blacktriangle\,3.3\times $}}\end{tabular} & \hspace*{1pt}\begin{tabular}[t]{@{}c@{}}\rule{0pt}{12pt}3.5s\\[-.2pt]{\fontsize{5.8}{6.4}\selectfont \textcolor{ForestGreen}{$ \blacktriangle\,3.1\times $}}\\[-.2pt]2.7s\\[-.2pt]{\fontsize{5.8}{6.4}\selectfont \textcolor{ForestGreen}{$ \blacktriangle\,3.2\times $}}\end{tabular} & \hspace*{1pt}\begin{tabular}[t]{@{}c@{}}\rule{0pt}{12pt}6.5s\\[-.2pt]{\fontsize{5.8}{6.4}\selectfont \textcolor{ForestGreen}{$ \blacktriangle\,3.4\times $}}\\[-.2pt]4.8s\\[-.2pt]{\fontsize{5.8}{6.4}\selectfont \textcolor{ForestGreen}{$ \blacktriangle\,3.9\times $}}\end{tabular} & \hspace*{1pt}\begin{tabular}[t]{@{}c@{}}\rule{0pt}{12pt}3.4s\\[-.2pt]{\fontsize{5.8}{6.4}\selectfont \textcolor{ForestGreen}{$ \blacktriangle\,21.8\times $}}\\[-.2pt]2.5s\\[-.2pt]{\fontsize{5.8}{6.4}\selectfont \textcolor{ForestGreen}{$ \blacktriangle\,26.2\times $}}\end{tabular} & \textcolor{gray}{\ensuremath{\cdot}} & \textcolor{gray}{\ensuremath{\cdot}} & \textcolor{gray}{\ensuremath{\cdot}} & \textcolor{gray}{\ensuremath{\cdot}} & \hspace*{1pt}\begin{tabular}[t]{@{}c@{}}\rule{0pt}{12pt}2m\,51.4s\\[-.2pt]{\fontsize{5.8}{6.4}\selectfont \textcolor{BrickRed}{$ \blacktriangledown\,3.2\times $}}\\[-.2pt]2m\,29.1s\\[-.2pt]{\fontsize{5.8}{6.4}\selectfont \textcolor{BrickRed}{$ \blacktriangledown\,2.9\times $}}\end{tabular} & \hspace*{1pt}\begin{tabular}[t]{@{}c@{}}\rule{0pt}{12pt}2m\,59.3s\\[-.2pt]{\fontsize{5.8}{6.4}\selectfont \textcolor{BrickRed}{$ \blacktriangledown\,3.2\times $}}\\[-.2pt]2m\,39.2s\\[-.2pt]{\fontsize{5.8}{6.4}\selectfont \textcolor{BrickRed}{$ \blacktriangledown\,2.9\times $}}\end{tabular} & \textcolor{gray}{\ensuremath{\cdot}} & \textcolor{gray}{\ensuremath{\cdot}} & \textcolor{gray}{\ensuremath{\cdot}} & \textcolor{gray}{\ensuremath{\cdot}} & \textcolor{gray}{\ensuremath{\cdot}} & \textcolor{gray}{\ensuremath{\cdot}}\\[3pt]\hline
 & \begin{tabular}[t]{@{}l@{}}\rule{0pt}{12pt}MPF / Scalar\\Scalar-DP / KSW2\end{tabular} & \hspace*{1pt}\begin{tabular}[t]{@{}c@{}}\rule{0pt}{12pt}5.2s\\[-.2pt]{\fontsize{5.8}{6.4}\selectfont \textcolor{ForestGreen}{$ \blacktriangle\,2.1\times $}}\\[-.2pt]2.9s\\[-.2pt]{\fontsize{5.8}{6.4}\selectfont \textcolor{ForestGreen}{$ \blacktriangle\,3.5\times $}}\end{tabular} & \hspace*{1pt}\begin{tabular}[t]{@{}c@{}}\rule{0pt}{12pt}4.8s\\[-.2pt]{\fontsize{5.8}{6.4}\selectfont \textcolor{ForestGreen}{$ \blacktriangle\,2.2\times $}}\\[-.2pt]2.6s\\[-.2pt]{\fontsize{5.8}{6.4}\selectfont \textcolor{ForestGreen}{$ \blacktriangle\,3.3\times $}}\end{tabular} & \hspace*{1pt}\begin{tabular}[t]{@{}c@{}}\rule{0pt}{12pt}9.5s\\[-.2pt]{\fontsize{5.8}{6.4}\selectfont \textcolor{ForestGreen}{$ \blacktriangle\,2.4\times $}}\\[-.2pt]4.5s\\[-.2pt]{\fontsize{5.8}{6.4}\selectfont \textcolor{ForestGreen}{$ \blacktriangle\,4.2\times $}}\end{tabular} & \hspace*{1pt}\begin{tabular}[t]{@{}c@{}}\rule{0pt}{12pt}8.6s\\[-.2pt]{\fontsize{5.8}{6.4}\selectfont \textcolor{ForestGreen}{$ \blacktriangle\,8.5\times $}}\\[-.2pt]5.7s\\[-.2pt]{\fontsize{5.8}{6.4}\selectfont \textcolor{ForestGreen}{$ \blacktriangle\,11.6\times $}}\end{tabular} & \textcolor{gray}{\ensuremath{\cdot}} & \textcolor{gray}{\ensuremath{\cdot}} & \textcolor{gray}{\ensuremath{\cdot}} & \textcolor{gray}{\ensuremath{\cdot}} & \hspace*{1pt}\begin{tabular}[t]{@{}c@{}}\rule{0pt}{12pt}2m\,52.3s\\[-.2pt]{\fontsize{5.8}{6.4}\selectfont \textcolor{BrickRed}{$ \blacktriangledown\,3.3\times $}}\\[-.2pt]2m\,28.6s\\[-.2pt]{\fontsize{5.8}{6.4}\selectfont \textcolor{BrickRed}{$ \blacktriangledown\,2.9\times $}}\end{tabular} & \hspace*{1pt}\begin{tabular}[t]{@{}c@{}}\rule{0pt}{12pt}3m\,4.9s\\[-.2pt]{\fontsize{5.8}{6.4}\selectfont \textcolor{BrickRed}{$ \blacktriangledown\,3.3\times $}}\\[-.2pt]2m\,43.4s\\[-.2pt]{\fontsize{5.8}{6.4}\selectfont \textcolor{BrickRed}{$ \blacktriangledown\,3.0\times $}}\end{tabular} & \textcolor{gray}{\ensuremath{\cdot}} & \textcolor{gray}{\ensuremath{\cdot}} & \textcolor{gray}{\ensuremath{\cdot}} & \textcolor{gray}{\ensuremath{\cdot}} & \textcolor{gray}{\ensuremath{\cdot}} & \textcolor{gray}{\ensuremath{\cdot}}\\[3pt]\hline
 & \begin{tabular}[t]{@{}l@{}}\rule{0pt}{12pt}MPF / Scalar\\SIMD-DP / Block\end{tabular} & \hspace*{1pt}\begin{tabular}[t]{@{}c@{}}\rule{0pt}{12pt}3.9s\\[-.2pt]{\fontsize{5.8}{6.4}\selectfont \textcolor{ForestGreen}{$ \blacktriangle\,2.9\times $}}\\[-.2pt]3.1s\\[-.2pt]{\fontsize{5.8}{6.4}\selectfont \textcolor{ForestGreen}{$ \blacktriangle\,3.3\times $}}\end{tabular} & \hspace*{1pt}\begin{tabular}[t]{@{}c@{}}\rule{0pt}{12pt}4.5s\\[-.2pt]{\fontsize{5.8}{6.4}\selectfont \textcolor{ForestGreen}{$ \blacktriangle\,2.4\times $}}\\[-.2pt]3.6s\\[-.2pt]{\fontsize{5.8}{6.4}\selectfont \textcolor{ForestGreen}{$ \blacktriangle\,2.4\times $}}\end{tabular} & \hspace*{1pt}\begin{tabular}[t]{@{}c@{}}\rule{0pt}{12pt}6.6s\\[-.2pt]{\fontsize{5.8}{6.4}\selectfont \textcolor{ForestGreen}{$ \blacktriangle\,3.4\times $}}\\[-.2pt]4.7s\\[-.2pt]{\fontsize{5.8}{6.4}\selectfont \textcolor{ForestGreen}{$ \blacktriangle\,4.0\times $}}\end{tabular} & \hspace*{1pt}\begin{tabular}[t]{@{}c@{}}\rule{0pt}{12pt}9.1s\\[-.2pt]{\fontsize{5.8}{6.4}\selectfont \textcolor{ForestGreen}{$ \blacktriangle\,8.1\times $}}\\[-.2pt]8.1s\\[-.2pt]{\fontsize{5.8}{6.4}\selectfont \textcolor{ForestGreen}{$ \blacktriangle\,8.1\times $}}\end{tabular} & \textcolor{gray}{\ensuremath{\cdot}} & \textcolor{gray}{\ensuremath{\cdot}} & \textcolor{gray}{\ensuremath{\cdot}} & \hspace*{1pt}\begin{tabular}[t]{@{}c@{}}\rule{0pt}{12pt}14.3s\\[-.2pt]{\fontsize{5.8}{6.4}\selectfont \textcolor{ForestGreen}{$ \blacktriangle\,2.0\times $}}\\[-.2pt]12.1s\\[-.2pt]{\fontsize{5.8}{6.4}\selectfont \textcolor{ForestGreen}{$ \blacktriangle\,2.2\times $}}\end{tabular} & \hspace*{1pt}\begin{tabular}[t]{@{}c@{}}\rule{0pt}{12pt}3m\,8.5s\\[-.2pt]{\fontsize{5.8}{6.4}\selectfont \textcolor{BrickRed}{$ \blacktriangledown\,3.6\times $}}\\[-.2pt]2m\,46.3s\\[-.2pt]{\fontsize{5.8}{6.4}\selectfont \textcolor{BrickRed}{$ \blacktriangledown\,3.2\times $}}\end{tabular} & \hspace*{1pt}\begin{tabular}[t]{@{}c@{}}\rule{0pt}{12pt}3m\,1.2s\\[-.2pt]{\fontsize{5.8}{6.4}\selectfont \textcolor{BrickRed}{$ \blacktriangledown\,3.2\times $}}\\[-.2pt]2m\,41.3s\\[-.2pt]{\fontsize{5.8}{6.4}\selectfont \textcolor{BrickRed}{$ \blacktriangledown\,3.0\times $}}\end{tabular} & \textcolor{gray}{\ensuremath{\cdot}} & \textcolor{gray}{\ensuremath{\cdot}} & \textcolor{gray}{\ensuremath{\cdot}} & \textcolor{gray}{\ensuremath{\cdot}} & \textcolor{gray}{\ensuremath{\cdot}} & \textcolor{gray}{\ensuremath{\cdot}}\\[3pt]\hline
 & \begin{tabular}[t]{@{}l@{}}\rule{0pt}{12pt}MPF / Scalar\\SIMD-DP / KSW2\end{tabular} & \hspace*{1pt}\begin{tabular}[t]{@{}c@{}}\rule{0pt}{12pt}5.4s\\[-.2pt]{\fontsize{5.8}{6.4}\selectfont \textcolor{ForestGreen}{$ \blacktriangle\,2.1\times $}}\\[-.2pt]2.9s\\[-.2pt]{\fontsize{5.8}{6.4}\selectfont \textcolor{ForestGreen}{$ \blacktriangle\,3.5\times $}}\end{tabular} & \hspace*{1pt}\begin{tabular}[t]{@{}c@{}}\rule{0pt}{12pt}4.9s\\[-.2pt]{\fontsize{5.8}{6.4}\selectfont \textcolor{ForestGreen}{$ \blacktriangle\,2.2\times $}}\\[-.2pt]2.6s\\[-.2pt]{\fontsize{5.8}{6.4}\selectfont \textcolor{ForestGreen}{$ \blacktriangle\,3.3\times $}}\end{tabular} & \hspace*{1pt}\begin{tabular}[t]{@{}c@{}}\rule{0pt}{12pt}10.3s\\[-.2pt]{\fontsize{5.8}{6.4}\selectfont \textcolor{ForestGreen}{$ \blacktriangle\,2.2\times $}}\\[-.2pt]4.6s\\[-.2pt]{\fontsize{5.8}{6.4}\selectfont \textcolor{ForestGreen}{$ \blacktriangle\,4.1\times $}}\end{tabular} & \hspace*{1pt}\begin{tabular}[t]{@{}c@{}}\rule{0pt}{12pt}7.8s\\[-.2pt]{\fontsize{5.8}{6.4}\selectfont \textcolor{ForestGreen}{$ \blacktriangle\,9.3\times $}}\\[-.2pt]5s\\[-.2pt]{\fontsize{5.8}{6.4}\selectfont \textcolor{ForestGreen}{$ \blacktriangle\,13.2\times $}}\end{tabular} & \textcolor{gray}{\ensuremath{\cdot}} & \textcolor{gray}{\ensuremath{\cdot}} & \textcolor{gray}{\ensuremath{\cdot}} & \textcolor{gray}{\ensuremath{\cdot}} & \hspace*{1pt}\begin{tabular}[t]{@{}c@{}}\rule{0pt}{12pt}3m\,4.9s\\[-.2pt]{\fontsize{5.8}{6.4}\selectfont \textcolor{BrickRed}{$ \blacktriangledown\,3.5\times $}}\\[-.2pt]2m\,40.5s\\[-.2pt]{\fontsize{5.8}{6.4}\selectfont \textcolor{BrickRed}{$ \blacktriangledown\,3.1\times $}}\end{tabular} & \hspace*{1pt}\begin{tabular}[t]{@{}c@{}}\rule{0pt}{12pt}3m\,7.6s\\[-.2pt]{\fontsize{5.8}{6.4}\selectfont \textcolor{BrickRed}{$ \blacktriangledown\,3.3\times $}}\\[-.2pt]2m\,46s\\[-.2pt]{\fontsize{5.8}{6.4}\selectfont \textcolor{BrickRed}{$ \blacktriangledown\,3.0\times $}}\end{tabular} & \textcolor{gray}{\ensuremath{\cdot}} & \textcolor{gray}{\ensuremath{\cdot}} & \textcolor{gray}{\ensuremath{\cdot}} & \textcolor{gray}{\ensuremath{\cdot}} & \textcolor{gray}{\ensuremath{\cdot}} & \textcolor{gray}{\ensuremath{\cdot}}\\[3pt]\hline
 & \begin{tabular}[t]{@{}l@{}}\rule{0pt}{12pt}MPF / SIMD-ref\\RMQ / Block\end{tabular} & \hspace*{1pt}\begin{tabular}[t]{@{}c@{}}\rule{0pt}{12pt}3.7s\\[-.2pt]{\fontsize{5.8}{6.4}\selectfont \textcolor{ForestGreen}{$ \blacktriangle\,3.0\times $}}\\[-.2pt]3s\\[-.2pt]{\fontsize{5.8}{6.4}\selectfont \textcolor{ForestGreen}{$ \blacktriangle\,3.4\times $}}\end{tabular} & \hspace*{1pt}\begin{tabular}[t]{@{}c@{}}\rule{0pt}{12pt}3.5s\\[-.2pt]{\fontsize{5.8}{6.4}\selectfont \textcolor{ForestGreen}{$ \blacktriangle\,3.0\times $}}\\[-.2pt]2.7s\\[-.2pt]{\fontsize{5.8}{6.4}\selectfont \textcolor{ForestGreen}{$ \blacktriangle\,3.2\times $}}\end{tabular} & \hspace*{1pt}\begin{tabular}[t]{@{}c@{}}\rule{0pt}{12pt}6.7s\\[-.2pt]{\fontsize{5.8}{6.4}\selectfont \textcolor{ForestGreen}{$ \blacktriangle\,3.4\times $}}\\[-.2pt]4.6s\\[-.2pt]{\fontsize{5.8}{6.4}\selectfont \textcolor{ForestGreen}{$ \blacktriangle\,4.1\times $}}\end{tabular} & \hspace*{1pt}\begin{tabular}[t]{@{}c@{}}\rule{0pt}{12pt}3.5s\\[-.2pt]{\fontsize{5.8}{6.4}\selectfont \textcolor{ForestGreen}{$ \blacktriangle\,21.0\times $}}\\[-.2pt]2.6s\\[-.2pt]{\fontsize{5.8}{6.4}\selectfont \textcolor{ForestGreen}{$ \blacktriangle\,25.5\times $}}\end{tabular} & \textcolor{gray}{\ensuremath{\cdot}} & \textcolor{gray}{\ensuremath{\cdot}} & \textcolor{gray}{\ensuremath{\cdot}} & \textcolor{gray}{\ensuremath{\cdot}} & \textcolor{gray}{\ensuremath{\cdot}} & \textcolor{gray}{\ensuremath{\cdot}} & \textcolor{gray}{\ensuremath{\cdot}} & \textcolor{gray}{\ensuremath{\cdot}} & \textcolor{gray}{\ensuremath{\cdot}} & \textcolor{gray}{\ensuremath{\cdot}} & \textcolor{gray}{\ensuremath{\cdot}} & \textcolor{gray}{\ensuremath{\cdot}}\\[3pt]\hline
\end{tabular}\par}
\smallskip{\footnotesize\noindent Each cell: A, comparison, B, comparison. Times cover 1,024 queries. $\cdot$ = no approved observation. ---\textsuperscript{g} = GPU unavailable, not executed. n/a = no measured comparison.\par
\noindent \textsuperscript{f}SIMD-ref request observed as scalar, compared only as that effective path. \textsuperscript{p}Accepted by explicit preservation after wrapper postflight failure, no run receipt claimed.\par
\noindent MPF/MPS = minimizers per fragment/sequence. Row order: index representation / minimizer, then chaining / base alignment. SIMD-ref is the reference minimizer, Bounded is bounded ntHash.\par
\noindent Blue = one lowest A and one lowest B down each machine/workload column, across continuation pages. Exact ties select the first row. Repeated baselines are marked only on page 1.\par
\noindent Light gray = minimap2 baseline. Pale yellow = preferred MPF / Bounded / SIMD-DP / Block configuration, not a claim of winning every column.\par
\noindent When all Vulkan modes are unmeasured, they share one row. Baseline comparisons use the fastest measured BioseqDB effective path separately for A and B.\par}
\clearpage
\noindent{\large\bfseries AWS: 2 Gbp, k=15, w=10 (mapping interval)\hfill 3/7}\par\smallskip
{\fontsize{7.1}{8.1}\selectfont\setlength{\tabcolsep}{1pt}\setlength{\fboxsep}{2pt}\renewcommand{\arraystretch}{1.12}
\noindent\begin{tabular}{@{}p{14mm}p{27mm}|p{13.55mm}|p{13.55mm}|p{13.55mm}|p{13.55mm}|p{13.55mm}|p{13.55mm}|p{13.55mm}|p{13.55mm}|p{13.55mm}|p{13.55mm}|p{13.55mm}|p{13.55mm}|p{13.55mm}|p{13.55mm}|p{13.55mm}|p{13.55mm}|@{}}\hline
Engine & Execution & \multicolumn{8}{c|}{\textbf{Index}} & \multicolumn{8}{c|}{\textbf{No index}}\\
 &  & \multicolumn{8}{c|}{\textbf{2 Gbp}} & \multicolumn{8}{c|}{\textbf{2 Gbp}}\\
 &  & \multicolumn{8}{c|}{\textbf{k=15, w=10}} & \multicolumn{8}{c|}{\textbf{k=15, w=10}}\\
 &  & \multicolumn{4}{c|}{\textbf{DNA4}} & \multicolumn{4}{c|}{\textbf{DNA5}} & \multicolumn{4}{c|}{\textbf{DNA4}} & \multicolumn{4}{c|}{\textbf{DNA5}}\\
 &  & \multicolumn{1}{c|}{\textbf{M7A}} & \multicolumn{1}{c|}{\textbf{M7I}} & \multicolumn{1}{c|}{\textbf{M4}} & \multicolumn{1}{c|}{\textbf{C7G}} & \multicolumn{1}{c|}{\textbf{M7A}} & \multicolumn{1}{c|}{\textbf{M7I}} & \multicolumn{1}{c|}{\textbf{M4}} & \multicolumn{1}{c|}{\textbf{C7G}} & \multicolumn{1}{c|}{\textbf{M7A}} & \multicolumn{1}{c|}{\textbf{M7I}} & \multicolumn{1}{c|}{\textbf{M4}} & \multicolumn{1}{c|}{\textbf{C7G}} & \multicolumn{1}{c|}{\textbf{M7A}} & \multicolumn{1}{c|}{\textbf{M7I}} & \multicolumn{1}{c|}{\textbf{M4}} & \multicolumn{1}{c|}{\textbf{C7G}}\\
\hline
\rowcolor{BenchBaseline}[0.5pt][0.5pt]\textbf{minimap2} & \textbf{Baseline} & \hspace*{1pt}\begin{tabular}[t]{@{}c@{}}\rule{0pt}{12pt}11.2s\\[-.2pt]{\fontsize{5.8}{6.4}\selectfont \textcolor{BrickRed}{$ \blacktriangledown\,3.1\times $}}\\[-.2pt]10s\\[-.2pt]{\fontsize{5.8}{6.4}\selectfont \textcolor{BrickRed}{$ \blacktriangledown\,3.6\times $}}\end{tabular} & \hspace*{1pt}\begin{tabular}[t]{@{}c@{}}\rule{0pt}{12pt}10.7s\\[-.2pt]{\fontsize{5.8}{6.4}\selectfont \textcolor{BrickRed}{$ \blacktriangledown\,3.2\times $}}\\[-.2pt]8.6s\\[-.2pt]{\fontsize{5.8}{6.4}\selectfont \textcolor{BrickRed}{$ \blacktriangledown\,3.5\times $}}\end{tabular} & \hspace*{1pt}\begin{tabular}[t]{@{}c@{}}\rule{0pt}{12pt}22.5s\\[-.2pt]{\fontsize{5.8}{6.4}\selectfont \textcolor{BrickRed}{$ \blacktriangledown\,3.6\times $}}\\[-.2pt]18.8s\\[-.2pt]{\fontsize{5.8}{6.4}\selectfont \textcolor{BrickRed}{$ \blacktriangledown\,4.4\times $}}\end{tabular} & \hspace*{1pt}\begin{tabular}[t]{@{}c@{}}\rule{0pt}{12pt}1m\,13.1s\\[-.2pt]{\fontsize{5.8}{6.4}\selectfont \textcolor{BrickRed}{$ \blacktriangledown\,22.5\times $}}\\[-.2pt]1m\,5.7s\\[-.2pt]{\fontsize{5.8}{6.4}\selectfont \textcolor{BrickRed}{$ \blacktriangledown\,27.7\times $}}\end{tabular} & \hspace*{1pt}\begin{tabular}[t]{@{}c@{}}\rule{0pt}{12pt}9.9s\\[-.2pt]{\fontsize{5.8}{6.4}\selectfont \textcolor{BrickRed}{$ \blacktriangledown\,21\% $}}\\[-.2pt]8.8s\\[-.2pt]{\fontsize{5.8}{6.4}\selectfont \textcolor{BrickRed}{$ \blacktriangledown\,1.5\times $}}\end{tabular} & \hspace*{1pt}\begin{tabular}[t]{@{}c@{}}\rule{0pt}{12pt}9.5s\\[-.2pt]{\fontsize{5.8}{6.4}\selectfont \textcolor{BrickRed}{$ \blacktriangledown\,1.4\times $}}\\[-.2pt]7.6s\\[-.2pt]{\fontsize{5.8}{6.4}\selectfont \textcolor{BrickRed}{$ \blacktriangledown\,1.6\times $}}\end{tabular} & \hspace*{1pt}\begin{tabular}[t]{@{}c@{}}\rule{0pt}{12pt}16.5s\\[-.2pt]{\fontsize{5.8}{6.4}\selectfont \textcolor{BrickRed}{$ \blacktriangledown\,1.7\times $}}\\[-.2pt]12.9s\\[-.2pt]{\fontsize{5.8}{6.4}\selectfont \textcolor{BrickRed}{$ \blacktriangledown\,2.0\times $}}\end{tabular} & \hspace*{1pt}\begin{tabular}[t]{@{}c@{}}\rule{0pt}{12pt}28.2s\\[-.2pt]{\fontsize{5.8}{6.4}\selectfont \textcolor{BrickRed}{$ \blacktriangledown\,2.1\times $}}\\[-.2pt]26.2s\\[-.2pt]{\fontsize{5.8}{6.4}\selectfont \textcolor{BrickRed}{$ \blacktriangledown\,2.4\times $}}\end{tabular} & \hspace*{1pt}\begin{tabular}[t]{@{}c@{}}\rule{0pt}{12pt}53s\\[-.2pt]{\fontsize{5.8}{6.4}\selectfont \textcolor{ForestGreen}{$ \blacktriangle\,2.2\times $}}\\[-.2pt]51.9s\\[-.2pt]{\fontsize{5.8}{6.4}\selectfont \textcolor{ForestGreen}{$ \blacktriangle\,2.1\times $}}\end{tabular} & \hspace*{1pt}\begin{tabular}[t]{@{}c@{}}\rule{0pt}{12pt}56.6s\\[-.2pt]{\fontsize{5.8}{6.4}\selectfont \textcolor{ForestGreen}{$ \blacktriangle\,2.1\times $}}\\[-.2pt]54.6s\\[-.2pt]{\fontsize{5.8}{6.4}\selectfont \textcolor{ForestGreen}{$ \blacktriangle\,1.9\times $}}\end{tabular} & \hspace*{1pt}\begin{tabular}[t]{@{}c@{}}\rule{0pt}{12pt}1m\,55.5s\\[-.2pt]{\fontsize{5.8}{6.4}\selectfont \textcolor{ForestGreen}{$ \blacktriangle\,1.8\times $}}\\[-.2pt]1m\,51.8s\\[-.2pt]{\fontsize{5.8}{6.4}\selectfont \textcolor{ForestGreen}{$ \blacktriangle\,1.5\times $}}\end{tabular} & \hspace*{1pt}\begin{tabular}[t]{@{}c@{}}\rule{0pt}{12pt}6m\,1.3s\\[-.2pt]{\fontsize{5.8}{6.4}\selectfont \textcolor{gray}{n/a}}\\[-.2pt]5m\,55.2s\\[-.2pt]{\fontsize{5.8}{6.4}\selectfont \textcolor{gray}{n/a}}\end{tabular} & \hspace*{1pt}\begin{tabular}[t]{@{}c@{}}\rule{0pt}{12pt}48.1s\\[-.2pt]{\fontsize{5.8}{6.4}\selectfont \textcolor{ForestGreen}{$ \blacktriangle\,7.3\times $}}\\[-.2pt]47s\\[-.2pt]{\fontsize{5.8}{6.4}\selectfont \textcolor{ForestGreen}{$ \blacktriangle\,2.4\times $}}\end{tabular} & \hspace*{1pt}\begin{tabular}[t]{@{}c@{}}\rule{0pt}{12pt}44.5s\\[-.2pt]{\fontsize{5.8}{6.4}\selectfont \textcolor{ForestGreen}{$ \blacktriangle\,6.5\times $}}\\[-.2pt]42.6s\\[-.2pt]{\fontsize{5.8}{6.4}\selectfont \textcolor{ForestGreen}{$ \blacktriangle\,2.7\times $}}\end{tabular} & \hspace*{1pt}\begin{tabular}[t]{@{}c@{}}\rule{0pt}{12pt}1m\,30.2s\\[-.2pt]{\fontsize{5.8}{6.4}\selectfont \textcolor{ForestGreen}{$ \blacktriangle\,6.4\times $}}\\[-.2pt]1m\,26.7s\\[-.2pt]{\fontsize{5.8}{6.4}\selectfont \textcolor{ForestGreen}{$ \blacktriangle\,2.2\times $}}\end{tabular} & \hspace*{1pt}\begin{tabular}[t]{@{}c@{}}\rule{0pt}{12pt}2m\,46.4s\\[-.2pt]{\fontsize{5.8}{6.4}\selectfont \textcolor{gray}{n/a}}\\[-.2pt]2m\,42.7s\\[-.2pt]{\fontsize{5.8}{6.4}\selectfont \textcolor{gray}{n/a}}\end{tabular}\\[3pt]\hline
\textbf{BioseqDB} & \begin{tabular}[t]{@{}l@{}}\rule{0pt}{12pt}MPF / SIMD-ref\\RMQ / KSW2\end{tabular} & \hspace*{1pt}\begin{tabular}[t]{@{}c@{}}\rule{0pt}{12pt}5.3s\\[-.2pt]{\fontsize{5.8}{6.4}\selectfont \textcolor{ForestGreen}{$ \blacktriangle\,2.1\times $}}\\[-.2pt]2.8s\\[-.2pt]{\fontsize{5.8}{6.4}\selectfont \textcolor{ForestGreen}{$ \blacktriangle\,3.6\times $}}\end{tabular} & \hspace*{1pt}\begin{tabular}[t]{@{}c@{}}\rule{0pt}{12pt}4.9s\\[-.2pt]{\fontsize{5.8}{6.4}\selectfont \textcolor{ForestGreen}{$ \blacktriangle\,2.2\times $}}\\[-.2pt]2.5s\\[-.2pt]{\fontsize{5.8}{6.4}\selectfont \textcolor{ForestGreen}{$ \blacktriangle\,3.4\times $}}\end{tabular} & \hspace*{1pt}\begin{tabular}[t]{@{}c@{}}\rule{0pt}{12pt}9.9s\\[-.2pt]{\fontsize{5.8}{6.4}\selectfont \textcolor{ForestGreen}{$ \blacktriangle\,2.3\times $}}\\[-.2pt]4.3s\\[-.2pt]{\fontsize{5.8}{6.4}\selectfont \textcolor{ForestGreen}{$ \blacktriangle\,4.4\times $}}\end{tabular} & \hspace*{1pt}\begin{tabular}[t]{@{}c@{}}\rule{0pt}{12pt}5.5s\\[-.2pt]{\fontsize{5.8}{6.4}\selectfont \textcolor{ForestGreen}{$ \blacktriangle\,13.3\times $}}\\[-.2pt]2.4s\\[-.2pt]{\fontsize{5.8}{6.4}\selectfont \textcolor{ForestGreen}{$ \blacktriangle\,26.8\times $}}\end{tabular} & \textcolor{gray}{\ensuremath{\cdot}} & \textcolor{gray}{\ensuremath{\cdot}} & \textcolor{gray}{\ensuremath{\cdot}} & \textcolor{gray}{\ensuremath{\cdot}} & \textcolor{gray}{\ensuremath{\cdot}} & \textcolor{gray}{\ensuremath{\cdot}} & \textcolor{gray}{\ensuremath{\cdot}} & \textcolor{gray}{\ensuremath{\cdot}} & \textcolor{gray}{\ensuremath{\cdot}} & \textcolor{gray}{\ensuremath{\cdot}} & \textcolor{gray}{\ensuremath{\cdot}} & \textcolor{gray}{\ensuremath{\cdot}}\\[3pt]\hline
 & \begin{tabular}[t]{@{}l@{}}\rule{0pt}{12pt}MPF / SIMD-ref\\Scalar-DP / Block\end{tabular} & \hspace*{1pt}\begin{tabular}[t]{@{}c@{}}\rule{0pt}{12pt}\colorbox{BenchBest}{\strut 3.6s}\\[-.2pt]{\fontsize{5.8}{6.4}\selectfont \textcolor{ForestGreen}{$ \blacktriangle\,3.1\times $}}\\[-.2pt]2.9s\\[-.2pt]{\fontsize{5.8}{6.4}\selectfont \textcolor{ForestGreen}{$ \blacktriangle\,3.5\times $}}\end{tabular} & \hspace*{1pt}\begin{tabular}[t]{@{}c@{}}\rule{0pt}{12pt}\colorbox{BenchBest}{\strut 3.4s}\\[-.2pt]{\fontsize{5.8}{6.4}\selectfont \textcolor{ForestGreen}{$ \blacktriangle\,3.2\times $}}\\[-.2pt]2.6s\\[-.2pt]{\fontsize{5.8}{6.4}\selectfont \textcolor{ForestGreen}{$ \blacktriangle\,3.3\times $}}\end{tabular} & \hspace*{1pt}\begin{tabular}[t]{@{}c@{}}\rule{0pt}{12pt}6.5s\\[-.2pt]{\fontsize{5.8}{6.4}\selectfont \textcolor{ForestGreen}{$ \blacktriangle\,3.5\times $}}\\[-.2pt]4.4s\\[-.2pt]{\fontsize{5.8}{6.4}\selectfont \textcolor{ForestGreen}{$ \blacktriangle\,4.2\times $}}\end{tabular} & \hspace*{1pt}\begin{tabular}[t]{@{}c@{}}\rule{0pt}{12pt}3.3s\\[-.2pt]{\fontsize{5.8}{6.4}\selectfont \textcolor{ForestGreen}{$ \blacktriangle\,22.1\times $}}\\[-.2pt]2.5s\\[-.2pt]{\fontsize{5.8}{6.4}\selectfont \textcolor{ForestGreen}{$ \blacktriangle\,26.2\times $}}\end{tabular} & \textcolor{gray}{\ensuremath{\cdot}} & \textcolor{gray}{\ensuremath{\cdot}} & \textcolor{gray}{\ensuremath{\cdot}} & \textcolor{gray}{\ensuremath{\cdot}} & \textcolor{gray}{\ensuremath{\cdot}} & \textcolor{gray}{\ensuremath{\cdot}} & \textcolor{gray}{\ensuremath{\cdot}} & \textcolor{gray}{\ensuremath{\cdot}} & \textcolor{gray}{\ensuremath{\cdot}} & \textcolor{gray}{\ensuremath{\cdot}} & \textcolor{gray}{\ensuremath{\cdot}} & \textcolor{gray}{\ensuremath{\cdot}}\\[3pt]\hline
 & \begin{tabular}[t]{@{}l@{}}\rule{0pt}{12pt}MPF / SIMD-ref\\Scalar-DP / KSW2\end{tabular} & \hspace*{1pt}\begin{tabular}[t]{@{}c@{}}\rule{0pt}{12pt}5.1s\\[-.2pt]{\fontsize{5.8}{6.4}\selectfont \textcolor{ForestGreen}{$ \blacktriangle\,2.2\times $}}\\[-.2pt]2.8s\\[-.2pt]{\fontsize{5.8}{6.4}\selectfont \textcolor{ForestGreen}{$ \blacktriangle\,3.6\times $}}\end{tabular} & \hspace*{1pt}\begin{tabular}[t]{@{}c@{}}\rule{0pt}{12pt}4.7s\\[-.2pt]{\fontsize{5.8}{6.4}\selectfont \textcolor{ForestGreen}{$ \blacktriangle\,2.3\times $}}\\[-.2pt]2.5s\\[-.2pt]{\fontsize{5.8}{6.4}\selectfont \textcolor{ForestGreen}{$ \blacktriangle\,3.4\times $}}\end{tabular} & \hspace*{1pt}\begin{tabular}[t]{@{}c@{}}\rule{0pt}{12pt}9.8s\\[-.2pt]{\fontsize{5.8}{6.4}\selectfont \textcolor{ForestGreen}{$ \blacktriangle\,2.3\times $}}\\[-.2pt]4.3s\\[-.2pt]{\fontsize{5.8}{6.4}\selectfont \textcolor{ForestGreen}{$ \blacktriangle\,4.3\times $}}\end{tabular} & \hspace*{1pt}\begin{tabular}[t]{@{}c@{}}\rule{0pt}{12pt}7.2s\\[-.2pt]{\fontsize{5.8}{6.4}\selectfont \textcolor{ForestGreen}{$ \blacktriangle\,10.1\times $}}\\[-.2pt]4.4s\\[-.2pt]{\fontsize{5.8}{6.4}\selectfont \textcolor{ForestGreen}{$ \blacktriangle\,15.1\times $}}\end{tabular} & \textcolor{gray}{\ensuremath{\cdot}} & \textcolor{gray}{\ensuremath{\cdot}} & \textcolor{gray}{\ensuremath{\cdot}} & \textcolor{gray}{\ensuremath{\cdot}} & \textcolor{gray}{\ensuremath{\cdot}} & \textcolor{gray}{\ensuremath{\cdot}} & \textcolor{gray}{\ensuremath{\cdot}} & \textcolor{gray}{\ensuremath{\cdot}} & \textcolor{gray}{\ensuremath{\cdot}} & \textcolor{gray}{\ensuremath{\cdot}} & \textcolor{gray}{\ensuremath{\cdot}} & \textcolor{gray}{\ensuremath{\cdot}}\\[3pt]\hline
 & \begin{tabular}[t]{@{}l@{}}\rule{0pt}{12pt}MPF / SIMD-ref\\SIMD-DP / Block\end{tabular} & \hspace*{1pt}\begin{tabular}[t]{@{}c@{}}\rule{0pt}{12pt}3.7s\\[-.2pt]{\fontsize{5.8}{6.4}\selectfont \textcolor{ForestGreen}{$ \blacktriangle\,3.0\times $}}\\[-.2pt]2.9s\\[-.2pt]{\fontsize{5.8}{6.4}\selectfont \textcolor{ForestGreen}{$ \blacktriangle\,3.4\times $}}\end{tabular} & \hspace*{1pt}\begin{tabular}[t]{@{}c@{}}\rule{0pt}{12pt}3.5s\\[-.2pt]{\fontsize{5.8}{6.4}\selectfont \textcolor{ForestGreen}{$ \blacktriangle\,3.0\times $}}\\[-.2pt]2.7s\\[-.2pt]{\fontsize{5.8}{6.4}\selectfont \textcolor{ForestGreen}{$ \blacktriangle\,3.2\times $}}\end{tabular} & \hspace*{1pt}\begin{tabular}[t]{@{}c@{}}\rule{0pt}{12pt}6.9s\\[-.2pt]{\fontsize{5.8}{6.4}\selectfont \textcolor{ForestGreen}{$ \blacktriangle\,3.3\times $}}\\[-.2pt]4.5s\\[-.2pt]{\fontsize{5.8}{6.4}\selectfont \textcolor{ForestGreen}{$ \blacktriangle\,4.2\times $}}\end{tabular} & \hspace*{1pt}\begin{tabular}[t]{@{}c@{}}\rule{0pt}{12pt}3.4s\\[-.2pt]{\fontsize{5.8}{6.4}\selectfont \textcolor{ForestGreen}{$ \blacktriangle\,21.4\times $}}\\[-.2pt]2.5s\\[-.2pt]{\fontsize{5.8}{6.4}\selectfont \textcolor{ForestGreen}{$ \blacktriangle\,26.3\times $}}\end{tabular} & \textcolor{gray}{\ensuremath{\cdot}} & \textcolor{gray}{\ensuremath{\cdot}} & \textcolor{gray}{\ensuremath{\cdot}} & \hspace*{1pt}\begin{tabular}[t]{@{}c@{}}\rule{0pt}{12pt}26.2s\\[-.2pt]{\fontsize{5.8}{6.4}\selectfont \textcolor{ForestGreen}{$ \blacktriangle\,7\% $}}\\[-.2pt]24.4s\\[-.2pt]{\fontsize{5.8}{6.4}\selectfont \textcolor{ForestGreen}{$ \blacktriangle\,7\% $}}\end{tabular} & \textcolor{gray}{\ensuremath{\cdot}} & \hspace*{1pt}\begin{tabular}[t]{@{}c@{}}\rule{0pt}{12pt}2m\,32.6s\\[-.2pt]{\fontsize{5.8}{6.4}\selectfont \textcolor{BrickRed}{$ \blacktriangledown\,2.7\times $}}\\[-.2pt]2m\,12.9s\\[-.2pt]{\fontsize{5.8}{6.4}\selectfont \textcolor{BrickRed}{$ \blacktriangledown\,2.4\times $}}\end{tabular} & \textcolor{gray}{\ensuremath{\cdot}} & \textcolor{gray}{\ensuremath{\cdot}} & \textcolor{gray}{\ensuremath{\cdot}} & \textcolor{gray}{\ensuremath{\cdot}} & \textcolor{gray}{\ensuremath{\cdot}} & \textcolor{gray}{\ensuremath{\cdot}}\\[3pt]\hline
 & \begin{tabular}[t]{@{}l@{}}\rule{0pt}{12pt}MPF / SIMD-ref\\SIMD-DP / KSW2\end{tabular} & \hspace*{1pt}\begin{tabular}[t]{@{}c@{}}\rule{0pt}{12pt}5.3s\\[-.2pt]{\fontsize{5.8}{6.4}\selectfont \textcolor{ForestGreen}{$ \blacktriangle\,2.1\times $}}\\[-.2pt]2.8s\\[-.2pt]{\fontsize{5.8}{6.4}\selectfont \textcolor{ForestGreen}{$ \blacktriangle\,3.6\times $}}\end{tabular} & \hspace*{1pt}\begin{tabular}[t]{@{}c@{}}\rule{0pt}{12pt}4.8s\\[-.2pt]{\fontsize{5.8}{6.4}\selectfont \textcolor{ForestGreen}{$ \blacktriangle\,2.2\times $}}\\[-.2pt]2.5s\\[-.2pt]{\fontsize{5.8}{6.4}\selectfont \textcolor{ForestGreen}{$ \blacktriangle\,3.4\times $}}\end{tabular} & \hspace*{1pt}\begin{tabular}[t]{@{}c@{}}\rule{0pt}{12pt}9.7s\\[-.2pt]{\fontsize{5.8}{6.4}\selectfont \textcolor{ForestGreen}{$ \blacktriangle\,2.3\times $}}\\[-.2pt]4.3s\\[-.2pt]{\fontsize{5.8}{6.4}\selectfont \textcolor{ForestGreen}{$ \blacktriangle\,4.4\times $}}\end{tabular} & \hspace*{1pt}\begin{tabular}[t]{@{}c@{}}\rule{0pt}{12pt}9s\\[-.2pt]{\fontsize{5.8}{6.4}\selectfont \textcolor{ForestGreen}{$ \blacktriangle\,8.1\times $}}\\[-.2pt]6.2s\\[-.2pt]{\fontsize{5.8}{6.4}\selectfont \textcolor{ForestGreen}{$ \blacktriangle\,10.6\times $}}\end{tabular} & \textcolor{gray}{\ensuremath{\cdot}} & \textcolor{gray}{\ensuremath{\cdot}} & \textcolor{gray}{\ensuremath{\cdot}} & \textcolor{gray}{\ensuremath{\cdot}} & \textcolor{gray}{\ensuremath{\cdot}} & \textcolor{gray}{\ensuremath{\cdot}} & \hspace*{1pt}\begin{tabular}[t]{@{}c@{}}\rule{0pt}{12pt}3m\,23.3s\\[-.2pt]{\fontsize{5.8}{6.4}\selectfont \textcolor{BrickRed}{$ \blacktriangledown\,1.8\times $}}\\[-.2pt]2m\,50.3s\\[-.2pt]{\fontsize{5.8}{6.4}\selectfont \textcolor{BrickRed}{$ \blacktriangledown\,1.5\times $}}\end{tabular} & \textcolor{gray}{\ensuremath{\cdot}} & \textcolor{gray}{\ensuremath{\cdot}} & \textcolor{gray}{\ensuremath{\cdot}} & \textcolor{gray}{\ensuremath{\cdot}} & \textcolor{gray}{\ensuremath{\cdot}}\\[3pt]\hline
 & \begin{tabular}[t]{@{}l@{}}\rule{0pt}{12pt}MPF / Two-stack\\RMQ / Block\end{tabular} & \hspace*{1pt}\begin{tabular}[t]{@{}c@{}}\rule{0pt}{12pt}3.8s\\[-.2pt]{\fontsize{5.8}{6.4}\selectfont \textcolor{ForestGreen}{$ \blacktriangle\,2.9\times $}}\\[-.2pt]3s\\[-.2pt]{\fontsize{5.8}{6.4}\selectfont \textcolor{ForestGreen}{$ \blacktriangle\,3.3\times $}}\end{tabular} & \hspace*{1pt}\begin{tabular}[t]{@{}c@{}}\rule{0pt}{12pt}3.6s\\[-.2pt]{\fontsize{5.8}{6.4}\selectfont \textcolor{ForestGreen}{$ \blacktriangle\,3.0\times $}}\\[-.2pt]2.7s\\[-.2pt]{\fontsize{5.8}{6.4}\selectfont \textcolor{ForestGreen}{$ \blacktriangle\,3.2\times $}}\end{tabular} & \hspace*{1pt}\begin{tabular}[t]{@{}c@{}}\rule{0pt}{12pt}6.8s\\[-.2pt]{\fontsize{5.8}{6.4}\selectfont \textcolor{ForestGreen}{$ \blacktriangle\,3.3\times $}}\\[-.2pt]4.7s\\[-.2pt]{\fontsize{5.8}{6.4}\selectfont \textcolor{ForestGreen}{$ \blacktriangle\,4.0\times $}}\end{tabular} & \hspace*{1pt}\begin{tabular}[t]{@{}c@{}}\rule{0pt}{12pt}3.6s\\[-.2pt]{\fontsize{5.8}{6.4}\selectfont \textcolor{ForestGreen}{$ \blacktriangle\,20.1\times $}}\\[-.2pt]2.7s\\[-.2pt]{\fontsize{5.8}{6.4}\selectfont \textcolor{ForestGreen}{$ \blacktriangle\,24.2\times $}}\end{tabular} & \textcolor{gray}{\ensuremath{\cdot}} & \textcolor{gray}{\ensuremath{\cdot}} & \hspace*{1pt}\begin{tabular}[t]{@{}c@{}}\rule{0pt}{12pt}12.8s\\[-.2pt]{\fontsize{5.8}{6.4}\selectfont \textcolor{ForestGreen}{$ \blacktriangle\,23\% $}}\\[-.2pt]8.6s\\[-.2pt]{\fontsize{5.8}{6.4}\selectfont \textcolor{ForestGreen}{$ \blacktriangle\,1.5\times $}}\end{tabular} & \textcolor{gray}{\ensuremath{\cdot}} & \hspace*{1pt}\begin{tabular}[t]{@{}c@{}}\rule{0pt}{12pt}2m\,38.4s\\[-.2pt]{\fontsize{5.8}{6.4}\selectfont \textcolor{BrickRed}{$ \blacktriangledown\,3.0\times $}}\\[-.2pt]2m\,16.3s\\[-.2pt]{\fontsize{5.8}{6.4}\selectfont \textcolor{BrickRed}{$ \blacktriangledown\,2.6\times $}}\end{tabular} & \hspace*{1pt}\begin{tabular}[t]{@{}c@{}}\rule{0pt}{12pt}2m\,38s\\[-.2pt]{\fontsize{5.8}{6.4}\selectfont \textcolor{BrickRed}{$ \blacktriangledown\,2.8\times $}}\\[-.2pt]2m\,17.8s\\[-.2pt]{\fontsize{5.8}{6.4}\selectfont \textcolor{BrickRed}{$ \blacktriangledown\,2.5\times $}}\end{tabular} & \textcolor{gray}{\ensuremath{\cdot}} & \textcolor{gray}{\ensuremath{\cdot}} & \textcolor{gray}{\ensuremath{\cdot}} & \textcolor{gray}{\ensuremath{\cdot}} & \textcolor{gray}{\ensuremath{\cdot}} & \textcolor{gray}{\ensuremath{\cdot}}\\[3pt]\hline
 & \begin{tabular}[t]{@{}l@{}}\rule{0pt}{12pt}MPF / Two-stack\\RMQ / KSW2\end{tabular} & \hspace*{1pt}\begin{tabular}[t]{@{}c@{}}\rule{0pt}{12pt}5.3s\\[-.2pt]{\fontsize{5.8}{6.4}\selectfont \textcolor{ForestGreen}{$ \blacktriangle\,2.1\times $}}\\[-.2pt]2.9s\\[-.2pt]{\fontsize{5.8}{6.4}\selectfont \textcolor{ForestGreen}{$ \blacktriangle\,3.5\times $}}\end{tabular} & \hspace*{1pt}\begin{tabular}[t]{@{}c@{}}\rule{0pt}{12pt}5s\\[-.2pt]{\fontsize{5.8}{6.4}\selectfont \textcolor{ForestGreen}{$ \blacktriangle\,2.1\times $}}\\[-.2pt]2.5s\\[-.2pt]{\fontsize{5.8}{6.4}\selectfont \textcolor{ForestGreen}{$ \blacktriangle\,3.4\times $}}\end{tabular} & \hspace*{1pt}\begin{tabular}[t]{@{}c@{}}\rule{0pt}{12pt}10.1s\\[-.2pt]{\fontsize{5.8}{6.4}\selectfont \textcolor{ForestGreen}{$ \blacktriangle\,2.2\times $}}\\[-.2pt]4.5s\\[-.2pt]{\fontsize{5.8}{6.4}\selectfont \textcolor{ForestGreen}{$ \blacktriangle\,4.2\times $}}\end{tabular} & \hspace*{1pt}\begin{tabular}[t]{@{}c@{}}\rule{0pt}{12pt}5.6s\\[-.2pt]{\fontsize{5.8}{6.4}\selectfont \textcolor{ForestGreen}{$ \blacktriangle\,12.9\times $}}\\[-.2pt]2.6s\\[-.2pt]{\fontsize{5.8}{6.4}\selectfont \textcolor{ForestGreen}{$ \blacktriangle\,25.5\times $}}\end{tabular} & \textcolor{gray}{\ensuremath{\cdot}} & \textcolor{gray}{\ensuremath{\cdot}} & \textcolor{gray}{\ensuremath{\cdot}} & \textcolor{gray}{\ensuremath{\cdot}} & \hspace*{1pt}\begin{tabular}[t]{@{}c@{}}\rule{0pt}{12pt}2m\,37s\\[-.2pt]{\fontsize{5.8}{6.4}\selectfont \textcolor{BrickRed}{$ \blacktriangledown\,3.0\times $}}\\[-.2pt]2m\,12.9s\\[-.2pt]{\fontsize{5.8}{6.4}\selectfont \textcolor{BrickRed}{$ \blacktriangledown\,2.6\times $}}\end{tabular} & \hspace*{1pt}\begin{tabular}[t]{@{}c@{}}\rule{0pt}{12pt}2m\,34.6s\\[-.2pt]{\fontsize{5.8}{6.4}\selectfont \textcolor{BrickRed}{$ \blacktriangledown\,2.7\times $}}\\[-.2pt]2m\,12.9s\\[-.2pt]{\fontsize{5.8}{6.4}\selectfont \textcolor{BrickRed}{$ \blacktriangledown\,2.4\times $}}\end{tabular} & \textcolor{gray}{\ensuremath{\cdot}} & \textcolor{gray}{\ensuremath{\cdot}} & \textcolor{gray}{\ensuremath{\cdot}} & \textcolor{gray}{\ensuremath{\cdot}} & \textcolor{gray}{\ensuremath{\cdot}} & \textcolor{gray}{\ensuremath{\cdot}}\\[3pt]\hline
\end{tabular}\par}
\smallskip{\footnotesize\noindent Each cell: A, comparison, B, comparison. Times cover 1,024 queries. $\cdot$ = no approved observation. ---\textsuperscript{g} = GPU unavailable, not executed. n/a = no measured comparison.\par
\noindent \textsuperscript{f}SIMD-ref request observed as scalar, compared only as that effective path. \textsuperscript{p}Accepted by explicit preservation after wrapper postflight failure, no run receipt claimed.\par
\noindent MPF/MPS = minimizers per fragment/sequence. Row order: index representation / minimizer, then chaining / base alignment. SIMD-ref is the reference minimizer, Bounded is bounded ntHash.\par
\noindent Blue = one lowest A and one lowest B down each machine/workload column, across continuation pages. Exact ties select the first row. Repeated baselines are marked only on page 1.\par
\noindent Light gray = minimap2 baseline. Pale yellow = preferred MPF / Bounded / SIMD-DP / Block configuration, not a claim of winning every column.\par
\noindent When all Vulkan modes are unmeasured, they share one row. Baseline comparisons use the fastest measured BioseqDB effective path separately for A and B.\par}
\clearpage
\noindent{\large\bfseries AWS: 2 Gbp, k=15, w=10 (mapping interval)\hfill 4/7}\par\smallskip
{\fontsize{7.1}{8.1}\selectfont\setlength{\tabcolsep}{1pt}\setlength{\fboxsep}{2pt}\renewcommand{\arraystretch}{1.12}
\noindent\begin{tabular}{@{}p{14mm}p{27mm}|p{13.55mm}|p{13.55mm}|p{13.55mm}|p{13.55mm}|p{13.55mm}|p{13.55mm}|p{13.55mm}|p{13.55mm}|p{13.55mm}|p{13.55mm}|p{13.55mm}|p{13.55mm}|p{13.55mm}|p{13.55mm}|p{13.55mm}|p{13.55mm}|@{}}\hline
Engine & Execution & \multicolumn{8}{c|}{\textbf{Index}} & \multicolumn{8}{c|}{\textbf{No index}}\\
 &  & \multicolumn{8}{c|}{\textbf{2 Gbp}} & \multicolumn{8}{c|}{\textbf{2 Gbp}}\\
 &  & \multicolumn{8}{c|}{\textbf{k=15, w=10}} & \multicolumn{8}{c|}{\textbf{k=15, w=10}}\\
 &  & \multicolumn{4}{c|}{\textbf{DNA4}} & \multicolumn{4}{c|}{\textbf{DNA5}} & \multicolumn{4}{c|}{\textbf{DNA4}} & \multicolumn{4}{c|}{\textbf{DNA5}}\\
 &  & \multicolumn{1}{c|}{\textbf{M7A}} & \multicolumn{1}{c|}{\textbf{M7I}} & \multicolumn{1}{c|}{\textbf{M4}} & \multicolumn{1}{c|}{\textbf{C7G}} & \multicolumn{1}{c|}{\textbf{M7A}} & \multicolumn{1}{c|}{\textbf{M7I}} & \multicolumn{1}{c|}{\textbf{M4}} & \multicolumn{1}{c|}{\textbf{C7G}} & \multicolumn{1}{c|}{\textbf{M7A}} & \multicolumn{1}{c|}{\textbf{M7I}} & \multicolumn{1}{c|}{\textbf{M4}} & \multicolumn{1}{c|}{\textbf{C7G}} & \multicolumn{1}{c|}{\textbf{M7A}} & \multicolumn{1}{c|}{\textbf{M7I}} & \multicolumn{1}{c|}{\textbf{M4}} & \multicolumn{1}{c|}{\textbf{C7G}}\\
\hline
\rowcolor{BenchBaseline}[0.5pt][0.5pt]\textbf{minimap2} & \textbf{Baseline} & \hspace*{1pt}\begin{tabular}[t]{@{}c@{}}\rule{0pt}{12pt}11.2s\\[-.2pt]{\fontsize{5.8}{6.4}\selectfont \textcolor{BrickRed}{$ \blacktriangledown\,3.1\times $}}\\[-.2pt]10s\\[-.2pt]{\fontsize{5.8}{6.4}\selectfont \textcolor{BrickRed}{$ \blacktriangledown\,3.6\times $}}\end{tabular} & \hspace*{1pt}\begin{tabular}[t]{@{}c@{}}\rule{0pt}{12pt}10.7s\\[-.2pt]{\fontsize{5.8}{6.4}\selectfont \textcolor{BrickRed}{$ \blacktriangledown\,3.2\times $}}\\[-.2pt]8.6s\\[-.2pt]{\fontsize{5.8}{6.4}\selectfont \textcolor{BrickRed}{$ \blacktriangledown\,3.5\times $}}\end{tabular} & \hspace*{1pt}\begin{tabular}[t]{@{}c@{}}\rule{0pt}{12pt}22.5s\\[-.2pt]{\fontsize{5.8}{6.4}\selectfont \textcolor{BrickRed}{$ \blacktriangledown\,3.6\times $}}\\[-.2pt]18.8s\\[-.2pt]{\fontsize{5.8}{6.4}\selectfont \textcolor{BrickRed}{$ \blacktriangledown\,4.4\times $}}\end{tabular} & \hspace*{1pt}\begin{tabular}[t]{@{}c@{}}\rule{0pt}{12pt}1m\,13.1s\\[-.2pt]{\fontsize{5.8}{6.4}\selectfont \textcolor{BrickRed}{$ \blacktriangledown\,22.5\times $}}\\[-.2pt]1m\,5.7s\\[-.2pt]{\fontsize{5.8}{6.4}\selectfont \textcolor{BrickRed}{$ \blacktriangledown\,27.7\times $}}\end{tabular} & \hspace*{1pt}\begin{tabular}[t]{@{}c@{}}\rule{0pt}{12pt}9.9s\\[-.2pt]{\fontsize{5.8}{6.4}\selectfont \textcolor{BrickRed}{$ \blacktriangledown\,21\% $}}\\[-.2pt]8.8s\\[-.2pt]{\fontsize{5.8}{6.4}\selectfont \textcolor{BrickRed}{$ \blacktriangledown\,1.5\times $}}\end{tabular} & \hspace*{1pt}\begin{tabular}[t]{@{}c@{}}\rule{0pt}{12pt}9.5s\\[-.2pt]{\fontsize{5.8}{6.4}\selectfont \textcolor{BrickRed}{$ \blacktriangledown\,1.4\times $}}\\[-.2pt]7.6s\\[-.2pt]{\fontsize{5.8}{6.4}\selectfont \textcolor{BrickRed}{$ \blacktriangledown\,1.6\times $}}\end{tabular} & \hspace*{1pt}\begin{tabular}[t]{@{}c@{}}\rule{0pt}{12pt}16.5s\\[-.2pt]{\fontsize{5.8}{6.4}\selectfont \textcolor{BrickRed}{$ \blacktriangledown\,1.7\times $}}\\[-.2pt]12.9s\\[-.2pt]{\fontsize{5.8}{6.4}\selectfont \textcolor{BrickRed}{$ \blacktriangledown\,2.0\times $}}\end{tabular} & \hspace*{1pt}\begin{tabular}[t]{@{}c@{}}\rule{0pt}{12pt}28.2s\\[-.2pt]{\fontsize{5.8}{6.4}\selectfont \textcolor{BrickRed}{$ \blacktriangledown\,2.1\times $}}\\[-.2pt]26.2s\\[-.2pt]{\fontsize{5.8}{6.4}\selectfont \textcolor{BrickRed}{$ \blacktriangledown\,2.4\times $}}\end{tabular} & \hspace*{1pt}\begin{tabular}[t]{@{}c@{}}\rule{0pt}{12pt}53s\\[-.2pt]{\fontsize{5.8}{6.4}\selectfont \textcolor{ForestGreen}{$ \blacktriangle\,2.2\times $}}\\[-.2pt]51.9s\\[-.2pt]{\fontsize{5.8}{6.4}\selectfont \textcolor{ForestGreen}{$ \blacktriangle\,2.1\times $}}\end{tabular} & \hspace*{1pt}\begin{tabular}[t]{@{}c@{}}\rule{0pt}{12pt}56.6s\\[-.2pt]{\fontsize{5.8}{6.4}\selectfont \textcolor{ForestGreen}{$ \blacktriangle\,2.1\times $}}\\[-.2pt]54.6s\\[-.2pt]{\fontsize{5.8}{6.4}\selectfont \textcolor{ForestGreen}{$ \blacktriangle\,1.9\times $}}\end{tabular} & \hspace*{1pt}\begin{tabular}[t]{@{}c@{}}\rule{0pt}{12pt}1m\,55.5s\\[-.2pt]{\fontsize{5.8}{6.4}\selectfont \textcolor{ForestGreen}{$ \blacktriangle\,1.8\times $}}\\[-.2pt]1m\,51.8s\\[-.2pt]{\fontsize{5.8}{6.4}\selectfont \textcolor{ForestGreen}{$ \blacktriangle\,1.5\times $}}\end{tabular} & \hspace*{1pt}\begin{tabular}[t]{@{}c@{}}\rule{0pt}{12pt}6m\,1.3s\\[-.2pt]{\fontsize{5.8}{6.4}\selectfont \textcolor{gray}{n/a}}\\[-.2pt]5m\,55.2s\\[-.2pt]{\fontsize{5.8}{6.4}\selectfont \textcolor{gray}{n/a}}\end{tabular} & \hspace*{1pt}\begin{tabular}[t]{@{}c@{}}\rule{0pt}{12pt}48.1s\\[-.2pt]{\fontsize{5.8}{6.4}\selectfont \textcolor{ForestGreen}{$ \blacktriangle\,7.3\times $}}\\[-.2pt]47s\\[-.2pt]{\fontsize{5.8}{6.4}\selectfont \textcolor{ForestGreen}{$ \blacktriangle\,2.4\times $}}\end{tabular} & \hspace*{1pt}\begin{tabular}[t]{@{}c@{}}\rule{0pt}{12pt}44.5s\\[-.2pt]{\fontsize{5.8}{6.4}\selectfont \textcolor{ForestGreen}{$ \blacktriangle\,6.5\times $}}\\[-.2pt]42.6s\\[-.2pt]{\fontsize{5.8}{6.4}\selectfont \textcolor{ForestGreen}{$ \blacktriangle\,2.7\times $}}\end{tabular} & \hspace*{1pt}\begin{tabular}[t]{@{}c@{}}\rule{0pt}{12pt}1m\,30.2s\\[-.2pt]{\fontsize{5.8}{6.4}\selectfont \textcolor{ForestGreen}{$ \blacktriangle\,6.4\times $}}\\[-.2pt]1m\,26.7s\\[-.2pt]{\fontsize{5.8}{6.4}\selectfont \textcolor{ForestGreen}{$ \blacktriangle\,2.2\times $}}\end{tabular} & \hspace*{1pt}\begin{tabular}[t]{@{}c@{}}\rule{0pt}{12pt}2m\,46.4s\\[-.2pt]{\fontsize{5.8}{6.4}\selectfont \textcolor{gray}{n/a}}\\[-.2pt]2m\,42.7s\\[-.2pt]{\fontsize{5.8}{6.4}\selectfont \textcolor{gray}{n/a}}\end{tabular}\\[3pt]\hline
\textbf{BioseqDB} & \begin{tabular}[t]{@{}l@{}}\rule{0pt}{12pt}MPF / Two-stack\\Scalar-DP / Block\end{tabular} & \hspace*{1pt}\begin{tabular}[t]{@{}c@{}}\rule{0pt}{12pt}3.6s\\[-.2pt]{\fontsize{5.8}{6.4}\selectfont \textcolor{ForestGreen}{$ \blacktriangle\,3.1\times $}}\\[-.2pt]2.9s\\[-.2pt]{\fontsize{5.8}{6.4}\selectfont \textcolor{ForestGreen}{$ \blacktriangle\,3.4\times $}}\end{tabular} & \hspace*{1pt}\begin{tabular}[t]{@{}c@{}}\rule{0pt}{12pt}3.4s\\[-.2pt]{\fontsize{5.8}{6.4}\selectfont \textcolor{ForestGreen}{$ \blacktriangle\,3.2\times $}}\\[-.2pt]2.6s\\[-.2pt]{\fontsize{5.8}{6.4}\selectfont \textcolor{ForestGreen}{$ \blacktriangle\,3.3\times $}}\end{tabular} & \hspace*{1pt}\begin{tabular}[t]{@{}c@{}}\rule{0pt}{12pt}\colorbox{BenchBest}{\strut 6.3s}\\[-.2pt]{\fontsize{5.8}{6.4}\selectfont \textcolor{ForestGreen}{$ \blacktriangle\,3.6\times $}}\\[-.2pt]4.6s\\[-.2pt]{\fontsize{5.8}{6.4}\selectfont \textcolor{ForestGreen}{$ \blacktriangle\,4.1\times $}}\end{tabular} & \hspace*{1pt}\begin{tabular}[t]{@{}c@{}}\rule{0pt}{12pt}7.9s\\[-.2pt]{\fontsize{5.8}{6.4}\selectfont \textcolor{ForestGreen}{$ \blacktriangle\,9.2\times $}}\\[-.2pt]7.1s\\[-.2pt]{\fontsize{5.8}{6.4}\selectfont \textcolor{ForestGreen}{$ \blacktriangle\,9.3\times $}}\end{tabular} & \textcolor{gray}{\ensuremath{\cdot}} & \textcolor{gray}{\ensuremath{\cdot}} & \textcolor{gray}{\ensuremath{\cdot}} & \textcolor{gray}{\ensuremath{\cdot}} & \textcolor{gray}{\ensuremath{\cdot}} & \textcolor{gray}{\ensuremath{\cdot}} & \textcolor{gray}{\ensuremath{\cdot}} & \textcolor{gray}{\ensuremath{\cdot}} & \textcolor{gray}{\ensuremath{\cdot}} & \textcolor{gray}{\ensuremath{\cdot}} & \textcolor{gray}{\ensuremath{\cdot}} & \textcolor{gray}{\ensuremath{\cdot}}\\[3pt]\hline
 & \begin{tabular}[t]{@{}l@{}}\rule{0pt}{12pt}MPF / Two-stack\\Scalar-DP / KSW2\end{tabular} & \hspace*{1pt}\begin{tabular}[t]{@{}c@{}}\rule{0pt}{12pt}5.2s\\[-.2pt]{\fontsize{5.8}{6.4}\selectfont \textcolor{ForestGreen}{$ \blacktriangle\,2.1\times $}}\\[-.2pt]2.8s\\[-.2pt]{\fontsize{5.8}{6.4}\selectfont \textcolor{ForestGreen}{$ \blacktriangle\,3.5\times $}}\end{tabular} & \hspace*{1pt}\begin{tabular}[t]{@{}c@{}}\rule{0pt}{12pt}4.7s\\[-.2pt]{\fontsize{5.8}{6.4}\selectfont \textcolor{ForestGreen}{$ \blacktriangle\,2.3\times $}}\\[-.2pt]2.5s\\[-.2pt]{\fontsize{5.8}{6.4}\selectfont \textcolor{ForestGreen}{$ \blacktriangle\,3.4\times $}}\end{tabular} & \hspace*{1pt}\begin{tabular}[t]{@{}c@{}}\rule{0pt}{12pt}9.5s\\[-.2pt]{\fontsize{5.8}{6.4}\selectfont \textcolor{ForestGreen}{$ \blacktriangle\,2.4\times $}}\\[-.2pt]4.5s\\[-.2pt]{\fontsize{5.8}{6.4}\selectfont \textcolor{ForestGreen}{$ \blacktriangle\,4.2\times $}}\end{tabular} & \hspace*{1pt}\begin{tabular}[t]{@{}c@{}}\rule{0pt}{12pt}7.2s\\[-.2pt]{\fontsize{5.8}{6.4}\selectfont \textcolor{ForestGreen}{$ \blacktriangle\,10.1\times $}}\\[-.2pt]4.4s\\[-.2pt]{\fontsize{5.8}{6.4}\selectfont \textcolor{ForestGreen}{$ \blacktriangle\,14.9\times $}}\end{tabular} & \textcolor{gray}{\ensuremath{\cdot}} & \textcolor{gray}{\ensuremath{\cdot}} & \textcolor{gray}{\ensuremath{\cdot}} & \textcolor{gray}{\ensuremath{\cdot}} & \textcolor{gray}{\ensuremath{\cdot}} & \textcolor{gray}{\ensuremath{\cdot}} & \textcolor{gray}{\ensuremath{\cdot}} & \textcolor{gray}{\ensuremath{\cdot}} & \textcolor{gray}{\ensuremath{\cdot}} & \textcolor{gray}{\ensuremath{\cdot}} & \textcolor{gray}{\ensuremath{\cdot}} & \textcolor{gray}{\ensuremath{\cdot}}\\[3pt]\hline
 & \begin{tabular}[t]{@{}l@{}}\rule{0pt}{12pt}MPF / Two-stack\\SIMD-DP / Block\end{tabular} & \hspace*{1pt}\begin{tabular}[t]{@{}c@{}}\rule{0pt}{12pt}3.7s\\[-.2pt]{\fontsize{5.8}{6.4}\selectfont \textcolor{ForestGreen}{$ \blacktriangle\,3.0\times $}}\\[-.2pt]3s\\[-.2pt]{\fontsize{5.8}{6.4}\selectfont \textcolor{ForestGreen}{$ \blacktriangle\,3.4\times $}}\end{tabular} & \hspace*{1pt}\begin{tabular}[t]{@{}c@{}}\rule{0pt}{12pt}3.5s\\[-.2pt]{\fontsize{5.8}{6.4}\selectfont \textcolor{ForestGreen}{$ \blacktriangle\,3.0\times $}}\\[-.2pt]2.6s\\[-.2pt]{\fontsize{5.8}{6.4}\selectfont \textcolor{ForestGreen}{$ \blacktriangle\,3.3\times $}}\end{tabular} & \hspace*{1pt}\begin{tabular}[t]{@{}c@{}}\rule{0pt}{12pt}6.7s\\[-.2pt]{\fontsize{5.8}{6.4}\selectfont \textcolor{ForestGreen}{$ \blacktriangle\,3.4\times $}}\\[-.2pt]4.7s\\[-.2pt]{\fontsize{5.8}{6.4}\selectfont \textcolor{ForestGreen}{$ \blacktriangle\,4.0\times $}}\end{tabular} & \hspace*{1pt}\begin{tabular}[t]{@{}c@{}}\rule{0pt}{12pt}3.4s\\[-.2pt]{\fontsize{5.8}{6.4}\selectfont \textcolor{ForestGreen}{$ \blacktriangle\,21.5\times $}}\\[-.2pt]2.5s\\[-.2pt]{\fontsize{5.8}{6.4}\selectfont \textcolor{ForestGreen}{$ \blacktriangle\,25.8\times $}}\end{tabular} & \textcolor{gray}{\ensuremath{\cdot}} & \textcolor{gray}{\ensuremath{\cdot}} & \textcolor{gray}{\ensuremath{\cdot}} & \hspace*{1pt}\begin{tabular}[t]{@{}c@{}}\rule{0pt}{12pt}\colorbox{BenchBest}{\strut 13.5s}\\[-.2pt]{\fontsize{5.8}{6.4}\selectfont \textcolor{ForestGreen}{$ \blacktriangle\,2.1\times $}}\\[-.2pt]\colorbox{BenchBest}{\strut 11s}\\[-.2pt]{\fontsize{5.8}{6.4}\selectfont \textcolor{ForestGreen}{$ \blacktriangle\,2.4\times $}}\end{tabular} & \hspace*{1pt}\begin{tabular}[t]{@{}c@{}}\rule{0pt}{12pt}2m\,38.5s\\[-.2pt]{\fontsize{5.8}{6.4}\selectfont \textcolor{BrickRed}{$ \blacktriangledown\,3.0\times $}}\\[-.2pt]2m\,16s\\[-.2pt]{\fontsize{5.8}{6.4}\selectfont \textcolor{BrickRed}{$ \blacktriangledown\,2.6\times $}}\end{tabular} & \hspace*{1pt}\begin{tabular}[t]{@{}c@{}}\rule{0pt}{12pt}2m\,34.1s\\[-.2pt]{\fontsize{5.8}{6.4}\selectfont \textcolor{BrickRed}{$ \blacktriangledown\,2.7\times $}}\\[-.2pt]2m\,13.1s\\[-.2pt]{\fontsize{5.8}{6.4}\selectfont \textcolor{BrickRed}{$ \blacktriangledown\,2.4\times $}}\end{tabular} & \textcolor{gray}{\ensuremath{\cdot}} & \textcolor{gray}{\ensuremath{\cdot}} & \textcolor{gray}{\ensuremath{\cdot}} & \textcolor{gray}{\ensuremath{\cdot}} & \textcolor{gray}{\ensuremath{\cdot}} & \textcolor{gray}{\ensuremath{\cdot}}\\[3pt]\hline
 & \begin{tabular}[t]{@{}l@{}}\rule{0pt}{12pt}MPF / Two-stack\\SIMD-DP / KSW2\end{tabular} & \hspace*{1pt}\begin{tabular}[t]{@{}c@{}}\rule{0pt}{12pt}5.3s\\[-.2pt]{\fontsize{5.8}{6.4}\selectfont \textcolor{ForestGreen}{$ \blacktriangle\,2.1\times $}}\\[-.2pt]2.9s\\[-.2pt]{\fontsize{5.8}{6.4}\selectfont \textcolor{ForestGreen}{$ \blacktriangle\,3.5\times $}}\end{tabular} & \hspace*{1pt}\begin{tabular}[t]{@{}c@{}}\rule{0pt}{12pt}4.8s\\[-.2pt]{\fontsize{5.8}{6.4}\selectfont \textcolor{ForestGreen}{$ \blacktriangle\,2.2\times $}}\\[-.2pt]2.5s\\[-.2pt]{\fontsize{5.8}{6.4}\selectfont \textcolor{ForestGreen}{$ \blacktriangle\,3.4\times $}}\end{tabular} & \hspace*{1pt}\begin{tabular}[t]{@{}c@{}}\rule{0pt}{12pt}9.8s\\[-.2pt]{\fontsize{5.8}{6.4}\selectfont \textcolor{ForestGreen}{$ \blacktriangle\,2.3\times $}}\\[-.2pt]4.5s\\[-.2pt]{\fontsize{5.8}{6.4}\selectfont \textcolor{ForestGreen}{$ \blacktriangle\,4.2\times $}}\end{tabular} & \hspace*{1pt}\begin{tabular}[t]{@{}c@{}}\rule{0pt}{12pt}7.1s\\[-.2pt]{\fontsize{5.8}{6.4}\selectfont \textcolor{ForestGreen}{$ \blacktriangle\,10.3\times $}}\\[-.2pt]4.2s\\[-.2pt]{\fontsize{5.8}{6.4}\selectfont \textcolor{ForestGreen}{$ \blacktriangle\,15.6\times $}}\end{tabular} & \textcolor{gray}{\ensuremath{\cdot}} & \textcolor{gray}{\ensuremath{\cdot}} & \textcolor{gray}{\ensuremath{\cdot}} & \textcolor{gray}{\ensuremath{\cdot}} & \hspace*{1pt}\begin{tabular}[t]{@{}c@{}}\rule{0pt}{12pt}2m\,27.4s\\[-.2pt]{\fontsize{5.8}{6.4}\selectfont \textcolor{BrickRed}{$ \blacktriangledown\,2.8\times $}}\\[-.2pt]2m\,3.1s\\[-.2pt]{\fontsize{5.8}{6.4}\selectfont \textcolor{BrickRed}{$ \blacktriangledown\,2.4\times $}}\end{tabular} & \hspace*{1pt}\begin{tabular}[t]{@{}c@{}}\rule{0pt}{12pt}2m\,26.8s\\[-.2pt]{\fontsize{5.8}{6.4}\selectfont \textcolor{BrickRed}{$ \blacktriangledown\,2.6\times $}}\\[-.2pt]2m\,5.2s\\[-.2pt]{\fontsize{5.8}{6.4}\selectfont \textcolor{BrickRed}{$ \blacktriangledown\,2.3\times $}}\end{tabular} & \textcolor{gray}{\ensuremath{\cdot}} & \textcolor{gray}{\ensuremath{\cdot}} & \textcolor{gray}{\ensuremath{\cdot}} & \textcolor{gray}{\ensuremath{\cdot}} & \textcolor{gray}{\ensuremath{\cdot}} & \textcolor{gray}{\ensuremath{\cdot}}\\[3pt]\hline
 & \begin{tabular}[t]{@{}l@{}}\rule{0pt}{12pt}MPS / Bounded\\RMQ / Block\end{tabular} & \hspace*{1pt}\begin{tabular}[t]{@{}c@{}}\rule{0pt}{12pt}6.5s\\[-.2pt]{\fontsize{5.8}{6.4}\selectfont \textcolor{ForestGreen}{$ \blacktriangle\,1.7\times $}}\\[-.2pt]5.7s\\[-.2pt]{\fontsize{5.8}{6.4}\selectfont \textcolor{ForestGreen}{$ \blacktriangle\,1.8\times $}}\end{tabular} & \hspace*{1pt}\begin{tabular}[t]{@{}c@{}}\rule{0pt}{12pt}6.5s\\[-.2pt]{\fontsize{5.8}{6.4}\selectfont \textcolor{ForestGreen}{$ \blacktriangle\,1.6\times $}}\\[-.2pt]5.8s\\[-.2pt]{\fontsize{5.8}{6.4}\selectfont \textcolor{ForestGreen}{$ \blacktriangle\,1.5\times $}}\end{tabular} & \textcolor{gray}{\ensuremath{\cdot}} & \hspace*{1pt}\begin{tabular}[t]{@{}c@{}}\rule{0pt}{12pt}6.1s\\[-.2pt]{\fontsize{5.8}{6.4}\selectfont \textcolor{ForestGreen}{$ \blacktriangle\,12.0\times $}}\\[-.2pt]5.1s\\[-.2pt]{\fontsize{5.8}{6.4}\selectfont \textcolor{ForestGreen}{$ \blacktriangle\,12.8\times $}}\end{tabular} & \textcolor{gray}{\ensuremath{\cdot}} & \textcolor{gray}{\ensuremath{\cdot}} & \textcolor{gray}{\ensuremath{\cdot}} & \textcolor{gray}{\ensuremath{\cdot}} & \hspace*{1pt}\begin{tabular}[t]{@{}c@{}}\rule{0pt}{12pt}2m\,31.3s\\[-.2pt]{\fontsize{5.8}{6.4}\selectfont \textcolor{BrickRed}{$ \blacktriangledown\,2.9\times $}}\\[-.2pt]2m\,9.3s\\[-.2pt]{\fontsize{5.8}{6.4}\selectfont \textcolor{BrickRed}{$ \blacktriangledown\,2.5\times $}}\end{tabular} & \hspace*{1pt}\begin{tabular}[t]{@{}c@{}}\rule{0pt}{12pt}2m\,16.4s\\[-.2pt]{\fontsize{5.8}{6.4}\selectfont \textcolor{BrickRed}{$ \blacktriangledown\,2.4\times $}}\\[-.2pt]1m\,56.6s\\[-.2pt]{\fontsize{5.8}{6.4}\selectfont \textcolor{BrickRed}{$ \blacktriangledown\,2.1\times $}}\end{tabular} & \textcolor{gray}{\ensuremath{\cdot}} & \textcolor{gray}{\ensuremath{\cdot}} & \textcolor{gray}{\ensuremath{\cdot}} & \textcolor{gray}{\ensuremath{\cdot}} & \textcolor{gray}{\ensuremath{\cdot}} & \textcolor{gray}{\ensuremath{\cdot}}\\[3pt]\hline
 & \begin{tabular}[t]{@{}l@{}}\rule{0pt}{12pt}MPS / Bounded\\RMQ / KSW2\end{tabular} & \hspace*{1pt}\begin{tabular}[t]{@{}c@{}}\rule{0pt}{12pt}7.9s\\[-.2pt]{\fontsize{5.8}{6.4}\selectfont \textcolor{ForestGreen}{$ \blacktriangle\,29\% $}}\\[-.2pt]4.9s\\[-.2pt]{\fontsize{5.8}{6.4}\selectfont \textcolor{ForestGreen}{$ \blacktriangle\,2.0\times $}}\end{tabular} & \hspace*{1pt}\begin{tabular}[t]{@{}c@{}}\rule{0pt}{12pt}8.5s\\[-.2pt]{\fontsize{5.8}{6.4}\selectfont \textcolor{ForestGreen}{$ \blacktriangle\,20\% $}}\\[-.2pt]4.9s\\[-.2pt]{\fontsize{5.8}{6.4}\selectfont \textcolor{ForestGreen}{$ \blacktriangle\,1.7\times $}}\end{tabular} & \textcolor{gray}{\ensuremath{\cdot}} & \hspace*{1pt}\begin{tabular}[t]{@{}c@{}}\rule{0pt}{12pt}8.1s\\[-.2pt]{\fontsize{5.8}{6.4}\selectfont \textcolor{ForestGreen}{$ \blacktriangle\,9.1\times $}}\\[-.2pt]4.6s\\[-.2pt]{\fontsize{5.8}{6.4}\selectfont \textcolor{ForestGreen}{$ \blacktriangle\,14.2\times $}}\end{tabular} & \textcolor{gray}{\ensuremath{\cdot}} & \textcolor{gray}{\ensuremath{\cdot}} & \textcolor{gray}{\ensuremath{\cdot}} & \textcolor{gray}{\ensuremath{\cdot}} & \hspace*{1pt}\begin{tabular}[t]{@{}c@{}}\rule{0pt}{12pt}2m\,21.9s\\[-.2pt]{\fontsize{5.8}{6.4}\selectfont \textcolor{BrickRed}{$ \blacktriangledown\,2.7\times $}}\\[-.2pt]1m\,57.6s\\[-.2pt]{\fontsize{5.8}{6.4}\selectfont \textcolor{BrickRed}{$ \blacktriangledown\,2.3\times $}}\end{tabular} & \hspace*{1pt}\begin{tabular}[t]{@{}c@{}}\rule{0pt}{12pt}2m\,27.9s\\[-.2pt]{\fontsize{5.8}{6.4}\selectfont \textcolor{BrickRed}{$ \blacktriangledown\,2.6\times $}}\\[-.2pt]2m\,5.4s\\[-.2pt]{\fontsize{5.8}{6.4}\selectfont \textcolor{BrickRed}{$ \blacktriangledown\,2.3\times $}}\end{tabular} & \textcolor{gray}{\ensuremath{\cdot}} & \textcolor{gray}{\ensuremath{\cdot}} & \textcolor{gray}{\ensuremath{\cdot}} & \textcolor{gray}{\ensuremath{\cdot}} & \textcolor{gray}{\ensuremath{\cdot}} & \textcolor{gray}{\ensuremath{\cdot}}\\[3pt]\hline
\end{tabular}\par}
\smallskip{\footnotesize\noindent Each cell: A, comparison, B, comparison. Times cover 1,024 queries. $\cdot$ = no approved observation. ---\textsuperscript{g} = GPU unavailable, not executed. n/a = no measured comparison.\par
\noindent \textsuperscript{f}SIMD-ref request observed as scalar, compared only as that effective path. \textsuperscript{p}Accepted by explicit preservation after wrapper postflight failure, no run receipt claimed.\par
\noindent MPF/MPS = minimizers per fragment/sequence. Row order: index representation / minimizer, then chaining / base alignment. SIMD-ref is the reference minimizer, Bounded is bounded ntHash.\par
\noindent Blue = one lowest A and one lowest B down each machine/workload column, across continuation pages. Exact ties select the first row. Repeated baselines are marked only on page 1.\par
\noindent Light gray = minimap2 baseline. Pale yellow = preferred MPF / Bounded / SIMD-DP / Block configuration, not a claim of winning every column.\par
\noindent When all Vulkan modes are unmeasured, they share one row. Baseline comparisons use the fastest measured BioseqDB effective path separately for A and B.\par}
\clearpage
\noindent{\large\bfseries AWS: 2 Gbp, k=15, w=10 (mapping interval)\hfill 5/7}\par\smallskip
{\fontsize{7.1}{8.1}\selectfont\setlength{\tabcolsep}{1pt}\setlength{\fboxsep}{2pt}\renewcommand{\arraystretch}{1.12}
\noindent\begin{tabular}{@{}p{14mm}p{27mm}|p{13.55mm}|p{13.55mm}|p{13.55mm}|p{13.55mm}|p{13.55mm}|p{13.55mm}|p{13.55mm}|p{13.55mm}|p{13.55mm}|p{13.55mm}|p{13.55mm}|p{13.55mm}|p{13.55mm}|p{13.55mm}|p{13.55mm}|p{13.55mm}|@{}}\hline
Engine & Execution & \multicolumn{8}{c|}{\textbf{Index}} & \multicolumn{8}{c|}{\textbf{No index}}\\
 &  & \multicolumn{8}{c|}{\textbf{2 Gbp}} & \multicolumn{8}{c|}{\textbf{2 Gbp}}\\
 &  & \multicolumn{8}{c|}{\textbf{k=15, w=10}} & \multicolumn{8}{c|}{\textbf{k=15, w=10}}\\
 &  & \multicolumn{4}{c|}{\textbf{DNA4}} & \multicolumn{4}{c|}{\textbf{DNA5}} & \multicolumn{4}{c|}{\textbf{DNA4}} & \multicolumn{4}{c|}{\textbf{DNA5}}\\
 &  & \multicolumn{1}{c|}{\textbf{M7A}} & \multicolumn{1}{c|}{\textbf{M7I}} & \multicolumn{1}{c|}{\textbf{M4}} & \multicolumn{1}{c|}{\textbf{C7G}} & \multicolumn{1}{c|}{\textbf{M7A}} & \multicolumn{1}{c|}{\textbf{M7I}} & \multicolumn{1}{c|}{\textbf{M4}} & \multicolumn{1}{c|}{\textbf{C7G}} & \multicolumn{1}{c|}{\textbf{M7A}} & \multicolumn{1}{c|}{\textbf{M7I}} & \multicolumn{1}{c|}{\textbf{M4}} & \multicolumn{1}{c|}{\textbf{C7G}} & \multicolumn{1}{c|}{\textbf{M7A}} & \multicolumn{1}{c|}{\textbf{M7I}} & \multicolumn{1}{c|}{\textbf{M4}} & \multicolumn{1}{c|}{\textbf{C7G}}\\
\hline
\rowcolor{BenchBaseline}[0.5pt][0.5pt]\textbf{minimap2} & \textbf{Baseline} & \hspace*{1pt}\begin{tabular}[t]{@{}c@{}}\rule{0pt}{12pt}11.2s\\[-.2pt]{\fontsize{5.8}{6.4}\selectfont \textcolor{BrickRed}{$ \blacktriangledown\,3.1\times $}}\\[-.2pt]10s\\[-.2pt]{\fontsize{5.8}{6.4}\selectfont \textcolor{BrickRed}{$ \blacktriangledown\,3.6\times $}}\end{tabular} & \hspace*{1pt}\begin{tabular}[t]{@{}c@{}}\rule{0pt}{12pt}10.7s\\[-.2pt]{\fontsize{5.8}{6.4}\selectfont \textcolor{BrickRed}{$ \blacktriangledown\,3.2\times $}}\\[-.2pt]8.6s\\[-.2pt]{\fontsize{5.8}{6.4}\selectfont \textcolor{BrickRed}{$ \blacktriangledown\,3.5\times $}}\end{tabular} & \hspace*{1pt}\begin{tabular}[t]{@{}c@{}}\rule{0pt}{12pt}22.5s\\[-.2pt]{\fontsize{5.8}{6.4}\selectfont \textcolor{BrickRed}{$ \blacktriangledown\,3.6\times $}}\\[-.2pt]18.8s\\[-.2pt]{\fontsize{5.8}{6.4}\selectfont \textcolor{BrickRed}{$ \blacktriangledown\,4.4\times $}}\end{tabular} & \hspace*{1pt}\begin{tabular}[t]{@{}c@{}}\rule{0pt}{12pt}1m\,13.1s\\[-.2pt]{\fontsize{5.8}{6.4}\selectfont \textcolor{BrickRed}{$ \blacktriangledown\,22.5\times $}}\\[-.2pt]1m\,5.7s\\[-.2pt]{\fontsize{5.8}{6.4}\selectfont \textcolor{BrickRed}{$ \blacktriangledown\,27.7\times $}}\end{tabular} & \hspace*{1pt}\begin{tabular}[t]{@{}c@{}}\rule{0pt}{12pt}9.9s\\[-.2pt]{\fontsize{5.8}{6.4}\selectfont \textcolor{BrickRed}{$ \blacktriangledown\,21\% $}}\\[-.2pt]8.8s\\[-.2pt]{\fontsize{5.8}{6.4}\selectfont \textcolor{BrickRed}{$ \blacktriangledown\,1.5\times $}}\end{tabular} & \hspace*{1pt}\begin{tabular}[t]{@{}c@{}}\rule{0pt}{12pt}9.5s\\[-.2pt]{\fontsize{5.8}{6.4}\selectfont \textcolor{BrickRed}{$ \blacktriangledown\,1.4\times $}}\\[-.2pt]7.6s\\[-.2pt]{\fontsize{5.8}{6.4}\selectfont \textcolor{BrickRed}{$ \blacktriangledown\,1.6\times $}}\end{tabular} & \hspace*{1pt}\begin{tabular}[t]{@{}c@{}}\rule{0pt}{12pt}16.5s\\[-.2pt]{\fontsize{5.8}{6.4}\selectfont \textcolor{BrickRed}{$ \blacktriangledown\,1.7\times $}}\\[-.2pt]12.9s\\[-.2pt]{\fontsize{5.8}{6.4}\selectfont \textcolor{BrickRed}{$ \blacktriangledown\,2.0\times $}}\end{tabular} & \hspace*{1pt}\begin{tabular}[t]{@{}c@{}}\rule{0pt}{12pt}28.2s\\[-.2pt]{\fontsize{5.8}{6.4}\selectfont \textcolor{BrickRed}{$ \blacktriangledown\,2.1\times $}}\\[-.2pt]26.2s\\[-.2pt]{\fontsize{5.8}{6.4}\selectfont \textcolor{BrickRed}{$ \blacktriangledown\,2.4\times $}}\end{tabular} & \hspace*{1pt}\begin{tabular}[t]{@{}c@{}}\rule{0pt}{12pt}53s\\[-.2pt]{\fontsize{5.8}{6.4}\selectfont \textcolor{ForestGreen}{$ \blacktriangle\,2.2\times $}}\\[-.2pt]51.9s\\[-.2pt]{\fontsize{5.8}{6.4}\selectfont \textcolor{ForestGreen}{$ \blacktriangle\,2.1\times $}}\end{tabular} & \hspace*{1pt}\begin{tabular}[t]{@{}c@{}}\rule{0pt}{12pt}56.6s\\[-.2pt]{\fontsize{5.8}{6.4}\selectfont \textcolor{ForestGreen}{$ \blacktriangle\,2.1\times $}}\\[-.2pt]54.6s\\[-.2pt]{\fontsize{5.8}{6.4}\selectfont \textcolor{ForestGreen}{$ \blacktriangle\,1.9\times $}}\end{tabular} & \hspace*{1pt}\begin{tabular}[t]{@{}c@{}}\rule{0pt}{12pt}1m\,55.5s\\[-.2pt]{\fontsize{5.8}{6.4}\selectfont \textcolor{ForestGreen}{$ \blacktriangle\,1.8\times $}}\\[-.2pt]1m\,51.8s\\[-.2pt]{\fontsize{5.8}{6.4}\selectfont \textcolor{ForestGreen}{$ \blacktriangle\,1.5\times $}}\end{tabular} & \hspace*{1pt}\begin{tabular}[t]{@{}c@{}}\rule{0pt}{12pt}6m\,1.3s\\[-.2pt]{\fontsize{5.8}{6.4}\selectfont \textcolor{gray}{n/a}}\\[-.2pt]5m\,55.2s\\[-.2pt]{\fontsize{5.8}{6.4}\selectfont \textcolor{gray}{n/a}}\end{tabular} & \hspace*{1pt}\begin{tabular}[t]{@{}c@{}}\rule{0pt}{12pt}48.1s\\[-.2pt]{\fontsize{5.8}{6.4}\selectfont \textcolor{ForestGreen}{$ \blacktriangle\,7.3\times $}}\\[-.2pt]47s\\[-.2pt]{\fontsize{5.8}{6.4}\selectfont \textcolor{ForestGreen}{$ \blacktriangle\,2.4\times $}}\end{tabular} & \hspace*{1pt}\begin{tabular}[t]{@{}c@{}}\rule{0pt}{12pt}44.5s\\[-.2pt]{\fontsize{5.8}{6.4}\selectfont \textcolor{ForestGreen}{$ \blacktriangle\,6.5\times $}}\\[-.2pt]42.6s\\[-.2pt]{\fontsize{5.8}{6.4}\selectfont \textcolor{ForestGreen}{$ \blacktriangle\,2.7\times $}}\end{tabular} & \hspace*{1pt}\begin{tabular}[t]{@{}c@{}}\rule{0pt}{12pt}1m\,30.2s\\[-.2pt]{\fontsize{5.8}{6.4}\selectfont \textcolor{ForestGreen}{$ \blacktriangle\,6.4\times $}}\\[-.2pt]1m\,26.7s\\[-.2pt]{\fontsize{5.8}{6.4}\selectfont \textcolor{ForestGreen}{$ \blacktriangle\,2.2\times $}}\end{tabular} & \hspace*{1pt}\begin{tabular}[t]{@{}c@{}}\rule{0pt}{12pt}2m\,46.4s\\[-.2pt]{\fontsize{5.8}{6.4}\selectfont \textcolor{gray}{n/a}}\\[-.2pt]2m\,42.7s\\[-.2pt]{\fontsize{5.8}{6.4}\selectfont \textcolor{gray}{n/a}}\end{tabular}\\[3pt]\hline
\textbf{BioseqDB} & \begin{tabular}[t]{@{}l@{}}\rule{0pt}{12pt}MPS / Bounded\\Scalar-DP / Block\end{tabular} & \hspace*{1pt}\begin{tabular}[t]{@{}c@{}}\rule{0pt}{12pt}6.4s\\[-.2pt]{\fontsize{5.8}{6.4}\selectfont \textcolor{ForestGreen}{$ \blacktriangle\,1.7\times $}}\\[-.2pt]5.7s\\[-.2pt]{\fontsize{5.8}{6.4}\selectfont \textcolor{ForestGreen}{$ \blacktriangle\,1.8\times $}}\end{tabular} & \hspace*{1pt}\begin{tabular}[t]{@{}c@{}}\rule{0pt}{12pt}6.6s\\[-.2pt]{\fontsize{5.8}{6.4}\selectfont \textcolor{ForestGreen}{$ \blacktriangle\,1.6\times $}}\\[-.2pt]5.7s\\[-.2pt]{\fontsize{5.8}{6.4}\selectfont \textcolor{ForestGreen}{$ \blacktriangle\,1.5\times $}}\end{tabular} & \textcolor{gray}{\ensuremath{\cdot}} & \hspace*{1pt}\begin{tabular}[t]{@{}c@{}}\rule{0pt}{12pt}5.9s\\[-.2pt]{\fontsize{5.8}{6.4}\selectfont \textcolor{ForestGreen}{$ \blacktriangle\,12.3\times $}}\\[-.2pt]4.9s\\[-.2pt]{\fontsize{5.8}{6.4}\selectfont \textcolor{ForestGreen}{$ \blacktriangle\,13.4\times $}}\end{tabular} & \textcolor{gray}{\ensuremath{\cdot}} & \textcolor{gray}{\ensuremath{\cdot}} & \textcolor{gray}{\ensuremath{\cdot}} & \textcolor{gray}{\ensuremath{\cdot}} & \hspace*{1pt}\begin{tabular}[t]{@{}c@{}}\rule{0pt}{12pt}2m\,24.9s\\[-.2pt]{\fontsize{5.8}{6.4}\selectfont \textcolor{BrickRed}{$ \blacktriangledown\,2.7\times $}}\\[-.2pt]2m\,2.4s\\[-.2pt]{\fontsize{5.8}{6.4}\selectfont \textcolor{BrickRed}{$ \blacktriangledown\,2.4\times $}}\end{tabular} & \hspace*{1pt}\begin{tabular}[t]{@{}c@{}}\rule{0pt}{12pt}2m\,39.7s\\[-.2pt]{\fontsize{5.8}{6.4}\selectfont \textcolor{BrickRed}{$ \blacktriangledown\,2.8\times $}}\\[-.2pt]2m\,18.8s\\[-.2pt]{\fontsize{5.8}{6.4}\selectfont \textcolor{BrickRed}{$ \blacktriangledown\,2.5\times $}}\end{tabular} & \textcolor{gray}{\ensuremath{\cdot}} & \textcolor{gray}{\ensuremath{\cdot}} & \textcolor{gray}{\ensuremath{\cdot}} & \textcolor{gray}{\ensuremath{\cdot}} & \textcolor{gray}{\ensuremath{\cdot}} & \textcolor{gray}{\ensuremath{\cdot}}\\[3pt]\hline
 & \begin{tabular}[t]{@{}l@{}}\rule{0pt}{12pt}MPS / Bounded\\Scalar-DP / KSW2\end{tabular} & \hspace*{1pt}\begin{tabular}[t]{@{}c@{}}\rule{0pt}{12pt}7.9s\\[-.2pt]{\fontsize{5.8}{6.4}\selectfont \textcolor{ForestGreen}{$ \blacktriangle\,29\% $}}\\[-.2pt]4.9s\\[-.2pt]{\fontsize{5.8}{6.4}\selectfont \textcolor{ForestGreen}{$ \blacktriangle\,2.0\times $}}\end{tabular} & \hspace*{1pt}\begin{tabular}[t]{@{}c@{}}\rule{0pt}{12pt}8.6s\\[-.2pt]{\fontsize{5.8}{6.4}\selectfont \textcolor{ForestGreen}{$ \blacktriangle\,20\% $}}\\[-.2pt]5s\\[-.2pt]{\fontsize{5.8}{6.4}\selectfont \textcolor{ForestGreen}{$ \blacktriangle\,1.7\times $}}\end{tabular} & \textcolor{gray}{\ensuremath{\cdot}} & \hspace*{1pt}\begin{tabular}[t]{@{}c@{}}\rule{0pt}{12pt}8s\\[-.2pt]{\fontsize{5.8}{6.4}\selectfont \textcolor{ForestGreen}{$ \blacktriangle\,9.2\times $}}\\[-.2pt]4.6s\\[-.2pt]{\fontsize{5.8}{6.4}\selectfont \textcolor{ForestGreen}{$ \blacktriangle\,14.2\times $}}\end{tabular} & \textcolor{gray}{\ensuremath{\cdot}} & \textcolor{gray}{\ensuremath{\cdot}} & \textcolor{gray}{\ensuremath{\cdot}} & \textcolor{gray}{\ensuremath{\cdot}} & \hspace*{1pt}\begin{tabular}[t]{@{}c@{}}\rule{0pt}{12pt}2m\,37s\\[-.2pt]{\fontsize{5.8}{6.4}\selectfont \textcolor{BrickRed}{$ \blacktriangledown\,3.0\times $}}\\[-.2pt]2m\,12s\\[-.2pt]{\fontsize{5.8}{6.4}\selectfont \textcolor{BrickRed}{$ \blacktriangledown\,2.5\times $}}\end{tabular} & \hspace*{1pt}\begin{tabular}[t]{@{}c@{}}\rule{0pt}{12pt}2m\,30.6s\\[-.2pt]{\fontsize{5.8}{6.4}\selectfont \textcolor{BrickRed}{$ \blacktriangledown\,2.7\times $}}\\[-.2pt]2m\,7.1s\\[-.2pt]{\fontsize{5.8}{6.4}\selectfont \textcolor{BrickRed}{$ \blacktriangledown\,2.3\times $}}\end{tabular} & \textcolor{gray}{\ensuremath{\cdot}} & \textcolor{gray}{\ensuremath{\cdot}} & \textcolor{gray}{\ensuremath{\cdot}} & \textcolor{gray}{\ensuremath{\cdot}} & \textcolor{gray}{\ensuremath{\cdot}} & \textcolor{gray}{\ensuremath{\cdot}}\\[3pt]\hline
 & \begin{tabular}[t]{@{}l@{}}\rule{0pt}{12pt}MPS / Bounded\\SIMD-DP / Block\end{tabular} & \hspace*{1pt}\begin{tabular}[t]{@{}c@{}}\rule{0pt}{12pt}6.7s\\[-.2pt]{\fontsize{5.8}{6.4}\selectfont \textcolor{ForestGreen}{$ \blacktriangle\,1.7\times $}}\\[-.2pt]6s\\[-.2pt]{\fontsize{5.8}{6.4}\selectfont \textcolor{ForestGreen}{$ \blacktriangle\,1.7\times $}}\end{tabular} & \hspace*{1pt}\begin{tabular}[t]{@{}c@{}}\rule{0pt}{12pt}6.9s\\[-.2pt]{\fontsize{5.8}{6.4}\selectfont \textcolor{ForestGreen}{$ \blacktriangle\,1.5\times $}}\\[-.2pt]6s\\[-.2pt]{\fontsize{5.8}{6.4}\selectfont \textcolor{ForestGreen}{$ \blacktriangle\,1.4\times $}}\end{tabular} & \hspace*{1pt}\begin{tabular}[t]{@{}c@{}}\rule{0pt}{12pt}12.5s\\[-.2pt]{\fontsize{5.8}{6.4}\selectfont \textcolor{ForestGreen}{$ \blacktriangle\,1.8\times $}}\\[-.2pt]10.4s\\[-.2pt]{\fontsize{5.8}{6.4}\selectfont \textcolor{ForestGreen}{$ \blacktriangle\,1.8\times $}}\end{tabular} & \hspace*{1pt}\begin{tabular}[t]{@{}c@{}}\rule{0pt}{12pt}6s\\[-.2pt]{\fontsize{5.8}{6.4}\selectfont \textcolor{ForestGreen}{$ \blacktriangle\,12.1\times $}}\\[-.2pt]5s\\[-.2pt]{\fontsize{5.8}{6.4}\selectfont \textcolor{ForestGreen}{$ \blacktriangle\,13.0\times $}}\end{tabular} & \textcolor{gray}{\ensuremath{\cdot}} & \textcolor{gray}{\ensuremath{\cdot}} & \textcolor{gray}{\ensuremath{\cdot}} & \textcolor{gray}{\ensuremath{\cdot}} & \hspace*{1pt}\begin{tabular}[t]{@{}c@{}}\rule{0pt}{12pt}2m\,20.8s\\[-.2pt]{\fontsize{5.8}{6.4}\selectfont \textcolor{BrickRed}{$ \blacktriangledown\,2.7\times $}}\\[-.2pt]1m\,57.4s\\[-.2pt]{\fontsize{5.8}{6.4}\selectfont \textcolor{BrickRed}{$ \blacktriangledown\,2.3\times $}}\end{tabular} & \hspace*{1pt}\begin{tabular}[t]{@{}c@{}}\rule{0pt}{12pt}2m\,25.9s\\[-.2pt]{\fontsize{5.8}{6.4}\selectfont \textcolor{BrickRed}{$ \blacktriangledown\,2.6\times $}}\\[-.2pt]2m\,5.5s\\[-.2pt]{\fontsize{5.8}{6.4}\selectfont \textcolor{BrickRed}{$ \blacktriangledown\,2.3\times $}}\end{tabular} & \textcolor{gray}{\ensuremath{\cdot}} & \textcolor{gray}{\ensuremath{\cdot}} & \textcolor{gray}{\ensuremath{\cdot}} & \textcolor{gray}{\ensuremath{\cdot}} & \textcolor{gray}{\ensuremath{\cdot}} & \textcolor{gray}{\ensuremath{\cdot}}\\[3pt]\hline
 & \begin{tabular}[t]{@{}l@{}}\rule{0pt}{12pt}MPS / Bounded\\SIMD-DP / KSW2\end{tabular} & \hspace*{1pt}\begin{tabular}[t]{@{}c@{}}\rule{0pt}{12pt}8.2s\\[-.2pt]{\fontsize{5.8}{6.4}\selectfont \textcolor{ForestGreen}{$ \blacktriangle\,26\% $}}\\[-.2pt]5.4s\\[-.2pt]{\fontsize{5.8}{6.4}\selectfont \textcolor{ForestGreen}{$ \blacktriangle\,1.8\times $}}\end{tabular} & \hspace*{1pt}\begin{tabular}[t]{@{}c@{}}\rule{0pt}{12pt}8.4s\\[-.2pt]{\fontsize{5.8}{6.4}\selectfont \textcolor{ForestGreen}{$ \blacktriangle\,21\% $}}\\[-.2pt]4.9s\\[-.2pt]{\fontsize{5.8}{6.4}\selectfont \textcolor{ForestGreen}{$ \blacktriangle\,1.8\times $}}\end{tabular} & \textcolor{gray}{\ensuremath{\cdot}} & \hspace*{1pt}\begin{tabular}[t]{@{}c@{}}\rule{0pt}{12pt}9.6s\\[-.2pt]{\fontsize{5.8}{6.4}\selectfont \textcolor{ForestGreen}{$ \blacktriangle\,7.6\times $}}\\[-.2pt]6s\\[-.2pt]{\fontsize{5.8}{6.4}\selectfont \textcolor{ForestGreen}{$ \blacktriangle\,10.9\times $}}\end{tabular} & \textcolor{gray}{\ensuremath{\cdot}} & \textcolor{gray}{\ensuremath{\cdot}} & \textcolor{gray}{\ensuremath{\cdot}} & \textcolor{gray}{\ensuremath{\cdot}} & \hspace*{1pt}\begin{tabular}[t]{@{}c@{}}\rule{0pt}{12pt}2m\,20.7s\\[-.2pt]{\fontsize{5.8}{6.4}\selectfont \textcolor{BrickRed}{$ \blacktriangledown\,2.7\times $}}\\[-.2pt]1m\,56.8s\\[-.2pt]{\fontsize{5.8}{6.4}\selectfont \textcolor{BrickRed}{$ \blacktriangledown\,2.3\times $}}\end{tabular} & \hspace*{1pt}\begin{tabular}[t]{@{}c@{}}\rule{0pt}{12pt}2m\,19.7s\\[-.2pt]{\fontsize{5.8}{6.4}\selectfont \textcolor{BrickRed}{$ \blacktriangledown\,2.5\times $}}\\[-.2pt]1m\,57.9s\\[-.2pt]{\fontsize{5.8}{6.4}\selectfont \textcolor{BrickRed}{$ \blacktriangledown\,2.2\times $}}\end{tabular} & \textcolor{gray}{\ensuremath{\cdot}} & \textcolor{gray}{\ensuremath{\cdot}} & \textcolor{gray}{\ensuremath{\cdot}} & \textcolor{gray}{\ensuremath{\cdot}} & \textcolor{gray}{\ensuremath{\cdot}} & \textcolor{gray}{\ensuremath{\cdot}}\\[3pt]\hline
 & \begin{tabular}[t]{@{}l@{}}\rule{0pt}{12pt}MPS / Scalar\\RMQ / Block\end{tabular} & \hspace*{1pt}\begin{tabular}[t]{@{}c@{}}\rule{0pt}{12pt}6.9s\\[-.2pt]{\fontsize{5.8}{6.4}\selectfont \textcolor{ForestGreen}{$ \blacktriangle\,1.6\times $}}\\[-.2pt]6.2s\\[-.2pt]{\fontsize{5.8}{6.4}\selectfont \textcolor{ForestGreen}{$ \blacktriangle\,1.6\times $}}\end{tabular} & \hspace*{1pt}\begin{tabular}[t]{@{}c@{}}\rule{0pt}{12pt}7s\\[-.2pt]{\fontsize{5.8}{6.4}\selectfont \textcolor{ForestGreen}{$ \blacktriangle\,1.5\times $}}\\[-.2pt]6.2s\\[-.2pt]{\fontsize{5.8}{6.4}\selectfont \textcolor{ForestGreen}{$ \blacktriangle\,28\% $}}\end{tabular} & \textcolor{gray}{\ensuremath{\cdot}} & \textcolor{gray}{\ensuremath{\cdot}} & \textcolor{gray}{\ensuremath{\cdot}} & \textcolor{gray}{\ensuremath{\cdot}} & \textcolor{gray}{\ensuremath{\cdot}} & \textcolor{gray}{\ensuremath{\cdot}} & \hspace*{1pt}\begin{tabular}[t]{@{}c@{}}\rule{0pt}{12pt}3m\,15.8s\\[-.2pt]{\fontsize{5.8}{6.4}\selectfont \textcolor{BrickRed}{$ \blacktriangledown\,3.7\times $}}\\[-.2pt]2m\,53.2s\\[-.2pt]{\fontsize{5.8}{6.4}\selectfont \textcolor{BrickRed}{$ \blacktriangledown\,3.3\times $}}\end{tabular} & \hspace*{1pt}\begin{tabular}[t]{@{}c@{}}\rule{0pt}{12pt}3m\,13.2s\\[-.2pt]{\fontsize{5.8}{6.4}\selectfont \textcolor{BrickRed}{$ \blacktriangledown\,3.4\times $}}\\[-.2pt]2m\,53s\\[-.2pt]{\fontsize{5.8}{6.4}\selectfont \textcolor{BrickRed}{$ \blacktriangledown\,3.2\times $}}\end{tabular} & \textcolor{gray}{\ensuremath{\cdot}} & \textcolor{gray}{\ensuremath{\cdot}} & \textcolor{gray}{\ensuremath{\cdot}} & \textcolor{gray}{\ensuremath{\cdot}} & \textcolor{gray}{\ensuremath{\cdot}} & \textcolor{gray}{\ensuremath{\cdot}}\\[3pt]\hline
 & \begin{tabular}[t]{@{}l@{}}\rule{0pt}{12pt}MPS / Scalar\\RMQ / KSW2\end{tabular} & \hspace*{1pt}\begin{tabular}[t]{@{}c@{}}\rule{0pt}{12pt}8.2s\\[-.2pt]{\fontsize{5.8}{6.4}\selectfont \textcolor{ForestGreen}{$ \blacktriangle\,26\% $}}\\[-.2pt]5.4s\\[-.2pt]{\fontsize{5.8}{6.4}\selectfont \textcolor{ForestGreen}{$ \blacktriangle\,1.9\times $}}\end{tabular} & \hspace*{1pt}\begin{tabular}[t]{@{}c@{}}\rule{0pt}{12pt}8.7s\\[-.2pt]{\fontsize{5.8}{6.4}\selectfont \textcolor{ForestGreen}{$ \blacktriangle\,18\% $}}\\[-.2pt]5.4s\\[-.2pt]{\fontsize{5.8}{6.4}\selectfont \textcolor{ForestGreen}{$ \blacktriangle\,1.6\times $}}\end{tabular} & \textcolor{gray}{\ensuremath{\cdot}} & \textcolor{gray}{\ensuremath{\cdot}} & \textcolor{gray}{\ensuremath{\cdot}} & \textcolor{gray}{\ensuremath{\cdot}} & \textcolor{gray}{\ensuremath{\cdot}} & \textcolor{gray}{\ensuremath{\cdot}} & \hspace*{1pt}\begin{tabular}[t]{@{}c@{}}\rule{0pt}{12pt}3m\,14.2s\\[-.2pt]{\fontsize{5.8}{6.4}\selectfont \textcolor{BrickRed}{$ \blacktriangledown\,3.7\times $}}\\[-.2pt]2m\,49.7s\\[-.2pt]{\fontsize{5.8}{6.4}\selectfont \textcolor{BrickRed}{$ \blacktriangledown\,3.3\times $}}\end{tabular} & \textcolor{gray}{\ensuremath{\cdot}} & \textcolor{gray}{\ensuremath{\cdot}} & \textcolor{gray}{\ensuremath{\cdot}} & \textcolor{gray}{\ensuremath{\cdot}} & \textcolor{gray}{\ensuremath{\cdot}} & \textcolor{gray}{\ensuremath{\cdot}} & \textcolor{gray}{\ensuremath{\cdot}}\\[3pt]\hline
 & \begin{tabular}[t]{@{}l@{}}\rule{0pt}{12pt}MPS / Scalar\\Scalar-DP / Block\end{tabular} & \hspace*{1pt}\begin{tabular}[t]{@{}c@{}}\rule{0pt}{12pt}7.5s\\[-.2pt]{\fontsize{5.8}{6.4}\selectfont \textcolor{ForestGreen}{$ \blacktriangle\,1.5\times $}}\\[-.2pt]6.3s\\[-.2pt]{\fontsize{5.8}{6.4}\selectfont \textcolor{ForestGreen}{$ \blacktriangle\,1.6\times $}}\end{tabular} & \hspace*{1pt}\begin{tabular}[t]{@{}c@{}}\rule{0pt}{12pt}7.2s\\[-.2pt]{\fontsize{5.8}{6.4}\selectfont \textcolor{ForestGreen}{$ \blacktriangle\,1.5\times $}}\\[-.2pt]6.4s\\[-.2pt]{\fontsize{5.8}{6.4}\selectfont \textcolor{ForestGreen}{$ \blacktriangle\,26\% $}}\end{tabular} & \textcolor{gray}{\ensuremath{\cdot}} & \textcolor{gray}{\ensuremath{\cdot}} & \textcolor{gray}{\ensuremath{\cdot}} & \textcolor{gray}{\ensuremath{\cdot}} & \textcolor{gray}{\ensuremath{\cdot}} & \textcolor{gray}{\ensuremath{\cdot}} & \textcolor{gray}{\ensuremath{\cdot}} & \textcolor{gray}{\ensuremath{\cdot}} & \textcolor{gray}{\ensuremath{\cdot}} & \textcolor{gray}{\ensuremath{\cdot}} & \textcolor{gray}{\ensuremath{\cdot}} & \textcolor{gray}{\ensuremath{\cdot}} & \textcolor{gray}{\ensuremath{\cdot}} & \textcolor{gray}{\ensuremath{\cdot}}\\[3pt]\hline
\end{tabular}\par}
\smallskip{\footnotesize\noindent Each cell: A, comparison, B, comparison. Times cover 1,024 queries. $\cdot$ = no approved observation. ---\textsuperscript{g} = GPU unavailable, not executed. n/a = no measured comparison.\par
\noindent \textsuperscript{f}SIMD-ref request observed as scalar, compared only as that effective path. \textsuperscript{p}Accepted by explicit preservation after wrapper postflight failure, no run receipt claimed.\par
\noindent MPF/MPS = minimizers per fragment/sequence. Row order: index representation / minimizer, then chaining / base alignment. SIMD-ref is the reference minimizer, Bounded is bounded ntHash.\par
\noindent Blue = one lowest A and one lowest B down each machine/workload column, across continuation pages. Exact ties select the first row. Repeated baselines are marked only on page 1.\par
\noindent Light gray = minimap2 baseline. Pale yellow = preferred MPF / Bounded / SIMD-DP / Block configuration, not a claim of winning every column.\par
\noindent When all Vulkan modes are unmeasured, they share one row. Baseline comparisons use the fastest measured BioseqDB effective path separately for A and B.\par}
\clearpage
\noindent{\large\bfseries AWS: 2 Gbp, k=15, w=10 (mapping interval)\hfill 6/7}\par\smallskip
{\fontsize{7.1}{8.1}\selectfont\setlength{\tabcolsep}{1pt}\setlength{\fboxsep}{2pt}\renewcommand{\arraystretch}{1.12}
\noindent\begin{tabular}{@{}p{14mm}p{27mm}|p{13.55mm}|p{13.55mm}|p{13.55mm}|p{13.55mm}|p{13.55mm}|p{13.55mm}|p{13.55mm}|p{13.55mm}|p{13.55mm}|p{13.55mm}|p{13.55mm}|p{13.55mm}|p{13.55mm}|p{13.55mm}|p{13.55mm}|p{13.55mm}|@{}}\hline
Engine & Execution & \multicolumn{8}{c|}{\textbf{Index}} & \multicolumn{8}{c|}{\textbf{No index}}\\
 &  & \multicolumn{8}{c|}{\textbf{2 Gbp}} & \multicolumn{8}{c|}{\textbf{2 Gbp}}\\
 &  & \multicolumn{8}{c|}{\textbf{k=15, w=10}} & \multicolumn{8}{c|}{\textbf{k=15, w=10}}\\
 &  & \multicolumn{4}{c|}{\textbf{DNA4}} & \multicolumn{4}{c|}{\textbf{DNA5}} & \multicolumn{4}{c|}{\textbf{DNA4}} & \multicolumn{4}{c|}{\textbf{DNA5}}\\
 &  & \multicolumn{1}{c|}{\textbf{M7A}} & \multicolumn{1}{c|}{\textbf{M7I}} & \multicolumn{1}{c|}{\textbf{M4}} & \multicolumn{1}{c|}{\textbf{C7G}} & \multicolumn{1}{c|}{\textbf{M7A}} & \multicolumn{1}{c|}{\textbf{M7I}} & \multicolumn{1}{c|}{\textbf{M4}} & \multicolumn{1}{c|}{\textbf{C7G}} & \multicolumn{1}{c|}{\textbf{M7A}} & \multicolumn{1}{c|}{\textbf{M7I}} & \multicolumn{1}{c|}{\textbf{M4}} & \multicolumn{1}{c|}{\textbf{C7G}} & \multicolumn{1}{c|}{\textbf{M7A}} & \multicolumn{1}{c|}{\textbf{M7I}} & \multicolumn{1}{c|}{\textbf{M4}} & \multicolumn{1}{c|}{\textbf{C7G}}\\
\hline
\rowcolor{BenchBaseline}[0.5pt][0.5pt]\textbf{minimap2} & \textbf{Baseline} & \hspace*{1pt}\begin{tabular}[t]{@{}c@{}}\rule{0pt}{12pt}11.2s\\[-.2pt]{\fontsize{5.8}{6.4}\selectfont \textcolor{BrickRed}{$ \blacktriangledown\,3.1\times $}}\\[-.2pt]10s\\[-.2pt]{\fontsize{5.8}{6.4}\selectfont \textcolor{BrickRed}{$ \blacktriangledown\,3.6\times $}}\end{tabular} & \hspace*{1pt}\begin{tabular}[t]{@{}c@{}}\rule{0pt}{12pt}10.7s\\[-.2pt]{\fontsize{5.8}{6.4}\selectfont \textcolor{BrickRed}{$ \blacktriangledown\,3.2\times $}}\\[-.2pt]8.6s\\[-.2pt]{\fontsize{5.8}{6.4}\selectfont \textcolor{BrickRed}{$ \blacktriangledown\,3.5\times $}}\end{tabular} & \hspace*{1pt}\begin{tabular}[t]{@{}c@{}}\rule{0pt}{12pt}22.5s\\[-.2pt]{\fontsize{5.8}{6.4}\selectfont \textcolor{BrickRed}{$ \blacktriangledown\,3.6\times $}}\\[-.2pt]18.8s\\[-.2pt]{\fontsize{5.8}{6.4}\selectfont \textcolor{BrickRed}{$ \blacktriangledown\,4.4\times $}}\end{tabular} & \hspace*{1pt}\begin{tabular}[t]{@{}c@{}}\rule{0pt}{12pt}1m\,13.1s\\[-.2pt]{\fontsize{5.8}{6.4}\selectfont \textcolor{BrickRed}{$ \blacktriangledown\,22.5\times $}}\\[-.2pt]1m\,5.7s\\[-.2pt]{\fontsize{5.8}{6.4}\selectfont \textcolor{BrickRed}{$ \blacktriangledown\,27.7\times $}}\end{tabular} & \hspace*{1pt}\begin{tabular}[t]{@{}c@{}}\rule{0pt}{12pt}9.9s\\[-.2pt]{\fontsize{5.8}{6.4}\selectfont \textcolor{BrickRed}{$ \blacktriangledown\,21\% $}}\\[-.2pt]8.8s\\[-.2pt]{\fontsize{5.8}{6.4}\selectfont \textcolor{BrickRed}{$ \blacktriangledown\,1.5\times $}}\end{tabular} & \hspace*{1pt}\begin{tabular}[t]{@{}c@{}}\rule{0pt}{12pt}9.5s\\[-.2pt]{\fontsize{5.8}{6.4}\selectfont \textcolor{BrickRed}{$ \blacktriangledown\,1.4\times $}}\\[-.2pt]7.6s\\[-.2pt]{\fontsize{5.8}{6.4}\selectfont \textcolor{BrickRed}{$ \blacktriangledown\,1.6\times $}}\end{tabular} & \hspace*{1pt}\begin{tabular}[t]{@{}c@{}}\rule{0pt}{12pt}16.5s\\[-.2pt]{\fontsize{5.8}{6.4}\selectfont \textcolor{BrickRed}{$ \blacktriangledown\,1.7\times $}}\\[-.2pt]12.9s\\[-.2pt]{\fontsize{5.8}{6.4}\selectfont \textcolor{BrickRed}{$ \blacktriangledown\,2.0\times $}}\end{tabular} & \hspace*{1pt}\begin{tabular}[t]{@{}c@{}}\rule{0pt}{12pt}28.2s\\[-.2pt]{\fontsize{5.8}{6.4}\selectfont \textcolor{BrickRed}{$ \blacktriangledown\,2.1\times $}}\\[-.2pt]26.2s\\[-.2pt]{\fontsize{5.8}{6.4}\selectfont \textcolor{BrickRed}{$ \blacktriangledown\,2.4\times $}}\end{tabular} & \hspace*{1pt}\begin{tabular}[t]{@{}c@{}}\rule{0pt}{12pt}53s\\[-.2pt]{\fontsize{5.8}{6.4}\selectfont \textcolor{ForestGreen}{$ \blacktriangle\,2.2\times $}}\\[-.2pt]51.9s\\[-.2pt]{\fontsize{5.8}{6.4}\selectfont \textcolor{ForestGreen}{$ \blacktriangle\,2.1\times $}}\end{tabular} & \hspace*{1pt}\begin{tabular}[t]{@{}c@{}}\rule{0pt}{12pt}56.6s\\[-.2pt]{\fontsize{5.8}{6.4}\selectfont \textcolor{ForestGreen}{$ \blacktriangle\,2.1\times $}}\\[-.2pt]54.6s\\[-.2pt]{\fontsize{5.8}{6.4}\selectfont \textcolor{ForestGreen}{$ \blacktriangle\,1.9\times $}}\end{tabular} & \hspace*{1pt}\begin{tabular}[t]{@{}c@{}}\rule{0pt}{12pt}1m\,55.5s\\[-.2pt]{\fontsize{5.8}{6.4}\selectfont \textcolor{ForestGreen}{$ \blacktriangle\,1.8\times $}}\\[-.2pt]1m\,51.8s\\[-.2pt]{\fontsize{5.8}{6.4}\selectfont \textcolor{ForestGreen}{$ \blacktriangle\,1.5\times $}}\end{tabular} & \hspace*{1pt}\begin{tabular}[t]{@{}c@{}}\rule{0pt}{12pt}6m\,1.3s\\[-.2pt]{\fontsize{5.8}{6.4}\selectfont \textcolor{gray}{n/a}}\\[-.2pt]5m\,55.2s\\[-.2pt]{\fontsize{5.8}{6.4}\selectfont \textcolor{gray}{n/a}}\end{tabular} & \hspace*{1pt}\begin{tabular}[t]{@{}c@{}}\rule{0pt}{12pt}48.1s\\[-.2pt]{\fontsize{5.8}{6.4}\selectfont \textcolor{ForestGreen}{$ \blacktriangle\,7.3\times $}}\\[-.2pt]47s\\[-.2pt]{\fontsize{5.8}{6.4}\selectfont \textcolor{ForestGreen}{$ \blacktriangle\,2.4\times $}}\end{tabular} & \hspace*{1pt}\begin{tabular}[t]{@{}c@{}}\rule{0pt}{12pt}44.5s\\[-.2pt]{\fontsize{5.8}{6.4}\selectfont \textcolor{ForestGreen}{$ \blacktriangle\,6.5\times $}}\\[-.2pt]42.6s\\[-.2pt]{\fontsize{5.8}{6.4}\selectfont \textcolor{ForestGreen}{$ \blacktriangle\,2.7\times $}}\end{tabular} & \hspace*{1pt}\begin{tabular}[t]{@{}c@{}}\rule{0pt}{12pt}1m\,30.2s\\[-.2pt]{\fontsize{5.8}{6.4}\selectfont \textcolor{ForestGreen}{$ \blacktriangle\,6.4\times $}}\\[-.2pt]1m\,26.7s\\[-.2pt]{\fontsize{5.8}{6.4}\selectfont \textcolor{ForestGreen}{$ \blacktriangle\,2.2\times $}}\end{tabular} & \hspace*{1pt}\begin{tabular}[t]{@{}c@{}}\rule{0pt}{12pt}2m\,46.4s\\[-.2pt]{\fontsize{5.8}{6.4}\selectfont \textcolor{gray}{n/a}}\\[-.2pt]2m\,42.7s\\[-.2pt]{\fontsize{5.8}{6.4}\selectfont \textcolor{gray}{n/a}}\end{tabular}\\[3pt]\hline
\textbf{BioseqDB} & \begin{tabular}[t]{@{}l@{}}\rule{0pt}{12pt}MPS / Scalar\\Scalar-DP / KSW2\end{tabular} & \hspace*{1pt}\begin{tabular}[t]{@{}c@{}}\rule{0pt}{12pt}8.3s\\[-.2pt]{\fontsize{5.8}{6.4}\selectfont \textcolor{ForestGreen}{$ \blacktriangle\,26\% $}}\\[-.2pt]5.4s\\[-.2pt]{\fontsize{5.8}{6.4}\selectfont \textcolor{ForestGreen}{$ \blacktriangle\,1.8\times $}}\end{tabular} & \hspace*{1pt}\begin{tabular}[t]{@{}c@{}}\rule{0pt}{12pt}8.7s\\[-.2pt]{\fontsize{5.8}{6.4}\selectfont \textcolor{ForestGreen}{$ \blacktriangle\,18\% $}}\\[-.2pt]5.3s\\[-.2pt]{\fontsize{5.8}{6.4}\selectfont \textcolor{ForestGreen}{$ \blacktriangle\,1.6\times $}}\end{tabular} & \textcolor{gray}{\ensuremath{\cdot}} & \textcolor{gray}{\ensuremath{\cdot}} & \textcolor{gray}{\ensuremath{\cdot}} & \textcolor{gray}{\ensuremath{\cdot}} & \textcolor{gray}{\ensuremath{\cdot}} & \textcolor{gray}{\ensuremath{\cdot}} & \textcolor{gray}{\ensuremath{\cdot}} & \textcolor{gray}{\ensuremath{\cdot}} & \textcolor{gray}{\ensuremath{\cdot}} & \textcolor{gray}{\ensuremath{\cdot}} & \textcolor{gray}{\ensuremath{\cdot}} & \textcolor{gray}{\ensuremath{\cdot}} & \textcolor{gray}{\ensuremath{\cdot}} & \textcolor{gray}{\ensuremath{\cdot}}\\[3pt]\hline
 & \begin{tabular}[t]{@{}l@{}}\rule{0pt}{12pt}MPS / Scalar\\SIMD-DP / Block\end{tabular} & \hspace*{1pt}\begin{tabular}[t]{@{}c@{}}\rule{0pt}{12pt}6.8s\\[-.2pt]{\fontsize{5.8}{6.4}\selectfont \textcolor{ForestGreen}{$ \blacktriangle\,1.6\times $}}\\[-.2pt]6.2s\\[-.2pt]{\fontsize{5.8}{6.4}\selectfont \textcolor{ForestGreen}{$ \blacktriangle\,1.6\times $}}\end{tabular} & \hspace*{1pt}\begin{tabular}[t]{@{}c@{}}\rule{0pt}{12pt}7.1s\\[-.2pt]{\fontsize{5.8}{6.4}\selectfont \textcolor{ForestGreen}{$ \blacktriangle\,1.5\times $}}\\[-.2pt]6.2s\\[-.2pt]{\fontsize{5.8}{6.4}\selectfont \textcolor{ForestGreen}{$ \blacktriangle\,28\% $}}\end{tabular} & \textcolor{gray}{\ensuremath{\cdot}} & \hspace*{1pt}\begin{tabular}[t]{@{}c@{}}\rule{0pt}{12pt}6.6s\\[-.2pt]{\fontsize{5.8}{6.4}\selectfont \textcolor{ForestGreen}{$ \blacktriangle\,11.1\times $}}\\[-.2pt]5.6s\\[-.2pt]{\fontsize{5.8}{6.4}\selectfont \textcolor{ForestGreen}{$ \blacktriangle\,11.7\times $}}\end{tabular} & \textcolor{gray}{\ensuremath{\cdot}} & \textcolor{gray}{\ensuremath{\cdot}} & \textcolor{gray}{\ensuremath{\cdot}} & \textcolor{gray}{\ensuremath{\cdot}} & \hspace*{1pt}\begin{tabular}[t]{@{}c@{}}\rule{0pt}{12pt}3m\,9.8s\\[-.2pt]{\fontsize{5.8}{6.4}\selectfont \textcolor{BrickRed}{$ \blacktriangledown\,3.6\times $}}\\[-.2pt]2m\,47.5s\\[-.2pt]{\fontsize{5.8}{6.4}\selectfont \textcolor{BrickRed}{$ \blacktriangledown\,3.2\times $}}\end{tabular} & \hspace*{1pt}\begin{tabular}[t]{@{}c@{}}\rule{0pt}{12pt}3m\,6.4s\\[-.2pt]{\fontsize{5.8}{6.4}\selectfont \textcolor{BrickRed}{$ \blacktriangledown\,3.3\times $}}\\[-.2pt]2m\,46.4s\\[-.2pt]{\fontsize{5.8}{6.4}\selectfont \textcolor{BrickRed}{$ \blacktriangledown\,3.0\times $}}\end{tabular} & \textcolor{gray}{\ensuremath{\cdot}} & \textcolor{gray}{\ensuremath{\cdot}} & \textcolor{gray}{\ensuremath{\cdot}} & \textcolor{gray}{\ensuremath{\cdot}} & \textcolor{gray}{\ensuremath{\cdot}} & \textcolor{gray}{\ensuremath{\cdot}}\\[3pt]\hline
 & \begin{tabular}[t]{@{}l@{}}\rule{0pt}{12pt}MPS / Scalar\\SIMD-DP / KSW2\end{tabular} & \hspace*{1pt}\begin{tabular}[t]{@{}c@{}}\rule{0pt}{12pt}8.3s\\[-.2pt]{\fontsize{5.8}{6.4}\selectfont \textcolor{ForestGreen}{$ \blacktriangle\,26\% $}}\\[-.2pt]5.5s\\[-.2pt]{\fontsize{5.8}{6.4}\selectfont \textcolor{ForestGreen}{$ \blacktriangle\,1.8\times $}}\end{tabular} & \hspace*{1pt}\begin{tabular}[t]{@{}c@{}}\rule{0pt}{12pt}8.3s\\[-.2pt]{\fontsize{5.8}{6.4}\selectfont \textcolor{ForestGreen}{$ \blacktriangle\,22\% $}}\\[-.2pt]5.2s\\[-.2pt]{\fontsize{5.8}{6.4}\selectfont \textcolor{ForestGreen}{$ \blacktriangle\,1.6\times $}}\end{tabular} & \textcolor{gray}{\ensuremath{\cdot}} & \textcolor{gray}{\ensuremath{\cdot}} & \textcolor{gray}{\ensuremath{\cdot}} & \textcolor{gray}{\ensuremath{\cdot}} & \textcolor{gray}{\ensuremath{\cdot}} & \textcolor{gray}{\ensuremath{\cdot}} & \hspace*{1pt}\begin{tabular}[t]{@{}c@{}}\rule{0pt}{12pt}3m\,1.5s\\[-.2pt]{\fontsize{5.8}{6.4}\selectfont \textcolor{BrickRed}{$ \blacktriangledown\,3.4\times $}}\\[-.2pt]2m\,37.6s\\[-.2pt]{\fontsize{5.8}{6.4}\selectfont \textcolor{BrickRed}{$ \blacktriangledown\,3.0\times $}}\end{tabular} & \hspace*{1pt}\begin{tabular}[t]{@{}c@{}}\rule{0pt}{12pt}3m\,21.1s\\[-.2pt]{\fontsize{5.8}{6.4}\selectfont \textcolor{BrickRed}{$ \blacktriangledown\,3.6\times $}}\\[-.2pt]2m\,58.7s\\[-.2pt]{\fontsize{5.8}{6.4}\selectfont \textcolor{BrickRed}{$ \blacktriangledown\,3.3\times $}}\end{tabular} & \textcolor{gray}{\ensuremath{\cdot}} & \textcolor{gray}{\ensuremath{\cdot}} & \textcolor{gray}{\ensuremath{\cdot}} & \textcolor{gray}{\ensuremath{\cdot}} & \textcolor{gray}{\ensuremath{\cdot}} & \textcolor{gray}{\ensuremath{\cdot}}\\[3pt]\hline
 & \begin{tabular}[t]{@{}l@{}}\rule{0pt}{12pt}MPS / SIMD-ref\\RMQ / Block\end{tabular} & \textcolor{gray}{\ensuremath{\cdot}} & \textcolor{gray}{\ensuremath{\cdot}} & \textcolor{gray}{\ensuremath{\cdot}} & \textcolor{gray}{\ensuremath{\cdot}} & \textcolor{gray}{\ensuremath{\cdot}} & \textcolor{gray}{\ensuremath{\cdot}} & \textcolor{gray}{\ensuremath{\cdot}} & \textcolor{gray}{\ensuremath{\cdot}} & \textcolor{gray}{\ensuremath{\cdot}} & \textcolor{gray}{\ensuremath{\cdot}} & \hspace*{1pt}\begin{tabular}[t]{@{}c@{}}\rule{0pt}{12pt}3m\,24.8s\\[-.2pt]{\fontsize{5.8}{6.4}\selectfont \textcolor{BrickRed}{$ \blacktriangledown\,1.8\times $}}\\[-.2pt]2m\,59.9s\\[-.2pt]{\fontsize{5.8}{6.4}\selectfont \textcolor{BrickRed}{$ \blacktriangledown\,1.6\times $}}\end{tabular} & \textcolor{gray}{\ensuremath{\cdot}} & \textcolor{gray}{\ensuremath{\cdot}} & \textcolor{gray}{\ensuremath{\cdot}} & \textcolor{gray}{\ensuremath{\cdot}} & \textcolor{gray}{\ensuremath{\cdot}}\\[3pt]\hline
 & \begin{tabular}[t]{@{}l@{}}\rule{0pt}{12pt}MPS / SIMD-ref\\Scalar-DP / KSW2\end{tabular} & \textcolor{gray}{\ensuremath{\cdot}} & \textcolor{gray}{\ensuremath{\cdot}} & \hspace*{1pt}\begin{tabular}[t]{@{}c@{}}\rule{0pt}{12pt}16.4s\\[-.2pt]{\fontsize{5.8}{6.4}\selectfont \textcolor{ForestGreen}{$ \blacktriangle\,27\% $}}\\[-.2pt]8.9s\\[-.2pt]{\fontsize{5.8}{6.4}\selectfont \textcolor{ForestGreen}{$ \blacktriangle\,2.1\times $}}\end{tabular} & \textcolor{gray}{\ensuremath{\cdot}} & \textcolor{gray}{\ensuremath{\cdot}} & \textcolor{gray}{\ensuremath{\cdot}} & \textcolor{gray}{\ensuremath{\cdot}} & \textcolor{gray}{\ensuremath{\cdot}} & \textcolor{gray}{\ensuremath{\cdot}} & \textcolor{gray}{\ensuremath{\cdot}} & \textcolor{gray}{\ensuremath{\cdot}} & \textcolor{gray}{\ensuremath{\cdot}} & \textcolor{gray}{\ensuremath{\cdot}} & \textcolor{gray}{\ensuremath{\cdot}} & \textcolor{gray}{\ensuremath{\cdot}} & \textcolor{gray}{\ensuremath{\cdot}}\\[3pt]\hline
 & \begin{tabular}[t]{@{}l@{}}\rule{0pt}{12pt}MPS / SIMD-ref\\SIMD-DP / Block\end{tabular} & \textcolor{gray}{\ensuremath{\cdot}} & \textcolor{gray}{\ensuremath{\cdot}} & \textcolor{gray}{\ensuremath{\cdot}} & \hspace*{1pt}\begin{tabular}[t]{@{}c@{}}\rule{0pt}{12pt}6.4s\\[-.2pt]{\fontsize{5.8}{6.4}\selectfont \textcolor{ForestGreen}{$ \blacktriangle\,11.4\times $}}\\[-.2pt]5.5s\\[-.2pt]{\fontsize{5.8}{6.4}\selectfont \textcolor{ForestGreen}{$ \blacktriangle\,12.0\times $}}\end{tabular} & \textcolor{gray}{\ensuremath{\cdot}} & \textcolor{gray}{\ensuremath{\cdot}} & \textcolor{gray}{\ensuremath{\cdot}} & \textcolor{gray}{\ensuremath{\cdot}} & \textcolor{gray}{\ensuremath{\cdot}} & \textcolor{gray}{\ensuremath{\cdot}} & \hspace*{1pt}\begin{tabular}[t]{@{}c@{}}\rule{0pt}{12pt}4m\,15.9s\\[-.2pt]{\fontsize{5.8}{6.4}\selectfont \textcolor{BrickRed}{$ \blacktriangledown\,2.2\times $}}\\[-.2pt]3m\,41.9s\\[-.2pt]{\fontsize{5.8}{6.4}\selectfont \textcolor{BrickRed}{$ \blacktriangledown\,2.0\times $}}\end{tabular} & \textcolor{gray}{\ensuremath{\cdot}} & \textcolor{gray}{\ensuremath{\cdot}} & \textcolor{gray}{\ensuremath{\cdot}} & \textcolor{gray}{\ensuremath{\cdot}} & \textcolor{gray}{\ensuremath{\cdot}}\\[3pt]\hline
 & \begin{tabular}[t]{@{}l@{}}\rule{0pt}{12pt}MPS / SIMD-ref\\SIMD-DP / KSW2\end{tabular} & \textcolor{gray}{\ensuremath{\cdot}} & \textcolor{gray}{\ensuremath{\cdot}} & \hspace*{1pt}\begin{tabular}[t]{@{}c@{}}\rule{0pt}{12pt}15.7s\\[-.2pt]{\fontsize{5.8}{6.4}\selectfont \textcolor{ForestGreen}{$ \blacktriangle\,1.4\times $}}\\[-.2pt]9.3s\\[-.2pt]{\fontsize{5.8}{6.4}\selectfont \textcolor{ForestGreen}{$ \blacktriangle\,2.0\times $}}\end{tabular} & \textcolor{gray}{\ensuremath{\cdot}} & \textcolor{gray}{\ensuremath{\cdot}} & \textcolor{gray}{\ensuremath{\cdot}} & \textcolor{gray}{\ensuremath{\cdot}} & \textcolor{gray}{\ensuremath{\cdot}} & \textcolor{gray}{\ensuremath{\cdot}} & \textcolor{gray}{\ensuremath{\cdot}} & \textcolor{gray}{\ensuremath{\cdot}} & \textcolor{gray}{\ensuremath{\cdot}} & \textcolor{gray}{\ensuremath{\cdot}} & \textcolor{gray}{\ensuremath{\cdot}} & \textcolor{gray}{\ensuremath{\cdot}} & \textcolor{gray}{\ensuremath{\cdot}}\\[3pt]\hline
\end{tabular}\par}
\smallskip{\footnotesize\noindent Each cell: A, comparison, B, comparison. Times cover 1,024 queries. $\cdot$ = no approved observation. ---\textsuperscript{g} = GPU unavailable, not executed. n/a = no measured comparison.\par
\noindent \textsuperscript{f}SIMD-ref request observed as scalar, compared only as that effective path. \textsuperscript{p}Accepted by explicit preservation after wrapper postflight failure, no run receipt claimed.\par
\noindent MPF/MPS = minimizers per fragment/sequence. Row order: index representation / minimizer, then chaining / base alignment. SIMD-ref is the reference minimizer, Bounded is bounded ntHash.\par
\noindent Blue = one lowest A and one lowest B down each machine/workload column, across continuation pages. Exact ties select the first row. Repeated baselines are marked only on page 1.\par
\noindent Light gray = minimap2 baseline. Pale yellow = preferred MPF / Bounded / SIMD-DP / Block configuration, not a claim of winning every column.\par
\noindent When all Vulkan modes are unmeasured, they share one row. Baseline comparisons use the fastest measured BioseqDB effective path separately for A and B.\par}
\clearpage
\noindent{\large\bfseries AWS: 2 Gbp, k=15, w=10 (mapping interval)\hfill 7/7}\par\smallskip
{\fontsize{7.1}{8.1}\selectfont\setlength{\tabcolsep}{1pt}\setlength{\fboxsep}{2pt}\renewcommand{\arraystretch}{1.12}
\noindent\begin{tabular}{@{}p{14mm}p{27mm}|p{13.55mm}|p{13.55mm}|p{13.55mm}|p{13.55mm}|p{13.55mm}|p{13.55mm}|p{13.55mm}|p{13.55mm}|p{13.55mm}|p{13.55mm}|p{13.55mm}|p{13.55mm}|p{13.55mm}|p{13.55mm}|p{13.55mm}|p{13.55mm}|@{}}\hline
Engine & Execution & \multicolumn{8}{c|}{\textbf{Index}} & \multicolumn{8}{c|}{\textbf{No index}}\\
 &  & \multicolumn{8}{c|}{\textbf{2 Gbp}} & \multicolumn{8}{c|}{\textbf{2 Gbp}}\\
 &  & \multicolumn{8}{c|}{\textbf{k=15, w=10}} & \multicolumn{8}{c|}{\textbf{k=15, w=10}}\\
 &  & \multicolumn{4}{c|}{\textbf{DNA4}} & \multicolumn{4}{c|}{\textbf{DNA5}} & \multicolumn{4}{c|}{\textbf{DNA4}} & \multicolumn{4}{c|}{\textbf{DNA5}}\\
 &  & \multicolumn{1}{c|}{\textbf{M7A}} & \multicolumn{1}{c|}{\textbf{M7I}} & \multicolumn{1}{c|}{\textbf{M4}} & \multicolumn{1}{c|}{\textbf{C7G}} & \multicolumn{1}{c|}{\textbf{M7A}} & \multicolumn{1}{c|}{\textbf{M7I}} & \multicolumn{1}{c|}{\textbf{M4}} & \multicolumn{1}{c|}{\textbf{C7G}} & \multicolumn{1}{c|}{\textbf{M7A}} & \multicolumn{1}{c|}{\textbf{M7I}} & \multicolumn{1}{c|}{\textbf{M4}} & \multicolumn{1}{c|}{\textbf{C7G}} & \multicolumn{1}{c|}{\textbf{M7A}} & \multicolumn{1}{c|}{\textbf{M7I}} & \multicolumn{1}{c|}{\textbf{M4}} & \multicolumn{1}{c|}{\textbf{C7G}}\\
\hline
\rowcolor{BenchBaseline}[0.5pt][0.5pt]\textbf{minimap2} & \textbf{Baseline} & \hspace*{1pt}\begin{tabular}[t]{@{}c@{}}\rule{0pt}{12pt}11.2s\\[-.2pt]{\fontsize{5.8}{6.4}\selectfont \textcolor{BrickRed}{$ \blacktriangledown\,3.1\times $}}\\[-.2pt]10s\\[-.2pt]{\fontsize{5.8}{6.4}\selectfont \textcolor{BrickRed}{$ \blacktriangledown\,3.6\times $}}\end{tabular} & \hspace*{1pt}\begin{tabular}[t]{@{}c@{}}\rule{0pt}{12pt}10.7s\\[-.2pt]{\fontsize{5.8}{6.4}\selectfont \textcolor{BrickRed}{$ \blacktriangledown\,3.2\times $}}\\[-.2pt]8.6s\\[-.2pt]{\fontsize{5.8}{6.4}\selectfont \textcolor{BrickRed}{$ \blacktriangledown\,3.5\times $}}\end{tabular} & \hspace*{1pt}\begin{tabular}[t]{@{}c@{}}\rule{0pt}{12pt}22.5s\\[-.2pt]{\fontsize{5.8}{6.4}\selectfont \textcolor{BrickRed}{$ \blacktriangledown\,3.6\times $}}\\[-.2pt]18.8s\\[-.2pt]{\fontsize{5.8}{6.4}\selectfont \textcolor{BrickRed}{$ \blacktriangledown\,4.4\times $}}\end{tabular} & \hspace*{1pt}\begin{tabular}[t]{@{}c@{}}\rule{0pt}{12pt}1m\,13.1s\\[-.2pt]{\fontsize{5.8}{6.4}\selectfont \textcolor{BrickRed}{$ \blacktriangledown\,22.5\times $}}\\[-.2pt]1m\,5.7s\\[-.2pt]{\fontsize{5.8}{6.4}\selectfont \textcolor{BrickRed}{$ \blacktriangledown\,27.7\times $}}\end{tabular} & \hspace*{1pt}\begin{tabular}[t]{@{}c@{}}\rule{0pt}{12pt}9.9s\\[-.2pt]{\fontsize{5.8}{6.4}\selectfont \textcolor{BrickRed}{$ \blacktriangledown\,21\% $}}\\[-.2pt]8.8s\\[-.2pt]{\fontsize{5.8}{6.4}\selectfont \textcolor{BrickRed}{$ \blacktriangledown\,1.5\times $}}\end{tabular} & \hspace*{1pt}\begin{tabular}[t]{@{}c@{}}\rule{0pt}{12pt}9.5s\\[-.2pt]{\fontsize{5.8}{6.4}\selectfont \textcolor{BrickRed}{$ \blacktriangledown\,1.4\times $}}\\[-.2pt]7.6s\\[-.2pt]{\fontsize{5.8}{6.4}\selectfont \textcolor{BrickRed}{$ \blacktriangledown\,1.6\times $}}\end{tabular} & \hspace*{1pt}\begin{tabular}[t]{@{}c@{}}\rule{0pt}{12pt}16.5s\\[-.2pt]{\fontsize{5.8}{6.4}\selectfont \textcolor{BrickRed}{$ \blacktriangledown\,1.7\times $}}\\[-.2pt]12.9s\\[-.2pt]{\fontsize{5.8}{6.4}\selectfont \textcolor{BrickRed}{$ \blacktriangledown\,2.0\times $}}\end{tabular} & \hspace*{1pt}\begin{tabular}[t]{@{}c@{}}\rule{0pt}{12pt}28.2s\\[-.2pt]{\fontsize{5.8}{6.4}\selectfont \textcolor{BrickRed}{$ \blacktriangledown\,2.1\times $}}\\[-.2pt]26.2s\\[-.2pt]{\fontsize{5.8}{6.4}\selectfont \textcolor{BrickRed}{$ \blacktriangledown\,2.4\times $}}\end{tabular} & \hspace*{1pt}\begin{tabular}[t]{@{}c@{}}\rule{0pt}{12pt}53s\\[-.2pt]{\fontsize{5.8}{6.4}\selectfont \textcolor{ForestGreen}{$ \blacktriangle\,2.2\times $}}\\[-.2pt]51.9s\\[-.2pt]{\fontsize{5.8}{6.4}\selectfont \textcolor{ForestGreen}{$ \blacktriangle\,2.1\times $}}\end{tabular} & \hspace*{1pt}\begin{tabular}[t]{@{}c@{}}\rule{0pt}{12pt}56.6s\\[-.2pt]{\fontsize{5.8}{6.4}\selectfont \textcolor{ForestGreen}{$ \blacktriangle\,2.1\times $}}\\[-.2pt]54.6s\\[-.2pt]{\fontsize{5.8}{6.4}\selectfont \textcolor{ForestGreen}{$ \blacktriangle\,1.9\times $}}\end{tabular} & \hspace*{1pt}\begin{tabular}[t]{@{}c@{}}\rule{0pt}{12pt}1m\,55.5s\\[-.2pt]{\fontsize{5.8}{6.4}\selectfont \textcolor{ForestGreen}{$ \blacktriangle\,1.8\times $}}\\[-.2pt]1m\,51.8s\\[-.2pt]{\fontsize{5.8}{6.4}\selectfont \textcolor{ForestGreen}{$ \blacktriangle\,1.5\times $}}\end{tabular} & \hspace*{1pt}\begin{tabular}[t]{@{}c@{}}\rule{0pt}{12pt}6m\,1.3s\\[-.2pt]{\fontsize{5.8}{6.4}\selectfont \textcolor{gray}{n/a}}\\[-.2pt]5m\,55.2s\\[-.2pt]{\fontsize{5.8}{6.4}\selectfont \textcolor{gray}{n/a}}\end{tabular} & \hspace*{1pt}\begin{tabular}[t]{@{}c@{}}\rule{0pt}{12pt}48.1s\\[-.2pt]{\fontsize{5.8}{6.4}\selectfont \textcolor{ForestGreen}{$ \blacktriangle\,7.3\times $}}\\[-.2pt]47s\\[-.2pt]{\fontsize{5.8}{6.4}\selectfont \textcolor{ForestGreen}{$ \blacktriangle\,2.4\times $}}\end{tabular} & \hspace*{1pt}\begin{tabular}[t]{@{}c@{}}\rule{0pt}{12pt}44.5s\\[-.2pt]{\fontsize{5.8}{6.4}\selectfont \textcolor{ForestGreen}{$ \blacktriangle\,6.5\times $}}\\[-.2pt]42.6s\\[-.2pt]{\fontsize{5.8}{6.4}\selectfont \textcolor{ForestGreen}{$ \blacktriangle\,2.7\times $}}\end{tabular} & \hspace*{1pt}\begin{tabular}[t]{@{}c@{}}\rule{0pt}{12pt}1m\,30.2s\\[-.2pt]{\fontsize{5.8}{6.4}\selectfont \textcolor{ForestGreen}{$ \blacktriangle\,6.4\times $}}\\[-.2pt]1m\,26.7s\\[-.2pt]{\fontsize{5.8}{6.4}\selectfont \textcolor{ForestGreen}{$ \blacktriangle\,2.2\times $}}\end{tabular} & \hspace*{1pt}\begin{tabular}[t]{@{}c@{}}\rule{0pt}{12pt}2m\,46.4s\\[-.2pt]{\fontsize{5.8}{6.4}\selectfont \textcolor{gray}{n/a}}\\[-.2pt]2m\,42.7s\\[-.2pt]{\fontsize{5.8}{6.4}\selectfont \textcolor{gray}{n/a}}\end{tabular}\\[3pt]\hline
\textbf{BioseqDB} & \begin{tabular}[t]{@{}l@{}}\rule{0pt}{12pt}MPS / Two-stack\\RMQ / Block\end{tabular} & \hspace*{1pt}\begin{tabular}[t]{@{}c@{}}\rule{0pt}{12pt}6.9s\\[-.2pt]{\fontsize{5.8}{6.4}\selectfont \textcolor{ForestGreen}{$ \blacktriangle\,1.6\times $}}\\[-.2pt]6.2s\\[-.2pt]{\fontsize{5.8}{6.4}\selectfont \textcolor{ForestGreen}{$ \blacktriangle\,1.6\times $}}\end{tabular} & \hspace*{1pt}\begin{tabular}[t]{@{}c@{}}\rule{0pt}{12pt}7.5s\\[-.2pt]{\fontsize{5.8}{6.4}\selectfont \textcolor{ForestGreen}{$ \blacktriangle\,29\% $}}\\[-.2pt]6.6s\\[-.2pt]{\fontsize{5.8}{6.4}\selectfont \textcolor{ForestGreen}{$ \blacktriangle\,23\% $}}\end{tabular} & \textcolor{gray}{\ensuremath{\cdot}} & \textcolor{gray}{\ensuremath{\cdot}} & \textcolor{gray}{\ensuremath{\cdot}} & \textcolor{gray}{\ensuremath{\cdot}} & \textcolor{gray}{\ensuremath{\cdot}} & \textcolor{gray}{\ensuremath{\cdot}} & \hspace*{1pt}\begin{tabular}[t]{@{}c@{}}\rule{0pt}{12pt}2m\,47.3s\\[-.2pt]{\fontsize{5.8}{6.4}\selectfont \textcolor{BrickRed}{$ \blacktriangledown\,3.2\times $}}\\[-.2pt]2m\,24.6s\\[-.2pt]{\fontsize{5.8}{6.4}\selectfont \textcolor{BrickRed}{$ \blacktriangledown\,2.8\times $}}\end{tabular} & \hspace*{1pt}\begin{tabular}[t]{@{}c@{}}\rule{0pt}{12pt}2m\,30.3s\\[-.2pt]{\fontsize{5.8}{6.4}\selectfont \textcolor{BrickRed}{$ \blacktriangledown\,2.7\times $}}\\[-.2pt]2m\,10.8s\\[-.2pt]{\fontsize{5.8}{6.4}\selectfont \textcolor{BrickRed}{$ \blacktriangledown\,2.4\times $}}\end{tabular} & \hspace*{1pt}\begin{tabular}[t]{@{}c@{}}\rule{0pt}{12pt}4m\,45s\\[-.2pt]{\fontsize{5.8}{6.4}\selectfont \textcolor{BrickRed}{$ \blacktriangledown\,2.5\times $}}\\[-.2pt]4m\,11.3s\\[-.2pt]{\fontsize{5.8}{6.4}\selectfont \textcolor{BrickRed}{$ \blacktriangledown\,2.2\times $}}\end{tabular} & \textcolor{gray}{\ensuremath{\cdot}} & \textcolor{gray}{\ensuremath{\cdot}} & \textcolor{gray}{\ensuremath{\cdot}} & \textcolor{gray}{\ensuremath{\cdot}} & \textcolor{gray}{\ensuremath{\cdot}}\\[3pt]\hline
 & \begin{tabular}[t]{@{}l@{}}\rule{0pt}{12pt}MPS / Two-stack\\RMQ / KSW2\end{tabular} & \hspace*{1pt}\begin{tabular}[t]{@{}c@{}}\rule{0pt}{12pt}8.3s\\[-.2pt]{\fontsize{5.8}{6.4}\selectfont \textcolor{ForestGreen}{$ \blacktriangle\,26\% $}}\\[-.2pt]5.3s\\[-.2pt]{\fontsize{5.8}{6.4}\selectfont \textcolor{ForestGreen}{$ \blacktriangle\,1.9\times $}}\end{tabular} & \hspace*{1pt}\begin{tabular}[t]{@{}c@{}}\rule{0pt}{12pt}8.1s\\[-.2pt]{\fontsize{5.8}{6.4}\selectfont \textcolor{ForestGreen}{$ \blacktriangle\,23\% $}}\\[-.2pt]5.1s\\[-.2pt]{\fontsize{5.8}{6.4}\selectfont \textcolor{ForestGreen}{$ \blacktriangle\,1.7\times $}}\end{tabular} & \textcolor{gray}{\ensuremath{\cdot}} & \textcolor{gray}{\ensuremath{\cdot}} & \textcolor{gray}{\ensuremath{\cdot}} & \textcolor{gray}{\ensuremath{\cdot}} & \textcolor{gray}{\ensuremath{\cdot}} & \textcolor{gray}{\ensuremath{\cdot}} & \hspace*{1pt}\begin{tabular}[t]{@{}c@{}}\rule{0pt}{12pt}2m\,51.1s\\[-.2pt]{\fontsize{5.8}{6.4}\selectfont \textcolor{BrickRed}{$ \blacktriangledown\,3.2\times $}}\\[-.2pt]2m\,26.7s\\[-.2pt]{\fontsize{5.8}{6.4}\selectfont \textcolor{BrickRed}{$ \blacktriangledown\,2.8\times $}}\end{tabular} & \textcolor{gray}{\ensuremath{\cdot}} & \textcolor{gray}{\ensuremath{\cdot}} & \textcolor{gray}{\ensuremath{\cdot}} & \textcolor{gray}{\ensuremath{\cdot}} & \textcolor{gray}{\ensuremath{\cdot}} & \textcolor{gray}{\ensuremath{\cdot}} & \textcolor{gray}{\ensuremath{\cdot}}\\[3pt]\hline
 & \begin{tabular}[t]{@{}l@{}}\rule{0pt}{12pt}MPS / Two-stack\\Scalar-DP / Block\end{tabular} & \hspace*{1pt}\begin{tabular}[t]{@{}c@{}}\rule{0pt}{12pt}6.9s\\[-.2pt]{\fontsize{5.8}{6.4}\selectfont \textcolor{ForestGreen}{$ \blacktriangle\,1.6\times $}}\\[-.2pt]6.2s\\[-.2pt]{\fontsize{5.8}{6.4}\selectfont \textcolor{ForestGreen}{$ \blacktriangle\,1.6\times $}}\end{tabular} & \hspace*{1pt}\begin{tabular}[t]{@{}c@{}}\rule{0pt}{12pt}7.2s\\[-.2pt]{\fontsize{5.8}{6.4}\selectfont \textcolor{ForestGreen}{$ \blacktriangle\,1.5\times $}}\\[-.2pt]6.3s\\[-.2pt]{\fontsize{5.8}{6.4}\selectfont \textcolor{ForestGreen}{$ \blacktriangle\,26\% $}}\end{tabular} & \textcolor{gray}{\ensuremath{\cdot}} & \textcolor{gray}{\ensuremath{\cdot}} & \textcolor{gray}{\ensuremath{\cdot}} & \textcolor{gray}{\ensuremath{\cdot}} & \textcolor{gray}{\ensuremath{\cdot}} & \textcolor{gray}{\ensuremath{\cdot}} & \textcolor{gray}{\ensuremath{\cdot}} & \textcolor{gray}{\ensuremath{\cdot}} & \textcolor{gray}{\ensuremath{\cdot}} & \textcolor{gray}{\ensuremath{\cdot}} & \textcolor{gray}{\ensuremath{\cdot}} & \textcolor{gray}{\ensuremath{\cdot}} & \textcolor{gray}{\ensuremath{\cdot}} & \textcolor{gray}{\ensuremath{\cdot}}\\[3pt]\hline
 & \begin{tabular}[t]{@{}l@{}}\rule{0pt}{12pt}MPS / Two-stack\\Scalar-DP / KSW2\end{tabular} & \hspace*{1pt}\begin{tabular}[t]{@{}c@{}}\rule{0pt}{12pt}8.4s\\[-.2pt]{\fontsize{5.8}{6.4}\selectfont \textcolor{ForestGreen}{$ \blacktriangle\,25\% $}}\\[-.2pt]5.4s\\[-.2pt]{\fontsize{5.8}{6.4}\selectfont \textcolor{ForestGreen}{$ \blacktriangle\,1.8\times $}}\end{tabular} & \hspace*{1pt}\begin{tabular}[t]{@{}c@{}}\rule{0pt}{12pt}9.4s\\[-.2pt]{\fontsize{5.8}{6.4}\selectfont \textcolor{ForestGreen}{$ \blacktriangle\,12\% $}}\\[-.2pt]5.5s\\[-.2pt]{\fontsize{5.8}{6.4}\selectfont \textcolor{ForestGreen}{$ \blacktriangle\,1.6\times $}}\end{tabular} & \textcolor{gray}{\ensuremath{\cdot}} & \textcolor{gray}{\ensuremath{\cdot}} & \textcolor{gray}{\ensuremath{\cdot}} & \textcolor{gray}{\ensuremath{\cdot}} & \textcolor{gray}{\ensuremath{\cdot}} & \textcolor{gray}{\ensuremath{\cdot}} & \textcolor{gray}{\ensuremath{\cdot}} & \textcolor{gray}{\ensuremath{\cdot}} & \textcolor{gray}{\ensuremath{\cdot}} & \textcolor{gray}{\ensuremath{\cdot}} & \textcolor{gray}{\ensuremath{\cdot}} & \textcolor{gray}{\ensuremath{\cdot}} & \textcolor{gray}{\ensuremath{\cdot}} & \textcolor{gray}{\ensuremath{\cdot}}\\[3pt]\hline
 & \begin{tabular}[t]{@{}l@{}}\rule{0pt}{12pt}MPS / Two-stack\\SIMD-DP / Block\end{tabular} & \hspace*{1pt}\begin{tabular}[t]{@{}c@{}}\rule{0pt}{12pt}6.8s\\[-.2pt]{\fontsize{5.8}{6.4}\selectfont \textcolor{ForestGreen}{$ \blacktriangle\,1.6\times $}}\\[-.2pt]6.1s\\[-.2pt]{\fontsize{5.8}{6.4}\selectfont \textcolor{ForestGreen}{$ \blacktriangle\,1.7\times $}}\end{tabular} & \hspace*{1pt}\begin{tabular}[t]{@{}c@{}}\rule{0pt}{12pt}6.8s\\[-.2pt]{\fontsize{5.8}{6.4}\selectfont \textcolor{ForestGreen}{$ \blacktriangle\,1.6\times $}}\\[-.2pt]6s\\[-.2pt]{\fontsize{5.8}{6.4}\selectfont \textcolor{ForestGreen}{$ \blacktriangle\,1.4\times $}}\end{tabular} & \hspace*{1pt}\begin{tabular}[t]{@{}c@{}}\rule{0pt}{12pt}12.6s\\[-.2pt]{\fontsize{5.8}{6.4}\selectfont \textcolor{ForestGreen}{$ \blacktriangle\,1.8\times $}}\\[-.2pt]10.6s\\[-.2pt]{\fontsize{5.8}{6.4}\selectfont \textcolor{ForestGreen}{$ \blacktriangle\,1.8\times $}}\end{tabular} & \hspace*{1pt}\begin{tabular}[t]{@{}c@{}}\rule{0pt}{12pt}6.5s\\[-.2pt]{\fontsize{5.8}{6.4}\selectfont \textcolor{ForestGreen}{$ \blacktriangle\,11.2\times $}}\\[-.2pt]5.7s\\[-.2pt]{\fontsize{5.8}{6.4}\selectfont \textcolor{ForestGreen}{$ \blacktriangle\,11.6\times $}}\end{tabular} & \textcolor{gray}{\ensuremath{\cdot}} & \textcolor{gray}{\ensuremath{\cdot}} & \textcolor{gray}{\ensuremath{\cdot}} & \textcolor{gray}{\ensuremath{\cdot}} & \hspace*{1pt}\begin{tabular}[t]{@{}c@{}}\rule{0pt}{12pt}2m\,46.7s\\[-.2pt]{\fontsize{5.8}{6.4}\selectfont \textcolor{BrickRed}{$ \blacktriangledown\,3.1\times $}}\\[-.2pt]2m\,24.6s\\[-.2pt]{\fontsize{5.8}{6.4}\selectfont \textcolor{BrickRed}{$ \blacktriangledown\,2.8\times $}}\end{tabular} & \hspace*{1pt}\begin{tabular}[t]{@{}c@{}}\rule{0pt}{12pt}2m\,47.4s\\[-.2pt]{\fontsize{5.8}{6.4}\selectfont \textcolor{BrickRed}{$ \blacktriangledown\,3.0\times $}}\\[-.2pt]2m\,27.3s\\[-.2pt]{\fontsize{5.8}{6.4}\selectfont \textcolor{BrickRed}{$ \blacktriangledown\,2.7\times $}}\end{tabular} & \textcolor{gray}{\ensuremath{\cdot}} & \textcolor{gray}{\ensuremath{\cdot}} & \textcolor{gray}{\ensuremath{\cdot}} & \textcolor{gray}{\ensuremath{\cdot}} & \textcolor{gray}{\ensuremath{\cdot}} & \textcolor{gray}{\ensuremath{\cdot}}\\[3pt]\hline
 & \begin{tabular}[t]{@{}l@{}}\rule{0pt}{12pt}MPS / Two-stack\\SIMD-DP / KSW2\end{tabular} & \hspace*{1pt}\begin{tabular}[t]{@{}c@{}}\rule{0pt}{12pt}8.3s\\[-.2pt]{\fontsize{5.8}{6.4}\selectfont \textcolor{ForestGreen}{$ \blacktriangle\,26\% $}}\\[-.2pt]5.4s\\[-.2pt]{\fontsize{5.8}{6.4}\selectfont \textcolor{ForestGreen}{$ \blacktriangle\,1.8\times $}}\end{tabular} & \hspace*{1pt}\begin{tabular}[t]{@{}c@{}}\rule{0pt}{12pt}8.4s\\[-.2pt]{\fontsize{5.8}{6.4}\selectfont \textcolor{ForestGreen}{$ \blacktriangle\,21\% $}}\\[-.2pt]5.2s\\[-.2pt]{\fontsize{5.8}{6.4}\selectfont \textcolor{ForestGreen}{$ \blacktriangle\,1.7\times $}}\end{tabular} & \hspace*{1pt}\begin{tabular}[t]{@{}c@{}}\rule{0pt}{12pt}16.2s\\[-.2pt]{\fontsize{5.8}{6.4}\selectfont \textcolor{ForestGreen}{$ \blacktriangle\,28\% $}}\\[-.2pt]9.2s\\[-.2pt]{\fontsize{5.8}{6.4}\selectfont \textcolor{ForestGreen}{$ \blacktriangle\,2.0\times $}}\end{tabular} & \textcolor{gray}{\ensuremath{\cdot}} & \textcolor{gray}{\ensuremath{\cdot}} & \textcolor{gray}{\ensuremath{\cdot}} & \textcolor{gray}{\ensuremath{\cdot}} & \textcolor{gray}{\ensuremath{\cdot}} & \hspace*{1pt}\begin{tabular}[t]{@{}c@{}}\rule{0pt}{12pt}2m\,40.2s\\[-.2pt]{\fontsize{5.8}{6.4}\selectfont \textcolor{BrickRed}{$ \blacktriangledown\,3.0\times $}}\\[-.2pt]2m\,15.4s\\[-.2pt]{\fontsize{5.8}{6.4}\selectfont \textcolor{BrickRed}{$ \blacktriangledown\,2.6\times $}}\end{tabular} & \hspace*{1pt}\begin{tabular}[t]{@{}c@{}}\rule{0pt}{12pt}2m\,42.5s\\[-.2pt]{\fontsize{5.8}{6.4}\selectfont \textcolor{BrickRed}{$ \blacktriangledown\,2.9\times $}}\\[-.2pt]2m\,20.9s\\[-.2pt]{\fontsize{5.8}{6.4}\selectfont \textcolor{BrickRed}{$ \blacktriangledown\,2.6\times $}}\end{tabular} & \textcolor{gray}{\ensuremath{\cdot}} & \textcolor{gray}{\ensuremath{\cdot}} & \textcolor{gray}{\ensuremath{\cdot}} & \textcolor{gray}{\ensuremath{\cdot}} & \textcolor{gray}{\ensuremath{\cdot}} & \textcolor{gray}{\ensuremath{\cdot}}\\[3pt]\hline
\rule{0pt}{12pt}\textbf{BioseqDB} & Vulkan (all modes) & \textcolor{gray}{---\textsuperscript{g}} & \textcolor{gray}{---\textsuperscript{g}} & \textcolor{gray}{---\textsuperscript{g}} & \textcolor{gray}{---\textsuperscript{g}} & \textcolor{gray}{---\textsuperscript{g}} & \textcolor{gray}{---\textsuperscript{g}} & \textcolor{gray}{---\textsuperscript{g}} & \textcolor{gray}{---\textsuperscript{g}} & \textcolor{gray}{---\textsuperscript{g}} & \textcolor{gray}{---\textsuperscript{g}} & \textcolor{gray}{---\textsuperscript{g}} & \textcolor{gray}{---\textsuperscript{g}} & \textcolor{gray}{---\textsuperscript{g}} & \textcolor{gray}{---\textsuperscript{g}} & \textcolor{gray}{---\textsuperscript{g}} & \textcolor{gray}{---\textsuperscript{g}}\\[3pt]\hline
\end{tabular}\par}
\smallskip{\footnotesize\noindent Each cell: A, comparison, B, comparison. Times cover 1,024 queries. $\cdot$ = no approved observation. ---\textsuperscript{g} = GPU unavailable, not executed. n/a = no measured comparison.\par
\noindent \textsuperscript{f}SIMD-ref request observed as scalar, compared only as that effective path. \textsuperscript{p}Accepted by explicit preservation after wrapper postflight failure, no run receipt claimed.\par
\noindent MPF/MPS = minimizers per fragment/sequence. Row order: index representation / minimizer, then chaining / base alignment. SIMD-ref is the reference minimizer, Bounded is bounded ntHash.\par
\noindent Blue = one lowest A and one lowest B down each machine/workload column, across continuation pages. Exact ties select the first row. Repeated baselines are marked only on page 1.\par
\noindent Light gray = minimap2 baseline. Pale yellow = preferred MPF / Bounded / SIMD-DP / Block configuration, not a claim of winning every column.\par
\noindent When all Vulkan modes are unmeasured, they share one row. Baseline comparisons use the fastest measured BioseqDB effective path separately for A and B.\par}
\clearpage
\noindent{\large\bfseries AWS: 2 Gbp, k=20, w=15 (mapping interval)\hfill 1/6}\par\smallskip
{\fontsize{7.1}{8.1}\selectfont\setlength{\tabcolsep}{1pt}\setlength{\fboxsep}{2pt}\renewcommand{\arraystretch}{1.12}
\noindent\begin{tabular}{@{}p{14mm}p{27mm}|p{13.55mm}|p{13.55mm}|p{13.55mm}|p{13.55mm}|p{13.55mm}|p{13.55mm}|p{13.55mm}|p{13.55mm}|p{13.55mm}|p{13.55mm}|p{13.55mm}|p{13.55mm}|p{13.55mm}|p{13.55mm}|p{13.55mm}|p{13.55mm}|@{}}\hline
Engine & Execution & \multicolumn{8}{c|}{\textbf{Index}} & \multicolumn{8}{c|}{\textbf{No index}}\\
 &  & \multicolumn{8}{c|}{\textbf{2 Gbp}} & \multicolumn{8}{c|}{\textbf{2 Gbp}}\\
 &  & \multicolumn{8}{c|}{\textbf{k=20, w=15}} & \multicolumn{8}{c|}{\textbf{k=20, w=15}}\\
 &  & \multicolumn{4}{c|}{\textbf{DNA4}} & \multicolumn{4}{c|}{\textbf{DNA5}} & \multicolumn{4}{c|}{\textbf{DNA4}} & \multicolumn{4}{c|}{\textbf{DNA5}}\\
 &  & \multicolumn{1}{c|}{\textbf{M7A}} & \multicolumn{1}{c|}{\textbf{M7I}} & \multicolumn{1}{c|}{\textbf{M4}} & \multicolumn{1}{c|}{\textbf{C7G}} & \multicolumn{1}{c|}{\textbf{M7A}} & \multicolumn{1}{c|}{\textbf{M7I}} & \multicolumn{1}{c|}{\textbf{M4}} & \multicolumn{1}{c|}{\textbf{C7G}} & \multicolumn{1}{c|}{\textbf{M7A}} & \multicolumn{1}{c|}{\textbf{M7I}} & \multicolumn{1}{c|}{\textbf{M4}} & \multicolumn{1}{c|}{\textbf{C7G}} & \multicolumn{1}{c|}{\textbf{M7A}} & \multicolumn{1}{c|}{\textbf{M7I}} & \multicolumn{1}{c|}{\textbf{M4}} & \multicolumn{1}{c|}{\textbf{C7G}}\\
\hline
\rowcolor{BenchBaseline}[0.5pt][0.5pt]\textbf{minimap2} & \textbf{Baseline} & \hspace*{1pt}\begin{tabular}[t]{@{}c@{}}\rule{0pt}{12pt}19.6s\\[-.2pt]{\fontsize{5.8}{6.4}\selectfont \textcolor{BrickRed}{$ \blacktriangledown\,6.7\times $}}\\[-.2pt]18.6s\\[-.2pt]{\fontsize{5.8}{6.4}\selectfont \textcolor{BrickRed}{$ \blacktriangledown\,8.4\times $}}\end{tabular} & \hspace*{1pt}\begin{tabular}[t]{@{}c@{}}\rule{0pt}{12pt}19.4s\\[-.2pt]{\fontsize{5.8}{6.4}\selectfont \textcolor{BrickRed}{$ \blacktriangledown\,7.2\times $}}\\[-.2pt]17.5s\\[-.2pt]{\fontsize{5.8}{6.4}\selectfont \textcolor{BrickRed}{$ \blacktriangledown\,9.1\times $}}\end{tabular} & \hspace*{1pt}\begin{tabular}[t]{@{}c@{}}\rule{0pt}{12pt}1m\,12.7s\\[-.2pt]{\fontsize{5.8}{6.4}\selectfont \textcolor{BrickRed}{$ \blacktriangledown\,13.6\times $}}\\[-.2pt]1m\,8.1s\\[-.2pt]{\fontsize{5.8}{6.4}\selectfont \textcolor{BrickRed}{$ \blacktriangledown\,20.2\times $}}\end{tabular} & \hspace*{1pt}\begin{tabular}[t]{@{}c@{}}\rule{0pt}{12pt}7m\,25.4s\\[-.2pt]{\fontsize{5.8}{6.4}\selectfont \textcolor{BrickRed}{$ \blacktriangledown\,106.2\times $}}\\[-.2pt]6m\,50.8s\\[-.2pt]{\fontsize{5.8}{6.4}\selectfont \textcolor{BrickRed}{$ \blacktriangledown\,212.9\times $}}\end{tabular} & \hspace*{1pt}\begin{tabular}[t]{@{}c@{}}\rule{0pt}{12pt}9.3s\\[-.2pt]{\fontsize{5.8}{6.4}\selectfont \textcolor{BrickRed}{$ \blacktriangledown\,24\% $}}\\[-.2pt]8.4s\\[-.2pt]{\fontsize{5.8}{6.4}\selectfont \textcolor{BrickRed}{$ \blacktriangledown\,1.6\times $}}\end{tabular} & \hspace*{1pt}\begin{tabular}[t]{@{}c@{}}\rule{0pt}{12pt}8.7s\\[-.2pt]{\fontsize{5.8}{6.4}\selectfont \textcolor{BrickRed}{$ \blacktriangledown\,1.5\times $}}\\[-.2pt]7s\\[-.2pt]{\fontsize{5.8}{6.4}\selectfont \textcolor{BrickRed}{$ \blacktriangledown\,1.7\times $}}\end{tabular} & \hspace*{1pt}\begin{tabular}[t]{@{}c@{}}\rule{0pt}{12pt}15.8s\\[-.2pt]{\fontsize{5.8}{6.4}\selectfont \textcolor{BrickRed}{$ \blacktriangledown\,26\% $}}\\[-.2pt]12.7s\\[-.2pt]{\fontsize{5.8}{6.4}\selectfont \textcolor{BrickRed}{$ \blacktriangledown\,1.6\times $}}\end{tabular} & \hspace*{1pt}\begin{tabular}[t]{@{}c@{}}\rule{0pt}{12pt}18.5s\\[-.2pt]{\fontsize{5.8}{6.4}\selectfont \textcolor{BrickRed}{$ \blacktriangledown\,1.6\times $}}\\[-.2pt]17s\\[-.2pt]{\fontsize{5.8}{6.4}\selectfont \textcolor{BrickRed}{$ \blacktriangledown\,1.8\times $}}\end{tabular} & \hspace*{1pt}\begin{tabular}[t]{@{}c@{}}\rule{0pt}{12pt}\colorbox{BenchBest}{\strut 42.4s}\\[-.2pt]{\fontsize{5.8}{6.4}\selectfont \textcolor{ForestGreen}{$ \blacktriangle\,2.0\times $}}\\[-.2pt]\colorbox{BenchBest}{\strut 41.4s}\\[-.2pt]{\fontsize{5.8}{6.4}\selectfont \textcolor{ForestGreen}{$ \blacktriangle\,1.9\times $}}\end{tabular} & \hspace*{1pt}\begin{tabular}[t]{@{}c@{}}\rule{0pt}{12pt}\colorbox{BenchBest}{\strut 42.9s}\\[-.2pt]{\fontsize{5.8}{6.4}\selectfont \textcolor{ForestGreen}{$ \blacktriangle\,2.0\times $}}\\[-.2pt]\colorbox{BenchBest}{\strut 41.1s}\\[-.2pt]{\fontsize{5.8}{6.4}\selectfont \textcolor{ForestGreen}{$ \blacktriangle\,1.8\times $}}\end{tabular} & \hspace*{1pt}\begin{tabular}[t]{@{}c@{}}\rule{0pt}{12pt}\colorbox{BenchBest}{\strut 1m\,26.9s}\\[-.2pt]{\fontsize{5.8}{6.4}\selectfont \textcolor{ForestGreen}{$ \blacktriangle\,1.9\times $}}\\[-.2pt]\colorbox{BenchBest}{\strut 1m\,23.7s}\\[-.2pt]{\fontsize{5.8}{6.4}\selectfont \textcolor{ForestGreen}{$ \blacktriangle\,1.7\times $}}\end{tabular} & \hspace*{1pt}\begin{tabular}[t]{@{}c@{}}\rule{0pt}{12pt}\colorbox{BenchBest}{\strut 4m\,51.6s}\\[-.2pt]{\fontsize{5.8}{6.4}\selectfont \textcolor{gray}{n/a}}\\[-.2pt]\colorbox{BenchBest}{\strut 4m\,46.6s}\\[-.2pt]{\fontsize{5.8}{6.4}\selectfont \textcolor{gray}{n/a}}\end{tabular} & \hspace*{1pt}\begin{tabular}[t]{@{}c@{}}\rule{0pt}{12pt}\colorbox{BenchBest}{\strut 35s}\\[-.2pt]{\fontsize{5.8}{6.4}\selectfont \textcolor{ForestGreen}{$ \blacktriangle\,8.6\times $}}\\[-.2pt]\colorbox{BenchBest}{\strut 34s}\\[-.2pt]{\fontsize{5.8}{6.4}\selectfont \textcolor{ForestGreen}{$ \blacktriangle\,2.5\times $}}\end{tabular} & \hspace*{1pt}\begin{tabular}[t]{@{}c@{}}\rule{0pt}{12pt}\colorbox{BenchBest}{\strut 35.9s}\\[-.2pt]{\fontsize{5.8}{6.4}\selectfont \textcolor{ForestGreen}{$ \blacktriangle\,6.7\times $}}\\[-.2pt]\colorbox{BenchBest}{\strut 34.2s}\\[-.2pt]{\fontsize{5.8}{6.4}\selectfont \textcolor{ForestGreen}{$ \blacktriangle\,2.5\times $}}\end{tabular} & \hspace*{1pt}\begin{tabular}[t]{@{}c@{}}\rule{0pt}{12pt}\colorbox{BenchBest}{\strut 1m\,14.4s}\\[-.2pt]{\fontsize{5.8}{6.4}\selectfont \textcolor{ForestGreen}{$ \blacktriangle\,6.9\times $}}\\[-.2pt]\colorbox{BenchBest}{\strut 1m\,11.3s}\\[-.2pt]{\fontsize{5.8}{6.4}\selectfont \textcolor{ForestGreen}{$ \blacktriangle\,2.2\times $}}\end{tabular} & \hspace*{1pt}\begin{tabular}[t]{@{}c@{}}\rule{0pt}{12pt}\colorbox{BenchBest}{\strut 1m\,0.1s}\\[-.2pt]{\fontsize{5.8}{6.4}\selectfont \textcolor{gray}{n/a}}\\[-.2pt]\colorbox{BenchBest}{\strut 58.6s}\\[-.2pt]{\fontsize{5.8}{6.4}\selectfont \textcolor{gray}{n/a}}\end{tabular}\\[3pt]\hline
\textbf{BioseqDB} & \begin{tabular}[t]{@{}l@{}}\rule{0pt}{12pt}MPF / Bounded\\RMQ / Block\end{tabular} & \hspace*{1pt}\begin{tabular}[t]{@{}c@{}}\rule{0pt}{12pt}\colorbox{BenchBest}{\strut 2.9s}\\[-.2pt]{\fontsize{5.8}{6.4}\selectfont \textcolor{ForestGreen}{$ \blacktriangle\,6.7\times $}}\\[-.2pt]2.2s\\[-.2pt]{\fontsize{5.8}{6.4}\selectfont \textcolor{ForestGreen}{$ \blacktriangle\,8.3\times $}}\end{tabular} & \hspace*{1pt}\begin{tabular}[t]{@{}c@{}}\rule{0pt}{12pt}2.7s\\[-.2pt]{\fontsize{5.8}{6.4}\selectfont \textcolor{ForestGreen}{$ \blacktriangle\,7.1\times $}}\\[-.2pt]2s\\[-.2pt]{\fontsize{5.8}{6.4}\selectfont \textcolor{ForestGreen}{$ \blacktriangle\,8.8\times $}}\end{tabular} & \textcolor{gray}{\ensuremath{\cdot}} & \hspace*{1pt}\begin{tabular}[t]{@{}c@{}}\rule{0pt}{12pt}4.7s\\[-.2pt]{\fontsize{5.8}{6.4}\selectfont \textcolor{ForestGreen}{$ \blacktriangle\,94.5\times $}}\\[-.2pt]3.9s\\[-.2pt]{\fontsize{5.8}{6.4}\selectfont \textcolor{ForestGreen}{$ \blacktriangle\,106.2\times $}}\end{tabular} & \textcolor{gray}{\ensuremath{\cdot}} & \textcolor{gray}{\ensuremath{\cdot}} & \textcolor{gray}{\ensuremath{\cdot}} & \hspace*{1pt}\begin{tabular}[t]{@{}c@{}}\rule{0pt}{12pt}\colorbox{BenchBest}{\strut 11.6s}\\[-.2pt]{\fontsize{5.8}{6.4}\selectfont \textcolor{ForestGreen}{$ \blacktriangle\,1.6\times $}}\\[-.2pt]\colorbox{BenchBest}{\strut 9.4s}\\[-.2pt]{\fontsize{5.8}{6.4}\selectfont \textcolor{ForestGreen}{$ \blacktriangle\,1.8\times $}}\end{tabular} & \hspace*{1pt}\begin{tabular}[t]{@{}c@{}}\rule{0pt}{12pt}1m\,34s\\[-.2pt]{\fontsize{5.8}{6.4}\selectfont \textcolor{BrickRed}{$ \blacktriangledown\,2.2\times $}}\\[-.2pt]1m\,24.9s\\[-.2pt]{\fontsize{5.8}{6.4}\selectfont \textcolor{BrickRed}{$ \blacktriangledown\,2.1\times $}}\end{tabular} & \hspace*{1pt}\begin{tabular}[t]{@{}c@{}}\rule{0pt}{12pt}1m\,33s\\[-.2pt]{\fontsize{5.8}{6.4}\selectfont \textcolor{BrickRed}{$ \blacktriangledown\,2.2\times $}}\\[-.2pt]1m\,22.3s\\[-.2pt]{\fontsize{5.8}{6.4}\selectfont \textcolor{BrickRed}{$ \blacktriangledown\,2.0\times $}}\end{tabular} & \textcolor{gray}{\ensuremath{\cdot}} & \textcolor{gray}{\ensuremath{\cdot}} & \textcolor{gray}{\ensuremath{\cdot}} & \textcolor{gray}{\ensuremath{\cdot}} & \textcolor{gray}{\ensuremath{\cdot}} & \textcolor{gray}{\ensuremath{\cdot}}\\[3pt]\hline
 & \begin{tabular}[t]{@{}l@{}}\rule{0pt}{12pt}MPF / Bounded\\RMQ / KSW2\end{tabular} & \hspace*{1pt}\begin{tabular}[t]{@{}c@{}}\rule{0pt}{12pt}4.5s\\[-.2pt]{\fontsize{5.8}{6.4}\selectfont \textcolor{ForestGreen}{$ \blacktriangle\,4.3\times $}}\\[-.2pt]\colorbox{BenchBest}{\strut 2.2s}\\[-.2pt]{\fontsize{5.8}{6.4}\selectfont \textcolor{ForestGreen}{$ \blacktriangle\,8.4\times $}}\end{tabular} & \hspace*{1pt}\begin{tabular}[t]{@{}c@{}}\rule{0pt}{12pt}4.2s\\[-.2pt]{\fontsize{5.8}{6.4}\selectfont \textcolor{ForestGreen}{$ \blacktriangle\,4.6\times $}}\\[-.2pt]2s\\[-.2pt]{\fontsize{5.8}{6.4}\selectfont \textcolor{ForestGreen}{$ \blacktriangle\,8.9\times $}}\end{tabular} & \textcolor{gray}{\ensuremath{\cdot}} & \hspace*{1pt}\begin{tabular}[t]{@{}c@{}}\rule{0pt}{12pt}6.2s\\[-.2pt]{\fontsize{5.8}{6.4}\selectfont \textcolor{ForestGreen}{$ \blacktriangle\,72.1\times $}}\\[-.2pt]3.4s\\[-.2pt]{\fontsize{5.8}{6.4}\selectfont \textcolor{ForestGreen}{$ \blacktriangle\,120.4\times $}}\end{tabular} & \textcolor{gray}{\ensuremath{\cdot}} & \textcolor{gray}{\ensuremath{\cdot}} & \textcolor{gray}{\ensuremath{\cdot}} & \hspace*{1pt}\begin{tabular}[t]{@{}c@{}}\rule{0pt}{12pt}12.6s\\[-.2pt]{\fontsize{5.8}{6.4}\selectfont \textcolor{ForestGreen}{$ \blacktriangle\,1.5\times $}}\\[-.2pt]9.8s\\[-.2pt]{\fontsize{5.8}{6.4}\selectfont \textcolor{ForestGreen}{$ \blacktriangle\,1.7\times $}}\end{tabular} & \hspace*{1pt}\begin{tabular}[t]{@{}c@{}}\rule{0pt}{12pt}1m\,32.9s\\[-.2pt]{\fontsize{5.8}{6.4}\selectfont \textcolor{BrickRed}{$ \blacktriangledown\,2.2\times $}}\\[-.2pt]1m\,23.1s\\[-.2pt]{\fontsize{5.8}{6.4}\selectfont \textcolor{BrickRed}{$ \blacktriangledown\,2.0\times $}}\end{tabular} & \hspace*{1pt}\begin{tabular}[t]{@{}c@{}}\rule{0pt}{12pt}1m\,41.2s\\[-.2pt]{\fontsize{5.8}{6.4}\selectfont \textcolor{BrickRed}{$ \blacktriangledown\,2.4\times $}}\\[-.2pt]1m\,28.4s\\[-.2pt]{\fontsize{5.8}{6.4}\selectfont \textcolor{BrickRed}{$ \blacktriangledown\,2.2\times $}}\end{tabular} & \textcolor{gray}{\ensuremath{\cdot}} & \textcolor{gray}{\ensuremath{\cdot}} & \textcolor{gray}{\ensuremath{\cdot}} & \textcolor{gray}{\ensuremath{\cdot}} & \textcolor{gray}{\ensuremath{\cdot}} & \textcolor{gray}{\ensuremath{\cdot}}\\[3pt]\hline
 & \begin{tabular}[t]{@{}l@{}}\rule{0pt}{12pt}MPF / Bounded\\Scalar-DP / Block\end{tabular} & \hspace*{1pt}\begin{tabular}[t]{@{}c@{}}\rule{0pt}{12pt}3s\\[-.2pt]{\fontsize{5.8}{6.4}\selectfont \textcolor{ForestGreen}{$ \blacktriangle\,6.6\times $}}\\[-.2pt]2.3s\\[-.2pt]{\fontsize{5.8}{6.4}\selectfont \textcolor{ForestGreen}{$ \blacktriangle\,8.1\times $}}\end{tabular} & \hspace*{1pt}\begin{tabular}[t]{@{}c@{}}\rule{0pt}{12pt}2.8s\\[-.2pt]{\fontsize{5.8}{6.4}\selectfont \textcolor{ForestGreen}{$ \blacktriangle\,6.9\times $}}\\[-.2pt]2s\\[-.2pt]{\fontsize{5.8}{6.4}\selectfont \textcolor{ForestGreen}{$ \blacktriangle\,8.7\times $}}\end{tabular} & \textcolor{gray}{\ensuremath{\cdot}} & \hspace*{1pt}\begin{tabular}[t]{@{}c@{}}\rule{0pt}{12pt}\colorbox{BenchBest}{\strut 4.2s}\\[-.2pt]{\fontsize{5.8}{6.4}\selectfont \textcolor{ForestGreen}{$ \blacktriangle\,106.2\times $}}\\[-.2pt]3.3s\\[-.2pt]{\fontsize{5.8}{6.4}\selectfont \textcolor{ForestGreen}{$ \blacktriangle\,122.9\times $}}\end{tabular} & \textcolor{gray}{\ensuremath{\cdot}} & \textcolor{gray}{\ensuremath{\cdot}} & \textcolor{gray}{\ensuremath{\cdot}} & \hspace*{1pt}\begin{tabular}[t]{@{}c@{}}\rule{0pt}{12pt}12.8s\\[-.2pt]{\fontsize{5.8}{6.4}\selectfont \textcolor{ForestGreen}{$ \blacktriangle\,1.4\times $}}\\[-.2pt]10.6s\\[-.2pt]{\fontsize{5.8}{6.4}\selectfont \textcolor{ForestGreen}{$ \blacktriangle\,1.6\times $}}\end{tabular} & \hspace*{1pt}\begin{tabular}[t]{@{}c@{}}\rule{0pt}{12pt}1m\,32.9s\\[-.2pt]{\fontsize{5.8}{6.4}\selectfont \textcolor{BrickRed}{$ \blacktriangledown\,2.2\times $}}\\[-.2pt]1m\,23.4s\\[-.2pt]{\fontsize{5.8}{6.4}\selectfont \textcolor{BrickRed}{$ \blacktriangledown\,2.0\times $}}\end{tabular} & \hspace*{1pt}\begin{tabular}[t]{@{}c@{}}\rule{0pt}{12pt}1m\,31.2s\\[-.2pt]{\fontsize{5.8}{6.4}\selectfont \textcolor{BrickRed}{$ \blacktriangledown\,2.1\times $}}\\[-.2pt]1m\,21s\\[-.2pt]{\fontsize{5.8}{6.4}\selectfont \textcolor{BrickRed}{$ \blacktriangledown\,2.0\times $}}\end{tabular} & \textcolor{gray}{\ensuremath{\cdot}} & \textcolor{gray}{\ensuremath{\cdot}} & \textcolor{gray}{\ensuremath{\cdot}} & \textcolor{gray}{\ensuremath{\cdot}} & \textcolor{gray}{\ensuremath{\cdot}} & \textcolor{gray}{\ensuremath{\cdot}}\\[3pt]\hline
 & \begin{tabular}[t]{@{}l@{}}\rule{0pt}{12pt}MPF / Bounded\\Scalar-DP / KSW2\end{tabular} & \hspace*{1pt}\begin{tabular}[t]{@{}c@{}}\rule{0pt}{12pt}4.6s\\[-.2pt]{\fontsize{5.8}{6.4}\selectfont \textcolor{ForestGreen}{$ \blacktriangle\,4.3\times $}}\\[-.2pt]2.2s\\[-.2pt]{\fontsize{5.8}{6.4}\selectfont \textcolor{ForestGreen}{$ \blacktriangle\,8.3\times $}}\end{tabular} & \hspace*{1pt}\begin{tabular}[t]{@{}c@{}}\rule{0pt}{12pt}4.1s\\[-.2pt]{\fontsize{5.8}{6.4}\selectfont \textcolor{ForestGreen}{$ \blacktriangle\,4.7\times $}}\\[-.2pt]\colorbox{BenchBest}{\strut 1.9s}\\[-.2pt]{\fontsize{5.8}{6.4}\selectfont \textcolor{ForestGreen}{$ \blacktriangle\,9.1\times $}}\end{tabular} & \textcolor{gray}{\ensuremath{\cdot}} & \hspace*{1pt}\begin{tabular}[t]{@{}c@{}}\rule{0pt}{12pt}5.8s\\[-.2pt]{\fontsize{5.8}{6.4}\selectfont \textcolor{ForestGreen}{$ \blacktriangle\,76.9\times $}}\\[-.2pt]3s\\[-.2pt]{\fontsize{5.8}{6.4}\selectfont \textcolor{ForestGreen}{$ \blacktriangle\,137.6\times $}}\end{tabular} & \textcolor{gray}{\ensuremath{\cdot}} & \textcolor{gray}{\ensuremath{\cdot}} & \textcolor{gray}{\ensuremath{\cdot}} & \hspace*{1pt}\begin{tabular}[t]{@{}c@{}}\rule{0pt}{12pt}13.3s\\[-.2pt]{\fontsize{5.8}{6.4}\selectfont \textcolor{ForestGreen}{$ \blacktriangle\,28\% $}}\\[-.2pt]10.5s\\[-.2pt]{\fontsize{5.8}{6.4}\selectfont \textcolor{ForestGreen}{$ \blacktriangle\,1.6\times $}}\end{tabular} & \hspace*{1pt}\begin{tabular}[t]{@{}c@{}}\rule{0pt}{12pt}1m\,27.2s\\[-.2pt]{\fontsize{5.8}{6.4}\selectfont \textcolor{BrickRed}{$ \blacktriangledown\,2.1\times $}}\\[-.2pt]1m\,17s\\[-.2pt]{\fontsize{5.8}{6.4}\selectfont \textcolor{BrickRed}{$ \blacktriangledown\,1.9\times $}}\end{tabular} & \hspace*{1pt}\begin{tabular}[t]{@{}c@{}}\rule{0pt}{12pt}1m\,39.6s\\[-.2pt]{\fontsize{5.8}{6.4}\selectfont \textcolor{BrickRed}{$ \blacktriangledown\,2.3\times $}}\\[-.2pt]1m\,26.7s\\[-.2pt]{\fontsize{5.8}{6.4}\selectfont \textcolor{BrickRed}{$ \blacktriangledown\,2.1\times $}}\end{tabular} & \textcolor{gray}{\ensuremath{\cdot}} & \textcolor{gray}{\ensuremath{\cdot}} & \textcolor{gray}{\ensuremath{\cdot}} & \textcolor{gray}{\ensuremath{\cdot}} & \textcolor{gray}{\ensuremath{\cdot}} & \textcolor{gray}{\ensuremath{\cdot}}\\[3pt]\hline
\rowcolor{BenchPreferred}[0.5pt][0.5pt] & \begin{tabular}[t]{@{}l@{}}\rule{0pt}{12pt}MPF / Bounded\\SIMD-DP / Block\end{tabular} & \hspace*{1pt}\begin{tabular}[t]{@{}c@{}}\rule{0pt}{12pt}3s\\[-.2pt]{\fontsize{5.8}{6.4}\selectfont \textcolor{ForestGreen}{$ \blacktriangle\,6.5\times $}}\\[-.2pt]2.3s\\[-.2pt]{\fontsize{5.8}{6.4}\selectfont \textcolor{ForestGreen}{$ \blacktriangle\,8.1\times $}}\end{tabular} & \hspace*{1pt}\begin{tabular}[t]{@{}c@{}}\rule{0pt}{12pt}\colorbox{BenchBest}{\strut 2.7s}\\[-.2pt]{\fontsize{5.8}{6.4}\selectfont \textcolor{ForestGreen}{$ \blacktriangle\,7.2\times $}}\\[-.2pt]2s\\[-.2pt]{\fontsize{5.8}{6.4}\selectfont \textcolor{ForestGreen}{$ \blacktriangle\,8.9\times $}}\end{tabular} & \hspace*{1pt}\begin{tabular}[t]{@{}c@{}}\rule{0pt}{12pt}\colorbox{BenchBest}{\strut 5.4s}\\[-.2pt]{\fontsize{5.8}{6.4}\selectfont \textcolor{ForestGreen}{$ \blacktriangle\,13.6\times $}}\\[-.2pt]\colorbox{BenchBest}{\strut 3.4s}\\[-.2pt]{\fontsize{5.8}{6.4}\selectfont \textcolor{ForestGreen}{$ \blacktriangle\,20.2\times $}}\end{tabular} & \hspace*{1pt}\begin{tabular}[t]{@{}c@{}}\rule{0pt}{12pt}4.5s\\[-.2pt]{\fontsize{5.8}{6.4}\selectfont \textcolor{ForestGreen}{$ \blacktriangle\,99.2\times $}}\\[-.2pt]3.6s\\[-.2pt]{\fontsize{5.8}{6.4}\selectfont \textcolor{ForestGreen}{$ \blacktriangle\,112.7\times $}}\end{tabular} & \hspace*{1pt}\begin{tabular}[t]{@{}c@{}}\rule{0pt}{12pt}\colorbox{BenchBest}{\strut 7.5s}\\[-.2pt]{\fontsize{5.8}{6.4}\selectfont \textcolor{ForestGreen}{$ \blacktriangle\,19\% $}}\\[-.2pt]\colorbox{BenchBest}{\strut 5.2s}\\[-.2pt]{\fontsize{5.8}{6.4}\selectfont \textcolor{ForestGreen}{$ \blacktriangle\,1.6\times $}}\end{tabular} & \hspace*{1pt}\begin{tabular}[t]{@{}c@{}}\rule{0pt}{12pt}\colorbox{BenchBest}{\strut 5.8s}\\[-.2pt]{\fontsize{5.8}{6.4}\selectfont \textcolor{ForestGreen}{$ \blacktriangle\,1.5\times $}}\\[-.2pt]\colorbox{BenchBest}{\strut 4s}\\[-.2pt]{\fontsize{5.8}{6.4}\selectfont \textcolor{ForestGreen}{$ \blacktriangle\,1.7\times $}}\end{tabular} & \hspace*{1pt}\begin{tabular}[t]{@{}c@{}}\rule{0pt}{12pt}\colorbox{BenchBest}{\strut 12.6s}\\[-.2pt]{\fontsize{5.8}{6.4}\selectfont \textcolor{ForestGreen}{$ \blacktriangle\,21\% $}}\\[-.2pt]\colorbox{BenchBest}{\strut 7.9s}\\[-.2pt]{\fontsize{5.8}{6.4}\selectfont \textcolor{ForestGreen}{$ \blacktriangle\,1.6\times $}}\end{tabular} & \hspace*{1pt}\begin{tabular}[t]{@{}c@{}}\rule{0pt}{12pt}12.9s\\[-.2pt]{\fontsize{5.8}{6.4}\selectfont \textcolor{ForestGreen}{$ \blacktriangle\,1.4\times $}}\\[-.2pt]10.6s\\[-.2pt]{\fontsize{5.8}{6.4}\selectfont \textcolor{ForestGreen}{$ \blacktriangle\,1.6\times $}}\end{tabular} & \hspace*{1pt}\begin{tabular}[t]{@{}c@{}}\rule{0pt}{12pt}1m\,24.8s\\[-.2pt]{\fontsize{5.8}{6.4}\selectfont \textcolor{BrickRed}{$ \blacktriangledown\,2.0\times $}}\\[-.2pt]1m\,16.9s\\[-.2pt]{\fontsize{5.8}{6.4}\selectfont \textcolor{BrickRed}{$ \blacktriangledown\,1.9\times $}}\end{tabular} & \hspace*{1pt}\begin{tabular}[t]{@{}c@{}}\rule{0pt}{12pt}1m\,25.4s\\[-.2pt]{\fontsize{5.8}{6.4}\selectfont \textcolor{BrickRed}{$ \blacktriangledown\,2.0\times $}}\\[-.2pt]1m\,15.7s\\[-.2pt]{\fontsize{5.8}{6.4}\selectfont \textcolor{BrickRed}{$ \blacktriangledown\,1.8\times $}}\end{tabular} & \hspace*{1pt}\begin{tabular}[t]{@{}c@{}}\rule{0pt}{12pt}2m\,44s\\[-.2pt]{\fontsize{5.8}{6.4}\selectfont \textcolor{BrickRed}{$ \blacktriangledown\,1.9\times $}}\\[-.2pt]2m\,26.2s\\[-.2pt]{\fontsize{5.8}{6.4}\selectfont \textcolor{BrickRed}{$ \blacktriangledown\,1.7\times $}}\end{tabular} & \textcolor{gray}{\ensuremath{\cdot}} & \hspace*{1pt}\begin{tabular}[t]{@{}c@{}}\rule{0pt}{12pt}5m\,1.5s\\[-.2pt]{\fontsize{5.8}{6.4}\selectfont \textcolor{BrickRed}{$ \blacktriangledown\,8.6\times $}}\\[-.2pt]1m\,24.2s\\[-.2pt]{\fontsize{5.8}{6.4}\selectfont \textcolor{BrickRed}{$ \blacktriangledown\,2.5\times $}}\end{tabular} & \hspace*{1pt}\begin{tabular}[t]{@{}c@{}}\rule{0pt}{12pt}4m\,0.6s\\[-.2pt]{\fontsize{5.8}{6.4}\selectfont \textcolor{BrickRed}{$ \blacktriangledown\,6.7\times $}}\\[-.2pt]1m\,24.3s\\[-.2pt]{\fontsize{5.8}{6.4}\selectfont \textcolor{BrickRed}{$ \blacktriangledown\,2.5\times $}}\end{tabular} & \hspace*{1pt}\begin{tabular}[t]{@{}c@{}}\rule{0pt}{12pt}8m\,31.7s\\[-.2pt]{\fontsize{5.8}{6.4}\selectfont \textcolor{BrickRed}{$ \blacktriangledown\,6.9\times $}}\\[-.2pt]2m\,34.1s\\[-.2pt]{\fontsize{5.8}{6.4}\selectfont \textcolor{BrickRed}{$ \blacktriangledown\,2.2\times $}}\end{tabular} & \textcolor{gray}{\ensuremath{\cdot}}\\[3pt]\hline
 & \begin{tabular}[t]{@{}l@{}}\rule{0pt}{12pt}MPF / Bounded\\SIMD-DP / KSW2\end{tabular} & \hspace*{1pt}\begin{tabular}[t]{@{}c@{}}\rule{0pt}{12pt}4.6s\\[-.2pt]{\fontsize{5.8}{6.4}\selectfont \textcolor{ForestGreen}{$ \blacktriangle\,4.3\times $}}\\[-.2pt]2.2s\\[-.2pt]{\fontsize{5.8}{6.4}\selectfont \textcolor{ForestGreen}{$ \blacktriangle\,8.3\times $}}\end{tabular} & \hspace*{1pt}\begin{tabular}[t]{@{}c@{}}\rule{0pt}{12pt}4.1s\\[-.2pt]{\fontsize{5.8}{6.4}\selectfont \textcolor{ForestGreen}{$ \blacktriangle\,4.8\times $}}\\[-.2pt]1.9s\\[-.2pt]{\fontsize{5.8}{6.4}\selectfont \textcolor{ForestGreen}{$ \blacktriangle\,9.0\times $}}\end{tabular} & \textcolor{gray}{\ensuremath{\cdot}} & \hspace*{1pt}\begin{tabular}[t]{@{}c@{}}\rule{0pt}{12pt}4.7s\\[-.2pt]{\fontsize{5.8}{6.4}\selectfont \textcolor{ForestGreen}{$ \blacktriangle\,95.2\times $}}\\[-.2pt]\colorbox{BenchBest}{\strut 1.9s}\\[-.2pt]{\fontsize{5.8}{6.4}\selectfont \textcolor{ForestGreen}{$ \blacktriangle\,212.9\times $}}\end{tabular} & \textcolor{gray}{\ensuremath{\cdot}} & \textcolor{gray}{\ensuremath{\cdot}} & \textcolor{gray}{\ensuremath{\cdot}} & \hspace*{1pt}\begin{tabular}[t]{@{}c@{}}\rule{0pt}{12pt}12.5s\\[-.2pt]{\fontsize{5.8}{6.4}\selectfont \textcolor{ForestGreen}{$ \blacktriangle\,1.5\times $}}\\[-.2pt]9.7s\\[-.2pt]{\fontsize{5.8}{6.4}\selectfont \textcolor{ForestGreen}{$ \blacktriangle\,1.8\times $}}\end{tabular} & \hspace*{1pt}\begin{tabular}[t]{@{}c@{}}\rule{0pt}{12pt}1m\,34.1s\\[-.2pt]{\fontsize{5.8}{6.4}\selectfont \textcolor{BrickRed}{$ \blacktriangledown\,2.2\times $}}\\[-.2pt]1m\,24s\\[-.2pt]{\fontsize{5.8}{6.4}\selectfont \textcolor{BrickRed}{$ \blacktriangledown\,2.0\times $}}\end{tabular} & \hspace*{1pt}\begin{tabular}[t]{@{}c@{}}\rule{0pt}{12pt}1m\,32.2s\\[-.2pt]{\fontsize{5.8}{6.4}\selectfont \textcolor{BrickRed}{$ \blacktriangledown\,2.2\times $}}\\[-.2pt]1m\,20s\\[-.2pt]{\fontsize{5.8}{6.4}\selectfont \textcolor{BrickRed}{$ \blacktriangledown\,1.9\times $}}\end{tabular} & \textcolor{gray}{\ensuremath{\cdot}} & \textcolor{gray}{\ensuremath{\cdot}} & \textcolor{gray}{\ensuremath{\cdot}} & \textcolor{gray}{\ensuremath{\cdot}} & \textcolor{gray}{\ensuremath{\cdot}} & \textcolor{gray}{\ensuremath{\cdot}}\\[3pt]\hline
 & \begin{tabular}[t]{@{}l@{}}\rule{0pt}{12pt}MPF / Scalar\\RMQ / Block\end{tabular} & \hspace*{1pt}\begin{tabular}[t]{@{}c@{}}\rule{0pt}{12pt}5m\,39.9s\\[-.2pt]{\fontsize{5.8}{6.4}\selectfont \textcolor{BrickRed}{$ \blacktriangledown\,17.3\times $}}\\[-.2pt]5m\,39.5s\\[-.2pt]{\fontsize{5.8}{6.4}\selectfont \textcolor{BrickRed}{$ \blacktriangledown\,18.2\times $}}\end{tabular} & \hspace*{1pt}\begin{tabular}[t]{@{}c@{}}\rule{0pt}{12pt}12m\,19.6s\\[-.2pt]{\fontsize{5.8}{6.4}\selectfont \textcolor{BrickRed}{$ \blacktriangledown\,38.2\times $}}\\[-.2pt]12m\,19.1s\\[-.2pt]{\fontsize{5.8}{6.4}\selectfont \textcolor{BrickRed}{$ \blacktriangledown\,42.3\times $}}\end{tabular} & \hspace*{1pt}\begin{tabular}[t]{@{}c@{}}\rule{0pt}{12pt}22m\,41.7s\\[-.2pt]{\fontsize{5.8}{6.4}\selectfont \textcolor{BrickRed}{$ \blacktriangledown\,18.7\times $}}\\[-.2pt]22m\,40.7s\\[-.2pt]{\fontsize{5.8}{6.4}\selectfont \textcolor{BrickRed}{$ \blacktriangledown\,20.0\times $}}\end{tabular} & \hspace*{1pt}\begin{tabular}[t]{@{}c@{}}\rule{0pt}{12pt}12m\,57s\\[-.2pt]{\fontsize{5.8}{6.4}\selectfont \textcolor{BrickRed}{$ \blacktriangledown\,1.7\times $}}\\[-.2pt]12m\,56.4s\\[-.2pt]{\fontsize{5.8}{6.4}\selectfont \textcolor{BrickRed}{$ \blacktriangledown\,1.9\times $}}\end{tabular} & \textcolor{gray}{\ensuremath{\cdot}} & \textcolor{gray}{\ensuremath{\cdot}} & \textcolor{gray}{\ensuremath{\cdot}} & \textcolor{gray}{\ensuremath{\cdot}} & \hspace*{1pt}\begin{tabular}[t]{@{}c@{}}\rule{0pt}{12pt}2m\,49.9s\\[-.2pt]{\fontsize{5.8}{6.4}\selectfont \textcolor{BrickRed}{$ \blacktriangledown\,4.0\times $}}\\[-.2pt]2m\,13.4s\\[-.2pt]{\fontsize{5.8}{6.4}\selectfont \textcolor{BrickRed}{$ \blacktriangledown\,3.2\times $}}\end{tabular} & \hspace*{1pt}\begin{tabular}[t]{@{}c@{}}\rule{0pt}{12pt}2m\,29.9s\\[-.2pt]{\fontsize{5.8}{6.4}\selectfont \textcolor{BrickRed}{$ \blacktriangledown\,3.5\times $}}\\[-.2pt]2m\,1.1s\\[-.2pt]{\fontsize{5.8}{6.4}\selectfont \textcolor{BrickRed}{$ \blacktriangledown\,2.9\times $}}\end{tabular} & \textcolor{gray}{\ensuremath{\cdot}} & \textcolor{gray}{\ensuremath{\cdot}} & \textcolor{gray}{\ensuremath{\cdot}} & \textcolor{gray}{\ensuremath{\cdot}} & \textcolor{gray}{\ensuremath{\cdot}} & \textcolor{gray}{\ensuremath{\cdot}}\\[3pt]\hline
\end{tabular}\par}
\smallskip{\footnotesize\noindent Each cell: A, comparison, B, comparison. Times cover 1,024 queries. $\cdot$ = no approved observation. ---\textsuperscript{g} = GPU unavailable, not executed. n/a = no measured comparison.\par
\noindent \textsuperscript{f}SIMD-ref request observed as scalar, compared only as that effective path. \textsuperscript{p}Accepted by explicit preservation after wrapper postflight failure, no run receipt claimed.\par
\noindent MPF/MPS = minimizers per fragment/sequence. Row order: index representation / minimizer, then chaining / base alignment. SIMD-ref is the reference minimizer, Bounded is bounded ntHash.\par
\noindent Blue = one lowest A and one lowest B down each machine/workload column, across continuation pages. Exact ties select the first row. Repeated baselines are marked only on page 1.\par
\noindent Light gray = minimap2 baseline. Pale yellow = preferred MPF / Bounded / SIMD-DP / Block configuration, not a claim of winning every column.\par
\noindent When all Vulkan modes are unmeasured, they share one row. Baseline comparisons use the fastest measured BioseqDB effective path separately for A and B.\par}
\clearpage
\noindent{\large\bfseries AWS: 2 Gbp, k=20, w=15 (mapping interval)\hfill 2/6}\par\smallskip
{\fontsize{7.1}{8.1}\selectfont\setlength{\tabcolsep}{1pt}\setlength{\fboxsep}{2pt}\renewcommand{\arraystretch}{1.12}
\noindent\begin{tabular}{@{}p{14mm}p{27mm}|p{13.55mm}|p{13.55mm}|p{13.55mm}|p{13.55mm}|p{13.55mm}|p{13.55mm}|p{13.55mm}|p{13.55mm}|p{13.55mm}|p{13.55mm}|p{13.55mm}|p{13.55mm}|p{13.55mm}|p{13.55mm}|p{13.55mm}|p{13.55mm}|@{}}\hline
Engine & Execution & \multicolumn{8}{c|}{\textbf{Index}} & \multicolumn{8}{c|}{\textbf{No index}}\\
 &  & \multicolumn{8}{c|}{\textbf{2 Gbp}} & \multicolumn{8}{c|}{\textbf{2 Gbp}}\\
 &  & \multicolumn{8}{c|}{\textbf{k=20, w=15}} & \multicolumn{8}{c|}{\textbf{k=20, w=15}}\\
 &  & \multicolumn{4}{c|}{\textbf{DNA4}} & \multicolumn{4}{c|}{\textbf{DNA5}} & \multicolumn{4}{c|}{\textbf{DNA4}} & \multicolumn{4}{c|}{\textbf{DNA5}}\\
 &  & \multicolumn{1}{c|}{\textbf{M7A}} & \multicolumn{1}{c|}{\textbf{M7I}} & \multicolumn{1}{c|}{\textbf{M4}} & \multicolumn{1}{c|}{\textbf{C7G}} & \multicolumn{1}{c|}{\textbf{M7A}} & \multicolumn{1}{c|}{\textbf{M7I}} & \multicolumn{1}{c|}{\textbf{M4}} & \multicolumn{1}{c|}{\textbf{C7G}} & \multicolumn{1}{c|}{\textbf{M7A}} & \multicolumn{1}{c|}{\textbf{M7I}} & \multicolumn{1}{c|}{\textbf{M4}} & \multicolumn{1}{c|}{\textbf{C7G}} & \multicolumn{1}{c|}{\textbf{M7A}} & \multicolumn{1}{c|}{\textbf{M7I}} & \multicolumn{1}{c|}{\textbf{M4}} & \multicolumn{1}{c|}{\textbf{C7G}}\\
\hline
\rowcolor{BenchBaseline}[0.5pt][0.5pt]\textbf{minimap2} & \textbf{Baseline} & \hspace*{1pt}\begin{tabular}[t]{@{}c@{}}\rule{0pt}{12pt}19.6s\\[-.2pt]{\fontsize{5.8}{6.4}\selectfont \textcolor{BrickRed}{$ \blacktriangledown\,6.7\times $}}\\[-.2pt]18.6s\\[-.2pt]{\fontsize{5.8}{6.4}\selectfont \textcolor{BrickRed}{$ \blacktriangledown\,8.4\times $}}\end{tabular} & \hspace*{1pt}\begin{tabular}[t]{@{}c@{}}\rule{0pt}{12pt}19.4s\\[-.2pt]{\fontsize{5.8}{6.4}\selectfont \textcolor{BrickRed}{$ \blacktriangledown\,7.2\times $}}\\[-.2pt]17.5s\\[-.2pt]{\fontsize{5.8}{6.4}\selectfont \textcolor{BrickRed}{$ \blacktriangledown\,9.1\times $}}\end{tabular} & \hspace*{1pt}\begin{tabular}[t]{@{}c@{}}\rule{0pt}{12pt}1m\,12.7s\\[-.2pt]{\fontsize{5.8}{6.4}\selectfont \textcolor{BrickRed}{$ \blacktriangledown\,13.6\times $}}\\[-.2pt]1m\,8.1s\\[-.2pt]{\fontsize{5.8}{6.4}\selectfont \textcolor{BrickRed}{$ \blacktriangledown\,20.2\times $}}\end{tabular} & \hspace*{1pt}\begin{tabular}[t]{@{}c@{}}\rule{0pt}{12pt}7m\,25.4s\\[-.2pt]{\fontsize{5.8}{6.4}\selectfont \textcolor{BrickRed}{$ \blacktriangledown\,106.2\times $}}\\[-.2pt]6m\,50.8s\\[-.2pt]{\fontsize{5.8}{6.4}\selectfont \textcolor{BrickRed}{$ \blacktriangledown\,212.9\times $}}\end{tabular} & \hspace*{1pt}\begin{tabular}[t]{@{}c@{}}\rule{0pt}{12pt}9.3s\\[-.2pt]{\fontsize{5.8}{6.4}\selectfont \textcolor{BrickRed}{$ \blacktriangledown\,24\% $}}\\[-.2pt]8.4s\\[-.2pt]{\fontsize{5.8}{6.4}\selectfont \textcolor{BrickRed}{$ \blacktriangledown\,1.6\times $}}\end{tabular} & \hspace*{1pt}\begin{tabular}[t]{@{}c@{}}\rule{0pt}{12pt}8.7s\\[-.2pt]{\fontsize{5.8}{6.4}\selectfont \textcolor{BrickRed}{$ \blacktriangledown\,1.5\times $}}\\[-.2pt]7s\\[-.2pt]{\fontsize{5.8}{6.4}\selectfont \textcolor{BrickRed}{$ \blacktriangledown\,1.7\times $}}\end{tabular} & \hspace*{1pt}\begin{tabular}[t]{@{}c@{}}\rule{0pt}{12pt}15.8s\\[-.2pt]{\fontsize{5.8}{6.4}\selectfont \textcolor{BrickRed}{$ \blacktriangledown\,26\% $}}\\[-.2pt]12.7s\\[-.2pt]{\fontsize{5.8}{6.4}\selectfont \textcolor{BrickRed}{$ \blacktriangledown\,1.6\times $}}\end{tabular} & \hspace*{1pt}\begin{tabular}[t]{@{}c@{}}\rule{0pt}{12pt}18.5s\\[-.2pt]{\fontsize{5.8}{6.4}\selectfont \textcolor{BrickRed}{$ \blacktriangledown\,1.6\times $}}\\[-.2pt]17s\\[-.2pt]{\fontsize{5.8}{6.4}\selectfont \textcolor{BrickRed}{$ \blacktriangledown\,1.8\times $}}\end{tabular} & \hspace*{1pt}\begin{tabular}[t]{@{}c@{}}\rule{0pt}{12pt}42.4s\\[-.2pt]{\fontsize{5.8}{6.4}\selectfont \textcolor{ForestGreen}{$ \blacktriangle\,2.0\times $}}\\[-.2pt]41.4s\\[-.2pt]{\fontsize{5.8}{6.4}\selectfont \textcolor{ForestGreen}{$ \blacktriangle\,1.9\times $}}\end{tabular} & \hspace*{1pt}\begin{tabular}[t]{@{}c@{}}\rule{0pt}{12pt}42.9s\\[-.2pt]{\fontsize{5.8}{6.4}\selectfont \textcolor{ForestGreen}{$ \blacktriangle\,2.0\times $}}\\[-.2pt]41.1s\\[-.2pt]{\fontsize{5.8}{6.4}\selectfont \textcolor{ForestGreen}{$ \blacktriangle\,1.8\times $}}\end{tabular} & \hspace*{1pt}\begin{tabular}[t]{@{}c@{}}\rule{0pt}{12pt}1m\,26.9s\\[-.2pt]{\fontsize{5.8}{6.4}\selectfont \textcolor{ForestGreen}{$ \blacktriangle\,1.9\times $}}\\[-.2pt]1m\,23.7s\\[-.2pt]{\fontsize{5.8}{6.4}\selectfont \textcolor{ForestGreen}{$ \blacktriangle\,1.7\times $}}\end{tabular} & \hspace*{1pt}\begin{tabular}[t]{@{}c@{}}\rule{0pt}{12pt}4m\,51.6s\\[-.2pt]{\fontsize{5.8}{6.4}\selectfont \textcolor{gray}{n/a}}\\[-.2pt]4m\,46.6s\\[-.2pt]{\fontsize{5.8}{6.4}\selectfont \textcolor{gray}{n/a}}\end{tabular} & \hspace*{1pt}\begin{tabular}[t]{@{}c@{}}\rule{0pt}{12pt}35s\\[-.2pt]{\fontsize{5.8}{6.4}\selectfont \textcolor{ForestGreen}{$ \blacktriangle\,8.6\times $}}\\[-.2pt]34s\\[-.2pt]{\fontsize{5.8}{6.4}\selectfont \textcolor{ForestGreen}{$ \blacktriangle\,2.5\times $}}\end{tabular} & \hspace*{1pt}\begin{tabular}[t]{@{}c@{}}\rule{0pt}{12pt}35.9s\\[-.2pt]{\fontsize{5.8}{6.4}\selectfont \textcolor{ForestGreen}{$ \blacktriangle\,6.7\times $}}\\[-.2pt]34.2s\\[-.2pt]{\fontsize{5.8}{6.4}\selectfont \textcolor{ForestGreen}{$ \blacktriangle\,2.5\times $}}\end{tabular} & \hspace*{1pt}\begin{tabular}[t]{@{}c@{}}\rule{0pt}{12pt}1m\,14.4s\\[-.2pt]{\fontsize{5.8}{6.4}\selectfont \textcolor{ForestGreen}{$ \blacktriangle\,6.9\times $}}\\[-.2pt]1m\,11.3s\\[-.2pt]{\fontsize{5.8}{6.4}\selectfont \textcolor{ForestGreen}{$ \blacktriangle\,2.2\times $}}\end{tabular} & \hspace*{1pt}\begin{tabular}[t]{@{}c@{}}\rule{0pt}{12pt}1m\,0.1s\\[-.2pt]{\fontsize{5.8}{6.4}\selectfont \textcolor{gray}{n/a}}\\[-.2pt]58.6s\\[-.2pt]{\fontsize{5.8}{6.4}\selectfont \textcolor{gray}{n/a}}\end{tabular}\\[3pt]\hline
\textbf{BioseqDB} & \begin{tabular}[t]{@{}l@{}}\rule{0pt}{12pt}MPF / Scalar\\RMQ / KSW2\end{tabular} & \hspace*{1pt}\begin{tabular}[t]{@{}c@{}}\rule{0pt}{12pt}5m\,46.6s\\[-.2pt]{\fontsize{5.8}{6.4}\selectfont \textcolor{BrickRed}{$ \blacktriangledown\,17.7\times $}}\\[-.2pt]5m\,44.5s\\[-.2pt]{\fontsize{5.8}{6.4}\selectfont \textcolor{BrickRed}{$ \blacktriangledown\,18.5\times $}}\end{tabular} & \hspace*{1pt}\begin{tabular}[t]{@{}c@{}}\rule{0pt}{12pt}12m\,26.1s\\[-.2pt]{\fontsize{5.8}{6.4}\selectfont \textcolor{BrickRed}{$ \blacktriangledown\,38.5\times $}}\\[-.2pt]12m\,24.3s\\[-.2pt]{\fontsize{5.8}{6.4}\selectfont \textcolor{BrickRed}{$ \blacktriangledown\,42.6\times $}}\end{tabular} & \textcolor{gray}{\ensuremath{\cdot}} & \hspace*{1pt}\begin{tabular}[t]{@{}c@{}}\rule{0pt}{12pt}11m\,45.8s\\[-.2pt]{\fontsize{5.8}{6.4}\selectfont \textcolor{BrickRed}{$ \blacktriangledown\,1.6\times $}}\\[-.2pt]11m\,43.3s\\[-.2pt]{\fontsize{5.8}{6.4}\selectfont \textcolor{BrickRed}{$ \blacktriangledown\,1.7\times $}}\end{tabular} & \textcolor{gray}{\ensuremath{\cdot}} & \textcolor{gray}{\ensuremath{\cdot}} & \textcolor{gray}{\ensuremath{\cdot}} & \textcolor{gray}{\ensuremath{\cdot}} & \hspace*{1pt}\begin{tabular}[t]{@{}c@{}}\rule{0pt}{12pt}2m\,53.6s\\[-.2pt]{\fontsize{5.8}{6.4}\selectfont \textcolor{BrickRed}{$ \blacktriangledown\,4.1\times $}}\\[-.2pt]2m\,15.9s\\[-.2pt]{\fontsize{5.8}{6.4}\selectfont \textcolor{BrickRed}{$ \blacktriangledown\,3.3\times $}}\end{tabular} & \hspace*{1pt}\begin{tabular}[t]{@{}c@{}}\rule{0pt}{12pt}2m\,44s\\[-.2pt]{\fontsize{5.8}{6.4}\selectfont \textcolor{BrickRed}{$ \blacktriangledown\,3.8\times $}}\\[-.2pt]2m\,12.1s\\[-.2pt]{\fontsize{5.8}{6.4}\selectfont \textcolor{BrickRed}{$ \blacktriangledown\,3.2\times $}}\end{tabular} & \textcolor{gray}{\ensuremath{\cdot}} & \textcolor{gray}{\ensuremath{\cdot}} & \textcolor{gray}{\ensuremath{\cdot}} & \textcolor{gray}{\ensuremath{\cdot}} & \textcolor{gray}{\ensuremath{\cdot}} & \textcolor{gray}{\ensuremath{\cdot}}\\[3pt]\hline
 & \begin{tabular}[t]{@{}l@{}}\rule{0pt}{12pt}MPF / Scalar\\Scalar-DP / Block\end{tabular} & \hspace*{1pt}\begin{tabular}[t]{@{}c@{}}\rule{0pt}{12pt}5m\,39.5s\\[-.2pt]{\fontsize{5.8}{6.4}\selectfont \textcolor{BrickRed}{$ \blacktriangledown\,17.3\times $}}\\[-.2pt]5m\,39s\\[-.2pt]{\fontsize{5.8}{6.4}\selectfont \textcolor{BrickRed}{$ \blacktriangledown\,18.2\times $}}\end{tabular} & \hspace*{1pt}\begin{tabular}[t]{@{}c@{}}\rule{0pt}{12pt}12m\,20.4s\\[-.2pt]{\fontsize{5.8}{6.4}\selectfont \textcolor{BrickRed}{$ \blacktriangledown\,38.2\times $}}\\[-.2pt]12m\,19.8s\\[-.2pt]{\fontsize{5.8}{6.4}\selectfont \textcolor{BrickRed}{$ \blacktriangledown\,42.3\times $}}\end{tabular} & \hspace*{1pt}\begin{tabular}[t]{@{}c@{}}\rule{0pt}{12pt}22m\,35.5s\\[-.2pt]{\fontsize{5.8}{6.4}\selectfont \textcolor{BrickRed}{$ \blacktriangledown\,18.7\times $}}\\[-.2pt]22m\,34.5s\\[-.2pt]{\fontsize{5.8}{6.4}\selectfont \textcolor{BrickRed}{$ \blacktriangledown\,19.9\times $}}\end{tabular} & \hspace*{1pt}\begin{tabular}[t]{@{}c@{}}\rule{0pt}{12pt}12m\,31.2s\\[-.2pt]{\fontsize{5.8}{6.4}\selectfont \textcolor{BrickRed}{$ \blacktriangledown\,1.7\times $}}\\[-.2pt]12m\,30.6s\\[-.2pt]{\fontsize{5.8}{6.4}\selectfont \textcolor{BrickRed}{$ \blacktriangledown\,1.8\times $}}\end{tabular} & \textcolor{gray}{\ensuremath{\cdot}} & \textcolor{gray}{\ensuremath{\cdot}} & \textcolor{gray}{\ensuremath{\cdot}} & \textcolor{gray}{\ensuremath{\cdot}} & \textcolor{gray}{\ensuremath{\cdot}} & \textcolor{gray}{\ensuremath{\cdot}} & \textcolor{gray}{\ensuremath{\cdot}} & \textcolor{gray}{\ensuremath{\cdot}} & \textcolor{gray}{\ensuremath{\cdot}} & \textcolor{gray}{\ensuremath{\cdot}} & \textcolor{gray}{\ensuremath{\cdot}} & \textcolor{gray}{\ensuremath{\cdot}}\\[3pt]\hline
 & \begin{tabular}[t]{@{}l@{}}\rule{0pt}{12pt}MPF / Scalar\\Scalar-DP / KSW2\end{tabular} & \hspace*{1pt}\begin{tabular}[t]{@{}c@{}}\rule{0pt}{12pt}5m\,39.7s\\[-.2pt]{\fontsize{5.8}{6.4}\selectfont \textcolor{BrickRed}{$ \blacktriangledown\,17.3\times $}}\\[-.2pt]5m\,37.2s\\[-.2pt]{\fontsize{5.8}{6.4}\selectfont \textcolor{BrickRed}{$ \blacktriangledown\,18.1\times $}}\end{tabular} & \hspace*{1pt}\begin{tabular}[t]{@{}c@{}}\rule{0pt}{12pt}12m\,8.9s\\[-.2pt]{\fontsize{5.8}{6.4}\selectfont \textcolor{BrickRed}{$ \blacktriangledown\,37.7\times $}}\\[-.2pt]12m\,7.2s\\[-.2pt]{\fontsize{5.8}{6.4}\selectfont \textcolor{BrickRed}{$ \blacktriangledown\,41.6\times $}}\end{tabular} & \textcolor{gray}{\ensuremath{\cdot}} & \hspace*{1pt}\begin{tabular}[t]{@{}c@{}}\rule{0pt}{12pt}10m\,39.3s\\[-.2pt]{\fontsize{5.8}{6.4}\selectfont \textcolor{BrickRed}{$ \blacktriangledown\,1.4\times $}}\\[-.2pt]10m\,36.8s\\[-.2pt]{\fontsize{5.8}{6.4}\selectfont \textcolor{BrickRed}{$ \blacktriangledown\,1.5\times $}}\end{tabular} & \textcolor{gray}{\ensuremath{\cdot}} & \textcolor{gray}{\ensuremath{\cdot}} & \textcolor{gray}{\ensuremath{\cdot}} & \textcolor{gray}{\ensuremath{\cdot}} & \textcolor{gray}{\ensuremath{\cdot}} & \textcolor{gray}{\ensuremath{\cdot}} & \textcolor{gray}{\ensuremath{\cdot}} & \textcolor{gray}{\ensuremath{\cdot}} & \textcolor{gray}{\ensuremath{\cdot}} & \textcolor{gray}{\ensuremath{\cdot}} & \textcolor{gray}{\ensuremath{\cdot}} & \textcolor{gray}{\ensuremath{\cdot}}\\[3pt]\hline
 & \begin{tabular}[t]{@{}l@{}}\rule{0pt}{12pt}MPF / Scalar\\SIMD-DP / Block\end{tabular} & \hspace*{1pt}\begin{tabular}[t]{@{}c@{}}\rule{0pt}{12pt}5m\,39.7s\\[-.2pt]{\fontsize{5.8}{6.4}\selectfont \textcolor{BrickRed}{$ \blacktriangledown\,17.3\times $}}\\[-.2pt]5m\,39.2s\\[-.2pt]{\fontsize{5.8}{6.4}\selectfont \textcolor{BrickRed}{$ \blacktriangledown\,18.2\times $}}\end{tabular} & \hspace*{1pt}\begin{tabular}[t]{@{}c@{}}\rule{0pt}{12pt}12m\,5.2s\\[-.2pt]{\fontsize{5.8}{6.4}\selectfont \textcolor{BrickRed}{$ \blacktriangledown\,37.5\times $}}\\[-.2pt]12m\,4.6s\\[-.2pt]{\fontsize{5.8}{6.4}\selectfont \textcolor{BrickRed}{$ \blacktriangledown\,41.4\times $}}\end{tabular} & \hspace*{1pt}\begin{tabular}[t]{@{}c@{}}\rule{0pt}{12pt}22m\,32.6s\\[-.2pt]{\fontsize{5.8}{6.4}\selectfont \textcolor{BrickRed}{$ \blacktriangledown\,18.6\times $}}\\[-.2pt]22m\,31.6s\\[-.2pt]{\fontsize{5.8}{6.4}\selectfont \textcolor{BrickRed}{$ \blacktriangledown\,19.8\times $}}\end{tabular} & \hspace*{1pt}\begin{tabular}[t]{@{}c@{}}\rule{0pt}{12pt}12m\,48.6s\\[-.2pt]{\fontsize{5.8}{6.4}\selectfont \textcolor{BrickRed}{$ \blacktriangledown\,1.7\times $}}\\[-.2pt]12m\,47.9s\\[-.2pt]{\fontsize{5.8}{6.4}\selectfont \textcolor{BrickRed}{$ \blacktriangledown\,1.9\times $}}\end{tabular} & \textcolor{gray}{\ensuremath{\cdot}} & \textcolor{gray}{\ensuremath{\cdot}} & \textcolor{gray}{\ensuremath{\cdot}} & \textcolor{gray}{\ensuremath{\cdot}} & \hspace*{1pt}\begin{tabular}[t]{@{}c@{}}\rule{0pt}{12pt}2m\,45.2s\\[-.2pt]{\fontsize{5.8}{6.4}\selectfont \textcolor{BrickRed}{$ \blacktriangledown\,3.9\times $}}\\[-.2pt]2m\,8.9s\\[-.2pt]{\fontsize{5.8}{6.4}\selectfont \textcolor{BrickRed}{$ \blacktriangledown\,3.1\times $}}\end{tabular} & \hspace*{1pt}\begin{tabular}[t]{@{}c@{}}\rule{0pt}{12pt}2m\,44.3s\\[-.2pt]{\fontsize{5.8}{6.4}\selectfont \textcolor{BrickRed}{$ \blacktriangledown\,3.8\times $}}\\[-.2pt]2m\,14.4s\\[-.2pt]{\fontsize{5.8}{6.4}\selectfont \textcolor{BrickRed}{$ \blacktriangledown\,3.3\times $}}\end{tabular} & \textcolor{gray}{\ensuremath{\cdot}} & \textcolor{gray}{\ensuremath{\cdot}} & \textcolor{gray}{\ensuremath{\cdot}} & \textcolor{gray}{\ensuremath{\cdot}} & \textcolor{gray}{\ensuremath{\cdot}} & \textcolor{gray}{\ensuremath{\cdot}}\\[3pt]\hline
 & \begin{tabular}[t]{@{}l@{}}\rule{0pt}{12pt}MPF / Scalar\\SIMD-DP / KSW2\end{tabular} & \hspace*{1pt}\begin{tabular}[t]{@{}c@{}}\rule{0pt}{12pt}5m\,42.8s\\[-.2pt]{\fontsize{5.8}{6.4}\selectfont \textcolor{BrickRed}{$ \blacktriangledown\,17.5\times $}}\\[-.2pt]5m\,40.6s\\[-.2pt]{\fontsize{5.8}{6.4}\selectfont \textcolor{BrickRed}{$ \blacktriangledown\,18.3\times $}}\end{tabular} & \hspace*{1pt}\begin{tabular}[t]{@{}c@{}}\rule{0pt}{12pt}12m\,19.8s\\[-.2pt]{\fontsize{5.8}{6.4}\selectfont \textcolor{BrickRed}{$ \blacktriangledown\,38.2\times $}}\\[-.2pt]12m\,18s\\[-.2pt]{\fontsize{5.8}{6.4}\selectfont \textcolor{BrickRed}{$ \blacktriangledown\,42.2\times $}}\end{tabular} & \hspace*{1pt}\begin{tabular}[t]{@{}c@{}}\rule{0pt}{12pt}22m\,43.4s\\[-.2pt]{\fontsize{5.8}{6.4}\selectfont \textcolor{BrickRed}{$ \blacktriangledown\,18.8\times $}}\\[-.2pt]22m\,39.5s\\[-.2pt]{\fontsize{5.8}{6.4}\selectfont \textcolor{BrickRed}{$ \blacktriangledown\,20.0\times $}}\end{tabular} & \hspace*{1pt}\begin{tabular}[t]{@{}c@{}}\rule{0pt}{12pt}11m\,32.1s\\[-.2pt]{\fontsize{5.8}{6.4}\selectfont \textcolor{BrickRed}{$ \blacktriangledown\,1.6\times $}}\\[-.2pt]11m\,29.4s\\[-.2pt]{\fontsize{5.8}{6.4}\selectfont \textcolor{BrickRed}{$ \blacktriangledown\,1.7\times $}}\end{tabular} & \textcolor{gray}{\ensuremath{\cdot}} & \textcolor{gray}{\ensuremath{\cdot}} & \textcolor{gray}{\ensuremath{\cdot}} & \textcolor{gray}{\ensuremath{\cdot}} & \hspace*{1pt}\begin{tabular}[t]{@{}c@{}}\rule{0pt}{12pt}2m\,52.8s\\[-.2pt]{\fontsize{5.8}{6.4}\selectfont \textcolor{BrickRed}{$ \blacktriangledown\,4.1\times $}}\\[-.2pt]2m\,15.7s\\[-.2pt]{\fontsize{5.8}{6.4}\selectfont \textcolor{BrickRed}{$ \blacktriangledown\,3.3\times $}}\end{tabular} & \hspace*{1pt}\begin{tabular}[t]{@{}c@{}}\rule{0pt}{12pt}2m\,32.2s\\[-.2pt]{\fontsize{5.8}{6.4}\selectfont \textcolor{BrickRed}{$ \blacktriangledown\,3.6\times $}}\\[-.2pt]2m\,0.8s\\[-.2pt]{\fontsize{5.8}{6.4}\selectfont \textcolor{BrickRed}{$ \blacktriangledown\,2.9\times $}}\end{tabular} & \textcolor{gray}{\ensuremath{\cdot}} & \textcolor{gray}{\ensuremath{\cdot}} & \textcolor{gray}{\ensuremath{\cdot}} & \textcolor{gray}{\ensuremath{\cdot}} & \textcolor{gray}{\ensuremath{\cdot}} & \textcolor{gray}{\ensuremath{\cdot}}\\[3pt]\hline
 & \begin{tabular}[t]{@{}l@{}}\rule{0pt}{12pt}MPF / SIMD-ref\\RMQ / KSW2\end{tabular} & \textcolor{gray}{\ensuremath{\cdot}} & \textcolor{gray}{\ensuremath{\cdot}} & \textcolor{gray}{\ensuremath{\cdot}} & \hspace*{1pt}\begin{tabular}[t]{@{}c@{}}\rule{0pt}{12pt}10m\,52.1s\\[-.2pt]{\fontsize{5.8}{6.4}\selectfont \textcolor{BrickRed}{$ \blacktriangledown\,1.5\times $}}\\[-.2pt]10m\,49.7s\\[-.2pt]{\fontsize{5.8}{6.4}\selectfont \textcolor{BrickRed}{$ \blacktriangledown\,1.6\times $}}\end{tabular} & \textcolor{gray}{\ensuremath{\cdot}} & \textcolor{gray}{\ensuremath{\cdot}} & \textcolor{gray}{\ensuremath{\cdot}} & \textcolor{gray}{\ensuremath{\cdot}} & \textcolor{gray}{\ensuremath{\cdot}} & \textcolor{gray}{\ensuremath{\cdot}} & \textcolor{gray}{\ensuremath{\cdot}} & \textcolor{gray}{\ensuremath{\cdot}} & \textcolor{gray}{\ensuremath{\cdot}} & \textcolor{gray}{\ensuremath{\cdot}} & \textcolor{gray}{\ensuremath{\cdot}} & \textcolor{gray}{\ensuremath{\cdot}}\\[3pt]\hline
 & \begin{tabular}[t]{@{}l@{}}\rule{0pt}{12pt}MPF / SIMD-ref\\SIMD-DP / Block\end{tabular} & \textcolor{gray}{\ensuremath{\cdot}} & \textcolor{gray}{\ensuremath{\cdot}} & \textcolor{gray}{\ensuremath{\cdot}} & \hspace*{1pt}\begin{tabular}[t]{@{}c@{}}\rule{0pt}{12pt}13m\,2.2s\\[-.2pt]{\fontsize{5.8}{6.4}\selectfont \textcolor{BrickRed}{$ \blacktriangledown\,1.8\times $}}\\[-.2pt]13m\,1.5s\\[-.2pt]{\fontsize{5.8}{6.4}\selectfont \textcolor{BrickRed}{$ \blacktriangledown\,1.9\times $}}\end{tabular} & \textcolor{gray}{\ensuremath{\cdot}} & \textcolor{gray}{\ensuremath{\cdot}} & \textcolor{gray}{\ensuremath{\cdot}} & \textcolor{gray}{\ensuremath{\cdot}} & \textcolor{gray}{\ensuremath{\cdot}} & \textcolor{gray}{\ensuremath{\cdot}} & \textcolor{gray}{\ensuremath{\cdot}} & \textcolor{gray}{\ensuremath{\cdot}} & \textcolor{gray}{\ensuremath{\cdot}} & \textcolor{gray}{\ensuremath{\cdot}} & \textcolor{gray}{\ensuremath{\cdot}} & \textcolor{gray}{\ensuremath{\cdot}}\\[3pt]\hline
\end{tabular}\par}
\smallskip{\footnotesize\noindent Each cell: A, comparison, B, comparison. Times cover 1,024 queries. $\cdot$ = no approved observation. ---\textsuperscript{g} = GPU unavailable, not executed. n/a = no measured comparison.\par
\noindent \textsuperscript{f}SIMD-ref request observed as scalar, compared only as that effective path. \textsuperscript{p}Accepted by explicit preservation after wrapper postflight failure, no run receipt claimed.\par
\noindent MPF/MPS = minimizers per fragment/sequence. Row order: index representation / minimizer, then chaining / base alignment. SIMD-ref is the reference minimizer, Bounded is bounded ntHash.\par
\noindent Blue = one lowest A and one lowest B down each machine/workload column, across continuation pages. Exact ties select the first row. Repeated baselines are marked only on page 1.\par
\noindent Light gray = minimap2 baseline. Pale yellow = preferred MPF / Bounded / SIMD-DP / Block configuration, not a claim of winning every column.\par
\noindent When all Vulkan modes are unmeasured, they share one row. Baseline comparisons use the fastest measured BioseqDB effective path separately for A and B.\par}
\clearpage
\noindent{\large\bfseries AWS: 2 Gbp, k=20, w=15 (mapping interval)\hfill 3/6}\par\smallskip
{\fontsize{7.1}{8.1}\selectfont\setlength{\tabcolsep}{1pt}\setlength{\fboxsep}{2pt}\renewcommand{\arraystretch}{1.12}
\noindent\begin{tabular}{@{}p{14mm}p{27mm}|p{13.55mm}|p{13.55mm}|p{13.55mm}|p{13.55mm}|p{13.55mm}|p{13.55mm}|p{13.55mm}|p{13.55mm}|p{13.55mm}|p{13.55mm}|p{13.55mm}|p{13.55mm}|p{13.55mm}|p{13.55mm}|p{13.55mm}|p{13.55mm}|@{}}\hline
Engine & Execution & \multicolumn{8}{c|}{\textbf{Index}} & \multicolumn{8}{c|}{\textbf{No index}}\\
 &  & \multicolumn{8}{c|}{\textbf{2 Gbp}} & \multicolumn{8}{c|}{\textbf{2 Gbp}}\\
 &  & \multicolumn{8}{c|}{\textbf{k=20, w=15}} & \multicolumn{8}{c|}{\textbf{k=20, w=15}}\\
 &  & \multicolumn{4}{c|}{\textbf{DNA4}} & \multicolumn{4}{c|}{\textbf{DNA5}} & \multicolumn{4}{c|}{\textbf{DNA4}} & \multicolumn{4}{c|}{\textbf{DNA5}}\\
 &  & \multicolumn{1}{c|}{\textbf{M7A}} & \multicolumn{1}{c|}{\textbf{M7I}} & \multicolumn{1}{c|}{\textbf{M4}} & \multicolumn{1}{c|}{\textbf{C7G}} & \multicolumn{1}{c|}{\textbf{M7A}} & \multicolumn{1}{c|}{\textbf{M7I}} & \multicolumn{1}{c|}{\textbf{M4}} & \multicolumn{1}{c|}{\textbf{C7G}} & \multicolumn{1}{c|}{\textbf{M7A}} & \multicolumn{1}{c|}{\textbf{M7I}} & \multicolumn{1}{c|}{\textbf{M4}} & \multicolumn{1}{c|}{\textbf{C7G}} & \multicolumn{1}{c|}{\textbf{M7A}} & \multicolumn{1}{c|}{\textbf{M7I}} & \multicolumn{1}{c|}{\textbf{M4}} & \multicolumn{1}{c|}{\textbf{C7G}}\\
\hline
\rowcolor{BenchBaseline}[0.5pt][0.5pt]\textbf{minimap2} & \textbf{Baseline} & \hspace*{1pt}\begin{tabular}[t]{@{}c@{}}\rule{0pt}{12pt}19.6s\\[-.2pt]{\fontsize{5.8}{6.4}\selectfont \textcolor{BrickRed}{$ \blacktriangledown\,6.7\times $}}\\[-.2pt]18.6s\\[-.2pt]{\fontsize{5.8}{6.4}\selectfont \textcolor{BrickRed}{$ \blacktriangledown\,8.4\times $}}\end{tabular} & \hspace*{1pt}\begin{tabular}[t]{@{}c@{}}\rule{0pt}{12pt}19.4s\\[-.2pt]{\fontsize{5.8}{6.4}\selectfont \textcolor{BrickRed}{$ \blacktriangledown\,7.2\times $}}\\[-.2pt]17.5s\\[-.2pt]{\fontsize{5.8}{6.4}\selectfont \textcolor{BrickRed}{$ \blacktriangledown\,9.1\times $}}\end{tabular} & \hspace*{1pt}\begin{tabular}[t]{@{}c@{}}\rule{0pt}{12pt}1m\,12.7s\\[-.2pt]{\fontsize{5.8}{6.4}\selectfont \textcolor{BrickRed}{$ \blacktriangledown\,13.6\times $}}\\[-.2pt]1m\,8.1s\\[-.2pt]{\fontsize{5.8}{6.4}\selectfont \textcolor{BrickRed}{$ \blacktriangledown\,20.2\times $}}\end{tabular} & \hspace*{1pt}\begin{tabular}[t]{@{}c@{}}\rule{0pt}{12pt}7m\,25.4s\\[-.2pt]{\fontsize{5.8}{6.4}\selectfont \textcolor{BrickRed}{$ \blacktriangledown\,106.2\times $}}\\[-.2pt]6m\,50.8s\\[-.2pt]{\fontsize{5.8}{6.4}\selectfont \textcolor{BrickRed}{$ \blacktriangledown\,212.9\times $}}\end{tabular} & \hspace*{1pt}\begin{tabular}[t]{@{}c@{}}\rule{0pt}{12pt}9.3s\\[-.2pt]{\fontsize{5.8}{6.4}\selectfont \textcolor{BrickRed}{$ \blacktriangledown\,24\% $}}\\[-.2pt]8.4s\\[-.2pt]{\fontsize{5.8}{6.4}\selectfont \textcolor{BrickRed}{$ \blacktriangledown\,1.6\times $}}\end{tabular} & \hspace*{1pt}\begin{tabular}[t]{@{}c@{}}\rule{0pt}{12pt}8.7s\\[-.2pt]{\fontsize{5.8}{6.4}\selectfont \textcolor{BrickRed}{$ \blacktriangledown\,1.5\times $}}\\[-.2pt]7s\\[-.2pt]{\fontsize{5.8}{6.4}\selectfont \textcolor{BrickRed}{$ \blacktriangledown\,1.7\times $}}\end{tabular} & \hspace*{1pt}\begin{tabular}[t]{@{}c@{}}\rule{0pt}{12pt}15.8s\\[-.2pt]{\fontsize{5.8}{6.4}\selectfont \textcolor{BrickRed}{$ \blacktriangledown\,26\% $}}\\[-.2pt]12.7s\\[-.2pt]{\fontsize{5.8}{6.4}\selectfont \textcolor{BrickRed}{$ \blacktriangledown\,1.6\times $}}\end{tabular} & \hspace*{1pt}\begin{tabular}[t]{@{}c@{}}\rule{0pt}{12pt}18.5s\\[-.2pt]{\fontsize{5.8}{6.4}\selectfont \textcolor{BrickRed}{$ \blacktriangledown\,1.6\times $}}\\[-.2pt]17s\\[-.2pt]{\fontsize{5.8}{6.4}\selectfont \textcolor{BrickRed}{$ \blacktriangledown\,1.8\times $}}\end{tabular} & \hspace*{1pt}\begin{tabular}[t]{@{}c@{}}\rule{0pt}{12pt}42.4s\\[-.2pt]{\fontsize{5.8}{6.4}\selectfont \textcolor{ForestGreen}{$ \blacktriangle\,2.0\times $}}\\[-.2pt]41.4s\\[-.2pt]{\fontsize{5.8}{6.4}\selectfont \textcolor{ForestGreen}{$ \blacktriangle\,1.9\times $}}\end{tabular} & \hspace*{1pt}\begin{tabular}[t]{@{}c@{}}\rule{0pt}{12pt}42.9s\\[-.2pt]{\fontsize{5.8}{6.4}\selectfont \textcolor{ForestGreen}{$ \blacktriangle\,2.0\times $}}\\[-.2pt]41.1s\\[-.2pt]{\fontsize{5.8}{6.4}\selectfont \textcolor{ForestGreen}{$ \blacktriangle\,1.8\times $}}\end{tabular} & \hspace*{1pt}\begin{tabular}[t]{@{}c@{}}\rule{0pt}{12pt}1m\,26.9s\\[-.2pt]{\fontsize{5.8}{6.4}\selectfont \textcolor{ForestGreen}{$ \blacktriangle\,1.9\times $}}\\[-.2pt]1m\,23.7s\\[-.2pt]{\fontsize{5.8}{6.4}\selectfont \textcolor{ForestGreen}{$ \blacktriangle\,1.7\times $}}\end{tabular} & \hspace*{1pt}\begin{tabular}[t]{@{}c@{}}\rule{0pt}{12pt}4m\,51.6s\\[-.2pt]{\fontsize{5.8}{6.4}\selectfont \textcolor{gray}{n/a}}\\[-.2pt]4m\,46.6s\\[-.2pt]{\fontsize{5.8}{6.4}\selectfont \textcolor{gray}{n/a}}\end{tabular} & \hspace*{1pt}\begin{tabular}[t]{@{}c@{}}\rule{0pt}{12pt}35s\\[-.2pt]{\fontsize{5.8}{6.4}\selectfont \textcolor{ForestGreen}{$ \blacktriangle\,8.6\times $}}\\[-.2pt]34s\\[-.2pt]{\fontsize{5.8}{6.4}\selectfont \textcolor{ForestGreen}{$ \blacktriangle\,2.5\times $}}\end{tabular} & \hspace*{1pt}\begin{tabular}[t]{@{}c@{}}\rule{0pt}{12pt}35.9s\\[-.2pt]{\fontsize{5.8}{6.4}\selectfont \textcolor{ForestGreen}{$ \blacktriangle\,6.7\times $}}\\[-.2pt]34.2s\\[-.2pt]{\fontsize{5.8}{6.4}\selectfont \textcolor{ForestGreen}{$ \blacktriangle\,2.5\times $}}\end{tabular} & \hspace*{1pt}\begin{tabular}[t]{@{}c@{}}\rule{0pt}{12pt}1m\,14.4s\\[-.2pt]{\fontsize{5.8}{6.4}\selectfont \textcolor{ForestGreen}{$ \blacktriangle\,6.9\times $}}\\[-.2pt]1m\,11.3s\\[-.2pt]{\fontsize{5.8}{6.4}\selectfont \textcolor{ForestGreen}{$ \blacktriangle\,2.2\times $}}\end{tabular} & \hspace*{1pt}\begin{tabular}[t]{@{}c@{}}\rule{0pt}{12pt}1m\,0.1s\\[-.2pt]{\fontsize{5.8}{6.4}\selectfont \textcolor{gray}{n/a}}\\[-.2pt]58.6s\\[-.2pt]{\fontsize{5.8}{6.4}\selectfont \textcolor{gray}{n/a}}\end{tabular}\\[3pt]\hline
\textbf{BioseqDB} & \begin{tabular}[t]{@{}l@{}}\rule{0pt}{12pt}MPF / Two-stack\\RMQ / Block\end{tabular} & \hspace*{1pt}\begin{tabular}[t]{@{}c@{}}\rule{0pt}{12pt}4m\,37.6s\\[-.2pt]{\fontsize{5.8}{6.4}\selectfont \textcolor{BrickRed}{$ \blacktriangledown\,14.2\times $}}\\[-.2pt]4m\,37.1s\\[-.2pt]{\fontsize{5.8}{6.4}\selectfont \textcolor{BrickRed}{$ \blacktriangledown\,14.9\times $}}\end{tabular} & \hspace*{1pt}\begin{tabular}[t]{@{}c@{}}\rule{0pt}{12pt}7m\,30s\\[-.2pt]{\fontsize{5.8}{6.4}\selectfont \textcolor{BrickRed}{$ \blacktriangledown\,23.2\times $}}\\[-.2pt]7m\,29.2s\\[-.2pt]{\fontsize{5.8}{6.4}\selectfont \textcolor{BrickRed}{$ \blacktriangledown\,25.7\times $}}\end{tabular} & \textcolor{gray}{\ensuremath{\cdot}} & \textcolor{gray}{\ensuremath{\cdot}} & \textcolor{gray}{\ensuremath{\cdot}} & \textcolor{gray}{\ensuremath{\cdot}} & \textcolor{gray}{\ensuremath{\cdot}} & \textcolor{gray}{\ensuremath{\cdot}} & \textcolor{gray}{\ensuremath{\cdot}} & \textcolor{gray}{\ensuremath{\cdot}} & \textcolor{gray}{\ensuremath{\cdot}} & \textcolor{gray}{\ensuremath{\cdot}} & \textcolor{gray}{\ensuremath{\cdot}} & \textcolor{gray}{\ensuremath{\cdot}} & \textcolor{gray}{\ensuremath{\cdot}} & \textcolor{gray}{\ensuremath{\cdot}}\\[3pt]\hline
 & \begin{tabular}[t]{@{}l@{}}\rule{0pt}{12pt}MPF / Two-stack\\RMQ / KSW2\end{tabular} & \hspace*{1pt}\begin{tabular}[t]{@{}c@{}}\rule{0pt}{12pt}4m\,34.2s\\[-.2pt]{\fontsize{5.8}{6.4}\selectfont \textcolor{BrickRed}{$ \blacktriangledown\,14.0\times $}}\\[-.2pt]4m\,32.2s\\[-.2pt]{\fontsize{5.8}{6.4}\selectfont \textcolor{BrickRed}{$ \blacktriangledown\,14.6\times $}}\end{tabular} & \hspace*{1pt}\begin{tabular}[t]{@{}c@{}}\rule{0pt}{12pt}6m\,52.3s\\[-.2pt]{\fontsize{5.8}{6.4}\selectfont \textcolor{BrickRed}{$ \blacktriangledown\,21.3\times $}}\\[-.2pt]6m\,49.6s\\[-.2pt]{\fontsize{5.8}{6.4}\selectfont \textcolor{BrickRed}{$ \blacktriangledown\,23.4\times $}}\end{tabular} & \textcolor{gray}{\ensuremath{\cdot}} & \textcolor{gray}{\ensuremath{\cdot}} & \textcolor{gray}{\ensuremath{\cdot}} & \textcolor{gray}{\ensuremath{\cdot}} & \textcolor{gray}{\ensuremath{\cdot}} & \textcolor{gray}{\ensuremath{\cdot}} & \textcolor{gray}{\ensuremath{\cdot}} & \textcolor{gray}{\ensuremath{\cdot}} & \textcolor{gray}{\ensuremath{\cdot}} & \textcolor{gray}{\ensuremath{\cdot}} & \textcolor{gray}{\ensuremath{\cdot}} & \textcolor{gray}{\ensuremath{\cdot}} & \textcolor{gray}{\ensuremath{\cdot}} & \textcolor{gray}{\ensuremath{\cdot}}\\[3pt]\hline
 & \begin{tabular}[t]{@{}l@{}}\rule{0pt}{12pt}MPF / Two-stack\\Scalar-DP / Block\end{tabular} & \hspace*{1pt}\begin{tabular}[t]{@{}c@{}}\rule{0pt}{12pt}4m\,36s\\[-.2pt]{\fontsize{5.8}{6.4}\selectfont \textcolor{BrickRed}{$ \blacktriangledown\,14.1\times $}}\\[-.2pt]4m\,35.5s\\[-.2pt]{\fontsize{5.8}{6.4}\selectfont \textcolor{BrickRed}{$ \blacktriangledown\,14.8\times $}}\end{tabular} & \hspace*{1pt}\begin{tabular}[t]{@{}c@{}}\rule{0pt}{12pt}6m\,53.8s\\[-.2pt]{\fontsize{5.8}{6.4}\selectfont \textcolor{BrickRed}{$ \blacktriangledown\,21.4\times $}}\\[-.2pt]6m\,52.9s\\[-.2pt]{\fontsize{5.8}{6.4}\selectfont \textcolor{BrickRed}{$ \blacktriangledown\,23.6\times $}}\end{tabular} & \textcolor{gray}{\ensuremath{\cdot}} & \textcolor{gray}{\ensuremath{\cdot}} & \textcolor{gray}{\ensuremath{\cdot}} & \textcolor{gray}{\ensuremath{\cdot}} & \textcolor{gray}{\ensuremath{\cdot}} & \textcolor{gray}{\ensuremath{\cdot}} & \textcolor{gray}{\ensuremath{\cdot}} & \textcolor{gray}{\ensuremath{\cdot}} & \textcolor{gray}{\ensuremath{\cdot}} & \textcolor{gray}{\ensuremath{\cdot}} & \textcolor{gray}{\ensuremath{\cdot}} & \textcolor{gray}{\ensuremath{\cdot}} & \textcolor{gray}{\ensuremath{\cdot}} & \textcolor{gray}{\ensuremath{\cdot}}\\[3pt]\hline
 & \begin{tabular}[t]{@{}l@{}}\rule{0pt}{12pt}MPF / Two-stack\\Scalar-DP / KSW2\end{tabular} & \hspace*{1pt}\begin{tabular}[t]{@{}c@{}}\rule{0pt}{12pt}4m\,38.1s\\[-.2pt]{\fontsize{5.8}{6.4}\selectfont \textcolor{BrickRed}{$ \blacktriangledown\,14.2\times $}}\\[-.2pt]4m\,36.2s\\[-.2pt]{\fontsize{5.8}{6.4}\selectfont \textcolor{BrickRed}{$ \blacktriangledown\,14.8\times $}}\end{tabular} & \hspace*{1pt}\begin{tabular}[t]{@{}c@{}}\rule{0pt}{12pt}7m\,34.7s\\[-.2pt]{\fontsize{5.8}{6.4}\selectfont \textcolor{BrickRed}{$ \blacktriangledown\,23.5\times $}}\\[-.2pt]7m\,31.6s\\[-.2pt]{\fontsize{5.8}{6.4}\selectfont \textcolor{BrickRed}{$ \blacktriangledown\,25.8\times $}}\end{tabular} & \textcolor{gray}{\ensuremath{\cdot}} & \textcolor{gray}{\ensuremath{\cdot}} & \textcolor{gray}{\ensuremath{\cdot}} & \textcolor{gray}{\ensuremath{\cdot}} & \textcolor{gray}{\ensuremath{\cdot}} & \textcolor{gray}{\ensuremath{\cdot}} & \textcolor{gray}{\ensuremath{\cdot}} & \textcolor{gray}{\ensuremath{\cdot}} & \textcolor{gray}{\ensuremath{\cdot}} & \textcolor{gray}{\ensuremath{\cdot}} & \textcolor{gray}{\ensuremath{\cdot}} & \textcolor{gray}{\ensuremath{\cdot}} & \textcolor{gray}{\ensuremath{\cdot}} & \textcolor{gray}{\ensuremath{\cdot}}\\[3pt]\hline
 & \begin{tabular}[t]{@{}l@{}}\rule{0pt}{12pt}MPF / Two-stack\\SIMD-DP / Block\end{tabular} & \hspace*{1pt}\begin{tabular}[t]{@{}c@{}}\rule{0pt}{12pt}4m\,36.1s\\[-.2pt]{\fontsize{5.8}{6.4}\selectfont \textcolor{BrickRed}{$ \blacktriangledown\,14.1\times $}}\\[-.2pt]4m\,35.5s\\[-.2pt]{\fontsize{5.8}{6.4}\selectfont \textcolor{BrickRed}{$ \blacktriangledown\,14.8\times $}}\end{tabular} & \hspace*{1pt}\begin{tabular}[t]{@{}c@{}}\rule{0pt}{12pt}7m\,3.9s\\[-.2pt]{\fontsize{5.8}{6.4}\selectfont \textcolor{BrickRed}{$ \blacktriangledown\,21.9\times $}}\\[-.2pt]7m\,3s\\[-.2pt]{\fontsize{5.8}{6.4}\selectfont \textcolor{BrickRed}{$ \blacktriangledown\,24.2\times $}}\end{tabular} & \textcolor{gray}{\ensuremath{\cdot}} & \hspace*{1pt}\begin{tabular}[t]{@{}c@{}}\rule{0pt}{12pt}5m\,27s\\[-.2pt]{\fontsize{5.8}{6.4}\selectfont \textcolor{ForestGreen}{$ \blacktriangle\,27\% $}}\\[-.2pt]5m\,26s\\[-.2pt]{\fontsize{5.8}{6.4}\selectfont \textcolor{ForestGreen}{$ \blacktriangle\,21\% $}}\end{tabular} & \textcolor{gray}{\ensuremath{\cdot}} & \textcolor{gray}{\ensuremath{\cdot}} & \textcolor{gray}{\ensuremath{\cdot}} & \hspace*{1pt}\begin{tabular}[t]{@{}c@{}}\rule{0pt}{12pt}4m\,0.4s\\[-.2pt]{\fontsize{5.8}{6.4}\selectfont \textcolor{BrickRed}{$ \blacktriangledown\,13.0\times $}}\\[-.2pt]3m\,59s\\[-.2pt]{\fontsize{5.8}{6.4}\selectfont \textcolor{BrickRed}{$ \blacktriangledown\,14.0\times $}}\end{tabular} & \textcolor{gray}{\ensuremath{\cdot}} & \textcolor{gray}{\ensuremath{\cdot}} & \textcolor{gray}{\ensuremath{\cdot}} & \textcolor{gray}{\ensuremath{\cdot}} & \textcolor{gray}{\ensuremath{\cdot}} & \textcolor{gray}{\ensuremath{\cdot}} & \textcolor{gray}{\ensuremath{\cdot}} & \textcolor{gray}{\ensuremath{\cdot}}\\[3pt]\hline
 & \begin{tabular}[t]{@{}l@{}}\rule{0pt}{12pt}MPF / Two-stack\\SIMD-DP / KSW2\end{tabular} & \hspace*{1pt}\begin{tabular}[t]{@{}c@{}}\rule{0pt}{12pt}4m\,37.8s\\[-.2pt]{\fontsize{5.8}{6.4}\selectfont \textcolor{BrickRed}{$ \blacktriangledown\,14.2\times $}}\\[-.2pt]4m\,35.9s\\[-.2pt]{\fontsize{5.8}{6.4}\selectfont \textcolor{BrickRed}{$ \blacktriangledown\,14.8\times $}}\end{tabular} & \hspace*{1pt}\begin{tabular}[t]{@{}c@{}}\rule{0pt}{12pt}7m\,37.6s\\[-.2pt]{\fontsize{5.8}{6.4}\selectfont \textcolor{BrickRed}{$ \blacktriangledown\,23.6\times $}}\\[-.2pt]7m\,34.4s\\[-.2pt]{\fontsize{5.8}{6.4}\selectfont \textcolor{BrickRed}{$ \blacktriangledown\,26.0\times $}}\end{tabular} & \textcolor{gray}{\ensuremath{\cdot}} & \textcolor{gray}{\ensuremath{\cdot}} & \textcolor{gray}{\ensuremath{\cdot}} & \textcolor{gray}{\ensuremath{\cdot}} & \textcolor{gray}{\ensuremath{\cdot}} & \textcolor{gray}{\ensuremath{\cdot}} & \textcolor{gray}{\ensuremath{\cdot}} & \textcolor{gray}{\ensuremath{\cdot}} & \textcolor{gray}{\ensuremath{\cdot}} & \textcolor{gray}{\ensuremath{\cdot}} & \textcolor{gray}{\ensuremath{\cdot}} & \textcolor{gray}{\ensuremath{\cdot}} & \textcolor{gray}{\ensuremath{\cdot}} & \textcolor{gray}{\ensuremath{\cdot}}\\[3pt]\hline
 & \begin{tabular}[t]{@{}l@{}}\rule{0pt}{12pt}MPS / Bounded\\RMQ / Block\end{tabular} & \hspace*{1pt}\begin{tabular}[t]{@{}c@{}}\rule{0pt}{12pt}6.6s\\[-.2pt]{\fontsize{5.8}{6.4}\selectfont \textcolor{ForestGreen}{$ \blacktriangle\,3.0\times $}}\\[-.2pt]5.9s\\[-.2pt]{\fontsize{5.8}{6.4}\selectfont \textcolor{ForestGreen}{$ \blacktriangle\,3.2\times $}}\end{tabular} & \hspace*{1pt}\begin{tabular}[t]{@{}c@{}}\rule{0pt}{12pt}6.5s\\[-.2pt]{\fontsize{5.8}{6.4}\selectfont \textcolor{ForestGreen}{$ \blacktriangle\,3.0\times $}}\\[-.2pt]5.8s\\[-.2pt]{\fontsize{5.8}{6.4}\selectfont \textcolor{ForestGreen}{$ \blacktriangle\,3.0\times $}}\end{tabular} & \textcolor{gray}{\ensuremath{\cdot}} & \hspace*{1pt}\begin{tabular}[t]{@{}c@{}}\rule{0pt}{12pt}6.4s\\[-.2pt]{\fontsize{5.8}{6.4}\selectfont \textcolor{ForestGreen}{$ \blacktriangle\,69.8\times $}}\\[-.2pt]5.5s\\[-.2pt]{\fontsize{5.8}{6.4}\selectfont \textcolor{ForestGreen}{$ \blacktriangle\,75.0\times $}}\end{tabular} & \textcolor{gray}{\ensuremath{\cdot}} & \textcolor{gray}{\ensuremath{\cdot}} & \textcolor{gray}{\ensuremath{\cdot}} & \textcolor{gray}{\ensuremath{\cdot}} & \hspace*{1pt}\begin{tabular}[t]{@{}c@{}}\rule{0pt}{12pt}1m\,43.2s\\[-.2pt]{\fontsize{5.8}{6.4}\selectfont \textcolor{BrickRed}{$ \blacktriangledown\,2.4\times $}}\\[-.2pt]1m\,21s\\[-.2pt]{\fontsize{5.8}{6.4}\selectfont \textcolor{BrickRed}{$ \blacktriangledown\,2.0\times $}}\end{tabular} & \hspace*{1pt}\begin{tabular}[t]{@{}c@{}}\rule{0pt}{12pt}1m\,41.8s\\[-.2pt]{\fontsize{5.8}{6.4}\selectfont \textcolor{BrickRed}{$ \blacktriangledown\,2.4\times $}}\\[-.2pt]1m\,21.9s\\[-.2pt]{\fontsize{5.8}{6.4}\selectfont \textcolor{BrickRed}{$ \blacktriangledown\,2.0\times $}}\end{tabular} & \hspace*{1pt}\begin{tabular}[t]{@{}c@{}}\rule{0pt}{12pt}2m\,55s\\[-.2pt]{\fontsize{5.8}{6.4}\selectfont \textcolor{BrickRed}{$ \blacktriangledown\,2.0\times $}}\\[-.2pt]2m\,21s\\[-.2pt]{\fontsize{5.8}{6.4}\selectfont \textcolor{BrickRed}{$ \blacktriangledown\,1.7\times $}}\end{tabular} & \textcolor{gray}{\ensuremath{\cdot}} & \textcolor{gray}{\ensuremath{\cdot}} & \textcolor{gray}{\ensuremath{\cdot}} & \textcolor{gray}{\ensuremath{\cdot}} & \textcolor{gray}{\ensuremath{\cdot}}\\[3pt]\hline
\end{tabular}\par}
\smallskip{\footnotesize\noindent Each cell: A, comparison, B, comparison. Times cover 1,024 queries. $\cdot$ = no approved observation. ---\textsuperscript{g} = GPU unavailable, not executed. n/a = no measured comparison.\par
\noindent \textsuperscript{f}SIMD-ref request observed as scalar, compared only as that effective path. \textsuperscript{p}Accepted by explicit preservation after wrapper postflight failure, no run receipt claimed.\par
\noindent MPF/MPS = minimizers per fragment/sequence. Row order: index representation / minimizer, then chaining / base alignment. SIMD-ref is the reference minimizer, Bounded is bounded ntHash.\par
\noindent Blue = one lowest A and one lowest B down each machine/workload column, across continuation pages. Exact ties select the first row. Repeated baselines are marked only on page 1.\par
\noindent Light gray = minimap2 baseline. Pale yellow = preferred MPF / Bounded / SIMD-DP / Block configuration, not a claim of winning every column.\par
\noindent When all Vulkan modes are unmeasured, they share one row. Baseline comparisons use the fastest measured BioseqDB effective path separately for A and B.\par}
\clearpage
\noindent{\large\bfseries AWS: 2 Gbp, k=20, w=15 (mapping interval)\hfill 4/6}\par\smallskip
{\fontsize{7.1}{8.1}\selectfont\setlength{\tabcolsep}{1pt}\setlength{\fboxsep}{2pt}\renewcommand{\arraystretch}{1.12}
\noindent\begin{tabular}{@{}p{14mm}p{27mm}|p{13.55mm}|p{13.55mm}|p{13.55mm}|p{13.55mm}|p{13.55mm}|p{13.55mm}|p{13.55mm}|p{13.55mm}|p{13.55mm}|p{13.55mm}|p{13.55mm}|p{13.55mm}|p{13.55mm}|p{13.55mm}|p{13.55mm}|p{13.55mm}|@{}}\hline
Engine & Execution & \multicolumn{8}{c|}{\textbf{Index}} & \multicolumn{8}{c|}{\textbf{No index}}\\
 &  & \multicolumn{8}{c|}{\textbf{2 Gbp}} & \multicolumn{8}{c|}{\textbf{2 Gbp}}\\
 &  & \multicolumn{8}{c|}{\textbf{k=20, w=15}} & \multicolumn{8}{c|}{\textbf{k=20, w=15}}\\
 &  & \multicolumn{4}{c|}{\textbf{DNA4}} & \multicolumn{4}{c|}{\textbf{DNA5}} & \multicolumn{4}{c|}{\textbf{DNA4}} & \multicolumn{4}{c|}{\textbf{DNA5}}\\
 &  & \multicolumn{1}{c|}{\textbf{M7A}} & \multicolumn{1}{c|}{\textbf{M7I}} & \multicolumn{1}{c|}{\textbf{M4}} & \multicolumn{1}{c|}{\textbf{C7G}} & \multicolumn{1}{c|}{\textbf{M7A}} & \multicolumn{1}{c|}{\textbf{M7I}} & \multicolumn{1}{c|}{\textbf{M4}} & \multicolumn{1}{c|}{\textbf{C7G}} & \multicolumn{1}{c|}{\textbf{M7A}} & \multicolumn{1}{c|}{\textbf{M7I}} & \multicolumn{1}{c|}{\textbf{M4}} & \multicolumn{1}{c|}{\textbf{C7G}} & \multicolumn{1}{c|}{\textbf{M7A}} & \multicolumn{1}{c|}{\textbf{M7I}} & \multicolumn{1}{c|}{\textbf{M4}} & \multicolumn{1}{c|}{\textbf{C7G}}\\
\hline
\rowcolor{BenchBaseline}[0.5pt][0.5pt]\textbf{minimap2} & \textbf{Baseline} & \hspace*{1pt}\begin{tabular}[t]{@{}c@{}}\rule{0pt}{12pt}19.6s\\[-.2pt]{\fontsize{5.8}{6.4}\selectfont \textcolor{BrickRed}{$ \blacktriangledown\,6.7\times $}}\\[-.2pt]18.6s\\[-.2pt]{\fontsize{5.8}{6.4}\selectfont \textcolor{BrickRed}{$ \blacktriangledown\,8.4\times $}}\end{tabular} & \hspace*{1pt}\begin{tabular}[t]{@{}c@{}}\rule{0pt}{12pt}19.4s\\[-.2pt]{\fontsize{5.8}{6.4}\selectfont \textcolor{BrickRed}{$ \blacktriangledown\,7.2\times $}}\\[-.2pt]17.5s\\[-.2pt]{\fontsize{5.8}{6.4}\selectfont \textcolor{BrickRed}{$ \blacktriangledown\,9.1\times $}}\end{tabular} & \hspace*{1pt}\begin{tabular}[t]{@{}c@{}}\rule{0pt}{12pt}1m\,12.7s\\[-.2pt]{\fontsize{5.8}{6.4}\selectfont \textcolor{BrickRed}{$ \blacktriangledown\,13.6\times $}}\\[-.2pt]1m\,8.1s\\[-.2pt]{\fontsize{5.8}{6.4}\selectfont \textcolor{BrickRed}{$ \blacktriangledown\,20.2\times $}}\end{tabular} & \hspace*{1pt}\begin{tabular}[t]{@{}c@{}}\rule{0pt}{12pt}7m\,25.4s\\[-.2pt]{\fontsize{5.8}{6.4}\selectfont \textcolor{BrickRed}{$ \blacktriangledown\,106.2\times $}}\\[-.2pt]6m\,50.8s\\[-.2pt]{\fontsize{5.8}{6.4}\selectfont \textcolor{BrickRed}{$ \blacktriangledown\,212.9\times $}}\end{tabular} & \hspace*{1pt}\begin{tabular}[t]{@{}c@{}}\rule{0pt}{12pt}9.3s\\[-.2pt]{\fontsize{5.8}{6.4}\selectfont \textcolor{BrickRed}{$ \blacktriangledown\,24\% $}}\\[-.2pt]8.4s\\[-.2pt]{\fontsize{5.8}{6.4}\selectfont \textcolor{BrickRed}{$ \blacktriangledown\,1.6\times $}}\end{tabular} & \hspace*{1pt}\begin{tabular}[t]{@{}c@{}}\rule{0pt}{12pt}8.7s\\[-.2pt]{\fontsize{5.8}{6.4}\selectfont \textcolor{BrickRed}{$ \blacktriangledown\,1.5\times $}}\\[-.2pt]7s\\[-.2pt]{\fontsize{5.8}{6.4}\selectfont \textcolor{BrickRed}{$ \blacktriangledown\,1.7\times $}}\end{tabular} & \hspace*{1pt}\begin{tabular}[t]{@{}c@{}}\rule{0pt}{12pt}15.8s\\[-.2pt]{\fontsize{5.8}{6.4}\selectfont \textcolor{BrickRed}{$ \blacktriangledown\,26\% $}}\\[-.2pt]12.7s\\[-.2pt]{\fontsize{5.8}{6.4}\selectfont \textcolor{BrickRed}{$ \blacktriangledown\,1.6\times $}}\end{tabular} & \hspace*{1pt}\begin{tabular}[t]{@{}c@{}}\rule{0pt}{12pt}18.5s\\[-.2pt]{\fontsize{5.8}{6.4}\selectfont \textcolor{BrickRed}{$ \blacktriangledown\,1.6\times $}}\\[-.2pt]17s\\[-.2pt]{\fontsize{5.8}{6.4}\selectfont \textcolor{BrickRed}{$ \blacktriangledown\,1.8\times $}}\end{tabular} & \hspace*{1pt}\begin{tabular}[t]{@{}c@{}}\rule{0pt}{12pt}42.4s\\[-.2pt]{\fontsize{5.8}{6.4}\selectfont \textcolor{ForestGreen}{$ \blacktriangle\,2.0\times $}}\\[-.2pt]41.4s\\[-.2pt]{\fontsize{5.8}{6.4}\selectfont \textcolor{ForestGreen}{$ \blacktriangle\,1.9\times $}}\end{tabular} & \hspace*{1pt}\begin{tabular}[t]{@{}c@{}}\rule{0pt}{12pt}42.9s\\[-.2pt]{\fontsize{5.8}{6.4}\selectfont \textcolor{ForestGreen}{$ \blacktriangle\,2.0\times $}}\\[-.2pt]41.1s\\[-.2pt]{\fontsize{5.8}{6.4}\selectfont \textcolor{ForestGreen}{$ \blacktriangle\,1.8\times $}}\end{tabular} & \hspace*{1pt}\begin{tabular}[t]{@{}c@{}}\rule{0pt}{12pt}1m\,26.9s\\[-.2pt]{\fontsize{5.8}{6.4}\selectfont \textcolor{ForestGreen}{$ \blacktriangle\,1.9\times $}}\\[-.2pt]1m\,23.7s\\[-.2pt]{\fontsize{5.8}{6.4}\selectfont \textcolor{ForestGreen}{$ \blacktriangle\,1.7\times $}}\end{tabular} & \hspace*{1pt}\begin{tabular}[t]{@{}c@{}}\rule{0pt}{12pt}4m\,51.6s\\[-.2pt]{\fontsize{5.8}{6.4}\selectfont \textcolor{gray}{n/a}}\\[-.2pt]4m\,46.6s\\[-.2pt]{\fontsize{5.8}{6.4}\selectfont \textcolor{gray}{n/a}}\end{tabular} & \hspace*{1pt}\begin{tabular}[t]{@{}c@{}}\rule{0pt}{12pt}35s\\[-.2pt]{\fontsize{5.8}{6.4}\selectfont \textcolor{ForestGreen}{$ \blacktriangle\,8.6\times $}}\\[-.2pt]34s\\[-.2pt]{\fontsize{5.8}{6.4}\selectfont \textcolor{ForestGreen}{$ \blacktriangle\,2.5\times $}}\end{tabular} & \hspace*{1pt}\begin{tabular}[t]{@{}c@{}}\rule{0pt}{12pt}35.9s\\[-.2pt]{\fontsize{5.8}{6.4}\selectfont \textcolor{ForestGreen}{$ \blacktriangle\,6.7\times $}}\\[-.2pt]34.2s\\[-.2pt]{\fontsize{5.8}{6.4}\selectfont \textcolor{ForestGreen}{$ \blacktriangle\,2.5\times $}}\end{tabular} & \hspace*{1pt}\begin{tabular}[t]{@{}c@{}}\rule{0pt}{12pt}1m\,14.4s\\[-.2pt]{\fontsize{5.8}{6.4}\selectfont \textcolor{ForestGreen}{$ \blacktriangle\,6.9\times $}}\\[-.2pt]1m\,11.3s\\[-.2pt]{\fontsize{5.8}{6.4}\selectfont \textcolor{ForestGreen}{$ \blacktriangle\,2.2\times $}}\end{tabular} & \hspace*{1pt}\begin{tabular}[t]{@{}c@{}}\rule{0pt}{12pt}1m\,0.1s\\[-.2pt]{\fontsize{5.8}{6.4}\selectfont \textcolor{gray}{n/a}}\\[-.2pt]58.6s\\[-.2pt]{\fontsize{5.8}{6.4}\selectfont \textcolor{gray}{n/a}}\end{tabular}\\[3pt]\hline
\textbf{BioseqDB} & \begin{tabular}[t]{@{}l@{}}\rule{0pt}{12pt}MPS / Bounded\\RMQ / KSW2\end{tabular} & \hspace*{1pt}\begin{tabular}[t]{@{}c@{}}\rule{0pt}{12pt}8.2s\\[-.2pt]{\fontsize{5.8}{6.4}\selectfont \textcolor{ForestGreen}{$ \blacktriangle\,2.4\times $}}\\[-.2pt]5.7s\\[-.2pt]{\fontsize{5.8}{6.4}\selectfont \textcolor{ForestGreen}{$ \blacktriangle\,3.3\times $}}\end{tabular} & \hspace*{1pt}\begin{tabular}[t]{@{}c@{}}\rule{0pt}{12pt}8s\\[-.2pt]{\fontsize{5.8}{6.4}\selectfont \textcolor{ForestGreen}{$ \blacktriangle\,2.4\times $}}\\[-.2pt]5.6s\\[-.2pt]{\fontsize{5.8}{6.4}\selectfont \textcolor{ForestGreen}{$ \blacktriangle\,3.1\times $}}\end{tabular} & \textcolor{gray}{\ensuremath{\cdot}} & \hspace*{1pt}\begin{tabular}[t]{@{}c@{}}\rule{0pt}{12pt}8.3s\\[-.2pt]{\fontsize{5.8}{6.4}\selectfont \textcolor{ForestGreen}{$ \blacktriangle\,53.4\times $}}\\[-.2pt]5.3s\\[-.2pt]{\fontsize{5.8}{6.4}\selectfont \textcolor{ForestGreen}{$ \blacktriangle\,77.4\times $}}\end{tabular} & \textcolor{gray}{\ensuremath{\cdot}} & \textcolor{gray}{\ensuremath{\cdot}} & \textcolor{gray}{\ensuremath{\cdot}} & \textcolor{gray}{\ensuremath{\cdot}} & \hspace*{1pt}\begin{tabular}[t]{@{}c@{}}\rule{0pt}{12pt}1m\,48.4s\\[-.2pt]{\fontsize{5.8}{6.4}\selectfont \textcolor{BrickRed}{$ \blacktriangledown\,2.6\times $}}\\[-.2pt]1m\,24.7s\\[-.2pt]{\fontsize{5.8}{6.4}\selectfont \textcolor{BrickRed}{$ \blacktriangledown\,2.0\times $}}\end{tabular} & \hspace*{1pt}\begin{tabular}[t]{@{}c@{}}\rule{0pt}{12pt}1m\,39.1s\\[-.2pt]{\fontsize{5.8}{6.4}\selectfont \textcolor{BrickRed}{$ \blacktriangledown\,2.3\times $}}\\[-.2pt]1m\,18s\\[-.2pt]{\fontsize{5.8}{6.4}\selectfont \textcolor{BrickRed}{$ \blacktriangledown\,1.9\times $}}\end{tabular} & \textcolor{gray}{\ensuremath{\cdot}} & \textcolor{gray}{\ensuremath{\cdot}} & \textcolor{gray}{\ensuremath{\cdot}} & \textcolor{gray}{\ensuremath{\cdot}} & \textcolor{gray}{\ensuremath{\cdot}} & \textcolor{gray}{\ensuremath{\cdot}}\\[3pt]\hline
 & \begin{tabular}[t]{@{}l@{}}\rule{0pt}{12pt}MPS / Bounded\\Scalar-DP / Block\end{tabular} & \hspace*{1pt}\begin{tabular}[t]{@{}c@{}}\rule{0pt}{12pt}6.6s\\[-.2pt]{\fontsize{5.8}{6.4}\selectfont \textcolor{ForestGreen}{$ \blacktriangle\,3.0\times $}}\\[-.2pt]5.9s\\[-.2pt]{\fontsize{5.8}{6.4}\selectfont \textcolor{ForestGreen}{$ \blacktriangle\,3.2\times $}}\end{tabular} & \hspace*{1pt}\begin{tabular}[t]{@{}c@{}}\rule{0pt}{12pt}6.6s\\[-.2pt]{\fontsize{5.8}{6.4}\selectfont \textcolor{ForestGreen}{$ \blacktriangle\,2.9\times $}}\\[-.2pt]5.9s\\[-.2pt]{\fontsize{5.8}{6.4}\selectfont \textcolor{ForestGreen}{$ \blacktriangle\,2.9\times $}}\end{tabular} & \textcolor{gray}{\ensuremath{\cdot}} & \hspace*{1pt}\begin{tabular}[t]{@{}c@{}}\rule{0pt}{12pt}6.3s\\[-.2pt]{\fontsize{5.8}{6.4}\selectfont \textcolor{ForestGreen}{$ \blacktriangle\,70.5\times $}}\\[-.2pt]5.4s\\[-.2pt]{\fontsize{5.8}{6.4}\selectfont \textcolor{ForestGreen}{$ \blacktriangle\,75.7\times $}}\end{tabular} & \textcolor{gray}{\ensuremath{\cdot}} & \textcolor{gray}{\ensuremath{\cdot}} & \textcolor{gray}{\ensuremath{\cdot}} & \textcolor{gray}{\ensuremath{\cdot}} & \hspace*{1pt}\begin{tabular}[t]{@{}c@{}}\rule{0pt}{12pt}1m\,42.1s\\[-.2pt]{\fontsize{5.8}{6.4}\selectfont \textcolor{BrickRed}{$ \blacktriangledown\,2.4\times $}}\\[-.2pt]1m\,20s\\[-.2pt]{\fontsize{5.8}{6.4}\selectfont \textcolor{BrickRed}{$ \blacktriangledown\,1.9\times $}}\end{tabular} & \hspace*{1pt}\begin{tabular}[t]{@{}c@{}}\rule{0pt}{12pt}1m\,37.4s\\[-.2pt]{\fontsize{5.8}{6.4}\selectfont \textcolor{BrickRed}{$ \blacktriangledown\,2.3\times $}}\\[-.2pt]1m\,18s\\[-.2pt]{\fontsize{5.8}{6.4}\selectfont \textcolor{BrickRed}{$ \blacktriangledown\,1.9\times $}}\end{tabular} & \textcolor{gray}{\ensuremath{\cdot}} & \textcolor{gray}{\ensuremath{\cdot}} & \textcolor{gray}{\ensuremath{\cdot}} & \textcolor{gray}{\ensuremath{\cdot}} & \textcolor{gray}{\ensuremath{\cdot}} & \textcolor{gray}{\ensuremath{\cdot}}\\[3pt]\hline
 & \begin{tabular}[t]{@{}l@{}}\rule{0pt}{12pt}MPS / Bounded\\Scalar-DP / KSW2\end{tabular} & \hspace*{1pt}\begin{tabular}[t]{@{}c@{}}\rule{0pt}{12pt}8.2s\\[-.2pt]{\fontsize{5.8}{6.4}\selectfont \textcolor{ForestGreen}{$ \blacktriangle\,2.4\times $}}\\[-.2pt]5.7s\\[-.2pt]{\fontsize{5.8}{6.4}\selectfont \textcolor{ForestGreen}{$ \blacktriangle\,3.3\times $}}\end{tabular} & \hspace*{1pt}\begin{tabular}[t]{@{}c@{}}\rule{0pt}{12pt}8.1s\\[-.2pt]{\fontsize{5.8}{6.4}\selectfont \textcolor{ForestGreen}{$ \blacktriangle\,2.4\times $}}\\[-.2pt]5.6s\\[-.2pt]{\fontsize{5.8}{6.4}\selectfont \textcolor{ForestGreen}{$ \blacktriangle\,3.1\times $}}\end{tabular} & \textcolor{gray}{\ensuremath{\cdot}} & \hspace*{1pt}\begin{tabular}[t]{@{}c@{}}\rule{0pt}{12pt}8.4s\\[-.2pt]{\fontsize{5.8}{6.4}\selectfont \textcolor{ForestGreen}{$ \blacktriangle\,53.2\times $}}\\[-.2pt]5.3s\\[-.2pt]{\fontsize{5.8}{6.4}\selectfont \textcolor{ForestGreen}{$ \blacktriangle\,77.0\times $}}\end{tabular} & \textcolor{gray}{\ensuremath{\cdot}} & \textcolor{gray}{\ensuremath{\cdot}} & \textcolor{gray}{\ensuremath{\cdot}} & \textcolor{gray}{\ensuremath{\cdot}} & \hspace*{1pt}\begin{tabular}[t]{@{}c@{}}\rule{0pt}{12pt}1m\,50.8s\\[-.2pt]{\fontsize{5.8}{6.4}\selectfont \textcolor{BrickRed}{$ \blacktriangledown\,2.6\times $}}\\[-.2pt]1m\,27.1s\\[-.2pt]{\fontsize{5.8}{6.4}\selectfont \textcolor{BrickRed}{$ \blacktriangledown\,2.1\times $}}\end{tabular} & \hspace*{1pt}\begin{tabular}[t]{@{}c@{}}\rule{0pt}{12pt}1m\,43.6s\\[-.2pt]{\fontsize{5.8}{6.4}\selectfont \textcolor{BrickRed}{$ \blacktriangledown\,2.4\times $}}\\[-.2pt]1m\,21.8s\\[-.2pt]{\fontsize{5.8}{6.4}\selectfont \textcolor{BrickRed}{$ \blacktriangledown\,2.0\times $}}\end{tabular} & \textcolor{gray}{\ensuremath{\cdot}} & \textcolor{gray}{\ensuremath{\cdot}} & \textcolor{gray}{\ensuremath{\cdot}} & \textcolor{gray}{\ensuremath{\cdot}} & \textcolor{gray}{\ensuremath{\cdot}} & \textcolor{gray}{\ensuremath{\cdot}}\\[3pt]\hline
 & \begin{tabular}[t]{@{}l@{}}\rule{0pt}{12pt}MPS / Bounded\\SIMD-DP / Block\end{tabular} & \hspace*{1pt}\begin{tabular}[t]{@{}c@{}}\rule{0pt}{12pt}6.5s\\[-.2pt]{\fontsize{5.8}{6.4}\selectfont \textcolor{ForestGreen}{$ \blacktriangle\,3.0\times $}}\\[-.2pt]5.8s\\[-.2pt]{\fontsize{5.8}{6.4}\selectfont \textcolor{ForestGreen}{$ \blacktriangle\,3.2\times $}}\end{tabular} & \hspace*{1pt}\begin{tabular}[t]{@{}c@{}}\rule{0pt}{12pt}6.5s\\[-.2pt]{\fontsize{5.8}{6.4}\selectfont \textcolor{ForestGreen}{$ \blacktriangle\,3.0\times $}}\\[-.2pt]5.8s\\[-.2pt]{\fontsize{5.8}{6.4}\selectfont \textcolor{ForestGreen}{$ \blacktriangle\,3.0\times $}}\end{tabular} & \textcolor{gray}{\ensuremath{\cdot}} & \hspace*{1pt}\begin{tabular}[t]{@{}c@{}}\rule{0pt}{12pt}6.4s\\[-.2pt]{\fontsize{5.8}{6.4}\selectfont \textcolor{ForestGreen}{$ \blacktriangle\,69.2\times $}}\\[-.2pt]5.5s\\[-.2pt]{\fontsize{5.8}{6.4}\selectfont \textcolor{ForestGreen}{$ \blacktriangle\,74.5\times $}}\end{tabular} & \textcolor{gray}{\ensuremath{\cdot}} & \textcolor{gray}{\ensuremath{\cdot}} & \textcolor{gray}{\ensuremath{\cdot}} & \textcolor{gray}{\ensuremath{\cdot}} & \textcolor{gray}{\ensuremath{\cdot}} & \textcolor{gray}{\ensuremath{\cdot}} & \textcolor{gray}{\ensuremath{\cdot}} & \textcolor{gray}{\ensuremath{\cdot}} & \textcolor{gray}{\ensuremath{\cdot}} & \textcolor{gray}{\ensuremath{\cdot}} & \hspace*{1pt}\begin{tabular}[t]{@{}c@{}}\rule{0pt}{12pt}8m\,50.5s\\[-.2pt]{\fontsize{5.8}{6.4}\selectfont \textcolor{BrickRed}{$ \blacktriangledown\,7.1\times $}}\\[-.2pt]2m\,35s\\[-.2pt]{\fontsize{5.8}{6.4}\selectfont \textcolor{BrickRed}{$ \blacktriangledown\,2.2\times $}}\end{tabular} & \textcolor{gray}{\ensuremath{\cdot}}\\[3pt]\hline
 & \begin{tabular}[t]{@{}l@{}}\rule{0pt}{12pt}MPS / Bounded\\SIMD-DP / KSW2\end{tabular} & \hspace*{1pt}\begin{tabular}[t]{@{}c@{}}\rule{0pt}{12pt}8.4s\\[-.2pt]{\fontsize{5.8}{6.4}\selectfont \textcolor{ForestGreen}{$ \blacktriangle\,2.3\times $}}\\[-.2pt]5.7s\\[-.2pt]{\fontsize{5.8}{6.4}\selectfont \textcolor{ForestGreen}{$ \blacktriangle\,3.3\times $}}\end{tabular} & \hspace*{1pt}\begin{tabular}[t]{@{}c@{}}\rule{0pt}{12pt}8s\\[-.2pt]{\fontsize{5.8}{6.4}\selectfont \textcolor{ForestGreen}{$ \blacktriangle\,2.4\times $}}\\[-.2pt]5.6s\\[-.2pt]{\fontsize{5.8}{6.4}\selectfont \textcolor{ForestGreen}{$ \blacktriangle\,3.1\times $}}\end{tabular} & \textcolor{gray}{\ensuremath{\cdot}} & \hspace*{1pt}\begin{tabular}[t]{@{}c@{}}\rule{0pt}{12pt}8.2s\\[-.2pt]{\fontsize{5.8}{6.4}\selectfont \textcolor{ForestGreen}{$ \blacktriangle\,54.1\times $}}\\[-.2pt]5.3s\\[-.2pt]{\fontsize{5.8}{6.4}\selectfont \textcolor{ForestGreen}{$ \blacktriangle\,77.9\times $}}\end{tabular} & \textcolor{gray}{\ensuremath{\cdot}} & \textcolor{gray}{\ensuremath{\cdot}} & \textcolor{gray}{\ensuremath{\cdot}} & \textcolor{gray}{\ensuremath{\cdot}} & \hspace*{1pt}\begin{tabular}[t]{@{}c@{}}\rule{0pt}{12pt}1m\,43.1s\\[-.2pt]{\fontsize{5.8}{6.4}\selectfont \textcolor{BrickRed}{$ \blacktriangledown\,2.4\times $}}\\[-.2pt]1m\,19.2s\\[-.2pt]{\fontsize{5.8}{6.4}\selectfont \textcolor{BrickRed}{$ \blacktriangledown\,1.9\times $}}\end{tabular} & \hspace*{1pt}\begin{tabular}[t]{@{}c@{}}\rule{0pt}{12pt}1m\,39.2s\\[-.2pt]{\fontsize{5.8}{6.4}\selectfont \textcolor{BrickRed}{$ \blacktriangledown\,2.3\times $}}\\[-.2pt]1m\,17.9s\\[-.2pt]{\fontsize{5.8}{6.4}\selectfont \textcolor{BrickRed}{$ \blacktriangledown\,1.9\times $}}\end{tabular} & \textcolor{gray}{\ensuremath{\cdot}} & \textcolor{gray}{\ensuremath{\cdot}} & \textcolor{gray}{\ensuremath{\cdot}} & \textcolor{gray}{\ensuremath{\cdot}} & \textcolor{gray}{\ensuremath{\cdot}} & \textcolor{gray}{\ensuremath{\cdot}}\\[3pt]\hline
 & \begin{tabular}[t]{@{}l@{}}\rule{0pt}{12pt}MPS / Scalar\\RMQ / Block\end{tabular} & \hspace*{1pt}\begin{tabular}[t]{@{}c@{}}\rule{0pt}{12pt}43.6s\\[-.2pt]{\fontsize{5.8}{6.4}\selectfont \textcolor{BrickRed}{$ \blacktriangledown\,2.2\times $}}\\[-.2pt]8.8s\\[-.2pt]{\fontsize{5.8}{6.4}\selectfont \textcolor{ForestGreen}{$ \blacktriangle\,2.1\times $}}\end{tabular} & \hspace*{1pt}\begin{tabular}[t]{@{}c@{}}\rule{0pt}{12pt}37.4s\\[-.2pt]{\fontsize{5.8}{6.4}\selectfont \textcolor{BrickRed}{$ \blacktriangledown\,1.9\times $}}\\[-.2pt]8.8s\\[-.2pt]{\fontsize{5.8}{6.4}\selectfont \textcolor{ForestGreen}{$ \blacktriangle\,2.0\times $}}\end{tabular} & \textcolor{gray}{\ensuremath{\cdot}} & \textcolor{gray}{\ensuremath{\cdot}} & \textcolor{gray}{\ensuremath{\cdot}} & \textcolor{gray}{\ensuremath{\cdot}} & \textcolor{gray}{\ensuremath{\cdot}} & \textcolor{gray}{\ensuremath{\cdot}} & \hspace*{1pt}\begin{tabular}[t]{@{}c@{}}\rule{0pt}{12pt}2m\,53.3s\\[-.2pt]{\fontsize{5.8}{6.4}\selectfont \textcolor{BrickRed}{$ \blacktriangledown\,4.1\times $}}\\[-.2pt]2m\,17.7s\\[-.2pt]{\fontsize{5.8}{6.4}\selectfont \textcolor{BrickRed}{$ \blacktriangledown\,3.3\times $}}\end{tabular} & \hspace*{1pt}\begin{tabular}[t]{@{}c@{}}\rule{0pt}{12pt}2m\,48s\\[-.2pt]{\fontsize{5.8}{6.4}\selectfont \textcolor{BrickRed}{$ \blacktriangledown\,3.9\times $}}\\[-.2pt]2m\,18.9s\\[-.2pt]{\fontsize{5.8}{6.4}\selectfont \textcolor{BrickRed}{$ \blacktriangledown\,3.4\times $}}\end{tabular} & \hspace*{1pt}\begin{tabular}[t]{@{}c@{}}\rule{0pt}{12pt}3m\,24s\\[-.2pt]{\fontsize{5.8}{6.4}\selectfont \textcolor{BrickRed}{$ \blacktriangledown\,2.3\times $}}\\[-.2pt]2m\,47.8s\\[-.2pt]{\fontsize{5.8}{6.4}\selectfont \textcolor{BrickRed}{$ \blacktriangledown\,2.0\times $}}\end{tabular} & \textcolor{gray}{\ensuremath{\cdot}} & \textcolor{gray}{\ensuremath{\cdot}} & \textcolor{gray}{\ensuremath{\cdot}} & \textcolor{gray}{\ensuremath{\cdot}} & \textcolor{gray}{\ensuremath{\cdot}}\\[3pt]\hline
 & \begin{tabular}[t]{@{}l@{}}\rule{0pt}{12pt}MPS / Scalar\\RMQ / KSW2\end{tabular} & \hspace*{1pt}\begin{tabular}[t]{@{}c@{}}\rule{0pt}{12pt}45.4s\\[-.2pt]{\fontsize{5.8}{6.4}\selectfont \textcolor{BrickRed}{$ \blacktriangledown\,2.3\times $}}\\[-.2pt]8.8s\\[-.2pt]{\fontsize{5.8}{6.4}\selectfont \textcolor{ForestGreen}{$ \blacktriangle\,2.1\times $}}\end{tabular} & \hspace*{1pt}\begin{tabular}[t]{@{}c@{}}\rule{0pt}{12pt}39.7s\\[-.2pt]{\fontsize{5.8}{6.4}\selectfont \textcolor{BrickRed}{$ \blacktriangledown\,2.1\times $}}\\[-.2pt]9s\\[-.2pt]{\fontsize{5.8}{6.4}\selectfont \textcolor{ForestGreen}{$ \blacktriangle\,2.0\times $}}\end{tabular} & \textcolor{gray}{\ensuremath{\cdot}} & \textcolor{gray}{\ensuremath{\cdot}} & \textcolor{gray}{\ensuremath{\cdot}} & \textcolor{gray}{\ensuremath{\cdot}} & \textcolor{gray}{\ensuremath{\cdot}} & \textcolor{gray}{\ensuremath{\cdot}} & \hspace*{1pt}\begin{tabular}[t]{@{}c@{}}\rule{0pt}{12pt}2m\,52.9s\\[-.2pt]{\fontsize{5.8}{6.4}\selectfont \textcolor{BrickRed}{$ \blacktriangledown\,4.1\times $}}\\[-.2pt]2m\,15.4s\\[-.2pt]{\fontsize{5.8}{6.4}\selectfont \textcolor{BrickRed}{$ \blacktriangledown\,3.3\times $}}\end{tabular} & \textcolor{gray}{\ensuremath{\cdot}} & \textcolor{gray}{\ensuremath{\cdot}} & \textcolor{gray}{\ensuremath{\cdot}} & \textcolor{gray}{\ensuremath{\cdot}} & \textcolor{gray}{\ensuremath{\cdot}} & \textcolor{gray}{\ensuremath{\cdot}} & \textcolor{gray}{\ensuremath{\cdot}}\\[3pt]\hline
\end{tabular}\par}
\smallskip{\footnotesize\noindent Each cell: A, comparison, B, comparison. Times cover 1,024 queries. $\cdot$ = no approved observation. ---\textsuperscript{g} = GPU unavailable, not executed. n/a = no measured comparison.\par
\noindent \textsuperscript{f}SIMD-ref request observed as scalar, compared only as that effective path. \textsuperscript{p}Accepted by explicit preservation after wrapper postflight failure, no run receipt claimed.\par
\noindent MPF/MPS = minimizers per fragment/sequence. Row order: index representation / minimizer, then chaining / base alignment. SIMD-ref is the reference minimizer, Bounded is bounded ntHash.\par
\noindent Blue = one lowest A and one lowest B down each machine/workload column, across continuation pages. Exact ties select the first row. Repeated baselines are marked only on page 1.\par
\noindent Light gray = minimap2 baseline. Pale yellow = preferred MPF / Bounded / SIMD-DP / Block configuration, not a claim of winning every column.\par
\noindent When all Vulkan modes are unmeasured, they share one row. Baseline comparisons use the fastest measured BioseqDB effective path separately for A and B.\par}
\clearpage
\noindent{\large\bfseries AWS: 2 Gbp, k=20, w=15 (mapping interval)\hfill 5/6}\par\smallskip
{\fontsize{7.1}{8.1}\selectfont\setlength{\tabcolsep}{1pt}\setlength{\fboxsep}{2pt}\renewcommand{\arraystretch}{1.12}
\noindent\begin{tabular}{@{}p{14mm}p{27mm}|p{13.55mm}|p{13.55mm}|p{13.55mm}|p{13.55mm}|p{13.55mm}|p{13.55mm}|p{13.55mm}|p{13.55mm}|p{13.55mm}|p{13.55mm}|p{13.55mm}|p{13.55mm}|p{13.55mm}|p{13.55mm}|p{13.55mm}|p{13.55mm}|@{}}\hline
Engine & Execution & \multicolumn{8}{c|}{\textbf{Index}} & \multicolumn{8}{c|}{\textbf{No index}}\\
 &  & \multicolumn{8}{c|}{\textbf{2 Gbp}} & \multicolumn{8}{c|}{\textbf{2 Gbp}}\\
 &  & \multicolumn{8}{c|}{\textbf{k=20, w=15}} & \multicolumn{8}{c|}{\textbf{k=20, w=15}}\\
 &  & \multicolumn{4}{c|}{\textbf{DNA4}} & \multicolumn{4}{c|}{\textbf{DNA5}} & \multicolumn{4}{c|}{\textbf{DNA4}} & \multicolumn{4}{c|}{\textbf{DNA5}}\\
 &  & \multicolumn{1}{c|}{\textbf{M7A}} & \multicolumn{1}{c|}{\textbf{M7I}} & \multicolumn{1}{c|}{\textbf{M4}} & \multicolumn{1}{c|}{\textbf{C7G}} & \multicolumn{1}{c|}{\textbf{M7A}} & \multicolumn{1}{c|}{\textbf{M7I}} & \multicolumn{1}{c|}{\textbf{M4}} & \multicolumn{1}{c|}{\textbf{C7G}} & \multicolumn{1}{c|}{\textbf{M7A}} & \multicolumn{1}{c|}{\textbf{M7I}} & \multicolumn{1}{c|}{\textbf{M4}} & \multicolumn{1}{c|}{\textbf{C7G}} & \multicolumn{1}{c|}{\textbf{M7A}} & \multicolumn{1}{c|}{\textbf{M7I}} & \multicolumn{1}{c|}{\textbf{M4}} & \multicolumn{1}{c|}{\textbf{C7G}}\\
\hline
\rowcolor{BenchBaseline}[0.5pt][0.5pt]\textbf{minimap2} & \textbf{Baseline} & \hspace*{1pt}\begin{tabular}[t]{@{}c@{}}\rule{0pt}{12pt}19.6s\\[-.2pt]{\fontsize{5.8}{6.4}\selectfont \textcolor{BrickRed}{$ \blacktriangledown\,6.7\times $}}\\[-.2pt]18.6s\\[-.2pt]{\fontsize{5.8}{6.4}\selectfont \textcolor{BrickRed}{$ \blacktriangledown\,8.4\times $}}\end{tabular} & \hspace*{1pt}\begin{tabular}[t]{@{}c@{}}\rule{0pt}{12pt}19.4s\\[-.2pt]{\fontsize{5.8}{6.4}\selectfont \textcolor{BrickRed}{$ \blacktriangledown\,7.2\times $}}\\[-.2pt]17.5s\\[-.2pt]{\fontsize{5.8}{6.4}\selectfont \textcolor{BrickRed}{$ \blacktriangledown\,9.1\times $}}\end{tabular} & \hspace*{1pt}\begin{tabular}[t]{@{}c@{}}\rule{0pt}{12pt}1m\,12.7s\\[-.2pt]{\fontsize{5.8}{6.4}\selectfont \textcolor{BrickRed}{$ \blacktriangledown\,13.6\times $}}\\[-.2pt]1m\,8.1s\\[-.2pt]{\fontsize{5.8}{6.4}\selectfont \textcolor{BrickRed}{$ \blacktriangledown\,20.2\times $}}\end{tabular} & \hspace*{1pt}\begin{tabular}[t]{@{}c@{}}\rule{0pt}{12pt}7m\,25.4s\\[-.2pt]{\fontsize{5.8}{6.4}\selectfont \textcolor{BrickRed}{$ \blacktriangledown\,106.2\times $}}\\[-.2pt]6m\,50.8s\\[-.2pt]{\fontsize{5.8}{6.4}\selectfont \textcolor{BrickRed}{$ \blacktriangledown\,212.9\times $}}\end{tabular} & \hspace*{1pt}\begin{tabular}[t]{@{}c@{}}\rule{0pt}{12pt}9.3s\\[-.2pt]{\fontsize{5.8}{6.4}\selectfont \textcolor{BrickRed}{$ \blacktriangledown\,24\% $}}\\[-.2pt]8.4s\\[-.2pt]{\fontsize{5.8}{6.4}\selectfont \textcolor{BrickRed}{$ \blacktriangledown\,1.6\times $}}\end{tabular} & \hspace*{1pt}\begin{tabular}[t]{@{}c@{}}\rule{0pt}{12pt}8.7s\\[-.2pt]{\fontsize{5.8}{6.4}\selectfont \textcolor{BrickRed}{$ \blacktriangledown\,1.5\times $}}\\[-.2pt]7s\\[-.2pt]{\fontsize{5.8}{6.4}\selectfont \textcolor{BrickRed}{$ \blacktriangledown\,1.7\times $}}\end{tabular} & \hspace*{1pt}\begin{tabular}[t]{@{}c@{}}\rule{0pt}{12pt}15.8s\\[-.2pt]{\fontsize{5.8}{6.4}\selectfont \textcolor{BrickRed}{$ \blacktriangledown\,26\% $}}\\[-.2pt]12.7s\\[-.2pt]{\fontsize{5.8}{6.4}\selectfont \textcolor{BrickRed}{$ \blacktriangledown\,1.6\times $}}\end{tabular} & \hspace*{1pt}\begin{tabular}[t]{@{}c@{}}\rule{0pt}{12pt}18.5s\\[-.2pt]{\fontsize{5.8}{6.4}\selectfont \textcolor{BrickRed}{$ \blacktriangledown\,1.6\times $}}\\[-.2pt]17s\\[-.2pt]{\fontsize{5.8}{6.4}\selectfont \textcolor{BrickRed}{$ \blacktriangledown\,1.8\times $}}\end{tabular} & \hspace*{1pt}\begin{tabular}[t]{@{}c@{}}\rule{0pt}{12pt}42.4s\\[-.2pt]{\fontsize{5.8}{6.4}\selectfont \textcolor{ForestGreen}{$ \blacktriangle\,2.0\times $}}\\[-.2pt]41.4s\\[-.2pt]{\fontsize{5.8}{6.4}\selectfont \textcolor{ForestGreen}{$ \blacktriangle\,1.9\times $}}\end{tabular} & \hspace*{1pt}\begin{tabular}[t]{@{}c@{}}\rule{0pt}{12pt}42.9s\\[-.2pt]{\fontsize{5.8}{6.4}\selectfont \textcolor{ForestGreen}{$ \blacktriangle\,2.0\times $}}\\[-.2pt]41.1s\\[-.2pt]{\fontsize{5.8}{6.4}\selectfont \textcolor{ForestGreen}{$ \blacktriangle\,1.8\times $}}\end{tabular} & \hspace*{1pt}\begin{tabular}[t]{@{}c@{}}\rule{0pt}{12pt}1m\,26.9s\\[-.2pt]{\fontsize{5.8}{6.4}\selectfont \textcolor{ForestGreen}{$ \blacktriangle\,1.9\times $}}\\[-.2pt]1m\,23.7s\\[-.2pt]{\fontsize{5.8}{6.4}\selectfont \textcolor{ForestGreen}{$ \blacktriangle\,1.7\times $}}\end{tabular} & \hspace*{1pt}\begin{tabular}[t]{@{}c@{}}\rule{0pt}{12pt}4m\,51.6s\\[-.2pt]{\fontsize{5.8}{6.4}\selectfont \textcolor{gray}{n/a}}\\[-.2pt]4m\,46.6s\\[-.2pt]{\fontsize{5.8}{6.4}\selectfont \textcolor{gray}{n/a}}\end{tabular} & \hspace*{1pt}\begin{tabular}[t]{@{}c@{}}\rule{0pt}{12pt}35s\\[-.2pt]{\fontsize{5.8}{6.4}\selectfont \textcolor{ForestGreen}{$ \blacktriangle\,8.6\times $}}\\[-.2pt]34s\\[-.2pt]{\fontsize{5.8}{6.4}\selectfont \textcolor{ForestGreen}{$ \blacktriangle\,2.5\times $}}\end{tabular} & \hspace*{1pt}\begin{tabular}[t]{@{}c@{}}\rule{0pt}{12pt}35.9s\\[-.2pt]{\fontsize{5.8}{6.4}\selectfont \textcolor{ForestGreen}{$ \blacktriangle\,6.7\times $}}\\[-.2pt]34.2s\\[-.2pt]{\fontsize{5.8}{6.4}\selectfont \textcolor{ForestGreen}{$ \blacktriangle\,2.5\times $}}\end{tabular} & \hspace*{1pt}\begin{tabular}[t]{@{}c@{}}\rule{0pt}{12pt}1m\,14.4s\\[-.2pt]{\fontsize{5.8}{6.4}\selectfont \textcolor{ForestGreen}{$ \blacktriangle\,6.9\times $}}\\[-.2pt]1m\,11.3s\\[-.2pt]{\fontsize{5.8}{6.4}\selectfont \textcolor{ForestGreen}{$ \blacktriangle\,2.2\times $}}\end{tabular} & \hspace*{1pt}\begin{tabular}[t]{@{}c@{}}\rule{0pt}{12pt}1m\,0.1s\\[-.2pt]{\fontsize{5.8}{6.4}\selectfont \textcolor{gray}{n/a}}\\[-.2pt]58.6s\\[-.2pt]{\fontsize{5.8}{6.4}\selectfont \textcolor{gray}{n/a}}\end{tabular}\\[3pt]\hline
\textbf{BioseqDB} & \begin{tabular}[t]{@{}l@{}}\rule{0pt}{12pt}MPS / Scalar\\Scalar-DP / Block\end{tabular} & \hspace*{1pt}\begin{tabular}[t]{@{}c@{}}\rule{0pt}{12pt}43.3s\\[-.2pt]{\fontsize{5.8}{6.4}\selectfont \textcolor{BrickRed}{$ \blacktriangledown\,2.2\times $}}\\[-.2pt]8.7s\\[-.2pt]{\fontsize{5.8}{6.4}\selectfont \textcolor{ForestGreen}{$ \blacktriangle\,2.2\times $}}\end{tabular} & \hspace*{1pt}\begin{tabular}[t]{@{}c@{}}\rule{0pt}{12pt}37.5s\\[-.2pt]{\fontsize{5.8}{6.4}\selectfont \textcolor{BrickRed}{$ \blacktriangledown\,1.9\times $}}\\[-.2pt]8.9s\\[-.2pt]{\fontsize{5.8}{6.4}\selectfont \textcolor{ForestGreen}{$ \blacktriangle\,2.0\times $}}\end{tabular} & \textcolor{gray}{\ensuremath{\cdot}} & \textcolor{gray}{\ensuremath{\cdot}} & \textcolor{gray}{\ensuremath{\cdot}} & \textcolor{gray}{\ensuremath{\cdot}} & \textcolor{gray}{\ensuremath{\cdot}} & \textcolor{gray}{\ensuremath{\cdot}} & \textcolor{gray}{\ensuremath{\cdot}} & \textcolor{gray}{\ensuremath{\cdot}} & \textcolor{gray}{\ensuremath{\cdot}} & \textcolor{gray}{\ensuremath{\cdot}} & \textcolor{gray}{\ensuremath{\cdot}} & \textcolor{gray}{\ensuremath{\cdot}} & \textcolor{gray}{\ensuremath{\cdot}} & \textcolor{gray}{\ensuremath{\cdot}}\\[3pt]\hline
 & \begin{tabular}[t]{@{}l@{}}\rule{0pt}{12pt}MPS / Scalar\\Scalar-DP / KSW2\end{tabular} & \hspace*{1pt}\begin{tabular}[t]{@{}c@{}}\rule{0pt}{12pt}45.2s\\[-.2pt]{\fontsize{5.8}{6.4}\selectfont \textcolor{BrickRed}{$ \blacktriangledown\,2.3\times $}}\\[-.2pt]8.7s\\[-.2pt]{\fontsize{5.8}{6.4}\selectfont \textcolor{ForestGreen}{$ \blacktriangle\,2.1\times $}}\end{tabular} & \hspace*{1pt}\begin{tabular}[t]{@{}c@{}}\rule{0pt}{12pt}39.2s\\[-.2pt]{\fontsize{5.8}{6.4}\selectfont \textcolor{BrickRed}{$ \blacktriangledown\,2.0\times $}}\\[-.2pt]8.8s\\[-.2pt]{\fontsize{5.8}{6.4}\selectfont \textcolor{ForestGreen}{$ \blacktriangle\,2.0\times $}}\end{tabular} & \hspace*{1pt}\begin{tabular}[t]{@{}c@{}}\rule{0pt}{12pt}1m\,6.9s\\[-.2pt]{\fontsize{5.8}{6.4}\selectfont \textcolor{ForestGreen}{$ \blacktriangle\,8\% $}}\\[-.2pt]14.8s\\[-.2pt]{\fontsize{5.8}{6.4}\selectfont \textcolor{ForestGreen}{$ \blacktriangle\,4.6\times $}}\end{tabular} & \textcolor{gray}{\ensuremath{\cdot}} & \textcolor{gray}{\ensuremath{\cdot}} & \textcolor{gray}{\ensuremath{\cdot}} & \textcolor{gray}{\ensuremath{\cdot}} & \textcolor{gray}{\ensuremath{\cdot}} & \textcolor{gray}{\ensuremath{\cdot}} & \textcolor{gray}{\ensuremath{\cdot}} & \textcolor{gray}{\ensuremath{\cdot}} & \textcolor{gray}{\ensuremath{\cdot}} & \textcolor{gray}{\ensuremath{\cdot}} & \textcolor{gray}{\ensuremath{\cdot}} & \textcolor{gray}{\ensuremath{\cdot}} & \textcolor{gray}{\ensuremath{\cdot}}\\[3pt]\hline
 & \begin{tabular}[t]{@{}l@{}}\rule{0pt}{12pt}MPS / Scalar\\SIMD-DP / Block\end{tabular} & \hspace*{1pt}\begin{tabular}[t]{@{}c@{}}\rule{0pt}{12pt}43.8s\\[-.2pt]{\fontsize{5.8}{6.4}\selectfont \textcolor{BrickRed}{$ \blacktriangledown\,2.2\times $}}\\[-.2pt]8.6s\\[-.2pt]{\fontsize{5.8}{6.4}\selectfont \textcolor{ForestGreen}{$ \blacktriangle\,2.2\times $}}\end{tabular} & \hspace*{1pt}\begin{tabular}[t]{@{}c@{}}\rule{0pt}{12pt}38s\\[-.2pt]{\fontsize{5.8}{6.4}\selectfont \textcolor{BrickRed}{$ \blacktriangledown\,2.0\times $}}\\[-.2pt]8.7s\\[-.2pt]{\fontsize{5.8}{6.4}\selectfont \textcolor{ForestGreen}{$ \blacktriangle\,2.0\times $}}\end{tabular} & \textcolor{gray}{\ensuremath{\cdot}} & \hspace*{1pt}\begin{tabular}[t]{@{}c@{}}\rule{0pt}{12pt}36.6s\\[-.2pt]{\fontsize{5.8}{6.4}\selectfont \textcolor{ForestGreen}{$ \blacktriangle\,12.2\times $}}\\[-.2pt]8.2s\\[-.2pt]{\fontsize{5.8}{6.4}\selectfont \textcolor{ForestGreen}{$ \blacktriangle\,49.9\times $}}\end{tabular} & \textcolor{gray}{\ensuremath{\cdot}} & \textcolor{gray}{\ensuremath{\cdot}} & \textcolor{gray}{\ensuremath{\cdot}} & \textcolor{gray}{\ensuremath{\cdot}} & \hspace*{1pt}\begin{tabular}[t]{@{}c@{}}\rule{0pt}{12pt}2m\,52.9s\\[-.2pt]{\fontsize{5.8}{6.4}\selectfont \textcolor{BrickRed}{$ \blacktriangledown\,4.1\times $}}\\[-.2pt]2m\,16.2s\\[-.2pt]{\fontsize{5.8}{6.4}\selectfont \textcolor{BrickRed}{$ \blacktriangledown\,3.3\times $}}\end{tabular} & \hspace*{1pt}\begin{tabular}[t]{@{}c@{}}\rule{0pt}{12pt}2m\,34.3s\\[-.2pt]{\fontsize{5.8}{6.4}\selectfont \textcolor{BrickRed}{$ \blacktriangledown\,3.6\times $}}\\[-.2pt]2m\,5.6s\\[-.2pt]{\fontsize{5.8}{6.4}\selectfont \textcolor{BrickRed}{$ \blacktriangledown\,3.1\times $}}\end{tabular} & \textcolor{gray}{\ensuremath{\cdot}} & \textcolor{gray}{\ensuremath{\cdot}} & \textcolor{gray}{\ensuremath{\cdot}} & \textcolor{gray}{\ensuremath{\cdot}} & \textcolor{gray}{\ensuremath{\cdot}} & \textcolor{gray}{\ensuremath{\cdot}}\\[3pt]\hline
 & \begin{tabular}[t]{@{}l@{}}\rule{0pt}{12pt}MPS / Scalar\\SIMD-DP / KSW2\end{tabular} & \hspace*{1pt}\begin{tabular}[t]{@{}c@{}}\rule{0pt}{12pt}45.1s\\[-.2pt]{\fontsize{5.8}{6.4}\selectfont \textcolor{BrickRed}{$ \blacktriangledown\,2.3\times $}}\\[-.2pt]8.4s\\[-.2pt]{\fontsize{5.8}{6.4}\selectfont \textcolor{ForestGreen}{$ \blacktriangle\,2.2\times $}}\end{tabular} & \hspace*{1pt}\begin{tabular}[t]{@{}c@{}}\rule{0pt}{12pt}37.9s\\[-.2pt]{\fontsize{5.8}{6.4}\selectfont \textcolor{BrickRed}{$ \blacktriangledown\,2.0\times $}}\\[-.2pt]8.2s\\[-.2pt]{\fontsize{5.8}{6.4}\selectfont \textcolor{ForestGreen}{$ \blacktriangle\,2.1\times $}}\end{tabular} & \hspace*{1pt}\begin{tabular}[t]{@{}c@{}}\rule{0pt}{12pt}1m\,6.6s\\[-.2pt]{\fontsize{5.8}{6.4}\selectfont \textcolor{ForestGreen}{$ \blacktriangle\,8\% $}}\\[-.2pt]14.5s\\[-.2pt]{\fontsize{5.8}{6.4}\selectfont \textcolor{ForestGreen}{$ \blacktriangle\,4.7\times $}}\end{tabular} & \hspace*{1pt}\begin{tabular}[t]{@{}c@{}}\rule{0pt}{12pt}40.5s\\[-.2pt]{\fontsize{5.8}{6.4}\selectfont \textcolor{ForestGreen}{$ \blacktriangle\,11.0\times $}}\\[-.2pt]8.3s\\[-.2pt]{\fontsize{5.8}{6.4}\selectfont \textcolor{ForestGreen}{$ \blacktriangle\,49.8\times $}}\end{tabular} & \textcolor{gray}{\ensuremath{\cdot}} & \textcolor{gray}{\ensuremath{\cdot}} & \textcolor{gray}{\ensuremath{\cdot}} & \textcolor{gray}{\ensuremath{\cdot}} & \hspace*{1pt}\begin{tabular}[t]{@{}c@{}}\rule{0pt}{12pt}2m\,54.8s\\[-.2pt]{\fontsize{5.8}{6.4}\selectfont \textcolor{BrickRed}{$ \blacktriangledown\,4.1\times $}}\\[-.2pt]2m\,18.1s\\[-.2pt]{\fontsize{5.8}{6.4}\selectfont \textcolor{BrickRed}{$ \blacktriangledown\,3.3\times $}}\end{tabular} & \hspace*{1pt}\begin{tabular}[t]{@{}c@{}}\rule{0pt}{12pt}2m\,39.7s\\[-.2pt]{\fontsize{5.8}{6.4}\selectfont \textcolor{BrickRed}{$ \blacktriangledown\,3.7\times $}}\\[-.2pt]2m\,8.9s\\[-.2pt]{\fontsize{5.8}{6.4}\selectfont \textcolor{BrickRed}{$ \blacktriangledown\,3.1\times $}}\end{tabular} & \hspace*{1pt}\begin{tabular}[t]{@{}c@{}}\rule{0pt}{12pt}4m\,27.6s\\[-.2pt]{\fontsize{5.8}{6.4}\selectfont \textcolor{BrickRed}{$ \blacktriangledown\,3.1\times $}}\\[-.2pt]3m\,34.6s\\[-.2pt]{\fontsize{5.8}{6.4}\selectfont \textcolor{BrickRed}{$ \blacktriangledown\,2.6\times $}}\end{tabular} & \textcolor{gray}{\ensuremath{\cdot}} & \textcolor{gray}{\ensuremath{\cdot}} & \textcolor{gray}{\ensuremath{\cdot}} & \textcolor{gray}{\ensuremath{\cdot}} & \textcolor{gray}{\ensuremath{\cdot}}\\[3pt]\hline
 & \begin{tabular}[t]{@{}l@{}}\rule{0pt}{12pt}MPS / SIMD-ref\\RMQ / Block\end{tabular} & \textcolor{gray}{\ensuremath{\cdot}} & \textcolor{gray}{\ensuremath{\cdot}} & \textcolor{gray}{\ensuremath{\cdot}} & \textcolor{gray}{\ensuremath{\cdot}} & \textcolor{gray}{\ensuremath{\cdot}} & \textcolor{gray}{\ensuremath{\cdot}} & \textcolor{gray}{\ensuremath{\cdot}} & \textcolor{gray}{\ensuremath{\cdot}} & \textcolor{gray}{\ensuremath{\cdot}} & \textcolor{gray}{\ensuremath{\cdot}} & \hspace*{1pt}\begin{tabular}[t]{@{}c@{}}\rule{0pt}{12pt}3m\,51.4s\\[-.2pt]{\fontsize{5.8}{6.4}\selectfont \textcolor{BrickRed}{$ \blacktriangledown\,2.7\times $}}\\[-.2pt]3m\,1.9s\\[-.2pt]{\fontsize{5.8}{6.4}\selectfont \textcolor{BrickRed}{$ \blacktriangledown\,2.2\times $}}\end{tabular} & \textcolor{gray}{\ensuremath{\cdot}} & \textcolor{gray}{\ensuremath{\cdot}} & \textcolor{gray}{\ensuremath{\cdot}} & \textcolor{gray}{\ensuremath{\cdot}} & \textcolor{gray}{\ensuremath{\cdot}}\\[3pt]\hline
 & \begin{tabular}[t]{@{}l@{}}\rule{0pt}{12pt}MPS / SIMD-ref\\SIMD-DP / Block\end{tabular} & \textcolor{gray}{\ensuremath{\cdot}} & \textcolor{gray}{\ensuremath{\cdot}} & \textcolor{gray}{\ensuremath{\cdot}} & \hspace*{1pt}\begin{tabular}[t]{@{}c@{}}\rule{0pt}{12pt}34.8s\\[-.2pt]{\fontsize{5.8}{6.4}\selectfont \textcolor{ForestGreen}{$ \blacktriangle\,12.8\times $}}\\[-.2pt]8.2s\\[-.2pt]{\fontsize{5.8}{6.4}\selectfont \textcolor{ForestGreen}{$ \blacktriangle\,50.2\times $}}\end{tabular} & \textcolor{gray}{\ensuremath{\cdot}} & \textcolor{gray}{\ensuremath{\cdot}} & \textcolor{gray}{\ensuremath{\cdot}} & \textcolor{gray}{\ensuremath{\cdot}} & \textcolor{gray}{\ensuremath{\cdot}} & \textcolor{gray}{\ensuremath{\cdot}} & \textcolor{gray}{\ensuremath{\cdot}} & \textcolor{gray}{\ensuremath{\cdot}} & \textcolor{gray}{\ensuremath{\cdot}} & \textcolor{gray}{\ensuremath{\cdot}} & \textcolor{gray}{\ensuremath{\cdot}} & \textcolor{gray}{\ensuremath{\cdot}}\\[3pt]\hline
 & \begin{tabular}[t]{@{}l@{}}\rule{0pt}{12pt}MPS / SIMD-ref\\SIMD-DP / KSW2\end{tabular} & \textcolor{gray}{\ensuremath{\cdot}} & \textcolor{gray}{\ensuremath{\cdot}} & \hspace*{1pt}\begin{tabular}[t]{@{}c@{}}\rule{0pt}{12pt}1m\,6.7s\\[-.2pt]{\fontsize{5.8}{6.4}\selectfont \textcolor{ForestGreen}{$ \blacktriangle\,8\% $}}\\[-.2pt]14.4s\\[-.2pt]{\fontsize{5.8}{6.4}\selectfont \textcolor{ForestGreen}{$ \blacktriangle\,4.7\times $}}\end{tabular} & \textcolor{gray}{\ensuremath{\cdot}} & \textcolor{gray}{\ensuremath{\cdot}} & \textcolor{gray}{\ensuremath{\cdot}} & \textcolor{gray}{\ensuremath{\cdot}} & \textcolor{gray}{\ensuremath{\cdot}} & \textcolor{gray}{\ensuremath{\cdot}} & \textcolor{gray}{\ensuremath{\cdot}} & \textcolor{gray}{\ensuremath{\cdot}} & \textcolor{gray}{\ensuremath{\cdot}} & \textcolor{gray}{\ensuremath{\cdot}} & \textcolor{gray}{\ensuremath{\cdot}} & \textcolor{gray}{\ensuremath{\cdot}} & \textcolor{gray}{\ensuremath{\cdot}}\\[3pt]\hline
\end{tabular}\par}
\smallskip{\footnotesize\noindent Each cell: A, comparison, B, comparison. Times cover 1,024 queries. $\cdot$ = no approved observation. ---\textsuperscript{g} = GPU unavailable, not executed. n/a = no measured comparison.\par
\noindent \textsuperscript{f}SIMD-ref request observed as scalar, compared only as that effective path. \textsuperscript{p}Accepted by explicit preservation after wrapper postflight failure, no run receipt claimed.\par
\noindent MPF/MPS = minimizers per fragment/sequence. Row order: index representation / minimizer, then chaining / base alignment. SIMD-ref is the reference minimizer, Bounded is bounded ntHash.\par
\noindent Blue = one lowest A and one lowest B down each machine/workload column, across continuation pages. Exact ties select the first row. Repeated baselines are marked only on page 1.\par
\noindent Light gray = minimap2 baseline. Pale yellow = preferred MPF / Bounded / SIMD-DP / Block configuration, not a claim of winning every column.\par
\noindent When all Vulkan modes are unmeasured, they share one row. Baseline comparisons use the fastest measured BioseqDB effective path separately for A and B.\par}
\clearpage
\noindent{\large\bfseries AWS: 2 Gbp, k=20, w=15 (mapping interval)\hfill 6/6}\par\smallskip
{\fontsize{7.1}{8.1}\selectfont\setlength{\tabcolsep}{1pt}\setlength{\fboxsep}{2pt}\renewcommand{\arraystretch}{1.12}
\noindent\begin{tabular}{@{}p{14mm}p{27mm}|p{13.55mm}|p{13.55mm}|p{13.55mm}|p{13.55mm}|p{13.55mm}|p{13.55mm}|p{13.55mm}|p{13.55mm}|p{13.55mm}|p{13.55mm}|p{13.55mm}|p{13.55mm}|p{13.55mm}|p{13.55mm}|p{13.55mm}|p{13.55mm}|@{}}\hline
Engine & Execution & \multicolumn{8}{c|}{\textbf{Index}} & \multicolumn{8}{c|}{\textbf{No index}}\\
 &  & \multicolumn{8}{c|}{\textbf{2 Gbp}} & \multicolumn{8}{c|}{\textbf{2 Gbp}}\\
 &  & \multicolumn{8}{c|}{\textbf{k=20, w=15}} & \multicolumn{8}{c|}{\textbf{k=20, w=15}}\\
 &  & \multicolumn{4}{c|}{\textbf{DNA4}} & \multicolumn{4}{c|}{\textbf{DNA5}} & \multicolumn{4}{c|}{\textbf{DNA4}} & \multicolumn{4}{c|}{\textbf{DNA5}}\\
 &  & \multicolumn{1}{c|}{\textbf{M7A}} & \multicolumn{1}{c|}{\textbf{M7I}} & \multicolumn{1}{c|}{\textbf{M4}} & \multicolumn{1}{c|}{\textbf{C7G}} & \multicolumn{1}{c|}{\textbf{M7A}} & \multicolumn{1}{c|}{\textbf{M7I}} & \multicolumn{1}{c|}{\textbf{M4}} & \multicolumn{1}{c|}{\textbf{C7G}} & \multicolumn{1}{c|}{\textbf{M7A}} & \multicolumn{1}{c|}{\textbf{M7I}} & \multicolumn{1}{c|}{\textbf{M4}} & \multicolumn{1}{c|}{\textbf{C7G}} & \multicolumn{1}{c|}{\textbf{M7A}} & \multicolumn{1}{c|}{\textbf{M7I}} & \multicolumn{1}{c|}{\textbf{M4}} & \multicolumn{1}{c|}{\textbf{C7G}}\\
\hline
\rowcolor{BenchBaseline}[0.5pt][0.5pt]\textbf{minimap2} & \textbf{Baseline} & \hspace*{1pt}\begin{tabular}[t]{@{}c@{}}\rule{0pt}{12pt}19.6s\\[-.2pt]{\fontsize{5.8}{6.4}\selectfont \textcolor{BrickRed}{$ \blacktriangledown\,6.7\times $}}\\[-.2pt]18.6s\\[-.2pt]{\fontsize{5.8}{6.4}\selectfont \textcolor{BrickRed}{$ \blacktriangledown\,8.4\times $}}\end{tabular} & \hspace*{1pt}\begin{tabular}[t]{@{}c@{}}\rule{0pt}{12pt}19.4s\\[-.2pt]{\fontsize{5.8}{6.4}\selectfont \textcolor{BrickRed}{$ \blacktriangledown\,7.2\times $}}\\[-.2pt]17.5s\\[-.2pt]{\fontsize{5.8}{6.4}\selectfont \textcolor{BrickRed}{$ \blacktriangledown\,9.1\times $}}\end{tabular} & \hspace*{1pt}\begin{tabular}[t]{@{}c@{}}\rule{0pt}{12pt}1m\,12.7s\\[-.2pt]{\fontsize{5.8}{6.4}\selectfont \textcolor{BrickRed}{$ \blacktriangledown\,13.6\times $}}\\[-.2pt]1m\,8.1s\\[-.2pt]{\fontsize{5.8}{6.4}\selectfont \textcolor{BrickRed}{$ \blacktriangledown\,20.2\times $}}\end{tabular} & \hspace*{1pt}\begin{tabular}[t]{@{}c@{}}\rule{0pt}{12pt}7m\,25.4s\\[-.2pt]{\fontsize{5.8}{6.4}\selectfont \textcolor{BrickRed}{$ \blacktriangledown\,106.2\times $}}\\[-.2pt]6m\,50.8s\\[-.2pt]{\fontsize{5.8}{6.4}\selectfont \textcolor{BrickRed}{$ \blacktriangledown\,212.9\times $}}\end{tabular} & \hspace*{1pt}\begin{tabular}[t]{@{}c@{}}\rule{0pt}{12pt}9.3s\\[-.2pt]{\fontsize{5.8}{6.4}\selectfont \textcolor{BrickRed}{$ \blacktriangledown\,24\% $}}\\[-.2pt]8.4s\\[-.2pt]{\fontsize{5.8}{6.4}\selectfont \textcolor{BrickRed}{$ \blacktriangledown\,1.6\times $}}\end{tabular} & \hspace*{1pt}\begin{tabular}[t]{@{}c@{}}\rule{0pt}{12pt}8.7s\\[-.2pt]{\fontsize{5.8}{6.4}\selectfont \textcolor{BrickRed}{$ \blacktriangledown\,1.5\times $}}\\[-.2pt]7s\\[-.2pt]{\fontsize{5.8}{6.4}\selectfont \textcolor{BrickRed}{$ \blacktriangledown\,1.7\times $}}\end{tabular} & \hspace*{1pt}\begin{tabular}[t]{@{}c@{}}\rule{0pt}{12pt}15.8s\\[-.2pt]{\fontsize{5.8}{6.4}\selectfont \textcolor{BrickRed}{$ \blacktriangledown\,26\% $}}\\[-.2pt]12.7s\\[-.2pt]{\fontsize{5.8}{6.4}\selectfont \textcolor{BrickRed}{$ \blacktriangledown\,1.6\times $}}\end{tabular} & \hspace*{1pt}\begin{tabular}[t]{@{}c@{}}\rule{0pt}{12pt}18.5s\\[-.2pt]{\fontsize{5.8}{6.4}\selectfont \textcolor{BrickRed}{$ \blacktriangledown\,1.6\times $}}\\[-.2pt]17s\\[-.2pt]{\fontsize{5.8}{6.4}\selectfont \textcolor{BrickRed}{$ \blacktriangledown\,1.8\times $}}\end{tabular} & \hspace*{1pt}\begin{tabular}[t]{@{}c@{}}\rule{0pt}{12pt}42.4s\\[-.2pt]{\fontsize{5.8}{6.4}\selectfont \textcolor{ForestGreen}{$ \blacktriangle\,2.0\times $}}\\[-.2pt]41.4s\\[-.2pt]{\fontsize{5.8}{6.4}\selectfont \textcolor{ForestGreen}{$ \blacktriangle\,1.9\times $}}\end{tabular} & \hspace*{1pt}\begin{tabular}[t]{@{}c@{}}\rule{0pt}{12pt}42.9s\\[-.2pt]{\fontsize{5.8}{6.4}\selectfont \textcolor{ForestGreen}{$ \blacktriangle\,2.0\times $}}\\[-.2pt]41.1s\\[-.2pt]{\fontsize{5.8}{6.4}\selectfont \textcolor{ForestGreen}{$ \blacktriangle\,1.8\times $}}\end{tabular} & \hspace*{1pt}\begin{tabular}[t]{@{}c@{}}\rule{0pt}{12pt}1m\,26.9s\\[-.2pt]{\fontsize{5.8}{6.4}\selectfont \textcolor{ForestGreen}{$ \blacktriangle\,1.9\times $}}\\[-.2pt]1m\,23.7s\\[-.2pt]{\fontsize{5.8}{6.4}\selectfont \textcolor{ForestGreen}{$ \blacktriangle\,1.7\times $}}\end{tabular} & \hspace*{1pt}\begin{tabular}[t]{@{}c@{}}\rule{0pt}{12pt}4m\,51.6s\\[-.2pt]{\fontsize{5.8}{6.4}\selectfont \textcolor{gray}{n/a}}\\[-.2pt]4m\,46.6s\\[-.2pt]{\fontsize{5.8}{6.4}\selectfont \textcolor{gray}{n/a}}\end{tabular} & \hspace*{1pt}\begin{tabular}[t]{@{}c@{}}\rule{0pt}{12pt}35s\\[-.2pt]{\fontsize{5.8}{6.4}\selectfont \textcolor{ForestGreen}{$ \blacktriangle\,8.6\times $}}\\[-.2pt]34s\\[-.2pt]{\fontsize{5.8}{6.4}\selectfont \textcolor{ForestGreen}{$ \blacktriangle\,2.5\times $}}\end{tabular} & \hspace*{1pt}\begin{tabular}[t]{@{}c@{}}\rule{0pt}{12pt}35.9s\\[-.2pt]{\fontsize{5.8}{6.4}\selectfont \textcolor{ForestGreen}{$ \blacktriangle\,6.7\times $}}\\[-.2pt]34.2s\\[-.2pt]{\fontsize{5.8}{6.4}\selectfont \textcolor{ForestGreen}{$ \blacktriangle\,2.5\times $}}\end{tabular} & \hspace*{1pt}\begin{tabular}[t]{@{}c@{}}\rule{0pt}{12pt}1m\,14.4s\\[-.2pt]{\fontsize{5.8}{6.4}\selectfont \textcolor{ForestGreen}{$ \blacktriangle\,6.9\times $}}\\[-.2pt]1m\,11.3s\\[-.2pt]{\fontsize{5.8}{6.4}\selectfont \textcolor{ForestGreen}{$ \blacktriangle\,2.2\times $}}\end{tabular} & \hspace*{1pt}\begin{tabular}[t]{@{}c@{}}\rule{0pt}{12pt}1m\,0.1s\\[-.2pt]{\fontsize{5.8}{6.4}\selectfont \textcolor{gray}{n/a}}\\[-.2pt]58.6s\\[-.2pt]{\fontsize{5.8}{6.4}\selectfont \textcolor{gray}{n/a}}\end{tabular}\\[3pt]\hline
\textbf{BioseqDB} & \begin{tabular}[t]{@{}l@{}}\rule{0pt}{12pt}MPS / Two-stack\\RMQ / Block\end{tabular} & \hspace*{1pt}\begin{tabular}[t]{@{}c@{}}\rule{0pt}{12pt}40.6s\\[-.2pt]{\fontsize{5.8}{6.4}\selectfont \textcolor{BrickRed}{$ \blacktriangledown\,2.1\times $}}\\[-.2pt]6.8s\\[-.2pt]{\fontsize{5.8}{6.4}\selectfont \textcolor{ForestGreen}{$ \blacktriangle\,2.7\times $}}\end{tabular} & \hspace*{1pt}\begin{tabular}[t]{@{}c@{}}\rule{0pt}{12pt}35.5s\\[-.2pt]{\fontsize{5.8}{6.4}\selectfont \textcolor{BrickRed}{$ \blacktriangledown\,1.8\times $}}\\[-.2pt]7s\\[-.2pt]{\fontsize{5.8}{6.4}\selectfont \textcolor{ForestGreen}{$ \blacktriangle\,2.5\times $}}\end{tabular} & \textcolor{gray}{\ensuremath{\cdot}} & \textcolor{gray}{\ensuremath{\cdot}} & \textcolor{gray}{\ensuremath{\cdot}} & \textcolor{gray}{\ensuremath{\cdot}} & \textcolor{gray}{\ensuremath{\cdot}} & \textcolor{gray}{\ensuremath{\cdot}} & \textcolor{gray}{\ensuremath{\cdot}} & \textcolor{gray}{\ensuremath{\cdot}} & \textcolor{gray}{\ensuremath{\cdot}} & \textcolor{gray}{\ensuremath{\cdot}} & \textcolor{gray}{\ensuremath{\cdot}} & \textcolor{gray}{\ensuremath{\cdot}} & \textcolor{gray}{\ensuremath{\cdot}} & \textcolor{gray}{\ensuremath{\cdot}}\\[3pt]\hline
 & \begin{tabular}[t]{@{}l@{}}\rule{0pt}{12pt}MPS / Two-stack\\RMQ / KSW2\end{tabular} & \hspace*{1pt}\begin{tabular}[t]{@{}c@{}}\rule{0pt}{12pt}43.9s\\[-.2pt]{\fontsize{5.8}{6.4}\selectfont \textcolor{BrickRed}{$ \blacktriangledown\,2.2\times $}}\\[-.2pt]6.8s\\[-.2pt]{\fontsize{5.8}{6.4}\selectfont \textcolor{ForestGreen}{$ \blacktriangle\,2.7\times $}}\end{tabular} & \hspace*{1pt}\begin{tabular}[t]{@{}c@{}}\rule{0pt}{12pt}39.6s\\[-.2pt]{\fontsize{5.8}{6.4}\selectfont \textcolor{BrickRed}{$ \blacktriangledown\,2.0\times $}}\\[-.2pt]6.7s\\[-.2pt]{\fontsize{5.8}{6.4}\selectfont \textcolor{ForestGreen}{$ \blacktriangle\,2.6\times $}}\end{tabular} & \textcolor{gray}{\ensuremath{\cdot}} & \textcolor{gray}{\ensuremath{\cdot}} & \textcolor{gray}{\ensuremath{\cdot}} & \textcolor{gray}{\ensuremath{\cdot}} & \textcolor{gray}{\ensuremath{\cdot}} & \textcolor{gray}{\ensuremath{\cdot}} & \textcolor{gray}{\ensuremath{\cdot}} & \textcolor{gray}{\ensuremath{\cdot}} & \textcolor{gray}{\ensuremath{\cdot}} & \textcolor{gray}{\ensuremath{\cdot}} & \textcolor{gray}{\ensuremath{\cdot}} & \textcolor{gray}{\ensuremath{\cdot}} & \textcolor{gray}{\ensuremath{\cdot}} & \textcolor{gray}{\ensuremath{\cdot}}\\[3pt]\hline
 & \begin{tabular}[t]{@{}l@{}}\rule{0pt}{12pt}MPS / Two-stack\\Scalar-DP / Block\end{tabular} & \hspace*{1pt}\begin{tabular}[t]{@{}c@{}}\rule{0pt}{12pt}39.8s\\[-.2pt]{\fontsize{5.8}{6.4}\selectfont \textcolor{BrickRed}{$ \blacktriangledown\,2.0\times $}}\\[-.2pt]6.8s\\[-.2pt]{\fontsize{5.8}{6.4}\selectfont \textcolor{ForestGreen}{$ \blacktriangle\,2.7\times $}}\end{tabular} & \hspace*{1pt}\begin{tabular}[t]{@{}c@{}}\rule{0pt}{12pt}35.8s\\[-.2pt]{\fontsize{5.8}{6.4}\selectfont \textcolor{BrickRed}{$ \blacktriangledown\,1.8\times $}}\\[-.2pt]6.5s\\[-.2pt]{\fontsize{5.8}{6.4}\selectfont \textcolor{ForestGreen}{$ \blacktriangle\,2.7\times $}}\end{tabular} & \textcolor{gray}{\ensuremath{\cdot}} & \textcolor{gray}{\ensuremath{\cdot}} & \textcolor{gray}{\ensuremath{\cdot}} & \textcolor{gray}{\ensuremath{\cdot}} & \textcolor{gray}{\ensuremath{\cdot}} & \textcolor{gray}{\ensuremath{\cdot}} & \textcolor{gray}{\ensuremath{\cdot}} & \textcolor{gray}{\ensuremath{\cdot}} & \textcolor{gray}{\ensuremath{\cdot}} & \textcolor{gray}{\ensuremath{\cdot}} & \textcolor{gray}{\ensuremath{\cdot}} & \textcolor{gray}{\ensuremath{\cdot}} & \textcolor{gray}{\ensuremath{\cdot}} & \textcolor{gray}{\ensuremath{\cdot}}\\[3pt]\hline
 & \begin{tabular}[t]{@{}l@{}}\rule{0pt}{12pt}MPS / Two-stack\\Scalar-DP / KSW2\end{tabular} & \hspace*{1pt}\begin{tabular}[t]{@{}c@{}}\rule{0pt}{12pt}43.6s\\[-.2pt]{\fontsize{5.8}{6.4}\selectfont \textcolor{BrickRed}{$ \blacktriangledown\,2.2\times $}}\\[-.2pt]6.8s\\[-.2pt]{\fontsize{5.8}{6.4}\selectfont \textcolor{ForestGreen}{$ \blacktriangle\,2.7\times $}}\end{tabular} & \hspace*{1pt}\begin{tabular}[t]{@{}c@{}}\rule{0pt}{12pt}39.5s\\[-.2pt]{\fontsize{5.8}{6.4}\selectfont \textcolor{BrickRed}{$ \blacktriangledown\,2.0\times $}}\\[-.2pt]7.2s\\[-.2pt]{\fontsize{5.8}{6.4}\selectfont \textcolor{ForestGreen}{$ \blacktriangle\,2.4\times $}}\end{tabular} & \textcolor{gray}{\ensuremath{\cdot}} & \textcolor{gray}{\ensuremath{\cdot}} & \textcolor{gray}{\ensuremath{\cdot}} & \textcolor{gray}{\ensuremath{\cdot}} & \textcolor{gray}{\ensuremath{\cdot}} & \textcolor{gray}{\ensuremath{\cdot}} & \textcolor{gray}{\ensuremath{\cdot}} & \textcolor{gray}{\ensuremath{\cdot}} & \textcolor{gray}{\ensuremath{\cdot}} & \textcolor{gray}{\ensuremath{\cdot}} & \textcolor{gray}{\ensuremath{\cdot}} & \textcolor{gray}{\ensuremath{\cdot}} & \textcolor{gray}{\ensuremath{\cdot}} & \textcolor{gray}{\ensuremath{\cdot}}\\[3pt]\hline
 & \begin{tabular}[t]{@{}l@{}}\rule{0pt}{12pt}MPS / Two-stack\\SIMD-DP / Block\end{tabular} & \hspace*{1pt}\begin{tabular}[t]{@{}c@{}}\rule{0pt}{12pt}39.5s\\[-.2pt]{\fontsize{5.8}{6.4}\selectfont \textcolor{BrickRed}{$ \blacktriangledown\,2.0\times $}}\\[-.2pt]6.8s\\[-.2pt]{\fontsize{5.8}{6.4}\selectfont \textcolor{ForestGreen}{$ \blacktriangle\,2.7\times $}}\end{tabular} & \hspace*{1pt}\begin{tabular}[t]{@{}c@{}}\rule{0pt}{12pt}35.7s\\[-.2pt]{\fontsize{5.8}{6.4}\selectfont \textcolor{BrickRed}{$ \blacktriangledown\,1.8\times $}}\\[-.2pt]7.2s\\[-.2pt]{\fontsize{5.8}{6.4}\selectfont \textcolor{ForestGreen}{$ \blacktriangle\,2.4\times $}}\end{tabular} & \textcolor{gray}{\ensuremath{\cdot}} & \hspace*{1pt}\begin{tabular}[t]{@{}c@{}}\rule{0pt}{12pt}33.9s\\[-.2pt]{\fontsize{5.8}{6.4}\selectfont \textcolor{ForestGreen}{$ \blacktriangle\,13.1\times $}}\\[-.2pt]6.9s\\[-.2pt]{\fontsize{5.8}{6.4}\selectfont \textcolor{ForestGreen}{$ \blacktriangle\,59.6\times $}}\end{tabular} & \textcolor{gray}{\ensuremath{\cdot}} & \textcolor{gray}{\ensuremath{\cdot}} & \textcolor{gray}{\ensuremath{\cdot}} & \textcolor{gray}{\ensuremath{\cdot}} & \textcolor{gray}{\ensuremath{\cdot}} & \textcolor{gray}{\ensuremath{\cdot}} & \textcolor{gray}{\ensuremath{\cdot}} & \textcolor{gray}{\ensuremath{\cdot}} & \textcolor{gray}{\ensuremath{\cdot}} & \textcolor{gray}{\ensuremath{\cdot}} & \textcolor{gray}{\ensuremath{\cdot}} & \textcolor{gray}{\ensuremath{\cdot}}\\[3pt]\hline
 & \begin{tabular}[t]{@{}l@{}}\rule{0pt}{12pt}MPS / Two-stack\\SIMD-DP / KSW2\end{tabular} & \hspace*{1pt}\begin{tabular}[t]{@{}c@{}}\rule{0pt}{12pt}44s\\[-.2pt]{\fontsize{5.8}{6.4}\selectfont \textcolor{BrickRed}{$ \blacktriangledown\,2.2\times $}}\\[-.2pt]6.8s\\[-.2pt]{\fontsize{5.8}{6.4}\selectfont \textcolor{ForestGreen}{$ \blacktriangle\,2.8\times $}}\end{tabular} & \hspace*{1pt}\begin{tabular}[t]{@{}c@{}}\rule{0pt}{12pt}39.6s\\[-.2pt]{\fontsize{5.8}{6.4}\selectfont \textcolor{BrickRed}{$ \blacktriangledown\,2.0\times $}}\\[-.2pt]7.2s\\[-.2pt]{\fontsize{5.8}{6.4}\selectfont \textcolor{ForestGreen}{$ \blacktriangle\,2.4\times $}}\end{tabular} & \textcolor{gray}{\ensuremath{\cdot}} & \textcolor{gray}{\ensuremath{\cdot}} & \textcolor{gray}{\ensuremath{\cdot}} & \textcolor{gray}{\ensuremath{\cdot}} & \textcolor{gray}{\ensuremath{\cdot}} & \textcolor{gray}{\ensuremath{\cdot}} & \textcolor{gray}{\ensuremath{\cdot}} & \textcolor{gray}{\ensuremath{\cdot}} & \textcolor{gray}{\ensuremath{\cdot}} & \textcolor{gray}{\ensuremath{\cdot}} & \textcolor{gray}{\ensuremath{\cdot}} & \textcolor{gray}{\ensuremath{\cdot}} & \textcolor{gray}{\ensuremath{\cdot}} & \textcolor{gray}{\ensuremath{\cdot}}\\[3pt]\hline
\rule{0pt}{12pt}\textbf{BioseqDB} & Vulkan (all modes) & \textcolor{gray}{---\textsuperscript{g}} & \textcolor{gray}{---\textsuperscript{g}} & \textcolor{gray}{---\textsuperscript{g}} & \textcolor{gray}{---\textsuperscript{g}} & \textcolor{gray}{---\textsuperscript{g}} & \textcolor{gray}{---\textsuperscript{g}} & \textcolor{gray}{---\textsuperscript{g}} & \textcolor{gray}{---\textsuperscript{g}} & \textcolor{gray}{---\textsuperscript{g}} & \textcolor{gray}{---\textsuperscript{g}} & \textcolor{gray}{---\textsuperscript{g}} & \textcolor{gray}{---\textsuperscript{g}} & \textcolor{gray}{---\textsuperscript{g}} & \textcolor{gray}{---\textsuperscript{g}} & \textcolor{gray}{---\textsuperscript{g}} & \textcolor{gray}{---\textsuperscript{g}}\\[3pt]\hline
\end{tabular}\par}
\smallskip{\footnotesize\noindent Each cell: A, comparison, B, comparison. Times cover 1,024 queries. $\cdot$ = no approved observation. ---\textsuperscript{g} = GPU unavailable, not executed. n/a = no measured comparison.\par
\noindent \textsuperscript{f}SIMD-ref request observed as scalar, compared only as that effective path. \textsuperscript{p}Accepted by explicit preservation after wrapper postflight failure, no run receipt claimed.\par
\noindent MPF/MPS = minimizers per fragment/sequence. Row order: index representation / minimizer, then chaining / base alignment. SIMD-ref is the reference minimizer, Bounded is bounded ntHash.\par
\noindent Blue = one lowest A and one lowest B down each machine/workload column, across continuation pages. Exact ties select the first row. Repeated baselines are marked only on page 1.\par
\noindent Light gray = minimap2 baseline. Pale yellow = preferred MPF / Bounded / SIMD-DP / Block configuration, not a claim of winning every column.\par
\noindent When all Vulkan modes are unmeasured, they share one row. Baseline comparisons use the fastest measured BioseqDB effective path separately for A and B.\par}
\end{landscape}
\restoregeometry

\clearpage\subsection{Comparison with other aligners}\label{sec:external-aligner-results}
The final table places minimap2 and one fixed BioseqDB configuration alongside blastn, Winnowmap and LexicMap. Both minimap2 and BioseqDB timings come from the corresponding machine\textquotesingle s non-external remote summary, matched by machine class, target size, alphabet, $(k,w)$ and index mode. External-campaign minimap2 timings are excluded. BioseqDB always uses minimizers per fragment (MPF), bounded ntHash SIMD minimizers, SIMD-DP chaining and Block Aligner. Missing observations are marked with a dot. The same default / TEST3 / TEST4 / TEST7 / TEST2 / FAST priority applies to both rows. FAST entries carry a single-sample footnote. No other BioseqDB configuration or historical timing scope is substituted.\par
The table contains 302 distinct observations and 1,510 accepted repeats. Only approved results are included. Qualification runs, interrupted attempts and the small-query checks used before the full runs are excluded.\par
\paragraph{Measurement procedure.}
\begin{itemize}
\item Each external observation has five fresh-process repeats of 1,024 queries. A is the median wall time, including executable startup, index loading, mapping and writing output. Fixture generation and validation are outside the timer. The index mode prepares the native on-disk index beforehand. No-index mode creates it inside the measured interval, including tool-specific preparation such as counting frequent k-mers or partitioning the reference. Operating-system caches are not purged.
\item The external campaign checks one primary occurrence per query, its position and strand, and an edit script reconstructed from the gapped output. These files include the five raw timing samples and returned-count vectors. This is additional evidence beyond the median-only summaries used for the earlier tables. MAPQ is checked only where the tool provides it. No MAPQ value is invented for blastn or LexicMap.
\item B is computed separately in each repeat by subtracting the union of the tool\textquotesingle s instrumented alignment and available MAPQ intervals, then taking the median. BLAST subtracts traceback-search regions. LexicMap subtracts wavefront-alignment calls. Winnowmap subtracts its native alignment and MAPQ intervals. These are different boundaries, so B is a tool-specific remaining time, not the duration of a common seeding-and-chaining stage.
\end{itemize}
\paragraph{How the tools were configured.}
\begin{itemize}
\item \textbf{blastn 2.17.0+.} The word size equals the requested $k$. There is no minimizer-window setting, so $w$ labels the fixture only. Low-complexity masking is disabled, and each reported alignment must cover at least 80\% of the query to remove short chance matches. Match reward 2, mismatch penalty 4, gap-open penalty 4 and gap-extension penalty 2 reproduce the baseline\textquotesingle s first affine-gap branch. BLAST has no second branch. The default E-value threshold of 10 is retained. DNA5 cells are marked because BLAST scores N differently.
\item \textbf{Winnowmap.} Both $k$ and $w$ are retained. Frequent k-mers are counted with Meryl to build weighted minimizers. This preparation is included in no-index time. Baseline scoring and first-stage bandwidth are retained, while Winnowmap\textquotesingle s additional structural-variant-aware stage keeps its native settings. Secondary alignments are suppressed. The benchmark build restores persisted-index output and reconstructs the weighting filter on loading. It also initializes a repeat-length variable to prevent undefined MAPQ/output behaviour. Memory-limited and multipart runs carry cell-specific notes.
\item \textbf{LexicMap.} The requested $k$ is retained, but LexicHash masks replace minimizer windows, so $w$ is inapplicable. The native 20,000 masks, gap filling between sparsely seeded regions and low-complexity rejection remain enabled. Query coverage is at least 80\%. Its native wavefront-alignment costs are match 0, mismatch 4 and gap open/extension 6/2. Its separate reported-score calculation uses match 2, mismatch -3 and gaps 5/2. These settings are not rewritten to mimic minimap2. References are divided into records of at most 268,435,455 bases with overlap equal to the longest query, and output coordinates are restored to the logical reference. DNA5 cells are marked because its native reference encoding replaces N with A.
\end{itemize}
\paragraph{What the comparison establishes.}
The external runs used new instances of the same EC2 types, not the physical instances used for the retained minimap2 and BioseqDB results. External volumes are 300 GiB unencrypted gp3, compared with 500 GiB encrypted gp3 in the earlier campaign. Both provision 6,000 IOPS and 500 MiB/s. The external host bindings record four benchmark threads, nice $-20$ and Clang 18.1.3, whereas the earlier campaign used Clang 22.1.8. Compressed-RAM swap is recorded for many external runs. Its presence does not prove how much swapping occurred during each measurement.\par
The workload dimensions agree, but the timing boundaries, reference layouts, scoring and resource conditions are not identical. In particular, the retained framework timings exclude process-startup setup, and BioseqDB no-index mode constructs a transient index rather than the external tools\textquotesingle{} on-disk indexes. The summaries do not establish cross-campaign byte-for-byte fixture identity. Ratios therefore describe the published elapsed times under these qualified conditions, rather than an isolated causal effect of the aligner.\par
The table retains the green/red/gray triangles, 5\% tie band and $\times$ notation above 30\%. Every non-baseline row compares with minimap2 in the same machine-class/workload column. Minimap2 compares only with the displayed fixed BioseqDB configuration, not another faster BioseqDB row. Blue marks one lowest reported A and one lowest reported B down the five rows in each column. Exact ties select the first displayed row. Light gray identifies the minimap2 baseline. Pale yellow identifies the preferred MPF / Bounded / SIMD-DP / Block configuration, independently of the measured winners. For B, both the ratio and the highlight compare native residual times with different work removed, and do not establish superiority at an equivalent processing stage. Dots mean no approved timing, not failure or hardware unavailability.\par
\newgeometry{left=10mm,right=10mm,top=12mm,bottom=12mm}
\begin{landscape}
\refstepcounter{table}\label{tab:external-aligners}
\clearpage
\noindent{\large\bfseries Table \thetable: aligner comparison, 4 Gbp, k=15, w=10\hfill 1/4}\par\smallskip
\noindent{\footnotesize Different campaigns and native timing boundaries. A/B ratios are descriptive, subject to the measurement notes.}\par\smallskip
{\fontsize{7.1}{8.1}\selectfont\setlength{\tabcolsep}{1pt}\setlength{\fboxsep}{2pt}\renewcommand{\arraystretch}{1.12}
\noindent\begin{tabular}{@{}p{15mm}p{26mm}|p{13.55mm}|p{13.55mm}|p{13.55mm}|p{13.55mm}|p{13.55mm}|p{13.55mm}|p{13.55mm}|p{13.55mm}|p{13.55mm}|p{13.55mm}|p{13.55mm}|p{13.55mm}|p{13.55mm}|p{13.55mm}|p{13.55mm}|p{13.55mm}|@{}}\hline
Aligner & Configuration & \multicolumn{8}{c|}{\textbf{Index}} & \multicolumn{8}{c|}{\textbf{No index}}\\
 &  & \multicolumn{8}{c|}{\textbf{4 Gbp}} & \multicolumn{8}{c|}{\textbf{4 Gbp}}\\
 &  & \multicolumn{8}{c|}{\textbf{k=15, w=10}} & \multicolumn{8}{c|}{\textbf{k=15, w=10}}\\
 &  & \multicolumn{4}{c|}{\textbf{DNA4}} & \multicolumn{4}{c|}{\textbf{DNA5}} & \multicolumn{4}{c|}{\textbf{DNA4}} & \multicolumn{4}{c|}{\textbf{DNA5}}\\
 &  & \multicolumn{1}{c|}{\textbf{M7A}} & \multicolumn{1}{c|}{\textbf{M7I}} & \multicolumn{1}{c|}{\textbf{M4}} & \multicolumn{1}{c|}{\textbf{C7G}} & \multicolumn{1}{c|}{\textbf{M7A}} & \multicolumn{1}{c|}{\textbf{M7I}} & \multicolumn{1}{c|}{\textbf{M4}} & \multicolumn{1}{c|}{\textbf{C7G}} & \multicolumn{1}{c|}{\textbf{M7A}} & \multicolumn{1}{c|}{\textbf{M7I}} & \multicolumn{1}{c|}{\textbf{M4}} & \multicolumn{1}{c|}{\textbf{C7G}} & \multicolumn{1}{c|}{\textbf{M7A}} & \multicolumn{1}{c|}{\textbf{M7I}} & \multicolumn{1}{c|}{\textbf{M4}} & \multicolumn{1}{c|}{\textbf{C7G}}\\
\hline
\rowcolor{BenchBaseline}[0.5pt][0.5pt]minimap2 & Retained baseline & \hspace*{1pt}\begin{tabular}[t]{@{}c@{}}\rule{0pt}{12pt}36.2s\\[-.2pt]{\fontsize{5.8}{6.4}\selectfont \textcolor{BrickRed}{$ \blacktriangledown\,3.6\times $}}\\[-.2pt]35s\\[-.2pt]{\fontsize{5.8}{6.4}\selectfont \textcolor{BrickRed}{$ \blacktriangledown\,3.9\times $}}\end{tabular} & \hspace*{1pt}\begin{tabular}[t]{@{}c@{}}\rule{0pt}{12pt}38.1s\\[-.2pt]{\fontsize{5.8}{6.4}\selectfont \textcolor{BrickRed}{$ \blacktriangledown\,2.9\times $}}\\[-.2pt]36s\\[-.2pt]{\fontsize{5.8}{6.4}\selectfont \textcolor{BrickRed}{$ \blacktriangledown\,3.0\times $}}\end{tabular} & \hspace*{1pt}\begin{tabular}[t]{@{}c@{}}\rule{0pt}{12pt}1m\,33.5s\\[-.2pt]{\fontsize{5.8}{6.4}\selectfont \textcolor{BrickRed}{$ \blacktriangledown\,9.9\times $}}\\[-.2pt]1m\,29.7s\\[-.2pt]{\fontsize{5.8}{6.4}\selectfont \textcolor{BrickRed}{$ \blacktriangledown\,12.8\times $}}\end{tabular} & \hspace*{1pt}\begin{tabular}[t]{@{}c@{}}\rule{0pt}{12pt}6m\,31.5s\\[-.2pt]{\fontsize{5.8}{6.4}\selectfont \textcolor{gray}{n/a}\textsuperscript{g}}\\[-.2pt]6m\,26.3s\\[-.2pt]{\fontsize{5.8}{6.4}\selectfont \textcolor{gray}{n/a}}\end{tabular} & \hspace*{1pt}\begin{tabular}[t]{@{}c@{}}\rule{0pt}{12pt}31.6s\\[-.2pt]{\fontsize{5.8}{6.4}\selectfont \textcolor{gray}{n/a}}\\[-.2pt]30.3s\\[-.2pt]{\fontsize{5.8}{6.4}\selectfont \textcolor{gray}{n/a}}\end{tabular} & \hspace*{1pt}\begin{tabular}[t]{@{}c@{}}\rule{0pt}{12pt}\colorbox{BenchBest}{\strut 34.1s}\\[-.2pt]{\fontsize{5.8}{6.4}\selectfont \textcolor{gray}{n/a}}\\[-.2pt]\colorbox{BenchBest}{\strut 32s}\\[-.2pt]{\fontsize{5.8}{6.4}\selectfont \textcolor{gray}{n/a}}\end{tabular} & \hspace*{1pt}\begin{tabular}[t]{@{}c@{}}\rule{0pt}{12pt}1m\,9.5s\\[-.2pt]{\fontsize{5.8}{6.4}\selectfont \textcolor{gray}{n/a}}\\[-.2pt]1m\,6.1s\\[-.2pt]{\fontsize{5.8}{6.4}\selectfont \textcolor{gray}{n/a}}\end{tabular} & \hspace*{1pt}\begin{tabular}[t]{@{}c@{}}\rule{0pt}{12pt}3m\,1s\\[-.2pt]{\fontsize{5.8}{6.4}\selectfont \textcolor{gray}{n/a}\textsuperscript{g}}\\[-.2pt]2m\,51.5s\\[-.2pt]{\fontsize{5.8}{6.4}\selectfont \textcolor{gray}{n/a}}\end{tabular} & \hspace*{1pt}\begin{tabular}[t]{@{}c@{}}\rule{0pt}{12pt}\colorbox{BenchBest}{\strut 2m\,55.5s}\\[-.2pt]{\fontsize{5.8}{6.4}\selectfont \textcolor{ForestGreen}{$ \blacktriangle\,22\% $}}\\[-.2pt]\colorbox{BenchBest}{\strut 2m\,54s}\\[-.2pt]{\fontsize{5.8}{6.4}\selectfont \textcolor{ForestGreen}{$ \blacktriangle\,19\% $}}\end{tabular} & \hspace*{1pt}\begin{tabular}[t]{@{}c@{}}\rule{0pt}{12pt}\colorbox{BenchBest}{\strut 2m\,50.4s}\\[-.2pt]{\fontsize{5.8}{6.4}\selectfont \textcolor{ForestGreen}{$ \blacktriangle\,28\% $}}\\[-.2pt]\colorbox{BenchBest}{\strut 2m\,48.2s}\\[-.2pt]{\fontsize{5.8}{6.4}\selectfont \textcolor{ForestGreen}{$ \blacktriangle\,25\% $}}\end{tabular} & \hspace*{1pt}\begin{tabular}[t]{@{}c@{}}\rule{0pt}{12pt}\colorbox{BenchBest}{\strut 6m\,12.2s}\\[-.2pt]{\fontsize{5.8}{6.4}\selectfont \textcolor{ForestGreen}{$ \blacktriangle\,9\% $}}\\[-.2pt]\colorbox{BenchBest}{\strut 6m\,8.2s}\\[-.2pt]{\fontsize{5.8}{6.4}\selectfont \textcolor{ForestGreen}{$ \blacktriangle\,6\% $}}\end{tabular} & \hspace*{1pt}\begin{tabular}[t]{@{}c@{}}\rule{0pt}{12pt}44m\,20.5s\\[-.2pt]{\fontsize{5.8}{6.4}\selectfont \textcolor{gray}{n/a}\textsuperscript{g}}\\[-.2pt]44m\,15.6s\\[-.2pt]{\fontsize{5.8}{6.4}\selectfont \textcolor{gray}{n/a}}\end{tabular} & \hspace*{1pt}\begin{tabular}[t]{@{}c@{}}\rule{0pt}{12pt}\colorbox{BenchBest}{\strut 1m\,47.1s}\\[-.2pt]{\fontsize{5.8}{6.4}\selectfont \textcolor{ForestGreen}{$ \blacktriangle\,5.5\times $}}\\[-.2pt]\colorbox{BenchBest}{\strut 1m\,45.8s}\\[-.2pt]{\fontsize{5.8}{6.4}\selectfont \textcolor{ForestGreen}{$ \blacktriangle\,2.1\times $}}\end{tabular} & \hspace*{1pt}\begin{tabular}[t]{@{}c@{}}\rule{0pt}{12pt}\colorbox{BenchBest}{\strut 1m\,47.2s}\\[-.2pt]{\fontsize{5.8}{6.4}\selectfont \textcolor{ForestGreen}{$ \blacktriangle\,4.7\times $}}\\[-.2pt]\colorbox{BenchBest}{\strut 1m\,45.1s}\\[-.2pt]{\fontsize{5.8}{6.4}\selectfont \textcolor{ForestGreen}{$ \blacktriangle\,2.3\times $}}\end{tabular} & \hspace*{1pt}\begin{tabular}[t]{@{}c@{}}\rule{0pt}{12pt}\colorbox{BenchBest}{\strut 3m\,40.3s}\\[-.2pt]{\fontsize{5.8}{6.4}\selectfont \textcolor{ForestGreen}{$ \blacktriangle\,4.2\times $}}\\[-.2pt]\colorbox{BenchBest}{\strut 3m\,36.6s}\\[-.2pt]{\fontsize{5.8}{6.4}\selectfont \textcolor{ForestGreen}{$ \blacktriangle\,1.7\times $}}\end{tabular} & \hspace*{1pt}\begin{tabular}[t]{@{}c@{}}\rule{0pt}{12pt}29m\,43.4s\\[-.2pt]{\fontsize{5.8}{6.4}\selectfont \textcolor{gray}{n/a}\textsuperscript{g}}\\[-.2pt]29m\,38.5s\\[-.2pt]{\fontsize{5.8}{6.4}\selectfont \textcolor{gray}{n/a}}\end{tabular}\\[3pt]\hline
\rowcolor{BenchPreferred}[0.5pt][0.5pt]BioseqDB & \begin{tabular}[t]{@{}l@{}}\rule{0pt}{12pt}MPF / Bounded\\SIMD-DP / Block\end{tabular} & \hspace*{1pt}\begin{tabular}[t]{@{}c@{}}\rule{0pt}{12pt}\colorbox{BenchBest}{\strut 9.9s}\\[-.2pt]{\fontsize{5.8}{6.4}\selectfont \textcolor{ForestGreen}{$ \blacktriangle\,3.6\times $}}\\[-.2pt]\colorbox{BenchBest}{\strut 9s}\\[-.2pt]{\fontsize{5.8}{6.4}\selectfont \textcolor{ForestGreen}{$ \blacktriangle\,3.9\times $}}\end{tabular} & \hspace*{1pt}\begin{tabular}[t]{@{}c@{}}\rule{0pt}{12pt}\colorbox{BenchBest}{\strut 13s}\\[-.2pt]{\fontsize{5.8}{6.4}\selectfont \textcolor{ForestGreen}{$ \blacktriangle\,2.9\times $}}\\[-.2pt]\colorbox{BenchBest}{\strut 11.9s}\\[-.2pt]{\fontsize{5.8}{6.4}\selectfont \textcolor{ForestGreen}{$ \blacktriangle\,3.0\times $}}\end{tabular} & \hspace*{1pt}\begin{tabular}[t]{@{}c@{}}\rule{0pt}{12pt}\colorbox{BenchBest}{\strut 9.4s}\\[-.2pt]{\fontsize{5.8}{6.4}\selectfont \textcolor{ForestGreen}{$ \blacktriangle\,9.9\times $}}\\[-.2pt]\colorbox{BenchBest}{\strut 7s}\\[-.2pt]{\fontsize{5.8}{6.4}\selectfont \textcolor{ForestGreen}{$ \blacktriangle\,12.8\times $}}\end{tabular} & \textcolor{gray}{\ensuremath{\cdot}} & \textcolor{gray}{\ensuremath{\cdot}} & \textcolor{gray}{\ensuremath{\cdot}} & \textcolor{gray}{\ensuremath{\cdot}} & \textcolor{gray}{\ensuremath{\cdot}} & \hspace*{1pt}\begin{tabular}[t]{@{}c@{}}\rule{0pt}{12pt}3m\,44.5s\\[-.2pt]{\fontsize{5.8}{6.4}\selectfont \textcolor{BrickRed}{$ \blacktriangledown\,28\% $}}\\[-.2pt]3m\,35.9s\\[-.2pt]{\fontsize{5.8}{6.4}\selectfont \textcolor{BrickRed}{$ \blacktriangledown\,24\% $}}\end{tabular} & \hspace*{1pt}\begin{tabular}[t]{@{}c@{}}\rule{0pt}{12pt}3m\,57.4s\\[-.2pt]{\fontsize{5.8}{6.4}\selectfont \textcolor{BrickRed}{$ \blacktriangledown\,1.4\times $}}\\[-.2pt]3m\,45.7s\\[-.2pt]{\fontsize{5.8}{6.4}\selectfont \textcolor{BrickRed}{$ \blacktriangledown\,1.3\times $}}\end{tabular} & \hspace*{1pt}\begin{tabular}[t]{@{}c@{}}\rule{0pt}{12pt}6m\,50.6s\\[-.2pt]{\fontsize{5.8}{6.4}\selectfont \textcolor{BrickRed}{$ \blacktriangledown\,10\% $}}\\[-.2pt]6m\,30.7s\\[-.2pt]{\fontsize{5.8}{6.4}\selectfont \textcolor{BrickRed}{$ \blacktriangledown\,6\% $}}\end{tabular} & \textcolor{gray}{\ensuremath{\cdot}} & \hspace*{1pt}\begin{tabular}[t]{@{}c@{}}\rule{0pt}{12pt}9m\,45.5s\\[-.2pt]{\fontsize{5.8}{6.4}\selectfont \textcolor{BrickRed}{$ \blacktriangledown\,5.5\times $}}\\[-.2pt]3m\,41.5s\\[-.2pt]{\fontsize{5.8}{6.4}\selectfont \textcolor{BrickRed}{$ \blacktriangledown\,2.1\times $}}\end{tabular} & \hspace*{1pt}\begin{tabular}[t]{@{}c@{}}\rule{0pt}{12pt}8m\,25.6s\\[-.2pt]{\fontsize{5.8}{6.4}\selectfont \textcolor{BrickRed}{$ \blacktriangledown\,4.7\times $}}\\[-.2pt]3m\,56.7s\\[-.2pt]{\fontsize{5.8}{6.4}\selectfont \textcolor{BrickRed}{$ \blacktriangledown\,2.3\times $}}\end{tabular} & \hspace*{1pt}\begin{tabular}[t]{@{}c@{}}\rule{0pt}{12pt}15m\,24.9s\\[-.2pt]{\fontsize{5.8}{6.4}\selectfont \textcolor{BrickRed}{$ \blacktriangledown\,4.2\times $}}\\[-.2pt]6m\,0.2s\\[-.2pt]{\fontsize{5.8}{6.4}\selectfont \textcolor{BrickRed}{$ \blacktriangledown\,1.7\times $}}\end{tabular} & \textcolor{gray}{\ensuremath{\cdot}}\\[3pt]\hline
blastn & Native\textsuperscript{*} & \hspace*{1pt}\begin{tabular}[t]{@{}c@{}}\rule{0pt}{12pt}11m\,14.7s\\[-.2pt]{\fontsize{5.8}{6.4}\selectfont \textcolor{BrickRed}{$ \blacktriangledown\,18.6\times $}\textsuperscript{c}}\\[-.2pt]11m\,9.5s\\[-.2pt]{\fontsize{5.8}{6.4}\selectfont \textcolor{BrickRed}{$ \blacktriangledown\,19.1\times $}}\end{tabular} & \hspace*{1pt}\begin{tabular}[t]{@{}c@{}}\rule{0pt}{12pt}16m\,40.6s\\[-.2pt]{\fontsize{5.8}{6.4}\selectfont \textcolor{BrickRed}{$ \blacktriangledown\,26.3\times $}\textsuperscript{c}}\\[-.2pt]16m\,29.1s\\[-.2pt]{\fontsize{5.8}{6.4}\selectfont \textcolor{BrickRed}{$ \blacktriangledown\,27.5\times $}}\end{tabular} & \hspace*{1pt}\begin{tabular}[t]{@{}c@{}}\rule{0pt}{12pt}22m\,6.7s\\[-.2pt]{\fontsize{5.8}{6.4}\selectfont \textcolor{BrickRed}{$ \blacktriangledown\,14.2\times $}\textsuperscript{c}}\\[-.2pt]21m\,51.5s\\[-.2pt]{\fontsize{5.8}{6.4}\selectfont \textcolor{BrickRed}{$ \blacktriangledown\,14.6\times $}}\end{tabular} & \hspace*{1pt}\begin{tabular}[t]{@{}c@{}}\rule{0pt}{12pt}12m\,25.4s\\[-.2pt]{\fontsize{5.8}{6.4}\selectfont \textcolor{BrickRed}{$ \blacktriangledown\,1.9\times $}\textsuperscript{c}}\\[-.2pt]12m\,19.3s\\[-.2pt]{\fontsize{5.8}{6.4}\selectfont \textcolor{BrickRed}{$ \blacktriangledown\,1.9\times $}}\end{tabular} & \hspace*{1pt}\begin{tabular}[t]{@{}c@{}}\rule{0pt}{12pt}11m\,40.1s\\[-.2pt]{\fontsize{5.8}{6.4}\selectfont \textcolor{BrickRed}{$ \blacktriangledown\,22.2\times $}\textsuperscript{a,c}}\\[-.2pt]9m\,50.3s\\[-.2pt]{\fontsize{5.8}{6.4}\selectfont \textcolor{BrickRed}{$ \blacktriangledown\,19.5\times $}}\end{tabular} & \hspace*{1pt}\begin{tabular}[t]{@{}c@{}}\rule{0pt}{12pt}16m\,52.8s\\[-.2pt]{\fontsize{5.8}{6.4}\selectfont \textcolor{BrickRed}{$ \blacktriangledown\,29.7\times $}\textsuperscript{a,c}}\\[-.2pt]13m\,48.2s\\[-.2pt]{\fontsize{5.8}{6.4}\selectfont \textcolor{BrickRed}{$ \blacktriangledown\,25.9\times $}}\end{tabular} & \hspace*{1pt}\begin{tabular}[t]{@{}c@{}}\rule{0pt}{12pt}23m\,30.3s\\[-.2pt]{\fontsize{5.8}{6.4}\selectfont \textcolor{BrickRed}{$ \blacktriangledown\,20.3\times $}\textsuperscript{a,c}}\\[-.2pt]18m\,49.3s\\[-.2pt]{\fontsize{5.8}{6.4}\selectfont \textcolor{BrickRed}{$ \blacktriangledown\,17.1\times $}}\end{tabular} & \hspace*{1pt}\begin{tabular}[t]{@{}c@{}}\rule{0pt}{12pt}12m\,50.3s\\[-.2pt]{\fontsize{5.8}{6.4}\selectfont \textcolor{BrickRed}{$ \blacktriangledown\,4.3\times $}\textsuperscript{a,c}}\\[-.2pt]10m\,46.6s\\[-.2pt]{\fontsize{5.8}{6.4}\selectfont \textcolor{BrickRed}{$ \blacktriangledown\,3.8\times $}}\end{tabular} & \hspace*{1pt}\begin{tabular}[t]{@{}c@{}}\rule{0pt}{12pt}11m\,28s\\[-.2pt]{\fontsize{5.8}{6.4}\selectfont \textcolor{BrickRed}{$ \blacktriangledown\,3.9\times $}\textsuperscript{c}}\\[-.2pt]11m\,22.8s\\[-.2pt]{\fontsize{5.8}{6.4}\selectfont \textcolor{BrickRed}{$ \blacktriangledown\,3.9\times $}}\end{tabular} & \hspace*{1pt}\begin{tabular}[t]{@{}c@{}}\rule{0pt}{12pt}16m\,55s\\[-.2pt]{\fontsize{5.8}{6.4}\selectfont \textcolor{BrickRed}{$ \blacktriangledown\,6.0\times $}\textsuperscript{c}}\\[-.2pt]16m\,43.4s\\[-.2pt]{\fontsize{5.8}{6.4}\selectfont \textcolor{BrickRed}{$ \blacktriangledown\,6.0\times $}}\end{tabular} & \hspace*{1pt}\begin{tabular}[t]{@{}c@{}}\rule{0pt}{12pt}22m\,36.2s\\[-.2pt]{\fontsize{5.8}{6.4}\selectfont \textcolor{BrickRed}{$ \blacktriangledown\,3.6\times $}\textsuperscript{c}}\\[-.2pt]22m\,20.8s\\[-.2pt]{\fontsize{5.8}{6.4}\selectfont \textcolor{BrickRed}{$ \blacktriangledown\,3.6\times $}}\end{tabular} & \hspace*{1pt}\begin{tabular}[t]{@{}c@{}}\rule{0pt}{12pt}13m\,1.6s\\[-.2pt]{\fontsize{5.8}{6.4}\selectfont \textcolor{ForestGreen}{$ \blacktriangle\,3.4\times $}\textsuperscript{c}}\\[-.2pt]12m\,55.8s\\[-.2pt]{\fontsize{5.8}{6.4}\selectfont \textcolor{ForestGreen}{$ \blacktriangle\,3.4\times $}}\end{tabular} & \hspace*{1pt}\begin{tabular}[t]{@{}c@{}}\rule{0pt}{12pt}12m\,18.6s\\[-.2pt]{\fontsize{5.8}{6.4}\selectfont \textcolor{BrickRed}{$ \blacktriangledown\,6.9\times $}\textsuperscript{a,c}}\\[-.2pt]10m\,29.2s\\[-.2pt]{\fontsize{5.8}{6.4}\selectfont \textcolor{BrickRed}{$ \blacktriangledown\,5.9\times $}}\end{tabular} & \hspace*{1pt}\begin{tabular}[t]{@{}c@{}}\rule{0pt}{12pt}17m\,25.7s\\[-.2pt]{\fontsize{5.8}{6.4}\selectfont \textcolor{BrickRed}{$ \blacktriangledown\,9.8\times $}\textsuperscript{a,c}}\\[-.2pt]14m\,22.8s\\[-.2pt]{\fontsize{5.8}{6.4}\selectfont \textcolor{BrickRed}{$ \blacktriangledown\,8.2\times $}}\end{tabular} & \hspace*{1pt}\begin{tabular}[t]{@{}c@{}}\rule{0pt}{12pt}24m\,29.4s\\[-.2pt]{\fontsize{5.8}{6.4}\selectfont \textcolor{BrickRed}{$ \blacktriangledown\,6.7\times $}\textsuperscript{a,c}}\\[-.2pt]19m\,49.6s\\[-.2pt]{\fontsize{5.8}{6.4}\selectfont \textcolor{BrickRed}{$ \blacktriangledown\,5.5\times $}}\end{tabular} & \hspace*{1pt}\begin{tabular}[t]{@{}c@{}}\rule{0pt}{12pt}13m\,43.9s\\[-.2pt]{\fontsize{5.8}{6.4}\selectfont \textcolor{ForestGreen}{$ \blacktriangle\,2.2\times $}\textsuperscript{a,c}}\\[-.2pt]11m\,39.5s\\[-.2pt]{\fontsize{5.8}{6.4}\selectfont \textcolor{ForestGreen}{$ \blacktriangle\,2.5\times $}}\end{tabular}\\[3pt]\hline
winnowmap & Native\textsuperscript{*} & \hspace*{1pt}\begin{tabular}[t]{@{}c@{}}\rule{0pt}{12pt}52.9s\\[-.2pt]{\fontsize{5.8}{6.4}\selectfont \textcolor{BrickRed}{$ \blacktriangledown\,1.5\times $}\textsuperscript{c,e}}\\[-.2pt]50.6s\\[-.2pt]{\fontsize{5.8}{6.4}\selectfont \textcolor{BrickRed}{$ \blacktriangledown\,1.4\times $}}\end{tabular} & \hspace*{1pt}\begin{tabular}[t]{@{}c@{}}\rule{0pt}{12pt}59.4s\\[-.2pt]{\fontsize{5.8}{6.4}\selectfont \textcolor{BrickRed}{$ \blacktriangledown\,1.6\times $}\textsuperscript{c,e}}\\[-.2pt]55.2s\\[-.2pt]{\fontsize{5.8}{6.4}\selectfont \textcolor{BrickRed}{$ \blacktriangledown\,1.5\times $}}\end{tabular} & \hspace*{1pt}\begin{tabular}[t]{@{}c@{}}\rule{0pt}{12pt}1m\,49.1s\\[-.2pt]{\fontsize{5.8}{6.4}\selectfont \textcolor{BrickRed}{$ \blacktriangledown\,17\% $}\textsuperscript{c,e}}\\[-.2pt]1m\,41.3s\\[-.2pt]{\fontsize{5.8}{6.4}\selectfont \textcolor{BrickRed}{$ \blacktriangledown\,13\% $}}\end{tabular} & \hspace*{1pt}\begin{tabular}[t]{@{}c@{}}\rule{0pt}{12pt}\colorbox{BenchBest}{\strut 1m\,16.1s}\\[-.2pt]{\fontsize{5.8}{6.4}\selectfont \textcolor{ForestGreen}{$ \blacktriangle\,5.1\times $}\textsuperscript{c,d}}\\[-.2pt]\colorbox{BenchBest}{\strut 1m\,8.2s}\\[-.2pt]{\fontsize{5.8}{6.4}\selectfont \textcolor{ForestGreen}{$ \blacktriangle\,5.7\times $}}\end{tabular} & \hspace*{1pt}\begin{tabular}[t]{@{}c@{}}\rule{0pt}{12pt}\colorbox{BenchBest}{\strut 31.6s}\\[-.2pt]{\fontsize{5.8}{6.4}\selectfont \textcolor{gray}{$ \blacktriangleright\,0\% $}\textsuperscript{c,e}}\\[-.2pt]\colorbox{BenchBest}{\strut 29.4s}\\[-.2pt]{\fontsize{5.8}{6.4}\selectfont \textcolor{gray}{$ \blacktriangleright\,3\% $}}\end{tabular} & \hspace*{1pt}\begin{tabular}[t]{@{}c@{}}\rule{0pt}{12pt}44.2s\\[-.2pt]{\fontsize{5.8}{6.4}\selectfont \textcolor{BrickRed}{$ \blacktriangledown\,30\% $}\textsuperscript{c,e}}\\[-.2pt]40.3s\\[-.2pt]{\fontsize{5.8}{6.4}\selectfont \textcolor{BrickRed}{$ \blacktriangledown\,26\% $}}\end{tabular} & \hspace*{1pt}\begin{tabular}[t]{@{}c@{}}\rule{0pt}{12pt}\colorbox{BenchBest}{\strut 59s}\\[-.2pt]{\fontsize{5.8}{6.4}\selectfont \textcolor{ForestGreen}{$ \blacktriangle\,15\% $}\textsuperscript{c,e}}\\[-.2pt]\colorbox{BenchBest}{\strut 51.8s}\\[-.2pt]{\fontsize{5.8}{6.4}\selectfont \textcolor{ForestGreen}{$ \blacktriangle\,22\% $}}\end{tabular} & \hspace*{1pt}\begin{tabular}[t]{@{}c@{}}\rule{0pt}{12pt}\colorbox{BenchBest}{\strut 1m\,8s}\\[-.2pt]{\fontsize{5.8}{6.4}\selectfont \textcolor{ForestGreen}{$ \blacktriangle\,2.7\times $}\textsuperscript{c,d}}\\[-.2pt]\colorbox{BenchBest}{\strut 1m\,0.3s}\\[-.2pt]{\fontsize{5.8}{6.4}\selectfont \textcolor{ForestGreen}{$ \blacktriangle\,2.8\times $}}\end{tabular} & \hspace*{1pt}\begin{tabular}[t]{@{}c@{}}\rule{0pt}{12pt}8m\,6s\\[-.2pt]{\fontsize{5.8}{6.4}\selectfont \textcolor{BrickRed}{$ \blacktriangledown\,2.8\times $}\textsuperscript{c,e}}\\[-.2pt]8m\,3.7s\\[-.2pt]{\fontsize{5.8}{6.4}\selectfont \textcolor{BrickRed}{$ \blacktriangledown\,2.8\times $}}\end{tabular} & \hspace*{1pt}\begin{tabular}[t]{@{}c@{}}\rule{0pt}{12pt}7m\,59.6s\\[-.2pt]{\fontsize{5.8}{6.4}\selectfont \textcolor{BrickRed}{$ \blacktriangledown\,2.8\times $}\textsuperscript{c,e}}\\[-.2pt]7m\,55.3s\\[-.2pt]{\fontsize{5.8}{6.4}\selectfont \textcolor{BrickRed}{$ \blacktriangledown\,2.8\times $}}\end{tabular} & \hspace*{1pt}\begin{tabular}[t]{@{}c@{}}\rule{0pt}{12pt}15m\,50.3s\\[-.2pt]{\fontsize{5.8}{6.4}\selectfont \textcolor{BrickRed}{$ \blacktriangledown\,2.6\times $}\textsuperscript{c,e}}\\[-.2pt]15m\,42.3s\\[-.2pt]{\fontsize{5.8}{6.4}\selectfont \textcolor{BrickRed}{$ \blacktriangledown\,2.6\times $}}\end{tabular} & \hspace*{1pt}\begin{tabular}[t]{@{}c@{}}\rule{0pt}{12pt}\colorbox{BenchBest}{\strut 7m\,57s}\\[-.2pt]{\fontsize{5.8}{6.4}\selectfont \textcolor{ForestGreen}{$ \blacktriangle\,5.6\times $}\textsuperscript{c,d}}\\[-.2pt]\colorbox{BenchBest}{\strut 7m\,48.9s}\\[-.2pt]{\fontsize{5.8}{6.4}\selectfont \textcolor{ForestGreen}{$ \blacktriangle\,5.7\times $}}\end{tabular} & \hspace*{1pt}\begin{tabular}[t]{@{}c@{}}\rule{0pt}{12pt}6m\,51.2s\\[-.2pt]{\fontsize{5.8}{6.4}\selectfont \textcolor{BrickRed}{$ \blacktriangledown\,3.8\times $}\textsuperscript{c,e}}\\[-.2pt]6m\,49.1s\\[-.2pt]{\fontsize{5.8}{6.4}\selectfont \textcolor{BrickRed}{$ \blacktriangledown\,3.9\times $}}\end{tabular} & \hspace*{1pt}\begin{tabular}[t]{@{}c@{}}\rule{0pt}{12pt}6m\,41.1s\\[-.2pt]{\fontsize{5.8}{6.4}\selectfont \textcolor{BrickRed}{$ \blacktriangledown\,3.7\times $}\textsuperscript{c,e}}\\[-.2pt]6m\,37.2s\\[-.2pt]{\fontsize{5.8}{6.4}\selectfont \textcolor{BrickRed}{$ \blacktriangledown\,3.8\times $}}\end{tabular} & \hspace*{1pt}\begin{tabular}[t]{@{}c@{}}\rule{0pt}{12pt}12m\,42.6s\\[-.2pt]{\fontsize{5.8}{6.4}\selectfont \textcolor{BrickRed}{$ \blacktriangledown\,3.5\times $}\textsuperscript{c,e}}\\[-.2pt]12m\,35.1s\\[-.2pt]{\fontsize{5.8}{6.4}\selectfont \textcolor{BrickRed}{$ \blacktriangledown\,3.5\times $}}\end{tabular} & \hspace*{1pt}\begin{tabular}[t]{@{}c@{}}\rule{0pt}{12pt}\colorbox{BenchBest}{\strut 7m\,10.6s}\\[-.2pt]{\fontsize{5.8}{6.4}\selectfont \textcolor{ForestGreen}{$ \blacktriangle\,4.1\times $}\textsuperscript{c,d}}\\[-.2pt]\colorbox{BenchBest}{\strut 7m\,3.1s}\\[-.2pt]{\fontsize{5.8}{6.4}\selectfont \textcolor{ForestGreen}{$ \blacktriangle\,4.2\times $}}\end{tabular}\\[3pt]\hline
lexicmap & Native\textsuperscript{*} & \hspace*{1pt}\begin{tabular}[t]{@{}c@{}}\rule{0pt}{12pt}4m\,7.1s\\[-.2pt]{\fontsize{5.8}{6.4}\selectfont \textcolor{BrickRed}{$ \blacktriangledown\,6.8\times $}\textsuperscript{c}}\\[-.2pt]4m\,7s\\[-.2pt]{\fontsize{5.8}{6.4}\selectfont \textcolor{BrickRed}{$ \blacktriangledown\,7.0\times $}}\end{tabular} & \hspace*{1pt}\begin{tabular}[t]{@{}c@{}}\rule{0pt}{12pt}6m\,29.9s\\[-.2pt]{\fontsize{5.8}{6.4}\selectfont \textcolor{BrickRed}{$ \blacktriangledown\,10.2\times $}\textsuperscript{c}}\\[-.2pt]6m\,29.7s\\[-.2pt]{\fontsize{5.8}{6.4}\selectfont \textcolor{BrickRed}{$ \blacktriangledown\,10.8\times $}}\end{tabular} & \hspace*{1pt}\begin{tabular}[t]{@{}c@{}}\rule{0pt}{12pt}9m\,24.6s\\[-.2pt]{\fontsize{5.8}{6.4}\selectfont \textcolor{BrickRed}{$ \blacktriangledown\,6.0\times $}\textsuperscript{c}}\\[-.2pt]9m\,24.4s\\[-.2pt]{\fontsize{5.8}{6.4}\selectfont \textcolor{BrickRed}{$ \blacktriangledown\,6.3\times $}}\end{tabular} & \hspace*{1pt}\begin{tabular}[t]{@{}c@{}}\rule{0pt}{12pt}7m\,16.7s\\[-.2pt]{\fontsize{5.8}{6.4}\selectfont \textcolor{BrickRed}{$ \blacktriangledown\,12\% $}\textsuperscript{c}}\\[-.2pt]7m\,16.4s\\[-.2pt]{\fontsize{5.8}{6.4}\selectfont \textcolor{BrickRed}{$ \blacktriangledown\,13\% $}}\end{tabular} & \hspace*{1pt}\begin{tabular}[t]{@{}c@{}}\rule{0pt}{12pt}4m\,7.8s\\[-.2pt]{\fontsize{5.8}{6.4}\selectfont \textcolor{BrickRed}{$ \blacktriangledown\,7.9\times $}\textsuperscript{b,c}}\\[-.2pt]4m\,7.5s\\[-.2pt]{\fontsize{5.8}{6.4}\selectfont \textcolor{BrickRed}{$ \blacktriangledown\,8.2\times $}}\end{tabular} & \hspace*{1pt}\begin{tabular}[t]{@{}c@{}}\rule{0pt}{12pt}6m\,29.2s\\[-.2pt]{\fontsize{5.8}{6.4}\selectfont \textcolor{BrickRed}{$ \blacktriangledown\,11.4\times $}\textsuperscript{b,c}}\\[-.2pt]6m\,28.8s\\[-.2pt]{\fontsize{5.8}{6.4}\selectfont \textcolor{BrickRed}{$ \blacktriangledown\,12.1\times $}}\end{tabular} & \hspace*{1pt}\begin{tabular}[t]{@{}c@{}}\rule{0pt}{12pt}11m\,26.4s\\[-.2pt]{\fontsize{5.8}{6.4}\selectfont \textcolor{BrickRed}{$ \blacktriangledown\,9.9\times $}\textsuperscript{b,c}}\\[-.2pt]11m\,25.5s\\[-.2pt]{\fontsize{5.8}{6.4}\selectfont \textcolor{BrickRed}{$ \blacktriangledown\,10.4\times $}}\end{tabular} & \hspace*{1pt}\begin{tabular}[t]{@{}c@{}}\rule{0pt}{12pt}7m\,15.7s\\[-.2pt]{\fontsize{5.8}{6.4}\selectfont \textcolor{BrickRed}{$ \blacktriangledown\,2.4\times $}\textsuperscript{b,c}}\\[-.2pt]7m\,13s\\[-.2pt]{\fontsize{5.8}{6.4}\selectfont \textcolor{BrickRed}{$ \blacktriangledown\,2.5\times $}}\end{tabular} & \hspace*{1pt}\begin{tabular}[t]{@{}c@{}}\rule{0pt}{12pt}7m\,34.2s\\[-.2pt]{\fontsize{5.8}{6.4}\selectfont \textcolor{BrickRed}{$ \blacktriangledown\,2.6\times $}\textsuperscript{c}}\\[-.2pt]7m\,34.1s\\[-.2pt]{\fontsize{5.8}{6.4}\selectfont \textcolor{BrickRed}{$ \blacktriangledown\,2.6\times $}}\end{tabular} & \hspace*{1pt}\begin{tabular}[t]{@{}c@{}}\rule{0pt}{12pt}11m\,38.9s\\[-.2pt]{\fontsize{5.8}{6.4}\selectfont \textcolor{BrickRed}{$ \blacktriangledown\,4.1\times $}\textsuperscript{c}}\\[-.2pt]11m\,38.7s\\[-.2pt]{\fontsize{5.8}{6.4}\selectfont \textcolor{BrickRed}{$ \blacktriangledown\,4.2\times $}}\end{tabular} & \hspace*{1pt}\begin{tabular}[t]{@{}c@{}}\rule{0pt}{12pt}17m\,11.5s\\[-.2pt]{\fontsize{5.8}{6.4}\selectfont \textcolor{BrickRed}{$ \blacktriangledown\,2.8\times $}\textsuperscript{c}}\\[-.2pt]17m\,11.2s\\[-.2pt]{\fontsize{5.8}{6.4}\selectfont \textcolor{BrickRed}{$ \blacktriangledown\,2.8\times $}}\end{tabular} & \hspace*{1pt}\begin{tabular}[t]{@{}c@{}}\rule{0pt}{12pt}15m\,28s\\[-.2pt]{\fontsize{5.8}{6.4}\selectfont \textcolor{ForestGreen}{$ \blacktriangle\,2.9\times $}\textsuperscript{c}}\\[-.2pt]15m\,27.2s\\[-.2pt]{\fontsize{5.8}{6.4}\selectfont \textcolor{ForestGreen}{$ \blacktriangle\,2.9\times $}}\end{tabular} & \hspace*{1pt}\begin{tabular}[t]{@{}c@{}}\rule{0pt}{12pt}7m\,48.7s\\[-.2pt]{\fontsize{5.8}{6.4}\selectfont \textcolor{BrickRed}{$ \blacktriangledown\,4.4\times $}\textsuperscript{b,c}}\\[-.2pt]7m\,48.4s\\[-.2pt]{\fontsize{5.8}{6.4}\selectfont \textcolor{BrickRed}{$ \blacktriangledown\,4.4\times $}}\end{tabular} & \hspace*{1pt}\begin{tabular}[t]{@{}c@{}}\rule{0pt}{12pt}11m\,34.4s\\[-.2pt]{\fontsize{5.8}{6.4}\selectfont \textcolor{BrickRed}{$ \blacktriangledown\,6.5\times $}\textsuperscript{b,c}}\\[-.2pt]11m\,34s\\[-.2pt]{\fontsize{5.8}{6.4}\selectfont \textcolor{BrickRed}{$ \blacktriangledown\,6.6\times $}}\end{tabular} & \hspace*{1pt}\begin{tabular}[t]{@{}c@{}}\rule{0pt}{12pt}17m\,36.1s\\[-.2pt]{\fontsize{5.8}{6.4}\selectfont \textcolor{BrickRed}{$ \blacktriangledown\,4.8\times $}\textsuperscript{b,c}}\\[-.2pt]17m\,35.3s\\[-.2pt]{\fontsize{5.8}{6.4}\selectfont \textcolor{BrickRed}{$ \blacktriangledown\,4.9\times $}}\end{tabular} & \hspace*{1pt}\begin{tabular}[t]{@{}c@{}}\rule{0pt}{12pt}16m\,6s\\[-.2pt]{\fontsize{5.8}{6.4}\selectfont \textcolor{ForestGreen}{$ \blacktriangle\,1.8\times $}\textsuperscript{b,c}}\\[-.2pt]16m\,3.7s\\[-.2pt]{\fontsize{5.8}{6.4}\selectfont \textcolor{ForestGreen}{$ \blacktriangle\,1.8\times $}}\end{tabular}\\[3pt]\hline
\end{tabular}\par}
\smallskip{\footnotesize\noindent Each cell: A, comparison, B, comparison. Batches of 1,024 queries. Five repeats unless marked s (FAST: one sample). Blue = one lowest A and one lowest B down each machine/workload column.\par
\noindent Light gray = minimap2 baseline. Pale yellow = preferred MPF / Bounded / SIMD-DP / Block configuration, not a claim of winning every column. Exact ties select the first row.\par
\noindent Ratios use minimap2 as reference, except minimap2 itself uses the fixed BioseqDB row. $\cdot$ = no approved timing. n/a = no measured counterpart.\par
\noindent Both minimap2 and BioseqDB use the matching non-external remote measurements.\par
\noindent\textsuperscript{*}External A includes process startup and output. B removes different native stages and is not an equivalent-work comparison. BLAST and LexicMap have no $w$ or MAPQ equivalent.\par
\noindent\textsuperscript{a}BLAST DNA5 uses an N score of -3, compared with -1 in the baseline.\par
\noindent\textsuperscript{b}LexicMap stores reference N as A. Queries retain their bases, and edit-script checks use that native reference encoding.\par
\noindent\textsuperscript{c}External input uses two reference records for the 4 Gbp sequence to stay within signed 32-bit coordinates.\par
\noindent\textsuperscript{d}Winnowmap uses overlapping 250 Mb records and 499M index batches after an out-of-memory failure. Results are merged globally. This can affect MAPQ.\par
\noindent\textsuperscript{e}Winnowmap uses 1999M index batches and merges the results from multiple index parts.\par
\noindent\textsuperscript{g}Retained C7G minimap2 measurement comes from the earlier campaign with documented paging pressure.\par
}
\clearpage
\noindent{\large\bfseries Table \thetable: aligner comparison, 4 Gbp, k=20, w=15\hfill 2/4}\par\smallskip
\noindent{\footnotesize Different campaigns and native timing boundaries. A/B ratios are descriptive, subject to the measurement notes.}\par\smallskip
{\fontsize{7.1}{8.1}\selectfont\setlength{\tabcolsep}{1pt}\setlength{\fboxsep}{2pt}\renewcommand{\arraystretch}{1.12}
\noindent\begin{tabular}{@{}p{15mm}p{26mm}|p{13.55mm}|p{13.55mm}|p{13.55mm}|p{13.55mm}|p{13.55mm}|p{13.55mm}|p{13.55mm}|p{13.55mm}|p{13.55mm}|p{13.55mm}|p{13.55mm}|p{13.55mm}|p{13.55mm}|p{13.55mm}|p{13.55mm}|p{13.55mm}|@{}}\hline
Aligner & Configuration & \multicolumn{8}{c|}{\textbf{Index}} & \multicolumn{8}{c|}{\textbf{No index}}\\
 &  & \multicolumn{8}{c|}{\textbf{4 Gbp}} & \multicolumn{8}{c|}{\textbf{4 Gbp}}\\
 &  & \multicolumn{8}{c|}{\textbf{k=20, w=15}} & \multicolumn{8}{c|}{\textbf{k=20, w=15}}\\
 &  & \multicolumn{4}{c|}{\textbf{DNA4}} & \multicolumn{4}{c|}{\textbf{DNA5}} & \multicolumn{4}{c|}{\textbf{DNA4}} & \multicolumn{4}{c|}{\textbf{DNA5}}\\
 &  & \multicolumn{1}{c|}{\textbf{M7A}} & \multicolumn{1}{c|}{\textbf{M7I}} & \multicolumn{1}{c|}{\textbf{M4}} & \multicolumn{1}{c|}{\textbf{C7G}} & \multicolumn{1}{c|}{\textbf{M7A}} & \multicolumn{1}{c|}{\textbf{M7I}} & \multicolumn{1}{c|}{\textbf{M4}} & \multicolumn{1}{c|}{\textbf{C7G}} & \multicolumn{1}{c|}{\textbf{M7A}} & \multicolumn{1}{c|}{\textbf{M7I}} & \multicolumn{1}{c|}{\textbf{M4}} & \multicolumn{1}{c|}{\textbf{C7G}} & \multicolumn{1}{c|}{\textbf{M7A}} & \multicolumn{1}{c|}{\textbf{M7I}} & \multicolumn{1}{c|}{\textbf{M4}} & \multicolumn{1}{c|}{\textbf{C7G}}\\
\hline
\rowcolor{BenchBaseline}[0.5pt][0.5pt]minimap2 & Retained baseline & \hspace*{1pt}\begin{tabular}[t]{@{}c@{}}\rule{0pt}{12pt}2m\,48.1s\\[-.2pt]{\fontsize{5.8}{6.4}\selectfont \textcolor{BrickRed}{$ \blacktriangledown\,18.3\times $}}\\[-.2pt]2m\,41.7s\\[-.2pt]{\fontsize{5.8}{6.4}\selectfont \textcolor{BrickRed}{$ \blacktriangledown\,19.0\times $}}\end{tabular} & \hspace*{1pt}\begin{tabular}[t]{@{}c@{}}\rule{0pt}{12pt}2m\,50.3s\\[-.2pt]{\fontsize{5.8}{6.4}\selectfont \textcolor{BrickRed}{$ \blacktriangledown\,15.0\times $}}\\[-.2pt]2m\,43s\\[-.2pt]{\fontsize{5.8}{6.4}\selectfont \textcolor{BrickRed}{$ \blacktriangledown\,15.3\times $}}\end{tabular} & \hspace*{1pt}\begin{tabular}[t]{@{}c@{}}\rule{0pt}{12pt}4m\,54.3s\\[-.2pt]{\fontsize{5.8}{6.4}\selectfont \textcolor{BrickRed}{$ \blacktriangledown\,43.2\times $}}\\[-.2pt]4m\,39s\\[-.2pt]{\fontsize{5.8}{6.4}\selectfont \textcolor{BrickRed}{$ \blacktriangledown\,54.6\times $}}\end{tabular} & \hspace*{1pt}\begin{tabular}[t]{@{}c@{}}\rule{0pt}{12pt}25m\,5.5s\\[-.2pt]{\fontsize{5.8}{6.4}\selectfont \textcolor{gray}{n/a}\textsuperscript{g}}\\[-.2pt]25m\,0.7s\\[-.2pt]{\fontsize{5.8}{6.4}\selectfont \textcolor{gray}{n/a}}\end{tabular} & \hspace*{1pt}\begin{tabular}[t]{@{}c@{}}\rule{0pt}{12pt}33.9s\\[-.2pt]{\fontsize{5.8}{6.4}\selectfont \textcolor{BrickRed}{$ \blacktriangledown\,1.9\times $}}\\[-.2pt]32.9s\\[-.2pt]{\fontsize{5.8}{6.4}\selectfont \textcolor{BrickRed}{$ \blacktriangledown\,2.1\times $}}\end{tabular} & \hspace*{1pt}\begin{tabular}[t]{@{}c@{}}\rule{0pt}{12pt}34.4s\\[-.2pt]{\fontsize{5.8}{6.4}\selectfont \textcolor{BrickRed}{$ \blacktriangledown\,2.6\times $}}\\[-.2pt]32.7s\\[-.2pt]{\fontsize{5.8}{6.4}\selectfont \textcolor{BrickRed}{$ \blacktriangledown\,2.9\times $}}\end{tabular} & \hspace*{1pt}\begin{tabular}[t]{@{}c@{}}\rule{0pt}{12pt}1m\,8.2s\\[-.2pt]{\fontsize{5.8}{6.4}\selectfont \textcolor{BrickRed}{$ \blacktriangledown\,3.7\times $}}\\[-.2pt]1m\,5.1s\\[-.2pt]{\fontsize{5.8}{6.4}\selectfont \textcolor{BrickRed}{$ \blacktriangledown\,4.6\times $}}\end{tabular} & \hspace*{1pt}\begin{tabular}[t]{@{}c@{}}\rule{0pt}{12pt}3m\,9.8s\\[-.2pt]{\fontsize{5.8}{6.4}\selectfont \textcolor{gray}{n/a}\textsuperscript{g}}\\[-.2pt]3m\,3.7s\\[-.2pt]{\fontsize{5.8}{6.4}\selectfont \textcolor{gray}{n/a}}\end{tabular} & \hspace*{1pt}\begin{tabular}[t]{@{}c@{}}\rule{0pt}{12pt}3m\,28.3s\\[-.2pt]{\fontsize{5.8}{6.4}\selectfont \textcolor{BrickRed}{$ \blacktriangledown\,1.3\times $}}\\[-.2pt]3m\,22.8s\\[-.2pt]{\fontsize{5.8}{6.4}\selectfont \textcolor{BrickRed}{$ \blacktriangledown\,1.4\times $}}\end{tabular} & \hspace*{1pt}\begin{tabular}[t]{@{}c@{}}\rule{0pt}{12pt}3m\,53.2s\\[-.2pt]{\fontsize{5.8}{6.4}\selectfont \textcolor{BrickRed}{$ \blacktriangledown\,1.4\times $}}\\[-.2pt]3m\,43.5s\\[-.2pt]{\fontsize{5.8}{6.4}\selectfont \textcolor{BrickRed}{$ \blacktriangledown\,1.4\times $}}\end{tabular} & \hspace*{1pt}\begin{tabular}[t]{@{}c@{}}\rule{0pt}{12pt}7m\,14.6s\\[-.2pt]{\fontsize{5.8}{6.4}\selectfont \textcolor{BrickRed}{$ \blacktriangledown\,1.4\times $}}\\[-.2pt]6m\,57.2s\\[-.2pt]{\fontsize{5.8}{6.4}\selectfont \textcolor{BrickRed}{$ \blacktriangledown\,1.5\times $}}\end{tabular} & \hspace*{1pt}\begin{tabular}[t]{@{}c@{}}\rule{0pt}{12pt}51m\,22s\\[-.2pt]{\fontsize{5.8}{6.4}\selectfont \textcolor{gray}{n/a}\textsuperscript{g}}\\[-.2pt]51m\,17.2s\\[-.2pt]{\fontsize{5.8}{6.4}\selectfont \textcolor{gray}{n/a}}\end{tabular} & \hspace*{1pt}\begin{tabular}[t]{@{}c@{}}\rule{0pt}{12pt}\colorbox{BenchBest}{\strut 1m\,10.3s}\\[-.2pt]{\fontsize{5.8}{6.4}\selectfont \textcolor{ForestGreen}{$ \blacktriangle\,7.6\times $}}\\[-.2pt]\colorbox{BenchBest}{\strut 1m\,9s}\\[-.2pt]{\fontsize{5.8}{6.4}\selectfont \textcolor{ForestGreen}{$ \blacktriangle\,2.4\times $}}\end{tabular} & \hspace*{1pt}\begin{tabular}[t]{@{}c@{}}\rule{0pt}{12pt}\colorbox{BenchBest}{\strut 1m\,8.6s}\\[-.2pt]{\fontsize{5.8}{6.4}\selectfont \textcolor{ForestGreen}{$ \blacktriangle\,6.4\times $}}\\[-.2pt]\colorbox{BenchBest}{\strut 1m\,6.9s}\\[-.2pt]{\fontsize{5.8}{6.4}\selectfont \textcolor{ForestGreen}{$ \blacktriangle\,2.5\times $}}\end{tabular} & \hspace*{1pt}\begin{tabular}[t]{@{}c@{}}\rule{0pt}{12pt}\colorbox{BenchBest}{\strut 2m\,21.5s}\\[-.2pt]{\fontsize{5.8}{6.4}\selectfont \textcolor{ForestGreen}{$ \blacktriangle\,6.2\times $}}\\[-.2pt]\colorbox{BenchBest}{\strut 2m\,18.4s}\\[-.2pt]{\fontsize{5.8}{6.4}\selectfont \textcolor{ForestGreen}{$ \blacktriangle\,2.2\times $}}\end{tabular} & \hspace*{1pt}\begin{tabular}[t]{@{}c@{}}\rule{0pt}{12pt}15m\,56.8s\\[-.2pt]{\fontsize{5.8}{6.4}\selectfont \textcolor{gray}{n/a}\textsuperscript{g}}\\[-.2pt]15m\,51.9s\\[-.2pt]{\fontsize{5.8}{6.4}\selectfont \textcolor{gray}{n/a}}\end{tabular}\\[3pt]\hline
\rowcolor{BenchPreferred}[0.5pt][0.5pt]BioseqDB & \begin{tabular}[t]{@{}l@{}}\rule{0pt}{12pt}MPF / Bounded\\SIMD-DP / Block\end{tabular} & \hspace*{1pt}\begin{tabular}[t]{@{}c@{}}\rule{0pt}{12pt}\colorbox{BenchBest}{\strut 9.2s}\\[-.2pt]{\fontsize{5.8}{6.4}\selectfont \textcolor{ForestGreen}{$ \blacktriangle\,18.3\times $}}\\[-.2pt]\colorbox{BenchBest}{\strut 8.5s}\\[-.2pt]{\fontsize{5.8}{6.4}\selectfont \textcolor{ForestGreen}{$ \blacktriangle\,19.0\times $}}\end{tabular} & \hspace*{1pt}\begin{tabular}[t]{@{}c@{}}\rule{0pt}{12pt}\colorbox{BenchBest}{\strut 11.4s}\\[-.2pt]{\fontsize{5.8}{6.4}\selectfont \textcolor{ForestGreen}{$ \blacktriangle\,15.0\times $}}\\[-.2pt]\colorbox{BenchBest}{\strut 10.7s}\\[-.2pt]{\fontsize{5.8}{6.4}\selectfont \textcolor{ForestGreen}{$ \blacktriangle\,15.3\times $}}\end{tabular} & \hspace*{1pt}\begin{tabular}[t]{@{}c@{}}\rule{0pt}{12pt}\colorbox{BenchBest}{\strut 6.8s}\\[-.2pt]{\fontsize{5.8}{6.4}\selectfont \textcolor{ForestGreen}{$ \blacktriangle\,43.2\times $}}\\[-.2pt]\colorbox{BenchBest}{\strut 5.1s}\\[-.2pt]{\fontsize{5.8}{6.4}\selectfont \textcolor{ForestGreen}{$ \blacktriangle\,54.6\times $}}\end{tabular} & \textcolor{gray}{\ensuremath{\cdot}} & \hspace*{1pt}\begin{tabular}[t]{@{}c@{}}\rule{0pt}{12pt}\colorbox{BenchBest}{\strut 18.2s}\\[-.2pt]{\fontsize{5.8}{6.4}\selectfont \textcolor{ForestGreen}{$ \blacktriangle\,1.9\times $}}\\[-.2pt]\colorbox{BenchBest}{\strut 16s}\\[-.2pt]{\fontsize{5.8}{6.4}\selectfont \textcolor{ForestGreen}{$ \blacktriangle\,2.1\times $}}\end{tabular} & \hspace*{1pt}\begin{tabular}[t]{@{}c@{}}\rule{0pt}{12pt}\colorbox{BenchBest}{\strut 13.2s}\\[-.2pt]{\fontsize{5.8}{6.4}\selectfont \textcolor{ForestGreen}{$ \blacktriangle\,2.6\times $}}\\[-.2pt]\colorbox{BenchBest}{\strut 11.4s}\\[-.2pt]{\fontsize{5.8}{6.4}\selectfont \textcolor{ForestGreen}{$ \blacktriangle\,2.9\times $}}\end{tabular} & \hspace*{1pt}\begin{tabular}[t]{@{}c@{}}\rule{0pt}{12pt}\colorbox{BenchBest}{\strut 18.6s}\\[-.2pt]{\fontsize{5.8}{6.4}\selectfont \textcolor{ForestGreen}{$ \blacktriangle\,3.7\times $}}\\[-.2pt]\colorbox{BenchBest}{\strut 14.3s}\\[-.2pt]{\fontsize{5.8}{6.4}\selectfont \textcolor{ForestGreen}{$ \blacktriangle\,4.6\times $}}\end{tabular} & \textcolor{gray}{\ensuremath{\cdot}} & \hspace*{1pt}\begin{tabular}[t]{@{}c@{}}\rule{0pt}{12pt}\colorbox{BenchBest}{\strut 2m\,38.5s}\\[-.2pt]{\fontsize{5.8}{6.4}\selectfont \textcolor{ForestGreen}{$ \blacktriangle\,24\% $}}\\[-.2pt]\colorbox{BenchBest}{\strut 2m\,29.8s}\\[-.2pt]{\fontsize{5.8}{6.4}\selectfont \textcolor{ForestGreen}{$ \blacktriangle\,26\% $}}\end{tabular} & \hspace*{1pt}\begin{tabular}[t]{@{}c@{}}\rule{0pt}{12pt}\colorbox{BenchBest}{\strut 2m\,49s}\\[-.2pt]{\fontsize{5.8}{6.4}\selectfont \textcolor{ForestGreen}{$ \blacktriangle\,28\% $}}\\[-.2pt]\colorbox{BenchBest}{\strut 2m\,37s}\\[-.2pt]{\fontsize{5.8}{6.4}\selectfont \textcolor{ForestGreen}{$ \blacktriangle\,30\% $}}\end{tabular} & \hspace*{1pt}\begin{tabular}[t]{@{}c@{}}\rule{0pt}{12pt}\colorbox{BenchBest}{\strut 5m\,3.3s}\\[-.2pt]{\fontsize{5.8}{6.4}\selectfont \textcolor{ForestGreen}{$ \blacktriangle\,1.4\times $}}\\[-.2pt]\colorbox{BenchBest}{\strut 4m\,42.7s}\\[-.2pt]{\fontsize{5.8}{6.4}\selectfont \textcolor{ForestGreen}{$ \blacktriangle\,1.5\times $}}\end{tabular} & \textcolor{gray}{\ensuremath{\cdot}} & \hspace*{1pt}\begin{tabular}[t]{@{}c@{}}\rule{0pt}{12pt}8m\,52.3s\\[-.2pt]{\fontsize{5.8}{6.4}\selectfont \textcolor{BrickRed}{$ \blacktriangledown\,7.6\times $}}\\[-.2pt]2m\,45.5s\\[-.2pt]{\fontsize{5.8}{6.4}\selectfont \textcolor{BrickRed}{$ \blacktriangledown\,2.4\times $}}\end{tabular} & \hspace*{1pt}\begin{tabular}[t]{@{}c@{}}\rule{0pt}{12pt}7m\,22.2s\\[-.2pt]{\fontsize{5.8}{6.4}\selectfont \textcolor{BrickRed}{$ \blacktriangledown\,6.4\times $}}\\[-.2pt]2m\,48.2s\\[-.2pt]{\fontsize{5.8}{6.4}\selectfont \textcolor{BrickRed}{$ \blacktriangledown\,2.5\times $}}\end{tabular} & \hspace*{1pt}\begin{tabular}[t]{@{}c@{}}\rule{0pt}{12pt}14m\,43.3s\\[-.2pt]{\fontsize{5.8}{6.4}\selectfont \textcolor{BrickRed}{$ \blacktriangledown\,6.2\times $}}\\[-.2pt]5m\,9.5s\\[-.2pt]{\fontsize{5.8}{6.4}\selectfont \textcolor{BrickRed}{$ \blacktriangledown\,2.2\times $}}\end{tabular} & \textcolor{gray}{\ensuremath{\cdot}}\\[3pt]\hline
blastn & Native\textsuperscript{*} & \hspace*{1pt}\begin{tabular}[t]{@{}c@{}}\rule{0pt}{12pt}4m\,35.4s\\[-.2pt]{\fontsize{5.8}{6.4}\selectfont \textcolor{BrickRed}{$ \blacktriangledown\,1.6\times $}\textsuperscript{c}}\\[-.2pt]4m\,30.3s\\[-.2pt]{\fontsize{5.8}{6.4}\selectfont \textcolor{BrickRed}{$ \blacktriangledown\,1.7\times $}}\end{tabular} & \hspace*{1pt}\begin{tabular}[t]{@{}c@{}}\rule{0pt}{12pt}6m\,21.2s\\[-.2pt]{\fontsize{5.8}{6.4}\selectfont \textcolor{BrickRed}{$ \blacktriangledown\,2.2\times $}\textsuperscript{c}}\\[-.2pt]6m\,10.8s\\[-.2pt]{\fontsize{5.8}{6.4}\selectfont \textcolor{BrickRed}{$ \blacktriangledown\,2.3\times $}}\end{tabular} & \hspace*{1pt}\begin{tabular}[t]{@{}c@{}}\rule{0pt}{12pt}8m\,55.3s\\[-.2pt]{\fontsize{5.8}{6.4}\selectfont \textcolor{BrickRed}{$ \blacktriangledown\,1.8\times $}\textsuperscript{c}}\\[-.2pt]8m\,41.3s\\[-.2pt]{\fontsize{5.8}{6.4}\selectfont \textcolor{BrickRed}{$ \blacktriangledown\,1.9\times $}}\end{tabular} & \hspace*{1pt}\begin{tabular}[t]{@{}c@{}}\rule{0pt}{12pt}5m\,9.4s\\[-.2pt]{\fontsize{5.8}{6.4}\selectfont \textcolor{ForestGreen}{$ \blacktriangle\,4.9\times $}\textsuperscript{c}}\\[-.2pt]5m\,4.2s\\[-.2pt]{\fontsize{5.8}{6.4}\selectfont \textcolor{ForestGreen}{$ \blacktriangle\,4.9\times $}}\end{tabular} & \hspace*{1pt}\begin{tabular}[t]{@{}c@{}}\rule{0pt}{12pt}5m\,39.4s\\[-.2pt]{\fontsize{5.8}{6.4}\selectfont \textcolor{BrickRed}{$ \blacktriangledown\,10.0\times $}\textsuperscript{a,c}}\\[-.2pt]4m\,13.6s\\[-.2pt]{\fontsize{5.8}{6.4}\selectfont \textcolor{BrickRed}{$ \blacktriangledown\,7.7\times $}}\end{tabular} & \hspace*{1pt}\begin{tabular}[t]{@{}c@{}}\rule{0pt}{12pt}7m\,32.4s\\[-.2pt]{\fontsize{5.8}{6.4}\selectfont \textcolor{BrickRed}{$ \blacktriangledown\,13.1\times $}\textsuperscript{a,c}}\\[-.2pt]5m\,27.1s\\[-.2pt]{\fontsize{5.8}{6.4}\selectfont \textcolor{BrickRed}{$ \blacktriangledown\,10.0\times $}}\end{tabular} & \hspace*{1pt}\begin{tabular}[t]{@{}c@{}}\rule{0pt}{12pt}11m\,29.3s\\[-.2pt]{\fontsize{5.8}{6.4}\selectfont \textcolor{BrickRed}{$ \blacktriangledown\,10.1\times $}\textsuperscript{a,c}}\\[-.2pt]8m\,4s\\[-.2pt]{\fontsize{5.8}{6.4}\selectfont \textcolor{BrickRed}{$ \blacktriangledown\,7.4\times $}}\end{tabular} & \hspace*{1pt}\begin{tabular}[t]{@{}c@{}}\rule{0pt}{12pt}6m\,10.2s\\[-.2pt]{\fontsize{5.8}{6.4}\selectfont \textcolor{BrickRed}{$ \blacktriangledown\,2.0\times $}\textsuperscript{a,c}}\\[-.2pt]4m\,39s\\[-.2pt]{\fontsize{5.8}{6.4}\selectfont \textcolor{BrickRed}{$ \blacktriangledown\,1.5\times $}}\end{tabular} & \hspace*{1pt}\begin{tabular}[t]{@{}c@{}}\rule{0pt}{12pt}4m\,58.3s\\[-.2pt]{\fontsize{5.8}{6.4}\selectfont \textcolor{BrickRed}{$ \blacktriangledown\,1.4\times $}\textsuperscript{c}}\\[-.2pt]4m\,53.3s\\[-.2pt]{\fontsize{5.8}{6.4}\selectfont \textcolor{BrickRed}{$ \blacktriangledown\,1.4\times $}}\end{tabular} & \hspace*{1pt}\begin{tabular}[t]{@{}c@{}}\rule{0pt}{12pt}6m\,40.8s\\[-.2pt]{\fontsize{5.8}{6.4}\selectfont \textcolor{BrickRed}{$ \blacktriangledown\,1.7\times $}\textsuperscript{c}}\\[-.2pt]6m\,30.3s\\[-.2pt]{\fontsize{5.8}{6.4}\selectfont \textcolor{BrickRed}{$ \blacktriangledown\,1.7\times $}}\end{tabular} & \hspace*{1pt}\begin{tabular}[t]{@{}c@{}}\rule{0pt}{12pt}9m\,34s\\[-.2pt]{\fontsize{5.8}{6.4}\selectfont \textcolor{BrickRed}{$ \blacktriangledown\,1.3\times $}\textsuperscript{c}}\\[-.2pt]9m\,19.8s\\[-.2pt]{\fontsize{5.8}{6.4}\selectfont \textcolor{BrickRed}{$ \blacktriangledown\,1.3\times $}}\end{tabular} & \hspace*{1pt}\begin{tabular}[t]{@{}c@{}}\rule{0pt}{12pt}\colorbox{BenchBest}{\strut 5m\,40.5s}\\[-.2pt]{\fontsize{5.8}{6.4}\selectfont \textcolor{ForestGreen}{$ \blacktriangle\,9.1\times $}\textsuperscript{c}}\\[-.2pt]\colorbox{BenchBest}{\strut 5m\,35.1s}\\[-.2pt]{\fontsize{5.8}{6.4}\selectfont \textcolor{ForestGreen}{$ \blacktriangle\,9.2\times $}}\end{tabular} & \hspace*{1pt}\begin{tabular}[t]{@{}c@{}}\rule{0pt}{12pt}6m\,15.6s\\[-.2pt]{\fontsize{5.8}{6.4}\selectfont \textcolor{BrickRed}{$ \blacktriangledown\,5.3\times $}\textsuperscript{a,c}}\\[-.2pt]4m\,49.2s\\[-.2pt]{\fontsize{5.8}{6.4}\selectfont \textcolor{BrickRed}{$ \blacktriangledown\,4.2\times $}}\end{tabular} & \hspace*{1pt}\begin{tabular}[t]{@{}c@{}}\rule{0pt}{12pt}8m\,11.1s\\[-.2pt]{\fontsize{5.8}{6.4}\selectfont \textcolor{BrickRed}{$ \blacktriangledown\,7.2\times $}\textsuperscript{a,c}}\\[-.2pt]6m\,2.7s\\[-.2pt]{\fontsize{5.8}{6.4}\selectfont \textcolor{BrickRed}{$ \blacktriangledown\,5.4\times $}}\end{tabular} & \hspace*{1pt}\begin{tabular}[t]{@{}c@{}}\rule{0pt}{12pt}12m\,29.6s\\[-.2pt]{\fontsize{5.8}{6.4}\selectfont \textcolor{BrickRed}{$ \blacktriangledown\,5.3\times $}\textsuperscript{a,c}}\\[-.2pt]9m\,3.4s\\[-.2pt]{\fontsize{5.8}{6.4}\selectfont \textcolor{BrickRed}{$ \blacktriangledown\,3.9\times $}}\end{tabular} & \hspace*{1pt}\begin{tabular}[t]{@{}c@{}}\rule{0pt}{12pt}\colorbox{BenchBest}{\strut 7m\,4.7s}\\[-.2pt]{\fontsize{5.8}{6.4}\selectfont \textcolor{ForestGreen}{$ \blacktriangle\,2.3\times $}\textsuperscript{a,c}}\\[-.2pt]\colorbox{BenchBest}{\strut 5m\,34.7s}\\[-.2pt]{\fontsize{5.8}{6.4}\selectfont \textcolor{ForestGreen}{$ \blacktriangle\,2.8\times $}}\end{tabular}\\[3pt]\hline
winnowmap & Native\textsuperscript{*} & \hspace*{1pt}\begin{tabular}[t]{@{}c@{}}\rule{0pt}{12pt}55.6s\\[-.2pt]{\fontsize{5.8}{6.4}\selectfont \textcolor{ForestGreen}{$ \blacktriangle\,3.0\times $}\textsuperscript{c,e}}\\[-.2pt]53.8s\\[-.2pt]{\fontsize{5.8}{6.4}\selectfont \textcolor{ForestGreen}{$ \blacktriangle\,3.0\times $}}\end{tabular} & \hspace*{1pt}\begin{tabular}[t]{@{}c@{}}\rule{0pt}{12pt}59.6s\\[-.2pt]{\fontsize{5.8}{6.4}\selectfont \textcolor{ForestGreen}{$ \blacktriangle\,2.9\times $}\textsuperscript{c,e}}\\[-.2pt]56.2s\\[-.2pt]{\fontsize{5.8}{6.4}\selectfont \textcolor{ForestGreen}{$ \blacktriangle\,2.9\times $}}\end{tabular} & \hspace*{1pt}\begin{tabular}[t]{@{}c@{}}\rule{0pt}{12pt}1m\,43s\\[-.2pt]{\fontsize{5.8}{6.4}\selectfont \textcolor{ForestGreen}{$ \blacktriangle\,2.9\times $}\textsuperscript{c,e}}\\[-.2pt]1m\,36.3s\\[-.2pt]{\fontsize{5.8}{6.4}\selectfont \textcolor{ForestGreen}{$ \blacktriangle\,2.9\times $}}\end{tabular} & \hspace*{1pt}\begin{tabular}[t]{@{}c@{}}\rule{0pt}{12pt}\colorbox{BenchBest}{\strut 1m\,0.7s}\\[-.2pt]{\fontsize{5.8}{6.4}\selectfont \textcolor{ForestGreen}{$ \blacktriangle\,24.8\times $}\textsuperscript{c,d}}\\[-.2pt]\colorbox{BenchBest}{\strut 54.1s}\\[-.2pt]{\fontsize{5.8}{6.4}\selectfont \textcolor{ForestGreen}{$ \blacktriangle\,27.7\times $}}\end{tabular} & \hspace*{1pt}\begin{tabular}[t]{@{}c@{}}\rule{0pt}{12pt}24.6s\\[-.2pt]{\fontsize{5.8}{6.4}\selectfont \textcolor{ForestGreen}{$ \blacktriangle\,28\% $}\textsuperscript{c,e}}\\[-.2pt]22.8s\\[-.2pt]{\fontsize{5.8}{6.4}\selectfont \textcolor{ForestGreen}{$ \blacktriangle\,1.4\times $}}\end{tabular} & \hspace*{1pt}\begin{tabular}[t]{@{}c@{}}\rule{0pt}{12pt}23.7s\\[-.2pt]{\fontsize{5.8}{6.4}\selectfont \textcolor{ForestGreen}{$ \blacktriangle\,1.4\times $}\textsuperscript{c,e}}\\[-.2pt]20.4s\\[-.2pt]{\fontsize{5.8}{6.4}\selectfont \textcolor{ForestGreen}{$ \blacktriangle\,1.6\times $}}\end{tabular} & \hspace*{1pt}\begin{tabular}[t]{@{}c@{}}\rule{0pt}{12pt}36.1s\\[-.2pt]{\fontsize{5.8}{6.4}\selectfont \textcolor{ForestGreen}{$ \blacktriangle\,1.9\times $}\textsuperscript{c,e}}\\[-.2pt]29.6s\\[-.2pt]{\fontsize{5.8}{6.4}\selectfont \textcolor{ForestGreen}{$ \blacktriangle\,2.2\times $}}\end{tabular} & \hspace*{1pt}\begin{tabular}[t]{@{}c@{}}\rule{0pt}{12pt}\colorbox{BenchBest}{\strut 41.4s}\\[-.2pt]{\fontsize{5.8}{6.4}\selectfont \textcolor{ForestGreen}{$ \blacktriangle\,4.6\times $}\textsuperscript{c,d}}\\[-.2pt]\colorbox{BenchBest}{\strut 34.7s}\\[-.2pt]{\fontsize{5.8}{6.4}\selectfont \textcolor{ForestGreen}{$ \blacktriangle\,5.3\times $}}\end{tabular} & \hspace*{1pt}\begin{tabular}[t]{@{}c@{}}\rule{0pt}{12pt}10m\,41.8s\\[-.2pt]{\fontsize{5.8}{6.4}\selectfont \textcolor{BrickRed}{$ \blacktriangledown\,3.1\times $}\textsuperscript{c,e}}\\[-.2pt]10m\,40s\\[-.2pt]{\fontsize{5.8}{6.4}\selectfont \textcolor{BrickRed}{$ \blacktriangledown\,3.2\times $}}\end{tabular} & \hspace*{1pt}\begin{tabular}[t]{@{}c@{}}\rule{0pt}{12pt}11m\,28.6s\\[-.2pt]{\fontsize{5.8}{6.4}\selectfont \textcolor{BrickRed}{$ \blacktriangledown\,3.0\times $}\textsuperscript{c,e}}\\[-.2pt]11m\,25.2s\\[-.2pt]{\fontsize{5.8}{6.4}\selectfont \textcolor{BrickRed}{$ \blacktriangledown\,3.1\times $}}\end{tabular} & \hspace*{1pt}\begin{tabular}[t]{@{}c@{}}\rule{0pt}{12pt}33m\,6.8s\\[-.2pt]{\fontsize{5.8}{6.4}\selectfont \textcolor{BrickRed}{$ \blacktriangledown\,4.6\times $}\textsuperscript{c,e}}\\[-.2pt]33m\,0.1s\\[-.2pt]{\fontsize{5.8}{6.4}\selectfont \textcolor{BrickRed}{$ \blacktriangledown\,4.7\times $}}\end{tabular} & \hspace*{1pt}\begin{tabular}[t]{@{}c@{}}\rule{0pt}{12pt}12m\,28.4s\\[-.2pt]{\fontsize{5.8}{6.4}\selectfont \textcolor{ForestGreen}{$ \blacktriangle\,4.1\times $}\textsuperscript{c,d}}\\[-.2pt]12m\,21.9s\\[-.2pt]{\fontsize{5.8}{6.4}\selectfont \textcolor{ForestGreen}{$ \blacktriangle\,4.1\times $}}\end{tabular} & \hspace*{1pt}\begin{tabular}[t]{@{}c@{}}\rule{0pt}{12pt}8m\,28.7s\\[-.2pt]{\fontsize{5.8}{6.4}\selectfont \textcolor{BrickRed}{$ \blacktriangledown\,7.2\times $}\textsuperscript{c,e}}\\[-.2pt]8m\,26.9s\\[-.2pt]{\fontsize{5.8}{6.4}\selectfont \textcolor{BrickRed}{$ \blacktriangledown\,7.3\times $}}\end{tabular} & \hspace*{1pt}\begin{tabular}[t]{@{}c@{}}\rule{0pt}{12pt}8m\,59.5s\\[-.2pt]{\fontsize{5.8}{6.4}\selectfont \textcolor{BrickRed}{$ \blacktriangledown\,7.9\times $}\textsuperscript{c,e}}\\[-.2pt]8m\,56.1s\\[-.2pt]{\fontsize{5.8}{6.4}\selectfont \textcolor{BrickRed}{$ \blacktriangledown\,8.0\times $}}\end{tabular} & \hspace*{1pt}\begin{tabular}[t]{@{}c@{}}\rule{0pt}{12pt}24m\,58.6s\\[-.2pt]{\fontsize{5.8}{6.4}\selectfont \textcolor{BrickRed}{$ \blacktriangledown\,10.6\times $}\textsuperscript{c,e}}\\[-.2pt]24m\,52s\\[-.2pt]{\fontsize{5.8}{6.4}\selectfont \textcolor{BrickRed}{$ \blacktriangledown\,10.8\times $}}\end{tabular} & \hspace*{1pt}\begin{tabular}[t]{@{}c@{}}\rule{0pt}{12pt}10m\,19s\\[-.2pt]{\fontsize{5.8}{6.4}\selectfont \textcolor{ForestGreen}{$ \blacktriangle\,1.5\times $}\textsuperscript{c,d}}\\[-.2pt]10m\,12.1s\\[-.2pt]{\fontsize{5.8}{6.4}\selectfont \textcolor{ForestGreen}{$ \blacktriangle\,1.6\times $}}\end{tabular}\\[3pt]\hline
lexicmap & Native\textsuperscript{*} & \hspace*{1pt}\begin{tabular}[t]{@{}c@{}}\rule{0pt}{12pt}56.9s\\[-.2pt]{\fontsize{5.8}{6.4}\selectfont \textcolor{ForestGreen}{$ \blacktriangle\,3.0\times $}\textsuperscript{c}}\\[-.2pt]56.7s\\[-.2pt]{\fontsize{5.8}{6.4}\selectfont \textcolor{ForestGreen}{$ \blacktriangle\,2.9\times $}}\end{tabular} & \hspace*{1pt}\begin{tabular}[t]{@{}c@{}}\rule{0pt}{12pt}1m\,21.3s\\[-.2pt]{\fontsize{5.8}{6.4}\selectfont \textcolor{ForestGreen}{$ \blacktriangle\,2.1\times $}\textsuperscript{c}}\\[-.2pt]1m\,21.1s\\[-.2pt]{\fontsize{5.8}{6.4}\selectfont \textcolor{ForestGreen}{$ \blacktriangle\,2.0\times $}}\end{tabular} & \hspace*{1pt}\begin{tabular}[t]{@{}c@{}}\rule{0pt}{12pt}2m\,1.2s\\[-.2pt]{\fontsize{5.8}{6.4}\selectfont \textcolor{ForestGreen}{$ \blacktriangle\,2.4\times $}\textsuperscript{c}}\\[-.2pt]2m\,0.9s\\[-.2pt]{\fontsize{5.8}{6.4}\selectfont \textcolor{ForestGreen}{$ \blacktriangle\,2.3\times $}}\end{tabular} & \hspace*{1pt}\begin{tabular}[t]{@{}c@{}}\rule{0pt}{12pt}1m\,16.5s\\[-.2pt]{\fontsize{5.8}{6.4}\selectfont \textcolor{ForestGreen}{$ \blacktriangle\,19.7\times $}\textsuperscript{c}}\\[-.2pt]1m\,16.1s\\[-.2pt]{\fontsize{5.8}{6.4}\selectfont \textcolor{ForestGreen}{$ \blacktriangle\,19.7\times $}}\end{tabular} & \hspace*{1pt}\begin{tabular}[t]{@{}c@{}}\rule{0pt}{12pt}56.4s\\[-.2pt]{\fontsize{5.8}{6.4}\selectfont \textcolor{BrickRed}{$ \blacktriangledown\,1.7\times $}\textsuperscript{b,c}}\\[-.2pt]55.8s\\[-.2pt]{\fontsize{5.8}{6.4}\selectfont \textcolor{BrickRed}{$ \blacktriangledown\,1.7\times $}}\end{tabular} & \hspace*{1pt}\begin{tabular}[t]{@{}c@{}}\rule{0pt}{12pt}1m\,20.1s\\[-.2pt]{\fontsize{5.8}{6.4}\selectfont \textcolor{BrickRed}{$ \blacktriangledown\,2.3\times $}\textsuperscript{b,c}}\\[-.2pt]1m\,19.6s\\[-.2pt]{\fontsize{5.8}{6.4}\selectfont \textcolor{BrickRed}{$ \blacktriangledown\,2.4\times $}}\end{tabular} & \hspace*{1pt}\begin{tabular}[t]{@{}c@{}}\rule{0pt}{12pt}2m\,9.7s\\[-.2pt]{\fontsize{5.8}{6.4}\selectfont \textcolor{BrickRed}{$ \blacktriangledown\,1.9\times $}\textsuperscript{b,c}}\\[-.2pt]2m\,8.5s\\[-.2pt]{\fontsize{5.8}{6.4}\selectfont \textcolor{BrickRed}{$ \blacktriangledown\,2.0\times $}}\end{tabular} & \hspace*{1pt}\begin{tabular}[t]{@{}c@{}}\rule{0pt}{12pt}1m\,14.6s\\[-.2pt]{\fontsize{5.8}{6.4}\selectfont \textcolor{ForestGreen}{$ \blacktriangle\,2.5\times $}\textsuperscript{b,c}}\\[-.2pt]1m\,12.8s\\[-.2pt]{\fontsize{5.8}{6.4}\selectfont \textcolor{ForestGreen}{$ \blacktriangle\,2.5\times $}}\end{tabular} & \hspace*{1pt}\begin{tabular}[t]{@{}c@{}}\rule{0pt}{12pt}4m\,52.5s\\[-.2pt]{\fontsize{5.8}{6.4}\selectfont \textcolor{BrickRed}{$ \blacktriangledown\,1.4\times $}\textsuperscript{c}}\\[-.2pt]4m\,52.3s\\[-.2pt]{\fontsize{5.8}{6.4}\selectfont \textcolor{BrickRed}{$ \blacktriangledown\,1.4\times $}}\end{tabular} & \hspace*{1pt}\begin{tabular}[t]{@{}c@{}}\rule{0pt}{12pt}6m\,45.3s\\[-.2pt]{\fontsize{5.8}{6.4}\selectfont \textcolor{BrickRed}{$ \blacktriangledown\,1.7\times $}\textsuperscript{c}}\\[-.2pt]6m\,45.1s\\[-.2pt]{\fontsize{5.8}{6.4}\selectfont \textcolor{BrickRed}{$ \blacktriangledown\,1.8\times $}}\end{tabular} & \hspace*{1pt}\begin{tabular}[t]{@{}c@{}}\rule{0pt}{12pt}10m\,57.1s\\[-.2pt]{\fontsize{5.8}{6.4}\selectfont \textcolor{BrickRed}{$ \blacktriangledown\,1.5\times $}\textsuperscript{c}}\\[-.2pt]10m\,56.8s\\[-.2pt]{\fontsize{5.8}{6.4}\selectfont \textcolor{BrickRed}{$ \blacktriangledown\,1.6\times $}}\end{tabular} & \hspace*{1pt}\begin{tabular}[t]{@{}c@{}}\rule{0pt}{12pt}9m\,46.1s\\[-.2pt]{\fontsize{5.8}{6.4}\selectfont \textcolor{ForestGreen}{$ \blacktriangle\,5.3\times $}\textsuperscript{c}}\\[-.2pt]9m\,45.6s\\[-.2pt]{\fontsize{5.8}{6.4}\selectfont \textcolor{ForestGreen}{$ \blacktriangle\,5.3\times $}}\end{tabular} & \hspace*{1pt}\begin{tabular}[t]{@{}c@{}}\rule{0pt}{12pt}4m\,52.2s\\[-.2pt]{\fontsize{5.8}{6.4}\selectfont \textcolor{BrickRed}{$ \blacktriangledown\,4.2\times $}\textsuperscript{b,c}}\\[-.2pt]4m\,51.8s\\[-.2pt]{\fontsize{5.8}{6.4}\selectfont \textcolor{BrickRed}{$ \blacktriangledown\,4.2\times $}}\end{tabular} & \hspace*{1pt}\begin{tabular}[t]{@{}c@{}}\rule{0pt}{12pt}6m\,46.8s\\[-.2pt]{\fontsize{5.8}{6.4}\selectfont \textcolor{BrickRed}{$ \blacktriangledown\,5.9\times $}\textsuperscript{b,c}}\\[-.2pt]6m\,46.3s\\[-.2pt]{\fontsize{5.8}{6.4}\selectfont \textcolor{BrickRed}{$ \blacktriangledown\,6.1\times $}}\end{tabular} & \textcolor{gray}{\ensuremath{\cdot}} & \textcolor{gray}{\ensuremath{\cdot}}\\[3pt]\hline
\end{tabular}\par}
\smallskip{\footnotesize\noindent Each cell: A, comparison, B, comparison. Batches of 1,024 queries. Five repeats unless marked s (FAST: one sample). Blue = one lowest A and one lowest B down each machine/workload column.\par
\noindent Light gray = minimap2 baseline. Pale yellow = preferred MPF / Bounded / SIMD-DP / Block configuration, not a claim of winning every column. Exact ties select the first row.\par
\noindent Ratios use minimap2 as reference, except minimap2 itself uses the fixed BioseqDB row. $\cdot$ = no approved timing. n/a = no measured counterpart.\par
\noindent Both minimap2 and BioseqDB use the matching non-external remote measurements.\par
\noindent\textsuperscript{*}External A includes process startup and output. B removes different native stages and is not an equivalent-work comparison. BLAST and LexicMap have no $w$ or MAPQ equivalent.\par
\noindent\textsuperscript{a}BLAST DNA5 uses an N score of -3, compared with -1 in the baseline.\par
\noindent\textsuperscript{b}LexicMap stores reference N as A. Queries retain their bases, and edit-script checks use that native reference encoding.\par
\noindent\textsuperscript{c}External input uses two reference records for the 4 Gbp sequence to stay within signed 32-bit coordinates.\par
\noindent\textsuperscript{d}Winnowmap uses overlapping 250 Mb records and 499M index batches after an out-of-memory failure. Results are merged globally. This can affect MAPQ.\par
\noindent\textsuperscript{e}Winnowmap uses 1999M index batches and merges the results from multiple index parts.\par
\noindent\textsuperscript{g}Retained C7G minimap2 measurement comes from the earlier campaign with documented paging pressure.\par
}
\clearpage
\noindent{\large\bfseries Table \thetable: aligner comparison, 2 Gbp, k=15, w=10\hfill 3/4}\par\smallskip
\noindent{\footnotesize Different campaigns and native timing boundaries. A/B ratios are descriptive, subject to the measurement notes.}\par\smallskip
{\fontsize{7.1}{8.1}\selectfont\setlength{\tabcolsep}{1pt}\setlength{\fboxsep}{2pt}\renewcommand{\arraystretch}{1.12}
\noindent\begin{tabular}{@{}p{15mm}p{26mm}|p{13.55mm}|p{13.55mm}|p{13.55mm}|p{13.55mm}|p{13.55mm}|p{13.55mm}|p{13.55mm}|p{13.55mm}|p{13.55mm}|p{13.55mm}|p{13.55mm}|p{13.55mm}|p{13.55mm}|p{13.55mm}|p{13.55mm}|p{13.55mm}|@{}}\hline
Aligner & Configuration & \multicolumn{8}{c|}{\textbf{Index}} & \multicolumn{8}{c|}{\textbf{No index}}\\
 &  & \multicolumn{8}{c|}{\textbf{2 Gbp}} & \multicolumn{8}{c|}{\textbf{2 Gbp}}\\
 &  & \multicolumn{8}{c|}{\textbf{k=15, w=10}} & \multicolumn{8}{c|}{\textbf{k=15, w=10}}\\
 &  & \multicolumn{4}{c|}{\textbf{DNA4}} & \multicolumn{4}{c|}{\textbf{DNA5}} & \multicolumn{4}{c|}{\textbf{DNA4}} & \multicolumn{4}{c|}{\textbf{DNA5}}\\
 &  & \multicolumn{1}{c|}{\textbf{M7A}} & \multicolumn{1}{c|}{\textbf{M7I}} & \multicolumn{1}{c|}{\textbf{M4}} & \multicolumn{1}{c|}{\textbf{C7G}} & \multicolumn{1}{c|}{\textbf{M7A}} & \multicolumn{1}{c|}{\textbf{M7I}} & \multicolumn{1}{c|}{\textbf{M4}} & \multicolumn{1}{c|}{\textbf{C7G}} & \multicolumn{1}{c|}{\textbf{M7A}} & \multicolumn{1}{c|}{\textbf{M7I}} & \multicolumn{1}{c|}{\textbf{M4}} & \multicolumn{1}{c|}{\textbf{C7G}} & \multicolumn{1}{c|}{\textbf{M7A}} & \multicolumn{1}{c|}{\textbf{M7I}} & \multicolumn{1}{c|}{\textbf{M4}} & \multicolumn{1}{c|}{\textbf{C7G}}\\
\hline
\rowcolor{BenchBaseline}[0.5pt][0.5pt]minimap2 & Retained baseline & \hspace*{1pt}\begin{tabular}[t]{@{}c@{}}\rule{0pt}{12pt}11.2s\\[-.2pt]{\fontsize{5.8}{6.4}\selectfont \textcolor{BrickRed}{$ \blacktriangledown\,3.0\times $}}\\[-.2pt]10s\\[-.2pt]{\fontsize{5.8}{6.4}\selectfont \textcolor{BrickRed}{$ \blacktriangledown\,3.4\times $}}\end{tabular} & \hspace*{1pt}\begin{tabular}[t]{@{}c@{}}\rule{0pt}{12pt}10.7s\\[-.2pt]{\fontsize{5.8}{6.4}\selectfont \textcolor{BrickRed}{$ \blacktriangledown\,3.0\times $}}\\[-.2pt]8.6s\\[-.2pt]{\fontsize{5.8}{6.4}\selectfont \textcolor{BrickRed}{$ \blacktriangledown\,3.2\times $}}\end{tabular} & \hspace*{1pt}\begin{tabular}[t]{@{}c@{}}\rule{0pt}{12pt}22.5s\\[-.2pt]{\fontsize{5.8}{6.4}\selectfont \textcolor{BrickRed}{$ \blacktriangledown\,3.3\times $}}\\[-.2pt]18.8s\\[-.2pt]{\fontsize{5.8}{6.4}\selectfont \textcolor{BrickRed}{$ \blacktriangledown\,4.2\times $}}\end{tabular} & \hspace*{1pt}\begin{tabular}[t]{@{}c@{}}\rule{0pt}{12pt}1m\,13.1s\\[-.2pt]{\fontsize{5.8}{6.4}\selectfont \textcolor{BrickRed}{$ \blacktriangledown\,8.3\times $}\textsuperscript{g}}\\[-.2pt]1m\,5.7s\\[-.2pt]{\fontsize{5.8}{6.4}\selectfont \textcolor{BrickRed}{$ \blacktriangledown\,8.3\times $}}\end{tabular} & \hspace*{1pt}\begin{tabular}[t]{@{}c@{}}\rule{0pt}{12pt}9.9s\\[-.2pt]{\fontsize{5.8}{6.4}\selectfont \textcolor{BrickRed}{$ \blacktriangledown\,21\% $}}\\[-.2pt]8.8s\\[-.2pt]{\fontsize{5.8}{6.4}\selectfont \textcolor{BrickRed}{$ \blacktriangledown\,1.5\times $}}\end{tabular} & \hspace*{1pt}\begin{tabular}[t]{@{}c@{}}\rule{0pt}{12pt}9.5s\\[-.2pt]{\fontsize{5.8}{6.4}\selectfont \textcolor{BrickRed}{$ \blacktriangledown\,1.4\times $}}\\[-.2pt]7.6s\\[-.2pt]{\fontsize{5.8}{6.4}\selectfont \textcolor{BrickRed}{$ \blacktriangledown\,1.6\times $}}\end{tabular} & \hspace*{1pt}\begin{tabular}[t]{@{}c@{}}\rule{0pt}{12pt}16.5s\\[-.2pt]{\fontsize{5.8}{6.4}\selectfont \textcolor{BrickRed}{$ \blacktriangledown\,20\% $}}\\[-.2pt]12.9s\\[-.2pt]{\fontsize{5.8}{6.4}\selectfont \textcolor{BrickRed}{$ \blacktriangledown\,1.5\times $}}\end{tabular} & \hspace*{1pt}\begin{tabular}[t]{@{}c@{}}\rule{0pt}{12pt}\colorbox{BenchBest}{\strut 28.2s}\\[-.2pt]{\fontsize{5.8}{6.4}\selectfont \textcolor{gray}{n/a}\textsuperscript{g}}\\[-.2pt]\colorbox{BenchBest}{\strut 26.2s}\\[-.2pt]{\fontsize{5.8}{6.4}\selectfont \textcolor{gray}{n/a}}\end{tabular} & \hspace*{1pt}\begin{tabular}[t]{@{}c@{}}\rule{0pt}{12pt}\colorbox{BenchBest}{\strut 53s}\\[-.2pt]{\fontsize{5.8}{6.4}\selectfont \textcolor{ForestGreen}{$ \blacktriangle\,2.2\times $}}\\[-.2pt]\colorbox{BenchBest}{\strut 51.9s}\\[-.2pt]{\fontsize{5.8}{6.4}\selectfont \textcolor{ForestGreen}{$ \blacktriangle\,2.1\times $}}\end{tabular} & \hspace*{1pt}\begin{tabular}[t]{@{}c@{}}\rule{0pt}{12pt}\colorbox{BenchBest}{\strut 56.6s}\\[-.2pt]{\fontsize{5.8}{6.4}\selectfont \textcolor{ForestGreen}{$ \blacktriangle\,2.1\times $}}\\[-.2pt]\colorbox{BenchBest}{\strut 54.6s}\\[-.2pt]{\fontsize{5.8}{6.4}\selectfont \textcolor{ForestGreen}{$ \blacktriangle\,1.9\times $}}\end{tabular} & \hspace*{1pt}\begin{tabular}[t]{@{}c@{}}\rule{0pt}{12pt}\colorbox{BenchBest}{\strut 1m\,55.5s}\\[-.2pt]{\fontsize{5.8}{6.4}\selectfont \textcolor{ForestGreen}{$ \blacktriangle\,1.9\times $}}\\[-.2pt]\colorbox{BenchBest}{\strut 1m\,51.8s}\\[-.2pt]{\fontsize{5.8}{6.4}\selectfont \textcolor{ForestGreen}{$ \blacktriangle\,1.8\times $}}\end{tabular} & \hspace*{1pt}\begin{tabular}[t]{@{}c@{}}\rule{0pt}{12pt}6m\,1.3s\\[-.2pt]{\fontsize{5.8}{6.4}\selectfont \textcolor{gray}{n/a}\textsuperscript{g}}\\[-.2pt]5m\,55.2s\\[-.2pt]{\fontsize{5.8}{6.4}\selectfont \textcolor{gray}{n/a}}\end{tabular} & \hspace*{1pt}\begin{tabular}[t]{@{}c@{}}\rule{0pt}{12pt}\colorbox{BenchBest}{\strut 48.1s}\\[-.2pt]{\fontsize{5.8}{6.4}\selectfont \textcolor{ForestGreen}{$ \blacktriangle\,7.3\times $}}\\[-.2pt]\colorbox{BenchBest}{\strut 47s}\\[-.2pt]{\fontsize{5.8}{6.4}\selectfont \textcolor{ForestGreen}{$ \blacktriangle\,2.4\times $}}\end{tabular} & \hspace*{1pt}\begin{tabular}[t]{@{}c@{}}\rule{0pt}{12pt}\colorbox{BenchBest}{\strut 44.5s}\\[-.2pt]{\fontsize{5.8}{6.4}\selectfont \textcolor{ForestGreen}{$ \blacktriangle\,6.5\times $}}\\[-.2pt]\colorbox{BenchBest}{\strut 42.6s}\\[-.2pt]{\fontsize{5.8}{6.4}\selectfont \textcolor{ForestGreen}{$ \blacktriangle\,2.7\times $}}\end{tabular} & \hspace*{1pt}\begin{tabular}[t]{@{}c@{}}\rule{0pt}{12pt}\colorbox{BenchBest}{\strut 1m\,30.2s}\\[-.2pt]{\fontsize{5.8}{6.4}\selectfont \textcolor{ForestGreen}{$ \blacktriangle\,6.4\times $}}\\[-.2pt]\colorbox{BenchBest}{\strut 1m\,26.7s}\\[-.2pt]{\fontsize{5.8}{6.4}\selectfont \textcolor{ForestGreen}{$ \blacktriangle\,2.3\times $}}\end{tabular} & \hspace*{1pt}\begin{tabular}[t]{@{}c@{}}\rule{0pt}{12pt}\colorbox{BenchBest}{\strut 2m\,46.4s}\\[-.2pt]{\fontsize{5.8}{6.4}\selectfont \textcolor{gray}{n/a}\textsuperscript{g}}\\[-.2pt]\colorbox{BenchBest}{\strut 2m\,42.7s}\\[-.2pt]{\fontsize{5.8}{6.4}\selectfont \textcolor{gray}{n/a}}\end{tabular}\\[3pt]\hline
\rowcolor{BenchPreferred}[0.5pt][0.5pt]BioseqDB & \begin{tabular}[t]{@{}l@{}}\rule{0pt}{12pt}MPF / Bounded\\SIMD-DP / Block\end{tabular} & \hspace*{1pt}\begin{tabular}[t]{@{}c@{}}\rule{0pt}{12pt}\colorbox{BenchBest}{\strut 3.7s}\\[-.2pt]{\fontsize{5.8}{6.4}\selectfont \textcolor{ForestGreen}{$ \blacktriangle\,3.0\times $}}\\[-.2pt]\colorbox{BenchBest}{\strut 3s}\\[-.2pt]{\fontsize{5.8}{6.4}\selectfont \textcolor{ForestGreen}{$ \blacktriangle\,3.4\times $}}\end{tabular} & \hspace*{1pt}\begin{tabular}[t]{@{}c@{}}\rule{0pt}{12pt}\colorbox{BenchBest}{\strut 3.5s}\\[-.2pt]{\fontsize{5.8}{6.4}\selectfont \textcolor{ForestGreen}{$ \blacktriangle\,3.0\times $}}\\[-.2pt]\colorbox{BenchBest}{\strut 2.7s}\\[-.2pt]{\fontsize{5.8}{6.4}\selectfont \textcolor{ForestGreen}{$ \blacktriangle\,3.2\times $}}\end{tabular} & \hspace*{1pt}\begin{tabular}[t]{@{}c@{}}\rule{0pt}{12pt}\colorbox{BenchBest}{\strut 6.7s}\\[-.2pt]{\fontsize{5.8}{6.4}\selectfont \textcolor{ForestGreen}{$ \blacktriangle\,3.3\times $}}\\[-.2pt]\colorbox{BenchBest}{\strut 4.5s}\\[-.2pt]{\fontsize{5.8}{6.4}\selectfont \textcolor{ForestGreen}{$ \blacktriangle\,4.2\times $}}\end{tabular} & \hspace*{1pt}\begin{tabular}[t]{@{}c@{}}\rule{0pt}{12pt}\colorbox{BenchBest}{\strut 8.8s}\\[-.2pt]{\fontsize{5.8}{6.4}\selectfont \textcolor{ForestGreen}{$ \blacktriangle\,8.3\times $}\textsuperscript{g}}\\[-.2pt]\colorbox{BenchBest}{\strut 7.9s}\\[-.2pt]{\fontsize{5.8}{6.4}\selectfont \textcolor{ForestGreen}{$ \blacktriangle\,8.3\times $}}\end{tabular} & \hspace*{1pt}\begin{tabular}[t]{@{}c@{}}\rule{0pt}{12pt}\colorbox{BenchBest}{\strut 8.2s}\\[-.2pt]{\fontsize{5.8}{6.4}\selectfont \textcolor{ForestGreen}{$ \blacktriangle\,17\% $}}\\[-.2pt]\colorbox{BenchBest}{\strut 5.8s}\\[-.2pt]{\fontsize{5.8}{6.4}\selectfont \textcolor{ForestGreen}{$ \blacktriangle\,1.5\times $}}\end{tabular} & \hspace*{1pt}\begin{tabular}[t]{@{}c@{}}\rule{0pt}{12pt}\colorbox{BenchBest}{\strut 6.6s}\\[-.2pt]{\fontsize{5.8}{6.4}\selectfont \textcolor{ForestGreen}{$ \blacktriangle\,1.4\times $}}\\[-.2pt]\colorbox{BenchBest}{\strut 4.6s}\\[-.2pt]{\fontsize{5.8}{6.4}\selectfont \textcolor{ForestGreen}{$ \blacktriangle\,1.6\times $}}\end{tabular} & \hspace*{1pt}\begin{tabular}[t]{@{}c@{}}\rule{0pt}{12pt}\colorbox{BenchBest}{\strut 13.7s}\\[-.2pt]{\fontsize{5.8}{6.4}\selectfont \textcolor{ForestGreen}{$ \blacktriangle\,17\% $}}\\[-.2pt]\colorbox{BenchBest}{\strut 8.8s}\\[-.2pt]{\fontsize{5.8}{6.4}\selectfont \textcolor{ForestGreen}{$ \blacktriangle\,1.5\times $}}\end{tabular} & \textcolor{gray}{\ensuremath{\cdot}} & \hspace*{1pt}\begin{tabular}[t]{@{}c@{}}\rule{0pt}{12pt}1m\,56.4s\\[-.2pt]{\fontsize{5.8}{6.4}\selectfont \textcolor{BrickRed}{$ \blacktriangledown\,2.2\times $}}\\[-.2pt]1m\,48.1s\\[-.2pt]{\fontsize{5.8}{6.4}\selectfont \textcolor{BrickRed}{$ \blacktriangledown\,2.1\times $}}\end{tabular} & \hspace*{1pt}\begin{tabular}[t]{@{}c@{}}\rule{0pt}{12pt}1m\,56.2s\\[-.2pt]{\fontsize{5.8}{6.4}\selectfont \textcolor{BrickRed}{$ \blacktriangledown\,2.1\times $}}\\[-.2pt]1m\,46.2s\\[-.2pt]{\fontsize{5.8}{6.4}\selectfont \textcolor{BrickRed}{$ \blacktriangledown\,1.9\times $}}\end{tabular} & \hspace*{1pt}\begin{tabular}[t]{@{}c@{}}\rule{0pt}{12pt}3m\,42.9s\\[-.2pt]{\fontsize{5.8}{6.4}\selectfont \textcolor{BrickRed}{$ \blacktriangledown\,1.9\times $}}\\[-.2pt]3m\,25s\\[-.2pt]{\fontsize{5.8}{6.4}\selectfont \textcolor{BrickRed}{$ \blacktriangledown\,1.8\times $}}\end{tabular} & \textcolor{gray}{\ensuremath{\cdot}} & \hspace*{1pt}\begin{tabular}[t]{@{}c@{}}\rule{0pt}{12pt}5m\,52.1s\\[-.2pt]{\fontsize{5.8}{6.4}\selectfont \textcolor{BrickRed}{$ \blacktriangledown\,7.3\times $}}\\[-.2pt]1m\,53.5s\\[-.2pt]{\fontsize{5.8}{6.4}\selectfont \textcolor{BrickRed}{$ \blacktriangledown\,2.4\times $}}\end{tabular} & \hspace*{1pt}\begin{tabular}[t]{@{}c@{}}\rule{0pt}{12pt}4m\,49.2s\\[-.2pt]{\fontsize{5.8}{6.4}\selectfont \textcolor{BrickRed}{$ \blacktriangledown\,6.5\times $}}\\[-.2pt]1m\,56.8s\\[-.2pt]{\fontsize{5.8}{6.4}\selectfont \textcolor{BrickRed}{$ \blacktriangledown\,2.7\times $}}\end{tabular} & \hspace*{1pt}\begin{tabular}[t]{@{}c@{}}\rule{0pt}{12pt}9m\,39.5s\\[-.2pt]{\fontsize{5.8}{6.4}\selectfont \textcolor{BrickRed}{$ \blacktriangledown\,6.4\times $}}\\[-.2pt]3m\,18.1s\\[-.2pt]{\fontsize{5.8}{6.4}\selectfont \textcolor{BrickRed}{$ \blacktriangledown\,2.3\times $}}\end{tabular} & \textcolor{gray}{\ensuremath{\cdot}}\\[3pt]\hline
blastn & Native\textsuperscript{*} & \hspace*{1pt}\begin{tabular}[t]{@{}c@{}}\rule{0pt}{12pt}5m\,36.2s\\[-.2pt]{\fontsize{5.8}{6.4}\selectfont \textcolor{BrickRed}{$ \blacktriangledown\,30.1\times $}\textsuperscript{f}}\\[-.2pt]5m\,30.7s\\[-.2pt]{\fontsize{5.8}{6.4}\selectfont \textcolor{BrickRed}{$ \blacktriangledown\,33.0\times $}}\end{tabular} & \hspace*{1pt}\begin{tabular}[t]{@{}c@{}}\rule{0pt}{12pt}8m\,14.3s\\[-.2pt]{\fontsize{5.8}{6.4}\selectfont \textcolor{BrickRed}{$ \blacktriangledown\,46.4\times $}\textsuperscript{f}}\\[-.2pt]8m\,2.4s\\[-.2pt]{\fontsize{5.8}{6.4}\selectfont \textcolor{BrickRed}{$ \blacktriangledown\,56.0\times $}}\end{tabular} & \hspace*{1pt}\begin{tabular}[t]{@{}c@{}}\rule{0pt}{12pt}10m\,51.3s\\[-.2pt]{\fontsize{5.8}{6.4}\selectfont \textcolor{BrickRed}{$ \blacktriangledown\,28.9\times $}}\\[-.2pt]10m\,35.3s\\[-.2pt]{\fontsize{5.8}{6.4}\selectfont \textcolor{BrickRed}{$ \blacktriangledown\,33.8\times $}}\end{tabular} & \hspace*{1pt}\begin{tabular}[t]{@{}c@{}}\rule{0pt}{12pt}5m\,54.4s\\[-.2pt]{\fontsize{5.8}{6.4}\selectfont \textcolor{BrickRed}{$ \blacktriangledown\,4.9\times $}\textsuperscript{f}}\\[-.2pt]5m\,48.6s\\[-.2pt]{\fontsize{5.8}{6.4}\selectfont \textcolor{BrickRed}{$ \blacktriangledown\,5.3\times $}}\end{tabular} & \hspace*{1pt}\begin{tabular}[t]{@{}c@{}}\rule{0pt}{12pt}6m\,37.9s\\[-.2pt]{\fontsize{5.8}{6.4}\selectfont \textcolor{BrickRed}{$ \blacktriangledown\,40.2\times $}\textsuperscript{a}}\\[-.2pt]4m\,57.9s\\[-.2pt]{\fontsize{5.8}{6.4}\selectfont \textcolor{BrickRed}{$ \blacktriangledown\,33.9\times $}}\end{tabular} & \hspace*{1pt}\begin{tabular}[t]{@{}c@{}}\rule{0pt}{12pt}9m\,32.3s\\[-.2pt]{\fontsize{5.8}{6.4}\selectfont \textcolor{BrickRed}{$ \blacktriangledown\,60.3\times $}\textsuperscript{a}}\\[-.2pt]6m\,18s\\[-.2pt]{\fontsize{5.8}{6.4}\selectfont \textcolor{BrickRed}{$ \blacktriangledown\,49.6\times $}}\end{tabular} & \hspace*{1pt}\begin{tabular}[t]{@{}c@{}}\rule{0pt}{12pt}13m\,27.4s\\[-.2pt]{\fontsize{5.8}{6.4}\selectfont \textcolor{BrickRed}{$ \blacktriangledown\,49.0\times $}\textsuperscript{a}}\\[-.2pt]8m\,52.1s\\[-.2pt]{\fontsize{5.8}{6.4}\selectfont \textcolor{BrickRed}{$ \blacktriangledown\,41.2\times $}}\end{tabular} & \hspace*{1pt}\begin{tabular}[t]{@{}c@{}}\rule{0pt}{12pt}6m\,53.6s\\[-.2pt]{\fontsize{5.8}{6.4}\selectfont \textcolor{BrickRed}{$ \blacktriangledown\,14.7\times $}\textsuperscript{a}}\\[-.2pt]4m\,51.8s\\[-.2pt]{\fontsize{5.8}{6.4}\selectfont \textcolor{BrickRed}{$ \blacktriangledown\,11.1\times $}}\end{tabular} & \hspace*{1pt}\begin{tabular}[t]{@{}c@{}}\rule{0pt}{12pt}5m\,39.9s\\[-.2pt]{\fontsize{5.8}{6.4}\selectfont \textcolor{BrickRed}{$ \blacktriangledown\,6.4\times $}}\\[-.2pt]5m\,34.5s\\[-.2pt]{\fontsize{5.8}{6.4}\selectfont \textcolor{BrickRed}{$ \blacktriangledown\,6.4\times $}}\end{tabular} & \hspace*{1pt}\begin{tabular}[t]{@{}c@{}}\rule{0pt}{12pt}8m\,22.6s\\[-.2pt]{\fontsize{5.8}{6.4}\selectfont \textcolor{BrickRed}{$ \blacktriangledown\,8.9\times $}}\\[-.2pt]8m\,10.6s\\[-.2pt]{\fontsize{5.8}{6.4}\selectfont \textcolor{BrickRed}{$ \blacktriangledown\,9.0\times $}}\end{tabular} & \hspace*{1pt}\begin{tabular}[t]{@{}c@{}}\rule{0pt}{12pt}11m\,6.9s\\[-.2pt]{\fontsize{5.8}{6.4}\selectfont \textcolor{BrickRed}{$ \blacktriangledown\,5.8\times $}}\\[-.2pt]10m\,50.9s\\[-.2pt]{\fontsize{5.8}{6.4}\selectfont \textcolor{BrickRed}{$ \blacktriangledown\,5.8\times $}}\end{tabular} & \hspace*{1pt}\begin{tabular}[t]{@{}c@{}}\rule{0pt}{12pt}6m\,6.4s\\[-.2pt]{\fontsize{5.8}{6.4}\selectfont \textcolor{gray}{$ \blacktriangleright\,1\% $}}\\[-.2pt]6m\,0.6s\\[-.2pt]{\fontsize{5.8}{6.4}\selectfont \textcolor{gray}{$ \blacktriangleright\,2\% $}}\end{tabular} & \hspace*{1pt}\begin{tabular}[t]{@{}c@{}}\rule{0pt}{12pt}6m\,51.8s\\[-.2pt]{\fontsize{5.8}{6.4}\selectfont \textcolor{BrickRed}{$ \blacktriangledown\,8.6\times $}\textsuperscript{a}}\\[-.2pt]5m\,12.2s\\[-.2pt]{\fontsize{5.8}{6.4}\selectfont \textcolor{BrickRed}{$ \blacktriangledown\,6.6\times $}}\end{tabular} & \hspace*{1pt}\begin{tabular}[t]{@{}c@{}}\rule{0pt}{12pt}9m\,48.4s\\[-.2pt]{\fontsize{5.8}{6.4}\selectfont \textcolor{BrickRed}{$ \blacktriangledown\,13.2\times $}\textsuperscript{a}}\\[-.2pt]6m\,35.3s\\[-.2pt]{\fontsize{5.8}{6.4}\selectfont \textcolor{BrickRed}{$ \blacktriangledown\,9.3\times $}}\end{tabular} & \hspace*{1pt}\begin{tabular}[t]{@{}c@{}}\rule{0pt}{12pt}13m\,55.3s\\[-.2pt]{\fontsize{5.8}{6.4}\selectfont \textcolor{BrickRed}{$ \blacktriangledown\,9.3\times $}\textsuperscript{a}}\\[-.2pt]9m\,23.3s\\[-.2pt]{\fontsize{5.8}{6.4}\selectfont \textcolor{BrickRed}{$ \blacktriangledown\,6.5\times $}}\end{tabular} & \hspace*{1pt}\begin{tabular}[t]{@{}c@{}}\rule{0pt}{12pt}7m\,14.5s\\[-.2pt]{\fontsize{5.8}{6.4}\selectfont \textcolor{BrickRed}{$ \blacktriangledown\,2.6\times $}\textsuperscript{a}}\\[-.2pt]5m\,12.1s\\[-.2pt]{\fontsize{5.8}{6.4}\selectfont \textcolor{BrickRed}{$ \blacktriangledown\,1.9\times $}}\end{tabular}\\[3pt]\hline
winnowmap & Native\textsuperscript{*} & \hspace*{1pt}\begin{tabular}[t]{@{}c@{}}\rule{0pt}{12pt}12.9s\\[-.2pt]{\fontsize{5.8}{6.4}\selectfont \textcolor{BrickRed}{$ \blacktriangledown\,15\% $}}\\[-.2pt]11.1s\\[-.2pt]{\fontsize{5.8}{6.4}\selectfont \textcolor{BrickRed}{$ \blacktriangledown\,11\% $}}\end{tabular} & \hspace*{1pt}\begin{tabular}[t]{@{}c@{}}\rule{0pt}{12pt}13.5s\\[-.2pt]{\fontsize{5.8}{6.4}\selectfont \textcolor{BrickRed}{$ \blacktriangledown\,27\% $}}\\[-.2pt]10.3s\\[-.2pt]{\fontsize{5.8}{6.4}\selectfont \textcolor{BrickRed}{$ \blacktriangledown\,20\% $}}\end{tabular} & \hspace*{1pt}\begin{tabular}[t]{@{}c@{}}\rule{0pt}{12pt}23.2s\\[-.2pt]{\fontsize{5.8}{6.4}\selectfont \textcolor{gray}{$ \blacktriangleright\,3\% $}}\\[-.2pt]17.1s\\[-.2pt]{\fontsize{5.8}{6.4}\selectfont \textcolor{ForestGreen}{$ \blacktriangle\,9\% $}}\end{tabular} & \hspace*{1pt}\begin{tabular}[t]{@{}c@{}}\rule{0pt}{12pt}44.3s\\[-.2pt]{\fontsize{5.8}{6.4}\selectfont \textcolor{ForestGreen}{$ \blacktriangle\,1.7\times $}\textsuperscript{d}}\\[-.2pt]37.8s\\[-.2pt]{\fontsize{5.8}{6.4}\selectfont \textcolor{ForestGreen}{$ \blacktriangle\,1.7\times $}}\end{tabular} & \hspace*{1pt}\begin{tabular}[t]{@{}c@{}}\rule{0pt}{12pt}12.3s\\[-.2pt]{\fontsize{5.8}{6.4}\selectfont \textcolor{BrickRed}{$ \blacktriangledown\,24\% $}}\\[-.2pt]10.6s\\[-.2pt]{\fontsize{5.8}{6.4}\selectfont \textcolor{BrickRed}{$ \blacktriangledown\,21\% $}}\end{tabular} & \hspace*{1pt}\begin{tabular}[t]{@{}c@{}}\rule{0pt}{12pt}13s\\[-.2pt]{\fontsize{5.8}{6.4}\selectfont \textcolor{BrickRed}{$ \blacktriangledown\,1.4\times $}}\\[-.2pt]10s\\[-.2pt]{\fontsize{5.8}{6.4}\selectfont \textcolor{BrickRed}{$ \blacktriangledown\,1.3\times $}}\end{tabular} & \hspace*{1pt}\begin{tabular}[t]{@{}c@{}}\rule{0pt}{12pt}23.8s\\[-.2pt]{\fontsize{5.8}{6.4}\selectfont \textcolor{BrickRed}{$ \blacktriangledown\,1.4\times $}}\\[-.2pt]17.6s\\[-.2pt]{\fontsize{5.8}{6.4}\selectfont \textcolor{BrickRed}{$ \blacktriangledown\,1.4\times $}}\end{tabular} & \hspace*{1pt}\begin{tabular}[t]{@{}c@{}}\rule{0pt}{12pt}37.3s\\[-.2pt]{\fontsize{5.8}{6.4}\selectfont \textcolor{BrickRed}{$ \blacktriangledown\,1.3\times $}\textsuperscript{d}}\\[-.2pt]31.4s\\[-.2pt]{\fontsize{5.8}{6.4}\selectfont \textcolor{BrickRed}{$ \blacktriangledown\,20\% $}}\end{tabular} & \hspace*{1pt}\begin{tabular}[t]{@{}c@{}}\rule{0pt}{12pt}3m\,33.2s\\[-.2pt]{\fontsize{5.8}{6.4}\selectfont \textcolor{BrickRed}{$ \blacktriangledown\,4.0\times $}}\\[-.2pt]3m\,31.5s\\[-.2pt]{\fontsize{5.8}{6.4}\selectfont \textcolor{BrickRed}{$ \blacktriangledown\,4.1\times $}}\end{tabular} & \hspace*{1pt}\begin{tabular}[t]{@{}c@{}}\rule{0pt}{12pt}3m\,22.6s\\[-.2pt]{\fontsize{5.8}{6.4}\selectfont \textcolor{BrickRed}{$ \blacktriangledown\,3.6\times $}}\\[-.2pt]3m\,19.3s\\[-.2pt]{\fontsize{5.8}{6.4}\selectfont \textcolor{BrickRed}{$ \blacktriangledown\,3.7\times $}}\end{tabular} & \hspace*{1pt}\begin{tabular}[t]{@{}c@{}}\rule{0pt}{12pt}6m\,21.3s\\[-.2pt]{\fontsize{5.8}{6.4}\selectfont \textcolor{BrickRed}{$ \blacktriangledown\,3.3\times $}}\\[-.2pt]6m\,15s\\[-.2pt]{\fontsize{5.8}{6.4}\selectfont \textcolor{BrickRed}{$ \blacktriangledown\,3.4\times $}}\end{tabular} & \hspace*{1pt}\begin{tabular}[t]{@{}c@{}}\rule{0pt}{12pt}\colorbox{BenchBest}{\strut 3m\,56.6s}\\[-.2pt]{\fontsize{5.8}{6.4}\selectfont \textcolor{ForestGreen}{$ \blacktriangle\,1.5\times $}\textsuperscript{d}}\\[-.2pt]\colorbox{BenchBest}{\strut 3m\,50.3s}\\[-.2pt]{\fontsize{5.8}{6.4}\selectfont \textcolor{ForestGreen}{$ \blacktriangle\,1.5\times $}}\end{tabular} & \hspace*{1pt}\begin{tabular}[t]{@{}c@{}}\rule{0pt}{12pt}3m\,7.4s\\[-.2pt]{\fontsize{5.8}{6.4}\selectfont \textcolor{BrickRed}{$ \blacktriangledown\,3.9\times $}}\\[-.2pt]3m\,5.8s\\[-.2pt]{\fontsize{5.8}{6.4}\selectfont \textcolor{BrickRed}{$ \blacktriangledown\,4.0\times $}}\end{tabular} & \hspace*{1pt}\begin{tabular}[t]{@{}c@{}}\rule{0pt}{12pt}3m\,2.2s\\[-.2pt]{\fontsize{5.8}{6.4}\selectfont \textcolor{BrickRed}{$ \blacktriangledown\,4.1\times $}}\\[-.2pt]2m\,59.3s\\[-.2pt]{\fontsize{5.8}{6.4}\selectfont \textcolor{BrickRed}{$ \blacktriangledown\,4.2\times $}}\end{tabular} & \hspace*{1pt}\begin{tabular}[t]{@{}c@{}}\rule{0pt}{12pt}6m\,35.9s\\[-.2pt]{\fontsize{5.8}{6.4}\selectfont \textcolor{BrickRed}{$ \blacktriangledown\,4.4\times $}}\\[-.2pt]6m\,29.8s\\[-.2pt]{\fontsize{5.8}{6.4}\selectfont \textcolor{BrickRed}{$ \blacktriangledown\,4.5\times $}}\end{tabular} & \hspace*{1pt}\begin{tabular}[t]{@{}c@{}}\rule{0pt}{12pt}3m\,33.9s\\[-.2pt]{\fontsize{5.8}{6.4}\selectfont \textcolor{BrickRed}{$ \blacktriangledown\,29\% $}\textsuperscript{d}}\\[-.2pt]3m\,28.1s\\[-.2pt]{\fontsize{5.8}{6.4}\selectfont \textcolor{BrickRed}{$ \blacktriangledown\,28\% $}}\end{tabular}\\[3pt]\hline
lexicmap & Native\textsuperscript{*} & \hspace*{1pt}\begin{tabular}[t]{@{}c@{}}\rule{0pt}{12pt}2m\,14.1s\\[-.2pt]{\fontsize{5.8}{6.4}\selectfont \textcolor{BrickRed}{$ \blacktriangledown\,12.0\times $}}\\[-.2pt]2m\,14s\\[-.2pt]{\fontsize{5.8}{6.4}\selectfont \textcolor{BrickRed}{$ \blacktriangledown\,13.4\times $}}\end{tabular} & \hspace*{1pt}\begin{tabular}[t]{@{}c@{}}\rule{0pt}{12pt}3m\,26.7s\\[-.2pt]{\fontsize{5.8}{6.4}\selectfont \textcolor{BrickRed}{$ \blacktriangledown\,19.4\times $}}\\[-.2pt]3m\,26.6s\\[-.2pt]{\fontsize{5.8}{6.4}\selectfont \textcolor{BrickRed}{$ \blacktriangledown\,24.0\times $}}\end{tabular} & \hspace*{1pt}\begin{tabular}[t]{@{}c@{}}\rule{0pt}{12pt}4m\,40.1s\\[-.2pt]{\fontsize{5.8}{6.4}\selectfont \textcolor{BrickRed}{$ \blacktriangledown\,12.4\times $}}\\[-.2pt]4m\,39.9s\\[-.2pt]{\fontsize{5.8}{6.4}\selectfont \textcolor{BrickRed}{$ \blacktriangledown\,14.9\times $}}\end{tabular} & \hspace*{1pt}\begin{tabular}[t]{@{}c@{}}\rule{0pt}{12pt}3m\,24.3s\\[-.2pt]{\fontsize{5.8}{6.4}\selectfont \textcolor{BrickRed}{$ \blacktriangledown\,2.8\times $}}\\[-.2pt]3m\,24.2s\\[-.2pt]{\fontsize{5.8}{6.4}\selectfont \textcolor{BrickRed}{$ \blacktriangledown\,3.1\times $}}\end{tabular} & \hspace*{1pt}\begin{tabular}[t]{@{}c@{}}\rule{0pt}{12pt}2m\,12.9s\\[-.2pt]{\fontsize{5.8}{6.4}\selectfont \textcolor{BrickRed}{$ \blacktriangledown\,13.4\times $}\textsuperscript{b}}\\[-.2pt]2m\,12.5s\\[-.2pt]{\fontsize{5.8}{6.4}\selectfont \textcolor{BrickRed}{$ \blacktriangledown\,15.1\times $}}\end{tabular} & \hspace*{1pt}\begin{tabular}[t]{@{}c@{}}\rule{0pt}{12pt}3m\,24.1s\\[-.2pt]{\fontsize{5.8}{6.4}\selectfont \textcolor{BrickRed}{$ \blacktriangledown\,21.5\times $}\textsuperscript{b}}\\[-.2pt]3m\,23.6s\\[-.2pt]{\fontsize{5.8}{6.4}\selectfont \textcolor{BrickRed}{$ \blacktriangledown\,26.7\times $}}\end{tabular} & \hspace*{1pt}\begin{tabular}[t]{@{}c@{}}\rule{0pt}{12pt}4m\,56.4s\\[-.2pt]{\fontsize{5.8}{6.4}\selectfont \textcolor{BrickRed}{$ \blacktriangledown\,18.0\times $}\textsuperscript{b}}\\[-.2pt]4m\,55.6s\\[-.2pt]{\fontsize{5.8}{6.4}\selectfont \textcolor{BrickRed}{$ \blacktriangledown\,22.9\times $}}\end{tabular} & \hspace*{1pt}\begin{tabular}[t]{@{}c@{}}\rule{0pt}{12pt}3m\,22.8s\\[-.2pt]{\fontsize{5.8}{6.4}\selectfont \textcolor{BrickRed}{$ \blacktriangledown\,7.2\times $}\textsuperscript{b}}\\[-.2pt]3m\,22.5s\\[-.2pt]{\fontsize{5.8}{6.4}\selectfont \textcolor{BrickRed}{$ \blacktriangledown\,7.7\times $}}\end{tabular} & \hspace*{1pt}\begin{tabular}[t]{@{}c@{}}\rule{0pt}{12pt}3m\,18.9s\\[-.2pt]{\fontsize{5.8}{6.4}\selectfont \textcolor{BrickRed}{$ \blacktriangledown\,3.8\times $}}\\[-.2pt]3m\,18.7s\\[-.2pt]{\fontsize{5.8}{6.4}\selectfont \textcolor{BrickRed}{$ \blacktriangledown\,3.8\times $}}\end{tabular} & \hspace*{1pt}\begin{tabular}[t]{@{}c@{}}\rule{0pt}{12pt}5m\,10.5s\\[-.2pt]{\fontsize{5.8}{6.4}\selectfont \textcolor{BrickRed}{$ \blacktriangledown\,5.5\times $}}\\[-.2pt]5m\,10.3s\\[-.2pt]{\fontsize{5.8}{6.4}\selectfont \textcolor{BrickRed}{$ \blacktriangledown\,5.7\times $}}\end{tabular} & \hspace*{1pt}\begin{tabular}[t]{@{}c@{}}\rule{0pt}{12pt}8m\,19.4s\\[-.2pt]{\fontsize{5.8}{6.4}\selectfont \textcolor{BrickRed}{$ \blacktriangledown\,4.3\times $}}\\[-.2pt]8m\,19.2s\\[-.2pt]{\fontsize{5.8}{6.4}\selectfont \textcolor{BrickRed}{$ \blacktriangledown\,4.5\times $}}\end{tabular} & \hspace*{1pt}\begin{tabular}[t]{@{}c@{}}\rule{0pt}{12pt}5m\,43.1s\\[-.2pt]{\fontsize{5.8}{6.4}\selectfont \textcolor{ForestGreen}{$ \blacktriangle\,5\% $}}\\[-.2pt]5m\,43s\\[-.2pt]{\fontsize{5.8}{6.4}\selectfont \textcolor{gray}{$ \blacktriangleright\,3\% $}}\end{tabular} & \hspace*{1pt}\begin{tabular}[t]{@{}c@{}}\rule{0pt}{12pt}3m\,16.5s\\[-.2pt]{\fontsize{5.8}{6.4}\selectfont \textcolor{BrickRed}{$ \blacktriangledown\,4.1\times $}\textsuperscript{b}}\\[-.2pt]3m\,16.1s\\[-.2pt]{\fontsize{5.8}{6.4}\selectfont \textcolor{BrickRed}{$ \blacktriangledown\,4.2\times $}}\end{tabular} & \hspace*{1pt}\begin{tabular}[t]{@{}c@{}}\rule{0pt}{12pt}5m\,5s\\[-.2pt]{\fontsize{5.8}{6.4}\selectfont \textcolor{BrickRed}{$ \blacktriangledown\,6.9\times $}\textsuperscript{b}}\\[-.2pt]5m\,4.5s\\[-.2pt]{\fontsize{5.8}{6.4}\selectfont \textcolor{BrickRed}{$ \blacktriangledown\,7.2\times $}}\end{tabular} & \hspace*{1pt}\begin{tabular}[t]{@{}c@{}}\rule{0pt}{12pt}7m\,40.8s\\[-.2pt]{\fontsize{5.8}{6.4}\selectfont \textcolor{BrickRed}{$ \blacktriangledown\,5.1\times $}\textsuperscript{b}}\\[-.2pt]7m\,39.9s\\[-.2pt]{\fontsize{5.8}{6.4}\selectfont \textcolor{BrickRed}{$ \blacktriangledown\,5.3\times $}}\end{tabular} & \hspace*{1pt}\begin{tabular}[t]{@{}c@{}}\rule{0pt}{12pt}5m\,52.3s\\[-.2pt]{\fontsize{5.8}{6.4}\selectfont \textcolor{BrickRed}{$ \blacktriangledown\,2.1\times $}\textsuperscript{b}}\\[-.2pt]5m\,51.9s\\[-.2pt]{\fontsize{5.8}{6.4}\selectfont \textcolor{BrickRed}{$ \blacktriangledown\,2.2\times $}}\end{tabular}\\[3pt]\hline
\end{tabular}\par}
\smallskip{\footnotesize\noindent Each cell: A, comparison, B, comparison. Batches of 1,024 queries. Five repeats unless marked s (FAST: one sample). Blue = one lowest A and one lowest B down each machine/workload column.\par
\noindent Light gray = minimap2 baseline. Pale yellow = preferred MPF / Bounded / SIMD-DP / Block configuration, not a claim of winning every column. Exact ties select the first row.\par
\noindent Ratios use minimap2 as reference, except minimap2 itself uses the fixed BioseqDB row. $\cdot$ = no approved timing. n/a = no measured counterpart.\par
\noindent Both minimap2 and BioseqDB use the matching non-external remote measurements.\par
\noindent\textsuperscript{*}External A includes process startup and output. B removes different native stages and is not an equivalent-work comparison. BLAST and LexicMap have no $w$ or MAPQ equivalent.\par
\noindent\textsuperscript{a}BLAST DNA5 uses an N score of -3, compared with -1 in the baseline.\par
\noindent\textsuperscript{b}LexicMap stores reference N as A. Queries retain their bases, and edit-script checks use that native reference encoding.\par
\noindent\textsuperscript{d}Winnowmap uses overlapping 250 Mb records and 499M index batches after an out-of-memory failure. Results are merged globally. This can affect MAPQ.\par
\noindent\textsuperscript{f}Swap configuration was not recorded before this external measurement. Later host snapshots do not establish its earlier swap state.\par
\noindent\textsuperscript{g}Retained C7G minimap2 measurement comes from the earlier campaign with documented paging pressure.\par
}
\clearpage
\noindent{\large\bfseries Table \thetable: aligner comparison, 2 Gbp, k=20, w=15\hfill 4/4}\par\smallskip
\noindent{\footnotesize Different campaigns and native timing boundaries. A/B ratios are descriptive, subject to the measurement notes.}\par\smallskip
{\fontsize{7.1}{8.1}\selectfont\setlength{\tabcolsep}{1pt}\setlength{\fboxsep}{2pt}\renewcommand{\arraystretch}{1.12}
\noindent\begin{tabular}{@{}p{15mm}p{26mm}|p{13.55mm}|p{13.55mm}|p{13.55mm}|p{13.55mm}|p{13.55mm}|p{13.55mm}|p{13.55mm}|p{13.55mm}|p{13.55mm}|p{13.55mm}|p{13.55mm}|p{13.55mm}|p{13.55mm}|p{13.55mm}|p{13.55mm}|p{13.55mm}|@{}}\hline
Aligner & Configuration & \multicolumn{8}{c|}{\textbf{Index}} & \multicolumn{8}{c|}{\textbf{No index}}\\
 &  & \multicolumn{8}{c|}{\textbf{2 Gbp}} & \multicolumn{8}{c|}{\textbf{2 Gbp}}\\
 &  & \multicolumn{8}{c|}{\textbf{k=20, w=15}} & \multicolumn{8}{c|}{\textbf{k=20, w=15}}\\
 &  & \multicolumn{4}{c|}{\textbf{DNA4}} & \multicolumn{4}{c|}{\textbf{DNA5}} & \multicolumn{4}{c|}{\textbf{DNA4}} & \multicolumn{4}{c|}{\textbf{DNA5}}\\
 &  & \multicolumn{1}{c|}{\textbf{M7A}} & \multicolumn{1}{c|}{\textbf{M7I}} & \multicolumn{1}{c|}{\textbf{M4}} & \multicolumn{1}{c|}{\textbf{C7G}} & \multicolumn{1}{c|}{\textbf{M7A}} & \multicolumn{1}{c|}{\textbf{M7I}} & \multicolumn{1}{c|}{\textbf{M4}} & \multicolumn{1}{c|}{\textbf{C7G}} & \multicolumn{1}{c|}{\textbf{M7A}} & \multicolumn{1}{c|}{\textbf{M7I}} & \multicolumn{1}{c|}{\textbf{M4}} & \multicolumn{1}{c|}{\textbf{C7G}} & \multicolumn{1}{c|}{\textbf{M7A}} & \multicolumn{1}{c|}{\textbf{M7I}} & \multicolumn{1}{c|}{\textbf{M4}} & \multicolumn{1}{c|}{\textbf{C7G}}\\
\hline
\rowcolor{BenchBaseline}[0.5pt][0.5pt]minimap2 & Retained baseline & \hspace*{1pt}\begin{tabular}[t]{@{}c@{}}\rule{0pt}{12pt}19.6s\\[-.2pt]{\fontsize{5.8}{6.4}\selectfont \textcolor{BrickRed}{$ \blacktriangledown\,6.5\times $}}\\[-.2pt]18.6s\\[-.2pt]{\fontsize{5.8}{6.4}\selectfont \textcolor{BrickRed}{$ \blacktriangledown\,8.1\times $}}\end{tabular} & \hspace*{1pt}\begin{tabular}[t]{@{}c@{}}\rule{0pt}{12pt}19.4s\\[-.2pt]{\fontsize{5.8}{6.4}\selectfont \textcolor{BrickRed}{$ \blacktriangledown\,7.2\times $}}\\[-.2pt]17.5s\\[-.2pt]{\fontsize{5.8}{6.4}\selectfont \textcolor{BrickRed}{$ \blacktriangledown\,8.9\times $}}\end{tabular} & \hspace*{1pt}\begin{tabular}[t]{@{}c@{}}\rule{0pt}{12pt}1m\,12.7s\\[-.2pt]{\fontsize{5.8}{6.4}\selectfont \textcolor{BrickRed}{$ \blacktriangledown\,13.6\times $}}\\[-.2pt]1m\,8.1s\\[-.2pt]{\fontsize{5.8}{6.4}\selectfont \textcolor{BrickRed}{$ \blacktriangledown\,20.2\times $}}\end{tabular} & \hspace*{1pt}\begin{tabular}[t]{@{}c@{}}\rule{0pt}{12pt}7m\,25.4s\\[-.2pt]{\fontsize{5.8}{6.4}\selectfont \textcolor{BrickRed}{$ \blacktriangledown\,99.2\times $}\textsuperscript{g}}\\[-.2pt]6m\,50.8s\\[-.2pt]{\fontsize{5.8}{6.4}\selectfont \textcolor{BrickRed}{$ \blacktriangledown\,112.7\times $}}\end{tabular} & \hspace*{1pt}\begin{tabular}[t]{@{}c@{}}\rule{0pt}{12pt}9.3s\\[-.2pt]{\fontsize{5.8}{6.4}\selectfont \textcolor{BrickRed}{$ \blacktriangledown\,24\% $}}\\[-.2pt]8.4s\\[-.2pt]{\fontsize{5.8}{6.4}\selectfont \textcolor{BrickRed}{$ \blacktriangledown\,1.6\times $}}\end{tabular} & \hspace*{1pt}\begin{tabular}[t]{@{}c@{}}\rule{0pt}{12pt}8.7s\\[-.2pt]{\fontsize{5.8}{6.4}\selectfont \textcolor{BrickRed}{$ \blacktriangledown\,1.5\times $}}\\[-.2pt]7s\\[-.2pt]{\fontsize{5.8}{6.4}\selectfont \textcolor{BrickRed}{$ \blacktriangledown\,1.7\times $}}\end{tabular} & \hspace*{1pt}\begin{tabular}[t]{@{}c@{}}\rule{0pt}{12pt}15.8s\\[-.2pt]{\fontsize{5.8}{6.4}\selectfont \textcolor{BrickRed}{$ \blacktriangledown\,26\% $}}\\[-.2pt]12.7s\\[-.2pt]{\fontsize{5.8}{6.4}\selectfont \textcolor{BrickRed}{$ \blacktriangledown\,1.6\times $}}\end{tabular} & \hspace*{1pt}\begin{tabular}[t]{@{}c@{}}\rule{0pt}{12pt}18.5s\\[-.2pt]{\fontsize{5.8}{6.4}\selectfont \textcolor{BrickRed}{$ \blacktriangledown\,1.4\times $}\textsuperscript{g}}\\[-.2pt]17s\\[-.2pt]{\fontsize{5.8}{6.4}\selectfont \textcolor{BrickRed}{$ \blacktriangledown\,1.6\times $}}\end{tabular} & \hspace*{1pt}\begin{tabular}[t]{@{}c@{}}\rule{0pt}{12pt}\colorbox{BenchBest}{\strut 42.4s}\\[-.2pt]{\fontsize{5.8}{6.4}\selectfont \textcolor{ForestGreen}{$ \blacktriangle\,2.0\times $}}\\[-.2pt]\colorbox{BenchBest}{\strut 41.4s}\\[-.2pt]{\fontsize{5.8}{6.4}\selectfont \textcolor{ForestGreen}{$ \blacktriangle\,1.9\times $}}\end{tabular} & \hspace*{1pt}\begin{tabular}[t]{@{}c@{}}\rule{0pt}{12pt}\colorbox{BenchBest}{\strut 42.9s}\\[-.2pt]{\fontsize{5.8}{6.4}\selectfont \textcolor{ForestGreen}{$ \blacktriangle\,2.0\times $}}\\[-.2pt]\colorbox{BenchBest}{\strut 41.1s}\\[-.2pt]{\fontsize{5.8}{6.4}\selectfont \textcolor{ForestGreen}{$ \blacktriangle\,1.8\times $}}\end{tabular} & \hspace*{1pt}\begin{tabular}[t]{@{}c@{}}\rule{0pt}{12pt}\colorbox{BenchBest}{\strut 1m\,26.9s}\\[-.2pt]{\fontsize{5.8}{6.4}\selectfont \textcolor{ForestGreen}{$ \blacktriangle\,1.9\times $}}\\[-.2pt]\colorbox{BenchBest}{\strut 1m\,23.7s}\\[-.2pt]{\fontsize{5.8}{6.4}\selectfont \textcolor{ForestGreen}{$ \blacktriangle\,1.7\times $}}\end{tabular} & \hspace*{1pt}\begin{tabular}[t]{@{}c@{}}\rule{0pt}{12pt}4m\,51.6s\\[-.2pt]{\fontsize{5.8}{6.4}\selectfont \textcolor{gray}{n/a}\textsuperscript{g}}\\[-.2pt]4m\,46.6s\\[-.2pt]{\fontsize{5.8}{6.4}\selectfont \textcolor{gray}{n/a}}\end{tabular} & \hspace*{1pt}\begin{tabular}[t]{@{}c@{}}\rule{0pt}{12pt}\colorbox{BenchBest}{\strut 35s}\\[-.2pt]{\fontsize{5.8}{6.4}\selectfont \textcolor{ForestGreen}{$ \blacktriangle\,8.6\times $}}\\[-.2pt]\colorbox{BenchBest}{\strut 34s}\\[-.2pt]{\fontsize{5.8}{6.4}\selectfont \textcolor{ForestGreen}{$ \blacktriangle\,2.5\times $}}\end{tabular} & \hspace*{1pt}\begin{tabular}[t]{@{}c@{}}\rule{0pt}{12pt}\colorbox{BenchBest}{\strut 35.9s}\\[-.2pt]{\fontsize{5.8}{6.4}\selectfont \textcolor{ForestGreen}{$ \blacktriangle\,6.7\times $}}\\[-.2pt]\colorbox{BenchBest}{\strut 34.2s}\\[-.2pt]{\fontsize{5.8}{6.4}\selectfont \textcolor{ForestGreen}{$ \blacktriangle\,2.5\times $}}\end{tabular} & \hspace*{1pt}\begin{tabular}[t]{@{}c@{}}\rule{0pt}{12pt}\colorbox{BenchBest}{\strut 1m\,14.4s}\\[-.2pt]{\fontsize{5.8}{6.4}\selectfont \textcolor{ForestGreen}{$ \blacktriangle\,6.9\times $}}\\[-.2pt]\colorbox{BenchBest}{\strut 1m\,11.3s}\\[-.2pt]{\fontsize{5.8}{6.4}\selectfont \textcolor{ForestGreen}{$ \blacktriangle\,2.2\times $}}\end{tabular} & \hspace*{1pt}\begin{tabular}[t]{@{}c@{}}\rule{0pt}{12pt}\colorbox{BenchBest}{\strut 1m\,0.1s}\\[-.2pt]{\fontsize{5.8}{6.4}\selectfont \textcolor{gray}{n/a}\textsuperscript{g}}\\[-.2pt]\colorbox{BenchBest}{\strut 58.6s}\\[-.2pt]{\fontsize{5.8}{6.4}\selectfont \textcolor{gray}{n/a}}\end{tabular}\\[3pt]\hline
\rowcolor{BenchPreferred}[0.5pt][0.5pt]BioseqDB & \begin{tabular}[t]{@{}l@{}}\rule{0pt}{12pt}MPF / Bounded\\SIMD-DP / Block\end{tabular} & \hspace*{1pt}\begin{tabular}[t]{@{}c@{}}\rule{0pt}{12pt}\colorbox{BenchBest}{\strut 3s}\\[-.2pt]{\fontsize{5.8}{6.4}\selectfont \textcolor{ForestGreen}{$ \blacktriangle\,6.5\times $}}\\[-.2pt]\colorbox{BenchBest}{\strut 2.3s}\\[-.2pt]{\fontsize{5.8}{6.4}\selectfont \textcolor{ForestGreen}{$ \blacktriangle\,8.1\times $}}\end{tabular} & \hspace*{1pt}\begin{tabular}[t]{@{}c@{}}\rule{0pt}{12pt}\colorbox{BenchBest}{\strut 2.7s}\\[-.2pt]{\fontsize{5.8}{6.4}\selectfont \textcolor{ForestGreen}{$ \blacktriangle\,7.2\times $}}\\[-.2pt]\colorbox{BenchBest}{\strut 2s}\\[-.2pt]{\fontsize{5.8}{6.4}\selectfont \textcolor{ForestGreen}{$ \blacktriangle\,8.9\times $}}\end{tabular} & \hspace*{1pt}\begin{tabular}[t]{@{}c@{}}\rule{0pt}{12pt}\colorbox{BenchBest}{\strut 5.4s}\\[-.2pt]{\fontsize{5.8}{6.4}\selectfont \textcolor{ForestGreen}{$ \blacktriangle\,13.6\times $}}\\[-.2pt]\colorbox{BenchBest}{\strut 3.4s}\\[-.2pt]{\fontsize{5.8}{6.4}\selectfont \textcolor{ForestGreen}{$ \blacktriangle\,20.2\times $}}\end{tabular} & \hspace*{1pt}\begin{tabular}[t]{@{}c@{}}\rule{0pt}{12pt}\colorbox{BenchBest}{\strut 4.5s}\\[-.2pt]{\fontsize{5.8}{6.4}\selectfont \textcolor{ForestGreen}{$ \blacktriangle\,99.2\times $}\textsuperscript{g}}\\[-.2pt]\colorbox{BenchBest}{\strut 3.6s}\\[-.2pt]{\fontsize{5.8}{6.4}\selectfont \textcolor{ForestGreen}{$ \blacktriangle\,112.7\times $}}\end{tabular} & \hspace*{1pt}\begin{tabular}[t]{@{}c@{}}\rule{0pt}{12pt}\colorbox{BenchBest}{\strut 7.5s}\\[-.2pt]{\fontsize{5.8}{6.4}\selectfont \textcolor{ForestGreen}{$ \blacktriangle\,19\% $}}\\[-.2pt]\colorbox{BenchBest}{\strut 5.2s}\\[-.2pt]{\fontsize{5.8}{6.4}\selectfont \textcolor{ForestGreen}{$ \blacktriangle\,1.6\times $}}\end{tabular} & \hspace*{1pt}\begin{tabular}[t]{@{}c@{}}\rule{0pt}{12pt}\colorbox{BenchBest}{\strut 5.8s}\\[-.2pt]{\fontsize{5.8}{6.4}\selectfont \textcolor{ForestGreen}{$ \blacktriangle\,1.5\times $}}\\[-.2pt]\colorbox{BenchBest}{\strut 4s}\\[-.2pt]{\fontsize{5.8}{6.4}\selectfont \textcolor{ForestGreen}{$ \blacktriangle\,1.7\times $}}\end{tabular} & \hspace*{1pt}\begin{tabular}[t]{@{}c@{}}\rule{0pt}{12pt}\colorbox{BenchBest}{\strut 12.6s}\\[-.2pt]{\fontsize{5.8}{6.4}\selectfont \textcolor{ForestGreen}{$ \blacktriangle\,21\% $}}\\[-.2pt]\colorbox{BenchBest}{\strut 7.9s}\\[-.2pt]{\fontsize{5.8}{6.4}\selectfont \textcolor{ForestGreen}{$ \blacktriangle\,1.6\times $}}\end{tabular} & \hspace*{1pt}\begin{tabular}[t]{@{}c@{}}\rule{0pt}{12pt}\colorbox{BenchBest}{\strut 12.9s}\\[-.2pt]{\fontsize{5.8}{6.4}\selectfont \textcolor{ForestGreen}{$ \blacktriangle\,1.4\times $}}\\[-.2pt]\colorbox{BenchBest}{\strut 10.6s}\\[-.2pt]{\fontsize{5.8}{6.4}\selectfont \textcolor{ForestGreen}{$ \blacktriangle\,1.6\times $}}\end{tabular} & \hspace*{1pt}\begin{tabular}[t]{@{}c@{}}\rule{0pt}{12pt}1m\,24.8s\\[-.2pt]{\fontsize{5.8}{6.4}\selectfont \textcolor{BrickRed}{$ \blacktriangledown\,2.0\times $}}\\[-.2pt]1m\,16.9s\\[-.2pt]{\fontsize{5.8}{6.4}\selectfont \textcolor{BrickRed}{$ \blacktriangledown\,1.9\times $}}\end{tabular} & \hspace*{1pt}\begin{tabular}[t]{@{}c@{}}\rule{0pt}{12pt}1m\,25.4s\\[-.2pt]{\fontsize{5.8}{6.4}\selectfont \textcolor{BrickRed}{$ \blacktriangledown\,2.0\times $}}\\[-.2pt]1m\,15.7s\\[-.2pt]{\fontsize{5.8}{6.4}\selectfont \textcolor{BrickRed}{$ \blacktriangledown\,1.8\times $}}\end{tabular} & \hspace*{1pt}\begin{tabular}[t]{@{}c@{}}\rule{0pt}{12pt}2m\,44s\\[-.2pt]{\fontsize{5.8}{6.4}\selectfont \textcolor{BrickRed}{$ \blacktriangledown\,1.9\times $}}\\[-.2pt]2m\,26.2s\\[-.2pt]{\fontsize{5.8}{6.4}\selectfont \textcolor{BrickRed}{$ \blacktriangledown\,1.7\times $}}\end{tabular} & \textcolor{gray}{\ensuremath{\cdot}} & \hspace*{1pt}\begin{tabular}[t]{@{}c@{}}\rule{0pt}{12pt}5m\,1.5s\\[-.2pt]{\fontsize{5.8}{6.4}\selectfont \textcolor{BrickRed}{$ \blacktriangledown\,8.6\times $}}\\[-.2pt]1m\,24.2s\\[-.2pt]{\fontsize{5.8}{6.4}\selectfont \textcolor{BrickRed}{$ \blacktriangledown\,2.5\times $}}\end{tabular} & \hspace*{1pt}\begin{tabular}[t]{@{}c@{}}\rule{0pt}{12pt}4m\,0.6s\\[-.2pt]{\fontsize{5.8}{6.4}\selectfont \textcolor{BrickRed}{$ \blacktriangledown\,6.7\times $}}\\[-.2pt]1m\,24.3s\\[-.2pt]{\fontsize{5.8}{6.4}\selectfont \textcolor{BrickRed}{$ \blacktriangledown\,2.5\times $}}\end{tabular} & \hspace*{1pt}\begin{tabular}[t]{@{}c@{}}\rule{0pt}{12pt}8m\,31.7s\\[-.2pt]{\fontsize{5.8}{6.4}\selectfont \textcolor{BrickRed}{$ \blacktriangledown\,6.9\times $}}\\[-.2pt]2m\,34.1s\\[-.2pt]{\fontsize{5.8}{6.4}\selectfont \textcolor{BrickRed}{$ \blacktriangledown\,2.2\times $}}\end{tabular} & \textcolor{gray}{\ensuremath{\cdot}}\\[3pt]\hline
blastn & Native\textsuperscript{*} & \hspace*{1pt}\begin{tabular}[t]{@{}c@{}}\rule{0pt}{12pt}2m\,19.5s\\[-.2pt]{\fontsize{5.8}{6.4}\selectfont \textcolor{BrickRed}{$ \blacktriangledown\,7.1\times $}}\\[-.2pt]2m\,14.5s\\[-.2pt]{\fontsize{5.8}{6.4}\selectfont \textcolor{BrickRed}{$ \blacktriangledown\,7.2\times $}}\end{tabular} & \hspace*{1pt}\begin{tabular}[t]{@{}c@{}}\rule{0pt}{12pt}3m\,12.3s\\[-.2pt]{\fontsize{5.8}{6.4}\selectfont \textcolor{BrickRed}{$ \blacktriangledown\,9.9\times $}}\\[-.2pt]3m\,1.8s\\[-.2pt]{\fontsize{5.8}{6.4}\selectfont \textcolor{BrickRed}{$ \blacktriangledown\,10.4\times $}}\end{tabular} & \hspace*{1pt}\begin{tabular}[t]{@{}c@{}}\rule{0pt}{12pt}4m\,30.1s\\[-.2pt]{\fontsize{5.8}{6.4}\selectfont \textcolor{BrickRed}{$ \blacktriangledown\,3.7\times $}}\\[-.2pt]4m\,16s\\[-.2pt]{\fontsize{5.8}{6.4}\selectfont \textcolor{BrickRed}{$ \blacktriangledown\,3.8\times $}}\end{tabular} & \hspace*{1pt}\begin{tabular}[t]{@{}c@{}}\rule{0pt}{12pt}2m\,27.9s\\[-.2pt]{\fontsize{5.8}{6.4}\selectfont \textcolor{ForestGreen}{$ \blacktriangle\,3.0\times $}}\\[-.2pt]2m\,22.7s\\[-.2pt]{\fontsize{5.8}{6.4}\selectfont \textcolor{ForestGreen}{$ \blacktriangle\,2.9\times $}}\end{tabular} & \hspace*{1pt}\begin{tabular}[t]{@{}c@{}}\rule{0pt}{12pt}3m\,27.6s\\[-.2pt]{\fontsize{5.8}{6.4}\selectfont \textcolor{BrickRed}{$ \blacktriangledown\,22.3\times $}\textsuperscript{a}}\\[-.2pt]2m\,22.7s\\[-.2pt]{\fontsize{5.8}{6.4}\selectfont \textcolor{BrickRed}{$ \blacktriangledown\,17.0\times $}}\end{tabular} & \hspace*{1pt}\begin{tabular}[t]{@{}c@{}}\rule{0pt}{12pt}4m\,37.9s\\[-.2pt]{\fontsize{5.8}{6.4}\selectfont \textcolor{BrickRed}{$ \blacktriangledown\,31.8\times $}\textsuperscript{a}}\\[-.2pt]2m\,47s\\[-.2pt]{\fontsize{5.8}{6.4}\selectfont \textcolor{BrickRed}{$ \blacktriangledown\,23.8\times $}}\end{tabular} & \hspace*{1pt}\begin{tabular}[t]{@{}c@{}}\rule{0pt}{12pt}7m\,14.2s\\[-.2pt]{\fontsize{5.8}{6.4}\selectfont \textcolor{BrickRed}{$ \blacktriangledown\,27.5\times $}\textsuperscript{a}}\\[-.2pt]4m\,26.5s\\[-.2pt]{\fontsize{5.8}{6.4}\selectfont \textcolor{BrickRed}{$ \blacktriangledown\,20.9\times $}}\end{tabular} & \hspace*{1pt}\begin{tabular}[t]{@{}c@{}}\rule{0pt}{12pt}3m\,30.4s\\[-.2pt]{\fontsize{5.8}{6.4}\selectfont \textcolor{BrickRed}{$ \blacktriangledown\,11.4\times $}\textsuperscript{a}}\\[-.2pt]2m\,14.4s\\[-.2pt]{\fontsize{5.8}{6.4}\selectfont \textcolor{BrickRed}{$ \blacktriangledown\,7.9\times $}}\end{tabular} & \hspace*{1pt}\begin{tabular}[t]{@{}c@{}}\rule{0pt}{12pt}2m\,28.5s\\[-.2pt]{\fontsize{5.8}{6.4}\selectfont \textcolor{BrickRed}{$ \blacktriangledown\,3.5\times $}}\\[-.2pt]2m\,23.4s\\[-.2pt]{\fontsize{5.8}{6.4}\selectfont \textcolor{BrickRed}{$ \blacktriangledown\,3.5\times $}}\end{tabular} & \hspace*{1pt}\begin{tabular}[t]{@{}c@{}}\rule{0pt}{12pt}3m\,20.9s\\[-.2pt]{\fontsize{5.8}{6.4}\selectfont \textcolor{BrickRed}{$ \blacktriangledown\,4.7\times $}}\\[-.2pt]3m\,10.4s\\[-.2pt]{\fontsize{5.8}{6.4}\selectfont \textcolor{BrickRed}{$ \blacktriangledown\,4.6\times $}}\end{tabular} & \hspace*{1pt}\begin{tabular}[t]{@{}c@{}}\rule{0pt}{12pt}4m\,46.7s\\[-.2pt]{\fontsize{5.8}{6.4}\selectfont \textcolor{BrickRed}{$ \blacktriangledown\,3.3\times $}}\\[-.2pt]4m\,32.6s\\[-.2pt]{\fontsize{5.8}{6.4}\selectfont \textcolor{BrickRed}{$ \blacktriangledown\,3.3\times $}}\end{tabular} & \hspace*{1pt}\begin{tabular}[t]{@{}c@{}}\rule{0pt}{12pt}\colorbox{BenchBest}{\strut 2m\,38.7s}\\[-.2pt]{\fontsize{5.8}{6.4}\selectfont \textcolor{ForestGreen}{$ \blacktriangle\,1.8\times $}}\\[-.2pt]\colorbox{BenchBest}{\strut 2m\,33.5s}\\[-.2pt]{\fontsize{5.8}{6.4}\selectfont \textcolor{ForestGreen}{$ \blacktriangle\,1.9\times $}}\end{tabular} & \hspace*{1pt}\begin{tabular}[t]{@{}c@{}}\rule{0pt}{12pt}3m\,43.6s\\[-.2pt]{\fontsize{5.8}{6.4}\selectfont \textcolor{BrickRed}{$ \blacktriangledown\,6.4\times $}\textsuperscript{a}}\\[-.2pt]2m\,37.7s\\[-.2pt]{\fontsize{5.8}{6.4}\selectfont \textcolor{BrickRed}{$ \blacktriangledown\,4.6\times $}}\end{tabular} & \hspace*{1pt}\begin{tabular}[t]{@{}c@{}}\rule{0pt}{12pt}4m\,51.4s\\[-.2pt]{\fontsize{5.8}{6.4}\selectfont \textcolor{BrickRed}{$ \blacktriangledown\,8.1\times $}\textsuperscript{a}}\\[-.2pt]2m\,59.6s\\[-.2pt]{\fontsize{5.8}{6.4}\selectfont \textcolor{BrickRed}{$ \blacktriangledown\,5.3\times $}}\end{tabular} & \hspace*{1pt}\begin{tabular}[t]{@{}c@{}}\rule{0pt}{12pt}7m\,39.4s\\[-.2pt]{\fontsize{5.8}{6.4}\selectfont \textcolor{BrickRed}{$ \blacktriangledown\,6.2\times $}\textsuperscript{a}}\\[-.2pt]4m\,49.6s\\[-.2pt]{\fontsize{5.8}{6.4}\selectfont \textcolor{BrickRed}{$ \blacktriangledown\,4.1\times $}}\end{tabular} & \hspace*{1pt}\begin{tabular}[t]{@{}c@{}}\rule{0pt}{12pt}3m\,50.8s\\[-.2pt]{\fontsize{5.8}{6.4}\selectfont \textcolor{BrickRed}{$ \blacktriangledown\,3.8\times $}\textsuperscript{a}}\\[-.2pt]2m\,36.2s\\[-.2pt]{\fontsize{5.8}{6.4}\selectfont \textcolor{BrickRed}{$ \blacktriangledown\,2.7\times $}}\end{tabular}\\[3pt]\hline
winnowmap & Native\textsuperscript{*} & \hspace*{1pt}\begin{tabular}[t]{@{}c@{}}\rule{0pt}{12pt}28.4s\\[-.2pt]{\fontsize{5.8}{6.4}\selectfont \textcolor{BrickRed}{$ \blacktriangledown\,1.4\times $}}\\[-.2pt]26.9s\\[-.2pt]{\fontsize{5.8}{6.4}\selectfont \textcolor{BrickRed}{$ \blacktriangledown\,1.4\times $}}\end{tabular} & \hspace*{1pt}\begin{tabular}[t]{@{}c@{}}\rule{0pt}{12pt}27.6s\\[-.2pt]{\fontsize{5.8}{6.4}\selectfont \textcolor{BrickRed}{$ \blacktriangledown\,1.4\times $}}\\[-.2pt]24.8s\\[-.2pt]{\fontsize{5.8}{6.4}\selectfont \textcolor{BrickRed}{$ \blacktriangledown\,1.4\times $}}\end{tabular} & \hspace*{1pt}\begin{tabular}[t]{@{}c@{}}\rule{0pt}{12pt}1m\,2.8s\\[-.2pt]{\fontsize{5.8}{6.4}\selectfont \textcolor{ForestGreen}{$ \blacktriangle\,14\% $}}\\[-.2pt]55s\\[-.2pt]{\fontsize{5.8}{6.4}\selectfont \textcolor{ForestGreen}{$ \blacktriangle\,19\% $}}\end{tabular} & \hspace*{1pt}\begin{tabular}[t]{@{}c@{}}\rule{0pt}{12pt}8m\,35.8s\\[-.2pt]{\fontsize{5.8}{6.4}\selectfont \textcolor{BrickRed}{$ \blacktriangledown\,16\% $}}\\[-.2pt]8m\,4.4s\\[-.2pt]{\fontsize{5.8}{6.4}\selectfont \textcolor{BrickRed}{$ \blacktriangledown\,18\% $}}\end{tabular} & \hspace*{1pt}\begin{tabular}[t]{@{}c@{}}\rule{0pt}{12pt}10.8s\\[-.2pt]{\fontsize{5.8}{6.4}\selectfont \textcolor{BrickRed}{$ \blacktriangledown\,16\% $}}\\[-.2pt]9.4s\\[-.2pt]{\fontsize{5.8}{6.4}\selectfont \textcolor{BrickRed}{$ \blacktriangledown\,12\% $}}\end{tabular} & \hspace*{1pt}\begin{tabular}[t]{@{}c@{}}\rule{0pt}{12pt}10.9s\\[-.2pt]{\fontsize{5.8}{6.4}\selectfont \textcolor{BrickRed}{$ \blacktriangledown\,24\% $}}\\[-.2pt]8s\\[-.2pt]{\fontsize{5.8}{6.4}\selectfont \textcolor{BrickRed}{$ \blacktriangledown\,14\% $}}\end{tabular} & \hspace*{1pt}\begin{tabular}[t]{@{}c@{}}\rule{0pt}{12pt}20.1s\\[-.2pt]{\fontsize{5.8}{6.4}\selectfont \textcolor{BrickRed}{$ \blacktriangledown\,27\% $}}\\[-.2pt]14.5s\\[-.2pt]{\fontsize{5.8}{6.4}\selectfont \textcolor{BrickRed}{$ \blacktriangledown\,14\% $}}\end{tabular} & \hspace*{1pt}\begin{tabular}[t]{@{}c@{}}\rule{0pt}{12pt}21s\\[-.2pt]{\fontsize{5.8}{6.4}\selectfont \textcolor{BrickRed}{$ \blacktriangledown\,14\% $}}\\[-.2pt]17.1s\\[-.2pt]{\fontsize{5.8}{6.4}\selectfont \textcolor{gray}{$ \blacktriangleright\,0\% $}}\end{tabular} & \hspace*{1pt}\begin{tabular}[t]{@{}c@{}}\rule{0pt}{12pt}4m\,16.8s\\[-.2pt]{\fontsize{5.8}{6.4}\selectfont \textcolor{BrickRed}{$ \blacktriangledown\,6.1\times $}}\\[-.2pt]4m\,15.2s\\[-.2pt]{\fontsize{5.8}{6.4}\selectfont \textcolor{BrickRed}{$ \blacktriangledown\,6.2\times $}}\end{tabular} & \hspace*{1pt}\begin{tabular}[t]{@{}c@{}}\rule{0pt}{12pt}5m\,0.6s\\[-.2pt]{\fontsize{5.8}{6.4}\selectfont \textcolor{BrickRed}{$ \blacktriangledown\,7.0\times $}}\\[-.2pt]4m\,57.9s\\[-.2pt]{\fontsize{5.8}{6.4}\selectfont \textcolor{BrickRed}{$ \blacktriangledown\,7.2\times $}}\end{tabular} & \hspace*{1pt}\begin{tabular}[t]{@{}c@{}}\rule{0pt}{12pt}10m\,23.6s\\[-.2pt]{\fontsize{5.8}{6.4}\selectfont \textcolor{BrickRed}{$ \blacktriangledown\,7.2\times $}}\\[-.2pt]10m\,18s\\[-.2pt]{\fontsize{5.8}{6.4}\selectfont \textcolor{BrickRed}{$ \blacktriangledown\,7.4\times $}}\end{tabular} & \hspace*{1pt}\begin{tabular}[t]{@{}c@{}}\rule{0pt}{12pt}16m\,47.8s\\[-.2pt]{\fontsize{5.8}{6.4}\selectfont \textcolor{BrickRed}{$ \blacktriangledown\,3.5\times $}}\\[-.2pt]16m\,11s\\[-.2pt]{\fontsize{5.8}{6.4}\selectfont \textcolor{BrickRed}{$ \blacktriangledown\,3.4\times $}}\end{tabular} & \hspace*{1pt}\begin{tabular}[t]{@{}c@{}}\rule{0pt}{12pt}3m\,21.4s\\[-.2pt]{\fontsize{5.8}{6.4}\selectfont \textcolor{BrickRed}{$ \blacktriangledown\,5.8\times $}}\\[-.2pt]3m\,19.9s\\[-.2pt]{\fontsize{5.8}{6.4}\selectfont \textcolor{BrickRed}{$ \blacktriangledown\,5.9\times $}}\end{tabular} & \hspace*{1pt}\begin{tabular}[t]{@{}c@{}}\rule{0pt}{12pt}4m\,0.5s\\[-.2pt]{\fontsize{5.8}{6.4}\selectfont \textcolor{BrickRed}{$ \blacktriangledown\,6.7\times $}}\\[-.2pt]3m\,57.6s\\[-.2pt]{\fontsize{5.8}{6.4}\selectfont \textcolor{BrickRed}{$ \blacktriangledown\,7.0\times $}}\end{tabular} & \hspace*{1pt}\begin{tabular}[t]{@{}c@{}}\rule{0pt}{12pt}8m\,56.2s\\[-.2pt]{\fontsize{5.8}{6.4}\selectfont \textcolor{BrickRed}{$ \blacktriangledown\,7.2\times $}}\\[-.2pt]8m\,50.6s\\[-.2pt]{\fontsize{5.8}{6.4}\selectfont \textcolor{BrickRed}{$ \blacktriangledown\,7.4\times $}}\end{tabular} & \hspace*{1pt}\begin{tabular}[t]{@{}c@{}}\rule{0pt}{12pt}4m\,43.2s\\[-.2pt]{\fontsize{5.8}{6.4}\selectfont \textcolor{BrickRed}{$ \blacktriangledown\,4.7\times $}}\\[-.2pt]4m\,39.3s\\[-.2pt]{\fontsize{5.8}{6.4}\selectfont \textcolor{BrickRed}{$ \blacktriangledown\,4.8\times $}}\end{tabular}\\[3pt]\hline
lexicmap & Native\textsuperscript{*} & \hspace*{1pt}\begin{tabular}[t]{@{}c@{}}\rule{0pt}{12pt}38.8s\\[-.2pt]{\fontsize{5.8}{6.4}\selectfont \textcolor{BrickRed}{$ \blacktriangledown\,2.0\times $}}\\[-.2pt]38.6s\\[-.2pt]{\fontsize{5.8}{6.4}\selectfont \textcolor{BrickRed}{$ \blacktriangledown\,2.1\times $}}\end{tabular} & \hspace*{1pt}\begin{tabular}[t]{@{}c@{}}\rule{0pt}{12pt}54.3s\\[-.2pt]{\fontsize{5.8}{6.4}\selectfont \textcolor{BrickRed}{$ \blacktriangledown\,2.8\times $}}\\[-.2pt]54s\\[-.2pt]{\fontsize{5.8}{6.4}\selectfont \textcolor{BrickRed}{$ \blacktriangledown\,3.1\times $}}\end{tabular} & \hspace*{1pt}\begin{tabular}[t]{@{}c@{}}\rule{0pt}{12pt}1m\,29.8s\\[-.2pt]{\fontsize{5.8}{6.4}\selectfont \textcolor{BrickRed}{$ \blacktriangledown\,24\% $}}\\[-.2pt]1m\,29.4s\\[-.2pt]{\fontsize{5.8}{6.4}\selectfont \textcolor{BrickRed}{$ \blacktriangledown\,1.3\times $}}\end{tabular} & \hspace*{1pt}\begin{tabular}[t]{@{}c@{}}\rule{0pt}{12pt}49.9s\\[-.2pt]{\fontsize{5.8}{6.4}\selectfont \textcolor{ForestGreen}{$ \blacktriangle\,8.9\times $}}\\[-.2pt]49.7s\\[-.2pt]{\fontsize{5.8}{6.4}\selectfont \textcolor{ForestGreen}{$ \blacktriangle\,8.3\times $}}\end{tabular} & \hspace*{1pt}\begin{tabular}[t]{@{}c@{}}\rule{0pt}{12pt}37.8s\\[-.2pt]{\fontsize{5.8}{6.4}\selectfont \textcolor{BrickRed}{$ \blacktriangledown\,4.1\times $}\textsuperscript{b}}\\[-.2pt]37.3s\\[-.2pt]{\fontsize{5.8}{6.4}\selectfont \textcolor{BrickRed}{$ \blacktriangledown\,4.4\times $}}\end{tabular} & \hspace*{1pt}\begin{tabular}[t]{@{}c@{}}\rule{0pt}{12pt}53.5s\\[-.2pt]{\fontsize{5.8}{6.4}\selectfont \textcolor{BrickRed}{$ \blacktriangledown\,6.1\times $}\textsuperscript{b}}\\[-.2pt]52.8s\\[-.2pt]{\fontsize{5.8}{6.4}\selectfont \textcolor{BrickRed}{$ \blacktriangledown\,7.5\times $}}\end{tabular} & \hspace*{1pt}\begin{tabular}[t]{@{}c@{}}\rule{0pt}{12pt}1m\,18.6s\\[-.2pt]{\fontsize{5.8}{6.4}\selectfont \textcolor{BrickRed}{$ \blacktriangledown\,5.0\times $}\textsuperscript{b}}\\[-.2pt]1m\,17.5s\\[-.2pt]{\fontsize{5.8}{6.4}\selectfont \textcolor{BrickRed}{$ \blacktriangledown\,6.1\times $}}\end{tabular} & \hspace*{1pt}\begin{tabular}[t]{@{}c@{}}\rule{0pt}{12pt}49.7s\\[-.2pt]{\fontsize{5.8}{6.4}\selectfont \textcolor{BrickRed}{$ \blacktriangledown\,2.7\times $}\textsuperscript{b}}\\[-.2pt]49.3s\\[-.2pt]{\fontsize{5.8}{6.4}\selectfont \textcolor{BrickRed}{$ \blacktriangledown\,2.9\times $}}\end{tabular} & \hspace*{1pt}\begin{tabular}[t]{@{}c@{}}\rule{0pt}{12pt}1m\,41.8s\\[-.2pt]{\fontsize{5.8}{6.4}\selectfont \textcolor{BrickRed}{$ \blacktriangledown\,2.4\times $}}\\[-.2pt]1m\,41.7s\\[-.2pt]{\fontsize{5.8}{6.4}\selectfont \textcolor{BrickRed}{$ \blacktriangledown\,2.5\times $}}\end{tabular} & \hspace*{1pt}\begin{tabular}[t]{@{}c@{}}\rule{0pt}{12pt}2m\,35.6s\\[-.2pt]{\fontsize{5.8}{6.4}\selectfont \textcolor{BrickRed}{$ \blacktriangledown\,3.6\times $}}\\[-.2pt]2m\,35.4s\\[-.2pt]{\fontsize{5.8}{6.4}\selectfont \textcolor{BrickRed}{$ \blacktriangledown\,3.8\times $}}\end{tabular} & \hspace*{1pt}\begin{tabular}[t]{@{}c@{}}\rule{0pt}{12pt}4m\,6.7s\\[-.2pt]{\fontsize{5.8}{6.4}\selectfont \textcolor{BrickRed}{$ \blacktriangledown\,2.8\times $}}\\[-.2pt]4m\,6.3s\\[-.2pt]{\fontsize{5.8}{6.4}\selectfont \textcolor{BrickRed}{$ \blacktriangledown\,2.9\times $}}\end{tabular} & \hspace*{1pt}\begin{tabular}[t]{@{}c@{}}\rule{0pt}{12pt}3m\,25.5s\\[-.2pt]{\fontsize{5.8}{6.4}\selectfont \textcolor{ForestGreen}{$ \blacktriangle\,30\% $}}\\[-.2pt]3m\,25.3s\\[-.2pt]{\fontsize{5.8}{6.4}\selectfont \textcolor{ForestGreen}{$ \blacktriangle\,28\% $}}\end{tabular} & \hspace*{1pt}\begin{tabular}[t]{@{}c@{}}\rule{0pt}{12pt}1m\,43.4s\\[-.2pt]{\fontsize{5.8}{6.4}\selectfont \textcolor{BrickRed}{$ \blacktriangledown\,3.0\times $}\textsuperscript{b}}\\[-.2pt]1m\,42.6s\\[-.2pt]{\fontsize{5.8}{6.4}\selectfont \textcolor{BrickRed}{$ \blacktriangledown\,3.0\times $}}\end{tabular} & \hspace*{1pt}\begin{tabular}[t]{@{}c@{}}\rule{0pt}{12pt}2m\,37.7s\\[-.2pt]{\fontsize{5.8}{6.4}\selectfont \textcolor{BrickRed}{$ \blacktriangledown\,4.4\times $}\textsuperscript{b}}\\[-.2pt]2m\,37.1s\\[-.2pt]{\fontsize{5.8}{6.4}\selectfont \textcolor{BrickRed}{$ \blacktriangledown\,4.6\times $}}\end{tabular} & \hspace*{1pt}\begin{tabular}[t]{@{}c@{}}\rule{0pt}{12pt}3m\,45.3s\\[-.2pt]{\fontsize{5.8}{6.4}\selectfont \textcolor{BrickRed}{$ \blacktriangledown\,3.0\times $}\textsuperscript{b}}\\[-.2pt]3m\,43.9s\\[-.2pt]{\fontsize{5.8}{6.4}\selectfont \textcolor{BrickRed}{$ \blacktriangledown\,3.1\times $}}\end{tabular} & \hspace*{1pt}\begin{tabular}[t]{@{}c@{}}\rule{0pt}{12pt}3m\,16.3s\\[-.2pt]{\fontsize{5.8}{6.4}\selectfont \textcolor{BrickRed}{$ \blacktriangledown\,3.3\times $}\textsuperscript{b}}\\[-.2pt]3m\,15.8s\\[-.2pt]{\fontsize{5.8}{6.4}\selectfont \textcolor{BrickRed}{$ \blacktriangledown\,3.3\times $}}\end{tabular}\\[3pt]\hline
\end{tabular}\par}
\smallskip{\footnotesize\noindent Each cell: A, comparison, B, comparison. Batches of 1,024 queries. Five repeats unless marked s (FAST: one sample). Blue = one lowest A and one lowest B down each machine/workload column.\par
\noindent Light gray = minimap2 baseline. Pale yellow = preferred MPF / Bounded / SIMD-DP / Block configuration, not a claim of winning every column. Exact ties select the first row.\par
\noindent Ratios use minimap2 as reference, except minimap2 itself uses the fixed BioseqDB row. $\cdot$ = no approved timing. n/a = no measured counterpart.\par
\noindent Both minimap2 and BioseqDB use the matching non-external remote measurements.\par
\noindent\textsuperscript{*}External A includes process startup and output. B removes different native stages and is not an equivalent-work comparison. BLAST and LexicMap have no $w$ or MAPQ equivalent.\par
\noindent\textsuperscript{a}BLAST DNA5 uses an N score of -3, compared with -1 in the baseline.\par
\noindent\textsuperscript{b}LexicMap stores reference N as A. Queries retain their bases, and edit-script checks use that native reference encoding.\par
\noindent\textsuperscript{g}Retained C7G minimap2 measurement comes from the earlier campaign with documented paging pressure.\par
}
\end{landscape}
\restoregeometry
